\documentclass{book}%

%\usepackage{fontspec}[greek]
\usepackage{fourier-otf} % pour les arcs (double brève qui compte comme une longue)
\usepackage{arcs} % pour les arcs (double brève qui compte comme une longue)
\usepackage{graphicx} % pour l'image de la coupe d'Onésime
\usepackage{fontspec}%
\usepackage{xunicode}%
%\usepackage{yfonts}%
\usepackage{hyperref}% pour le courriel
\usepackage{amsmath}%
\usepackage{cancel}% pour \inRas
\usepackage{quoting}%
%\usepackage{csquotes} % Pour remplacer les guillemets anglais par des français
%\MakeOuterQuote{"}
\usepackage{ulem}% pour souligner en pointillé
\usepackage{comment}%
\usepackage{xcolor}%
%scale = {hscale, vscale}
\usepackage[scale={0.75,0.8}]{geometry}%
%\setlength{\hoffset}{-18pt}        
\setlength{\oddsidemargin}{20pt} % Marge gauche sur pages impaires ATTENTION écrit 2 fois
%\setlength{\evensidemargin}{9pt} % Marge gauche sur pages paires
%\setlength{\marginparwidth}{54pt} % Largeur de note dans la marge
%\setlength{\textwidth}{481pt} % Largeur de la zone de texte (17cm)
%\setlength{\voffset}{-18pt} % Bon pour DOS
%\setlength{\marginparsep}{7pt} % Séparation de la marge
%\setlength{\topmargin}{0pt} % Pas de marge en haut
%\setlength{\headheight}{13pt} % Haut de page
%\setlength{\headsep}{10pt} % Entre le haut de page et le texte
%\setlength{\footskip}{27pt} % Bas de page + séparation
%\setlength{\textheight}{708pt} % Hauteur de la zone de texte (25cm)


%\newfontfamily\policeGr{Liberation Serif}[
%  Script=Greek,
%  NFSSFamily=liberationserif,
%
%]
%\newfontfamily\policeGr{Liberation Serif}%

%\newfontfamily{\policeFr}{Garamond Premier Pro}%
%\newfontfamily{\policeLat}{Book Antiqua}%
%\newfontfamily{\policeAn}{Comic Sans MS}%
%\newfontfamily{\policeAll}{Fraktur}%
%\newfontfamily{\policeIt}{BodoniFLF}%
%\newfontfamily{\policeEsp}{Great Moravia}%
% ...
%\setmainfont{Liberation Serif}%
%\setmainfont{Liberation Serif}%
%\newfontfamily{\policeFr}{Garamond Premier Pro}%
%\newfontfamily{\policeFr}{EB Garamond}%
\newfontfamily{\policeFr}{GFS Bodoni}%
%\newfontfamily{\policeLat}{Book Antiqua}%
\newfontfamily{\policeLat}{GFS Neohellenic}%
%\newfontfamily{\policeAn}{Comic Sans MS}%
\newfontfamily{\policeAn}{Courier 10 Pitch}%
\newfontfamily{\policeEsp}{Abyssinica SIL}%
%\newfontfamily{\policeAll}{Fraktur}%
%\newfontfamily{\policeAll}{fraktur-bt.ttf}%
\newfontfamily{\policeAll}{Bertholdr Mainzer Fraktur}%
%\newfontfamily{\policeIt}{BodoniFLF}%
\newfontfamily{\policeIt}{GFS Artemisia}%
%\newfontfamily{\policeEsp}{Great Moravia}%
%\newfontfamily\policeGr{GFSBodoni.otf}[BoldFont = GFSBodoniBold.otf, ItalicFont = GFSBodoniIt.otf, BoldItalicFont = GFSBodoniBoldIt.otf]%
%\setmainfont{GFSBodoni.otf}[BoldFont = GFSBodoniBold.otf, ItalicFont = GFSBodoniIt.otf, BoldItalicFont = GFSBodoniBoldIt.otf]
\newfontfamily\policeGr{GFS Bodoni}%
\setmainfont{GFS Bodoni}
%\setdefaultlanguage[variant=ancient]{greek}%
\usepackage{polyglossia}%
\setdefaultlanguage{french}%
%\setotherlanguage[variant=ancient]{greek}%
%\setmainlanguage[variant=ancient]{greek} 
%\setotherlanguage[variant=ancient]{greek} 
%\setmainfont{Libertinus Serif} 
%\newfontfamily{\greekfont}{Liberation Serif}[
%  Script=Greek,
%  NFSSFamily=liberationserif,
%]

\usepackage[maxlevel=3]{csquotes}%
%\MakeOuterQuote{"}% pour les guillemets français
\usepackage{verse}%
\usepackage[parapparatus]{reledmac}%
\usepackage[parledgroup, nomaxlines]{reledpar}% 
%\usepackage[parledgroup, nomaxlines, nopbinverse]{reledpar}% nopbinverse empêche de casser un groupe
%\usepackage[parledgroup]{reledpar}% nomaxlines permettait d'avoir un nombre de lignes différents sur des pages qui se font face.
%\parledgrouptrue
\setgoalfraction{0.9}% sept 2022

% Pour que les commentaires apparaissent sur les deux pages
\AtBeginDocument{%
\Xmaxhnotes[A,B]{0.38\textheight}%
\Xmaxhnotes[C]{0.19\textheight}%
%\Xmaxhnotes{0.9\textheight}%
%\maxhnotesX{0.45\textheight}%
\setRlineflag{}%
%\newcommand{\decoupeVers}[1]{\linenumannotation{#1}}%
\newcommand{\decoupeVers}[1]{}%
%\setlength{\oddsidemargin}{100pt} % Marge gauche sur pages impaires ATTENTION écrit 2 fois
% La Largeur du commentaire
%\Xwidth{\columnwidth+\marginparsep+\ledrsnotewidth}
}%

\newcommand{\italique}[1]{\textit{#1}}%
\newcommand{\citComm}[1]{\begin{quoting}[vskip=0pt] #1 \end{quoting}}% vskip = \topsep + \partopsep par défaut
\newcommand{\passeLigne}[0]{\skipnumbering}%
\newcommand{\ligne}[1]{\setline{#1}\hidenumbering}%
\newcommand{\espace}[1]{\hspace{#1}}%
\newcommand{\cesure}[0]{ [...] }%
\newcommand{\nieme}[1]{#1\italique{e}}%
\newcommand{\lettreBreve}[1]{$\breve{\mbox{#1}}$}
\newcommand{\synizese}[1]{$\widearc{\mbox{#1}}$}
\newcommand{\inRas}[2]{$\overset{\text{#2}}{\text{\sout{#1}}}$}
\newcommand{\auDessus}[2]{$\overset{\text{#2}}{\text{#1}}$}
%\newcommand{\etc}[0]{κτλ}%
%\newcommand{\cit}[0]{\espace{1cm}} % indent. pour une citation
\newcommand{\expo}[1]{\textsuperscript{#1}}%
\newcommand{\alinea}[0]{ \ \ }%
\newcommand{\Syn}[0]{ Συναγωγή }%
\newcommand{\SynB}[0]{ Συναγωγή\expo{B} }% la deuxième version pour la lettre A
\newcommand{\ClemAlex}{ Clemens Alexandrinus }
\newcommand{\Eustathe}{ Eustathius }
\newcommand{\AelDion}{ Aelius Dionysius }
\newcommand{\EtymMagnum}{ Etymologicum magnum }
\newcommand{\EtymGen}{ Etymologicum genuinum }
\newcommand{\EtymSym}{ Etymologicum Symeonis }
\newcommand{\Apostol}{ Michael Apostolius }
\newcommand{\Macaire}{ Macarius } 
\newcommand{\Phryn}{ Phrynichus \italique{Praep. soph.} }
\newcommand{\Stobee}{ Joannes Stobaeus }
\newcommand{\EpimAlphab}{ Epimerismi alphabetici in Homerum (ed. Cramer) }
\newcommand{\GregCor}{ Gregorius Corinthius }
\newcommand{\Antiatt}{ Antiatticista (ed. Bekker) }
\newcommand{\Harp}{ Harpocration }% Autre édition de référence : Dindorf
\newcommand{\Rpc}[0]{R\expo{pc}}%
\newcommand{\Rac}[0]{R\expo{ac}}%
\newcommand{\Rdeux}[0]{R\expo{2}}%
\newcommand{\Mon}[0]{ M }%
\newcommand{\PsiUn}[0]{PSI-1194}
\newcommand{\PsiDeux}[0]{POxy-3839}
\newcommand{\PsiTrois}[0]{POxy-3840}
\newcommand{\PsiQuatre}[0]{POxy-4935}
\newcommand{\PsiCinq}[0]{POxy-1176}%P50

\newcommand{\francais}[1]{{\policeFr #1}}%
\newcommand{\latin}[1]{{\policeLat #1}}%
\newcommand{\allemand}[1]{{\policeAll #1}}%
%\newcommand{\allemand}[1]{{\textfrak{#1}}}%
%\newcommand{\grec}[1]{{\policeGr #1}}%
\newcommand{\grec}[1]{#1}%
\newcommand{\anglais}[1]{{\policeAn \scriptsize{#1}}}%
\newcommand{\italien}[1]{{\policeIt #1}}%
\newcommand{\espagnol}[1]{{\policeEsp #1}}%

\newcommand{\Nauck}[0]{N\expo{2}}%
\newcommand*{\auteur}[1]{\textsc{#1}*}%
\newcommand{\cf}[0]{\italique{cf.}\xspace}%
\newcommand{\chap}[1]{\^{#1}}%
%\newcommand{\chap}[1]{$\hat{#1}$}


%\marginparwidth = 1.2\marginparwidth
%\ledrsnotewidth = 1.2\marginparwidth
\newlength{\larg}%
\newcommand{\decalage}[1]{%
\settowidth{\larg}{#1}%
\espace{\larg}%
}
\newcommand{\decalageN}[1]{%
\settowidth{\larg}{#1}%
\espace{-\larg}%
}
\newcommand{\apparat}[2]{\ledrightnote{\espace{-1.2\marginparwidth}\begin{minipage}{1.4\marginparwidth}\footnotesize{\textbf{#1} #2}\end{minipage}}}%
%\newcommand{\apparat}[2]{{\espace{-0.4\marginparwidth}\ledrightnote{\footnotesize{\textbf{#1} #2}}\espace{-0.4\marginparwidth}}}%
%\newcommand{\apparat}[2]{\ledrightnote{\footnotesize{\textbf{#1} #2}}}%
%\newcommand{\lecon}[1]{\dotuline{#1}} % la leçon retenue de l'apparat critiqe
\newcommand{\lecon}[1]{\textbf{#1}}
%\newcommand{\douteux}[1]{\dotuline{#1}} % pour les lettres douteuses : un point en-dessous
\newcommand{\douteux}[1]{$\underset{\cdot}{\text{#1}}$}
\newcommand{\difference}[1]{\dotuline{#1}}
\rightnoteupfalse 

%\setsidenotesep{ sep }%
\newcommand{\espApp}[0]{ || }%\espace{10pt}%
\newcommand{\espTest}[0]{ || }%\espace{10pt}%
\newcommand{\ligneApp}[1]{ \ \ \textbf{#1} }%
%% 3 arguments : les numéros de ligne, le lemme et la note elle-même
\newcommand{\commentaire}[3]{\edtext{}{\Bfootnote[nonum, nosep]{{\espace{-1cm}\textbf{#1~~#2} #3}}}}%
%\renewcommand{\footfudgefiddle}{68}%Note that you must call it before 
% \Xarrangement[C]{paragraph}% entrâine trop d'erreurs, cf. documentation
\Xgroupbyline[A,C]
\newcommand{\debutNote}[1]{\edtext{}{\Cfootnote[nonum, nosep]{\espace{-15pt}\textbf{\dotuline{\textbf{#1}}}}}}%
\newcommand{\note}[2]{\edtext{}{\Cfootnote[nonum, nosep]{\espace{-15pt}\textbf{#1} #2}}}%
\newcommand{\testimonia}[2]{\edtext{}{\Afootnote[nonum, nosep]{\espace{-15pt}\textbf{#1} #2}}}%
\Xonlyside[A]{R}% imprime les testimonia page de droite
\newcommand{\marqueur}[1]{\textsuperscript{#1}}%
%\newcommand{\cat}[0]{$""_"^"$}%
\newcommand{\cat}[0]{\texorpdfstring{\textsubscript{\textasciicircum}}{cat}}%
%\fnpos{critical-familiar}%
%\mpfnpos{familiar-critical}%

\newcommand{\eff}[1]{%
{\color{orange}#1}%
}%

\newlength{\textlarg}%
\newcommand{\typeVers}[1]{%
\settowidth{\textlarg}{\footnotesize{#1}}%
\hspace{-0.8cm}{\footnotesize{#1}}\hspace{-\textlarg}\hspace{0.8cm}}%
\newlength{\largtext}%
\newcommand{\strophe}[1]{%
\settowidth{\largtext}{~~\footnotesize{\francais{#1}}}%
~~\footnotesize{\francais{#1}}%\hspace{-\largtext}%
}%
%\newcommand{\strophe}[1]{%
%\settowidth{\largtext}{\footnotesize{\francais{#1}}}%
%	\ledleftnote{%
		%\hspace{-0.5cm}\footnotesize{\francais{#1}}\hspace{\largtext}\hspace{0.5cm%}%
%		\hspace{-0.5cm}\footnotesize{\francais{#1}}\hspace{\largtext}\hspace{0.5cm}%
%	}%
%}%
%\newcommand{\personnageGrec}[1]{\ledleftnote{\textbf{\large{#1}}\espace{-1.5cm}}}%
\newlength{\textelargeur}%
\newcommand{\personnageGrec}[1]{\ledleftnote{%
%\settowidth{\textelargeur}{\large{\textbf{#1}}}%
\settowidth{\textelargeur}{\large{\textbf{Μνησίλοχος }}}%
%\espace{-\textelargeur}
\begin{minipage}{\textelargeur}\begin{flushright}\large{\textbf{#1}}\end{flushright}\end{minipage}}}%
\newcommand{\Strophe}[1]{\ledleftnote{%
\settowidth{\textelargeur}{\footnotesize{\francais{#1}}}%
%\espace{-\textelargeur}
\begin{minipage}{\textelargeur}\footnotesize{\francais{#1}}\end{minipage}}}%
%\newcommand{\personnageFr}[1]{\ledleftnote{\textbf{\large{\francais{#1}}}\espace{-1.5cm}}}%
\newlength{\texto}%
\newcommand{\personnageFr}[1]{\ledleftnote{%
%\settowidth{\texto}{\large{\textbf{#1}}}%
\settowidth{\texto}{\large{\textbf{Marchande d}}}%
%\espace{-\textelargeur}
\begin{minipage}{\texto}\begin{flushright}\large{\textbf{#1}}\end{flushright}\end{minipage}}}%
%% Les différents types de vers
\newcommand{\trimIamb}[1]{%
  \hangindent=50pt%
  \espace{30pt} % la valeur précédente -20
  #1%
}%
\newcommand{\trimIambFr}[1]{%
  \hangindent=50pt%
  \espace{30pt} % la valeur précédente -20
  \francais{#1%
}}%
\newcommand{\trimIambCat}[1]{%
  \hangindent=53pt%
  \espace{33pt} % la valeur précédente -20
  #1%
}%
\newcommand{\trimIambCatFr}[1]{%
  \hangindent=53pt%
  \espace{33pt} % la valeur précédente -20
  \francais{#1%
}}%
\newcommand{\tetramIamb}[1]{%
  \hangindent=40pt%
  \espace{20pt} % la valeur précédente -20
  #1%
}%
\newcommand{\tetramIambFr}[1]{%
  \hangindent=40pt%
  \espace{20pt} % la valeur précédente -20
  \francais{#1%
}}%
\newcommand{\tetramTroch}[1]{%
  \hangindent=40pt%
  \espace{20pt} % la valeur précédente -20
  #1%
}%
\newcommand{\tetramTrochFr}[1]{%
  \hangindent=40pt%
  \espace{20pt} % la valeur précédente -20
  \francais{#1%
}}%
\newcommand{\tetramIambCat}[1]{%
  \hangindent=42pt%
  \espace{22pt} % la valeur précédente -20
  #1%
}%
\newcommand{\tetramIambCatFr}[1]{%
  \hangindent=42pt%
  \espace{22pt} % la valeur précédente -20
  \francais{#1%
}}%
\newcommand{\dimAnap}[1]{%
  \hangindent=70pt%
  \espace{50pt} % la valeur précédente -20
  #1%
}%
\newcommand{\dimAnapFr}[1]{%
  \hangindent=70pt%
  \espace{50pt} % la valeur précédente -20
  \francais{#1%
}}%
\newcommand{\anap}[1]{%
  \hangindent=70pt%
  \espace{50pt} % la valeur précédente -20
  #1%
}%
\newcommand{\anapFr}[1]{%
  \hangindent=70pt%
  \espace{50pt} % la valeur précédente -20
  \francais{#1%
}}%
% 2 dochmiaques
\newcommand{\diDoch}[1]{%
  \hangindent=60pt%
  \espace{40pt} % la valeur précédente -20
  #1%
}%
\newcommand{\diDochFr}[1]{%
  \hangindent=60pt%
  \espace{40pt} % la valeur précédente -20
  \francais{#1%
}}%
\newcommand{\tetramAnap}[1]{%
  \hangindent=10pt%
  \espace{-10pt} % la valeur précédente -20
  #1%
}%
\newcommand{\tetramAnapFr}[1]{%
  \hangindent=10pt%
  \espace{-10pt} % la valeur précédente -20
  \francais{#1%
}}%
\newcommand{\chantAgath}[1]{%
  \hangindent=90pt%
  \espace{70pt} % la valeur précédente -20
  #1%
}%
\newcommand{\chantAgathFr}[1]{%
  \hangindent=90pt%
  \espace{70pt} % la valeur précédente -20
  \francais{#1%
}}%
%% Pour les "vers" 295 à 311
\newcommand{\prose}[1]{%
  \hangindent=30pt%
  \espace{10pt} % la valeur précédente -20
  #1%
}%
\newcommand{\proseFr}[1]{%
  \hangindent=30pt%
  \espace{10pt} % la valeur précédente -20
  \francais{#1%
}}%
%% pour les vers courts de structure changeante dans un choeur
\newcommand{\chantChoeur}[1]{%
%{\centering
%	#1\par
%}%
\hangindent=50pt%
\espace{30pt} % la valeur précédente -20
#1%
}%
\newcommand{\chantChoeurFr}[1]{%
%{\centering
%	\francais{#1}\par
%}%
\hangindent=50pt%
\espace{30pt} % la valeur précédente -20
\francais{#1%
}}%
%% pour les vers 1015 à 1055
\newcommand{\paratragiques}[1]{%
%{\centering
%	#1\par
%}%
\hangindent=50pt%
\espace{30pt} % la valeur précédente -20
#1%
}%
\newcommand{\paratragiquesFr}[1]{%
%{\centering
%	\francais{#1}\par
%}%
\hangindent=30pt%
\espace{10pt} % la valeur précédente -20
\francais{#1%
}}%
% pour indiquer que les notes marginales s'alignent au niveau de la première ligne
%\leftnoteupfalse
%\rightnoteupfalse



\newenvironment{groupe}{%
\pstart%
\ledgroup%
}{%
\pend %
}%

\newcommand{\LSJ}[0]{LSJ}%

%% pour voir tous les numéros de ligne (version brouillonne)
\linenumincrement{1}%
\firstlinenum{1}%
%% Lignes à imprimer
\gdef\linenumberlist{2,6,10,15,20,24,30,36,40,45,50,55,60,65,69,74,80,87,91,94,102,105,110,115,120,125,131,135,140,145,150,155,160,165,170,174,180,185,190,196,201,205,215,220,225,229,234,239,245,250,261,265,270,276,280,285,290,295,300,305,310,315,320,325,330,335,340,345,350,355,360,365,370,375,379,385,390,395,400,405,410,415,420,425,430,435,440,445,450,455,460,465,470,475,480,485,490,495,499,501,505,510,515,521,526,530,535,541,545,550,556,561,565,572,575,580,585,591,596,600,609,614,620,624,629,635,640,650,656,660,665,670,675,680,685,690,695,700,708,713,720,725,729,735,741,745,750,755,761,766,770,775,780,786,790,795,800,805,810,815,820,825,831,835,840,845,850,856,860,865,870,875,880,884,890,894,900,905,915,921,930,935,940,946,950,955,960,965,970,975,980,985,990,996,1000,1010,1016,1020,1025,1030,1035,1040,1045,1050,1055,1060,1066,1071,1076,1100,1106,1116,1120,1124,1130,1141,1147,1150,1156,1161,1166,1171,1175,1180,1185,1190,1195,1200,1205,1211,1215,1221,1224,1230}

\begin{document} 
%\francais{Une phrase en français. (Garamond Premier Pro)}\\
%\latin{Une phrase en latin. (Book Antiqua)}\\
%\grec{Une phrase en grec. (Liberation serif)}\\
%\allemand{Une phrase en allemand. (Fraktur)}\\
%\italien{Une phrase en italien. (Bodoni)}\\
%\anglais{Une phrase en anglais. (Comic Sans MS)}\\
%\espagnol{Une phrase en espagnol. (Great Moravia)}\\
\title{Les \italique{Thesmophories} d'Aristophane.\\
\allemand{Urtext}\\
avec commentaire,\\
traduction,\\
apparats (critique \& témoins).}
\author{Louis Hauseux \\ courriel : \href{mailto:louis.hauseux@gmail.com}{louis.hauseux@gmail.com} }
\date{IX \oldstylenums{2022}}

\maketitle

\begin{pages}%

%\setlength{\Lcolwidth}{0.75\textwidth}%
%\setlength{\ledlsnotesep}{2\linenumsep}%
%\setlength{\ledrsnotesep}{2\linenumsep}%
%\setlength{\ledlsnotewidth }{2\marginparwidth}%
\begin{Leftside}%
\beginnumbering
%% v. 1-4
\begin{groupe}%
\personnageGrec{Μνησίλοχος}%
\typeVers{trim. iamb.}\trimIamb{ὦ Ζεῦ Χελιδὼν ἆρά ποτε φανήσεται;\apparat{1}{χελιδὼν R : \lecon{Χελιδὼν} Bibi ap. Clements \espApp ἄρα πότε R : \lecon{ἆρά ποτε} Biset \ligneApp{2} ἄνθρωπος R : \lecon{ἅνθρωπος} Boissonade \ligneApp{3} τι R : \lecon{τε} schol. (δυνατόν), Suid. οι 150}\\
ἀπολεῖ μ᾽ ἀλοῶν ἅνθρωπος ἐξ ἑωθινοῦ.\testimonia{2}{(ἀλοῶν) Suidas α 1330 ; Photius α 1029 $=$ \SynB α 984}%
\commentaire{1-10}{ὦ Ζεῦ ... οὐχ ἅ γ᾽ ἂν ἀκούειν δέῃ.}{Nous sommes au tout début de la pièce, il fait encore nuit, un personnage rentre d'un pas pressé ; il est immédiatement reconnaissable à son masque bien connu des spectateurs : c'est Euripide !\marqueur{1} Il est suivi d'un vieillard qui traîne la patte (on apprendra plus loin qu'il est un parent d'Euripide, la tradition l'identifie à Mnésiloque, le père de la première femme d'Euripide, Chœrine\marqueur{2}). Plainte du vieillard. Réplique fracassante d'Euripide, et bien selon ses façons : au lieu de répondre simplement à la question, il préfère proférer une de ces sentences générales (les fameuses γνῶμαι d'Euripide) fort goûtées des spectateurs : ἀλλ᾽ οὐκ ἀκούειν δεῖ σε πάνθ᾽ ὅσ᾽ αὐτίκα ὄψει παρεστώς, «~il n'est pas besoin d'entendre tout ce qui va bientôt se présenter à la vue~».
Hélas pour lui, Mnésiloque ne fait pas partie de ces Athéniens prisant les euripidéens aphorismes ; il y est même tout-à-fait sourd : il n'arrive pas à suivre Euripide, sur scène comme dans ses sentences aux constructions tortueuses\marqueur{3} ; la fatigue aidant, Mnésiloque s'est arrêté à la première partie de la phrase : ἀλλ᾽ οὐκ ἀκούειν δεῖ σε, «~il ne faut pas que tu entendes~». Il n'en croit pas ses oreilles et demande à Euripide de répéter. \\
C'est là que nous arrivons au vers 7b. Nouvelle amphibologie : Euripide ne fait que reformuler ce qu'il a déjà dit : οὐχ ἅ γ ἂν μέλῃς ὁρᾶν, « pas certes ce que tu es sur le point de voir ». Les commentateurs s'étaient jusqu'à présent contenté de cet unique sens, rendant incompréhensible la réplique de Mnésiloque : οὐδ ἆρ ὁρᾶν δεῖ μ , « Je ne dois pas non plus voir ? » (?! sic) Sans transition !.. on ne comprend vraiment pas pourquoi Mnésiloque passe du coq à l'âne, ou plutôt d'une défense d'entendre à une défense de voir. 
En fait, il ne fait que déformer encore une fois les propos d'Euripide : ici, dans `οὐχ ἅ γ ἂν μέλῃς ὁρᾶν', il commet une erreur d'interprétation : il entend `ἄγαν' à la place de `ἅ γ ̓ ἂν'\marqueur{4}. 
Mnésiloque comprend donc : οὐκ ἄγαν μέλῃς ὁρᾶν, « Tu n'as pas spécialement à voir ! »... d'où sa réaction étonnée !}%
\debutNote{1-10}%
\note{1}{À quoi pouvait ressembler ce masque d'Euripide ? On peut supposer qu'il présentait une barbe très touffue (cf. le vers 190) ainsi que de multiples taches de rousseur au niveau des yeux : Ἐλέγετο δὲ καὶ βαθὺν πώγωνα θρέψαι καὶ ἐπὶ τῆς ὄψεως φακοὺς ἐσχηκέναι. (Tiré d'une des \italique{Vies} d'Euripide, rassemblées par Richard Kannicht pour son édition des fragments d'Euripide, cinquième volume des \italique{Tragicorum Graecorum Fragmenta} (2004). Ce détail se trouve à la page 47.)}%
\note{2}{\italique{Confer} l'introduction sur la vie d'Euripide et de Mnésiloque. %\eff{Il n'y a aucune raison de mettre en doute l'existence historique de Mnésiloque ainsi que le lien de parenté qu'il entretenait avec lui. S'il est vrai que les indices internes au texte d'Aristophane sont ténus L'arbre généalogique d'Euripide chez Kannicht concorde avec l'usage d'appeler son premier fils du nom du grand-père paternel et second de celui maternel. Cf. Willems Démosthène ?}
}%
\note{3}{Le scholiaste en a fait la remarque : ὁ μὲν τραγικώτερον καὶ ὑψηλότερον φράζει, ὁ δὲ ταπεινότερον ἀκούει, « l'un s'exprime dans un style plus tragique et élevé, quand l'autre comprend de façon plus basse ».}%
\note{4}{Aristophane ayant mis à dessein ces particules dans un ordre inhabituel : le cas se rencontre une seule autre fois dans les onze comédies d'Aristophane... et c'est encore dans la bouche d'Euripide (celui des \italique{Grenouilles}, cette fois, vers 960) : ἐξ \textbf{ὧν γ᾽ ἂν} ἐξηλεγχόμην. Le véritable Euripide s'exprimait ainsi : (\italique{Oreste}, vers 81) : Ἑλένη, τί σοι λέγοιμ᾽ \textbf{ἂν ἅ γε} παροῦσ᾽ ὁρᾷς;}%
\\
παρὰ σοῦ πυθέσθαι ποῖ μ᾽ ἄγεις ωὖριπίδη;%
\commentaire{1}{Χελιδὼν.}{\italique{Hirondelle} avec un `H' majuscule ! Comme l'a très justement vu Ashley Clements\marqueur{1}, Χελιδών désigne une étoile (qu'évoque Hésiode aux vers 564-70 des \italique{Travaux et des jours}) et, par suite, l'heure de la nuit où cette étoile est visible, savoir : peu de temps avant que le jour ne blanchisse.
Signe du jour nouveau et, croit-il, de la fin de ses tribulations, Mnésiloque guette cette étoile et non pas, comme on peut le lire dans les différents commentaires, une \italique{hirondelle}, synonyme de printemps et de jours meilleurs (la période des Thesmophories, en plein mois de Pyanepsion, c'est-à-dire pendant notre seconde moitié d'octobre serait un peu prématurée pour une supplique au retour du printemps).}%
\debutNote{1}%
\note{1}{Dans l'article « Thesmophoriazusae's Two Dawns », in \italique{Mnemosyne} 62 (2009), pp. 535-547. A. Clements se base sur l'article de Mary Blomberg « The meaning of χελιδών in Hesiod », \italique{Opuscula Atheniensia} 19 (1992), pp. 49-57. \italique{Confer} le fragment 617 d'Aristophane (ed. Kassel-Austin) : πυθοῦ χελιδὼν πηνίκ᾽ ἄττα φαίνεται.}%
%\commentaire{2-3}{χελιδὼν ἄρα πότε φανήσεται; / ἀπολεῖ μ᾽ ἀλοῶν ἅνθρωπος ἐξ ἑωθινοῦ.}{Ce passage avait jusqu'ici été interprété dans les diverses éditions des \italique{Thesmophories} avec deux erreurs à la clef : \ \ \ 1°) la demande du protagoniste Mnésiloque pour qu'apparaisse une hirondelle était comprise, par métonymie, comme prière pressante pour qu'enfin revienne le printemps – synonyme de renaissance et de jours meilleurs –, car en grec comme en français, une hirondelle ne fait pas le printemps, mais du moins l'annonce-t-elle\marqueur{2}. \ \ \ 2°) ἐξ ἑωθινοῦ était traduit par « dès le matin » ou « au point du jour ».\\
%Pour le premier point, rendons à tout seigneur tout honneur : le mérite revient à Ashley Clements\marqueur{1} d'avoir correctement interprété le mot χελιδὼν : il s'agit non pas ici d'une hirondelle, symbole du retour du printemps – interprétation qui ne laissait pas de soulever quelque difficulté : la scène est censée se dérouler au beau milieu des Thesmophories, le 12e jour du mois de Pyanepsion, c'est-à-dire en octobre-novembre alors que la nuit triomphe chaque jour de plus en plus... période un peu prématurée pour une supplique au retour du printemps – en fait, Χελιδὼν désigne une étoile (qu'évoque Hésiode aux vers 564-70 des \italique{Travaux et des jours}\marqueur{2}) et, par suite, l'heure de la nuit où cette étoile est visible, savoir : peu de temps avant que le jour ne blanchisse. Mnésiloque guette cette étoile, signe du jour nouveau et, croit-il, de la fin de ses tribulations.Sitôt le sens de χελιδὼν rétabli, on comprend qu'il faut donner à ἐξ ἑωθινοῦ le sens (établi par le Bailly\marqueur{4}) : \italique{pendant la veille du matin}, c'est-à-dire pendant la \italique{quatrième veille} (en latin : \latin{quarta vigilia}). Car le soleil ne s'est pas encore levé : Mnésiloque parviendra à rejoindre la procession des Thesmophories à temps alors que les femmes se réunissent à l'aube (cf. v. 376). De plus, Mnésiloque et Euripide trouvent Agathon encore dans son lit alors que son serviteur vient de préciser qu'il avait besoin du soleil pour modeler ses chants (v. 69 ; ce qui laisse à penser que le soleil ne s'est pas encore levé à ce moment) et quand Agathon refuse de se rendre aux Thesmophories, c'est au prétexte de ne pas voler les œuvres \italique{nocturnes} (v. 204) des femmes. Tous ces éléments rendaient contradictoires les traductions parlant de lever du jour dès le second vers ; contradictions qui n'avaient d'ailleurs pas manqué d'être relevées, mais certains commentateurs préféraient en reprocher la faute à notre poète !.. Passons... Ou plutôt, ne passons pas : je profite de ce premier commentaire pour exposer la méthode suivie à chaque fois que j'achoppai sur quelque difficulté :\ \ \ 1° m'en prendre à mon ignorance profonde du grec, tant de sa langue que des us et coutumes de l'époque attique (et parfois aussi chez Aristophane la difficulté supplémentaire à saisir des allusions à des faits que la chronique du temps ne nous a pas conservés). Cette ignorance doit être contrebalancée par une consultation assidue de ce que les lexiques et l'archéologie peuvent nous apporter \ \ \ 2° m'en prendre au texte, mal transmis ; recopier un texte n'est pas tâche aisée : il n'y a qu'à regarder certains manuscrits anciens : au premier abord un bloc de pattes de mouches peu avenant et, en s'y penchant de plus près, on découvre tout un système d'abréviations, de renvois, \&c. inventé par un scribe n'ayant pas laissé de mode d'emploi, surajoutez à cela des gloses inutiles – donc nuisibles – intercalées çà-et-là pour embrouiller le tout et vous comprendrez que la – finalement bonne – conservation des textes tient du miracle et d'un aride travail de bénédictins et autres copistes \ \ \ 3° si toutes les sources, indépendantes et concordantes, établissent un texte sûr dont le sens est certain et qu'il n'en demeure pas moins des incongruités, alors il ne reste qu'à s'en prendre à l'auteur lui-même qui aura écrit quelque bêtise... Évidemment, en arriver-là avec Aristophane, le plus fin des attiques, relève du péché de présomption ! ce n'est pas avec mon grec anémié, malappris dans le peu qu'il nous reste de ses comédies que je vais en remontrer à ce maître de la langue. Pour le cas présent des \italique{Thesmophories}, c'est même méthodologiquement impossible : on ne dispose que d'un seul témoin complet, le codex 429 de la \italique{Biblioteca Classense} de Ravenne du \nieme{X} siècle et seulement quelques fragments épars de papyrus, l'on restera donc bloqué à la deuxième étape. Si les différents éditeurs d'Aristophane s'en étaient tenus à cette méthode ou quelqu'une analogue, cela nous aurait évité d'avoir à lire bien des sottises – que j'épargne dans la plupart des cas au lecteur, encore que l'exposition et le châtiment en place publique soient la plus sûre façon pour se prémunir de ce genre de comportement à l'avenir... voir !}%
%\debutNote{1-2}%
%\note{1}{Dans l'article \italique{Thesmophoriazusae's Two Dawns}, in \italique{Mnemosyne} 62, 2009, p. 535-547.}%
%\note{2}{Cf. les \italique{Oiseaux}, v. 1421, ou la réplique de ce personnage des \italique{Cavaliers} : « Voici le printemps : une hirondelle ! (v. 419) »}%
%\note{3}{A. Clements s'appuie pour l'exégèse du passage d'Hésiode sur l'article de M. Blomberg, \italique{The Meaning of ΧΕΛΙΔΩΝ in Hesiod}, \italique{Opuscula Atheniensia} 19, p. 49-57. On pourrait aussi citer le fragment 617 (601) d'Aristophane, dans l'édition de R. Kassel \& C. Austin, \italique{Poetae Comici Graeci}, vol. III 2, 1984 : πυθοΰ χελιδών πηνίκ' άττα φαίνεται;}%
%\note{4}{Le Bailly atteste de ce sens à l'appui des passages de l'historien Polybe de Megalopolis (\nieme{II} s. av. J.-C.) 3, 43, 1 ; 3, 67, 9 ; etc.}%
\\
\commentaire{2}{ἀλοῶν.}{Ἀλοᾶν est un terme agricole signifiant \italique{battre le grain}, selon une technique bien définie : sur l'aire à battre (ἡ ἅλως), on fait tourner des bœufs, attachés à un piquet, qui piétinent et qui broient la céréale jonchant le sol. De ce fait, ἁλοᾶν est devenu un synonyme expressif de \grec{τύπτειν} avec le sens de \italique{broyer de coups} (note de Taillardat\marqueur{1})\\
Ici, ἁλοᾶν a le sens moral de \italique{tourmenter}. Conformité du grec et du français -- ou du latin, plutôt : l'on possède la même image en avec notre \italique{tribulation}, le latin \latin{tribulare} signifie « battre le blé, presser avec la herse (\latin{tribulum}) ».}%
\debutNote{2}%
\note{1}{Je me référerai souvent – et avec force citations~!~– au magistral ouvrage que Jean Taillardat a consacré à la langue bigarrée et haute en couleurs d'Aristophane : \italique{Les images d'Aristophane} édité aux Belles Lettres (seconde édition de 1965). Je cite ici le §138 ; voir aussi les références du §606 : la métaphore était usuelle et pas de l'invention de notre poète : \italique{Grenouilles} 149, Sophocle fr. 19 \Nauck [$=$ fr. 20 Radt (dub.)] ; Hérondas 2, 34 ; Babrius, \italique{Le lion amoureux}, 98.15 ; καταλοᾶν Eschine 2, 140 ; Xénoph. \italique{Cyrop.} 7, 1, 31 ; Lucien \italique{Icarom.} 15 ; συναλοᾶν Hippon. 14AD\expo{3}.4 [$=$ fr. 92.4 Masson $=$ fr 95.4 Degani] ; Théocr. 22.128.}%
οἷόν τε, πρὶν τὸν σπλῆνα κομιδῇ μ᾽ ἐκβαλεῖν,%
\\\passeLigne}%
\end{groupe}%
%\ledpb page break
%% v. 5-6a
\begin{groupe}%
\personnageGrec{Εὐριπίδης}%
\trimIamb{ἀλλ᾽ οὐκ ἀκούειν δεῖ σε πάνθ᾽ ὅσ᾽ αὐτίκα\apparat{5}{ὅσ R : \lecon{ὅσ᾽} \Mon}%
\commentaire{3}{πρὶν τὸν σπλῆνα κομιδῇ μ᾽ ἐκβαλεῖν.}{Les dérangements (palpitations) dus à la course – d'autant plus aggravés par l'âge avancé de Mnésiloque – étaient imputés par les Anciens à des troubles de la rate (σπλήν, \anglais{spleen} en anglais)\marqueur{1}. En français, on dit de quelqu'un allant vite sans être pour autant arrêté par ce genre de dérangements qu'il court comme un «~dé-raté~».%\eff{ventre à terre, rendre les boyaux}
}%
\debutNote{3}%
\note{1}{Voir Pline l'Ancien, \italique{Histoire naturelle} 11,205 ; et l'esclave épuisé préposé à la course de Plaute (\italique{Mercator}, v. 124).}\\
ὄψει παρεστώς.\\\passeLigne}%
\end{groupe}%
%% v. 6b-7a
\begin{groupe}%
\personnageGrec{Μνησίλοχος}%
\trimIamb{\ligne{6}\decalage{ὄψει παρεστώς.}πῶς λέγεις; αὖθις φράσον.\\
οὐ δεῖ μ᾽ ἀκούειν;\\\passeLigne}%
\end{groupe}%
%% v. 7b
\begin{groupe}%
\personnageGrec{Εὐριπίδης}%
\trimIamb{\decalage{οὐ δεῖ μ᾽ ἀκούειν;}\ligne{7}οὐχ ἅ γ᾽ ἂν μέλλῃς ὁρᾶν.\\\passeLigne}%
\end{groupe}%
%% v. 8a
\begin{groupe}%
\personnageGrec{Μνησίλοχος}%
\trimIamb{οὐδ᾽ ἆρ᾽ ὁρᾶν δεῖ μ᾽;\\\passeLigne}%
\end{groupe}%
%% v. 8b
\begin{groupe}%
\personnageGrec{Εὐριπίδης}%
\trimIamb{\decalage{οὐδ᾽ ἆρ᾽ ὁρᾶν δεῖ μ᾽;}\ligne{8}οὐχ ἅ γ᾽ ἂν ἀκούειν δέῃ.\\\passeLigne}%
\end{groupe}%
%% v. 9-10
\begin{groupe}%
\ledpb%
\personnageGrec{Μνησίλοχος}%
\trimIamb{πῶς μοι παραινεῖς; δεξιῶς μέντοι λέγεις.\\
οὐ φῂς σὺ χρῆναί μ᾽ οὔτ᾽ ἀκούειν οὔθ᾽ ὁρᾶν.\apparat{10}{\lecon{ὁρᾶν.} R : ὁρᾶν; Gelenius}\\\passeLigne}%
\end{groupe}%
%% v. 11
\begin{groupe}%
\personnageGrec{Εὐριπίδης}%
\trimIamb{χωρὶς γὰρ αὐτοῖν ἑκατέρου 'στὶν ἡ φύσις.\\\passeLigne}%
\end{groupe}%
%% v. 12a
\begin{groupe}%
\personnageGrec{Μνησίλοχος}%
\trimIamb{τοῦ μήτ᾽ ἀκούειν μήθ᾽ ὁρᾶν;\\\passeLigne}%
\end{groupe}%
%% v. 12b
\begin{groupe}%
\personnageGrec{Εὐριπίδης}%
\trimIamb{\ligne{12}\decalage{τοῦ μήτ᾽ ἀκούειν μήθ᾽ ὁρᾶν;}εὖ ἴσθ᾽ ὅτι.\\\passeLigne}%
\end{groupe}%
%% v. 13a
\begin{groupe}%
\personnageGrec{Μνησίλοχος}%
\trimIamb{πῶς χωρίς;\\\passeLigne}%
\end{groupe}%
%% v. 13b-18
\begin{groupe}%
\personnageGrec{Εὐριπίδης}%
\trimIamb{\ligne{13}\decalage{πῶς χωρίς;}οὕτω ταῦτα διεκρίθη τότε.\\
αἰθὴρ γὰρ ὅτε τὰ πρῶτα διεχωρίζετο\\
καὶ ζᾦ᾽ ἐν αὑτῷ ξυνετέκνου κινούμενα,\\
ᾧ μὲν βλέπειν χρὴ πρῶτ᾽ ἐμηχανήσατο\apparat{16}{πρῶτα R : \lecon{πρῶτ᾽} Grynaeus}\\
ὀφθαλμὸν ἀντίμιμον ἡλίου τροχῷ,\\
ἀκοὴν δὲ χοάνης ὦτα διετετρήνατο.\testimonia{18-9}{(χοάνη = χώνη) Moeris χ 1 ; cf. Hesychius χ 576 χοάνη}\apparat{18}{\lecon{ἀκοὴν δὲ χοάνης} R, Willems : ἀκοὴς δὲ χοάνην Biset : ἀκοῇ δὲ χοάνην Boissonade}%
\commentaire{13-8}{}{Aristophane nous peint Euripide comme un de ces faiseurs de cosmogonies (tel le Socrate des \italique{Nuées} avec son Δῖνος). Le véritable Euripide avait effectivement un penchant pour ce genre de nuées\marqueur{1}. En l'occurrence, c'est un pastiche de cosmogonie qui semble inspiré d'Empédocle\marqueur{2}.}%
\debutNote{13-8}%
\note{1}{Voir par exemple le fragment 5 de \italique{Mélanippe la sage} dans l'édition de Jouan-van Looy aux Belles Lettres, tome VIII, \nieme{2} partie ($=$ fragment 484 de l'édition Kannicht).}%
\note{2}{Cf. notamment ce passage d'Eusèbe de Césarée (\italique{Préparation évangélique}, XV, 37 ; fragment A49 Diels-Kranz) : Ἐμπεδοκλῆς τὸν μὲν αἰθέρα πρῶτον διακριθῆναι et la création de l'œil par Aphrodite (B84.9 D.-K.) reconstitué par Blass (1883) : χοάνῃσι δίαντα τετρήατο θεσπεσίῃσιν.}%
\commentaire{16-8}{πρῶτα ἐμηχανήσατο ... ἀκοὴν δὲ χοάνης ὦτα διετετρήνατο.}{Pour l'explication de ce passage diversement retouché au gré des éditeurs, je ne peux que renvoyer au brillant exercice d'exégèse d'Alphonse Willems (ses notes sur les \italique{Thesmophories} sont reproduites en annexe, celle relative à ce passage en est la première. À propos d'Alphonse Willems, \italique{confer} l'introduction.) : Euripide recourt à une image tirée de la fonderie -- certainement plus parlante aux spectateurs de l'époque qu'à nous : l'oreille (ἀκοή) est comparée aux ὦτα (littéralement les \italique{oreilles} et, par analogie, les \italique{anses}, \italique{poignées} ou, comme ici, les \italique{godets} ou \italique{entonnoirs} d'un moule, χοάνη). Willems construit ainsi la phrase, en accord avec la sobre indication du scholiaste (λείπει ὡς) : διετετρήνατο δὲ ἀκοὴν (ὡς) χοάνης ὦτα.}%
\\\passeLigne}%
\end{groupe}%
%% v. 19-21
\begin{groupe}%
\personnageGrec{Μνησίλοχος}%
\trimIamb{διὰ τὴν χοάνην οὖν μήτ᾽ ἀκούω μήθ᾽ ὁρῶ;
\commentaire{19}{διὰ τὴν χοάνην οὖν μήτ᾽ ἀκούω μήθ᾽ ὁρῶ;}{Mnésiloque ne goûte guère aux subtilités des poètereaux et tord à plaisir leurs propos. Ses saillies fort désobligeantes se repèrent aisément puisqu'elles sont toujours superbement ignorées. C'en est ici le premier exemple ; il est évident que Mnésiloque se livre ici à un détournement, instantanément intelligible de l'Athénien du \nieme{V} siècle av. J.-C... mais qui nous est bien plus difficile à saisir.
On peut toutefois proposer une hypothèse.
Le tour `μήτ᾽ ἀκούω μήθ᾽ ὁρῶ'\marqueur{1} nous semble à rapprocher de l'expression proverbiale latine : \latin{Nec aures me credo habere nec tango}. La langue d'Aristophane, comme nous l'allons voir en chemin, est truffée d'échos aux proverbes grecs ; on ne peut que regretter le manque d'intérêt criant des commentateurs d'Aristophane pour ce domaine. Nous nous efforcerons au contraire, et selon nos moyens, de souligner chacun des emprunts d'Aristophane à cette langue proverbiale. Et ils sont nombreux ! Le cas présent est malheureusement moins probant puisque je n'ai pas pu trouver d'autres occurrences dans la littérature grecque. Revenons au proverbe latin\marqueur{2}, \latin{Nec aures me credo habere nec tango}, «~je crois n'avoir plus ni oreilles ni main~» ; il est employé pour signifier qu'on ne se sent plus la force de faire quoi que ce soit. Si donc Mnésiloque est complètement \italique{moulu}, c'est d'après les explications d'Euripide, la faute à ce \italique{moule} (διὰ τὴν χοάνην). Le jeu de mots exact en grec, qui tourne autour de χοάνη, nous échappe.
%\eff{Nοῦς ὁρᾷ, καὶ νοῦς ἀκούει: ?Érasme 2500 ; n° 1971 du Dictionnaire des sentences latines et grecques.}
%L'explication qui nous semble la plus simple est un jeu de mots sur les différents sens de διά + accusatif.\\
%Le premier : \italique{à cause de} ; si donc Mnésiloque ne peut ni voir ni entendre, c'est \italique{à cause de} cette fichue χοάνη... conclusion apparemment logique du long exposé justificatif d'Euripide.\\
%Second sens possible de διά + accusatif : \italique{à travers}. La χοάνη est à prendre dans un sens littéral (de \italique{tamis} ou de \italique{filtre} ?), opaque à la lumière et au son ; le vers suivant est d'une insolente ironie pour Euripide (il fallait au moins un Euripide pour parvenir à la démonstration qu'on ne peut ni voir ni entendre \italique{à travers} un tel objet !)
}%
\debutNote{19}%
\note{1}{«~Je ne vois plus Claire ; j'ai perdu Louis.~»}%
\note{2}{Je renvoie le lecteur à l'admirable \italique{Dictionnaire des sentences latines et grecques} de Renzo Tosi (traduit de l'italien par Rebecca Lenoir ; proverbe n° 1971, cité par saint Jérôme notamment). Érasme avait aussi recensé le proverbe dans ses \italique{Adages} (n° 2500). (Je cite les \italique{Adages} ainsi que leur traduction française selon l'édition de Jean-Christophe Saladin aux Belles Lettres (2019).)}\\
νὴ τὸν Δί᾽ ἥδομαί γε τουτὶ προσμαθών.\\
οἷόν γέ που 'στιν αἱ σοφαὶ ξυνουσίαι.\apparat{21}{τε R : \lecon{γε} Schol.}%
\commentaire{21}{οἷόν τέ που 'στιν αἱ σοφαὶ ξυνουσίαι.}{Le scholiaste note à ce vers une allusion à un mot célèbre : σοφοὶ τύραννοι τῶν σοφῶν ξυνουσίᾳ, qu'Aristophane dans ses \italique{Héros}\marqueur{1} mais aussi Platon\marqueur{2} et Anthistène\marqueur{3} attribuent à Euripide. Cela n'empêche pas moins notre scholiaste, bien sûr de lui, d'accorder la paternité de cette formule à Sophocle\marqueur{4} (dans son Αἴας Λοκρός) ; or il le fait avec une telle assurance et une telle érudition qu'on serait tenté de le croire\marqueur{5}. Après tout, Euripide est coutumier du fait : Aulu-Gelle\marqueur{6} nous rapporte entre autres exemples que le vers 193 des \italique{Bacchantes}, γέρων γέροντα παιδαγωγήσω σ᾽ ἐγώ, se retrouve aussi à l'identique chez Sophocle\marqueur{7} ; de même, Aristote\marqueur{8} cite une reprise d'Eschyle par Euripide. Toujours est-il qu'Euripide n'en passait certainement pas moins aux yeux d'Aristophane et de ses spectateurs pour le père putatif de la citation. Sur le sens de γέ που, on lira avec intérêt la seconde note de Willems reproduite ci-après en annexe.}%
\debutNote{21}%
\note{1}{Fragment 323 de l'édition Kassel-Austin.}%
\note{2}{\italique{Théagès}, 125b et la \italique{République}, VIII, 568b. Platon devait avoir un point de vue sur Euripide assez proche de celui d'Aristophane : il chasse joyeusement de sa République idéale Euripide, cet homme qui occupe une place prestigieuse parmi les semi-habiles de sa cité.}%
\note{3}{Fragment 59 de l'édition Decleva-Caizzi}%
\note{4}{Fragment 14 de l'édition de Stefan Radt, \italique{Tragicorum Graecorum Fragmenta}, vol. 4, édition corrigée et augmentée de 1999. Il figure aussi parmi les fragments incertains d'Euripide (n° 888b) dans l'édition Kannicht des \italique{Tragicorum Graecorum Fragmenta}, vol 5.% \eff{chez Jouan - van Looy ?}
}%
\note{5}{Preuve que la scholie est de Didyme, selon le critère de Pierre Boudreaux ? (\italique{Le texte d'Aristophane et ses commentateurs} (1920). Boudreaux, brillant helléniste fauché bien jeune par la Grande guerre, réussit le tour de force d'écrire un bouquin -- publié de façon posthume -- passionnant quoique traitant de sujets à priori fort arides.)}%
\note{6}{Dans ses \italique{Nuits attiques}, livre XIII, chapitre 18.}%
\note{7}{Fragment 695 chez Radt.}%
\note{8}{Aristote, Περὶ ποιητικῆς, XIII, 1458b22. À rebours de la querelle entre Eschyle et Euripide dans les \italique{Grenouilles}, Aristote cite un vers iambique d'Eschyle jugé plat en raison du recours à un mot courant tandis qu'Euripide reprit le vers en lui substituant un mot emprunté et rare : οἷον τὸ αὐτὸ ποιήσαντος ἰαμβεῖον Αἰσχύλου καὶ Εὐριπίδου, ἓν δὲ μόνον ὄνομα μεταθέντος, ἀντὶ κυρίου εἰωθότος γλῶτταν, τὸ μὲν φαίνεται καλὸν τὸ δ᾽ εὐτελές. Αἰσχύλος μὲν γὰρ ἐν τῷ Φιλοκτήτῃ ἐποίησε·\citComm{φαγέδαιναν ἥ μου σάρκας ἐσθίει ποδός,} ὁ δὲ ἀντὶ τοῦ `ἐσθίει' τὸ `θοινᾶται' μετέθηκεν.}%
\\\passeLigne}%
\end{groupe}%
%% v. 22a
\begin{groupe}%
\ledpb%
\personnageGrec{Εὐριπίδης}%
\trimIamb{πόλλ᾽ ἂν μάθοις τοιαῦτα παρ᾽ ἐμοῦ.\\\passeLigne}%
\end{groupe}%
%% v. 22b-24
\begin{groupe}%
\personnageGrec{Μνησίλοχος}%
\trimIamb{\ligne{23}\decalage{πόλλ᾽ ἂν μάθοις τοιαῦτα παρ᾽ ἐμοῦ.}πῶς ἂν οὖν\apparat{v. 23}{πρὸς τοῖς ἀγαθοῖς τούτοισιν ἐξεύροιμ᾽ ὅπως R : \lecon{del.} Willems \ligneApp{24} προσμάθοι· μὴ R : \lecon{προσμάθοιμι} Ellebodius}\\
ἔτι προσμάθοιμι χωλὸς εἶναι τὼ σκέλει;\testimonia{24}{\Macaire VI 90 ; \Apostol II 94}%\apparat{24}{προσμάθοι· μὴ R : \lecon{προσμάθοιμι} Ellebodius}%
\commentaire{22-4}{πῶς ἂν οὖν ... προσμάθοιμι χωλὸς εἶναι τὼ σκέλει;}{Aristophane joue sur un détournement de proverbe équivalent à notre «~Qui fréquente le boiteux apprend à boiter~» : Παρὰ χωλὸν οἰκῶν κἂν ἐπισκάζειν μαθήσῃ\marqueur{1}.
Le boiteux -- ou plutôt le \italique{faiseur de boiteux}, χωλοποιός\marqueur{2} -- c'est bien-sûr Euripide et son fâcheux penchant pour mettre en scène des personnages peu reluisants\marqueur{3}.
La pointe de la plaisanterie de Mnésiloque est rejetée à la toute fin ; quand il précise vouloir être boiteux τὼ σκέλει, \italique{des deux jambes}. 
% \eff{Érasme 2256 et 2584 et 422n.2 et 603: c'est-à-dire être versatile comme l'Injuste ... mais un examen approfondi des proverbes grecs tend à privilégier une troisième hypothèse : ...}
%Deux interprétations me semblent possibles et, en l'absence d'élément pouvant privilégier l'une sur l'autre, je te soumets les deux, ô lecteur :\\
%Soit le Mnésiloque historique était un véritable boiteux (comme on peut le penser à la peine qu'il a pour suivre le déjà vieil Euripide au début de la pièce) ; souhaiter claudiquer des deux côtés, c'est souhaiter ne plus claudiquer, comme dans la chanson de Nadaud : "Si je marche plus droit qu'un autre / C'est que je boite des deux pieds".\\
Il est peu vraisemblable de supposer que le Mnésiloque historique fût un véritable boiteux (seul son âge est en cause lorsqu'il peine à suivre Euripide). Nous ne sommes pas dans la chanson de Nadaud où souhaiter claudiquer des deux côtés, c'est souhaiter ne plus claudiquer :%
\citComm{Si je marche plus droit qu'un autre\\
C'est que je boite des deux pieds.}%
Le scholiaste a sans doute raison : %
\citComm{χωλὸς εἶναι· ἵνα μὴ περιπατήσῃ μετ᾽ αὐτοῦ.}%
Si Mnésiloque venait à boiter des deux jambes, il ne serait plus forcé de suivre Euripide.
Mais ce dernier ne s'en laisse pas compter ; s'il ne fait pas penser clair Mnésiloque, du moins le fait-il marcher droit (βάδιζε δευρὶ) !\\
De simples allusions suffisaient au public d'Aristophane, mais une moderne traduction devrait peut-être insister davantage et le rendre ainsi : «~Euripide, toi qui as une prédilection pour les éclopés, tu pourrais me faire deux belles jambes en coupant court à tous ces va-et-vient~», ou bien : «~m'apprendre à te désemboîter le pas~».\\
Willems a montré dans \nieme{3} note, reproduite en annexe, que le vers 23 n'était qu'une glose explicative du suivant.
%L'autre explication est celle du scholiaste : si Mnésiloque venait à boiter des deux jambes, il ne serait plus forcé de suivre Euripide.\\
%Ma préférence va évidemment à la première hypothèse... que je ne suis malheureusement pas parvenu à étayer : ni pour montrer que Mnésiloque soit boiteux, ni que "boiter des deux jambes" pût signifier pour un Grec "marcher droit" ; je m'en remets donc à toi, ô lecteur : puisses-tu m'apporter quelque éclaircissement !\\
}%
\debutNote{22-4}%
\note{1}{C'est le proverbe n° 1674 du \italique{Dictionnaire des sentences latines \& grecques}. Recensé par Macaire au \nieme{XIV} siècle dans son anthologie Ῥοδωνιά (VI, 90 de l'édition de Leutsch \& Schneidewim du \italique{Corpus Paroemiographorum Graecorum}, vol. II) ainsi que sous d'autres formes proches par Michel Apostolius, Plutarque et un scholiaste de Pindare.}%
\note{2}{\italique{Les Grenouilles}, vers 846.}%
\note{3}{Cf. les \italique{Acharniens} v. 411 sq. : Télèphe, Bellérophon, Philoctète le mendiant ; voir aussi la \italique{Paix} v. 147 ; les \italique{Oiseaux} v. 1293 et le fragment 53 de l'édition Kassel-Austin des \italique{Poetae Comici Graeci}, vol. III 2 (1984).}\\\passeLigne}%
\end{groupe}%
%% v. 25a
\begin{groupe}%
\personnageGrec{Εὐριπίδης}%
\trimIamb{\ligne{25}βάδιζε δευρὶ καὶ πρόσεχε τὸν νοῦν.\\\passeLigne}%
\end{groupe}%
%% v. 25b
\begin{groupe}%
\personnageGrec{Μνησίλοχος}%
\trimIamb{\ligne{25}\decalage{βάδιζε δευρὶ καὶ πρόσεχε τὸν νοῦν.}ἰδού.\\\passeLigne}%
\end{groupe}%
%% v. 26a
\begin{groupe}%
\personnageGrec{Εὐριπίδης}%
\trimIamb{ὁρᾷς τὸ θύριον τοῦτο;\\\passeLigne}%
\end{groupe}%
%% v. 26b-27a
\begin{groupe}%
\personnageGrec{Μνησίλοχος}%
\trimIamb{\ligne{26}\decalage{ὁρᾷς τὸ θύριον τοῦτο;}νὴ τὸν Ἡρακλέα\\
οἶμαί γε.\\\passeLigne}%
\end{groupe}%
%% v. 27b
\begin{groupe}%
\personnageGrec{Εὐριπίδης}%
\trimIamb{\ligne{27}\decalage{οἶμαί γε.}σίγα νυν.\apparat{27}{νῦν R : \lecon{νυν} Boissonade}\\\passeLigne}%
\end{groupe}%
%% v. 27c
\begin{groupe}%
\personnageGrec{Μνησίλοχος}%
\trimIamb{\ligne{27}\decalage{οἶμαί γε.σίγα νυν.}σιωπῶ τὸ θύριον;\apparat{27}{θύριον R : \lecon{θύριον;} Dindorf}%
%\commentaire{27}{σιωπῶ τὸ θύριον}{\eff{Lire σιω πῶ τὸ θύριον; ? du verbe σείω ? puis a-krouw ?}}
\\\passeLigne}%
\end{groupe}%
%% v. 28a
\begin{groupe}%
\personnageGrec{Εὐριπίδης}%
\trimIamb{ἄκου᾽.\\\passeLigne}%
\end{groupe}%
%% v. 28b
\begin{groupe}%
\personnageGrec{Μνησίλοχος}%
\trimIamb{\ligne{28}\decalage{ἄκου᾽.}ἀκούω καὶ σιωπῶ τὸ θύριον;\apparat{28}{ἀκούσω R : \lecon{ἀκούω} Brunck \espApp vers. 28 post 30 transp. \Rac : \lecon{corr.} \Rpc }\\\passeLigne}%
\end{groupe}%
%% v. 29-30a
\begin{groupe}%
\ledpb%
\personnageGrec{Εὐριπίδης}%
\trimIamb{ἐνταῦθ᾽ Ἀγάθων ὁ κλεινὸς οἰκῶν τυγχάνει
\commentaire{29}{Ἀγάθων.}{Voici enfin la réponse à la question de Mnésiloque sur la destination de leur expédition (v. 4) : Euripide l'a conduit jusqu'à la demeure d'Agathon, le fameux – sauf pour Mnésiloque – poète tragique\marqueur{1}, rendu célèbre par sa victoire dès la première participation à un concours tragique : celui des Lénéennes de 416\marqueur{2}. Platon immortalisa dans son \italique{Banquet} la réception qu'il fit chez lui le lendemain de cette victoire. Le rôle d'Agathon dans les \italique{Thesmophories} est semblable à celui d'Euripide dans les \italique{Acharniens} (v. 410-489) : il fournit le héros en matériel de travestissement. Euripide confirmera cette parenté aux vers 173-174\marqueur{3}.\\
Aristophane lui rendit un hommage posthume dans les \italique{Grenouilles} aux vers 83-85.}%
\debutNote{29}%
\note{1}{On n'a conservé d'Agathon que des fragments et quelques témoignages. Ils ont été regroupés par Snell \& Kannicht au premier volume des \italique{Tragicorum Graecorum Fragmenta} (1986).}
\note{2}{Athénée, \italique{Deipnosophistes}, V, 216f-217a.}%
\note{3}{Aristophane n'est d'ailleurs pas le seul à rapprocher Euripide \& Agathon : ainsi fait Aristote dans sa \italique{Poétique} (1456a25-32) : Agathon introduisit des interludes (ἐμβόλιμα) chantés par le chœur sans rapport avec l'action, suivant la voie ouverte par Euripide : %
\citComm{καὶ τὸν χορὸν δὲ ἕνα δεῖ ὑπολαμβάνειν τῶν ὑποκριτῶν, καὶ μόριον εἶναι τοῦ ὅλου καὶ συναγωνίζεσθαι μὴ ὥσπερ Εὐριπίδῃ ἀλλ᾽ ὥσπερ Σοφοκλεῖ. τοῖς δὲ λοιποῖς τὰ ᾀδόμενα <οὐδὲν> μᾶλλον τοῦ μύθου ἢ ἄλλης τραγῳδίας ἐστίν: διὸ ἐμβόλιμα ᾁδουσιν πρώτου ἄρξαντος Ἀγάθωνος τοῦ τοιούτου. καίτοι τί διαφέρει ἢ ἐμβόλιμα ᾁδειν ἢ εἰ ῥῆσιν ἐξ ἄλλου εἰς ἄλλο ἁρμόττοι ἢ ἐπεισόδιον ὅλον;\\
Le chœur doit être considéré comme l'un des acteurs ; il doit faire partie de l'ensemble et participer à l'action, non comme chez Euripide mais comme chez Sophocle. Chez tous les autres, les parties chantées n'ont pas plus de rapport avec l'histoire qu'avec une autre tragédie ; aussi chantent-ils des interludes, pratique dont l'origine remonte à Agathon. Et pourtant quelle différence y a-t-il entre chanter des interludes et adapter à une pièce une tirade ou un épisode tout entier tirés d'un autre ?}}\\
ὁ τραγῳδοποιός.\\\passeLigne}%
\end{groupe}%
%% v. 30b
\begin{groupe}%
%\lednopb%
\personnageGrec{Μνησίλοχος}%
\trimIamb{\ligne{30}\decalage{ὁ τραγῳδοποιός.}ποῖος οὗτος Ἁγάθων;\apparat{30}{Ἀγάθων R : \lecon{Ἁγάθων} Biset}\\
ἔστιν τις ἀγαθῶν; μῶν ὁ μέλας ὁ καρτερός;\apparat{31}{\lecon{(Μνη.) continuat} ἔστιν τις ἀγαθῶν R : (Εὐρ.) \Mon, edd. plerique \espApp Ἀγάθων R : \lecon{ἀγαθῶν;} Bibi}%
\commentaire{30-3}{ποῖος οὗτος Ἁγάθων; ... οὐχ ἑόρακας πώποτε.}{Les répartitions des répliques de ce passage fluctuent d'une édition à l'autre, hésitations du manuscrit de Ravenne lui-même : il comporte au niveau de ces vers de nombreuses corrections. Tâchons d'y voir plus clair !\\
La première erreur de répartition fut introduite par le scribe du manuscrit de Munich (copie du \italique{Ravennas} du \nieme{XV} siècle) et reprise depuis par la plupart des éditions : il attribua à Euripide le début du vers 31 : ἔστιν τις Ἀγάθων ---, avant que Mnésiloque n'enchaîne en demandant des précisions : μῶν ὁ μέλας ὁ καρτερός;\\
Laisser le vers entier ἔστιν τις Ἀγάθων μῶν ὁ μέλας ὁ καρτερός; dans la bouche de Mnésiloque relevait d'une construction douteuse... mais il y a une clef de compréhension beaucoup plus simple et qui dispense de découper arbitrairement les vers (d'autant qu'il n'est pas vraisemblable qu'Euripide se laisse ainsi couper le sifflet par Mnésiloque) : il suffit de lire :%
\citComm{ἔστιν τις ἀγαθῶν;}%
Mnésiloque joue sur le nom d'Agathon : \latin{nomen omen}, demande-t-il : Agathon est-il bien l'\italique{honnête homme} que suggère son nom ? Or, pour une nature rustre comme celle de Mnésiloque, un bon-homme se doit d'être un solide gaillard, à la peau bronzée par le travail en extérieur !.. et non pas au teint fromage-blanc comme ces désœuvrées de femmes ! Il continue donc sur sa lancée : μῶν ὁ μέλας ὁ καρτερός; «~n'est-ce pas le noiraud, le costaud ?~»
(Agathon est en réalité tout le contraire d'un rustaud ; Aristophane donne ainsi facilement à rire à son public, qui connaît bien l'original.)\\
Le vers suivant était initialement écrit ainsi dans le manuscrit de Ravenne :%
\citComm{οὔκ ἀλλ᾽ ἕτερός τις --- οὐχ ἑόρακα πώποτε}%
Le signe mis pour signifier un changement d'interlocuteur\marqueur{1} ayant par la suite été gratté (\latin{ut videtur}).\\
L'entièreté du vers est effectivement à restituer à Euripide : une scholie écrite juste au-dessus du vers l'atteste : %que telle était ainsi une ancienne leçon : 
\citComm{ἐπειδὴ οὐ πάλαι ἤρξατο διδάσκειν, ἀλλὰ τρισὶν πρὸ τούτων ἔτεσιν.\\
Puisqu'il n'a pas commencé récemment de faire jouer des pièces, mais depuis trois [six, en réalité] ans.}
Cette scholie est évidemment là pour expliquer l'étonnement d'Euripide : οὐχ ἑόρακας πώποτε; «~Tu ne l'as donc jamais vu ?! [alors que cela fait trois année qu'il monte des pièces !..]~» ; le scholiaste ne l'aurait pas ainsi rédigée si c'était effectivement Mnésiloque qui prononçait au vers 32 le `οὐχ ἑόρακα πώποτε'.\\
Cette interprétation, outre qu'elle est anciennement corroborée par la scholie, est aussi celle qui permet la lecture la plus drôle : οὐχ ἑόρακας πώποτε; --- μῶν ὁ δασυπώγων; --- οὐχ ἑόρακας πώποτε, «~Tu ne l'as jamais vu ?! --- N'est-il pas barbu ? --- Tu ne l'as jamais vu !..~», bien meilleure par le rythme et la répétition déconvenue d'Euripide que l'autre. 
}%
\debutNote{30-3}%
\note{1}{Non un tiret, \italique{paragraphos}, sur le manuscrit mais un double point ` : ' qui est d'usage sur le \italique{Ravennas} lorsque le changement d'interlocuteur intervient au milieu d'un vers.}\\\passeLigne}%
\end{groupe}%
%% v. 32
\begin{groupe}%
\personnageGrec{Εὐριπίδης}%
\trimIamb{οὔκ, ἀλλ᾽ ἕτερός τις· οὐχ ἑόρακας πώποτε;\apparat{32}{(Μνη.) οὐχ ... πώποτε \Rac : \lecon{(Εὐρ.) continuat} \Rpc, Schol. \espApp ἑόρακα R : \lecon{ἑόρακας} Bentley \ligneApp{v. 33} omm. \Rac, \Mon : in marg. \Rpc}% \ligneApp{33} (Μ) ($=$ (Μνη.) ?) \Rpc}
\\\passeLigne}%
\end{groupe}%
%% v. 33a
\begin{groupe}%
\personnageGrec{Μνησίλοχος}%
\trimIamb{μῶν ὁ δασυπώγων;\testimonia{33}{(δασυπώγων) Pollux II 88}\\\passeLigne}%
\end{groupe}%
%% v. 33b
\begin{groupe}%
\personnageGrec{Εὐριπίδης}%
\trimIamb{\ligne{33}\decalage{μῶν ὁ δασυπώγων;}οὐχ ἑόρακας πώποτε.\\\passeLigne}%
\end{groupe}%
%% v. 34
\begin{groupe}%
\personnageGrec{Μνησίλοχος}%
\trimIamb{μὰ τὸν Δί᾽ οὔτοι γ᾽ ὥστε κἀμέ γ᾽ εἰδέναι.\apparat{34}{καὶ μεγ᾽ \Rac (ὶ eras. \Rpc) : καὶ μεγ᾽ \Mon : \lecon{κἀμέ γ᾽} edd. omn.}\\\passeLigne}%
\end{groupe}%
%% v. 35-38
\begin{groupe}%
\ledpb%
\personnageGrec{Εὐριπίδης}%
\trimIamb{καὶ μὴν †βεβίνηκας† σύ γ᾽, ἀλλ᾽ οὐκ οἶσθ᾽ ἴσως.
\commentaire{35}{καὶ μὴν βεβίνηκας σύ γ᾽, ἀλλ᾽ οὐκ οἶσθ᾽ ἴσως.}{Le ton est inhabituel -- voire plutôt choquant -- dans la bouche d'Euripide.
Qu'à cela ne tienne : puisqu'il semble d'autre part qu'Agathon tourne le dos lorsqu'il accorde ses faveurs\marqueur{1}, nos commentateurs en concluent qu'il est logique que Mnésiloque ne le connaisse pas -- sauf éventuellement au sens biblique... de Sodome, s'entend. 
Et c'est le grave et docte Euripide qui colporterait ce genre de ragots.%nous renseigne de la sorte. 
Y a-t-il besoin de commenter ce genre de commentaires ?.. je précise : ces commentaires \italique{modernes} ; car le βεβίνηκας du \italique{Ravennas} n'a pas laissé derrière lui la moindre scholie ancienne ; chose étonnante lorsqu'on sait que les anciens commentateurs étaient presqu'autant portés au scabreux que nos modernes. Doublement étonnant même, car la forme βεβίνηκας méritait que les grammairiens s'y arrêtassent aussi : ce serait la seule attestation du verbe βινεῖν au parfait.
Il faut croire que le texte original d'Aristophane, l'\allemand{Urtext} des \italique{Thesmophories}, comportait une tout autre leçon ; une banalité sur laquelle aucun commentateur ancien ne se sera arrêté.
Las !.. on ne détruit que ce que l'on remplace... à défaut d'avoir une proposition alternative solidement étayée, il faut nous résoudre à \latin{crux}-désespérer...
%Ce †  †  
%Chaussons plutôt une bonne paire de lunettes et allons voir le texte d'un peu plus près. Von Velsen \marqueur{\italique{Aristophanis Thesmophoriazusae} (1883). Von Velsen a établi définitivement dans cette édition que le second manuscrit par lequel nous sont connues les \italique{Thesmophories}, le \italique{Monacensis gr. 492} (G dans son apparat critique, \Mon dans le nôtre), n'était qu'une copie du \nieme{XV} siècle du \italique{Ravennas}, apportant quelques émendations bien vues. Son édition compare exhaustivement le \italique{Ravennas} et le \italique{Monacensis} ; ce qui est d'un grand intérêt pour déterminer si une correction du \italique{Ravennas} est ancienne ou plus récente : il suffit d'aller comparer avec la leçon du manuscrit de Münich. Son apparat critique est le plus précis et permet d'avoir la meilleure idée possible du manuscrit de Ravenne -- avec celui de l'édition Prato ; en partie une reprise du premier mais revu et amélioré en de nombreux points. Pour qui voudrait pousser le scrupule plus loin, et puisque le manuscrit de Ravenne n'a toujours pas été numérisé -- ô scandale ! le manuscrit sans qui nous ne connaîtrions que neuf des onze pièces survivantes d'Aristophane !.. --, Van Leeuwen édita en fac-similé le manuscrit de Ravenne : \italique{Comoediae Undecim cum Scholiis : Codex Ravennas 137, 4, A} [$=$ \italique{Rav.} 429 (nouvelle cote)], \italique{phototypice editus} in \italique{Codices Graeci et Latini photographice depicti}, t. IX (1904), pp. XVIII-XIX.} remarqua que le `βεβίνηκας' du \italique{Ravennas} était la correction d'une première leçon : `βεαίνηκας' (l'alpha α ayant été gratté pour le transformer en bêta β dont la forme est très semblable dans notre manuscrit). Le scribe avait-il simplement mal formé son bêta ou bien avait-il omis cette lettre lors de la première rédaction ? 
%Si l'on fait la seconde hypothèse, il faut croire que la leçon recopiée par notre scribe était peu ou prou ΒΕΒΑΙΝΗΚΑC 
%Casanier sédentaire : pourquoi ne se croisent-ils pas ? car aucun des deux n'est un marcheur !
%Casanier -- même il ne quitte pas son lit !
%καὶ μὴ ἀναβαίνω
%καὶ συ γε (nombreuses références)
%
%    ΛΥ.                 Ἀλλ’, ὦ μέλ’, ὄψει τοι σφόδρ’ αὐτὰς Ἀττικάς,
%ἅπαντα δρώσας τοῦ δέοντος ὕστερον.
%Ἀλλ’ οὐδὲ Παράλων οὐδεμία γυνὴ πάρα,
%οὐδ’ ἐκ Σαλαμῖνος.
%ΚΛ.                                    Ἀλλ’ ἐκεῖναί γ’ οἶδ’ ὅτι   (59)
%ἐπὶ τῶν κελήτων διαβεβήκασ’ ὄρθριαι.   (60)
%ΛΥ.                 Οὐδ’ ἃς προσεδόκων κἀλογιζόμην ἐγὼ
%πρώτας παρέσεσθαι δεῦρο τὰς Ἀχαρνέων @1
%γυναῖκας, οὐχ ἥκουσιν.
%ΚΛ.                                       Ἡ γοῦν Θεογένους   (63)
%ὡς δεῦρ’ ἰοῦσα τἀκάτειον ᾔρετο.
%
% W.-Moellendorff, Zum Lexikon des Photios, SPAW 1907,2-14 (= Kleine Schriften, IV,. 528-541).
%Lysistrata 966 καὶ μὴ βινῶν τοὺς ὄρθρους;
%
%βινεῖν· παρὰ Σόλωνι τὸ βίᾳ μίγνυσθαι. τὸ δὲ κατὰ νόμον ὀπύειν + retrouver passage euripide. Cf. : Héraclides
% Χο.        εἰκὸς θεῶν ἱκτῆρας αἰδεῖσθαι, ξένε,
%καὶ μὴ βιαίῳ χερὶ δαιμόνων
%ἀπολείπειν [σ’] ἕδη· @1
%πότνια γὰρ Δίκα τάδ’ οὐ πείσεται.
%Κη.        ἔκπεμπέ νυν γῆς τούσδε τοὺς Εὐρυσθέως,    (105)
%
%
% Ποῖος γὰρ <ἔτ’> ἂν νέφρος ἀντίσχοι,
%ποία ψυχή, ποῖοι δ’ ὄρχεις,
%ποία δ’ ὀσφύς, ποῖος δ’ ὄρρος
%κατατεινόμενος   (965)
%καὶ μὴ βινῶν τοὺς ὄρθρους;
%
%Même problème de concentration qu'avec la belette / le calme ?
%Ecclesiazusae
%Νεις.        αἰαῖ, τί ποτε πείσομαι;
%οὐχ ἥκει μοὐταῖρος
%μόνη δ’ αὐτοῦ λείπομ’· ἡ   (912)
%γάρ μοι μήτηρ ἄλλῃ βέβηκεν.
%
%Ma ton Dia out'(e) oig(w) wste kai (e)me g' eidenai
%-- kai mhn bebhnai kai su g(e) allh\& ()
%kai mh βεβιακαs
%kai mhn dis enika
%Semble confirmer le commentaire du scholiaste au vers 258\marqueur{1} :  Profitons-en pour préciser que c'est bien le fait qu'Agathon se travestisse qui aurait pu l'empêcher d'être connu de Mnésiloque – sauf au sens biblique de Sodome – et non pas, comme l'ont soutenu nombre de commentateurs qui ont une vision bien réductrice de ce genre de rencontres, le simple fait qu'il tourne le dos lorsqu'il accorde ses faveurs.\eff{Le ton est inhabituel dans la bouche d'Euripide.... si bien que l'on se demande s'il ne vaudrait pas mieux lire bainw, piétiner, et le public d'interpréter différemment. Quoi qu'il en soit, ce vers semble confirmer}
}%
%anaphore allitération ἡνίκ’ ἂν σὺ νι-   (1210) κᾷς λέγων τὰς δίκας.
%\debutNote{35}%
%\note{1}{Commentaire paraissant de prime abord tout-à-fait fantaisiste et émanant d'une imagination délurée. Mais il semble corroboré par le }
\debutNote{35}%
\note{1}{Voir le vers 258 et sa scholie ; Agathon se livrerait au travestissement dans l'interlope \italique{Athens by night} pour y être éprouvé à la manière d'une femme -- mais le commentaire déluré du scholiaste n'est peut-être qu'une mauvaise interprétation du présent vers 35 adjointe à une forte dose d'imagination lubrique. Ce qui est sûr, c'est qu'au vers 205, Agathon dit craindre s'il était reconnu d'être lynché pour avoir chapardé nuitamment les œuvres de la Cypris : les Athéniennes voient en lui un concurrent irrégulier.}%
\\
\lednopb%
ἀλλ᾽ ἐκποδὼν πτήξωμεν, ὡς ἐξέρχεται\testimonia{36}{(ἀλλ᾽ -- πτήξωμεν) Suidas π 3024}\\
\lednopb%
θεράπων τις αὐτοῦ πῦρ ἔχων καὶ μυρρίνας\commentaire{37}{μυρρίνας}{désigne les \italique{rameaux de myrte} que l'on brûlait lors des sacrifices (cf. les \italique{Guêpes}, v. 860 ; les \italique{Oiseaux} au v. 43).}\\
\lednopb%
προθυσόμενος ἔοικε τῆς ποιήσεως.\\\passeLigne}%
\lednopb%
\end{groupe}%
%% v. 39-41
\begin{groupe}%
\ledpb%
\personnageGrec{Θεράπων}%
\typeVers{dim. anap.}\dimAnap{εὔφημος πᾶς ἔστω λαός,\testimonia{39-41}{Suidas ε 3791 (λεὼς codd. AF : λαὸς codd. GIVM, συγκλείσας, ἐνδημεῖ pro ἐπιδημεῖ )}\apparat{39}{λέως R, Suid. ε 3792 codd. AF : \lecon{λαός} Schol., Suid. codd. GIVM}\\
στόμα συγκλῄσας· ἐπιδημεῖ γὰρ\testimonia{40-2}{(ἐπιδημεῖ -- μελοποιῶν) Suidas θ 378}%
\commentaire{39-62}{Passage en dimètres anapestiques.}{(Ceux aux vers 42 et 62 sont catalectiques.) Ce mètre est d'usage courant chez Aristophane pour les prières (cf. les \italique{Guêpes} v. 323-333 et la \italique{Paix} v. 947-1015). Le serviteur d'Agathon est tout empreint de la préciosité du maître, comme, du reste, le serviteur de l'Euripide des \italique{Acharniens} (v. 396-400) l'était.}%
\commentaire{39-40}{εὔφημος πᾶς ἔστω λέως / στόμα συγκλῄσας·}{Ce genre de formule est classique au début d'une prière ou d'un sacrifice. On peut néanmoins constater ici un peu plus d'emphase avec le tour tragique\marqueur{1} et l'utilisation de l'adjectif εὔφημος plutôt que le classique verbe εὐφημεῖν\marqueur{2} ou substantif εὐφημία\marqueur{3}.}%
\debutNote{39-40}%
\note{1}{C'est ce que souligne le scholiaste : παρὰ τοῖς τραγικοῖς οὕτω λέγεται, ``<εὔφημος> πᾶς ἔστω λαός''.}%
\note{2}{\italique{Acharniens} v. 241 ; \italique{Cavaliers} v. 1316 εὐφημεῖν χρὴ καὶ στόμα κλείεν ; \italique{Nuées} v. 263 ; \italique{Paix} v. 434 et 1316 ; \italique{Grenouilles} v. 354.}%
\note{3}{\italique{Guêpes} v. 868 ; \italique{Oiseaux} v. 959.}%
\commentaire{40}{ἐπιδημεῖ.}{\italique{Séjourner}, tel me paraît bien le sens de ἐπιδημεῖ, \latin{peregre advenit et adest}. (Xénophon, \italique{Mémor.}, I, 2, 61.) Cf. ἐπιχωρτάζει, \latin{venit ibique commoratur}. Mελοποιῶν a le sens du futur (voir note sur les \italique{Grenouilles} au v. 148, t. III p. 109-110). (Note de Willems.)}\\
θίασος Μουσῶν ἔνδον μελάθρων}%
\end{groupe}%
%% v. 42-45a
\begin{groupe}%
\typeVers{(catalect.)}\dimAnap{τῶν δεσποσύνων μελοποιῶν.
\commentaire{41-2}{θίασος Μουσῶν ἔνδον μελάθρων / τῶν δεσποσύνων μελοποιῶν.}{Pour bien goûter ces vers, il faut savoir ce qu'est un \italique{thiase} : une bruyante troupe de gens célébrant des rites (bachiques la plupart du temps) et parcourant les rues, les femmes munies d'un tambourin en interprétant des danses scandées de cris rituels (évohé) avec une fougue poussée jusqu'à la frénésie (cf. les \italique{Bacchantes} d'Euripide). Par extension, ce mot désigne une foule extrêmement bruyante. Il est très prisé des tragiques ; Euripide l'emploie à propos d'une bande de centaures (\italique{Iphigénie à Aulis}, v. 1059) ; une troupe endiablée de jeunes danseuses (\italique{Iphigénie en Tauride}, v. 1146) ou encore une meute de guerriers (\italique{Phaéton}, v. 796). Ici, nous sommes en présence d'un thiase de... Muses !.. Muses qui viennent séjourner chez Agathon à l'occasion de sa composition. Comme le remarque Sommerstein, l'inversion est totale : au lieu d'une troupe d'humains honorant une divinité, c'est une troupe de Muses qui vient honorer Agathon\marqueur{1}.}%
\debutNote{41-2}%
\note{1}{L'image banale en grec étant celle de l'\italique{essaim}. Cf. \italique{Lysistrata} (vers 353), ἑσμὸς γυναικῶν ; \italique{Nuées} (vers 297), μέγα γάρ τι θεῶν κινεῖται σμῆνος. Voir Jean Taillardat, \italique{Les images d'Aristophane}, §662.}\\
ἐχέτω δὲ πνοὰς νήνεμος αἰθήρ,\testimonia{43-5}{(ἐχέτω -- γλαυκόν) Suidas ε 3792 ; Zonaras 941 ἐχέτω}%
\commentaire{43-7}{ἐχέτω δὲ πνοὰς νήνεμος αἰθήρ, ...}{De même que dans les \italique{Bacchantes} d'Euripide (v. 1084), à la voix de Dionysos « l'éther se tut, le ravin herbeux étouffa le murmure des feuilles, et l'on n'eût point entendu le cri des bêtes sauvages ».%
\citComm{σίγησε δ᾽ αἰθήρ, σῖγα δ᾽ ὕλιμος νάπη\\
φύλλ᾽ εἶχε, θηρῶν δ᾽ οὐκ ἂν ἤκουσας βοήν.}%
Et aussi dans Théocrite (II, 38), tout se tait à l'arrivée d'Hécate. (Note de Willems.)}%
\commentaire{43}{ἐχέτω}{ est mis pour ἐπεχέτω (ici : \italique{retenir}) comme l'indique la scholie intérieure au texte du \italique{Ravennas}.}\\
κῦμα δὲ πόντου μὴ κελαδείτω\testimonia{44}{(κελαδείτω) Suidas κ 1278}\apparat{44}{\lecon{δὲ} R : τε Austin, cf. v.46, 47}\\
γλαυκόν—\\\passeLigne}%
\end{groupe}%
%% v. 45b
\begin{groupe}%
\personnageGrec{Μνησίλοχος}%
\dimAnap{\ligne{45}\decalage{γλαυκόν—}βομβάξ.\testimonia{45}{(βομβάξ) Hesychius β 787 βομβάξ ; Suidas ε 2807 ; \Eustathe 1718.9}%
\commentaire{45}{βομβάξ.}{Forgée sur βόμβος, \italique{bruit sourd}, \italique{grondement}, \italique{bourdonnement}, βομβάξ est une interjection créée pour railler un style pompeux « ran plan plan » et avec redoublement (βομβαλοβομβάξ, v. 48). (Note de Chantraine.)\\
Il pourrait être plaisant de la traduire en citant la célèbre chanson de Dalida « Paroles \& paroles \& paroles » et au redoublement : « Et encore des paroles que tu sèmes dans l'Éther ! » }\\\passeLigne}%
\end{groupe}%
%% v. 45c
\begin{groupe}%
\personnageGrec{Εὐριπίδης}%
\dimAnap{\ligne{45}\decalage{γλαυκόν—βομβάξ.}σίγα.\\\passeLigne}%
\end{groupe}%
%% v. 45d
\begin{groupe}%
\personnageGrec{Μνησίλοχος}%
\dimAnap{\ligne{45}\decalage{γλαυκόν—βομβάξ.σίγα.}τι λέγει;\\\passeLigne}%
\end{groupe}%
%% v. 46-48a
\begin{groupe}%
\personnageGrec{Θεράπων}%
\dimAnap{πτηνῶν τε γένη κατακοιμάσθω,\\
θηρῶν τ᾽ ἀγρίων πόδες ὑλοδρόμων\\
μὴ λυέσθων.\\\passeLigne}%
\end{groupe}%
%% v. 48b
\begin{groupe}%
%\lednopb%
\personnageGrec{Μνησίλοχος}%
\dimAnap{\ligne{48}\decalage{μὴ λυέσθων.}βομβαλοβομβάξ.\testimonia{48}{(βομβαλοβομβάξ) Suidas ε 2807 ; \Eustathe 790.13, 1718.13}\\\passeLigne}%
\end{groupe}%
%% v. 49-50a
\begin{groupe}%
%\ledpb%
\personnageGrec{Θεράπων}%
\dimAnap{μέλλει γὰρ ὁ καλλιεπὴς Ἀγάθων\testimonia{49}{(καλλιεπὴς) Philostrate \italique{Vit. soph.} I 9.1 [$=$ II 12.8 Kayser]}\\
πραῖνοι ἡμέτερος—\testimonia{50}{(πραῖνοι) Hesychius π 3188 πραῖνοι· πρηνίζει, καταστρέφει}\apparat{50}{πράμος R : πρόμος Grynaeus : \lecon{πραῖνοι} Bibi ap. Hesych. π 3188}%
\commentaire{49-50}{ὁ καλλιεπὴς Ἀγάθων / πράμος ἡμέτερος.}{Deux hapax dans cette courte réplique du serviteur d'Agathon : καλλιεπὴς et πράμος ; préciosité de style qui permet un détournement obscène de la part de Mnésiloque (vers 50).\\
Si le sens de καλλιεπὴς se laisse saisir sans problème, ce n'est pas le cas de πράμος, ni autrement attesté ni même relevé par les grammairiens. Pour rentrer dans le mètre anapestique, la scansion est nécessairement $\cup \cup $, ce qui empêche tout rapprochement étymologique pertinent. 
Autant d'indices qui font davantage soupçonner une erreur textuelle plutôt qu'un hapax agathonien.\\
πρόμος, \italique{le chef}, fut proposé par Grynaeus et quasi-unanimement repris depuis. 
Pourquoi cette erreur aurait-elle été commise ? ne chicanons pas paléographiquement : même s'il est rare qu'un O se transforme en A, l'erreur serait tout-à-fait possible\marqueur{1}. Cette émendation contredit surtout le principe de la \italique{lectio difficilior} : un mot certes rare, πρόμος, mais bien attesté aurait été transformé en un hapax... 
Plus important encore : il manque au terme πρόμος une acception \italique{sensu obsceno} pour expliquer l'interjection de Mnésiloque : μῶν βινεῖσθαι;\marqueur{2}\\
Heureusement pour nous, Hésychius (π 3188) n'a pas manqué de relever le vocable d'Agathon, rare et à double sens :%
\citComm{πραῖνοι· πρηνίζει, καταστρέφει.}%
Paléographiquement, ΠΡΑΙΝΟΙ explique parfaitement le ΠΡΑΜΟC de notre texte : un μ a été substitué au digramme approchant `ιν' et le ι final, légèrement incurvé, a été lu en σ (lunaire). Toutes erreurs très courantes.\\
Pour le sens πρηνίζειν signifie \italique{se jeter tête baissée}, ou simplement \italique{pencher la tête en avant} -- image d'un Agathon se jetant à son travail dans la position du penseur de Rodin ? L'occasion est trop belle pour Mnésiloque : il feint avoir compris qu'Agathon est en position κύβδα (cf. le vers 489), qu'il se \italique{retourne complètement} (καταστρέφει) et \italique{tourne le dos} à d'éventuels assaillants --- μῶν βινεῖσθαι; Le véritable Agathon devait effectivement travailler dans une posture qui prêtait à rire aux mauvais esprits, puisqu'elle lui vaut d'être comparé à Cyrène \italique{aux douze postures} (δωδεκαμήχανος) au vers 98.\\
Les deux vers 49-50 étaient donc vraisemblablement ceux-ci :%
\citComm{μέλλει γὰρ ὁ καλλιεπὴς Ἀγάθων\\
πραῖν\synizese{οι ἡ}μέτερος --- μῶν βινεῖσθαι;}%
Noter la présence d'une synérèse\marqueur{3} entre πρανοῖ et ἡμέτερος.
}%
\debutNote{49-50}%
\note{1}{L'erreur inverse se produit en effet aux vers 117 ou 1174.}%
\note{2}{Sommerstein voit en πρόμος une légère déformation suggérant à l'oreille πόρνος... «~Notre chef de file --- Celui qu'on enfile ?~» Mais que c'est loin en grec !.. Il faut croire que Sommerstein a pris un peu trop au pied de la lettre la plaisanterie dans les \italique{Nuées} (v. 394) sur l'indissociabilité -- certes pas au niveau de la prononciation -- entre βροντή et πορδή. \italique{En passant}, profitons-en pour exégéter ce jeu de mots mécompris. Il y avait en grec un proverbe, βὴξ ἀντὶ πορδῆς, «~une toux à la place d'un pet~», qui servait à désigner quelqu'un tentant vainement de camoufler une chose honteuse. (Suidas (β 261), les parémiographes ainsi qu'Érasme, au n°563 de ses \italique{Adages}, l'ont relevé.) Il devait certainement aussi exister l'adage cousin βροντὴ ἀντὶ πορδῆς ; ainsi s'explique notre passage tout comme une anecdote amusante rapportée par Strabon (14, 5, 14) et par Macaire (4, 11) : à Tarse, au temps de la lutte entre Antoine et Octave, les partisans d'un certain Boëthus, mauvais poète et ami d'Antoine, banni de la cité suite à ses exactions, couvrirent les murs d'inscriptions injurieuses contre le vieil Athénodore, qui avait chassé Boëthus, dont celle-ci : `ἔργα νέων, βουλαὶ δὲ μέσων, πορδαὶ δὲ γερόντων', « Aux jeunes l'action, aux adultes le conseil, au vieux le pet », en substituant πορδαί aux \italique{prières} (εὐχαί) du vers original d'Hésiode, entré dans la langue proverbiale (Tosi n° 1081). Athénodore avait pris la chose en riant et avait fait inscrire à son tour: « au vieux, le tonnerre vengeur » (βρονταί), faisant évidemment allusion au pouvoir qu'il détenait. Enfin, il existait aussi l'expression «~péter au nez de quelqu'un~», Aristophane, dans un de ses éclairs de génie dont il est coutumier, a fait se rencontrer ces deux images (\italique{Nuées}, vers 293 sq.) : καὶ βούλομαι ἀνταποπαρδεῖν πρὸς τὰς βροντάς. Ce n'est pas autrement que s'explique ce passage aussi. (Noter le préfixe ἀντ- de l'hapax ἀνταποπέρδω, rappel de l'ἀντὶ du proverbe.)}%
\note{3}{Ou d'une \italique{synalèphe} pour reprendre le vocabulaire de Koster dans son \italique{Traité de métrique grecque}. D'autres parlent de \italique{synizèse} ; on croise même parfois le mot \italique{crase}... Que les grammairiens se mettent d'accord pour nous proposer un jargon qui soit sans ambiguïté.}%
\\\passeLigne}%
\end{groupe}%
%% v. 50b
\begin{groupe}%
%\lednopb%
\personnageGrec{Μνησίλοχος}%
\dimAnap{\ligne{50}\decalage{πραῖνοι ἡμέτερος—}μῶν βινεῖσθαι;\\\passeLigne}%
\end{groupe}%
%% v. 51a
\begin{groupe}%
\personnageGrec{Θεράπων}%
\dimAnap{τίς ὁ φωνήσας;\\\passeLigne}%
\end{groupe}%
%% v. 51b
\begin{groupe}%
\personnageGrec{Μνησίλοχος}%
\dimAnap{\ligne{51}\decalage{τίς ὁ φωνήσας;}νήνεμος αἰθήρ.\\\passeLigne}%
\end{groupe}%
%% v. 52-57a
\begin{groupe}%
\personnageGrec{Θεράπων}%
\dimAnap{δρυόχους τιθέναι δράματος ἀρχάς.\testimonia{52-7}{Suidas δ 1547 (sed τορεύει pro τορνεύει, χιανεύει pro χοανεύει)}\testimonia{52}{Zonaras 569 δρύοχοι}\\
κάμπτει δὲ νέας ἁψῖδας ἐπῶν,\testimonia{53}{(κάμπτει) Hesychius κ 617 κάμπτειν}\apparat{53}{ἀσπίδας R : \lecon{ἁψῖδας} Suid. δ 1547, λ 181}\\
τὰ δὲ τορνεύει, τὰ δὲ κολλομελεῖ,\testimonia{54}{(τορνεύει) Hesychius τ 1168 τορνεύει}\\
καὶ γνωμοτυπεῖ κἀντονομάζει\\
καὶ κηροχυτεῖ καὶ γογγύλλει\testimonia{56}{(γογγύλλει) Hesychius γ 764 γογγύλλειν \espTest Suidas γ 361 $=$ Zonaras 449 γογγυλίζει}\apparat{56}{γογγύλίζει R, Suid. γ 361 : \lecon{γογγύλλει} Porson : γογγυλίει Hermann (ad \italique{Nub.} 992)}\\
καὶ χοανεύει—\testimonia{57}{(χοανεύει) \Eustathe 1153.40}%
\commentaire{52-7}{δρυόχους τιθέναι δράματος ἀρχάς... καὶ χοανεύει.}{Agathon le tragique n'est pas un génie bâtisseur, un τέκτων ; il ne crée pas : il fabrique, il bricole. Il est donc tout-à-tour charpentier de marine, menuisier, modeleur : « il va poser des chantiers, fondements d'un drame. etc. » Ces métaphores sont évidemment burlesques ; elles avilissent délibérément l'art d'Agathon, ravalé au rang de métier et ne peuvent que tourner le poète en dérision...\\
Agathon est donc à mille lieues du mystère de l'inspiration, décrit par Platon (dans l'\italique{Ion}) : κοῦφον γὰρ χρῆμα ποιητής ἐστιν καὶ πτηνὸν καὶ ἱερόν, « le poète est chose légère, ailée, sacrée ». De fait, Platon connaît bien les images du menuisier et du modeleur\marqueur{1} mais jamais il ne les applique au poète (ἀποτορνεύειν est dit de Lysias, \italique{Phèdre} 243e : ἀκριβῶς ἕκαστα τῶν ὀνομάτων ἀποτετόρνευται et reste péjoratif. La même observation vaut pour κολλᾶν, en parlant d'un auteur, \italique{Ménexène} 263b ; \italique{Phèdre} 278e). À noter que ces images semblent banales ; ainsi, dans le \italique{Timée} (81b), on trouve : οἷον ἐκ δρυόχων, « comme depuis les longerons de la carcasse », c'est-à-dire \italique{entièrement}. κάμπτειν : \italique{courber}, \italique{plier} désigne en musique les courbes capricieuses d'un air ou les inflexions de la voix. (Notes de Taillardat\marqueur{2}.)}%
\debutNote{52-7}%
\note{1}{Voir Pierre Louis, \italique{Les Métaphores de Platon}, Belles Lettres, 1945, p. 98, 170, 204-5.}%
\note{2}{\italique{Les images d'Aristophane}, §758 et §784.}\\\passeLigne}%
\end{groupe}%
%% v. 57b
\begin{groupe}%
\personnageGrec{Μνησίλοχος}%
\dimAnap{\ligne{57}\decalage{καὶ χοανεύει—}καὶ λαικάζει.
\commentaire{56-7}{καὶ κηροχυτεῖ καὶ γογγύλίζει / καὶ χοανεύει — καὶ λαικάζει.}{La métaphore du serviteur est tirée de la technique de la \italique{cire perdue} ; un procédé de moulage pour obtenir une statue en métal à partir d'un modèle en cire : on commence par modeler à la cire la statue (κηροχυτεῖ), modèle qu'on recouvre d'une gangue d'argile (qui fera office de moule) en laissant des cheminées, certaines pour insérer le métal en fusion qui va prendre la place de la cire, et d'autres cheminées pour permettre à la cire fondue de s'échapper. Mnésiloque voit dans cette description un double sens érotique, et interrompt alors : « καὶ λαικάζει ! » Chantraine\marqueur{1} suppose que ce verbe a le sens de \latin{fellare}. Si tel est le cas, on comprend mieux l'interprétation que fait Mnésiloque des trois verbes précédents. Pour γογγύλλειν, les commentateurs \& traducteurs donnent le sens d'\italique{arrondir} ; mais si l'on garde \latin{fellare} en tête, il semble bien préférable de lui donner le sens de \italique{rendre ramassé}, \italique{compact}, \italique{serré}, autrement attesté\marqueur{2}.}%
\debutNote{56-7}%
\note{1}{\italique{Dictionnaire étymologique de la langue grecque}.}%
\note{2}{Cf. Hésychius (γ 764) : γογγύλλειν· συστρέφειν et la \nieme{3} note §166 de Taillardat : \italique{ramassé}, \italique{compact}, \italique{serré} sont les sens que Taillardat donne à l'adjectif γογγύλος aux vers 25-28 de la \italique{Paix}.}\\\passeLigne}%
\end{groupe}%
%% v. 58
\begin{groupe}%
\personnageGrec{Θεράπων}%
\dimAnap{τίς ἀγροιώτας πελάθει θριγκοῖς;\testimonia{58}{(θριγκοῖς) Hesychius θ 750 θριγκοῖς· φραγμοῖς}\apparat{58}{ἀγριώτας R : \lecon{ἀγροιώτας} anonym. Parisinus}\\\passeLigne}%
\end{groupe}%
%% v. 59-61
\begin{groupe}%
\personnageGrec{Μνησίλοχος}%
\dimAnap{ὃς ἕτοιμος σοῦ τοῦ τε ποιητοῦ\\
τοῦ καλλιεποῦς κατὰ τοῦ θριγκοῦ\apparat{60}{κατὰ om. \Rac : suppl. manus coaeva \Rpc}\\
συγγογγύλας καὶ συστρέψας\apparat{61}{γογγυλίσας R : συγγογγυλίσας Brunck : \lecon{συγγογγύλας} Enger}%
\commentaire{61}{γογγυλίσας καὶ συστρέψας.}{Vers retouché \italique{gratia metri} d'après un passage de \italique{Lysistrata} (v. 975), jouée probablement l'année précédente (cf. Introduction).}}%
\end{groupe}%
%% v. 62
\begin{groupe}%
\typeVers{(catalect.)}\dimAnap{τουτὶ τὸ πέος χοανεῦσαι.\testimonia{62}{(χοανεῦσαι $=$ χωνεῦσαι) Moeris χ 21 ; Hesychius χ 575 χοανεῦσαι}\\\passeLigne}%
\end{groupe}%
%% v. 63
\begin{groupe}%
\personnageGrec{Θεράπων}%
\typeVers{trim. iamb.}\trimIamb{ἦ που νέος γ᾽ ὢν ἦσθ᾽ ὑβριστὴς ὦ γέρον.\\\passeLigne}%
\end{groupe}%
%% v. 64-65
\begin{groupe}%
\personnageGrec{Εὐριπίδης}%
\trimIamb{ὦ δαιμόνιε τοῦτον μὲν ἔα χαίρειν, σὺ δὲ\\
Ἀγάθωνά μοι δεῦρ᾽ ἐκκάλεσον πάσῃ τέχνῃ.
%\commentaire{65}{πάσῃ τέχνῃ.}{\eff{\italique{Bon gré, mal gré} (traduction de Willems.)}}
\\\passeLigne}%
\end{groupe}%
%% v. 66-69
\begin{groupe}%
\personnageGrec{Θεράπων}%
\trimIamb{μηδὲν ἱκέτευ᾽· αὐτὸς γὰρ ἔξεισιν τάχα.\\
καὶ γὰρ μελοποιεῖν ἄρχεται· χειμῶνος οὖν\\
ὄντος κατακάμπτειν τὰς στροφὰς οὐ ῥᾴδιον,
\commentaire{67-8}{χειμῶνος οὖν / ὄντος κατακάμπτειν τὰς στροφὰς οὐ ῥᾴδιον.}{Les Thesmophories avaient lieu du 11 au 13 du mois de Pyanepsion, jours qui tombent en général pendant notre seconde moitié du mois d'octobre.
%Jeu de mots sur στροφὰς qui, littéralement, sont les \italique{tours}.
Le serviteur continue de filer la métaphore avec la cire d'abeille qu'il est plus facile de modeler une fois exposée au Soleil.
}\\
ἢν μὴ προίῃ θύρασι πρὸς τὸν ἥλιον.\apparat{69}{\lecon{θύρασι} R : θύραζε Zanetti}\\\passeLigne}%
\end{groupe}%
%% v. 70a
\begin{groupe}%
\personnageGrec{Μνησίλοχος}%
\trimIamb{τί οὖν ἐγὼ δρῶ;%
\commentaire{70-4}{τί οὖν ἐγὼ δρῶ; ... ὄντα κηδεστὴν ἐμόν.}{La répartition des ce passage diffère d'un éditeur à un autre ; chacun voulant apporter son grain de sel. Inscrivons-nous, nous aussi, dans cette grande tradition philologique et apportons notre pierre à l'édifice : une nouvelle répartition des répliques qui, espérons-le, clôturera un débat qui n'a que trop duré. La répartition des deux répliques du vers 70 ne pose aucun problème – ou plutôt n'aurait dû poser aucun problème : il suffit de songer aux caractères des trois personnages en présence : le ton empoté du « τί οὖν ἐγὼ δρῶ; » ne peut convenir qu'au personnage pour l'instant passif et qui subit l'action : Mnésiloque\marqueur{1} (Euripide ruisselle de défauts mais garde toujours une certaine dignité ; jamais il ne s'abaisserait à lancer un mot aussi trivial et décontenancé – ses plaintes ont davantage de lyrisme !) La réponse : « περίμεν᾽, ὡς ἐξερχεται » ne peut être attribuée qu'à Euripide : c'est lui qui a besoin de Mnésiloque ; de plus, le serviteur ne délivre ni conseil ni ordre : il se contente de répondre le plus sèchement possible aux questions (cette répartition est d'ailleurs celle du manuscrit de Ravenne). Les vers 72 à 74 sont évidemment prononcés par Mnésiloque qui cherche à percer la détresse d'Euripide et nous révèle par la même occasion son lien de parenté avec le poète. Reste le vers 71 : il est plus que tentant de l'attribuer à Mnésiloque et faire ainsi des quatre vers 71-74 l'exact pendant des quatre vers 1-4 : une invocation à Zeus suivie d'une question : là pour savoir où Euripide le mène et ici pour savoir ce qui le tourmente. Problème : la présence d'un \italique{signe} (σημεῖον), probablement un \italique{paragraphos}, sur d'anciens manuscrits a été interprétée comme indiquant un changement d'interlocuteurs entre le vers 71 et les vers 72-74. Mais tel n'est pas le cas ici, et c'est ce que nous indiquait une ancienne scholie, très précieuse et incomprise\marqueur{2} : 
\citComm{τὸ σημεῖον· οὐ μόνον κηδεστὴς ὁ ἐκδιδούς, ἀλλὰ καὶ ὁ λαμβάνων.\\
\francais{Le signe [de ponctuation] : non seulement le Parent est celui qui abandonne [la parole], mais aussi celui qui la [re]prend.}}%
La scholie ne précise pas à quoi servait ce σημεῖον ; serait-ce une pause ? quoi qu'il en soit, le protagoniste devait certainement faire quelque jeu de théâtre après l'invocation à Zeus du vers 71, identique à celui déjà fait après le premier vers et montrer ainsi aux spectateurs le parallélisme entre les deux quatrains.}%
\debutNote{70-4}%
\note{1}{Réplique analogue de Mnésiloque au vers 925 : ἐγὼ δ᾽ ὁ κακοδαίμων τί δρῶ;}
\note{2}{Incomprise déjà de Suidas – lexicographe hypothétique du \nieme{IX} ou \nieme{X} siècle, à qui l'on attribue le lexique de la \italique{Souda}. Ce dernier fourmille en références et citations d'Aristophane : la \italique{Souda} est un outil de recherche sur Aristophane indispensable. 
Il en existe une version numérique : \href{https://www.cs.uky.edu/~raphael/sol/sol-html/}{https://www.cs.uky.edu/~raphael/sol/sol-html/} – incomprise donc de Suidas ce qui tend à prouver son ancienneté : Suidas (κ 1487) biffe le lemme `τὸ σημεῖον', pourtant indispensable à la bonne compréhension et égare ainsi les différents éditeurs des scholies qui l'appliquent au vers 74, voire qui la retouchent pour finir de la rendre tout-à-fait inintelligible.}\\\passeLigne}%
\end{groupe}%
%% v. 70b
\begin{groupe}%
\personnageGrec{Εὐριπίδης}%
\trimIamb{\ligne{70}\decalage{τί οὖν ἐγὼ δρῶ;}περίμεν᾽, ὡς ἐξερχεται.\\\passeLigne}%
\end{groupe}%
%% v. 71
\begin{groupe}%
\personnageGrec{Μνησίλοχος}\trimIamb{ὦ Ζεῦ τί δρᾶσαι διανοεῖ με τήμερον;\apparat{71-4}{(Μνη.) (71) et (Εὐρ.) (72-4) R : \lecon{(Μνη.) (71-4)} Bibi ap. Schol. : alii alia}%
\commentaire{71}{ὦ Ζεῦ τί δρᾶσαι διανοεῖ με τήμερον;}{Formule courante ; se reporter, par exemple, aux \italique{Oiseaux}, vv. 58 \& 62.}}%
\end{groupe}%
%% v. 72-74
\begin{groupe}%
\trimIamb{νὴ τοὺς θεοὺς ἐγὼ πυθέσθαι βούλομαι\\
τί τὸ πρᾶγμα τουτί. τί στένεις; τί δυσφορεῖς;\\
οὐ χρῆν σε κρύπτειν ὄντα κηδεστὴν ἐμόν.\testimonia{74}{Suidas κ 1487, cf. Moeris κ 38}\apparat{74}{\lecon{χρῆν} R, Suid. κ 1487 codd. FV : χρή Suid. codd. AGM \espApp \lecon{ἐμόν} R : ἐμέ Valckenaer ap. Dobree (p. 237)}\\\passeLigne}%
\end{groupe}%
%% v. 75
\begin{groupe}%
\personnageGrec{Εὐριπίδης}%
\trimIamb{ἔστιν κακόν μοι μέγα τι προπεφυραμένον.
\commentaire{75}{προπεφυραμένον}{Pour comprendre cette image, l'on se rapportera à un passage des \italique{Oiseaux} (v. 462-463) et le commentaire qu'en fait Taillardat (§755) :
\citComm{καὶ προπεφύραται λόγος εἷς μοι, / ὃν διαμάττειν οὐ κωλύει\\
\francais{j'ai sous la main un discours déjà bien délayé et rien n'empêche que j'achève de le pétrir.}}%
« Φυρᾶν signifie proprement \italique{délayer} de la farine dans un liquide\marqueur{1}. » Est donc προ-πεφυραμένον ce qui a déjà été délayé à l'avance mais point encore pétri ; autrement dit : le complot est ourdi mais rien n'est encore joué. L'avantage de ce genre de métaphore boulangère est qu'elle se transpose assez aisément en français\marqueur{2}.}%
\debutNote{75}%
\note{1}{Cf. Hippocrate \italique{Du régime} (\italique{Περί διαίτης}) 2,40 et Thucydide 3,49.}%
\note{2}{Comme l'épisode raconté par Démosthène dans les \italique{Cavaliers} (v. 55-57) : il s'est fait voler par Cléon les pains qu'il avait distribués... alors que c'est lui qui avait mis la main à la pâte ! Le verbe \italique{pétrir} et le \italique{pétrin} ont donné nombre d'expressions en français.}\\\passeLigne}%
\end{groupe}%
%% v. 76a
\begin{groupe}%
\personnageGrec{Μνησίλοχος}%
\trimIamb{ποῖόν τι;\\\passeLigne}%
\end{groupe}%
%% v. 76b-77
\begin{groupe}%
\personnageGrec{Εὐριπίδης}%
\trimIamb{\ligne{76}\decalage{ποῖόν τι;}τῇδε θἠμέρᾳ κριθήσεται\\
εἴτ᾽ ἔστ᾽ ἔτι ζῶν εἴτ᾽ ἀπόλωλ᾽ Εὐριπίδης.
\commentaire{76-7}{τῇδε θἠμέρᾳ κριθήσεται / εἴτ᾽ ἔστ᾽ ἔτι ζῶν εἴτ᾽ ἀπόλωλ᾽ Εὐριπίδης.}{Le style est paratragique (avec l'expression poétique `τῇδε θἠμέρᾳ' plutôt que la forme plus courante τήμερον). Apprécier surtout à sa juste valeur l'allitération ! Lyrisme ? emphase ? Les `ttttttttt' saccadent-ils le vers ou lui donnent-ils du rythme ? sans doute un peu des deux. (Euripide se tient souvent à la frontière ténue du sublime et du ridicule, frontière qu'il passe et repasse allègrement).}\\\passeLigne}%
\end{groupe}%
%% v. 78-80
\begin{groupe}%
\personnageGrec{Μνησίλοχος}%
\trimIamb{καὶ πῶς; ἐπεὶ νῦν γ᾽ οὔτε τὰ δικαστήρια\\
μέλλει δικάζειν οὔτε βουλῆς ἐσθ᾽ ἕδρα,\\
ἐπείπερ ἐστὶ Θεσμοφορίων ἡ μέση.\testimonia{80}{Hesychius τ 1440 τρίτη Θεσμοφορίων}\apparat{80}{ἐπεὶ τρίτη 'στὶ R : \lecon{ἐπείπερ ἐστὶ} Nauck (\italique{Mél.} (1866), II pp.468-71)}%
\commentaire{80}{ἐπεὶ τρίτη 'στὶ Θεσμοφορίων ἡ μέση.}{Le manuscrit de Ravenne donne au vers 80 : « puisque c'est le troisième [jour] des Thesmophories, celui du milieu. » Problème : si les Thesmophories duraient bien trois jours, celui du milieu ne peut en aucune façon être le troisième. Cette contradiction eut la fâcheuse conséquence de déchaîner des torrents de commentaires (à commencer par celui du scholiaste) ; on a inventé, en contradiction avec tous les témoignages\marqueur{1}, des Thesmophories de quatre jours à grand renfort d'érudition et de savantes notes, et ce sous le prétexte qu'il ait pu y avoir avant ces trois jours d'autres festivités ailleurs (surprenant, non ? mais il suffit de prendre un calendrier des fêtes grecques pour constater qu'avec ce genre de raisonnements, on pourrait tout aussi bien \italique{prouver} que les Thesmophories duraient 365 jours par an ; au reste – mais nous ne sommes plus à ce genre de détails près –, il faudrait encore montrer que le milieu d'un ensemble de quatre éléments en est le \nieme{3} ; pour finir – \italique{last but not least} –, faire préciser à Aristophane quel jour tombe ἡ μέση serait un détail inepte tant cela était connu de tous les Athéniens). Le texte avait été pourtant brillamment rétabli par August Nauck\marqueur{2} et intégré avec explications dès l'édition aux Belles Lettres de Victor Coulon ; cela n'empêche pas la faute d'apparaître dans nombre d'éditions modernes. Venons-en donc à la bonne correction : l'apparition du τρίτη, \italique{troisième}, est due à la mauvaise lecture d'un Π' (abréviation de περί ou περ) en Γ' (\italique{troisième} lettre de l'alphabet) dans un ancien manuscrit\marqueur{3} – cette abréviation est fréquente, et source d'erreurs aussi dans l'autre sens ; ainsi a-t-il fallu en maintes occasions rétablir un `τρίτη' à partir d'un gamma\marqueur{4}. Bref, voici l'histoire de notre début de \nieme{80} vers, très ingénieusement reconstituée : $$\text{ἐπείπερ ἐστὶ}\xrightarrow{\text{onciales}}\text{ΕΠΕΙΠ'ΕCΤΙ}\xrightarrow{\text{corruption}}\text{ΕΠΕΙΓ'ΕCΤΙ}\xrightarrow{\text{interprétation}}\text{ἐπεὶ τρίτη ᾽στὶ.}$$}%
\debutNote{80}%
\note{1}{Outre le scholiaste, qui \italique{sait} que les Thesmophories ne duraient que trois jours – sinon il ne prendrait pas la peine de tant ergoter, citons : Alciphron 2,37.2 ; 3,39 ; et la « \anglais{speculative restauration} » du papyrus n°3540.1-2, \italique{The Oxyrhynchus Papyri}, L (1983), ed. E. W. Handley}%
\note{2}{in \italique{Mélanges gréco-romains}, St.-Pétersburg (1866), t. II, pp. 468-71 ; voir aussi pp. 344-5 pour des témoignages sur les \italique{Thesmophories}.}%
\note{3}{La corruption doit être ancienne puisqu'elle – conjointe au bavardage du scholiaste – rendait déjà perplexe Hésychius (τ 1440): τρίτη Θεσμοφορίων· ζητεῖται πῶς ἅμα μὲν λέγει τρίτην Θεσμοφορίων εἶναι, ἅμα δὲ μέσην, τεττάρων οὐσῶν ἡμερῶν.}%
\note{4}{Pour rester avec notre cher Aristophane, un `τρεῖς' a dû être restitué au vers 1507 des \italique{Guêpes} là où tous les manuscrits comportaient un `γε'.}\\\passeLigne}%
\end{groupe}%
%% v. 81-84a
\begin{groupe}%
\personnageGrec{Εὐριπίδης}%
\trimIamb{τοῦτ᾽ αὐτὸ γὰρ τοι κἀπολεῖν με προσδοκῶ.\\
αἱ γὰρ γυναῖκες ἐπιβεβουλεύκασί μοι\\
κἀν Θεσμοφόροιν μέλλουσι περί μου τήμερον\apparat{83}{ἐμοῦ R : \lecon{μου} Biset}%
\commentaire{83}{κἀν Θεσμοφόροιν.}{Le scholiaste précise que c'est une Ἀττικὴ σύνταξις καὶ δυϊκὴ, ἀντὶ τοῦ `ἐν τῷ τῶν θεσμοφόρων ἱερῷ' ; les \italique{deux} déesses thesmophores étant Déméter et sa fille Perséphone.}\\
ἐκκλησιάζειν ἐπ᾽ ὀλέθρῳ.\\\passeLigne}%
\end{groupe}%
%% v. 84b
\begin{groupe}%
\personnageGrec{Μνησίλοχος}%
\trimIamb{\ligne{84}\decalage{ἐκκλησιάζειν ἐπ᾽ ὀλέθρῳ.}τιὴ τί δή;\\\passeLigne}%
\end{groupe}%
%% v. 85
\begin{groupe}%
\personnageGrec{Εὐριπίδης}%
\trimIamb{ὁτιὴ τραγῳδῶ καὶ κακῶς αὐτὰς λέγω.\\\passeLigne}%
\end{groupe}%
%% v. 86-87
\begin{groupe}%
\personnageGrec{Μνησίλοχος}%
\trimIamb{νὴ τὸν Ποσειδῶ καὶ δικαίαν ἂν πάθοις.\apparat{86}{δί᾽ και᾽ R : δίκαιά γ᾽ Grynaeus : \lecon{δικαίαν} Wilamowitz}%
\commentaire{86-7}{νὴ τὸν Ποσειδῶ καὶ δί᾽ και᾽ ἂν πάθοις / ἀτὰρ τίν᾽ ἐκ ταύτης σὺ μηχανὴν ἔχεις;}{Ne voyant pas à quel nom féminin renvoyait `ἐκ ταύτης', on a retouché le texte du vers 87 à l'envi (ἐκ τούτων Bergler, εἰς ταύτας Sandbach, \&c.). 
C'est à mon avis aller un peu vite en besogne alors qu'il y a justement un problème textuel au vers précédent : il manque une syllabe, probablement au niveau du fautif `δί᾽ και᾽'. Je fais l'hypothèse que c'est justement à cet endroit que se trouvait ce nom féminin, repris par le `ἐκ ταύτης' du vers suivant. De quoi Euripide peut-il se mettre \italique{hors} en déployant quelque machine ? Réponse : tel le Strepsiade des \italique{Nuées}, c'est à la justice qu'il veut échapper. Au vu du texte conservé, la correction est évidente : il faut substituer au `δί᾽ και᾽' un `δικαίαν' ; la chute de l'`αν' s'expliquant par la courante erreur de «~saut du même au même~» (`δικαίαν' étant directement suivi de la particule `ἂν'). 
Lisons donc :
\citComm{νὴ τὸν Ποσειδῶ καὶ δικαίαν ἂν πάθοις.\\
ἀτὰρ τίν᾽ ἐκ ταύτης σὺ μηχανὴν ἔχεις;\\
\francais{Sainte Mer ! ce serait en justice que tu pâtirais.\\
Mais pour y échapper, de quelle machine disposes-tu ?}}
Cette correction est d'autant plus intéressante que ἡ δικαία, mot rare, est un synonyme poétique de δίκη ; Mnésiloque retourne-t-il contre Euripide une de ses propres sentences\marqueur{1} ?\\
Wilamowitz\marqueur{2} avait déjà proposé cette émendation... sans qu'elle soit jamais reprise dans les éditions publiées depuis -- à une notable exception : celle de Prato.%encore une fois, on regrettera que les erreurs soient continuellement recopiées quand les meilleures restitutions sont superbement ignorées !.. manque de travail des différents éditeurs d'Aristophane ? ou manque de discernement ? Aristophane n'a pas eu la chance d'avoir beaucoup de commentateurs du calibre d'Alphonse Willems.
}%
\debutNote{86-7}%
\note{1}{Procédé qu'Aristophane s'est amusé à répéter tout au long de la pièce : Euripide est pris au piège de ses propres γνῶμαι aux vers 194, 275-276, etc. Quoi qu'il en soit, on trouve des expressions très proches aux vers 1379 des \italique{Nuées} : %
\citComm{Φειδιππίδης -- νὴ τὸν Δί᾽ ἐν δίκῃ γ᾽ ἄν,}%
ou 403 de \italique{Lysistrata} :%
\citComm{Πρόβουλος -- νὴ τὸν Ποσειδῶ τὸν ἁλυκὸν δίκαιά [δικαία ?] γε.}%
}%
\note{2}{Au commentaire du vers 681 de \italique{La Folie d'Héraclès} : second volume de son \italique{Euripides Herakles erklärt von Ulrich von Wilamowitz-Moellendorff}. Wilamowitz cite \italique{en passant} le vers 86 des \italique{Thesmophories} ainsi corrigé sans même préciser la correction -- dont il semble bien être l'auteur -- tant elle devait lui sembler aller de soi.}\\
\lednopb%
ἀτὰρ τίν᾽ ἐκ ταύτης σὺ μηχανὴν ἔχεις;
\commentaire{87}{μηχανὴν.}{Au sens imagé, une μηχανή est un \italique{expédient}, un \italique{tour}, un \italique{stratagème}. Certes, Euripide est un roué, un filou, un retors... mais Euripide est avant tout un homme de théâtre ! Il dispose de \italique{machineries} dont il use et abuse\marqueur{1}. Ce sont ces deux sens de μηχανή, le propre et le figuré, qui étaient instantanément saisis du spectateur du \nieme{V} siècle av. J.-C. Conformité du langage français avec le grec : le mot « truc », qui désignait au \nieme{XVII} siècle une machine de théâtre, s'est approprié le même sens imagé ; seul hic : le sens originel de \italique{machine théâtrale} s'est davantage estompé en français qu'en grec.}%
\debutNote{87}%
\note{1}{Euripide recourut à la μηχανή massivement dans ses œuvres tardives. Aristote, dans sa \italique{Poétique} (XV, 1454a37), le lui reproche pour sa \italique{Médée} : Euripide privilégia un dénouement ἀπὸ μηχανῆς plutôt que par l'histoire elle-même : φανερὸν οὖν ὅτι καὶ τὰς λύσεις τῶν μύθων ἐξ αὐτοῦ δεῖ τοῦ μύθου συμβαίνειν, καὶ μὴ ὥσπερ ἐν τῇ Μηδείᾳ ἀπὸ μηχανῆς. Aristophane parodie en maints endroits cet artifice de théâtre. Pour un dossier plus complet, aller voir l'article «~Machines et \italique{deus ex machina} dans les spectacles comiques~» de Leonardo Fiorentini, in \italique{L'appareil scénique dans les spectacles de l'Antiquité} (2013).}\\\passeLigne}%
\end{groupe}%
%% v. 88-89a
\begin{groupe}%
\personnageGrec{Εὐριπίδης}%
\trimIamb{Ἀγάθωνα πεῖσαι τὸν τραγῳδοδιδάσκαλον\\
ἐς Θεσμοφόροιν ἐλθεῖν.\\\passeLigne}%
\end{groupe}%
%% v. 89b
\begin{groupe}%
\personnageGrec{Μνησίλοχος}%
\trimIamb{\ligne{89}\decalage{ἐς Θεσμοφόροιν ἐλθεῖν.}τί δράσοντ᾽; εἰπέ μοι.\\\passeLigne}%
\end{groupe}%
%% v. 90-91a
\begin{groupe}%
\personnageGrec{Εὐριπίδης}%
\trimIamb{ἐκκλησιάσοντ᾽ ἐν ταῖς γυναιξὶ χἂν δέῃ\apparat{90}{γυναιξίν R : \lecon{γυναιξὶ} \Mon \espApp κἂν R : \lecon{χἂν} Markland (ad Eur. \italique{Suppl.} 364, (1763))}\\
λέξονθ᾽ ὑπὲρ ἐμοῦ.
\commentaire{90-1}{ἐκκλησιάσοντ᾽ ἐν ταῖς γυναιξὶν κἂν δέῃ / λέξονθ᾽ ὑπὲρ ἐμοῦ.}{Je reprends la correction de Markland : substituer χἂν ($=$ καὶ ἃ ἂν) à κἂν ($=$ καὶ ἐὰν) qui présente un sens bien meilleur : Euripide ne dit certes pas « qu'au besoin » ou « s'il le faut » (κἂν) Agathon interviendra mais bien plutôt qu'il lui faut intervenir autant qu'il le pourra (χἂν). Euripide a sa vie en jeu – si l'on peut s'exprimer ainsi – et sait bien le courroux de la gent féminine ; l'inaction aurait forcément une issue fatale... heureuse pour Athènes et les Athéniennes mais un peu plus fâcheuse pour sa personne. 
Précisons que ἐν a ici le sens technique que donne le Bailly (A, III.10) : « lorsqu'il s'agit d'un orateur, d'un accusé, ἐν correspond au français \italique{devant} » (cf. les \italique{Acharniens} v. 630 ; les \italique{Guêpes} v. 1185 ; le fragment 161 K.-A.) ; il faut donner un sens actif au `ἐκκλησιάσοντ᾽ ἐν ταῖς γυναιξὶν' : « pour intervenir à l'assemblée devant les femmes ».}\\\passeLigne}%
\end{groupe}%
%% v. 91b
\begin{groupe}%
\personnageGrec{Μνησίλοχος}%
\trimIamb{\ligne{91}\decalage{λέξονθ᾽ ὑπὲρ ἐμοῦ.}πότερα φανερὸν ἢ λάθρᾳ;\apparat{91}{\lecon{φανερὸν} R : φανερῶς Meineke}\\\passeLigne}%
\end{groupe}%
%% v. 92
\begin{groupe}%
\personnageGrec{Εὐριπίδης}%
\trimIamb{λάθρᾳ, στολὴν γυναικὸς ἠμφιεσμένον.\\\passeLigne}%
\end{groupe}%
%% v. 93-94
\begin{groupe}%
\personnageGrec{Μνησίλοχος}%
\trimIamb{τὸ πρᾶγμα κομψὸν καὶ σφόδρ᾽ ἐκ τοῦ σοῦ τρόπου·\testimonia{93-4}{Suidas τ 439}\testimonia{93}{Suidas κ 2025}\\
τοῦ γὰρ τεχνάζειν ἡμέτερος ὁ πυραμοῦς.\testimonia{94}{(ἡμέτερος ὁ πυραμοῦς) Suidas π 3192, cf. Tzetzes \italique{Chil.} XIII 281-288}%
\commentaire{94}{ἡμέτερος ὁ πυραμοῦς.}{L'expression proverbiale signifie : « à nous la palme de la victoire ! » En effet, on offrait un gâteau de miel et de froment (πυραμοῦς), à titre de récompense, à celui qui avait le mieux résisté au sommeil dans une παννυχίς (scholiaste des \italique{Cavaliers}, au v. 277 où l'expression apparaît à l'identique ; Athénée 647b-c, 668c). (Note de Taillardat, §847.)}\\\passeLigne}%
\end{groupe}%
%% v. 95a
\begin{groupe}%
\personnageGrec{Εὐριπίδης}%
\trimIamb{σίγα.\\\passeLigne}%
\end{groupe}%
%% v. 95b
\begin{groupe}%
\personnageGrec{Μνησίλοχος}%
\trimIamb{\ligne{95}\decalage{σίγα.}τί δ᾽ ἔστιν;\apparat{95}{\lecon{τί δ᾽} R : τί Kaibel}\\\passeLigne}%
\end{groupe}%
%% v. 95c
\begin{groupe}%
\personnageGrec{Εὐριπίδης}%
\trimIamb{\ligne{95}\decalage{σίγα.τί δ᾽ ἔστιν;}Ἁγάθων ἐξέρχεται.\\\passeLigne}%
\end{groupe}%
%% v. 96a
\begin{groupe}%
\personnageGrec{Μνησίλοχος}%
\trimIamb{καὶ ποῖός ἐστιν;\apparat{96}{ποῖο ἐστὶν οὗτος — οὑκκυκλούμενος R : ποῖος ἐστὶν οὗτος — οὑκκυκλούμενος \Mon : ποῦ ᾽σθ᾽; — ὅπου ᾽στίν; οὗτος οὑκκυκλούμενος Meineke : \lecon{ποῖος ἐστὶν; — οὗτος οὑκκυκλούμενος} Tyrwhitt (p. 65)}\\\passeLigne}%
\end{groupe}%
%% v. 96b
\begin{groupe}%
\personnageGrec{Εὐριπίδης}%
\trimIamb{\ligne{96}\decalage{καὶ ποῖός ἐστιν;}οὗτος οὑκκυκλούμενος.\testimonia{96}{(οὑκκυκλούμενος) Zonaras 678 ἐκκυκλεῖ}%
\commentaire{96}{οὑκκυκλούμενος}{($=$ ὁ ἐκκυκλούμενος). Agathon apparaît sur l'eccyclème, « machine de théâtre munie de roues qui, en tournant, faisait voir aux spectateurs ce qui était censé se passer dans une maison. » (Chantraine.) 
Le scholiaste l'avait noté :%
\citComm{ἐπὶ ἐκκυκλήματος γὰρ φαίνεται.}%
%\eff{Houston ! we have a textual problem !}
}\\\passeLigne}%
\end{groupe}%
%% v. 97-98
\begin{groupe}%
\personnageGrec{Μνησίλοχος}%
\trimIamb{ἀλλ᾽ ἦ τυφλὸς μέν εἰμ᾽· ἐγὼ γὰρ οὐχ ὁρῶ\apparat{97}{ἢ R : \lecon{ἦ} Biset}\\
ἄνδρ᾽ οὐδέν᾽ ἐνθάδ᾽ ὄντα, Κυρήνην δ᾽ ὁρῶ.
\commentaire{98}{Κυρήνην δ᾽ ὁρῶ.}{Cyrène était une célèbre courtisane évoquée aussi aux vers 1327-1328 des \italique{Grenouilles} : 
\citComm{ἀνὰ τὸ δωδεκαμήχανον / Κυρήνης μελοποιῶν\\
\francais{$[$Euripide$]$ compose ses chants à la mode de Cyrène aux douze postures.}}%
Le scholiaste y est plus loquace que le nôtre et nous apprend le trait distinctif de cette Cyrène :%
\citComm{Κυρήνη τις ἑταίρα ἐπίσημος, δωδεκαμήχανος ἐπικαλουμένη διὰ τὸ τοσαῦτα σχήματα ἐν τῇ συνουσίᾳ ποιεῖν.}%
Taillardat (§734) rapproche pertinemment ce sobriquet du titre d'un ``traité'' de Paxamos : Δωδεκάτεχνον, à propos duquel la \italique{Souda} (à l'entrée Πάξαμος) dit : ἔστι δὲ περὶ αἰσχρῶν σχημάτων. Que moque donc Aristophane en rapprochant, non cette fois la muse du poète comme dans les \italique{Grenouilles}, mais Agathon lui-même de Cyrène ? Bien-sûr son air efféminé ! bien-sûr son travestissement !.. mais pourquoi Cyrène en particulier ? Réponse : pour ce qu'elle a de remarquable (ἐπίσημος) qui lui valut son surnom ; il ne me paraît pas douteux qu'Agathon, qui vient de se faire rouler par l'eccyclème, apparaît affalé sur son lit dans une position suggestive, ou du moins qui peut prêter à la confusion de Mnésiloque. (Cf. la plaisanterie au vers 50.)}\\\passeLigne}%
\end{groupe}%
%% v. 99
\begin{groupe}%
\personnageGrec{Εὐριπίδης}%
\trimIamb{σίγα· μελῳδεῖν γὰρ παρασκευάζεται.\apparat{99}{ἂν R : δὴ Sykutris (\italique{utique} Divus) : οὖν Ellebodius : \lecon{γὰρ} Burger : ἢν Austin (1987)}\\\passeLigne}%
\end{groupe}%
%% v. 100
\begin{groupe}%
\personnageGrec{Μνησίλοχος}%
\trimIamb{Μύρμηκος ἀτραπούς, ἢ τί διαμινύρεται;\testimonia{100}{Hesychius μ 1904 Μύρμηκος ἀτραπούς ; Photius μ 606 Μύρμηκος ἀτραπός ; Suidas μ 1445 }\apparat{100}{μύρμηκος R : \lecon{Μύρμηκος} Bibi ap Hesych. μ 1904}%
\commentaire{100}{μύρμηκος ἀτραπούς.}{Willems, en accord avec le scholiaste, se propose d'expliquer ainsi l'expression μύρμηκος ἀτραπούς : « par allusion à la marche saccadée et tortueuse des fourmis. Image familière aux Anciens. Dans l'\italique{Anthologie}, XI, 78, il est question des galeries obliques et droites, λόξα καὶ ὀρθά, d'une fourmilière. Voir aussi Plutarque, \italique{de Solertia animal.}, II, p. 968a. Plaute a dit : \latin{formicinum gradum movere}. » Taillardat (§784) abonde en son sens avec d'autres exemples à l'appui. Rapprocher les airs contournés d'Agathon de sentiers tortus était certainement le but d'Aristophane.
Cependant, pour suggérer l'image de \italique{sentiers de fourmis}, n'aurait-il pas fallu, en français comme en grec, un pluriel (μυρμήκ\textbf{ων} ἀτραπούς) ?\\
Hésychius ne le comprenait pas autrement et voyait dans `μύρμηκος' un nom propre (μ 1904) :
\citComm{Μύρμηκος ἀτραπούς· Ἀθήνησιν ἐν Σκαμβωνιδῶν ἔστι Μύρμηκος ἀτραπός, ἀπὸ ἥρωος Μύρμηκος ὀνομαζομένη.\\
\francais{Dans [le dème] athénien des Scambônides [σκαμβός signifie \italique{courbe}, \italique{tortu}] se trouve un sentier de Myrmex, qui tire son nom du héros Myrmex.}}
Que l'image soit tirée de la fourmilière ou d'un nom propre, la critique n'en reste pas moins d'ordre musical\marqueur{2}.
%Il y a donc plus vraisemblablement une évocation d'un véritable lieu, connu des Athéniens, plutôt qu'une image animalière.\\
%Voici pour le sens littéral. Venons-en au côté facétieux du commentaire de Mnésiloque !\\
% L'image des lignes brisées de la fourmilière, μυρμηκιά, servit en musique pour décrire des trilles ou des arpèges comme l'atteste un fragment de Phérécrate\marqueur{2}.\\
C'est sûrement ce genre de comparaison qui valut à Philoxène de Cythère, un représentant de cette «~nouvelle musique~» dénoncée par Aristophane, le surnom de Μύρμηξ\marqueur{3}. Le persiflage de Mnésiloque ne serait-il pas une allusion audit Philoxène, alias Myrmex ?
%\citComm{Μύρμηκος ἀτραπούς\\
%C'est la voie de Myrmex}
}%
\debutNote{100}%
\note{1}{Relevé par nombre de parémiographes ainsi que Suidas à l'entrée Ἀγαθῶν ἀγαθίδες.}
\note{2}{L'image des lignes brisées de la fourmilière, μυρμηκιά, servit en musique pour décrire des trilles ou des arpèges. On trouve l'expression `ἐκτραπέλους μυρμηκιάς' chez le Pseudo-Plutarque, \italique{De Musica}, 30 ; il s'agit du fragment 155(v.23) de Phérécrate dans l'édition Kassel-Austin. À noter que le même Phérécrate a composé une comédie intitulée les μυρμηκάνθρωποι.}%
\note{3}{Voir l'article de la \italique{Souda} (φ 393) qui lui est consacré, à l'entrée Φιλόξενος.}\\\passeLigne}%
\end{groupe}%
%% v. 101-129
\begin{groupe}%
\personnageGrec{Ἀγάθων}%
\typeVers{ion. min\cat + ion. min.}\chantAgath{ἱερὰν Χθονίαιν δε-\apparat{101}{χθονίαις R : \lecon{χθονίαιν} Meineke}%
\commentaire{101-29}{}{Agathon chante un fragment de tragédie. Sujet : un coryphée et un chœur de jeunes filles troyennes célèbrent en strophes alternées Léto et ses deux rejetons. Il est très évident qu'Agathon chante seul. Les arguments fournis par M. P. Mazon\marqueur{1} sont à cet égard sans réplique. (Note de Willems.)\\
Ajoutons que Suidas (χ 403) nous confirme qu'Agathon avait l'habitude de chanter de la sorte : 
\citComm{Χορικός· ὁ τοῦ χοροῦ ἐξάρχων. τοιοῦτος ἦν Ἀγάθων, ὃς ὑποκριτικὰ μέλη ἐποίει καὶ αὐτὸς ὑποκρίνεται\\
\francais{Chorique : celui qui mène le chœur. Tel Agathon, qui avait l'habitude de créer des chants dramatiques et de les interpréter lui-même.}}%
Les divers commentateurs n'ont pas manqué de remarquer en quoi les ruptures de rythmes (παράρυθμ' εὔρυθμα) ou le contexte ionique de cette parodie d'Agathon en font un morceau efféminé et suave, ce qui est tout-à-fait juste\marqueur{2} ; seulement quand ils affirment que le caractère efféminé du chant est bien davantage dû au rythme et à la musicalité qu'aux paroles mêmes prononcées par Agathon\marqueur{3}, alors là il faut mettre le holà ! Et pour cela, je fais appel à feu Alain Blanchard et à son excellent article : \italique{La parodie d'Agathon et la structure des }``Bucoliques''\italique{ de Virgile}\marqueur{4}. Blanchard décèla dans ce chant une structure concentrique, en cela analogue aux \italique{Bucoliques} de Virgile ou encore aux cantates de Bach ; le point important étant que « la reconnaissance de ce schéma concentrique » est d'une grande « valeur exégétique » (pour reprendre les expressions de Blanchard). Laissons-lui la parole : « la coryphée invite le chœur à célébrer Apollon, Artémis et leur mère Léto, et le chœur, docilement, obéit à cette injonction. En réalité, l'étude de composition, une fois admise l'existence d'une structure concentrique conduit à plus de subtilité.\\
Le parallélisme entre le premier élément (vv. 101-103) et le neuvième (vv. 126-129) réside dans le thème de la lumière (lumière des torches portées par le chœur, lumière qui brille en réponse dans les yeux divins), thème qui précède un appel à la célébration d'un dieu – l'appel étant plus solennel, plus brusque et plutôt à titre de rappel dans le neuvième. Le parallélisme entre le deuxième élément (104-106) et le huitième (123-125) réside dans le thème du σέβεσθαι, introduit en mineur à la fin du deuxième élément et qui éclate sur le mode majeur au début du huitième.
Le parallélisme entre le troisième élément (107-110) et le septième (120-122) réside dans le caractère national des divinités honorées : caractère national sensible surtout dans le troisième élément (Apollon constructeur des remparts de Troie), mais présent, plus discrètement, dans le septième avec les Charites phrygiennes.
Enfin, le thème de la majesté divine (majesté d'Apollon comme dispensateur de la gloire) entre le quatrième élément (111-113) et le sixième (117-119). Le cinquième (114-116), central, est un appel à honorer Artémis, divinité féminine. Dans la première partie de la parodie, les trois interventions du coryphée concernent des divinités armées de l'arc qu'il faut honorer avec un esprit libre et d'une voix forte (103), mais le coryphée s'adresse à un chœur de jeunes filles qui, pour ainsi dire, retourne complètement la nature de ces divinités dont on aperçoit alors le caractère ambigu. Déjà, aux appels à un culte viril qu'on leur faisait, elles avaient répondu, avec une soumission apparente bien féminine, par une profession de piété plus neutre (σεβίσαι, 106). Puis, avec elles, Apollon cesse d'être le dieu archer pour être le coryphée des Muses (112). Qu'on leur propose d'honorer Artémis, il n'est plus question avec elle de chasse – occupation virile – mais de maternité : maternité glorieuse de Léto (118), maternité refusée par elle-même par la déesse Vierge (119). Le ton est donné pour la seconde partie du poème où le chœur, à l'appel du coryphée, honorera, après Artémis, sa mère Léto et la cithare. Dans un cadre aussi favorable, le chœur peut même se permettre une concession – nous parlerions plutôt d'un appel spontané et aventureux du cœur – avec l'expression ἄρσενι βοᾷ, \italique{mâle accent}, en 125 qui rappelle en la précisant l'invitation faite en 103b par le coryphée. Mais qu'on invite de nouveau à honorer une divinité masculine comme Apollon : c'est la maternité de Léto qui est chantée (130). »\\
Fera-t-on la même erreur à propos de ce chant que les premiers auditeurs du sonnet d'Oronte en prétendant : « Il est bon » ? C'est en tout cas ce qu'il nous semble à l'analyse de ce petit bijou de malice au subtil renversement : d'un banal hommage au fondateur d'une cité, l'on passe doucement et insidieusement à un morceau d'érotisme sur fond de maternité frustrée, le tout échafaudé sur une structure solidement ancrée\marqueur{5}. Rien d'étonnant donc à ce que Mnésiloque en soit tout émoustillé ; rien d'étonnant non plus à ce que, par la bouche de Dionysos, Aristophane rende hommage à Agathon après sa mort : ἀγαθὸς ποιητὴς καὶ ποθεινὸς τοῖς φίλοις, « un bon poète apprécié de ses amis » (\italique{Grenouilles}, v. 84). Pour finir, l'élaboration complexe du morceau ainsi mise au jour par Blanchard fait voler en éclat la thèse fort répandue d'une improvisation de la part d'Agathon ; d'ailleurs l'on sait par Aristote qu'Agathon avait eu l'ambition dans une de ses tétralogies de traiter intégralement du sac de Troie, ambition qui l'avait mené à l'échec : (\italique{Poétique}, XVIII, 56a16-19) :
\citComm{ὅσοι πέρσιν Ἰλίου ὅλην ἐποίησαν καὶ μὴ κατὰ μέρος ὥσπερ Εὐριπίδης, <ἢ> Νιόβην καὶ μὴ ὥσπερ Αἰσχύλος, ἢ ἐκπίπτουσιν ἢ κακῶς ἀγωνίζονται, ἐπεὶ καὶ Ἀγάθων ἐξέπεσεν ἐν τούτῳ μόνῳ.\\
\francais{Tous les poètes qui veulent traiter la ruine de Troie entièrement et non pas par parties comme Euripide, ou la légende de Niobé tout entière au lieu de faire comme Eschyle, tantôt échouent complètement, tantôt font mauvaise figure dans les concours. Ainsi, Agathon échoua sur ce seul point.}}%
Je me plais à imaginer Aristophane, révolté devant cet échec injuste, remettre sur scène ce qu'il considérait comme la perle (érotique) du cycle troyen d'Agathon (c'est en tout cas ce sens de la justice qui l'avait poussé à vouloir faire rejouer ses \italique{Nuées}, mésestimées par les mauvais jurés athéniens).}%
\debutNote{101-29}%
\note{1}{in \italique{Essai sur la composition des comédies d'Aristophane} (1904), p. 127-128.}%
\note{2}{À ce propos, cf. les \italique{Grenouilles} v. 1299 où Eschyle dit s'être écarté du pré des Muses de Phrynichos – Muses par trop efféminées et ionisantes ; de même, dans l'\italique{Assemblée des femmes}, quand une vieille veut alpaguer les beaux jeunes hommes qu'elle est en droit de s'approprier par force de loi, elle se lance alors dans une ariette au goût ionien (v. 883).}%
\note{3}{L'internationale du commentaire s'est donné le mot : « \allemand{Vor allem aber wird die Musik selbst durch Karikatur verspottet} » (Rau, p. 106) ; « \italien{più agli aspetti ritmico-musicali dell'inno che a quelli stilistici} » (Prato, p. 177) ; « \anglais{are due to the metre and the music rather than the words} » (Austin \& Olson p. 87) ; légèrement plus nuancé chez Rosella Saetta Cottone : «~une partie essentielle de la parodie devait reposer sur les modulations vocales par lesquelles l'acteur imitait la voix féminine~» ; \&c... en dépouillant minutieusement tous les glosateurs d'Aristophane, on devrait pouvoir traduire cette phrase dans une bonne partie des idiomes de notre babélienne Terre. Cette opinion au demeurant est bien pratique puisqu'elle dispense d'avoir à porter trop d'attention sur le texte.
Le style d'Agathon est certes efféminé mais beaucoup moins caricatural qu'on a voulu le faire croire. Cf. le fragment 178 (Γηρυτάδης) de l'édition Kassel-Austin : Aristophane y emploie l'adjectif Ἀγαθώνειος dans un sens loin d'être péjoratif d'après les commentaires anciens (Zénobius I, 2 ; Hésych. α 281 ; Phot. α 83 ; Suid. α 125 ; \&c.) :%
\citComm{Ἀγαθώνιος αὔλησις: ἡ μαλακὴ καὶ ἐκλελυμένη: ἢ ἡ μήτε χαλαρὰ, μήτε πικρὰ, ἀλλ' εὔκρατος καὶ ἡδίστη.\\
Un air de flûte agathonien : doux et détendu ; ni mol ni strident, mais tempéré et très doux.}%
}%
\note{4}{in \italique{Dans l'ouvroir du poète}, \italique{Hellenica} (2008), pp. 35-45. J'eus la chance de recevoir d'Alain Blanchard peu de temps avant sa mort de vifs encouragements, source d'énergie nécessaire à l'abattage de la somme de travail que représente ce commentaire des \italique{Thesmophories}. Quel regret que ce grand helléniste n'ait pas consacré son temps de recherche à Aristophane plutôt qu'à son «~cher Ménandre~».}%
\note{5}{Cf. pour plus de détails l'article d'A. Blanchard avec les chiffres à la clef.}%
\commentaire{101}{Χθονίαιν}{ est un génitif se rapportant à ἱεράν. Cf. \italique{Ploutos}, 937, et \italique{Caval.}, 301. (Note de Willems.)}}%
\end{groupe}%
%% v. 101-129
\begin{groupe}%
\typeVers{2 ion. min.}\chantAgath{ξάμεναι λαμπάδα κοῦραι}%
\end{groupe}%
%% v. 101-129
\begin{groupe}%
\decoupeVers{a}\typeVers{ion. min\cat + bacch.}\chantAgath{ξὺν ἐλευθέρᾳ πατρί-\apparat{103}{\lecon{πατρίδι} R : πραπίδι Wecklein}}%
\end{groupe}%
%% v. 101-129
\begin{groupe}%
\decoupeVers{b}\typeVers{ion. min. + ion. min\cat}\chantAgath{\ligne{102}\alinea δι χορεύσασθε βοάν.\\}%
\end{groupe}%
%% v. 101-129
\begin{groupe}%
\typeVers{2 ion. min. anaclast.}\chantAgath{\ligne{104}τίνι δαιμόνων ὁ κῶμος;}%
\end{groupe}%
%% v. 101-129
\begin{groupe}%
\typeVers{2 troch.}\chantAgath{λέγε νιν εὐπείστως δὲ τοὐμὸν\testimonia{104-6}{(εὐπίστως -- σεβίσαι) Suidas σ 183}\apparat{105}{νυν R : \lecon{νιν} Dindorf, Meineke \espApp εὐπίστως R, Suid. σ 178, 183 : \lecon{εὐπείστως} Reiske \ligneApp{106} ἔχεις R : \lecon{ἔχει} Suid. \ligneApp{107} ὅπλιζε R : \lecon{ὄλβιζε} Bentley \espApp \lecon{Μοῦσα} R : μούσᾳ Bergk}}%
\end{groupe}%
%% v. 101-129
\begin{groupe}%
\typeVers{troch. + ion. min\cat}\chantAgath{δαίμονας ἔχει σεβίσαι.\\}%
\end{groupe}%
%% v. 101-129
\begin{groupe}%
\typeVers{ion. min. + troch.}\chantAgath{\ligne{107}ἄγε νῦν ὄλβιζε Μοῦσα}%
\end{groupe}%
%% v. 101-129
\begin{groupe}%
\typeVers{troch. + ion. min.}\chantAgath{χρυσέων ῥύτορα τόξων}%
\end{groupe}%
%% v. 101-129
\begin{groupe}%
\typeVers{2 ion. min.}\chantAgath{Φοῖβον, ὃς ἱδρύσατο χώρας}%
\end{groupe}%
%% v. 101-129
\begin{groupe}%
\typeVers{ion. min. + ion. min\cat}\chantAgath{γύαλα Σιμουντίδι γᾷ.\\}%
\end{groupe}%
%% v. 101-129
\begin{groupe}%
\typeVers{2 troch.}\chantAgath{χαῖρε καλλίσταις ἀοιδαῖς}%
\end{groupe}%
%% v. 101-129
\begin{groupe}%
\typeVers{2 troch.}\chantAgath{Φοῖβ᾽ ἐν εὐμούσοισι τιμαῖς}%
\end{groupe}%
%% v. 101-129
\begin{groupe}%
\typeVers{ion. min. + ion. min\cat}\chantAgath{\ligne{112}γέρας ἱερὸν προφέρων.\\}%
\end{groupe}%
%% v. 101-129
\begin{groupe}%
\typeVers{2 troch.}\chantAgath{τάν τ᾽ ἐν ὄρεσι δρυογόνοισι}%
\end{groupe}%
%% v. 101-129
\begin{groupe}%
\typeVers{2 ion. min. + ion. min\cat}\chantAgath{κόραν ἀείσατ᾽ Ἄρτεμιν ἀγροτέραν.\testimonia{115-6}{(Ἄρτεμιν ἀγροτέραν) \Eustathe 361.37}\apparat{115}{ἀεισαντ᾽ R : \lecon{ἀείσατ᾽} Zanetti (\italique{laudate} Divus) \ligneApp{117} σεμνὸν R : \lecon{σεμνὰν} Dindorf}\\}%
\end{groupe}%
%% v. 101-129
\begin{groupe}%
\typeVers{ion. min. + troch.}\chantAgath{\ligne{116}ἕπομαι κλῄζουσα σεμνὰν}%
\end{groupe}%
%% v. 101-129
\begin{groupe}%
\typeVers{ion. min. + troch.}\chantAgath{γόνον ὀλβίζουσα Λατοῦς}%
\end{groupe}%
%% v. 101-129
\begin{groupe}%
\typeVers{ion. min. + ion. min\cat}\chantAgath{\ligne{118}Ἄρτεμιν ἀπειρολεχῆ.\\}%
\end{groupe}%
%% v. 101-129
\begin{groupe}%
\typeVers{io. mi. + io. mi\cat + io. mi.}\chantAgath{Λατώ τε κρούματά τ' Ἀσιάδος ποδὶ\testimonia{120}{\EtymMagnum 153.21 Ἀσιάδος κρούματα· τῆς κιθάρας. οὔτως Αριστοφάνης εἶπε, παρῳδῶν τὸ ἐξ Ἐρεχθέως Εὐριπίδου ; \SynB α 2226 $=$ Photius α 2957 $=$ Suidas 4147 $=$ \EtymGen (AB) α 1271 ed. Lasserre-Livadaras}}%
\end{groupe}%
%% v. 101-129
\begin{groupe}%
\typeVers{2 ion. min.}\chantAgath{παράρυθμ' εὔρυθμα Φρυγίων}%
\end{groupe}%
%% v. 101-129
\begin{groupe}%
\typeVers{ion. min. + ion. min\cat}\chantAgath{διὰ νεύματα Χαρίτων.\apparat{122}{διανεύματα R : δινεύματα Bentley : \lecon{διὰ νεύματα} Fritzsche}\\}%
\end{groupe}%
%% v. 101-129
\begin{groupe}%
\typeVers{ion. min. + troch.}\chantAgath{\ligne{123}σέβομαι Λατώ τ' ἄνασσαν}%
\end{groupe}%
%% v. 101-129
\begin{groupe}%
\typeVers{2 ion. min. anacl.}\chantAgath{κίθαρίν τε ματέρ᾽ ὕμνων}%
\end{groupe}%
%% v. 101-129
\begin{groupe}%
\typeVers{troch. + ion. min\cat}\chantAgath{ἄρσενι βοᾷ δοκιμῳ.\testimonia{125}{(ἄρσενι βοᾷ δοκιμῳ) Suidas α 4016 $=$ Scholium}\apparat{125}{\lecon{δοκιμῳ} R, Schol.\expo{1}, Suid. α 4016 : δόκιμον Schol.\expo{2} : δοκίμων Schöne \ligneApp{126} φῶς R : \lecon{φάος} Burges}\\}%
\end{groupe}%
%% v. 101-129
\begin{groupe}%
\typeVers{4 dact\cat}\chantAgath{\ligne{126}τᾷ φάος ἔσσυτο δαιμονίοις}%
\end{groupe}%
%% v. 101-129
\begin{groupe}%
\typeVers{6 dact. acatal.}\chantAgath{ὄμμασιν ἁμετέρας τε δι᾽ αἰφνι-\\
\alinea δίου ὀπός. ὧν χάριν\apparat{127}{ἡμετέρας R : ὑμετέρας Nietzsche : \lecon{ἁμετέρας} Meineke \ligneApp{128} \lecon{ἄγαλλε} R : ἀγάλλετε Austin (1987) \espApp Φοῖβον τιμᾷ R : \lecon{Φοῖβον} Dindorf}}%
\end{groupe}%
%% v. 101-129
\begin{groupe}%
\typeVers{2 iamb\cat}\chantAgath{\ligne{128}ἄνακτ᾽ ἄγαλλε Φοῖβον.\testimonia{128}{(ἄγαλλε) Photius α 164 $=$ \SynB α 145}\\}%
\end{groupe}%
%% v. 101-129
\begin{groupe}%
\typeVers{telesil. scazon}\chantAgath{\ligne{129}χαῖρ᾽ ὄλβιε παῖ Λατοῦς.\apparat{post v. 129}{parepigraphe ὀλολύζει R : ὀλολύζει ὁ γέρων Suid. ο 194}\\\passeLigne}%
\end{groupe}%
%% v. 130-145
\begin{groupe}%
\personnageGrec{Μνησίλοχος}%
\typeVers{trim. iamb.}\trimIamb{ὡς ἡδὺ τὸ μέλος ὦ πότνιαι Γενετυλλίδες\testimonia{130-2}{(ἡδὺ -- μανδαλωτόν) Suidas γ 141}\testimonia{130}{(Γενετυλλίδες) Hesychius γ 343 Γενετυλλίς}\apparat{130}{πότνια R : \lecon{πότνιαι} Suid. γ 141}%
\commentaire{130}{ὦ πότνιαι Γενετυλλίδες.}{Déesses de la génération, les Génétyllides sont aussi citées par Aristophane au vers 52 des \italique{Nuées} et au second vers de \italique{Lysistrata} ; de ce dernier passage on peut inférer que les femmes étaient réputées célébrer des rites licencieux en leur honneur.}\\
καὶ θηλυδριῶδες καὶ κατεγλωττισμένον\testimonia{131-2}{(θηλυδριῶδες -- μανδαλωτόν) Suidas μ 134}\testimonia{131}{(θηλυδριῶδες καὶ κατεγλωττισμένον) Suidas κ 912 \espTest Hesychius θ 491 θηλυδριῶδες \espTest (κατεγλωττισμένον) Pollux II.109}%
\commentaire{131}{θηλυδριῶδες.}{Semble être un néologisme d'Aristophane. Chantraine\marqueur{1} suppose un relais *θηλύδριον qui expliquerait aussi θηλυδρίας, bien attesté. Aristophane a-t-il porté son choix sur θηλυδριώδης pour faire entendre θηριώδης et toute l'animalité du spasme de Mnésiloque ? (cf. κατεγλωττισμένον et μανδαλωτόν.)}%
\debutNote{131}%
\note{1}{\italique{Dictionnaire étymologique de la langue grecque}, \italique{sub voce} θηλή.}\\
καὶ μανδαλωτόν, ὥστ᾽ ἐμοῦ γ᾽ ἀκροωμένου\testimonia{132}{Hesychius μ 224 μανδαλωτόν $\simeq$ Photius μ 86, cf. Photius γ 115}\testimonia{132-3}{(ὥστ᾽ -- γάργαλος) Suidas γ 83 (sed τοῦ μέλους post ἀκροωμένου addito)}%
\commentaire{131-2}{καὶ κατεγλωττισμένον / καὶ μανδαλωτόν}{: ces mots renvoient tous deux à ce que les anglophones appellent – ne boudons pas l'hommage qui nous est ainsi rendu ! – \italique{french kiss}. À propos de μανδαλωτόν, je ne résiste pas à citer Taillardat (§195) qui met au jour la bigarrure de la langue aristophanienne : « le substantif, τὸ φίλημα, \italique{le baiser}, est déjà connu d'Eschyle ; il se trouve chez Sophocle, chez Euripide, chez les comiques Platon et Nicophon. En revanche, il n'est pas attesté dans les œuvres conservées d'Aristophane, bien que le poëte emploie φιλεῖν comme synonyme de κυνεῖν [...] Aristophane préfère des mots imagés et comme ``techniques'' : τὸ περιπεταστόν (\italique{Acharniens} 1201), τὸ ἐπιμανδαλωτόν, \italique{ib.}, τὸ μανδαλωτόν (\italique{Thesm.} 132), les deux derniers termes évoquant l'idée d'une porte fermée ; ὁ μάνδαλος, comme βάλανος, désigne la clavette assurant la barre de la porte. La clavette, c'est ici la langue. »}\\
ὑπὸ τὴν ἕδραν αὐτὴν ὑπῆλθε γάργαλος.\testimonia{133}{(γάργαλος $=$ γαργαλισμός) Moeris γ 23, cf. \Phryn 56.9}%
%\commentaire{133}{ὑπὸ τὴν ἕδραν αὐτὴν ὑπῆλθε γάργαλος.}{Rabelais, \italique{Gargantua}, XIII :\citComm{\francais{[sa soye] me causoit au fondement une volupté bien grande.}}}
\\
καί σ᾽ ὦ νεανίσκ᾽ εἴ τις εἶ, κατ᾽ Αἰσχύλον\testimonia{134}{Hesychius η 948 ἥτις· ὅστις (?)}%
\apparat{134-5}{νεανίσχ᾽ εἴ τις R : \lecon{νεανίσκ᾽ εἴ τις} \Mon : νεανίσχ᾽ ἥτις Gelenius \espApp \textbf{135} Λυκουργίας R : \lecon{Λυκουργείας} Dindorf, Dobree}%
\commentaire{134}{νεανίσχ᾽ εἴ τις εἶ.}{Le manuscrit de Ravenne comporte une coquille : le κ de νεανίσκ᾽ est aspiré alors qu'il est suivi de εἴ. Plusieurs solutions ont été apportées : lire ἥτις (justifiant le féminin par Hésychius (η 948) glosant ἥτις· ὅστις). D'autres se sont contentés de substituer un κ au χ. J'ai une préférence pour cette dernière correction\marqueur{1} qui me semble être une pique contre Agathon : « jeune homme... si tu en es un ! » car Agathon serait âgé au moment de la représentation de 35 ans ou environ ; à croire que son mode de vie le conserve bien... mais tout-de-même ! cette allusion permanente à la jeunesse d'Agathon – comme dans le \italique{Banquet} de Platon censé se passer en 416, cinq ans plus tôt – me paraît plutôt être une allusion au fait qu'Agathon pèche par certain attribut pour être tout-à-fait considéré comme un homme\marqueur{2}.}%
\debutNote{134}%
\note{1}{Signalons qu'une coquille analogue (οὐχ᾽ εἴρηχ᾽ R en place d'un οὐκ εἴρηχ᾽) se retrouve au vers 556. De plus, les anciens manuscrits sur lesquels on a copié les \italique{Thesophories} n'étaient pas systématiquement élidés ; ce genre de fautes s'explique donc aisément. Cf. par exemples v. 148 ἐσθητ' ἅμα ; v. 279 θραττ' έπου ; v. 1178 ἔρχετ' ὡς. }%
\note{2}{καὶ ποῦ πέος; demande Mnésiloque au vers 142.}%
\\
ἐκ τῆς Λυκουργείας ἐρέσθαι βούλομαι.\\
\commentaire{134-5}{κατ᾽ Αἰσχύλον / ἐκ τῆς Λυκουργείας.}{D'après le scholiaste, la \italique{Lycurgie} est le nom collectif d'une tétralogie d'Eschyle, composée des \italique{Édoniens} (fr.\marqueur{1} 57-67), des \italique{Bassarides} (fr. 23-25), des \italique{Jouvenceaux} (Νεανίσκοι, fr. 146-9) et du drame satyrique \italique{Lycurgue} (fr. 124-6). Elle compte les mésaventures du roi Lycurgue essayant de contrer l'invasion du culte dionysiaque en Thrace. Laissons la parole à François Jouan\marqueur{2} : « La légende de Lycurgue était déjà connue d'Homère. L'\italique{Iliade} rappelait sommairement que Lycurgue avait été aveuglé par Zeus pour avoir pourchassé les Nourrices de Dionysos et forcé le petit dieu à se jeter à la mer\marqueur{3}. La version eschyléenne est très différente, et comporte en particulier une localisation précise : le héros est roi des Édoniens et vit en Thrace, dans la région du Pangée. On admet en général que cette tétralogie [...] datait de la fin de la vie d'Eschyle, entre 466 et 459. Or dans cette période, des forces athéniennes commandées par Cimon guerroyaient précisément contre les Édoniens, autour d'Eion et d'Ennéahodoi. [...] Les \italique{Édoniens} sont la pièce la moins mal connue, en particulier du fait de l'existence de fragments d'un \italique{Lycurgue} du poète latin Naevius, qui semble avoir suivi d'assez près le drame d'Eschyle. Dans le prologue, un messager annonçait au roi thrace l'arrivée d'un personnage entouré d'un bruyant cortège de musiciens et de danseurs – Dionysos et son thiase. Le roi ordonnait de les arrêter et de les faire comparaître devant lui. Dans la parodos, le chœur, sans doute composé de sujets de Lycurgue, décrivait le cortège dionysiaque. Un garde rendait compte de l'arrestation des trublions et de leur docilité. Le roi interrogeait lui-même Dionysos, comme le Penthée des \italique{Bacchantes}. Aussi mécontent des réponses de Dionysos que le prince thébain, il l'emprisonnait avec sa suite dans son palais, où il devait trouver, lui disait-il ironiquement «~un logement gratuit [ἄδασμος, fr. 63]~». Mais la défaite du dieu n'était que provisoire. Bientôt les liens des captifs tombaient d'eux-mêmes, Dionysos apparaissait dans sa gloire sur le toit du palais, qui vacillait sur sa base et se nimbait de flammes : autre motif inclus dans les \italique{Bacchantes}\marqueur{5}. » \\
C'est donc bien à cette première tragédie des \italique{Édoniens} que semble se référer Mnésiloque, interrogeant à la manière de Lycurgue, associant ainsi Agathon à Dionysos (métaphore qui avait commencée d'être filée au vers 41 quand le serviteur d'Agathon parlait du thiase de Muses réunies chez son maître). Rappelons que, pour les Grecs de l'époque, Dionysos était un dieu jeune : « Les particularités de son physique et de son costume sont détaillées aussi bien par Eschyle\marqueur{6} que par Euripide. Il est grand et mince avec de longues jambes, une chevelure blonde, abondante et parfumée, de belles boucles. Il est beau, avec un teint vermeil et des traits de visage plutôt féminins, l'air souriant. Il est vêtu d'une robe asiatique de fin tissu, tombant jusqu'à ses pieds et tient en main le thyrse. »\marqueur{7} 
(Ajoutons à propos de cette \italique{crocote} au cœur des péripéties de Mnésiloque qu'elle est aussi associée à Dionysos dans les \italique{Grenouilles} (v.46) ; par Cratinos, fr. 40 de l'édition de Kassel-Austin et, comme on l'a vu, Eschyle, fr. 59.) 
Voilà qui explique un peu les associations d'idées de Mnésiloque !}%
\debutNote{134-5}%
\note{1}{Les fragments d'Eschyle sont cités selon l'édition de S. Radt, \italique{Tragicorum Graecorum Fragmenta}, vol. 3 (1985).}%
\note{2}{On lira avec intérêt deux de ses articles : « Dionysos chez Eschyle » (in \italique{Kernos} 5, 1992, p. 71-86) et « Une introduction au mythe des \italique{Bacchantes} » (\italique{Kentron} 14, 1998). C'est le premier que nous citons. }%
\note{3}{\italique{Iliade}, VI, 130-140.}%
\note{4}{Eur. \italique{Bacch}., 586-599 ; 622-624 ; 633.}%
\note{5}{Cf. Hérodote, II, 145.}%
\note{6}{Voir les fragments 59, 60, 62 des \italique{Édoniens}.}%
\note{7}{F. Jouan, \italique{Une introduction au mythe des} Bacchantes, p. 12.}%
ποδαπὸς ὁ γύννις; τίς πάτρα; τίς ἡ στολή;\testimonia{136}{(ποδαπὸς -- πάτρα) Suidas γ 504 \espTest (γύννις) Hesychius γ 1015 γύννις, cf. \Phryn 60.19}\testimonia{136-40}{(τίς ἡ στολή -- κοινωνία) Suidas β 110 (δὲ pro δαὶ)}}
\end{groupe}%
%% v. 130-145
\begin{groupe}%
\trimIamb{τίς ἡ τάραξις τοῦ βίου; τί βάρβιτος\testimonia{137}{(βάρβιτος) Hesychius β 211 βάρβιτος, cf. Photius β 62}\\
λαλεῖ κροκωτῷ; τί δὲ λύρα κεκρυφάλῳ;\\
%\commentaire{138}{τί δὲ λύρα κεκρυφάλῳ;}{Mnésiloque questionne Agathon sur les contrastes de son accoutrement. Il relève au vers précédent l'étrange vis-àv-is du βάρβιτος et du κροκωτός, attributs aussi portés par le Dionysos d'Eschyle ; il est donc probable qu'Aristophane soit toujours en train de pasticher la \italique{Lycurgie} (encore que la tournure avec le verbe λαλεῖν fasse davantage penser au style d'Aristophane qu'à celui d'Eschyle).\\
%Le vers 138,τί δὲ λύρα κεκρυφάλῳ, semble une redite du précédent, puisque λύρα est un synonyme de βάρβιτος et que κεκρύφαλος est un autre des atours d'une toilette soignée. En fait l'apparente répétition permet à Aristophane de jouer avec un proverbe en le détournant. (procédé identique au vers 140.)\\
%Aulu-Gelle nous a conservé la version latine de ce \italique{vieil adage} :
%\citComm{\latin{Nihil graculo cum fidibus}\\
%\francais{Une lyre et un choucas, ça fait deux}}%
%\eff{Equites 1021 : τί γάρ ἐστ᾽ Ἐρεχθεῖ καὶ κολοιοῖς καὶ κυνί;}
%Laissons la parole à Érasme, à qui nous devons grâce à ses précieux \italique{Adages}\marqueur{1} cette découverte :
%\citComm{\latin{Nihil graculo cum fidibus, id est nihil stolidis et imperitis cum bonis literis. Nam graculus avis est ineptae molestaeque garrulitatis atque obstrepera, adeo ut a voce qua sonat, nomen inditum putent auctore Quintiliano primo Institutionum libro. Tum sui generis conventus amat, quo mutuo garritu strepitus fiat odiosior. Porro cithara silentium postulat et aures attentas. A. Gellius Noctium Atticarum extrema lucubratione sane quam eleganter hoc adagium torquet in quosdam pinguiore ingenio homines stolidaque loquacitate praeditos, profanos, ἀμούσος, id est, a musarum sacris alienos, qui politiores literas elegantioremque doctrinam vel ridere possunt vel etiam contemnere, intelligere nequaquam possunt, quos ea demum juvant, quae ipsi scilicet didicerunt, sordida, muta, indocta. Hi pari importunitate et alienae doctrinae obturbant et suam inscitiam passim venditant et inculcant neminem prae se contemnentes. Hoc igitur hominum genus Gellius a suis commentariis, quibus nihil fieri potest neque tersius neque eruditius, quasi fuste quodam abigens, Vetus est, inquit, adagium, nihil graculo cum fidibus.}\\
%\francais{Une lyre et un choucas, ça fait deux. C'est-à-dire rien de commun entre les imbéciles ou les têtes creuses et les belles-lettres. En effet, le choucas est un oiseau qui casse les oreilles par son croassement ridicule et désagréable ; on pense d'ailleurs qu'il tient son nom du cri qu'il fait entendre, voyez Quintilien dans le livre I [6, 37] de l'\italique{Institution oratoire}. En outre, il aime faire bande avec ceux de son espèce [cf. \italique{Adages}, 123] : leur échange de criailleries en devient d'autant plus insupportable. À l'opposé, la cithare exige du silence et des oreilles attentives. À la toute fin de sa préface des \italique{Nuits attiques} (19-21), Aulu-Gelle avec beaucoup d'à-propos tourne cet adage contre ces hommes à l'esprit particulièrement épais, dotés d'un creux bagout, prêts à tout salir, \italique{amousoï}, c'est-à-dire étrangers aux Muses [cf. \italique{Adages}, 1518], qui peuvent se gausser d'une culture particulièrement raffinée ou d'un vaste savoir, et même en faire fi. Ils n'y peuvent rien comprendre, et ne trouvent de charme qu'aux grossièretés, aux inepties sans nom qu'ils ont eux-mêmes apprises. Tout aussi fielleux pour déranger qui se cultive, ils font étalage en tout lieu de leur ignorance et l'imposent à toute force, méprisant le monde entier sauf eux-mêmes. C'est ce genre d'hommes qu'Aulu-Gelle dans ses commentaires — d'une élégance et d'une érudition hors pair — malmène à coups de bâton pour ainsi dire : « Il existe un vieil adage : rien de commun entre un choucas et une lyre. »}}
%Si les \latin{fidibus} latines correspondent à la λύρα grecque ; le \latin{graculo}, \italique{choucas}, pourrait se traduire en grec par le mot κολοιός ou bien κόκκυξ. C'est le mot κόκκυξ qui devait être utilisé pour le proverbe, comme on peut l'inférer de la glose d'Hésychius (κ 3298) suivante :  
%\citComm{κόκκυξ· τὸ ζῷον. ἢ τὸ ὄρνεον [ ... ] καὶ λόφος, καὶ \textbf{περικεφαλαία}.\\
%coucou : l'animal [un poisson]. Ou l'oiseau. Ou une crête. Ou un couvre-chef.}
%Le mot κόκκυξ pouvait avoir le sens de περικεφαλαία ; Aristophane use donc de κεκρυφάλῳ comme "permutant méthaphorique" -- ainsi parlait Taillardat -- pour κόκκυγι.\\
%Il résulte que le \italique{vieil adage} dont parle Aulu-Gelle devait se rencontrer chez les Grecs, peu ou prou sous la forme :
%\citComm{τί λύρα λαλεῖ κόκκυγι;}}%
%\debutNote{138}%
%\note{1}{Je cite les \italique{Adages} ainsi que leur traduction française selon l'édition de Jean-Christophe Saladin aux Belles Lettres (2019). Pour l'adage qui nous concerne, il s'agit du \nieme{337}.}\\
\ligne{139}τί λήκυθος καὶ στρόφιον; ὡς οὐ ξύμφορον.\apparat{139}{\lecon{οὐ ξύμφορον} R, Suid. β 190 cod. V : ἀξυμφορον Suid. cod. M : ὀξύμφορον Suid. codd. AGI : ξύμφορον Suid. cod. F \ligneApp{140} \lecon{κατόπτρου} R, Suid. β 110 : \lecon{κατρόπτου} Dover ad \italique{Nub.} 752 et \italique{Ran.} 47 \ligneApp{141} \lecon{σύ τ᾽} \PsiUn, schol. varia lectio : τίς δ᾽ R \ligneApp{142} \lecon{πέος} \PsiUn, R : τὸ πέος Suid. λ 63, π 987}%
\commentaire{139-44}{Fragments d'un ancien papyrus.}{L'on possède quelques fragments d'un ancien papyrus (\PsiUn \ dans notre apparat) du premier ou second siècle ap. J.-C. pour les vers 133-144 ; 145-156 ; 237-245 ; 273-288 ; 594-596 ; 804-809. Il nous apprend en outre le \latin{numerus versuum} d'une ancienne édition des \italique{thesmophories} : la pièce comportait alors 1201 vers (contre 1231 pour celle de Brunck sur laquelle se basent tous les éditeurs modernes). La papyrus est reproduit en annexe. Cf. \italique{Papiri della Società Italiana} XI, n° 1194 et  \italique{Papiri della Società Italiana} XIV, p. XV.}%
\commentaire{139}{τί λήκυθος καὶ στρόφιον;}{Dans sa liste des attributs masculins et féminins qui jurent ensemble chez Agathon, Mnésiloque oppose ici le \italique{lécythe} et le \italique{strophion}. Si le cas du strophion ne pose pas de problème : il s'agit du \latin{strophium} ou \latin{fascia pectoralis}\marqueur{1} : bandelette que les femmes enroulaient au niveau de leur poitrine pour la soutenir ; bref, un soutien-gorge ! En revanche, le lécythe a étonnamment égaré tous les commentateurs qu'il m'ait été donné de lire : ils lui ont unanimement donné le sens de \italique{fiole à huile}. D'aucuns ont même déployé force énergie pour tenter de montrer en quoi cet accessoire était typiquement masculin (!?) Ils eussent plus vite – et mieux ! – fait d'ouvrir n'importe quel lexique\marqueur{2} : ils auraient appris que λήκυθος peut aussi désigner la \italique{pomme d'Adam}. On comprend mieux l'effarement de Mnésiloque face à cette bandelette toute féminine entourant la poitrine désespérément plate d'Agathon surmontée d'une proéminente et peu ragoûtante pomme d'Adam ; admettez que cela est d'assez mauvais goût\marqueur{3}.}%
\debutNote{139}%
\note{1}{Je renvoie ici aux articles correspondants de la monumentale somme qu'est le \italique{Dictionnaire des antiquités grecques et romaines} de Daremberg et Saglio}%
\note{2}{Les dictionnaires classiques : Bailly, \LSJ, \&c. Hésychius (λ 857): τὸ μεταξὺ τοῦ λαυκανίου καὶ αὐχένος ήχῶδες. Cf. le verbe ληκυθίζειν = λαρυγγίζειν = φαρυγγίζειν, \italique{donner de la voix}.}%
\note{3}{Faut-il en conclure que la présence d'une pomme d'Adam était caractéristique du masque comique représentant Agathon ?}\\
τίς δαὶ κατρόπτου καὶ ξίφους κοινωνία;\\
\commentaire{140}{κατόπτρου.}{La variante κάτροπτον (avec métathèse du ρ) est la seule attestée par l'épigraphie pour les inscriptions contemporaines de la représentation des \italique{Thesmophories}\marqueur{1}. Ce n'est qu'à partir de la fin du \nieme{IV} que la forme κάτοπτρον, plus conforme à l'étymologie, semble prendre le dessus.}%
\debutNote{140}%
\note{1}{Voir par exemple les \italique{Inscriptiones Graecae} XXII.1471.47, 1544.58.
Le lecteur pourra aisément vérifier ces dires par lui-même en papillonnant sur les \italique{Inscriptiones Graecae} à l'adresse suivante : https://epigraphy.packhum.org/}%
\commentaire{140}{τίς δαὶ κατόπτρου καὶ ξίφους κοινωνία;}{Le sel de ce vers ne peut se goûter que si l'on a en tête le proverbe grec\marqueur{1} :
\citComm{τίς γὰρ κατόπτρῳ καὶ τυφλῷ κοινωνία;\\
Qu'y a-t-il de commun entre un aveugle et un miroir ?}
Ce proverbe a été répertorié par les parémiographes\marqueur{2} avec sa version concise, κάτοπτρον ὁ τυφλός, aux côtés d'associations tout aussi inutiles comme φαλακρὸς κτένα , \italique{un chauve un peigne}.\\
L'on trouvait bien évidemment des proverbes équivalents dans notre vieux parler ; \italique{À l'aveugle ne duit} [$=$ convient]\italique{ peinture / Couleur, miroir ne figure}\marqueur{3}.\\
Aristophane joue par paronomase entre ξίφους et τυφλοῦ.\\
%\eff{Pour l'épée, Agathon l'a-t-elle car elle fait partie de l'attirail de Dionysos ? mais alors, et le reste ?}
Au vers 138, il devait sans doute se livrer à un détournement du même genre : on recense de nombreux proverbes mettant en relation une lyre avec un âne (ex. : τί γὰρ κοινὸν ... λύρᾳ καὶ ὄνῳ, voir Tosi n° 1934), animal lourd et rétif aux muses ; ou une lyre avec un choucas (\latin{Nihil graculo cum fidibus}, \italique{Adages} d'Érasme n° 337), oiseau criard.
}%
\debutNote{140}%
\note{1}{\nieme{3} vers d'un fragment comique adespote cité par Stobée (IV, 30, 6a) ; c'est le \nieme{903} fragment des \italique{Poetae Comici Graeci, Adespota}, VIII (1995).}%
\note{2}{\italique{Corpus paroemiographorum Graecorum}, II, p. 774 (\italique{Mantissae Proverbiorum} 100).}%
\note{3}{\italique{Le livre des proverbes français} de Le Roux de Lincy, Série n°V, à l'entrée \italique{Aveugle}. À l'initiative du génial Pierre Boutang, le livre a été réédité en 1996 chez Hachette, augmenté d'une sienne préface dans laquelle -- entre autres merveilles -- Boutang questionne avec bonheur l'étymologie du \italique{pro-verbe}.}%
%\commentaire{140}{κατόπτρου καὶ ξίφους}{Sur un fameux cratère\marqueur{1} (le « cratère de Würzburg », Martin-von-Wagner-Museum der Universitàt Wurzburg, H 5697 ) représentant la scène d’égorgement de la ‘petite’ de Mika, on voit un miroir suspenduest-ce le miroir que l'on voit sur le cratère de Würzburg ? Si oui, il resterait à élucider pourquoi il est représenté lors de la scène beaucoup plus lointaine de la parodie du \italique{Télèphe}.\eff{Pour ce qui est du rapport des miroirs et des épées – pourquoi Agathon aurait-il une épée ? – la réponse se trouvait peut-être dans les pièces perdues d'Eschyle... ou plutôt d'Agathon. Sinon, alors voir la justification de Dover, Grenouilles 47 pour changer κατόπτρου en κατρόπτου et si c'est au prétexte comme je le soupçonne d'une pureté étymologique, pousser un coup de gueule !!}}%
σύ τ᾽ αὐτὸς ὦ παῖ; πότερον ὡς ἀνὴρ τρέφει;\testimonia{141-2}{(πότερον -- πέος) Suidas π 987}\\
καὶ ποῦ πέος; ποῦ χλαῖνα; ποῦ Λακωνικαί;\testimonia{142}{(ποῦ τὸ πέος -- Λακωνικαί) Suidas λ 63 \espTest Moeris χ 19 χλαῖνα \espTest (Λακωνικαί) Hesychius λ 225 Λακωνικαί ; Photius λ 57 Λακωνικαί ; scholium \italique{Vespae} 1158-60}\\
ἀλλ᾽ ὡς γυνὴ δῆτ᾽ εἶτα ποῦ τὰ τιτθία;\\
τί φῄς; τί σιγᾷς; ἀλλὰ δῆτ᾽ ἐκ τοῦ μέλους\\%
%\commentaire{144}{ἐκ τοῦ μέλους}{« d'après ton chant », serait-on tenté de traduire ; mais le μέλος (litt. : \italique{membre}, \italique{partie}) peut désigner aussi une partie du corps. Y a-t-il un jeu de mots ?}
ζητῶ σ᾽, ἐπειδή γ᾽ αὐτὸς οὐ βούλει φράσαι;\\\passeLigne}%
\end{groupe}%
%% v. 146-152
\begin{groupe}%
\ledpb%
\personnageGrec{Ἀγάθων}%
\trimIamb{ὦ πρέσβυ πρέσβυ, τοῦ φθόνου μὲν τὸν ψόγον\testimonia{146-7}{Suidas α 1097 (μὲν omisso et δὲ pro δ')}\\
ἤκουσα, τὴν δ᾽ ἄλγησιν οὐ παρεσχόμην‎\\
ἐγὼ δὲ τὴν ἐσθῆθ᾽ ἅμα τῇ γνώμῃ φορῶ.\apparat{148}{γνώμη R : \lecon{τῇ γνώμῃ} Meineke}\\
χρὴ γὰρ ποιητὴν ἄνδρα πρὸς τὰ δράματα\testimonia{149-50}{\Apostol XVIII 41h}\\
ἃ δεῖ ποιεῖν πρὸς ταῦτα τοὺς τρόπους ἔχειν.\\
αὐτίκα γυναικεῖ᾽ ἢν ποιῇ τις δράματα,\\
μετουσίαν δεῖ τῶν τρόπων τὸ σῶμ᾽ ἔχειν.\\\passeLigne}%
\end{groupe}%
%% v. 153
\begin{groupe}%
\personnageGrec{Μνησίλοχος}%
\trimIamb{οὐκοῦν κελητίζεις, ὅταν Φαίδραν ποιῇς;
\commentaire{153}{κελητίζεις, ὅταν Φαίδραν ποιῇς.}{Allusion à la \italique{pendula Venus} des Latins. Cf. \italique{Guêpes}, 500 ; \italique{Lysistrata}, 60 et 773. Le poète Machon disait καθιππάζειν (Athénée, XIII, p. 581d). (Note d'Alphonse Willems. Phèdre était éprise d'Hippolyte, dont le nom signifie littéralement le «~délieur de chevaux~» ; c'est-à-dire celui dont les chevaux courent à bride abattue.)}\\\passeLigne}%
\end{groupe}%
%% v. 154-156
\begin{groupe}%
\personnageGrec{Ἀγάθων}%
\trimIamb{ἀνδρεῖα δ᾽ ἄν ποιῇ τις, ἐν τῷ σώματι\apparat{154}{\lecon{ἄν} \PsiUn, R : ἢν Dindorf}\\
ἔνεσθ᾽ ὑπάρχον τοῦθ᾽. ἃ δ᾽ οὐ κεκτήμεθα,\\
μίμησις ἤδη ταῦτα συνθηρεύεται.\\\passeLigne}%
\end{groupe}%
%% v. 157-158
\begin{groupe}%
\personnageGrec{Μνησίλοχος}%
\trimIamb{ὅταν σατύρους τοίνυν ποιῇς, καλεῖν ἐμέ,
\commentaire{157}{ὅταν σατύρους τοίνυν ποιῇς.}{Lors des Grandes Dionysies, chaque poète tragique en lice devait présenter une tétralogie composée de trois tragédies et d'un drame satyrique.}\\
ἵνα συμποιῶ σοὔπισθεν ἐστυκὼς ἐγώ.
\commentaire{158}{συμποιῶ σοὔπισθεν ἐστυκὼς}{ ; ἐστυκὼς de στύω, \italique{bander}, alors qu'on attendait ἑστηκώς, \italique{se tenir debout} ; σοὔπισθεν est mis ici pour σοί ὄπισθέν (le datif σοί réclamé par le συμ- de συμποιῶ).}\\\passeLigne}%
\end{groupe}%
%% v. 159-167
\begin{groupe}%
\ledpb%
\personnageGrec{Ἀγάθων}%
\trimIamb{ἄλλως τ᾽ ἄμουσόν ἐστι ποιητὴν ἰδεῖν\testimonia{159-61}{Suidas α 1633 (κ' Ἀνακρ.)}\\
ἀγρεῖον ὄντα καὶ δασύν σκέψαι δ᾽ ὅτι\\
Ἴβυκος ἐκεῖνος κἀνακρέων ὁ Τήιος\testimonia{161-2}{Suidas ε 4049 (Ἀνακρέων, Ἀλκαῖος, οἵπερ ἁρμονίαν ἐχύμισαν)}%
\commentaire{161-3}{Ἴβυκος ἐκεῖνος κἀνακρέων ὁ Τήιος / κἀλκαῖος, οἳ περὶ ἁρμονίαν ἐχύμησαν / ἐμιτροφόρουν τε καὶ διεκίνων Ἰωνικῶς,}{Les trois poètes lyriques Ibycos de Rhégium, Anacréon de Téos et Alcée de Lesbos sont aussi rapprochés par le scholiaste de Pindare (\italique{Isthmiques}, 2.1b\marqueur{1}) pour d'autres raisons que d'ordre harmonique : ce sont trois auteurs pédérastiques ; Cicéron les regroupe aussi sur le critère de la passion amoureuse (\italique{Tusculanes}, IV, 33) :
\citComm{\latin{quid denique homines doctissimi et summi poëtae de se ipsis et carminibus edant et cantibus ? Fortis vir in sua re publica cognitus quae de iuvenum amore scribit Alcaeus ! nam Anacreontis quidem tota poësis est amatoria. Maxume vero omnium flagrasse amore Reginum Ibycum apparet ex scriptis.}\\%
\francais{Avec quelle licence de très savants hommes et de grands poètes n'ont-ils pas chanté leurs galanteries ? Alcée, ce fameux guerrier de Mitylène, que n'a-t-il pas écrit de ces inclinations à la Grecque ? Anacréon respire-t-il autre chose que l'amour ? On voit la passion encore poussée bien plus loin dans les poésies d'Ibycus.}}%
Le spectateur du \nieme{V} siècle ne s'y trompait pas et savait à quoi s'en tenir avec ces trois bougres-là ; il pouvait ainsi interpréter correctement les euphémismes agathoniens.\\
Voici pour le rapprochement au point de vue littéraire. Maintenant, qu'est-ce qui permet de les rapprocher quant à leur mode de vie ? C'est au vers 163 qu'Agathon s'en explique.%
\citComm{ἐμιτροφόρουν τε καὶ διεκίνων Ἰωνικῶς}%
%\eff{Érasme 469}
Ce vers ne respecte pas le mètre iambique. Si l'on conserve l'adverbe Ἰωνικῶς -- d'ailleurs paraphrasé ἁβρῶς par le scholiaste\marqueur{2} --, l'erreur se trouve donc au niveau de †διεκίνων† (διεκίνουν dans la \italique{Souda}).\\
Comme toujours en présence d'erreurs textuelles chez Aristophane, je suis allé chercher dans Hésychius une possible solution. (J'espère démontrer de façon éclatante au fil de ce commentaire des \italique{Thesmophories} tout l'intérêt philologique -- encore largement inexploité -- du lexique d'Hésychius pour Aristophane. Hésychius a très largement repris le commentaire de Didyme pour l'intégrer à son lexique\marqueur{3}. Il est d'un secours d'autant plus indispensable dans le cas des \italique{Thesmophories} que fort peu de scholies nous sont parvenues via l'unicodex de Ravenne.)
Revenons à notre \nieme{163} vers :
\citComm{ἐμιτροφόρουν τε καὶ ........... Ἰωνικῶς}
Pour démontrer qu'un poète se doit de soigner son apparence, Agathon fait appel aux trois poètes Ibycos, Anacréon et Alcée qui portaient tous trois la \italique{mitrê}\marqueur{4} et ........... à l'ionienne.\\
Comme il est peu probable que le spectateur moyen d'Aristophane -- θεατής ὁδοῦ, dira-t-on -- ait eu à l'esprit des passages de chacun de ces poètes où ils décrivaient leur manière de s'apprêter ; c'est bien plutôt à l'iconographie qu'Agathon fait appel.\\
Faisons donc appel à l'archéologie ! Un article d'un certain Papaspyridi-Karouzou\marqueur{5} fait justement état des points communs entre les représentations sur vases d'Anacréon et d'Alcée : «~Les deux poètes lyriques ont en commun, comme traits caractéristiques, le chiton talaire, l'himation et la lyre, et, en outre, la \italique{mitrè}~». (Anacréon se distinguant par sa calvitie.) Outre la μίτρα -- et bien-sûr la lyre ! --, c'est donc à leur vêtement particulier que ces trois poètes se reconnaissent aux Athéniens du \nieme{V} siècle\marqueur{6}.\\
Passons à Hésychius\marqueur{7} (δ 1805) qui connaît un mot pour qualifier une telle tenue : 
\citComm{δίκανα· ποικίλα ἱμάτια\\
Dikana : habits bigarrés}
La substitution du δί- de δίκανον par le δια- (ou διε-) de διεκίνων peut aisément s'expliquer : ces préfixes s'abrégeant sur d'anciens manuscrits par la simple lettre Δ surmonté de points ; la confusion est donc fréquente\marqueur{8}.\\
Lisons donc :%
\citComm{ἐμιτροφόρουν τε καὶ δίκανον Ἰωνικῶς\\
Ils portaient la mitre et le dikanon à l'ionienne.}
}%
\debutNote{161-3}%
\note{1}{ταῦτα δὲ τείνει καὶ εἰς τοὺς περὶ Ἀλκαῖον καὶ Ἴβυκον καὶ Ἀνακρέοντα, καὶ εἴ τινες τῶν πρὸ αὐτοῦ δοκοῦσι περὶ τὰ παιδικὰ ἠσχολῆσθαι· οὗτοι γὰρ παλαιότεροι Πινδάρου· Ἀνακρέοντα γοῦν ἐρωτηθέντα, φασί, διατί οὐκ εἰς θεοὺς ἀλλ εἰς παῖδας γράφεις τοὺς ὕμνους; εἰπεῖν, ὅτι οὗτοι ἡμῶν θεοί εἰσιν.}%
\note{2}{Voir Érasme, \italique{Adages} (n° 469, \latin{Risus Ionicus}).}%
\note{3}{Le lecteur pourra très avantageusement se reporter au livre de Pierre Boudreaux, \italique{Le texte d'Aristophane et ses commentateurs} (1920).}%
\note{4}{La μίτρα est ici un bandeau servant de coiffure. (Bailly.)}%
\note{5}{Semni Papaspyridi-Karouzou, «~Anacréon à Athènes~». In \italique{Bulletin de correspondance hellénique} (1942), vol. 66-67, pp. 248-254.}%
\note{6}{C'est avec ces mêmes caractéristiques qu'est représenté Ibycus sur un monument à Reggio de Calabre.}%
\note{7}{L'\latin{Etymologicum magnum} donne une définition plus précise à l'entrée δείκανα :
\citComm{Δείκανα : Τὰ ποικίλα ὑφάσματα καὶ μορφὰς ἔχοντα.\\
Dikana : tissus chamarrés brodés de figures (?) [sic ; lire φορβειάς ?\marqueur{9}].}
Ce complément qu'on trouve dans la glose de l'\latin{Etymologicum magnum} est aussi présent chez Hésychius... mais ailleurs : à l'entrée δεικές (δ 450) et voisine de *δείκελον ($=$ δείκηλον) ; cette proximité est sans doute à l'origine de la confusion des deux termes dans l'\latin{Etymologicum magnum} -- car pour qui a un peu fréquenté ces deux lexiques, il est évident qu'il faut donner la primeur à Hésychius.
Si, comme je le crois vraisemblable, le texte original présentait ce terme, c'est bien avec la forme à vocalisme zéro d'Hésychius : δίκανον, et ce pour des raisons métriques.}%
\note{8}{Cf. l'article de Victor Coulon, «~Notes critiques et exégétiques sur Aristote, \italique{Poétique} , 1450a12, 1450a29, sur Thucydide I, 1, 2 et sur cinq passages d'Aristophane~», in \italique{Revue des Études Grecques} (1937), t. 50 fasc. 238, p. 458. (À propos du vers 122 et de la correction de Bentley δινεύματα à la place de διανεύματα R.)}%
\note{9}{Le manuscrit d'Hésychius -- encore un unicodex !.. -- porte φόρβας en place de μορφάς.}%
\\
κἀλκαῖος, οἳ περὶ ἁρμονίαν ἐχύμισαν,\testimonia{162-3}{(κ' Ἀλκαῖος -- Ἰωνικῶς) Suidas ε 989 $\simeq$ Zonaras 707 ἐμιτροφόρουν}\testimonia{162}{Zonaras 942 οἵπερ ἁρμονίαν ἐχύμισαν \espTest (ἐχύμισαν) Hesychius ε 7681 ἐχύμισαν}\apparat{162}{\lecon{κἀλκαῖος} R, Suid. : (κ)ἀχαιός Schol. \espApp \lecon{οἳ περὶ ἁρμονίαν} R : οἵπερ ἁρμονίαν Suid. ε 989 \espApp ἐχύμησαν R : \lecon{ἐχύμισαν} Suid. \ligneApp{163} διεκίνων R : διεκίνουν Suid. ε 989, ι 495 : διεκλῶντ᾽ Toup : \lecon{δίκανον} Bibi ap. Hesych. δ 1805 \ligneApp{165} \lecon{ἠμπέσχετο} R : ἠμπισχετο Elmsley ad Eur. \italique{Med.} 1128 (\italique{induebatur} Divus)}%
\commentaire{162}{οἳ περὶ ἁρμονίαν ἐχύμησαν.}{Certains éditeurs préfèrent suivre la leçon de la \italique{Souda} (ε 989) : \lecon{οἵπερ} ἁρμονίαν ἐχύμισαν. Les arguments de Victor Coulon\marqueur{1} me semblent l'emporter : « gardons-nous de suivre le lexicographe [Suidas], en remplaçant la leçon d[u] R[avennas] οἳ περὶ par le pronom οἵπερ formé de l'adjectif relatif οἵ et de la particule d'identité περ. Ici ce pronom n'est pas de mise, mais le verbe ἐχύμισαν a besoin d'un adverbe augmentant sa valeur. Le περί\marqueur{2} est donc à prendre au sens de περισσώς : \italique{au delà de la mesure}, \italique{supérieurement}, \italique{plus que}, comme par exemple dans Homère, \italique{Od}. I, 66-67 :
\citComm{ὅς περὶ μεν νόον ἐστὶ βροτῶν, περί δ᾽ ἱρὰ θεοῖσιν\\
ἀθανατοῖσιν ἔδωκε,}%
où de περὶ ἐστὶ dépend βροτῶν et le second περί se rapporte à ἔδωκε. Dans notre passage περὶ χυμίζειν signifie donc : \italique{rendre savoureux au-delà de la mesure} ou \italique{supérieurement savoureux}.
La faute du texte de Suidas est la même que celle de tous les manuscrits dans les \italique{Cavaliers}, v. 727 : οἷαπερ ὑβρίζομαι corrigée avec bonheur par Elmsley en οἷα περιυβρίζομαι. »}%
\debutNote{162}%
\note{1}{\italique{Notes critiques et exégétiques sur divers passages controversés d'Aristophane et sur Sophocle, Antigone 613-614}, in \italique{Revue des Études Grecques}, tome 66, fascicule 309-310, Janvier-juin 1953. p. 45-46}%
\note{2}{Cf. aussi l'adverbe ἀμφί dans \italique{Thesm}. 995, mais non pas dans \italique{Nuées} 595, comme le croit Todd, \italique{Index Aristophaneus}, s. v. ἀμφί. Là il faut sous-entendre avec ἀμφί μοι l'impératif épique ἔσσο.}\\
ἐμιτροφόρουν τε καὶ δίκανον Ἰωνικῶς,\testimonia{163}{Suidas ι 495 ἐμιτροφόρουν τε καὶ διεκίνουν Ἰωνικῶς}\\
καὶ Φρύνιχος, τοῦτον γὰρ οὖν ἀκήκοας,
\commentaire{164}{Φρύνιχος, τοῦτον γὰρ οὖν ἀκήκοας.}{Le poète Phrynicos, qui vivait au temps des guerres médiques portait naturellement le riche costume de son temps. Cf. les \italique{Cavaliers}, v. 1331-1332 :
\citComm{Ἀλλαντοπώλης — ὅδ᾽ ἐκεῖνος ὁρᾶν τεττιγοφόρας, ἀρχαίῳ σχήματι λαμπρός,\\
οὐ χοιρινῶν ὄζων ἀλλὰ σπονδῶν, σμύρνῃ κατάλειπτος.\\
\francais{Le charcutier — Le voilà qui apparait, la cigale à la chevelure, dans tout l'éclat du costume antique, fleurant, non les coquilles mais les libations de paix, et frotté avec de la myrrhe.}}%
Et Thucydide (I, 6) :%
\citComm{καὶ οἱ πρεσβύτεροι αὐτοῖς τῶν εὐδαιμόνων \{διὰ τὸ ἁβροδίαιτον\} οὐ πολὺς χρόνος ἐπειδὴ χιτῶνάς τε λινοῦς ἐπαύσαντο φοροῦντες καὶ χρυσῶν τεττίγων ἐνέρσει κρωβύλον ἀναδούμενοι τῶν ἐν τῇ κεφαλῇ τριχῶν.\\
Les vieillards athéniens de la classe aisée portaient encore, il n'y a pas longtemps, des tuniques (longues) de lin, et piquaient des cigales d'or dans leurs cheveux, ramassés et noués sur le sommet de la tête.}%
(Note de Willems.)\\
Le scholiaste explique l'incise d'Agathon (τοῦτον γὰρ οὖν ἀκήκοας) ainsi : ἐπεὶ Ἀθηναῖος ἦν.}\\
αὐτός τε καλὸς ἦν καὶ καλῶς ἠμπέσχετο\\%\apparat{165}{\lecon{ἠμπέσχετο} R : ἠμπισχετο Elmsley ad Eur. \italique{Med.} 1128 (\italique{induebatur} Divus)}\\
διὰ τοῦτ᾽ ἄρ᾽ αὐτοῦ καὶ κάλ᾽ ἦν τὰ δράματα.\\
ὅμοια γὰρ ποιεῖν ἀνάγκη τῇ φύσει.\\\passeLigne}%
\end{groupe}%
%% v. 168-170
\begin{groupe}%
\personnageGrec{Μνησίλοχος}%
\trimIamb{ταῦτ᾽ ἄρ᾽ ὁ Φιλοκλέης αἰσχρὸς ὢν αἰσχρῶς ποιεῖ,\testimonia{168}{(ὁ Φιλοκλέης αἰσχρὸς $=$ κορυδός) scholium \italique{Aves} 1295a.α}\apparat{168-9}{Φιλοκλῆς, Ξενοκλῆς R : \lecon{Φιλοκλέης}, \lecon{Ξενοκλέης} Bentley \espApp \textbf{169} \lecon{ὁ δ᾽ αὖ} R : ὁ δὲ Brunck}\\
\lednopb%
ὁ δ᾽ αὖ Ξενοκλέης ὢν κακὸς κακῶς ποιεῖ,\\
\lednopb%
ὁ δ᾽ αὖ Θέογνις ψυχρὸς ὢν ψυχρῶς ποιεῖ.
\commentaire{168-70}{Φιλοκλῆς, Ξενοκλῆς, Θέογνις.}{Trois poètereaux tragiques contemporains d'Aristophane.\\
À propos du vilain Philoclès, deux scholies nous ont été transmises dans le manuscrit de Ravenne\marqueur{1} :%
\citComm{Φιλοκλέης: τραγῳδίας ποιητὴς, αἰσχρὸς τὸ πρόσωπον καὶ σιμὸς καὶ μικρὸς τὸ σῶμα\marqueur{2}.\\
ἐκ τούτου δηλονότι καὶ ἐν τοῖς Ὄρνισιν (v. 1295) οὐκ ἄλλως κορυδῷ αὐτὸν ὡ[μοίου] ἀλλ᾽ ἢ ὅτι μικρὸς καὶ αἰσχρός. καὶ Τηλεκλείδης Ἡσιόδοις\marqueur{3}·
`ἀλλ᾽ ἡ τάλαινα Φιλοκλέα βδε\douteux{λ}\douteux{λ}...\douteux{λ}.οθεν οὖν·
εἰ δ᾽ ἐστὶν Αἰσχύλου φρόνημ᾽ ἔχων'.}%
Le scholiaste ne se trompe pas lorsqu'il prétend que la comparaison avec l'\italique{alouette huppée}, κορυδός, était d'ordre physique\marqueur{4} : on connaît par les parémiographes\marqueur{5} le proverbe grec «~être plus hideux que Corydée~» :%
\citComm{Κορυδέως εἰδεχθέστερος· ἐπὶ δυσμορφίᾳ ἐκωμῳδεῖτο καὶ οὗτος καὶ οἱ παῖδες αὐτοῦ ὧν ἕνα τὸν Ἀρχέστρατον φασιν.}%
La pauvre alouette huppée n'est pas bien lotie par la tradition populaire : en plus d'être laide, même dans la force de l'âge elle est chétive, plus que l'aigle vieillissant (ἀετοῦ γῆρας, κορύδου νεότης), et n'a pas pour autant le ramage du rossignol (ἐν ἀμούσοις καὶ κόρυδος φθέγγεται\marqueur{6})... «~aux pays des aveugles, les borgnes sont rois~».\\
(Voir le commentaire au vers 441 pour le cas Xénoclès.)\\
Jean Taillardat a consacré la note que voici à Théognis\marqueur{7} :\\
Théognis le Tragique est ennuyeux et \italique{froid}, ψυχρός, au point qu'on le surnomme la \italique{Neige} (Χιών, schol. \italique{Acharniens} 11). La représentation d'une de ses pièces suffit à déchaîner les rigueurs de l'hiver, même sous le ciel déjà peu clément de la Thrace (\italique{Ach.} 138-140). Pris de pitié pour un froid rimailleur, Pisthétaire lui offre un vêtement de peau (\italique{Oiseaux} 935) : 
\citComm{ἔχε τὴν σπολάδα· πάντως δέ μοι ῥιγῶν δοκεῖς.\\
\francais{Prends ce manteau ; tu me parais complètement gelé.}}
Cette plaisanterie est continuée aux vers 946 et 955 (ταυταγὶ | τὰ κρυερὰ, \italique{ces vers glacés} ; cf. νιφοβόλους ἀναβολάς (\italique{ib.} 1385), « des préludes battus de neige ». Cet emploi de ψυχρός est banal\marqueur{8}.)}
\debutNote{168-70}
\note{1}{Manuscrit qui possède l'unique exemplaire des \italique{Thesmophories} qui nous soit parvenue -- \italique{unicodex} dit-on. Le manuscrit de Munich (\nieme{XV} siècle) n'étant qu'une recopie du premier (voir l'édition de von Velsen des \italique{Aristophanis Thesmophoriazusae} (1883) pour une comparaison complète du \italique{Ravennas} (R) et du \italique{Monacensis} (G)).}%
\note{2}{Suidas (φ 379) reprend cette scholie qu'il joint à d'autres éléments tirés du scholiaste aux \italique{Oiseaux} (v. 281).}%
\note{3}{fr. 15 in \italique{Poetae Comici Graeci}, VII, edd. Kassel et Austin.}%
\note{4}{Ce qui explique le surnom du \italique{Coryde} dont il est question dans les \italique{Deipnosophistes} d'Athénée (VI, 241a \&c.) «~qui aimait à se faire persifler~».}%
\note{5}{Zénobius IV, 59 ; Diogenianus V, 57 ; Michel Apostolius IX, 92 ; \&c. à retrouver dans le \italique{Corpus paroemiographorum graecorum}, ed. Leutsch. Cf. aussi Hésychius (σ 3681). Kassel \& Austin l'ont inclus (n° 827) aux \italique{Comica Adespota} de leur édition des \italique{Poetae Comici Graeci}, VIII.}%
\note{6}{Voir ce qu'en dit Érasme, qui a recensé ce proverbe dans ses \italique{Adages} (n° 1192).}%
\note{7}{\italique{Les images d'Aristophane} (1965), §759.}%
\note{8}{Ψυχρός en parlant d'un \italique{poète} : Alexis 179 [$=$ 184 Kassel-Austin $=$ Athénée, III, 123e], Machon (\italique{ap.} Athénée 580a), d'un \italique{orateur} : Plat. \italique{Euthyd.} 284e ; Théophil. 4, 4, \&c. Ψυχρός se dit aussi d'une tragédie manquée (\italique{Thesm.} 848), d'une plaisanterie sans sel (Eupol. 244 [$=$ 261 Kassel-Austin]).}\\\passeLigne}%
\end{groupe}%
%% v. 171-172a
\begin{groupe}%
\personnageGrec{Ἀγάθων}%
\trimIamb{ἅπασ᾽ ἀνάγκη ταῦτα γάρ τοι γνοὺς ἐγὼ\\
ἐμαυτὸν ἐθεράπευσα.\\\passeLigne}%
\end{groupe}%
%% v. 172b
\begin{groupe}%
\personnageGrec{Μνησίλοχος}%
\trimIamb{\ligne{172}\decalage{ἐμαυτὸν ἐθεράπευσα.}πῶς πρὸς τῶν θεῶν;\\\passeLigne}%
\end{groupe}%
%% v. 173-174
\begin{groupe}%
\personnageGrec{Εὐριπίδης}%
\trimIamb{παῦσαι βαΰζων καὶ γὰρ ἐγὼ τοιοῦτος ἦν\testimonia{173}{(βαύ̈ζων -- ἦν) Suidas β 192 $\simeq$ Zonaras 379 (sed γὰρ καὶ pro καὶ γὰρ)}\apparat{173}{ἦν R, Suid. β 192 : \lecon{ἦ} Dindorf (1825)}%
\commentaire{172-3}{πῶς πρὸς τῶν θεῶν; — παῦσαι βαΰζων.}{Goûtons à l'allitération en `t', `p' et surtout en `o' aboyée par Mnésiloque : π\textbf{ῶ}ς πρ\textbf{ὸ}ς τ\textbf{ῶ}ν θε\textbf{ῶ}ν ; c'est cette dissonance que lui reproche Euripide – pour la forme, car la question de Mnésiloque est aussi tout-à-fait désobligeante et ironique... il est certainement à l'affût de quelque double sens au «~ἐμαυτὸν ἐθεράπευσα~» d'Agathon. 
Euripide interrompt donc : παῦσαι βαΰζων, « cesse d'aboyer ! » -- quoique le verbe «~aboyer~» ne rende pas le registre tragique du verbe βαΰζω, formé par onomatopée sur le cri βαῢ βαΰ du chien (αὗ αὗ au vers 903 des \italique{Guêpes}).}\\
ὢν τηλικοῦτος, ἡνίκ᾽ ἠρχόμην ποιεῖν.\\\passeLigne}%
\end{groupe}%
%% v. 175
\begin{groupe}%
\personnageGrec{Μνησίλοχος}%
\trimIamb{μὰ τὸν Δί᾽ οὐ ζηλῶ σε τῆς παιδεύσεως.\testimonia{175}{(οὐ -- παιδεύσεως) Suidas ζ 63}%
\commentaire{173-5}{καὶ γὰρ ἐγὼ τοιοῦτος ἦν ... — οὐ ζηλῶ σε τῆς παιδεύσεως.}{Euripide se reconnaît – non sans quelque condescendance – une parenté artistique\marqueur{1} quand Mnésiloque répond comme si Agathon était héritier des mœurs d'Euripide.}%
\debutNote{173-5}%
\note{1}{Confirmée par Aristote. (Voir le commentaire \italique{infra} au vers 29, n. 3.)}\\\passeLigne}%
\end{groupe}%
%% v. 176a
\begin{groupe}%
\personnageGrec{Εὐριπίδης}%
\trimIamb{ἀλλ᾽ ὧνπερ οὕνεκ᾽ ἦλθον, ἔα μ᾽ εἰπεῖν.\testimonia{176}{Scholium v. 500 (sed ἀλλ' εἰπεῖν βούλομαι pro ἔα μ' εἰπεῖν)}\\\passeLigne}%
\end{groupe}%
%% v. 176b
\begin{groupe}%
\personnageGrec{Ἀγάθων}%
\trimIamb{\ligne{176}\decalage{ἀλλ᾽ ὧνπερ οὕνεκ᾽ ἦλθον, ἔα μ᾽ εἰπεῖν.}λέγε.
\apparat{176}{\lecon{(Ἀγάθ.)} λέγε R, Fritzsche : (Μνη.) Küster}%
\commentaire{176}{λέγε.}{\latin{Nam neque Euripides neque Agatho Mnesilochi jocos solent curare. 
Apparet vero, orationem sic cohaerere},%
\citComm{\grec{παῦσαι βαΰζων — άλλ' ... ἔα μ᾽ εἰπεῖν}\\
\latin{desine me maledictis insequi, sed potius ea me patere loqui.}}%
\latin{Jam necessario is, quis rogatus est, respondere debet, estque certo verbum λέγε dic igitur Agathonis.} (Note de Fritzsche.)}\\\passeLigne}%
\end{groupe}%
%% v. 177-180a
\begin{groupe}%
\ledpb%
\personnageGrec{Εὐριπίδης}%
\trimIamb{Ἀγάθων, σοφοῦ πρὸς ἀνδρός, ὅστις ἐν βραχεῖ\testimonia{177-80}{(σοφοῦ -- πρὸς σέ) Suidas σ 792}\testimonia{177-8}{(σοφοῦ -- λόγους) Suidas οι 154}\\
πολλοὺς καλῶς οἷός τε συντέμνειν λόγους.
\apparat{178}{οἷόν R : \lecon{οἷός} Suid. σ 792, οι 154, Stob. III,35.3 \espApp \lecon{συντέμνειν} R, Suid. οι 154, Stob. : συντεμεῖν Suid. σ 792 \espApp \textbf{179} κοινῇ R, Suid. κ 2556, σ 792 : \lecon{καινῇ} Biset ap. Schol.}%
\commentaire{177-8}{Ἀγάθων σοφοῦ πρὸς ἀνδρός ὅστις ἐν βραχεῖ / πολλοὺς καλῶς οἷός τε συντέμνειν λόγους.}{Citation quasi-textuelle de l'\italique{Éole} (Αἰόλος) d'Euripide rapportée par Stobée (3, 35, 3)\marqueur{1} :\citComm{παῖδες, σοφοῦ πρὸς ἀνδρός ὅστις ἐν βραχεῖ,\\
πολλοὺς λόγους οἷός τε συντέμνειν καλῶς}%
P.-L. Courier, \italique{Pamphlet des pamphlets} :
\citComm{\francais{C'est un mérite non commun, ni facile, de clore en peu de mots beaucoup de sens.}}}%
\debutNote{177-8}
\note{1}{À retrouver dans l'édition de C. Wachsmuth \& O. Hense, \italique{Ioannis Stobaei Anthologium}, vol. 3, 1904, p. 688 l. 5-7. C'est le fragment 21 de l'édition Jouan-van Looy aux Belles Lettres, t. VIII, 1\italique{ère} partie ; il correspond respectivement aux fragments 28 et 33 des éditions Kannicht et Mette.}%
\\
ἐγὼ δὲ καινῇ ξυμφορᾷ πεπληγμένος\testimonia{179}{Suidas κ 2556}%
\commentaire{179}{ἐγὼ δὲ κοινῇ ξυμφορᾷ πεπληγμένος.}{L'émendation de Biset : substituer καινῇ à κοινῇ est reprise par quasiment tous les éditeurs. Elle a l'avantage de concorder avec la scholie : μεγάλῃ καὶ θαυμαστῇ (καινός désigne au propre ce qui est \italique{nouveau}, différent de ce qui s'est produit jusqu'à présent, et par suite, ce qui n'était pas attendu, \italique{étrange}, \italique{extraordinaire} – avec une évolution tout identique à notre \italique{in-ouï}). 
Nous sommes devant l'un des cas \marqueur{1} où la scholie rend compte d'un texte en meilleur état que celui auquel elle est attenante ; c'est donc cette leçon que nous retiendrons encore qu'elle aille à l'encontre de deux passages de la \italique{Souda} (aux articles Κοινολογία et Σοῦ) dont l'un (Κοινολογία) définit même Κοινῇ ainsi : ἀνθρωπίνῃ, μεγάλῃ καὶ θαυμαστῇ (!!) \\
Loin de nous paraître être une nouvelle acception de κοινῇ, non attestée ailleurs, il nous semble beaucoup plus probable qu'elle soit un emprunt direct à notre texte erroné et sa scholie. Suidas -- ou l'un des rédacteurs de la \italique{Souda} si l'on préfère --, disposant de la scholie et du texte fautif, aura ajouté le sens de κοινῇ par lui connu : ἀνθρωπίνῃ\marqueur{2}.\\
Pour ne pas laisser de doute quant au sens de l'épithète de ξυμφορᾷ, reportons-nous au style du véritable Euripide (\italique{Alceste}, 404-5) : ὥστ᾽ ἐγὼ / καὶ σφὼ \textbf{βαρείᾳ} συμφορᾷ πεπλήγμεθα et plus loin (v. 856) καίπερ \textbf{βαρείᾳ} συμφορᾷ πεπληγμένος. À noter qu'Euripide emploie l'expression καινάς συμφοράς aux vers 1226-7 de l'\italique{Oreste}.}%
\debutNote{179}%
\note{1}{Voir pour un autre exemple le commentaire au vers 1174.}%
\note{2}{Voir \italique{infra} le commentaire au vers 735 pour une erreur analogue de Suidas.}\\
ἱκέτης ἀφῖγμαι πρὸς σέ.\\\passeLigne}%
\end{groupe}%
%% v. 180b
\begin{groupe}%
\personnageGrec{Ἀγάθων}%
\trimIamb{\ligne{180}\decalage{ἱκέτης ἀφῖγμαι πρὸς σέ.}τοῦ χρείαν ἔχων;
\commentaire{180}{τοῦ χρείαν ἔχων; }{Euripide, \italique{Suppliantes}, 115 : Τί χρῆμα θηρῶν καὶ τίνος χρείαν ἔχων;}\\\passeLigne}%
\end{groupe}%
%% v. 181-182
\begin{groupe}%
\personnageGrec{Εὐριπίδης}%
\trimIamb{μέλλουσί μ᾽ αἱ γυναῖκες ἀπολεῖν τήμερον\\
\lednopb%
τοῖς Θεσμοφορίοις, ὅτι κακῶς αὐτὰς λέγω.\\\passeLigne}%
\end{groupe}%
%% v. 183
\begin{groupe}%
\personnageGrec{Ἀγάθων}%
\trimIamb{τίς οὖν παρ᾽ ἡμῶν ἐστιν ὠφέλειά σοι;\\\passeLigne}%
\end{groupe}%
%% v. 184-187
\begin{groupe}%
\personnageGrec{Εὐριπίδης}%
\trimIamb{ἡ πᾶσ᾽· ἐὰν γὰρ ἐγκαθεζόμενος λάθρᾳ\apparat{184}{ἡ πᾶσα R : \lecon{ἡ πᾶσ᾽·} Küster : ἅπασ' Scaliger}\\
ἐν ταῖς γυναιξίν, ὡς δοκῶν εἶναι γυνή,\\
ὑπεραποκρίνῃ μου, σαφῶς σώσεις ἐμέ.%
\commentaire{186}{σαφῶς σώσεις ἐμέ.}{Le sigmatisme tant reproché à Euripide ! toujours à cause du verbe σῴζειν, \italique{sauver} : le véritable Euripide avait sifflé au vers 476 de sa \italique{Médée} : ἔσῳσά σ' ὡς ἴσασιν Ἑλλήνων ὅσοι. «~Grand merci de nous avoir sauvé des sigmas d'Euripide~» avait persifflé Platon le Comique\marqueur{1} en reprenant le début de vers : %
\citComm{εὖ γέ σοι γενοιθ᾽ ἡμᾶς ὅτι\\
ἔσωσας ἐκ τῶν σίγμα τῶν Εὐριπίδου.}%
}%
\debutNote{186}%
\note{1}{Fragment 29(30) de l'édition Kassel-Austin, \italique{Poetae Comici Graeci}, vol. VII (1989) ; chez Meineke, \italique{Fragmenta Comicorum Graecorum}, il se trouve au tome II, p. 626.}\\
μόνος γὰρ ἂν λέξειας ἀξίως ἐμοῦ.\\\passeLigne}%
\end{groupe}%
%% v. 188
\begin{groupe}%
\personnageGrec{Ἀγάθων}%
\trimIamb{ἔπειτα πῶς οὐκ αὐτὸς ἀπολογεῖ παρών;\\\passeLigne}%
\end{groupe}%
%% v. 189-192
\begin{groupe}%
\personnageGrec{Εὐριπίδης}%
\trimIamb{ἐγὼ φράσω σοι. πρῶτα μὲν γιγνώσκομαι\\
ἔπειτα πολιός εἰμι καὶ πώγων᾽ ἔχω,\\
σὺ δ᾽ εὐπρόσωπος λευκὸς ἐξυρημένος\\
γυναικόφωνος ἁπαλὸς εὐπρεπὴς ἰδεῖν.\testimonia{192}{(γυναικόφωνος) Pollux II 111}\\\passeLigne}%
\end{groupe}%
%% v. 193a
\begin{groupe}%
\personnageGrec{Ἀγάθων}%
\trimIamb{Εὐριπίδη—\\\passeLigne}%
\end{groupe}%
%% v. 193b
\begin{groupe}%
\personnageGrec{Εὐριπίδης}%
\trimIamb{\ligne{193}\decalage{Εὐριπίδη—}τί ἔστιν;\\\passeLigne}%
\end{groupe}%
%% v. 193c-194
\begin{groupe}%
\personnageGrec{Ἀγάθων}%
\trimIamb{\ligne{193}\decalage{Εὐριπίδη—τί ἔστιν;}ἐποίησάς ποτε,\testimonia{193}{(ἐποίησάς ποτε) \Eustathe 4.9}\\
`χαίρεις ὁρῶν φῶς, πατέρα δ᾽ οὐ χαίρειν δοκεῖς;'
\commentaire{193}{χαίρεις ὁρῶν φῶς, πατέρα δ᾽ οὐ χαίρειν δοκεῖς;}{Citation tirée de l'\italique{Alceste} d'Euripide, v. 691 : Phérès envoie paître son fils Admète qui lui demande de mourir à sa place. C'est finalement Alceste qui se sacrifiera pour son mari ; tout comme Mnésiloque, parent par \italique{alliance} d'Euripide, ne va pas tarder à accepter ce qui ressemble fort à un suicide.}\\\passeLigne}%
\end{groupe}%
%% v. 195a
\begin{groupe}%
\personnageGrec{Εὐριπίδης}%
\trimIamb{ἔγωγε.\\\passeLigne}%
\end{groupe}%
%% v. 195b-199
\begin{groupe}%
\personnageGrec{Ἀγάθων}%
\trimIamb{\ligne{195}\decalage{ἔγωγε.}μή νυν ἐλπίσῃς τὸ σὸν κακὸν\testimonia{195-7}{(μή νυν -- φέρε) Suidas υ 709 (νῦν pro νυν et ἀλλ' omisso)}\\
ἡμᾶς ὑφέξειν. καὶ γὰρ ἂν μαινοίμεθ᾽ ἂν.\apparat{196}{γὰρ R : \lecon{γὰρ ἂν} Suid. υ 709}\\
ἀλλ᾽ αὐτὸς ὅ γε σόν ἐστιν οἰκείως φέρε.\\
τὰς συμφορὰς γὰρ οὐχὶ τοῖς τεχνάσμασιν\testimonia{198-201}{(τὰς συμφορὰς -- παθήμασιν) Suidas τ 440 (τεχνάσμασι et παθήμασι pro -σιν)}\\
φέρειν δίκαιον ἀλλὰ τοῖς παθήμασιν.\\\passeLigne}%
\end{groupe}%
%% v. 200-201
\begin{groupe}%
\personnageGrec{Μνησίλοχος}%
\trimIamb{καὶ μὴν σύ γ᾽ ὦ κατάπυγον εὐρύπρωκτος εἶ\\
οὐ τοῖς λόγοισιν ἀλλὰ τοῖς παθήμασιν.\commentaire{200-1}{εὐρύπρωκτος εἶ / οὐ τοῖς λόγοισιν ἀλλὰ τοῖς παθήμασιν.}{Mnésiloque reprend l'opposition τοῖς τεχνάσμασιν / τοῖς παθήμασιν d'Agathon en substituant τοῖς λόγοισιν à τοῖς τεχνάσμασιν.\\
Le spectateur attend donc l'opposition classique τοῖς λόγοις / τοῖς ἔργοις : \italique{en théorie} / \italique{en pratique}... sauf qu'ici les \italique{pratiques} d'Agathon sont jugées par trop passives.}\\\passeLigne}%
\end{groupe}%
%% v. 202
\begin{groupe}%
\personnageGrec{Εὐριπίδης}%
\trimIamb{τί δ᾽ ἔστιν ὅτι δέδοικας ἐλθεῖν αὐτόσε;\\\passeLigne}%
\end{groupe}%
%% v. 203a
\begin{groupe}%
\personnageGrec{Ἀγάθων}%
\trimIamb{κάκιον ἀπολοίμην ἂν ἢ σύ.\\\passeLigne}%
\end{groupe}%
%% v. 203b
\begin{groupe}%
\personnageGrec{Εὐριπίδης}%
\trimIamb{\ligne{203}\decalage{κάκιον ἀπολοίμην ἂν ἢ σύ.}πῶς;\\\passeLigne}%
\end{groupe}%
%% v. 203c-205
\begin{groupe}%
\personnageGrec{Ἀγάθων}%
\trimIamb{\ligne{203}\decalage{κάκιον ἀπολοίμην ἂν ἢ σύ.πῶς;}ὅπως;\\
δοκῶν γυναικῶν ἔργα νυκτερείσια\apparat{204}{\lecon{νυκτερείσια} R : νυκτερήσια Bothe}%
\commentaire{204}{νυκτερείσια.}{Bon nombre d'éditeurs corrigent en νυκτερήσια (\italique{nocturne}, attesté que tardivement ; par Lucien et Sextus Empiricus). S'il n'est pas douteux qu'Aristophane emploie cet adjectif en ce sens – ou bien le forge en miroir de ἡμερήσιος (Eschyle, \italique{Agamemnon}, 22) –, la phrase perdrait de son sel aristophanesque si l'on ne conservait pas la prononciation transmise par le codex de Ravenne et certainement voulue par Aristophane lui-même : νυκτ-ερείσια fait entendre le substantif du verbe ἐρείδειν, « \latin{quod in comoedia est, subagitare mulierem} » (Fritzsche).}\\
κλέπτειν ὑφαρπάζειν τε θήλειαν Κύπριν.%
\commentaire{205}{ὑφαρπάζειν τε θήλειαν Κύπριν.}{ L'expression ὑφαρπάζειν Κύπριν se retrouve au \nieme{722} vers de l'\italique{Assemblée des femmes}.\\
Les Athéniennes ont toujours une bonne raison d'en vouloir aux poètes tragiques : lorsqu'ils ne médisent pas d'elles, ils se livrent à une concurrence déloyale.%
\citComm{Mais les jeunes femmes d'Athènes\\
Privées d'leurs époux, d'leurs galants,\\
Accumulèrent la rancune patiemment...\\
Puis un jour, ivres de colère,\\
Elles s'armèrent de bâtons,\\
Et, farouches, elles immolèrent Agathon !}%
%Remarque en passant : le fait que l'action se déroule au \italique{jour du milieu} des Thesmophories, νηστεία, et qu'il fût théoriquement un jour d'abstinence sexuelle n'empêchera en rien les médisances... au contraire !
}\\\passeLigne}%
\end{groupe}%
%% v. 206-207
\begin{groupe}%
\personnageGrec{Μνησίλοχος}%
\trimIamb{ἰδού γε κλέπτειν· νὴ Δία βινεῖσθαι μὲν οὖν.\apparat{206}{\lecon{νὴ Δία} R : μà Δία Fritzsche}\\
ἀτὰρ ἡ πρόφασίς γε νὴ Δί᾽ εἰκότως ἔχει.
\commentaire{207}{ἀτὰρ ἡ πρόφασίς γε νὴ Δί᾽ εἰκότως ἔχει.}{Molière, \italique{Misanthrope}, IV, 3 :
\citComm{\francais{Ah! le détour est bon et l'excuse admirable.}}}\\\passeLigne}%
\end{groupe}%
%% v. 208a
\begin{groupe}%
\personnageGrec{Εὐριπίδης}%
\trimIamb{τί οὖν; ποιήσεις ταῦτα;\apparat{208}{οὖν et ταῦτα R : \lecon{οὖν;} et \lecon{ταῦτα;} Ellebodius}\\\passeLigne}%
\end{groupe}%
%% v. 208b
\begin{groupe}%
\personnageGrec{Ἀγάθων}%
\trimIamb{\ligne{208}\decalage{τί οὖν; ποιήσεις ταῦτα;}μὴ δόκει γε σύ.\\\passeLigne}%
\end{groupe}%
%% v. 209
\begin{groupe}%
\personnageGrec{Εὐριπίδης}%
\trimIamb{ὦ τρισκακοδαίμον, ὡς ἀπόλωλ᾽ Εὐριπίδης.
\apparat{209}{\lecon{τρισκακόδαιμον} R : τρισκακοδαίμων \Mon \espApp \lecon{(Εὐρ.)} R : (Μνη.) Elmsley ad \italique{Ach}. 475}%
\commentaire{209}{ὦ τρισκακόδαιμον ὡς ἀπόλωλ᾽ Εὐριπίδης.}{La plupart des éditeurs modernes découpent cette réplique, attribuant la première partie à Euripide, puis Mnésiloque de lancer l'apostrophe : Εὐριπίδη etc. Cela les oblige à retoucher le texte pour faire de ὤ τρισκακοδαίμων le sujet de ἀπόλωλ(α) et changer Εὐριπίδης en Εὐριπίδη... \\
Tout cela paraît bien tarabiscoté alors que le texte de notre unicodex se laisse comprendre facilement : « Ô triple malheureux ! ainsi périt Euripide », apostrophe qu'Euripide s'adresse à lui-même d'autant moins étonnante qu'il a déjà parlé de lui à la \nieme{3} personne (v. 77).
En outre, si l'attribution des répliques dans le manuscrit de Ravenne est souvent folklorique, la découpe des répliques, elle, est bien plus sûre ; les raisons invoquées sont donc bien trop faibles pour se permettre de tailler ainsi le vers.}\\\passeLigne}%
\end{groupe}%
%% v. 210
\begin{groupe}%
\personnageGrec{Μνησίλοχος}%
\trimIamb{ὦ φίλτατ᾽ ὦ κηδεστὰ μὴ σαυτὸν προδῷς.\\\passeLigne}%
\end{groupe}%
%% v. 211a
\begin{groupe}%
\personnageGrec{Εὐριπίδης}%
\trimIamb{πῶς οὖν ποιήσω δῆτα;\\\passeLigne}%
\end{groupe}%
%% v. 211b-212
\begin{groupe}%
\personnageGrec{Μνησίλοχος}%
\trimIamb{\ligne{211}\decalage{πῶς οὖν ποιήσω δῆτα;}τοῦτον μὲν μακρὰ\\
κλάειν κέλευ᾽, ἐμοὶ δ᾽ ὅ τι βούλει χρῶ λαβών.\\\passeLigne}%
\end{groupe}%
%% v. 213-214a
\begin{groupe}%
\personnageGrec{Εὐριπίδης}%
\trimIamb{ἄγε νυν ἐπειδὴ σαυτὸν ἐπιδίδως ἐμοί,\\
ἀπόδυθι τουτὶ θοἰμάτιον.\\\passeLigne}%
\end{groupe}%
%% v. 214b-215a
\begin{groupe}%
\personnageGrec{Μνησίλοχος}%
\trimIamb{\ligne{214}\decalage{ἀπόδυθι τουτὶ θοἰμάτιον.}καὶ δὴ χαμαί.\\
ἀτὰρ τί μέλλεις δρᾶν μ᾽;\apparat{215}{με R : \lecon{μ᾽} edd. omnes}\\\passeLigne}%
\end{groupe}%
%% v. 215b-216a
\begin{groupe}%
\personnageGrec{Εὐριπίδης}%
\trimIamb{\ligne{215}\decalage{ἀτὰρ τί μέλλεις δρᾶν μ᾽;}ἀποξυρεῖν ταδί,\testimonia{215-6}{(ἀποξυρεῖν -- ἀφεύειν) Suidas α 4565 (ἀφαύειν pro ἀφεύειν), cf. Scholium ταῦτα δὲ ἔλαβεν ἐκ τῶν Ἰδαίων Κρατίνου (fr. 90 Kassel-Austin) ; \ClemAlex, \italique{Στρώματα} VI 26.4 Ἀριστοφάνης δὲ ὁ κωμικὸς ἐν ταῖς πρώταις Θεσμοφοριαζούσαις τὰ ἐκ τῶν Κρατίνου Ἐμπιπραμένων μετήνεγκεν ἔπη.}\\
τὰ κάτω δ᾽ ἀφεύειν.
\apparat{216}{\lecon{ἀφεύειν} R : ἀφαύειν Suid. α 4565 \espApp είσω R : \lecon{εἴ σοι} Bonini}\\\passeLigne}%
\end{groupe}%
%% v. 216b-217
\begin{groupe}%
\personnageGrec{Μνησίλοχος}%
\trimIamb{\ligne{216}\decalage{τὰ κάτω δ᾽ ἀφεύειν.}ἀλλὰ πρᾶττ᾽, εἴ σοι δοκεῖ.\\
ἢ μὴ διδόναι γ' ἑαυτὸν ὤφελόν ποτε.
\apparat{217}{διδόναι γ' αὐτὸν R : διδόναι γ' ἑμαυτὸν Scaliger : 'πιδοῦναι 'μαυτὸν Dawes : \lecon{διδόναι γ' ἑαυτὸν} Bibi auctore Willems}%
\commentaire{217}{ἢ μὴ διδόναι γ' αὐτὸν ὤφελόν ποτε.}{Le vers pèche contre le mètre.\\
Willems (voir sa note 4 en annexe), en guerre perpétuelle contre les retoucheurs de texte, rendait justice à ce qu'il croyait être le texte original\marqueur{1}. Il s'en justifiait par des exemples qui ne laissent nul doute de l'epsilonesque correction qu'il aurait lui-même proposée s'il avait eu le texte des manuscrits sous les yeux : substituer γ' ἑαυτὸν à γ' αὐτὸν.\\
% Cette correction est d'autant plus facile à accepter si l'on s'imagine une version en onciales, ΓΕΑΥΤΟΝ ; un scribe ayant malencontreusement lu γε αυτον avant d'élider le γε. 
Au reste, comme l'avait souligné Willems, « διδόναι ἑαυτὸν τινι est la formule ordinaire, et qui se lit partout : \italique{Cavaliers}., 740 ; Sophocle, \italique{Philoct}., 84 ; \italique{Trach}., 1117; Hérodote, VI, 108; Thucydide, II, 68; Démosthène, \italique{Sur la couronne}, 503d.~»}%
\debutNote{217}%
\note{1}{En fait légèrement modifié : γ' ἑμαυτὸν à la place de γ' αὐτὸν. Willems a-t-il été trompé par Dindorf ? (qui prétendait -- à tort -- que γ' ἑμαυτὸν était la lecture originale du manuscrit de Münich.)}%
\\\passeLigne}%
\end{groupe}%
%% v. 218-219a
\begin{groupe}%
\personnageGrec{Εὐριπίδης}%
\trimIamb{Ἀγάθων σὺ μέντοι ξυροφορεῖς ἑκάστοτε,\\
χρῆσόν τι νῦν ἡμῖν ξυρόν.\testimonia{219-20}{(χρῆσόν -- ξυροδόκης) Pollux X 140 (νῦν omisso)}\\\passeLigne}%
\end{groupe}%
%% v. 219b-220a
\begin{groupe}%
\personnageGrec{Ἀγάθων}%
\trimIamb{\ligne{219}\decalage{χρῆσόν τί νυν ἡμῖν ξυρόν.}αὐτὸς λάμβανε\\
ἐντεῦθεν ἐκ τῆς ξυροδόκης.\testimonia{220}{(ξυροδόκης) Pollux II 32 ; \Eustathe 130.43, 1399.38 ; cf. Scholium}%
\apparat{220}{ξυροδίκης R : \lecon{ξυροδόκης} Schol., Pollux II.32, X.140}\\\passeLigne}%
\end{groupe}%
%% v. 220b-221
\begin{groupe}%
\personnageGrec{Εὐριπίδης}%
\trimIamb{\ligne{220}\decalage{ἐντεῦθεν ἐκ τῆς ξυροδόκης.}γενναῖος εἶ.\\
κάθιξε· φύσα τὴν γνάθον τὴν δεξιάν.\\\passeLigne}%
\end{groupe}%
%% v. 222a
\begin{groupe}%
\personnageGrec{Μνησίλοχος}%
\trimIamb{οἴμοι.
\apparat{222}{ὤιμοι R : \lecon{οἴμοι} Dindorf}\\\passeLigne}%
\end{groupe}%
%% v. 222b-223a
\begin{groupe}%
\personnageGrec{Εὐριπίδης}%
\trimIamb{\ligne{222}\decalage{οἴμοι.}τί κέκραγας; ἐμβαλῶ σοι πάτταλον,\\
ἢν μὴ σιωπᾷς.\\\passeLigne}%
\end{groupe}%
%% v. 223b
\begin{groupe}%
\personnageGrec{Μνησίλοχος}%
\trimIamb{\ligne{223}\decalage{ἢν μὴ σιωπᾷς.}ἀτταταῖ ἰατταταῖ.\testimonia{223}{(ἀτταταῖ ἰατταταῖ) Suidas ε 2807}\apparat{223}{ἀττατα· ατταται R : \lecon{ἀτταταῖ ἰατταταῖ} Biset}\\\passeLigne}%
\end{groupe}%
%% v. 224a
\begin{groupe}%
\personnageGrec{Εὐριπίδης}%
\trimIamb{οὗτος σὺ ποῖ θεῖς;\commentaire{224}{οὗτος σὺ ποῖ θεῖς;}{formule identique au vers 854 des \italique{Guêpes}.}\\\passeLigne}%
\end{groupe}%
%% v. 224b-226a
\begin{groupe}%
\personnageGrec{Μνησίλοχος}%
\trimIamb{\ligne{224}\decalage{οὗτος σὺ ποῖ θεῖς;}εἰς τὸ τῶν σεμνῶν θεῶν·
\commentaire{224}{εἰς τὸ τῶν σεμνῶν θεῶν·}{Le scholiaste :
\citComm{ἀντὶ τοῦ `εἰς <τὸ> τῶν ᾿Ερινύων'. ἄσυλον γὰρ εἶχον οἵ καταλαμβάνοντες τὰ ἱερὰ τῶν Ἐρινύων.}}\\
οὐ γὰρ μὰ τὴν Δήμητρά γ᾽ ἐνταυθοῖ μενῶ\apparat{225}{\lecon{Δήμητρά γ᾽} R : Δήμητρ᾽ ἔτ᾽ Porson, cf. ad \italique{Eq}. 698 : Δημήτερ᾽ Fritzsche}\\
τεμνόμενος.\\\passeLigne}%
\end{groupe}%
%% v. 226b-227
\begin{groupe}%
\personnageGrec{Εὐριπίδης}%
\trimIamb{\ligne{226}\decalage{τεμνόμενος.}οὔκουν καταγέλαστος δῆτ᾽ ἔσει\apparat{226}{οὐκ οὖν R : \lecon{οὔκουν} Brunck}\\
τὴν ἡμίκραιραν τὴν ἑτέραν ψιλὴν ἔχων;\testimonia{227}{Suidas η 334 \espTest (ἡμίκραιραν) Hesychius η 496 ἡμίκραιρα ; \EtymGen (AB) 150 ed. Miller ; \Eustathe 710.50, 1127.35 ; \AelDion η 9}\\\passeLigne}%
\end{groupe}%
%% v. 228a
\begin{groupe}%
\personnageGrec{Μνησίλοχος}%
\trimIamb{ὀλίγον μέλει μοι.\\\passeLigne}%
\end{groupe}%
%% v. 228b-229a
\begin{groupe}%
\personnageGrec{Εὐριπίδης}%
\trimIamb{\ligne{228}\decalage{ὀλίγον μέλει μοι.}μηδαμῶς πρὸς τῶν θεῶν\\
προδῷς με· χώρει δεῦρο.\\\passeLigne}%
\end{groupe}%
%% v. 229b
\begin{groupe}%
\personnageGrec{Μνησίλοχος}%
\trimIamb{\ligne{229}\decalage{προδῷς με· χώρει δεῦρο.}κακοδαίμων ἐγώ.\\\passeLigne}%
\end{groupe}%
%% v. 230
\begin{groupe}%
\personnageGrec{Εὐριπίδης}%
\trimIamb{ἔχ᾽ ἀτρέμα σαυτὸν κἀνάκυπτε· ποῖ στρέφει;
\commentaire{230}{ἔχ᾽ ἀτρέμα σαυτὸν κἀνάκυπτε·}{C'est la leçon du \italique{Ravennas}. 
M. van Leeuwen, d'après Dobree, corrige : ἔχ᾽ ἀτρέμας αὐτοῦ, \italique{quia ferri nequit accusativus} ; lequel αὐτοῦ était à sa place en \italique{Oiseaux}, 1200, mais pas du tout ici. 
Hérodote a pourtant dit, en parlant de l'armée athénienne : εἶχον ἀτρέμας σφέας αὐτοὺς ἵνα ἐτάχθησαν\marqueur{1}. Hérodote le pouvait, mais non Aristophane, qui ailleurs s'est borné à ἔχ' ἀτρέμα\marqueur{2}. Quoi ! voilà un comique qui, à quatre reprises (car on s'est donné la peine de compter), a écrit : \italique{reste coi}, et la cinquième fois il se permettrait d'écrire : \italique{tiens-toi coi}. 
À quoi servirait la statistique si l'on tolérait de pareilles licences ?
(Note vengeresse n°5 de Willems.)}%
\debutNote{230}%
\note{1}{Hérodote, IX, 54 (aussi III, 79). — De même Platon, \italique{Epist.}, 348 B, écrit : ἐκποδὼν ἑαυτὸν ἔσχεν et Euripide, \italique{Iph. Taur.}, 1129 : ἐκποδὼν τοῦδ' ἔχειν μιάσματος.}%
\note{2}{ll est vrai qu'une des quatre fois les bons manuscrits portent ἔχ' ἀτρεμί (\italique{Nuées}, 261). Était-ce assez dégoûtant ? Heureusement que M. van Leeuwen y a mis bon ordre.}\\\passeLigne}%
\end{groupe}%
%% v. 231a
\begin{groupe}%
\personnageGrec{Μνησίλοχος}%
\trimIamb{μυμῦ.\testimonia{231}{\EtymGen (AB) 218 ed. Miller \espTest (μῦ μῦ) Suidas ε 2807 ; cf. \Eustathe 723.9 \espTest (τί -- καλῶς) Suidas μ 1383 ; \Apostol XVI 60c \espTest (μύζεις) \Eustathe 440.23}\\\passeLigne}%
\end{groupe}%
%% v. 231b
\begin{groupe}%
\personnageGrec{Εὐριπίδης}%
\trimIamb{\ligne{231}\decalage{μυμῦ.}τί μύζεις; πάντα πεποίηται καλῶς.%\commentaire{231}{τί μύζεις;}{\eff{Érasme 702}}
\\\passeLigne}%
\end{groupe}%
%% v. 232
\begin{groupe}%
\personnageGrec{Μνησίλοχος}%
\trimIamb{οἴμοι κακοδαίμων, ψιλὸς αὖ στρατεύσομαι.\testimonia{232}{Suidas ψ 105}\apparat{232}{\lecon{αὖ} R, Suid. ψ 105 : οὖν von Velsen}%
\commentaire{232}{ψιλὸς αὖ στρατεύσομαι.}{ψιλός, comme \italique{levis} en latin, signifie à la fois \italique{lisse}, \italique{sans poil}, et \italique{léger}, \italique{armé à la légère}. (Note de Willems.)\\
% Il n'y a donc pas lieu de corriger le texte : le « αὖ » sert au jeu de mots.\\
Cette note sur ψιλός me sera l'occasion de digresser sur l'un des plus beaux échanges à bâtons rompus de la Tragédie grecque : chez Sophocle (\italique{Ajax}, v. 1120 sq.), entre Ménélas et Teucer :
\citComm{Μενέλαος — Ὁ τοξότης ἔοικεν οὐ σμικρὸν φρονεῖν.\\	
Τεύκρος — Οὐ γὰρ βάναυσον τὴν τέχνην ἐκτησάμην.\\
Μενέλαος — Μέγ᾽ ἄν τι κομπάσειας, ἀσπίδ᾽ εἰ λάβοις.\\
Τεύκρος — Κἂν \textbf{ψιλὸς} ἀρκέσαιμι σοί γ᾽ ὡπλισμένῳ.\\
Μενέλαος — Ἡ γλῶσσά σου τὸν θυμὸν ὡς δεινὸν τρέφει.\\
Τεύκρος — Ξὺν τῷ δικαίῳ γὰρ μέγ᾽ ἔξεστιν φρονεῖν.}%
Ménélas s'étonne de ce que Teucer se permette de lui adresser des reproches – à lui le roi de Sparte ! –, vil archer de son état — Mais l'archerie est loin d'être méprisable lui rétorque Teucer — Ménélas ironise : Que serait-ce s'il portait un bouclier — La réplique de Teucer fuse : « Même \italique{ψιλὸς}, je tiendrais bon devant toi et ton [royal] armement ». Comprendre : même \italique{fantassin de l'armée légère} ; c'est-à-dire l'échelon le plus bas chez les militaires – encore plus bas qu'archer –, poste souvent réservé à des esclaves et donc le plus opposé à la royale position de Ménélas et le plus inégal en terme d'armements. C'en est trop pour Ménélas qui s'indigne de cette langue par trop pendue. C'est alors que Teucer lance la fameuse et cinglante réplique : « La fierté est permise, quand on a pour soi la justice ! »}\\\passeLigne}%
\end{groupe}%
%% v. 233-234a
\begin{groupe}%
\personnageGrec{Εὐριπίδης}%
\trimIamb{μὴ φροντίσῃ· ὡς εὐπρεπὴς φανεῖ πάνυ.\\
βούλει θεάσασθ᾽ αὑτόν;\apparat{234}{θεάσασθαι σεαυτόν R : \lecon{θεάσασθ᾽ αὑτόν} Fritzsche : θεᾶσθαι σαυτόν Porson}%
\commentaire{234}{βούλει θεάσασθαι σεαυτόν; — εἰ δοκεῖ, φέρε.}{Le vers péchant contre le mètre, la correction la plus convaincante me semble être celle de Fritzsche qui substitue αὑτόν au σεαυτόν, faisant l'hypothèse que σεαυτόν aura été suggéré à un scribe inattentif par la proximité avec ὁρᾷς σεαυτόν du vers suivant.}%
\\\passeLigne}%
\end{groupe}%
%% v. 234b
\begin{groupe}%
\personnageGrec{Μνησίλοχος}%
\trimIamb{\ligne{234}\decalage{βούλει θεάσασθ᾽ αὑτόν;}εἰ δοκεῖ, φέρε.\\\passeLigne}%
\end{groupe}%
%% v. 235a
\begin{groupe}%
\personnageGrec{Εὐριπίδης}%
\trimIamb{ὁρᾷς σεαυτόν;%
%\commentaire{234-5}{φέρε — ὁρᾷς σεαυτόν;}{Sur le τόπος comique du miroir, \cf Platon, \italique{Alcibiade}, I, 132.\eff{Remarquons que ce doit être à ce moment que Mnésiloque saisit le miroir que l'on retrouve quelques scènes plus loin sur le cratère de Würtzburg. Et le poignard serait-il le rasoir ?}}
\\\passeLigne}%
\end{groupe}%
%% v. 235b
\begin{groupe}%
\personnageGrec{Μνησίλοχος}%
\trimIamb{\ligne{235}\decalage{ὁρᾷς σεαυτόν;}οὐ μὰ Δί᾽ ἀλλὰ Κλεισθένη.\testimonia{235}{Suidas κ 1756 Κλεισθένην ὁρῶ}\apparat{235}{Κλεισθένην R, Suid. κ 1756 : \lecon{Κλεισθένη} Brunck}%
\commentaire{235}{Κλεισθένην.}{Clisthène, l'efféminé notoire, personnage qui apparaîtra dans la pièce au vers 571 -- et sera d'ailleurs reconnue comme femme -- avant de prêter allégeance au genre féminin. 
Cette tête de Turc d'Aristophane est raillée comme : une pseudo-femme (\italique{Nuées} 355, \italique{Oiseaux} 829-831, \italique{Thesmophories} 571) ; en tout cas femme en amour (\italique{Lysistrata} 621, 1091-1092, \italique{Thesmophories} 574, \italique{Grenouilles} 48) ; imberbe comme femme (\italique{Acharniens} 117-121, \italique{Grenouilles} 422-424) ; Clisthène ne contrebalance même pas ces innombrables défauts – qui ne sont pas tous du ressort de sa volonté – par une misogynie affectée : il trouve encore le moyen d'être traître à son sexe en s'alliant au genre féminin (\italique{Lysistrata} 621, \italique{Thesmophories} 575-576, 602).}\\\passeLigne}%
\end{groupe}%
%% v. 236
\begin{groupe}%
\personnageGrec{Εὐριπίδης}%
\trimIamb{ἀνίστασ᾽, ἵν᾽ ἀφεύσω σε, κἀγκύψας ἔχε.\\\passeLigne}%
\end{groupe}%
%% v. 237
\begin{groupe}%
\personnageGrec{Μνησίλοχος}%
\trimIamb{οἴμοι κακοδαίμων δελφάκιον γενήσομαι.\testimonia{237}{Suidas δ 205 (sed γέγονα pro γενήσομαι)}%
\commentaire{237}{δελφάκιον}{ est un diminutif d'ἡ δέλφαξ, le porc \italique{mature}, opposé au χοῖρος, \italique{le porcelet}.
%Hésychius (δ 599) précise que le mot δελφάκιον, comme χοῖρος, il peut renvoyer au γυναικεῖον [αἰδοῖον]. Mais ne nous égarons pas ! 
N'en déplaise à certains commentateurs, il n'est aucunement question ici -- pas plus qu'au vers 1061 de \italique{Lysistrata} -- de sexe féminin. Et d'abord, ce serait absurde puisqu'il est question du côté pile ; côté épilé par flambage. 
L'image est ailleurs : c'est celle de de viande de cochon dont on retire les poils par le feu.
Cf. la \italique{Souda} (δ 205):%
\citComm{δελφάκιον : ὁ μικρὸς χοῖρος. Ἀριστοφάνης: οἴμοι κακοδαίμων, δελφάκιον γέγονα. τοῦτο εἴρηκεν ἀφευθεὶς τὸν πρωκτόν. ἀντὶ τοῦ φλογισθείς. μετὰ γὰρ τὸ τυθῆναι τὰ δελφάκια φλογίζονται, ἵνα ψιλωθῶσιν.}%
(On pourra aussi se reporter à Homère\marqueur{1} ou Eschyle\marqueur{2}.)}%
\debutNote{237}%
\note{1}{\italique{Iliade}, IX, 467-8 et \italique{Odyssée}, II, 300 ; XIV, 73-5}%
\note{2}{Fragment 310 de l'édition de S. Radt, \italique{Tragicorum Graecorum Fragmenta}, vol. 3 (1985).}%
\\\passeLigne}%
\end{groupe}%
%% v. 238-239
\begin{groupe}%
\personnageGrec{Εὐριπίδης}%
\trimIamb{ἐνεγκάτω τις ἔνδοθεν δᾷδ᾽ ἢ λύχνον.\\
ἐπίκυπτε· τὴν κέρκον φυλάττου νυν ἄκραν.
\commentaire{239}{τὴν κέρκον φυλάττου νυν ἄκραν.}{Euripide reprend l'image du δελφάκιον flambé ; il parle bien-sûr de l'\italique{olisbos} de Mnésiloque.}\\\passeLigne}%
\end{groupe}%
%% v. 240-242
\begin{groupe}%
\personnageGrec{Μνησίλοχος}%
\trimIamb{ἐμοὶ μελήσει νὴ Δία, πλήν γ᾽ ὅτι κάομαι.\\
οἴμοι τάλας. ὕδωρ ὕδωρ ὦ γείτονες%
\commentaire{241-2}{οἴμοι τάλας. ὕδωρ ὕδωρ ὦ γείτονες / πρὶν ἀντιλαβέσθαι πρωκτὸν ἓτερον τῆς φλογός.}{Mnésiloque est littéralement en train de prendre feu ; toute sa réplique est un pastiche d'un appel au secours, comme si sa propre maison s'embrasait\marqueur{1} et que celles, contiguës, de ses voisins allaient prendre feu à leur tour (sauf qu'il s'agit ici non d'οἶκός mais de πρωκτός)... il appelle donc à puiser de l'eau « ὕδωρ ! ὕδωρ ! », « À l'eau ! Non mais, Allô quoi ! » \\%j'suis un homme et j'ai plus de poils au c... »\\
Horace prodigue ce conseil dans ses \italique{Épîtres} (1, 18, 84) :
\citComm{\latin{Tua res agitur, paries cum priximus ardet}\\
\francais{C'est ton intérêt qui est en jeu, lorsque la maison de ton voisin brûle\marqueur{2}.}}
L'image de l'incendie qui se propage de maison en maison était si fréquente dans le monde latin et roman\marqueur{3} qu'on a peine à croire qu'elle ne le fût pas aussi dans le monde grec -- témoin ce passage d'Aristophane.}%
\debutNote{241-2}%
\note{1}{La scholie au vers 242 : (πρωκτὸν·) ὡς εἰ ἔλεγεν ``οἰκίαν''.}%
\note{2}{J'emprunte la traduction au magistral \italique{Dictionnaire des sentences latines et grecques} de Renzo Tosi (traduit de l'italien par Rebecca Lenoir). La consultation de ce dictionnaire me fut d'une grande aide pour découvrir, en ce passage comme en d'autres, des ponts entre la langue d'Aristophane et l'univers des proverbes.}%
\note{3}{Voir le \italique{Dictionnaire} sus-cité (§524).}\\
πρὶν ἀντιλαβέσθαι πρωκτὸν ἓτερον τῆς φλογός.\apparat{242}{\lecon{[ἓ]τερον} \PsiUn \ (restituit Medaglia) : ἓτερον om. R}%
\\\passeLigne}%
\end{groupe}%
%% v. 243a
\begin{groupe}%
\personnageGrec{Εὐριπίδης}%
\trimIamb{θάρρει.\\\passeLigne}%
\end{groupe}%
%% v. 243b
\begin{groupe}%
\personnageGrec{Μνησίλοχος}%
\trimIamb{\ligne{243}\decalage{θάρρει.}τί θαρρῶ καταπεπυρπολημένος;
\commentaire{243}{θάρρει — τί θαρρῶ;}{Plaute, \italique{Curculio}, IV, 3, 21 : \latin{bellator vale — quid valeam ?}}\\\passeLigne}%
\end{groupe}%
%% v. 244-245a
\begin{groupe}%
\personnageGrec{Εὐριπίδης}%
\trimIamb{ἀλλ᾽ οὐκ ἔτ᾽ οἰδὲν πρᾶγμά σοι· τὰ πλεῖστα γὰρ\testimonia{244-5}{(τὰ πλεῖστα -- ἀσβόλου) Photius α 2946 $=$ \SynB α 2213 ; \AelDion α 186}\\
ἀποπεπόνηκας.\\\passeLigne}%
\end{groupe}%
%% v. 245b-246
\begin{groupe}%
\personnageGrec{Μνησίλοχος}%
\trimIamb{\ligne{245}\decalage{ἀποπεπόνηκας.}φεῦ ἰοὺ τῆς ἀσβόλου.\testimonia{245}{(φεῦ -- ἀσβόλου) Suidas α 4132 \espTest (ἄσβολος $=$ ἀσβόλη) Moeris α 17 $\simeq$ Hesychius α 7634 ; cf. Phrynicus \italique{Praep. soph} 28.1 (ἄσβολος fem.)}\apparat{245}{\lecon{φεῦ} R, Suid. α 4132 : φῦ Dindorf, Wilamowitz \italique{Lysistr.} v.294 \ligneApp{246} γεγέννημαι R : \lecon{γεγένημαι} Suid. τ 903}\\
αἰθὸς γεγένημαι πάντα τὰ περὶ τὴν τράμιν.\testimonia{246}{Suidas τ 903 (sed τὸν pro τὴν) \espTest (αἰθὸς γεγένημαι) \Eustathe 83.6, 992.60, 1314.15 \espTest (αἰθὸς) \Eustathe 1385.59 \espTest (τράμιν) Hesychius τ 1243 ; Photius τ 415}\\\passeLigne}%
\end{groupe}%
%% v. 247
\begin{groupe}%
\personnageGrec{Εὐριπίδης}%
\trimIamb{μὴ φροντίσῃς· ἕτερος γὰρ αὐτὰ σπογγιεῖ.\apparat{247}{\lecon{σπογγιεῖ} R : σφογγιεῖ Elmsley ad \italique{Ach.} 463}%
\commentaire{247}{μὴ φροντίσῃς· ἕτερος γὰρ αὐτὰ σπογγιεῖ.}{Les Grecs se servaient d'éponges pour la toilette et, dans un usage détourné, pouvaient même les utiliser en guise de torche-cul comme cela ressort de la scène des \italique{Grenouilles} (vv. 481 sqq.)
\citComm{Διόνυσος — ἀλλ᾽ ὡρακιῶ.\\
ἀλλ᾽ οἶσε πρὸς τὴν καρδίαν μου σφογγιάν.\\
Ξανθίας — ἰδοὺ λαβέ, προσθοῦ.\\
Διόνυσος — ποῦ 'στιν;
Ξανθίας — ὦ χρυσοῖ θεοὶ\\
ἐνταῦθ᾽ ἔχεις τὴν καρδίαν;\\
Διόνυσος — δείσασα γὰρ\\
ἐς τὴν κάτω μου κοιλίαν καθείρπυσεν.\\
Dionysos — Je me sens défaillir ; apporte-moi une éponge à m'appliquer sur le cœur. — Tiens, prends, applique-là. — Où est-elle ? — Dieux d'or ! c'est là que tu loges ton cœur ? — La peur l'a fait descendre dans mon bas-ventre.}%
Retournons à notre vers : qui est cet ἕτερος ? Réponse : le voisin d'à côté : « La situation est bien celle qu'on attend : en cas d'incendie, les voisins, soucieux de leur propre intérêt, commençaient par faire la chaîne pour apporter de l'eau. Puis, le feu une fois éteint, il se trouvait presque toujours un voisin complaisant pour aider le propriétaire (ou le locataire) malchanceux à nettoyer sa maison et à y remettre de l'ordre. » (Taillardat\marqueur{1}.)%
%\eff{Ranae Ἐγώ;  (486)
%	Πῶς δειλὸς ὅστις σπογγιὰν ᾔτησά σε;
%	Οὐ τἂν ἕτερός γ᾽ αὔτ᾽ εἰργάσατ᾽ ἀνήρ.}
}%
\debutNote{247}%
\note{1}{Jean Taillardat, «~Aristophane, \italique{Thesmophories} v. 247, 289, 721 sqq. : essai d'interprétation du texte transmis~», in \italique{Revue des Études Grecques}, t. 85, fasc. 406-408, juillet-décembre 1972, p. XVII.
Le grand Taillardat avait donc compris le sens de ce passage avant même que le \nieme{242} vers, lacunaire, ne fût rétabli par la découverte d'un papyrus. 
Il n'en demeure pas moins encore énigmatique : dans les \italique{Grenouilles}, Dionysos s'écrie au vers 487, indigné qu'on l'ait pu traiter de lâche : Πῶς δειλὸς ὅστις σπογγιὰν ᾔτησά σε; / Οὐ τἂν ἕτερός γ᾽ αὔτ᾽ εἰργάσατ᾽ ἀνήρ. L'éponge et le courage semblent proverbialement liés (cf. μελάμπυγος). }%
\\\passeLigne}%
\end{groupe}%
%% v. 248
\begin{groupe}%
\personnageGrec{Μνησίλοχος}%
\trimIamb{οἰμώξετ᾽ ἆρ᾽ εἴ τις τὸν ἐμὸν πρωκτὸν πλυνεῖ.\testimonia{248}{(πλυνεῖ) Moeris π 79}\apparat{248}{οἰμώζετ᾽ R : \lecon{οἰμώξετ᾽} Zanetti \espApp εἰς R : \lecon{εἴ τις} Brunck}%
\commentaire{248}{οἰμώξετ᾽ ἄρ᾽ εἴ τις τὸν ἐμὸν πρωκτὸν πλυνεῖ.}{πλύνειν signifie au propre \italique{nettoyer} (et fait donc écho au σπογγιεῖ du vers précédent) et par suite \italique{dire du mal} (\italique{passer un savon} en français ; l'on manque toujours de la délicatesse requise lorsqu'on s'occupe de la toilette d'autrui).\\
Contrairement à ce qu'ont affirmé nombre de commentateurs \& traducteurs, ce n'est pas au sens dérivé mais bien au sens propre qu'il faut prendre ici πλύνειν. (On se demande bien pourquoi certains médisants s'aviseraient de dire du mal du derrière de Mnésiloque et plus encore pourquoi ce dernier, qui n'est pas un être d'une grande finesse, s'en offusquerait.)
En fait, Mnésiloque souhaite simplement bien du courage à qui voudrait \italique{s'attaquer} -- au prime sens de πλύνειν : \italique{commencer de nettoyer} -- à son fondement, dont la saleté égale celle des écuries d'Augias, et dont le \italique{travail} -- tel celui de l'hydre de Lerne -- serait à recommencer sitôt abattu !..
Aristophane fait allusion à un proverbe qu'il cite plus explicitement au vers 604 des \italique{Guêpes} -- le proverbe ou l'un de ses \italique{permutants métaphoriques}\marqueur{1} :%
\citComm{πρωκτὸς λουτροῦ περιγιγνόμενος\\
\francais{Le c.. l'emporte sur l'eau lustrale}}%
Bdélycléon lançait cette sentence à la barbe de Philocléon, saoul de discours et qui, quelque effort qu'il fît pour soigner sa glossolalite, n'en retomberait pas moins un jour dans son travers.
Le proverbe devait donc peu ou prou signifier qu'un naturel grossier, le πρωκτός, retournait inévitablement à sa fange ; en français : un \italique{indécrottable}.}%
\debutNote{248}%
\note{1}{Il est repris tel quel par les parémiographes, Hésychius (π 4144), Photius (π 1450), Suidas (π 2951). La source principale doit être d'anciens commentaires au vers 604 des \italique{Guêpes}. (Cf. les scholies aux manuscrits d'Aristophane ; on y apprend que Καλλίστρατος signalait l'expression comme proverbiale.) L'exégèse est de mon cru ; les anciens commentateurs n'ayant semble-t-il pas compris le sens du proverbe.}\\\passeLigne}%
\end{groupe}%
%% v. 249-251
\begin{groupe}%
\personnageGrec{Εὐριπίδης}%
\trimIamb{Ἀγάθων, ἐπειδὴ σαυτὸν ἐπιδοῦναι φθονεῖς,\\
ἀλλ᾽ ἱμάτιον γοῦν χρῆσον ἡμῖν τουτῳὶ\apparat{250}{τουτωί R : \lecon{τουτῳὶ} edd. plerique}\\
καὶ στρόφιον· οὐ γὰρ ταῦτά γ᾽ ὡς οὐκ ἔστ᾽ ἐρεῖς.\\\passeLigne}%
\end{groupe}%
%% v. 252a
\begin{groupe}%
\personnageGrec{Ἀγάθων}%
\trimIamb{λαμβάνετε καὶ χρῆσθ᾽· οὐ φθονῶ.\\\passeLigne}%
\end{groupe}%
%% v. 252b
\begin{groupe}%
\personnageGrec{Μνησίλοχος}%
\trimIamb{\ligne{252}\decalage{λαμβάνετε καὶ χρῆσθ᾽· οὐ φθονῶ.}τί οὖν λάβω;\apparat{252}{(Εὐρ.) R : \lecon{(Μνη.)} Bentley}\\\passeLigne}%
\end{groupe}%
%% v. 253
\begin{groupe}%
\personnageGrec{Εὐριπίδης}%
\trimIamb{ὅ τι; τὸν κροκωτὸν πρῶτον ἐνδύου λαβών.\apparat{253}{(Ἀγάθ.) R : \lecon{(Εὐρ.)} Dobree}\\\passeLigne}%
\end{groupe}%
%%% v. 254
\begin{groupe}%
\personnageGrec{Μνησίλοχος}%
\trimIamb{νὴ τὴν Ἀφροδίτην ἡδύ γ᾽ ὄζει ποσθίου.\testimonia{254}{Suidas π 2114 \espTest (ποσθίου) Photius π 1113 πόσθιον· αἰδοῖον· ἢ (sic) τὴν Ἀφροδίτην ; Hesychius π 3096}\apparat{254}{(Εὐρ.) R}\\\passeLigne}%
\end{groupe}%
%% v. 255a
\begin{groupe}%
\personnageGrec{Ἀγάθων}%
\trimIamb{σύζωσον ἀνύσας.\\\passeLigne}%
\end{groupe}%
%% v. 255b
\begin{groupe}%
\personnageGrec{Εὐριπίδης}%
\trimIamb{\ligne{255}\decalage{σύζωσον ἀνύσας.}αἶρε νῦν στρόφιον.\apparat{255}{(Ἀγάθ.) et dicolon ante αἶρε praefixit R}%
\commentaire{255}{σύζωσον ἀνύσας --- αἶρε νῦν στρόφιον.}{La répartition de ce vers comme de tout ce passage est très changeante d'un éditeur à un autre.
Pour ma part, j'ai respecté scrupuleusement le découpage des vers du manuscrit de Ravenne, sans tailler ni fusionner çà et là des répliques qui ne l'étaient pas originellement.
J'ai confié à Euripide les conseils d'habillement, Agathon continuant d'adopter une attitude froide et se contentant de tendre à Mnésiloque les vêtements pour ne pas que ce rustre aille directement puiser dans sa garde-robe.\\
La seule question qui se soit véritablement posée est celle de la première partie du vers 255 : σύζωσον ἀνύσας --- αἶρε νῦν στρόφιον.\\
Il comporte un \italique{dicôlon} entre ἀνύσας et αἶρε sans pour autant indiquer de changement de locuteur... Ce \italique{dicôlon} est-il là pour signifier une simple pause le temps que Mnésiloque passe la crocote ? ou y a-t-il véritablement changement de locuteur ?
`σύζωσον ἀνύσας' me semble être une réplique au ton sec, marquant la pointe d'agacement qu'Agathon aurait pu sentir à la remarque de Mnésiloque du vers 254. Aussi lui ai-je accordé cette réplique ; puis Euripide de demander le \italique{strophion} (αἶρε νῦν στρόφιον) et Agathon de le tendre à Mnésiloque (ἰδού).}\\\passeLigne}%
\end{groupe}%
%% v. 255c
\begin{groupe}%
\personnageGrec{Ἀγάθων}%
\trimIamb{\ligne{255}\decalage{σύζωσον ἀνύσας. αἶρε νῦν στρόφιον.}ἰδού.\apparat{255-6}{(Εὐρ.) ἰδού — (Μνη.) ἴθι ... σκέλη R}\\\passeLigne}%
\end{groupe}%
%% v. 256
\begin{groupe}%
\personnageGrec{Μνησίλοχος}%
\trimIamb{ἴθι νυν κατάστειλόν με τὰ περὶ τὼ σκέλει.\apparat{256}{σκέλη R : \lecon{σκέλει} van Leeuwen}\\\passeLigne}%
\end{groupe}%
%% v. 257a
\begin{groupe}%
\personnageGrec{Εὐριπίδης}%
\trimIamb{κεκρυφάλου δεῖ καὶ μίτρας.\\\passeLigne}%
\end{groupe}%
%% v. 257b-258
\begin{groupe}%
\personnageGrec{Ἀγάθων}%
\trimIamb{\ligne{257}\decalage{κεκρυφάλου δεῖ καὶ μίτρας.}ἡδὶ μὲν οὖν\testimonia{257-8}{(ἡδὶ -- φορῶ) Suidas π 1165 (sed ἥδιον pro ἡδὶ)}\\
κεφαλὴ περίθετος, ἣν ἐγὼ νύκτωρ φορῶ.\testimonia{258}{Photius π 690 περίθετος κεφαλή ; \Eustathe 1280.63 \espTest (κεφαλὴ περίθετος) Pollux II 35}%
\commentaire{258}{κεφαλὴ περίθετος, ἣν ἐγὼ νύκτωρ φορῶ.}{De tout l'attirail nécessaire pour se métamorphoser en femme, il est étonnant que ce soit le plus évident à priori qui ait rencontré le plus de difficultés : c'est Fritzsche qui le premier s'avisa que la \grec{κεφαλὴ περίθετος} désignait une \italique{perruque}. Hé oui ! car une résille et un bandeau ne suffiraient à cacher l'absence volumique d'une belle crinière, d'où la suggestion avisée d'Agathon – en homme habitué à de tels travestissements. Ce sens est clairement attesté par Pollux (X, 170) :%
\citComm{\grec{πηνήκην δὲ καὶ προκόμιον, καὶ ὡς Ἀριστφάνης ἐν Γηρυτάδῃ περίθεσιν· καὶ περιθέτην δέ, ὡς ἐν Ἄμφιδος Ἀλκμαίωνι καὶ Μενάνδρου Ὀλυνθίᾳ.}}%
Les synonymes \grec{πηνήκη} et \grec{προκόμιον} ne souffrent aucune ambiguïté.\\
Le tacle d'Aristophane envers les mœurs agathoniennes – qui lui fait préciser : ἣν ἐγὼ νύκτωρ φορῶ – n'est pas clair... le commentaire à l'imagination débridée du scholiaste, qui n'est peut-être pas plus renseigné que nous sur la question, voyant là une allusion aux chasses masquées contre-nature de notre poètereau dans l'interlope \italique{Athens by night}, ne trouve qu'un faible écho au vers 35 (et s'appuie peut-être sur une transmission fautive du texte, comme nous en avons fait l'hypothèse \italique{ad loc.}).}\\\passeLigne}%
\end{groupe}%
%% v. 259
\begin{groupe}%
\personnageGrec{Εὐριπίδης}%
\trimIamb{νὴ τὸν Δί᾽ ἀλλὰ κἀπιτηδεία πάνυ.\\\passeLigne}%
\end{groupe}%
%%% v. 260a
\begin{groupe}%
\personnageGrec{Μνησίλοχος}%
\trimIamb{ἆρ᾽ ἁρμόσει μοι;\apparat{260}{ἠραρμοσηι R : \lecon{ἆρ᾽ ἁρμόσει} Küster}\\\passeLigne}%
\end{groupe}%
%%% v. 260b-261a
\begin{groupe}%
\personnageGrec{Εὐριπίδης}%
\trimIamb{\ligne{261}\decalage{ἆρ᾽ ἁρμόσει μοι;}νὴ Δί᾽ ἀλλ᾽ ἄριστ᾽ ἔχει.\apparat{260}{(Ἀγάθ.) R : \lecon{(Εὐρ.)} Ellebodius}\\
φέρ᾽ ἔγκυκλον·\testimonia{261}{Suidas ε 135 (sed τουτὶ pro τι) \espTest (ἔγκυκλον -- κλινίδος) Suidas κ 1809}%
\commentaire{261}{φέρ᾽ ἔγκυκλον.}{Se reporter à la note au vers 643 pour un développement sur l'\italique{encycle}. Ce mot désigne une sorte de \italique{jupe courte}.}\\\passeLigne}%
\end{groupe}%
%%% v. 261b
\begin{groupe}%
\personnageGrec{Ἀγάθων}%
\trimIamb{\ligne{261}\decalage{φέρ᾽ ἔγκυκλον·}τουτὶ λάβ᾽ ἀπὸ τῆς κλινίδος.\apparat{261}{τοῦτὶ R : \lecon{τουτὶ} Suid. ε 135, κ 1809 \espApp λάμβαν᾽ R, Suid. : \lecon{λάβ᾽} Daubuz, Küster}\\\passeLigne}%
\end{groupe}%
%%% v. 262a
\begin{groupe}%
\personnageGrec{Εὐριπίδης}%
\trimIamb{ὑποδημάτων δεῖ.\\\passeLigne}%
\end{groupe}%
%%% v. 262b
\begin{groupe}%
\personnageGrec{Ἀγάθων}%
\trimIamb{\ligne{262}\decalage{ὑποδημάτων δεῖ.}τἀμὰ ταυτὶ λάμβανε.\\\passeLigne}%
\end{groupe}%
%% v. 263
\begin{groupe}%
\personnageGrec{Μνησίλοχος}%
\trimIamb{ἆρ᾽ ἁρμόσει μοι; χαλαρὰ γ᾽ οὐ χαίρεις φορῶν.\testimonia{263}{Suidas χ 7 χαλαρὰ γοῦν χαίρεις φορῶν ὑποδήματα (sed χαίροις codd. FV)}\apparat{263}{(Ἀγάθ.) χαλαρὰ ... φορῶν R : \lecon{(Μνη.)} Brunck \espApp μοι R : \lecon{μοι;} Küster \espApp \lecon{γ᾽ οὐ} \Rac, Willems : γοῦν \Rpc, Suid. χ 7 \espApp \lecon{χαίρεις} R, Suid. codd. AGM : χαίροις Suid. codd. FV}%
\commentaire{263}{χαλαρὰ γ᾽ οὐ χαίρεις φορῶν.}{Pour l'explication de ce vers, se reporter encore une fois au grand Alphonse Willems (\nieme{6} note, reproduite à la fin) : n'en déplaise aux fouines à ordure, il n'y a là aucune obscénité. (La meilleure preuve en est que si ç'avait été le cas, Agathon aurait superbement ignoré la remarque désobligeante de Mnésiloque.)
%\eff{En fait la correction est la bonne : Agathon est moqué ailleurs (Gérutades 178) comme mh xalara, c'est-à-dire non serré. Évidemment, les allusions scabreuses au trou-du-cul d'Agathon ne sont que pure imagination des fouines à ordure. Zénob. I 2 (cf. Diogenianus I 7)}
%Le sens n'était pourtant pas bien compliqué à saisir : 
De plus, Agathon « n'aurait eu garde de porter des souliers larges, car c'eût été faire preuve de d'inélégance et même de rusticité [...] Mnésiloque, en vieillard qui aime ses aises, est donc fondé à craindre que les souliers ne soient trop étroits pour lui. » À noter qu'une fois encore, le manuscrit de Ravenne était meilleur avant correction qu'après.}\\\passeLigne}%
\end{groupe}%
%% v. 264-265
\begin{groupe}%
\personnageGrec{Ἀγάθων}%
\trimIamb{σὺ τοῦτο γίγνωσκ᾽· ἀλλ᾽ ἔχεις γὰρ ὧν δέει,\\
εἴσω τις ὡς τάχιστά μ᾽ ἐσκυκλησάτω.%
\commentaire{265}{εἴσω τις ὡς τάχιστά μ᾽ ἐσκυκλησάτω.}{De même que l'Euripide des \italique{Acharniens} (vers 409) demande à être roulé par l'eccyclème pour sortir de la scène : il ne lui agrée pas de descendre (καταβαίνειν, cf. \italique{infra} v. 1215) :%
\citComm{ἀλλ' ἐκκυκλήσομαι· καταβαίνειν δ᾽ οὐ σχολή.}%
}\\\passeLigne}%
\end{groupe}%
%% v. 266-268a
\begin{groupe}%
\ledpb%
\personnageGrec{Εὐριπίδης}%
\trimIamb{ἁνὴρ μὲν ἡμῖν οὑτοσὶ καὶ δὴ γυνὴ\apparat{266}{ἀνὴρ R : \lecon{ἁνὴρ} Fritzsche}\\
τό γ᾽ εἶδος· ἢν λαλῇς δ᾽, ὅπως τῷ φθέγματι\apparat{267}{τόδ᾽ R : \lecon{τό γ᾽} Küster}\\
γυναικιεῖς εὖ καὶ πιθανῶς.\\\passeLigne}%
\end{groupe}%
%% v. 268b
\begin{groupe}%
\personnageGrec{Μνησίλοχος}%
\trimIamb{\ligne{268}\decalage{γυναικιεῖς εὖ καὶ πιθανῶς.}πειράσομαι.\\\passeLigne}%
\end{groupe}%
%% v. 269a
\begin{groupe}%
\personnageGrec{Εὐριπίδης}%
\trimIamb{βάδιζε τοίνυν.\\\passeLigne}%
\end{groupe}%
%% v. 269b-270a
\begin{groupe}%
\personnageGrec{Μνησίλοχος}%
\trimIamb{\ligne{269}\decalage{βάδιζε τοίνυν.}μὰ τὸν Ἀπόλλω οὔκ, ἤν γε μὴ\\
ὀμόσῃς ἐμοί—\\\passeLigne}%
\end{groupe}%
%% v. 270b
\begin{groupe}%
\personnageGrec{Εὐριπίδης}%
\trimIamb{\ligne{270}\decalage{ὀμόσῃς ἐμοί—}τί χρῆμα;\\\passeLigne}%
\end{groupe}%
%% v. 270c-271
\begin{groupe}%
\personnageGrec{Μνησίλοχος}%
\trimIamb{\ligne{270}\decalage{ὀμόσῃς ἐμοί—τί χρῆμα;}συσσώσειν ἐμὲ\\
πάσαις τέχναις, ἤν μοί τι περιπίπτῃ κακόν.\\\passeLigne}%
\end{groupe}%
%%% v. 272
\begin{groupe}%
\personnageGrec{Εὐριπίδης}%
\trimIamb{`ὄμνυμι τοίνυν αἰθέρ᾽ οἴκησιν Διός.'%
\commentaire{272}{ὄμνυμι τοίνυν αἰθέρ᾽ οἴκησιν Διός.}{Selon le scholiaste, ce vers est une citation de la \italique{Mélanippe [la sage]} d'Euripide\marqueur{1} (ἐκ τῆς Μελανίππης Εὐριπίδου) ; citation qui devait être fameuse à la façon dont Aristophane l'évoque encore aux vers 100 et 311 des \italique{Grenouilles}.}%
\debutNote{272}%
\note{1}{\nieme{8} fragment de l'édition Jouan-van Looy aux Belles Lettres ; fr. 487 de l'édition Kannicht ; 670 chez Mette.}\\\passeLigne}%
\end{groupe}%
%% v. 273
\begin{groupe}%
\personnageGrec{Μνησίλοχος}%
\trimIamb{τί μᾶλλον ἢ τὴν Ὑποκράτους ξυνοικίαν;\apparat{273}{\lecon{Ὑποκράτους} R : Ἱπποκράτους Schol.}%
\commentaire{273}{τί μᾶλλον ἢ τὴν Ὑποκράτους ξυνοικίαν;}{Au Ciel de qui se moque Euripide ?! Jurer sur l'Éther ? Et pourquoi pas `τὴν Ὑποκράτους ξυνοικίαν' (c'est-à-dire « par une porcherie ») ?\\
« Les fils d'Hippocratès étaient connus de tout le monde à Athènes dès l'époque des \italique{Nuées} (v. 1001). Car leur père étant fils d'Ariphron, propre frère de Périclès, ils étaient les petits neveux du fameux homme d'État. On leur avait donné le surnom de porcs à cause de leur \italique{balourdise} (ὑομουσία, \italique{Cavaliers}, v. 986). Les textes sont unanimes sur ce point. Athénée (III, 96e) : %
\citComm{ἑξῆς ἐπεισηνέχθη, μητρόπολίς τις ὡς ἀληθῶς οὖσα καὶ μήτηρ τῶν Ἱπποκράτους υἱῶν, οὓς εἰς ὑωδίαν κωμῳδουμένους οἶδα.\\
On servit ensuite une vulve de truie, mère des fils d'Hippocratès, lesquels, je le sais, ont été raillés par les comiques à cause de leur cochonnerie\marqueur{1}.}%
Cela se trouve confirmé par le scholiaste d'Aristophane sur les \italique{Nuées}, v. 1001, et par Suidas. De même Photius (σ 845), au mot σῦς :%
\citComm{ὗς τοὺς Ἱπποκράτους υἱοὺς ἔλεγον καὶ τοὺς Παναιτίου καὶ Μέμνονος εἰς ὑηνίαν κωμῳδοῦντες.\\
On nommait porcs les fils d'Hippocratès ainsi que ceux de Panaitios et de Memnon. Les comiques les raillaient à cause de leur cochonnerie.}%
(Autres textes dans Meineke, \italique{Fr. Com. Gr.}, t. II, p. 477. [Pour une bibliographie complète et mise à jour, voir les \italique{testimonia} de l'édition Theodoridis de Photius à l'adresse σ 845.])
Le logis de famille d'Hippocratès est donc une porcherie. Il est à croire que l'auteur ne se faisait pas faute de prononcer συῳκία. » (Note de Willems.) \\
Je conserve la leçon du manuscrit de Ravenne : \grec{Ὑποκράτους}, sans reprendre la correction communément admise \grec{Ἱπποκράτους}. En effet, le préfixe \grec{ὑ}- sert à construire des composés sur le mot \italique{porc} ; ainsi, pour se borner aux textes cités par Willems, on recense \grec{ὑωδία} que Willems rend par « cochonnerie » (littéralement, la \italique{ressemblance à un porc}) et \grec{ὑηνία}, littéralement \italique{nature de porc}, que Willems traduit encore par « cochonnerie ». 
En prononçant : \grec{τὴν ὑ-ποκράτου-ς-(σ)υ-ῳκίαν}, Mnésiloque insiste doublement sur la porcinité de l'hippocratique maison.
Voilà aussi pourquoi les enfants qui l'habitaient sont plus souvent qualifiés d'« \grec{υ-ἱοί Ἱπποκράτους} » plutôt que par l'expression plus commune « παῖδες Ἱπποκράτους »\marqueur{2}.}%
\debutNote{273}%
\note{1}{Comprendre : \italique{grossièreté}, \italique{lourdeur}. ὗς, le \italique{porc} est pour les Grecs le symbole de l'ignorance, partant, de la sottise. Voir à ce sujet le toujours bon Taillardat, §451. L'adjectif Ὑάντιος servait à désigner les Béotiens, peuples à l'éternelle réputation de lourdeur et d'ignorance, cf. l'adage n° 906 d'Érasme : ἡ βοιωτία ὗς.}%
\note{2}{Le scholiaste des \italique{Nuées} (v. 1001) confirme ce jeu de mots : l'Injuste s'adresse à Phidippide : «~Écoute-le [le Juste], jeune homme... et tu ressembleras aux fils d'Hippocratès~», τοῖς Ἱπποκράτους υἱέσιν εἴξεις. Tous les spectateurs comprenant ὑσίν.}\\\passeLigne}%
\end{groupe}%
%% v. 274
\begin{groupe}%
\personnageGrec{Εὐριπίδης}%
\trimIamb{ὄμνυμι τοίνυν πάντας ἄρδην τοὺς θεούς.\testimonia{274}{(ἄρδην) Hesychius α 7097}\\\passeLigne}%
\end{groupe}%
%% v. 275-276
\begin{groupe}%
\personnageGrec{Μνησίλοχος}%
\trimIamb{μέμνησο τοίνυν ταῦθ᾽, ὅτι ἡ φρὴν ὤμοσεν,\testimonia{275-6}{cf. \Eustathe 443.37}\\
ἡ γλῶττα δ᾽ οὐκ ὀμώμοκ᾽· οὐδ᾽ ὥρκωσ᾽ ἐγώ.
\commentaire{275-6}{ὅτι ἡ φρὴν ὤμοσεν / ἡ γλῶττα δ᾽ οὐκ ὀμώμοκ᾽· οὐδ᾽ ὥρκωσ᾽ ἐγώ.}{Mnésiloque se méfie avec raison d'Euripide-le-Retors, il a notamment en tête le vers 612 de son \italique{Hippolyte}\marqueur{1} : 
\citComm{ἡ γλῶσσ' ὀμώμοχ', ἡ δὲ φρὴν ἀνώμοτος.}%
(La citation d'Euripide comportait déjà 7 `o' (dont 4 d'affilée) ; avec Aristophane, le compteur explose !)}%
\debutNote{275-6}%
\note{1}{Cf. Érasme, \italique{Adages} n° 1441, \latin{Lingua juravit}.}%
\apparat{post 276}{parepigraphe ὀλολύζουσι τε ἱερὸν ὠθεῖται : ολολυ[ \PsiUn}%
\commentaire{post 276}{ὀλολύζουσι τε ἱερὸν ὠθεῖται.}{Le vers 276 est suivi de l'indication scénique : ὀλολύζουσι τε ἱερὸν ὠθεῖται ainsi que d'une κορωνίς en marge, qui marque le changement de scène. 
Le scholiaste confirme : παρεπιγραφή. ἐκκυκλεῖται ἐπὶ τὸ ἔξω τὸ θεσμοφόριον et l'on trouve déjà ολολυ[ sur un fragment de papyrus datant du second siècle après Jésus-Christ.\\
La maison d'Agathon est donc roulée par l'eccyclème au fond pour se transformer en temple des Thesmophores ; est hissé le \italique{signal}, le σημεῖον du vers 278 tandis que le chœur des femmes fait retentir l'ὀλολυγή ; c'est à-dire l'espèce de hurlement sacré qu'Hérodote (IV, 189) faisait remonter aux Libyens.
(Cette ululation dont il est question déjà dans Homère (\italique{Iliade}, VI, 301), et qui retentissait dans les temples en signe tantôt de douleur tantôt de joie.)
}\\\passeLigne}%
\end{groupe}%
%% v. 277-279a
\begin{groupe}%
\personnageGrec{Εὐριπίδης}%
\trimIamb{ἔκσπευδε ταχέως· ὡς τὸ τῆς ἐκκλησίας\testimonia{277-8}{(ὡς -- φαίνεται) Suidas σ 325 (sed Θεσμοφόρῳ pro Θεσμοφορίῳ)}\apparat{277}{\lecon{ἔκσπευδε} R : εἄ· σπεῦδε Maas}\\
σημεῖον ἐν τῷ Θεσμοφορίῳ φαίνεται.\apparat{278}{\lecon{Θεσμοφορίῳ} R, cf. \italique{IG} II\expo{2} 1177.24 ; 2498.12 : Θεσμοφόρῳ Suid. σ 325 : Θεσμοφορείῳ Scaliger}\\
ἐγὼ δ᾽ ἄπειμι.\\
\\\passeLigne}%
\end{groupe}%
%% v. 279b-294
\begin{groupe}%
\personnageGrec{Μνησίλοχος}%
\trimIamb{\ligne{279}\decalage{ἐγὼ δ᾽ ἄπειμι.}δεῦρό νυν ὦ Θρᾷτθ᾽ ἕπου.\\
ὦ Θρᾷττα θέασαι, καομένων τῶν λαμπάδων
\commentaire{280}{Θρᾷττα}{est un nom courant pour les esclaves Thraces (comme dans les \italique{Acharniens}).}\\
ὅσον τὸ χρῆμ᾽ ἀνέρχεθ᾽ ὑπὸ τῆς λιγνύος.
\commentaire{281}{ὑπὸ τῆς λιγνύος.}{Sur le sens de ὑπὸ : \italique{avec}, \italique{accompagné de} (souvent en parlant d'un accompagnement musical. Cf. Photius, \italique{Bibliothèque} 320b : οἱ παλαιοὶ τὴν ὑπό ἀντὶ τῆς μετά πολλάκις ἐλάμβανον), voir la note d'Alphonse Willems n°7, écrite d'une juste colère (et reproduite à la fin).}\\
ἀλλ᾽, ὦ περικαλλεῖ Θεσμοφόρω, δέξασθέ με\testimonia{282}{Photius α 1386 ; \AelDion α 113}\apparat{282}{περὶ καλλῆ R : \lecon{περικαλλεῖ} Kaibel, Hall \& Geldart \espApp Θεσμοφόρῳ R : \lecon{Θεσμοφόρω} Photius α 1386 \ligneApp{283}καὶ δεῦρο πάλιν R : \lecon{καὶ δεῦρο καὶ πάλιν} Bentley}\\
ἀγαθῇ τύχῃ καὶ δεῦρο καὶ πάλιν οἴκαδε.\\\passeLigne}
\end{groupe}%
%% v. 279b-294
\begin{groupe}%
\ledpb%
\trimIamb{ὦ Θρᾷττα τὴν κίστην κάθελε, κᾆτ᾽ ἔξελε\\
τὰ πόπαν᾽, ὅπως λαβοῦσα θύσω ταῖν θεαῖν.\testimonia{285}{(πόπαν᾽) \Eustathe 437.4}\apparat{285}{\lecon{τὰ πόπαν᾽} \PsiUn, Porson : τὸ πόπανον R}\\
δέσποινα πολυτίμητε Δήμητερ φίλη\\
καὶ Φερρέφαττα, πολλὰ πολλάκις μέ σοι\apparat{287}{\lecon{Φερρέφαττα} \PsiUn, Thiersch : Φερσέφαττα R}\\
θύειν ἔχουσαν, εἰ δὲ μή, ἀλλὰ νῦν λαθεῖν%
\commentaire{288}{εἰ δὲ μή, ἀλλὰ νῦν λαθεῖν.}{Mnésiloque, qui déclame sa prière à haute voix afin de mystifier l'assemblée des femmes, prononce en aparté cette réplique.}\\
καὶ τὴν θυγατέρα χοιρίον ἀνδρός μοι τυχεῖν\apparat{289}{\lecon{καὶ τὴν θυγατέρα} R : del. Willems \espApp χοῖρον R : \lecon{χοιρίον} Fritzsche ap. Schol.}\\
πλουτοῦντος, ἄλλως τ᾽ ἠλιθίου κἀβελτέρου,\\
καὶ πρὸς θάλικα νοῦν ἔχειν μοι καὶ φρένας.\apparat{291}{πρὸς θάληκον R : Ποσθαλίσκον Dindorf : \lecon{πρὸς θάλικα} Bibi}%
\commentaire{291}{καὶ πρὸς θάληκον νοῦν ἔχειν μοι καὶ φρένας.}{Le vers 291 a fait l'objet de nombreux débats ; et pour cause : le mot θάληκον (v. 291) est inconnu au bataillon des lexiques !..
Que signifie donc ce `θάληκον', inconnu autrement ?\\
Une scholie en marge du manuscrit de Ravenne se référant probablement au \nieme{291} vers indique :%
\citComm{τὸν παιδαρίσκον\marqueur{1}. ἴσως δὲ παρὰ τὴν πόσθην αὐτὸ συνέθηκεν,\\
le jeune garçon. Peut-être à rapprocher du pénis (πόσθη). }%
À l'unanimité, les commentateurs s'en sont remis à cette scholie pour en induire que πρὸς θάληκον était une transmission fautive d'un surnom de garçon, forgé sur le mot πόσθη, \italique{le membre}\marqueur{2}. 
La reconstruction encore la plus en vogue (et la plus proche du texte transmis) est celle de Dindorf : πρὸς θάληκον $ \to $ Ποσθαλίσκον, où Ποσθαλίσκον désignerait le fils (fictif) de Mnésiloque. 
À ce ποσθαλίσκον, Willems a opposé deux raisons\marqueur{3}, que je crois fortes et fondées : « il est plus que douteux que ce dernier mot [Ποσθαλίσκον] soit grec, les diminutifs se faisant en ισκος et non en αλισκος. (On cite, il est vrai, κωραλίσκος, παιδαρίσκος et μειρακίσκος. Ce serait concluant s'il n'y avait pas l'intermédiaire de κωράλιον, παιδάριον et μειράκιον\marqueur{4}. Si des diminutifs comme ποσθάλιον et ποσθαλίσκος avaient existé, ils n'auraient pu manquer d'être relevés par d'anciens lexicographes\marqueur{5}.) »
À l'école de Willems, j'ai aussi appris à me méfier du scholiaste\marqueur{6}, cet « ignare byzantin », et faire bien davantage confiance à l'« alexandrin » Hésychius. Autres lieux, autres temps et autres sciences !\\
Comme l'a déjà fait remarquer Alphonse Willems, Hésychius n'a jamais compulsé de ποσθαλίσκον ou autre composé de πόσθη pouvant expliquer notre πρὸς θάληκον. 
En revanche, l'on trouve chez Hésychius à l'adresse θ 50 :
\citComm{θάλλικα· σάκκου εἶδος\\
Thallika : sorte de sac.}
Substituer θάλικα à θάληκον semble à peine être une correction par rapport à ce qui était proposé jusqu'ici\marqueur{7}.
(La glose d'Hésychius sur θαλλικοποιοί laisse à penser que le sac en question était fait de cuir et qu'il ressemblait à un petit coffret.)\\
Pour le sens, j'ai reçu l'aide précieuse de Jean-Victor Vernhes qui « estime qu'il s'agit du sac où cette femme met des bijoux, la clef de sa maison, ses emplettes, sa carte de crédit (et bientôt, vous allez voir, celle de son mari), et qu'il symbolise d'une façon générale, ses affaires, ses intérêts, au profit desquels elle exploite son nigaud de conjoint. Elle finira par tout fourrer dans son sac ! »
Comprendre ce θάλλικα comme une métonymie du portefeuille du ménage\marqueur{8} (ou, pour rester plus proche de ce à quoi devait ressembler un véritable θάλιξ, une \italique{cassette}), portefeuille que la fille devra gérer avec économie et esprit contrairement à ces idiots d'hommes, me paraît être tout-à-fait dans le ton d'Aristophane ! 
(Thème qui lui cher : confier les clefs de la cité tout autant que celles du ménage à des femmes, femmes qui réussissent autrement mieux que les hommes à cette tâche ! Il y a de nombreux passages de l'\italique{Assemblée des Femmes} ou de discours de Lysistrata qui pourraient faire écho à ce genre de prière !)\\
Un point reste à élucider : comment expliquer la scholie : τὸν παιδαρίσκον. ἴσως δὲ παρὰ τὴν πόσθην αὐτὸ συνέθηκεν ? 
Si l'on admet que le scholiaste avait devant lui θάλικα ou θάλικον, et qu'il l'analysait non comme un dérivé de θάλιξ, qu'il ignorait, mais comme un mot formé sur τό θάλος, \italique{la jeune pousse} et, dans un sens métaphorique, \italique{rejeton}, \italique{enfant}, suivi d'un suffixe en ικος, l'on peut comprendre le quiproquo, facilité par l'évocation deux vers plus haut de la fille de Mnésiloque. De toute façon l'hésitation du scholiaste (ἴσως δὲ) interdit de penser qu'il ait pu avoir devant lui une évocation aussi claire de la πόσθη que Ποσθαλίσκον.\\
Lisons donc le vers 291 ainsi :
\citComm{καὶ πρὸς θάλικα νοῦν ἔχειν μοι καὶ φρένας.\\
\francais{et m'avoir l'esprit et l'intelligence appliqués à la cassette.}}
Je voudrais à présent aborder une question en rapport avec ce passage, celle de notre anti-héros : Mnésiloque. Dans les indications de personnage du manuscrit de Ravenne, le protagoniste des \italique{Thesmophories} est indifféremment nommé Μνησίλοχος ou κηδεστής, le \italique{parent par alliance}. 
Parent par alliance, il l'est assurément\marqueur{9} !.. mais est-il vraiment ce Mnésiloque que les \italique{Vies} d'Euripide\marqueur{10} présentent comme le père de la première femme d'Euripide, Χοιρίνη ou Χοιρίλη\marqueur{11} ? Si tel est le cas, alors il est certain qu'Aristophane ne pouvait se priver de cette prière factice de Mnésiloque à sa fille pour tacler Euripide ! 
Faut-il supposer que l'acteur mâchait `χοιρίον' au vers 289 pour faire entendre Χοιρίν'\marqueur{12} ? Ou, au contraire, faut-il arguer de l'absence d'allusion à Euripide dans ce passage pour aller dans le même sens que la majorité des éditeurs modernes, pour qui la précision de l'identité du κηδεστής est un ajout postérieur à Aristophane ? La question, pour moi, n'est toujours pas tranchée ; qu'un lecteur éclairé m'éclaire !}%
\debutNote{291}%
\note{1}{Le manuscrit de Ravenne, qu'on abrège `R', comporte en fait : πεδαρίσκον.}%
\note{2}{Comme le terme de caresse pour la jeune fille, χοιρίον, peut désigner – dans un autre contexte – le sexe féminin. Joie mauvaise des philologues qui ont ainsi pu laisser folâtrer leur imagination et plaquer leur saleté morale sur Aristophane. Juste courroux d'Alphonse Willems (\nieme{8} note) devant ces reconstructeurs de pièce qui nous inventent un personnage usant d'obscénités pour parler de sa propre fille (quand il ne lui souhaite pas les mœurs les plus infamantes), et ce au beau milieu d'une prière aux deux déesses parmi les plus révérées !! (Allez donc chercher dans tout Aristophane la moindre irrévérence à Déméter ou Coré.) On ne le répétera jamais assez : « Mais pour qui prennent-ils Aristophane ? » Las ! on ne détruit que ce que l'on remplace et le coup d'épée de Willems – qu'il avait aiguisée – avait porté dans l'eau car sa note était essentiellement aporétique. Prions les déesses thesmophores de conférer à la nôtre assez de force pour dissiper à tout jamais ces nuées !}%
\note{3}{Tome II de son édition posthume d'Aristophane (1919), pp. 548-549. Édition malheureusement désormais introuvable mais en accès libre chez les grandes bibliothèques américaines numériques – à condition de se faire passer pour américain.
Cette note sur ce passage des \italique{Thesmophories}, la \nieme{8} qu'il écrivit, est reproduite en annexe.}%
\note{4}{Chantraine, dans sa \italique{Formation des noms en grec ancien} aux pages 406 et suivantes, parle longuement du suffixe -ισκος sans même évoquer un possible -αλισκος ou -αρισκος ; Chantraine cite παιδαρίσκος et μειρακίσκος en expliquant que le Grec recourait à ce dernier diminutif parce que l'autre, μειράκιον, était devenu d'usage banal.}%
\note{5}{Willems pense notamment à Hésychius, qui semble avoir fait un dépouillement complet d'Aristophane (ou, du moins, a dépouillé le commentaire que Didyme en a fait, cf. P. Boudreaux, \italique{Le texte d'Aristophane et ses commentateurs}), à une époque où le texte transmis était bien meilleur que le nôtre, Hésychius étant lui-même un érudit comme il ne peut plus en exister ; il avait probablement lu davantage de grec que nous n'en avons conservé. Willems pense peut-être aussi au lexique, plus tardif, de l'hypothétique Suidas : la \italique{Souda} ; lexique qui fait la part belle à Aristophane et qui est une mine précieuse de renseignements car il s'appuie sur des leçons du texte parfois meilleures, notamment pour notre pièce des \italique{Thesmophories} qui n'est transmise intégralement que par un unicodex, le manuscrit de Ravenne. 
Suidas \& Hésychius n'ont effectivement pas manqué de relever le surnom πόσθων -- \italique{couillon}, traduira-t-on -- qui apparaît dans la \italique{Paix} (vers 1300). Ce surnom a aussi été relevé par le lexicographe Photius. De ποσθαλίσκον, par contre, point de trace !}%
\note{6}{Je parle toujours, par abus de langage et conformément à un usage trop bien établi, \italique{du} scholiaste. Il est bien évident qu'en réalité les commentaires en marge des manuscrits sont de différents auteurs, et ont été compulsés les uns à la suite des autres postérieurement. Certaines scholies font preuve d'une grande érudition et apportent parfois des leçons meilleures que celle du texte transmis, d'autres, par contre, – la plupart, hélas ! – ne sont, au mieux, d'aucun intérêt voire souvent sources d'erreurs pour qui y accorderait trop d'importance. Pour finir de gâter l'affaire, les \italique{Thesmophories} est la pièce d'Aristophane la moins pourvue en scholies.}%
\note{7}{Le remplacement du double lambda `λλ' par un simple `λ' est commandé par les lois de scansion, Aristophane prenant souvent des libertés avec le doublement des consonnes, et Hésychius n'étant pas bien sûr sur ce point (Hésychius, lui aussi, est conservé par un unique manuscrit bien postérieur : qu'on estime du \nieme{XIV} ou \nieme{XV} siècle), on ne peut même pas appeler ça une correction.
Idem pour le iota et le êta qui se prononçaient de la même façon et sont allègrement permutés sans scrupules dans les manuscrits tardifs. Seule la confusion de la terminaison de l'accusatif `α' en `ον' est réellement conjecturale, mais la faute est compréhensible pour un mot rare et certainement inconnu du scribe qui aura vu un suffixe en -ικος. 
Ajoutons qu'une confusion analogue, entre un neutre, singulier en `ον', et pluriel en `α', avait déjà eu lieu au vers 285 : le manuscrit de Ravenne comporte τὸ πόπανον ; le sens poussait Porson en son temps à réclamer la correction τὰ πόπαν(α) ; cette émendation fut ensuite confirmée par la découverte d'un papyrus.}%
\note{8}{ὁ θύλακος désigne en grec le \italique{sac à provisions}.}%
\note{9}{Les vers 74 et 210 l'attestent.}%
\note{10}{Voir, par exemple, la \italique{Souda} à l'entrée Εὐριπίδης.}%
\note{11}{L'épigraphie fait état de plusieurs Χοιρίνη tandis que le nom Χοιρίλη ne semble pas attesté.}%
\note{12}{Je n'exclue d'ailleurs pas que ce fût peut-être la leçon originale ; le iota de `χοιρίον' ayant été ajouté d'après la scholie. Voir le « dossier Mnésiloque » en Introduction.}\\
ποῦ ποῦ καθίζωμ᾽ ἐν καλῷ, τῶν ῥητόρων\testimonia{292}{(τῶν ῥητόρων) Photius ρ 106 ῥήτορα· καὶ τὴν γυναῖκα. Ἀριστοφάνης}\\
ἵν᾽ ἐξακούω; σὺ δ᾽ ἄπιθ᾽ ὦ Θρᾷττ᾽ ἐκποδών.\\
δούλοις γὰρ οὐκ ἔξεστ᾽ ἀκούειν τῶν λόγων.\\\passeLigne}%
\end{groupe}%
%% v. 295-311
\begin{groupe}%
\ledpb%
\personnageGrec{Κορ.-Ἱέρεια}%
\typeVers{prose}\prose{εὐφημία ἔστω, |\apparat{295}{\inRas{Κῆρυξ}{} \Rdeux : \auDessus{κη}{ρυ} \Mon : γυνὴ μιμουμένη κήρυκα Schol. : Κρίτυλλα Sommerstein \espApp εὐφημί᾽ ἔστω R : \lecon{εὐφημία ἔστω} Suid. ε 3797 \ligneApp{297} Θεσμοφόροιν τῇ Δήμητρῃ καὶ τῇ Κόρῃ R : \lecon{Θεσμοφόροιν} Blaydes \ligneApp{298} τῇ Κουροτρόφῳ τῇ Γῇ R : \lecon{τῇ Κουροτρόφῳ} Dobree ex Schol. \ligneApp{301} Χάρισιν R : \lecon{ταῖς Χάρισιν} Meineke}%
\commentaire{295-311}{}{La coryphée-prêtresse commence la prière, en prose selon l'usage : (scholie) :%
\citComm{πεζῇ εὐφημίᾳ χρῶντα οἱ κωμικοὶ, ἐπειδὰν εὐχὴν ἢ ψήφισμα εἰσάγωσιν.}%
Cf. en \italique{Oiseaux} (vers 864-888) pour une prière analogue.}%
\ligne{296} εὐφημία ἔστω. |\\
\ligne{297}εὔχεσθε ταῖν Θεσμοφόροιν |\\
\ligne{299}καὶ τῷ Πλούτῳ καὶ τῇ Καλλιγενείᾳ |\testimonia{299}{(Καλλιγενείᾳ $=$ Δημήτερι) Hesychius κ 472 Καλλιγένειαν ; cf. Photius κ 118 ($=$ fr. Aristophanes 331 Kassel-Austin)}\\
καὶ τῇ Κουροτρόφῳ |\testimonia{300}{(Κουροτρόφῳ $=$ Δημήτηρ) Hesychius κ 3865 κουροτρόφος ; scholium vet. in Hesiodus \italique{Opera et dies} 227-9}\\
\ligne{301}καὶ τῷ Ἑρμῇ καὶ ταῖς Χάρισιν |\\
\ligne{302}ἐκκλησίαν τήνδε καὶ σύνοδον τὴν νῦν |\\
\ligne{303}κάλλιστα καὶ ἄριστα ποῆσαι, |\\
\ligne{304}πολυωφελῶς μὲν τῇ πόλει τῇ Ἀθηναίων, |\apparat{304}{πόλει R : \lecon{τῇ πόλει} Meineke}\\
τυχηρῶς δ᾽ ἡμῖν αὐταῖς. |\apparat{305}{\lecon{ἡμῖν} R : ὑμῖν Fritzsche}\\
\ligne{306}καὶ τὴν δρῶσαν καὶ ἀγορεύουσαν |\apparat{306}{\expo{καὶ}τὴν ἀγορεύουσαν \Rdeux}\\
\ligne{307}τὰ βέλτιστα περὶ τὸν δῆμον τὸν Ἀθηναίων |\apparat{307}{τῶν Ἀθηναίων R : \lecon{τὸν Ἀθηναίων} Grynaeus}\\
\ligne{309}καὶ τὸν τῶν γυναικῶν |\ligne{309} ταύτην νικᾶν. |
\commentaire{307-8}{τὸν δῆμον τῶν Ἀθηναίων καὶ τὸν τῶν γυναικῶν.}{Plaisanterie \italique{contre l'attente} (παρὰ προσδοκίαν, comme on dit) : « Le peuple des Athéniens et celui... des femmes »... là où l'on attendait \grec{καὶ τὸν τῶν ξυμμάχων}, « et celui de ses alliés ».}\\
ταῦτ᾽ εὔχεσθε καὶ ὑμῖν αὐταῖς τἀγαθά. |\apparat{310}{\lecon{ὑμῖν} R : ἡμῖν \Mon}\\
ἰὴ παιὼν ἰὴ παιὼν ἰὴ παιών. χαίρωμεν.\\\passeLigne}%
\end{groupe}%
%% v. 312-330
\begin{groupe}%
\ledpb%
\personnageGrec{Χορός}%
\typeVers{$2$ iamb.}\chantChoeur{\ligne{312}δεχόμεθα καὶ θεῶν γένος\apparat{312}{δεχόμεσθα R : \lecon{δεχόμεθα} Dindorf}%
\commentaire{312-30}{Chant du chœur.}{Pour l'étude métrique de ce chant du chœur, l'on pourra consulter Lætitia Parker et son ouvrage de référence, \italique{The Songs of Aristophanes}\marqueur{1}. Elle parvient à faire apparaître comme une sorte de refrain scandé $- \cup \cup - \cup \cup - \cup \cup -$\marqueur{2} aux vers 315, 320 et 325.\\
Quoique la découverte soit extrêmement intéressante et encore à creuser, j'ai préféré pour ma part reprendre la colométrie de Prato, plus proche du manuscrit de Ravenne : «~\italien{Il Coro accoglie l'invito dell'Aralda e rivolge la sua preghiera alle divinità peculiari dell' ecclesia (maschile) escluse nella precedente invocazione in prosa e qui menzionate, non senza adattamenti, col consueto criterio selettivo, quasi per conferire un crisma di legalità maggiore alle decisioni finali.}~» (Prato.)}%
\debutNote{312-30}%
\note{1}{Lætitia P. E. Parker, \italique{The songs of Aristophanes}, Clarendon Press Oxford (1997), pp. 406-409.}%
\note{2}{Lætitia désigne ce vers comme un \italique{énoplien} (dactylo-épitrite prolongé) et Prato comme un \italique{alcméen}.}}%
\end{groupe}%
%% v. 312-330
\begin{groupe}%
\typeVers{$2$ iamb\cat}\chantChoeur{λιτόμεθα ταῖσδ᾽ ἐπ᾽ εὐχαῖς\apparat{313}{λιτόμεσθα R : \lecon{λιτόμεθα} Bourdin}}%
\end{groupe}%
%% v. 312-330
\begin{groupe}%
\typeVers{$2$ iamb\cat}\chantChoeur{φανέντας ἐπιχαρῆναι.}%
\end{groupe}%
%% v. 315
\begin{groupe}%
\typeVers{$2$ adon.}\chantChoeur{Ζεῦ μεγαλώνυμε χρυσολύρα τε}%
\end{groupe}%
%% v. 316
\begin{groupe}%
\typeVers{cr. $+$ choria.}\chantChoeur{Δῆλον ὃς ἔχεις ἱεράν,}%
\end{groupe}%
%% v. 317
\begin{groupe}%
\typeVers{lecyth.}\chantChoeur{καὶ σὺ παγκρατὲς κόρα\apparat{317}{παμκρατὲς R : \lecon{παγκρατὲς} ed. pr. \espApp κόρᾳ R : \lecon{κόρα} \Mon}}%
\end{groupe}%
%% v. 318
\begin{groupe}%
\typeVers{$2$ iamb.}\chantChoeur{γλαυκῶπι χρυσόλογχε πόλιν}%
\end{groupe}%
%% v. 319
\begin{groupe}%
\typeVers{$3$ iamb\cat}\chantChoeur{οἰκοῦσα περιμάχητον, ἐλθὲ δεῦρο.}%
\end{groupe}%
%% v. 320
\begin{groupe}%
\typeVers{alcm.}\chantChoeur{καὶ πολυώνυμε θηροφόνη\apparat{320}{θηροφ\inRas{ω}{ό}νε παῖ-- \Rpc : \lecon{θηροφόνη} Hermann}}%
\end{groupe}%
%% v. 321
\begin{groupe}%
\typeVers{paroem.}\chantChoeur{Λατοῦς χρυσώπιδος ἔρνος,}%
\end{groupe}%
%% v. 322
\begin{groupe}%
\typeVers{paroem.}\chantChoeur{σύ τε πόντιε σεμνὲ Πόσειδον\testimonia{322-4}{(σύ τε -- οἰστροδόνητον) Suidas οι 181 (sed σεμνὲ et ἁλιμέδον omissis)}}%
\end{groupe}%
%% v. 323
\begin{groupe}%
\typeVers{cret.}\chantChoeur{ἁλιμέδον}%
\end{groupe}%
%% v. 324a
\begin{groupe}%
\decoupeVers{a}\typeVers{paroem.}\chantChoeur{προλιπὼν μυχὸν ἰχθυόεντα,\apparat{324}{ἰχθυόεντ᾽ R, Suid. οι 181 : \lecon{ἰχθυόεντα} Wilamowitz \ligneApp{325} ἐναλίου R : ἐνάλιοί Fritzsche : \lecon{εἰναλίου} Bentley \ligneApp{326} \lecon{ὀρείπλαγκτοι} R : ὀρίπλαγκτοι Austin}}%
\end{groupe}%
%% v. 324b
\begin{groupe}%
\ligne{324}\decoupeVers{b}\typeVers{adon.}\chantChoeur{οἰστροδόνητον}%
\end{groupe}%
%% v. 325
\begin{groupe}%
\typeVers{alcm.}\chantChoeur{Νηρέος εἰναλίου τε κόραι}%
\end{groupe}%
%% v. 326
\begin{groupe}%
\typeVers{pros.}\chantChoeur{Νύμφαι τ᾽ ὀρείπλαγκτοι.}%
\end{groupe}%
%% v. 327
\begin{groupe}%
\typeVers{ith.}\chantChoeur{χρυσέα τε φόρμιγξ}%
\end{groupe}%
%% v. 328
\begin{groupe}%
\typeVers{ba. $+$ adon.}\chantChoeur{ἰαχήσειεν ἐπ᾽ εὐχαῖς}%
\end{groupe}%
%% v. 329
\begin{groupe}%
\typeVers{2 hemiepes}\chantChoeur{ἡμετέραις· τελέως δ᾽ ἐκκλησιάσαιμεν Ἀθηνῶν\apparat{329}{Ἀθηναίων R : \lecon{Ἀθηνῶν} Reisig}}%
\end{groupe}%
%% v. 330
\begin{groupe}%
\typeVers{ith.}\chantChoeurFr{εὐγενεῖς γυναῖκες.
\commentaire{330}{εὐγενεῖς γυναῖκες.}{Confirme, s'il était besoin, que seules les citoyennes athéniennes assistaient aux Thesmophories (cf. le vers 294 lorsque Mnésiloque renvoie Thratta).}\\\passeLigne}%
\end{groupe}%
%% v. 331-351
\begin{groupe}%
\ledpb%
\personnageGrec{Κορυφαῖος}%
\typeVers{trim. iamb.}\trimIamb{\ligne{331}εὔχεσθε τοῖς θεοῖσι τοῖς Ὀλυμπίοις\apparat{331}{(Κηρ.) R : \lecon{(Κορ.)} van Leeuwen}\\
καὶ ταῖς Ὀλυμπίαισι καὶ τοῖς Πυθίοις\\
καὶ ταῖσι Πυθίαισι καὶ τοῖς Δηλίοις\\
καὶ ταῖσι Δηλίαισι τοῖς τ᾽ ἄλλοις θεοῖς,\apparat{334}{καὶ τοῖς R : \lecon{τοῖς τ᾽} Grynaeus}\\
εἴ τις ἐπιβουλεύει τι τῷ δήμῳ κακὸν\testimonia{335-7}{Satyrus \italique{Vit. Euripid.} ($=$ \italique{The Oxyrhynchus Papyri} IX 1176) fr. 39 col. XII 8-16 (Εὐριπιδίῳ pro Εὐριπίδῃ (?))}\\
τῷ τῶν γυναικῶν ἢ 'πικηρυκεύεται\testimonia{336}{('πικηρυκεύεται) Photius ε 1580 ἐπικηρυκεύεται $=$ \Syn ε 682 $=$ Suidas ε 2375 $=$ \EtymMagnum 360.21 $=$ \EtymGen (B) ; cf. Hesychius ε 4853}\\
†Εὐριπίδῃ Μήδοις τ᾽† ἐπὶ βλάβῃ τινὶ\apparat{337}{ευριπιδιω\douteux{ι}[. .]\douteux{υ}τ \PsiCinq : Εὐριπίδῃ Μήδοις R : Εὐριπίδῃ Μήδοις τ᾽ Biset}% Biset : \lecon{Εὐριπίδῃ ᾧ ταῦτ᾽} Bibi}%
\commentaire{336-7}{ἢ 'πικηρυκεύεται / Εὐριπίδῃ Μήδοις ἐπὶ βλάβῃ τινὶ.}{Dans sa parodie d'ecclésie, Aristophane ne manque pas de reprendre les imprécations contre les Mèdes, qui étaient de règle comme on le sait par Isocrate (\italique{Panégyrique}, 157) :
\citComm{Καὶ πολλῶν μὲν οἱ πατέρες ἡμῶν μηδισμοῦ θάνατον κατέγνωσαν, ἐν δὲ τοῖς συλλόγοις ἔτι καὶ νῦν ἀρὰς ποιοῦνται, πρὶν ἄλλο τι χρηματίζειν, εἴ τις ἐπικηρυκεύεται Πέρσαις τῶν πολιτῶν.\\
\francais{Nos pères ont condamné à mort plusieurs citoyens pour leur adhésion au parti des Perses (μηδισμός). Encore aujourd'hui, dans nos assemblées, avant de traiter aucune affaire, on prononce des imprécations coutre celui des citoyens qui fait des ouvertures aux Mèdes.}}
L'ennemi héréditaire, pour la féminine assemblée, est remplacé -- on s'y attendait ! -- par Euripide !\\ 
Remplacé ? pas tout-à-fait ; car au côté d'Εὐριπίδῃ continuent de subsister étrangement les Μήδοις... détail qui fleure bon son interpolation : le spectateur athénien n'avait pas besoin qu'on lui souligne ainsi le trait comique...
% (Plutôt que d'une plaisanterie παρὰ προσδοκίαν, l'on devrait plus proprement parler d'une plaisanterie μετὰ προσδοκίας.)\\
%Interpolation, alors ? c'est ce que nous allons pouvoir essayer de montrer grâce à un papyrus d'Oxyrhynque !\\
Or nous possédons un papyrus d'Oxyrhynque\marqueur{1} du \nieme{II} siècle après Jésus-Christ\marqueur{2} qui comporte une citation des vers 335-337 des \italique{Thesmophories}. Malheureusement, le vers qui nous intéresse est lacunaire :%
\citComm{ευριπιδιω\douteux{ι}[. .]\douteux{υ}\marqueur{3}τ|επιβλαβη\douteux{ι}τι|νι}%
Hunt avait déjà remarqué en son temps que la lacune était bien insuffisante pour qu'ait pu s'y trouver Μήδοις τ᾽\marqueur{4}.%
%\eff{La lecture probable d'un υ\marqueur{4} avant le τ laisse à penser qu'il faut compléter ainsi la lacune : [τα]ῦτ᾽\marqueur{5}.\\
%Il suffit alors de faire la simple hypothèse de la chute d'un êta -- je la fais : ευριπιδ<η>ι pour que le passage devienne intelligible :%
%\citComm{ηεπικη|ρυκευετ[αι]ευρι|πιδ<η>ιω\douteux{ι}[τα]\douteux{υ}τ|επιβλαβη\douteux{ι}τι|νι [τῇ τῶν γυναικῶν]\\
%ἢ 'πικηρυκεύεται / Εὐριπίδῃ ᾧ ou bien ws ταῦτ᾽ ἐπὶ βλάβῃ τινὶ / τῇ τῶν γυναικῶν\\
%ou s'il envoie un ambassadeur à Euripide qui [utilise] ces informations au détriment des femmes}
%(Avec le redoublement Εὐριπίδῃ / ᾧ repris plus loin avec τινὶ / τῇ ; pastiche du style de la procédure ecclésiastique.)
}%
\debutNote{336-7}%
\note{1}{\italique{The Oxyrhynchus Papyri}, IX (1912), ed. Arthur S. Dunt. Il s'agit du papyrus n° 1176, fragment 39, colonne XII 8-16. Ce papyrus date du \nieme{II} siècle ap. J.-C. et comporte un fragment d'une \italique{Vie} d'Euripide par Satyros de Callatis. Je me réfère aussi à l'édition ainsi qu'au commentaire de ce fragment par Stefan Schorn, \italique{Satyros aus Kallatis} (2002).}%
\note{2}{Quand notre manuscrit de Ravenne est beaucoup plus tardif : \nieme{IX} ou \nieme{X} s. ap. J.-C.}%
\note{3}{Schorn indique que la présence d'un \douteux{c} est moins probable, et que ce n'est assurément pas un \douteux{ε}.}%
\note{4}{Ce qui n'a pas empêché les éditeurs depuis de prétendre dans leur apparat critique que ce papyrus confirmait l'addition du τ᾽ de Biset : Εὐριπίδῃ Μήδοις <τ᾽>... sauf que le τ identifié par Hunt ne pouvait en aucun cas se rattacher aux Μήδοις, absents du papyrus.}%
%\note{5}{On trouve l'expression ταῦτʼ ἐπικηρυκεύομαι au vers 1163, lorsqu'Euripide conclue la paix avec les femmes ; ταῦτα δι᾽ ἀγγέλων ἐπεκηρυκεύετο et ταῦτά οἱ πέμψας ὁ Σιτάλκης ἐπεκηρυκεύετο ou encore ταῦτα μὲν ἐπεκηρυκεύσατο τοῖσι Δηλίοισι chez Hérodote (I, 60 ; IV, 80 ; VI, 97).}
\\
τῇ τῶν γυναικῶν, ἢ τυραννεῖν ἐπινοεῖ\\
ἢ τὸν τύραννον συγκατάγειν, ἢ παιδίον
\commentaire{338-9}{ἢ τυραννεῖν ἐπινοεῖ ἢ τὸν τύραννον συγκατάγειν.}{Citation quasi-textuelle d'une ancienne loi, peut-être de Dracon, rapportée dans la \italique{Constitution d'Athènes} (XVI, 10) :
\citComm{ἦσαν δὲ καὶ τοῖς Ἀθηναίοις οἱ περὶ τῶν τυράννων νόμοι πρᾷοι κατ᾽ ἐκείνους τοὺς καιρούς, οἵ τ᾽ ἄλλοι καὶ δὴ καὶ ὁ μάλιστα καθήκων πρὸς τὴν τῆς τυραννίδος <κατάστασιν>. νόμος γὰρ αὐτοῖς ἦν ὅδε. `θέσμια τάδε Ἀθηναίων καὶ πάτρια· ἐάν τινες τυραννεῖν ἐπανιστῶνται ἐπὶ τυραννίδι, ἢ συγκαθιστῇ τὴν τυραννίδα, ἄτιμον εἶναι καὶ αὐτὸν καὶ γένος.'\\
\francais{Les lois des Athéniens sur la tyrannie étaient alors [à l'époque de Pisitrate] peu sévères, celle surtout qui visait les entreprises des tyrans et dont voici le texte : « Les lois athéniennes établies par nos pères portent que : quiconque aspire à la tyrannie ou forme un complot pour l'établir, sera frappé d'atimie dans sa personne et dans sa race. »}}%
Étant donnée l'époque troublée de la représentation des \italique{Thesmophories}, cette loi devait être sur toutes les bouches ; pour preuve, peu après, en 410/09\marqueur{1}, Démophante proposa un décret inspiré de cette dernière garantissant l'impunité à l'assassin de toute personne complice de crime de lèse-démocratie (cette loi fut gravée sur une stèle fameuse à Athènes qu'Andocide cite dans son \italique{Discours sur les mystères}, 97) :%
\citComm{\grec{καὶ ἐάν τις τυραννεῖν ἐπαναστῇ ἢ τὸν τύραννον συγκαταστήσῃ}.}%
}%
\debutNote{338-9}%
\note{1}{C. Carey, Ifigeneia G., Brenda G.-W., \italique{Use and abuse of law in the Athenian courts}, p. 295}\\
ὑποβαλλομένης κατεῖπέ τις, ἢ δούλη τινὸς\apparat{340}{κατεῖπεν τις R : \lecon{κατεῖπέ τις} Brunck : del. τις Bentley \ligneApp{341} ἐνετρύλλισεν R, Suid. ε 1271 (sed ἐνετρύλλισε) : \lecon{ἐνετρύλισεν} Dindorf}\\
προαγωγὸς οὖσ᾽ ἐνετρύλλισεν τῷ δεσπότῃ\testimonia{341}{Suidas ε 1271 \espTest (προαγωγὸς) Hesychius π 3308 προαγωγός \espTest (ἐνετρύλλισεν) Zonaras 745 οὔ σ' ἐνετρύλλησεν}\\
ἢ πεμπομένη τις ἀγγελίας ψευδεῖς φέρει,\\\passeLigne}
\end{groupe}%
%% v. 331-351
\begin{groupe}%
\ledpb%
\trimIamb{ἢ μοιχὸς εἴ τις ἐξαπατᾷ ψευδῆ λέγων\\
καὶ μὴ δίδωσιν ἃν ὑπόσχηταί ποτε,\apparat{344}{ἅ ἂν R : \lecon{ἃν} Faber}\\
ἢ δῶρά τις δίδωσι μοιχῷ γραῦς γυνή,\\
ἢ καὶ δέχεται προδιδοῦσ᾽ ἑταίρᾳ τὸν φίλον,\apparat{346}{ἑταίρα R : \lecon{ἑταίρᾳ} Fritzsche ap Schol.}%
\commentaire{346}{ἢ καὶ δέχεται προδιδοῦσ᾽ ἑταίρα τὸν φίλον.}{Je reprends la correction de Fritzsche qui substitue avec bonheur \grec{ἑταίρᾳ} à \grec{ἑταίρα}. Le scholiaste précise d'ailleurs qu'il y a différentes versions dans les manuscrits : certaines avec le iota souscrit, d'autres sans : 
\citComm{γρ(άφεται) ἔν τισιν ἑταίρα καὶ ἔστιν ἕτερος νοῦς [...] ἐὰν μὲν ἄνευ τοῦ ῑ γρ(άφηται) ἑταίρα [...] ἐὰν δὲ σὺν τῷ ῑ}
La leçon sans iota, qui est la leçon quasi-universellement acceptée, est absurde : en quoi une hétaïre, dont c'est le métier de recevoir des dons en échange de faveurs \italique{trahirait-elle} (\grec{προδιδοῦσα}) son amant ce faisant ? 
(Willems, sentant bien la difficulté, s'éloigne en traduisant : « à l'insu de son amant ».) 
Seule donc la première version est satisfaisante : il est question au vers précédent d'une vieille faisant des cadeaux pour s'attirer les faveurs d'un galant (thème qu'Aristophane développe dans le \italique{Ploutos}, vv. 959 sqq.) ; après la lubricité, le lucre : cette fois la vieille reçoit des cadeaux pour servir d'entremetteuse en livrant (προδιδοῦσ᾽) l'amant (τὸν φίλον) à une femme irrégulière (ἑταίρᾳ) (thème repris à la fin de la pièce lorsqu'Euripide se déguise et prend à sa suite une danseuse).% thème de la vieille procureuse qui sera repris à la fin de la pièce\eff{ (v. ....)}.
}\\
κεἴ τις κάπηλος ἢ καπηλὶς τοῦ χοὸς\testimonia{347-8}{Suidas χ 394 (sed εἴ pro κεἴ)}\apparat{347}{\lecon{χοὸς} R, Suid. λ 394 : χοῶς Brunck}\\
ἢ τῶν κοτυλῶν τὸ νόμισμα διαλυμαίνεται,%
\commentaire{347-8}{κεἴ τις κάπηλος ἢ καπηλὶς τοῦ χοὸς ἢ τῶν κοτυλῶν τὸ νόμισμα διαλυμαίνεται.}{Les débitantes de mauvais aloi sont brocardées de la même façon par Aristophane dans le \italique{Ploutos}, v. 435-6 :
\citComm{Βλεψίδημος — ἆρ᾽ ἐστὶν ἡ καπηλὶς ἡ ᾽κ τῶν γειτόνων,\\
ἣ ταῖς κοτύλαις ἀεί με διαλυμαίνεται;}%
}\\
κακῶς ἀπολέσθαι τοῦτον αὐτὸν κᾠκίαν\apparat{349}{κ' οἰκίαν R : \lecon{κᾠκίαν} Brunck}\\
ἀρᾶσθε, ταῖς δ᾽ ἄλλαισιν ὑμῖν τοὺς θεοὺς\apparat{350}{\lecon{ὑμῖν} R : ἡμῖν \Mon}\\
εὔχεσθε πάσαις πολλὰ δοῦναι κἀγαθά.\apparat{351}{\lecon{κἀγαθά} R : τἀγαθά Zanetti}\\\passeLigne}%
\end{groupe}%
%% v. 352-354
\begin{groupe}%
\ledpb%
\personnageGrec{Χορός}%
\typeVers{$2$ iamb.}\chantChoeur{ξυνευχόμεσθα τέλεα μὲν\apparat{352}{ξυνευχόμεθα R : \lecon{ξυνευχόμεσθα} Daubruz}%
\commentaire{352-71}{}{Je reprends encore une fois la colométrie de Prato pour ce chant du chœur.}}%
\end{groupe}%
%% v. 353
\begin{groupe}%
\typeVers{$2$ iamb\cat}\chantChoeur{πόλει τέλεα τε δήμῳ\apparat{353}{\lecon{τε} R : δὲ \Mon}}%
\end{groupe}%
%% v. 354
\begin{groupe}%
\typeVers{$2$ iamb\cat}\chantChoeur{τάδ᾽ εὔγματ᾽ ἀποτελεῖσθαι,\apparat{354}{γενέσθαι R : ἐκγενέσθαι Dindorf : \lecon{ἀποτελεῖσθαι} Willems}%
\commentaire{354}{τάδ᾽ εὔγματα γενέσθαι.}{J'adopte la correction de Willems qui substitue ἀποτελεῖσθαι à γενέσθαι (\nieme{9} note, reproduite en annexe).}}%
\end{groupe}%
%% v. 355
\begin{groupe}%
\typeVers{$2$ iamb\cat}\chantChoeur{τὰ δ᾽ ἄρισθ᾽ ὅσαις προσήκει}%
\end{groupe}%
%% v. 356
\begin{groupe}%
\typeVers{ia. $+$ choria.}\chantChoeur{νικᾶν λεγούσαις· ὁπόσαι δ᾽
\commentaire{356-67}{}{Sur tout ce passage, vers 356 à 367, pour le texte et l'interprétation, allez voir la \nieme{10} note d'Alphonse Willems qui fut le point de départ à l'excellent article de Taillardat \& Blanc reproduit ci-après.}}%
\end{groupe}%
%% v. 357-8
\begin{groupe}%
\ligne{357}\decoupeVers{8}\typeVers{$3$ choriamb.}\chantChoeur{ἐξαπατῶσιν παραβαίνουσί τε τοὺς\apparat{357}{ἐξαπατῶσι R : \lecon{ἐξαπατῶσιν} Bentley}}%
\end{groupe}%
%% v. 359
\begin{groupe}%
\ligne{359}\typeVers{glyc.}\chantChoeur{ὅρκους τοὺς νενομισμένους}%
\end{groupe}%
%% v. 360
\begin{groupe}%
\typeVers{glyc.}\chantChoeur{κερδῶν οὕνεκ᾽ ἐπὶ βλάβῃ,\apparat{360}{ἕνεκ᾽ R : \lecon{οὕνεκ᾽} Bentley}}%
\end{groupe}%
%% v. 361
\begin{groupe}%
\typeVers{glyc.}\chantChoeur{ἢ ψηφίσματα καὶ νόμον}%
\end{groupe}%
%% v. 362
\begin{groupe}%
\typeVers{glyc.}\chantChoeur{ζητοῦσ᾽ ἀντιμεθιστάναι,}%
\end{groupe}%
%% v. 363
\begin{groupe}%
\typeVers{glyc.}\chantChoeur{τἀπόρρητά τε τοῖσιν ἐχ-\testimonia{363-4}{(τἀπόρρητά -- λέγουσ᾽) Suidas α 3501 ἀπόρρητά τε τοῖσιν ἐχθροῖσι τοῖς ἡμετέροις λέγουσιν}}%
\end{groupe}%
%% v. 364
\begin{groupe}%
\typeVers{glyc.}\chantChoeur{\alinea θροῖς τοῖς ἡμετέροις λέγουσ᾽,\apparat{364}{λέγουσιν R, Suid. α 3501 : \lecon{λέγουσ᾽} Bentley}}%
\end{groupe}%
%% v. 365
\begin{groupe}%
\typeVers{glyc.}\chantChoeur{ἢ Μήδους ἐπάγουσι τῆς}%
\end{groupe}%
%% v. 366
\begin{groupe}%
\typeVers{glyc.}\chantChoeur{χώρας οὕνεκ᾽ ἐπὶ βλάβῃ,
\commentaire{365-6}{ἢ Μήδους ἐπάγουσι τῆς χώρας / οὕνεκ᾽ ἐπὶ βλάβῃ.}{Voilà deux vers bien malmenés et retouchés à l'envi par divers éditeurs d'Aristophane. On doit à Jean Taillardat et Alain Blanc d'avoir définitivement confirmé la leçon du manuscrit de Ravenne, tout en y relevant un jeu de mots jusqu'alors passé à l'as. La lecture de leur article\marqueur{1} est si rapicolante que je me ferais scrupule d'y taillarder, en voici donc un très copieux extrait : « Pendant plusieurs siècles, l'interprétation des vers 365-366 a été controversée. Les savants ont d'abord été d'accord pour chercher dans τῆς χώρας, v. 366, une désignation géographique (“l'Attique”) et pour comprendre, comme Andrea Divo, “\latin{quaecunque ... Persas inducunt in regionem nostram}”. Mais traduire τῆς χώρας οὕνεκα par “\latin{in regionem nostram}”\marqueur{2}, c'est prêter à οὕνεκα un sens qu'il n'a jamais. D'où diverses conjectures\marqueur{3} et, enfin, la \italique{crux philologorum} dont Van Leeuwen a marqué les vers 365-367. C'est à Willems\marqueur{4} que revient le mérite d'avoir fait progresser, d'une façon décisive, l'interprétation de ce texte : il a rendu leur juste sens aux leçons transmises, ἐπάγειν et οὕνεκα, en traduisant “toutes celles ... qui sollicitent les Mèdes en considération de notre pays, mais à notre dam”. Il commente ainsi sa version (t. II, p. 551-553) : “ Ici il s'agit de ceux qui sollicitent les Mèdes « dans l'intérêt du pays », τῆς χώρας οὕνεκα [ἐπὶ βλάβῃ]. « Dans l'intérêt du pays » est tout simplement un euphémisme équivalent à « contre les Péloponnésiens », ce qu'il fallait éviter de dire. On sait, en effet, que les oligarques intriguaient auprès des satrapes d'Asie Mineure pour détacher le Grand Roi de l'alliance de Sparte et le faire entrer dans celle d'Athènes (Thucydide VIII, 53). Mais l'on sait non moins bien que cette politique indignait Aristophane ». 
Que Μήδους ἐπάγουσι τῆς χώρας οὕνεκα puisse signifier « elles sollicitent l'alliance des Mèdes pour l'amour de notre pays », cela est sûr\marqueur{5} ; mais on ne voit pas qu'il y ait là « un euphémisme » lié à une sorte d'interdit. Mieux vaudrait dire que le chœur rapporte ironiquement les propos de toute femme « qui cherche à traiter avec lui (= l'ennemi) dans l'intérêt \italique{soi-disant d'Athènes} »\marqueur{5}. On peut admettre en effet que cette phrase transpose adroitement le prétexte donné par les intrigantes (Μήδους ἐπάγομεν τῆς χώρας οὕνεκα) : l'enjambement τῆς | χώρας οὕνεκα invitait le chœur à faire un sort au mot χώρας ainsi mis en relief, en sorte qu'une intonation particulière (ironique) des acteurs pouvait suffire à faire comprendre “pour l'amour de notre pays, disent-elles”.}%
\note{1}{«~Aristophane, ``Thesmophories'', v. 366 : τῆς χώρας οὕνεκα et l'ambition des oligarques dans l'hiver de 412-411~», in \italique{Revue de philologie, de littérature et d'histoire ancienne}, t. LX, 1986, fasc. 2, p. 183-186}%
\note{2}{Andreas Divus, \italique{Aristophanis comoediae undecim e graeco in latinum ad verbum translatae}, Basiliae, 1552. Brunck, en 1783, est même plus précis : « quaecunque ... Medos inducunt in regionem nostram, ad eam perdendam » (Brunck lit ἐπάγουσι τῆς χώρας ἐπὶ βλάβῃ, en supprimant οὕνεκα).}%
\note{3}{Par exemple : ἐπάγουσι τῆς | χώρας ἀσεβοῦσιν (Dobrée, suivi par Dindorf et Blaydes ; οὕνεκ᾽ ἐπὶ βλάβῃ disparaît) ; ἐπάγουσι τῆς | χώρας ἐπὶ βλάβῃ (Brunck) ; ἐπάγουσι χώρας | κερδῶν οὕνεκα ἐπὶ βλάβῃ (Reiske) ; ἐπάγουσι τῆς | ἔχθρας οὕνεκ᾽ ἐπὶ βλάβῃ (Hermann) ; ἐπάγουσι τῆς | ὥρας οὕνεκ᾽ ἐπὶ βλάβῃ (Fritzsche) ; ἐπάγουσι γῇ | κερδῶν οὕνεκ᾽ ἐπὶ βλάβῃ (van Velsen, suivi par Hall-Geldart et Rogers, ce dernier « with some hesitation ») ; τῆς | χώρας οὐ μεμελημένοι (C. Austin, voir ci-dessous, n. 5). Cet acharnement à modifier le texte transmis par le Ravennas (seul manuscrit des \italique{Thesmophories}) est d'autant plus étrange qu'il est métriquement sain : les vers 359-366 sont des glyconiens. }%
\note{4}{Il s'agit de la note 10 de Willems, reproduite en annexe.}%
\note{5}{Malgré C. Austin (\italique{P. C. Ph. S.}, 20, 1974, p. 1-2) qui refuse cette interprétation et corrige le texte (voir note 3).}%
\commentaire{}{}{\ \ \ Il y a cependant autre chose dans ce texte : un jeu de mots sur τῆς χώρας, plaisanterie dont personne ne s'est encore avisé. En effet, le parallélisme éclatant, et sûrement dû à Aristophane lui-même, de κερδῶν οὕνεκ᾽ ἐπὶ βλάβῃ (v. 360) et de τῆς χώρας οὕνεκ᾽ ἐπὶ βλάβῃ (v. 366) incite à chercher dans τῆς χώρας, autant que le nom de l'Attique, une variation de la notion exprimée par κερδῶν ; car le chœur peut faire allusion, ici comme là, à l'intérêt personnel et à l'avidité des gens qu'il maudit. Or le mot χώρα désigne, chez plusieurs auteurs, “la place” qu'on occupe dans la vie, par exemple “le rang”, qu'on a dans la hiérarchie sociale ou militaire\marqueur{7}. Dans les \italique{Thesmophories}, v. 366, c'est de rang dans la hiérarchie politique et sociale qu'il s'agit ; si bien que τῆς χώρας οὕνεκα, signifie à la fois “pour l'amour de notre pays (disent-elles)” et “pour se pousser aux (bonnes) places”\marqueur{8}, littéralement “en vue de leur place personnelle”\marqueur{9}. 
Plus que le cas général des σπουδαρχίδαι, “les sires de Brigue-charges”\marqueur{10}, Aristophane vise ici une circonstance particulière. Il pense, comme Willems l'a dit, aux intrigues ambitieuses d'Alcibiade\marqueur{11} et de Pisandre qui, pendant l'hiver 412-411, prétendaient arracher la Perse à son alliance spartiate et cherchaient à obtenir, pour Athènes, l'appui et l'or du Grand Roi\marqueur{12}, pour eux-mêmes, le pouvoir à Athènes. Et que fait Pisandre, à la fin de 412, quand le peuple le délègue auprès d'Alcibiade et de Tissapherne ? Avant de quitter l'Attique pour l'Asie Mineure, “il rend visite aux sociétés secrètes déjà établies à Athènes pour s'occuper des procès et des magistratures”, ξυνωμοσίας, αἵπερ ἐτύγχανον πρότερον ἐν τῇ πόλει οὖσαι ἐπὶ δίκαις καὶ ἀρχαῖς, ἁπάσας ἐπελθὼν... (Thuc. VIII, 54, 4). C'est dire que les comploteurs avaient ouvert la chasse aux places. Les oligarques se consultent ἐπὶ δίκαις καὶ ἀρχαῖς et tout le monde le sait à Athènes : deux mois avant les \italique{Thesmophories}, Aristophane a déjà dénoncé “ces gens qui s'agglomèrent et font bourre pour se pousser aux charges” τούς γε συνισταμένους τούτους καὶ τοὺς πιλοῦντας ἑαυτοὺς | ἐπὶ ταῖς ἀρχαῖσι (\italique{Lysistrata}, v. 577 sq.). Ces ambitieux sont concussionnaires et vénaux : “c'est pour être en mesure de voler que Pisandre et ceux qui ne pensent qu'aux charges créent à chaque instant grabuge et désordre. Qu'ils fassent donc, pour cela, tout ce qu'ils voudront ; mais cet argent-ci\marqueur{13}, il n'y a plus de danger qu'ils l'emportent”, 
\citComm{ἵνα γὰρ Πείσανδρος ἔχοι κλέπτειν χοἰ ταῖς ἀρχαῖς ἐπέχοντες,\\
ἀεί τινα κορκορυγὴν ἐκύκων. οἱ δ᾽ οὖν τοῦδ᾽ οὕνεκα δρώντων\\
ὅ τι βούλονται: τὸ γὰρ ἀργύριον τοῦτ᾽ οὐκέτι μὴ καθέλωσιν.\\
(\italique{Lysistrata}, v. 490-492).}%
Ἐπὶ (δίκαις καὶ) ἀρχαῖς, dit Thucydide ; ἐπὶ ταῖς ἀρχαῖσι, οἱ ταῖς ἀρχαῖς ἐπέχοντες, dit Aristophane. Peut-on rêver plus belles gloses au τῆς χώρας οὕνεκα des \italique{Thesmophories} ? Aux yeux d'Aristophane, tremper dans la grenouillère des combinaisons louches de l'hiver de 412-411, c'est donc prétendre à faire d'une pierre deux coups : devenir un personnage important à Athènes et s'enrichir en profitant de tous les avantages, licites ou non\marqueur{14}, attachés aux ἀρχαί ; c'est tout cela que résume τῆς χώρας οὕνεκ᾽ ἐπὶ βλάβῃ, car de telles ambitions ne réussissent qu'au mépris du bien public. }%
%\begin{comment}%
\note{6}{Willems, II, p. 552. C'est nous [Taillardat \& Blanc] qui soulignons soi-disant.}%
\note{7}{Xén., \italique{Anab.}, 5, 6, 13 (ἐν ἀνδραπόδων χώρᾳ ἐσόμεθα), \italique{Cyr.}, 2, 1, 18 (ἐν μισθοφόρου χώρᾳ εἶναι), Polybe, 1, 43, 1 (τῶν ἡγεμόνων τινὲς τῶν τὰς μεγίστας χώρας ἐχόντων ἐν τοῖς μισθοφόροις), Pausanias, 10, 1, 8 (ὀ δὲ χώραν ἐν τοῖς ἄρχουσιν ἔχων τὴν μεγίστην μάντις ἦν Τελλίας ὁ Ἠλεῖος). Avec le sens de « place dans la société », χώρα paraît déjà exister chez Théognis, v. 822. }%
\note{8}{On justifiera ci-dessous la traduction « aux \italique{bonnes} places ». }%
\note{9}{Tel est le sens de l'article, cf. Kühner-Gerth, \italique{Satzlehre}, I, p. 593, n°2. C'est l'emploi intentionnel du singulier τῆς χώρας (Kühner-Gerth, I, p. 14, n°4), au lieu du pluriel τῶν χωρῶν (plus fréquent dans ce cas, cf. Kühner-Gerth, I, p. 19, Anm. 3), qui permet le jeu de mots. Pour οὕνεκα, \cf \italique{Lysistrata}, v. 491 : τοῦδ᾽ οὕνεκα « en vue de cela (= en vue des magistratures et des vols qu'elles permettent) ». }%
\note{10}{\italique{Ach.}, 595 ; pour le fait, \cf \italique{Equ.}, 157-159, 164-165, et \italique{Aves}, 798-800.}%
\note{11}{Envisageant de provoquer une révolution à Athènes, il préparait alors son retour dans sa patrie (Thuc., 8, 47, 1-2).}%
\note{12}{Pour le rôle de l'or perse à ce moment, voir Thucydide, 8, 29, 1-2 ; 8, 45, 2-6; 8, 48, 2-3 ; 8, 53, 2.}%
\note{13}{Le trésor déposé à l'Acropole dont les Athéniennes barricadées interdisent l'entrée.}%
\note{14}{Voir J. Taillardat, \italique{Les images d'Aristophane}, pp. 408-423. Pour les avantages licites, \cf \italique{Acharniens}, v. 65 sqq. (indemnité de mission officielle et vie facile des ambassadeurs athéniens auprès du Grand Roi). }}%
\end{groupe}%
%% v. 367
\begin{groupe}%
\typeVers{reiz. $+$ ia.}\chantChoeur{ἀσεβοῦσ᾽ ἀδικοῦσί τε τὴν πόλιν.\apparat{367}{ἀσεβοῦσιν ἀδικοῦσίν R : \lecon{ἀσεβοῦσ᾽ ἀδικοῦσί} Bothe}}%
%\end{comment}%
\end{groupe}%
%% v. 368-371
\begin{groupe}%
\typeVers{dochm.}\chantChoeur{ἀλλ᾽ ὦ παγκρατὲς}%
\end{groupe}%
%% v. 369
\begin{groupe}%
\lednopb%
\typeVers{$2$ iamb.}\chantChoeur{Ζεῦ ταῦτα κυρώσειας, ὥσθ᾽\testimonia{369-71}{(ὥσθ᾽ -- οὔσαις) Suidas π 444 }\\\passeLigne}%
\end{groupe}%
%% v. 370
\begin{groupe}%
\ledpb%
\typeVers{$2$ iamb.}\chantChoeur{ἡμῖν θεοὺς παραστατεῖν}%
\end{groupe}%
%% v. 371
\begin{groupe}%
\typeVers{$2$ iamb\cat}\chantChoeur{καίπερ γυναιξὶν οὔσαις.\\\passeLigne}%
\end{groupe}%
%% v. 372-379
\begin{groupe}%
\lednopb%
\personnageGrec{Κορυφαῖος}%
\typeVers{trim. iamb.}\trimIamb{\ligne{372}ἄκουε πᾶσ᾽. ἔδοξε τῇ βουλῇ τάδε\apparat{372}{(Κηρ.) R : \lecon{(Κορ.)} van Leeuwen : (Κρίτ.) Sommerstein \espApp πᾶς R : \lecon{πᾶσ᾽} Voss \ligneApp{374} Ἀρ[. . . .]\douteux{λ}(vel \douteux{a}, \douteux{δ}, \douteux{μ})ειa \PsiCinq, Ἀρ[χίκλ]ει᾽ suppl. Hunt : \lecon{Τιμόκλει᾽} R, Schol., Suid. ε 2616}\\
τῇ τῶν γυναικῶν· Τιμόκλει᾽ ἐπεστάτει,\testimonia{373-4}{Satyrus \italique{Vit. Euripid.} ($=$ \italique{The Oxyrhynchus Papyri} IX 1176) fr. 39 col. XII 1 (vid. ad 335-7 ; sed Ἀρ[χίκλ]ει᾽ pro Τιμόκλει᾽, suppl. Hunt) \espTest (Τιμόκλει᾽ -- ἐγραμμάτευεν) Suidas ε 2616 (sed ἐπιστάτει pro ἐπεστάτει et ἐγραμμάτευε pro ἐγραμμάτευεν)}%
\commentaire{373}{Ἀρ[χίκλ]ει᾽ / Τιμόκλει᾽.}{Comme aux vers 336-337, le \italique{papyrus Oxyrhynchi} n° 1176\marqueur{1} diffère de la leçon du \italique{Ravennas} pour le nom de notre \italique{épistate}\marqueur{2} ; Hunt a suppléé aux quatre ou cinq lettres manquantes Ἀρ[. . . . .]ειa en proposant : Ἀρ[χίκλ]ει᾽. Pourquoi pas ? d'autres noms étaient cependant possibles.
Quoique le papyrus soit certainement plus fiable que notre \italique{Ravennas}, en l'absence de certitude sur la leçon originale du papyrus, je me bornerai à suivre celle de l'unicodex.}%
\debutNote{373}%
\note{1}{\italique{The Oxyrhynchus Papyri}, IX (1912), ed. Arthur S. Hunt. Il s'agit du papyrus n°1176, fragment 39, colonne XII 8-16 ; papyrus du \nieme{II} siècle ap. J.-C. comportant un fragment d'une \italique{Vie} d'Euripide par Satyros de Callatis. Voir aussi l'édition Schorn de Satyros, \italique{Satyros aus Kallatis} (2002).}%
\note{2}{On sait que c'est au niveau des noms propres que les textes s'altèrent le plus lors des copies successives, surtout si, comme probablement ici, il s'agit d'un personnage inconnu ailleurs.}\\
Λύσιλλ᾽ ἐγραμμάτευεν, εἶπε Ζωστράτη·\\
ἐκκλησίαν ποιεῖν ἕωθεν τῇ μέσῃ\\
τῶν Θεσμοφορίων, ᾗ μάλισθ᾽ ἡμῖν σχολή,\apparat{376}{σχολῇ R : \lecon{σχολή} Zanetti}\\
\commentaire{376-7}{ἐκκλησίαν ποιεῖν ἕωθεν τῇ μέσῃ / τῶν Θεσμοφορίων, ᾗ μάλισθ᾽ ἡμῖν σχολή.}{Remarquons au passage que les assemblées se tenaient à l'aube\marqueur{1}.< Le fait que le \italique{jour du milieu} soit celui du \italique{jeûne} (νηστεία) explique que les femmes aient tout le loisir pour tenir assemblée.}%
\debutNote{376-7}%
\note{1}{Coutume confirmée par d'autres passages d'Aristophane : \italique{Acharniens}, 19-20 ; \italique{Assemblée des femmes}, 20-21, 84-85, 739-741.}%
καὶ χρηματίζειν πρῶτα περὶ Εὐριπίδου,\testimonia{377-8}{(καὶ χρηματίζειν -- παθεῖν) Suidas χ 481 ; Lexicon Vindobonense 190.10}\\
ὅ τι χρὴ παθεῖν ἐκεῖνον· ἀδικεῖν γὰρ δοκεῖ\\
ἡμῖν ἁπάσαις. τίς ἀγορεύειν βούλεται;
\commentaire{373-9}{ἄκουε πᾶς ... τίς ἀγορεύειν βούλεται;}{L'assemblée à huis-clos féminin se donne les formes pour être en droit de juger Euripide : tout ce passage est un pastiche d'un \grec{προβούλευμα}, une \italique{résolution du Conseil}. En voici un exemple : le début du décret proposé par Démophante évoqué plus haut (Andocide, \italique{Discours sur les mystères}, 96) :
\citComm{ΝΟΜΟΣ\\
Ἔδοξε τῇ βουλῇ καὶ τῷ δήμῳ. Αἰαντὶς ἐπρυτάνευε, Κλειγένης ἐγραμμάτευε, Βοηθὸς ἐπεστάτει. Τάδε Δημόφαντος συνέγραψεν. Ἄρχει χρόνος τοῦδε τοῦ ψηφίσματος ἡ βουλὴ οἱ πεντακόσιοι λαχόντες τῷ κυάμῳ, οἷς Κλειγένης πρῶτος ἐγραμμάτευεν.\marqueur{1}}%
}%
\debutNote{373-9}%
\note{1}{Autres exemples : Thucydide, IV, 118 ; \italique{Inscriptiones Graecae} I\expo{3} 5, 9, 10, 12, 17 \&c. \&c.}\\\passeLigne}%
\end{groupe}%
%% v. 380a
\begin{groupe}%
\personnageGrec{Μίκα}%
\trimIamb{ἐγώ.\apparat{380}{(Γυνή τις) R : \lecon{(Μίκα)} Rogers, cf. v. 760}\\\passeLigne}%
\end{groupe}%
%% v. 380b-382
\begin{groupe}%
\personnageGrec{Κορυφαῖος}%
\trimIamb{\ligne{380}\decalage{ἐγώ.}περίθου νυν τόνδε πρῶτον πρὶν λέγειν.\testimonia{380}{Suid. π 1166 περίθου νῦν τόνδε πρότερον}\apparat{380}{\lecon{πρῶτον} R : πρότερον Suid. π 1166}}%
\end{groupe}%
\begin{groupe}%
\typeVers{tetr. iamb\cat}\tetramIamb{σίγα σιώπα, πρόσεχε τὸν νοῦν· χρέμπτεται γὰρ ἤδη\testimonia{381}{(χρέμπτεται) \Phryn 126.1}%
\commentaire{381}{χρέμπτεται.}{Elle \italique{tousse}, et non elle crache, comme traduisent les dictionnaires. L'ancien français avait un terme excellent, qui s'est perdu avec tant d'autres : « Il se lève, compose sa robe, sa barbe et ses yeux à la modestie, et ayant craché et toussi clair, il commença ». D'Aubigné, \italique{Avent. du baron de Foeneste}, I. IV, ch. 10. (Note de Willems.)}
\lednopb%
\\
\lednopb%
ὅπερ ποιοῦσ᾽ οἱ ῥήτορες. μακρὰν ἔοικε λέξειν.\\\passeLigne}%
\end{groupe}%
%% v. 383-432
\begin{groupe}%
\ledpb%
\personnageGrec{Μίκα}%
\typeVers{trim. iamb.}\trimIamb{φιλοτιμίᾳ μὲν οὐδεμιᾷ μὰ τὼ θεὼ\testimonia{383-4}{(φιλοτιμίᾳ -- γυναῖκες) Suidas φ 431}\apparat{383}{(\auDessus{γυ}{ν}) \Rac : (καλλιλεξία γυνή) \Rpc, \Mon : (\lecon{Μίκα}) Rogers}\\
λέξουσ᾽ ἀνέστην ὦ γυναῖκες· ἀλλὰ γὰρ\\
βαρέως φέρω τάλαινα πολὺν ἤδη χρόνον\\
προπηλακιζομένας ὁρῶς ἡμᾶς ὑπὸ\apparat{386}{ὑμᾶς R : \lecon{ἡμᾶς} Cobet}\\
Εὐριπίδου τοῦ τῆς λαχανοπωλητρίας\\
καὶ πολλὰ καὶ παντοῖ᾽ ἀκουούσας κακά.\\
τί γὰρ οὗτος ἡμᾶς οὐκ ἐπισμῇ τῶν κακῶν;\testimonia{389}{Suidas ε 2592 \espTest (ἐπισμῇ) Hesychius ε 5195 ἐπισμῇ}\apparat{389}{\lecon{ἡμᾶς} R, Suid. ε 2592 : ἡμῖν Brunck}\\
ποῦ δ᾽ οὐχὶ διαβέβληχ᾽, ὅπουπερ ἔμβραχυ\testimonia{390-4}{Suidas οι 128 (οἷς θέα pro εἰσὶν θεαταὶ et τραγῳδία pro τραγῳδοὶ)}\testimonia{390-1}{scholium Plato \italique{Theages} 127c $=$ scholium Plato \italique{Gorgias} 457b $=$ scholium Plato \italique{Hippias min.} 365d \espTest (ποῦ -- θεαταὶ) Suidas ε 955 (sed εἰσὶ pro εἰσὶν)}\apparat{390}{\lecon{ποῦ} R, Suid. : πῶς Schol. Plat. \italique{Theag}.127c \espApp ἐν βραχὺ R, Schol. Plat. \italique{Theag}. : \lecon{ἔμβραχυ} Suid., Schol. Plat. \italique{Gorg}.457b et \italique{Hipp.min}.365d \ligneApp{391} \lecon{τραγῳδοὶ καὶ} R : τραγῳδικοὶ Schol. Plat. \italique{Theag} : τραγῳδία καὶ Suid. οι 128 \ligneApp{392} μυχοτρόπους R : \lecon{μοιχοτρόπους} Suid. οι 128 \espApp ἀνδρεαστρίας R, Suid. AV : \lecon{ἀνδρεραστρίας} Suid. GFM \ligneApp{393} οἰνοπότιδας R : \lecon{οἰνοπίπας} Suid. : οἰνοπίπους Schol. \ligneApp{394} ἀνδράσι R : \lecon{ἀνδράσιν} Daubuz}\\% \ligneApp{398} ὥσπερ R : \lecon{ὧνπερ} Küster}\\
εἰσὶν θεαταὶ καὶ τραγῳδοὶ καὶ χοροί,\\
τὰς μοιχοτρόπους, τὰς ἀνδρεραστρίας καλῶν,\testimonia{392}{(ἀνδρεραστρίας) Pollux III 71 ; \Phryn 34.18 ; \Eustathe 380.15}\\
τὰς οἰνοπίπας, τὰς προδότιδας, τὰς λάλους,\testimonia{393}{(οἰνοπίπας) scholium \italique{Aves} 307}%
\commentaire{393}{τὰς οἰνοπίπας,}{Encore une fois, c'est le scholiaste, disposant de meilleurs manuscrits, qui nous permet de revenir à la leçon originale : après avoir évoqué Symmaque, il rapproche non sans raison \grec{οἰνοπίπης} de \grec{παιδοπίπης} (on aurait pu rajouter : \grec{γυναικοπίπης}, \grec{παρθενοπίπης} ou encore le \grec{πυροπίπης} / πυρροπίπης des \italique{Cavaliers}, v. 407) ; c'est-à-dire que le suffixe ne renverrait pas aux verbes \grec{πιπίζειν} ou \grec{πιπίσκειν} mais -\grec{οπίπης} serait de la famille de \grec{ὀπιπεύειν}, \italique{guigner}, \italique{guetter}, \italique{épier} (cf. aussi Hésychius (ο 1006) : \grec{ὀπῖπα· ἐξαπατᾷ, ἀπατεών, ἢ ἀπατῶν}). 
Aristophane ne va pas manquer de surenchérir : la difficile relation des femmes au vin sera allègrement mise en scène dans la suite de la pièce.}\\
τὰς οὐδὲν ὑγιές, τὰς μέγ᾽ ἀνδράσιν κακόν·\commentaire{394}{τὰς οὐδὲν ὑγιές.}{Érasme a compilé l'expression dans ses \italique{Adages} (n° 738, \latin{nihil sanum}).}\\
ὥστ᾽ εὐθὺς εἰσιόντες ἀπὸ τῶν ἰκρίων\testimonia{395-6}{(ὥστ᾽ εὐθὺς -- ἡμᾶς) Suidas ι 275}\testimonia{395}{(ἰκρίων) Hesychius ι 501 ἴκρια}\\
ὑποβλέπουσ᾽ ἡμᾶς σκοποῦνταί τ᾽ εὐθέως\\
μὴ μοιχὸς ἔνδον ᾖ τις ἀποκεκρυμμένος.\\
δρᾶσαι δ᾽ ἔθ᾽ ἡμῖν οὐδὲν ὧνπερ καὶ πρὸ τοῦ\apparat{398}{ὥσπερ R : \lecon{ὧνπερ} Küster}\\
ἔξεστι· τοιαῦθ᾽ οὗτος ἐδίδαξεν κακὰ\\\passeLigne}
\end{groupe}%
%% v. 383-432
\begin{groupe}%
\ledpb%
\trimIamb{τοὺς ἄνδρας ἡμῶν· ὥστ᾽ ἐάν τις ξυμπλέκῃ\apparat{400}{τις πλέκῃ R : \lecon{τις ξυμπλέκῃ} Willems : δή τις πλέκῃ Austin : alii alia}%
\commentaire{400}{τοὺς ἄνδρας ἡμῶν· ὥστ᾽ ἐάν τις πλέκῃ.}{Le vers 400 pèche contre le mètre. Willems y supplée en lisant <ξυμ>πλέκῃ ; il s'en justifie dans la \nieme{11} note, reproduite à la fin.}%
\commentaire{400-1}{ὥστ᾽ ἐάν τις πλέκῃ / γυνὴ στέφανον, ἐρᾶν δοκεῖ·}{Le scholiaste s'est montré court sur ce travers qu'Euripide impute aux femmes tressandières : l'on ne sait de laquelle de ses pièces (perdues) le reproche est tiré.
Toutefois, tresser une couronne pour son amant(e) est chose bien attestée ; et le scholiaste de citer un vers de Sappho\marqueur{1} :
\citComm{νεωτέρων καὶ ἐρωτικῶν τὸ στεφανοπλοκεῖν. πρὸς τὸ ἔθος, ὅτι ἐστεφανηπλόκουν αἱ παλαιαί.\\
Σαπφώ· αὐτὰ κόραις στεφαναπλόκεο\\
Sappho : «~toi qui te faisais tresser des couronnes par des jeunes filles~».}
Athénée (XV, 670d) confirme la coutume d'offrir des couronnes à son être aimé :
\citComm{στεφανοῦσιν δὲ τὰς τῶν ἐρωμένων θύρας ἤτοι τιμῆς χάριν καθαπερεί τινος θεοῦ τὰ πρόθυρα [στεφανοῦσιν αὐτῶν], ἢ οὐ τοῖς ἐρωμένοις ἀλλὰ τῷ Ἔρωτι ποιούμενοι τὴν τῶν στεφάνων ἀνάθεσιν τοῦ μὲν Ἔρωτος τὸν ἐρώμενον ἄγαλμα, τούτου δὲ ναὸν ὄντα τὴν οἴκησιν στεφανοῦσι [τὰ τῶν ἐρωμένων πρόθυρα]: διὰ ταῦτα δὲ καὶ θύουσιν ἔνιοι ἐπὶ ταῖς τῶν ἐρωμένων θύραις.\\
Quant aux couronnes et aux guirlandes dont les amants ornent les portes des personnes qu'ils aiment, ils les y mettent ou pour leur rendre des honneurs, comme couronnant le porche d'une divinité ; ou, c'est à l'amour qu'ils les offrent en les posant, et non aux personnes qui sont aimées ; car ces personnes sont pour les amants les images mêmes du dieu d'amour; et la maison où elles demeurent en étant le temple, c'en est aussi le porche qu'ils ornent de couronnes, en les mettant aux portes de ces personnes. Voilà aussi pourquoi on immole quelquefois des victimes aux portes des personnes qu'on aime.}}%
\debutNote{400-1}%
\note{1}{Citation corrompue ; le manuscrit de Ravenne donne : αυταόρα αἱστεφανηπλόκουν. La lecture proposée est de mon cru. Théodore Reinach, dubitatif, reprenait la conjecture de Hermann : αὐτὰρ ὀραῖαι στεφανηπλόκευν. Ce fragment de Sappho est le \nieme{70} de l'édition Reinach-Puech aux Belles Lettres (1937) et le \nieme{125} de l'édition d'Éva-Maria Voigt, \italique{Sappho et Alcaeus} (1971).}\\
γυνὴ στέφανον, ἐρᾶν δοκεῖ· κἂν ἐκβάλῃ\\
σκεῦός τι κατὰ τὴν οἰκίαν πλανωμένη,\\
ἁνὴρ ἐρωτᾷ· `τῷ κατέαγεν ἡ χύτρα;\apparat{403}{ἀνὴρ R : \lecon{ἁνὴρ} Brunck}\\
οὐκ ἔσθ᾽ ὅπως οὐ τῷ Κορινθίῳ ξένῳ.'\testimonia{404}{(τῷ Κορινθίῳ ξένῳ) Hesychius κ 3629 Κορίνθιος ξένος}%
\commentaire{401-4}{κἂν ἐκβάλῃ ... τῷ Κορινθίῳ ξένῳ.}{\italique{L'étranger de Corinthe}, c'est Bellérophon, objet d'une coupable passion de la part de la princesse Sthénébée dans la pièce éponyme d'Euripide. Cette dernière, voyant ses avances repoussées par le chaste héros, l'accuse devant le roi son mari d'avoir tenté de la violenter. Le roi ordonne à Bellérophon de partir pour la Carie. Il y reçoit pour mission d'affronter la Chimère. C'est pendant le temps que Bellérophon est en Carie, et que Sthénébée le croit mort, qu'elle verse des sortes de libations à la mémoire du héros, comme l'on fait au jeu du cottabe avec les personnes aimées. Athénée (X, 427e) nous a conservé ce fragment\marqueur{1} :
\citComm{πεσὸν δέ νιν λέληθεν οὐδὲν ἐκ χερός,\\
ἀλλ᾽ εὐθὺς αὐδᾷ· `τῷ Κορινθίῳ ξένῳ'\\
<La Nourrice>\\
\francais{Rien ne tombe, en lui échappant des mains,\\
qu'elle ne s'écrie aussitôt : « Pour l'hôte de Corinthe ! »}}%
Les maris athéniens prenaient leur femme pour autant de Sthénébées.\\
Ce passage est à rapprocher d'un fragment d'Aristophane, rapporté par Diogène Laërte\marqueur{2} :
\citComm{Τὰ δὲ πεσόντα μὴ ἀναιρεῖσθαι, ὑπὲρ τοῦ ἐθίζεσθαι μὴ ἀκολάστως ἐσθίειν ἢ ὅτι ἐπὶ τελευτῇ τινος· καὶ Ἀριστοφάνης δὲ τῶν ἡρώων φησὶν εἶναι τὰ πίπτοντα, λέγων ἐν τοῖς Ἥρωσι,\\
Μηδὲ γεύεσθ' ἅττ' ἂν ἐντὸς τῆς τραπέζης καταπέσῃ.\\
$[$ Pythagore $]$ interdit de ramasser ce qui tombe de la table pour habituer à manger avec tempérance, ou bien encore parce que cela est destiné aux morts. Ce qui tombe de la table est pour les héros, suivant Aristophane ; car il dit dans les \italique{Héros} :\\
\francais{Ne goûtez point à ce qui tombe de la table.}}}%
\debutNote{401-4}%
\note{1}{fragment 3, \italique{Sthénébée}, de l'édition Jouan-van Looy, Belles Lettres, t. VIII, \nieme{3} partie $=$ fr. 664 de l'édition Kannicht $=$ fr. 883 de l'édition Mette. Traduction Jouan \& van Looy.}%
\note{2}{\italique{Vies}, VIII, 1 (Pythagore), 34. Correspond au fragment 320(305) de l'édition Kassel-Austin (\italique{Poetae Comici Graeci}) des fragments d'Aristophane.}\\
κάμνει κόρη τις, εὐθὺς ἁδελφὸς λέγει·\apparat{405}{ἀδελφὸς R : \lecon{ἁδελφὸς} Biset}\\
`τὸ χρῶμα τοῦτό μ᾽ οὐκ ἀρέσκει τῆς κόρης.'\\
\commentaire{405-6}{εὐθὺς ἁδελφὸς λέγει· / `τὸ χρῶμα τοῦτό μ᾽ οὐκ ἀρέσκει τῆς κόρης.'}{Il la soupçonne d'être enceinte. La pièce d'Euripide à laquelle se réfère Aristophane n'est pas identifiée ; pourtant il semble clair qu'il s'agit d'une citation (ne serait-ce que par l'absence de résolution au vers 406).}%
εἶἑν, γυνή τις ὑποβαλέσθαι βούλεται\\
ἀποροῦσα παίδων, οὐδὲ τοῦτ᾽ ἔστιν λαθεῖν.
\commentaire{406-7}{γυνή τις ὑποβαλέσθαι βούλεται / ἀποροῦσα παίδων, οὐδὲ τοῦτ᾽ ἔστιν λαθεῖν.}{Le thème de la \italique{supposition d'enfants}\marqueur{1} est récurrent dans les \italique{Thesmophories} (cf. vv. 340, 502-516, 564-565). À en croire Satyros dans sa \italique{Vita Euripidis}\marqueur{2}, c'est à Euripide qu'on doit son omniprésence dans la Comédie Nouvelle... et le biographe d'Euripide de donner du grain à moudre à ses contempteurs :
\citComm{πρὸς γ[υ]ναῖκα καὶ πατρὶ πρὸς υ[ἱὸ]ν καὶ θερά[πον]τι πρὸς δ[εσ]πότην, ἢ τ[ὰ] κατὰ τὰς π[ερι]πετείας, β[ια]σμοὺς παρθένων, \textbf{ὑποβολὰς παιδίων}, ἀναγνωρισμοὺς διά τε δακτυλίων καὶ διὰ δεραίων· ταῦτα γάρ ἐστι δήπου τὰ συνέχοντα τα τὴν νεωτέραν κωμῳδίαν, ἃ πρὸς ἄκρον ἤγα[γ]εν Εὐρπίδης, Ὁμήρου [ὄ]ντος ἀρχῆς\\
\francais{$[$le mari$]$ contre sa femme, le père contre son fils, le serviteur contre son maître ; ou lors des péripéties : violences de jeunes filles, \textbf{substitutions d'enfants}, reconnaissances par des anneaux ou des colliers ; hé bien c'est tous ces éléments qui composent la trame de la Nouvelle Comédie, éléments qu'Euripide avait conduit au sommet, Homère en étant le point de départ}}%
La tradition ne nous a rapporté qu'une seule supposition d'enfants chez Euripide : dans les \italique{Phéniciennes} (v. 28-31) :
\citComm{Πολύβου δέ νιν λαβόντες ἱπποβουκόλοι\\
φέρουσ᾽ ἐς οἴκους ἔς τε δεσποίνης χέρας\\
ἔθηκαν. ἣ δὲ τὸν ἐμὸν ὠδίνων πόνον\\
μαστοῖς ὑφεῖτο καὶ πόσιν πείθει τεκεῖν,}%
ainsi qu'une allusion dans l'\italique{Alceste} (v. 639) : Admète, désapprouvant l'attitude de son père qui a refusé de donner sa vie pour lui, suppose qu'il a été supposé.}%
\debutNote{406-7}%
\note{1}{Quand les femmes font passer pour leur des enfants – mâles le plus souvent ! – auprès de leur mari.}%
\note{2}{\italique{The Oxyrhynchus Papyri}, IX.1776, fr. 39, Col. vii, 1-24}\\
ἅνδρες γὰρ ἤδη παρακάθηνται πλησίον·\apparat{409}{ἄνδρες R : \lecon{ἅνδρες} Brunck.}\\
πρὸς τοὺς γέροντάς θ᾽ οἳ πρὸ τοῦ τὰς μείρακας\\
ἤγοντο, διαβέβληκεν, ὥστ᾽ οὐδεὶς γέρων\\
γαμεῖν ἐθέλει γυναῖκα διὰ τοὔπος τοδὶ\apparat{412}{θέλει R : \lecon{ἐθέλει} Kappeyne van de Coppello \italique{Obs. crit.} I.2}\\
`δέσποινα γὰρ γέροντι νυμφίῳ γυνή'.\testimonia{413}{Menandros \italique{Monostich.} 191 ($=$ Euripides fr. 804.3 ed. Kannicht) ; \Apostol V 93b}%
\commentaire{413}{δέσποινα γὰρ γέροντι νυμφίῳ γυνή.}{Stobée (IV, 22a109) nous a conservé la fameuse citation d'Euripide : elle est tirée du \italique{Phoenix}\marqueur{1} :
\citComm{μοχθηρόν ἐστιν ἀνδρὶ πρεσβύτῃ τέκνα\\
δίδωσιν†, ὅστις οὐκέθ᾽ ὡραῖος γαμεῖ·\\
δέςποινα γὰρ γέροντι νυμφίῳ γυνή\\
\francais{C'est un tracas pour un vieillard <qui est doté d'enfants> de se marier alors qu'il n'en a plus l'âge ; avec un vieil époux, une femme est la maîtresse.}}%
\italique{Quand on est jeune, on aime en fou ; Quand on est vieux, qui aime est fou.}}%
\debutNote{413}%
\note{1}{fr. 3, \italique{Phoenix} de l'édition Jouan-van Looy aux Belles-lettres, tome VIII, \nieme{3} partie. Je reprends leur traduction. Ce fragment correspond au n°804 de l'édition Kannicht et 1139 de celle de Mette.}\\
εἶτα διὰ τοῦτον ταῖς γυναικωνίτισιν\\
σφραγῖδας ἐπιβάλλουσιν ἤδη καὶ μοχλοὺς\\\passeLigne}
\end{groupe}%
%% v. 383-432
\begin{groupe}%
\ledpb%
\trimIamb{τηροῦντες ἡμᾶς, καὶ προσέτι Μολοττικοὺς\\
τρέφουσι μορμολυκεῖα τοῖς μοιχοῖς κύνας.\testimonia{417}{(μορμολυκεῖα) \EtymMagnum 590.53}\apparat{417}{μορμολύκεια R : \lecon{μορμολυκεῖα} Dindorf (1830) \espApp μυχοῖς R : \lecon{μοιχοῖς} \Mon}%
\commentaire{416-7}{Μολοττικοὺς / τρέφουσι μορμολυκεῖα τοῖς μοιχοῖς κύνας.}{C'est le conseil et le vœu d'Hippolyte dans la pièce d'Euripide (vv. 645 et suiv. ; note de Willems) : 
\citComm{χρῆν δ᾽ ἐς γυναῖκα πρόσπολον μὲν οὐ περᾶν,\\
ἄφθογγα δ᾽ αὐταῖς συγκατοικίζειν δάκη\\
θηρῶν, ἵν᾽ εἶχον μήτε προσφωνεῖν τινα\\
μήτ᾽ ἐξ ἐκείνων φθέγμα δέξασθαι πάλιν.\\
\francais{Il faudrait que les femmes n'eussent point auprès d'elles de servantes,\\
mais qu'elles fussent servies par de muets animaux,\\
pour qu'elles n'eussent personne à qui parler,\\
ni qui pût à son tour leur adresser la parole.}}}\\
καὶ ταῦτα μὲν ξυγγνώσθ᾽. ἃ δ᾽ ἦν ἡμῖν πρὸ τοῦ\\
αὐταῖς ταμιεῦσαι καὶ προαιρούσαις λαβεῖν\apparat{419}{ταμιεύεσθαι R : \lecon{ταμιεῦσαι καὶ} Reiske}\\
ἄλφιτον ἔλαιον οἶνον, οὐδὲ ταῦτ᾽ ἔτι%
\commentaire{418-20}{ἃ δ᾽ ἦν ἡμῖν πρὸ τοῦ / αὐταῖς ταμιεύεσθαι προαιρούσαις λαβεῖν / ἄλφιτον ἔλαιον οἶνον.}{Je renvoie à la courte mais incisive note (la \nieme{12}) de Willems reproduite à la fin : il y fait un sort à la « déplorable conjecture de Scaliger » (qui avait substitué λαθεῖν à λαβεῖν au vers 419, donnant au passage un sens qui fait contre-sens ; contre-sens pourtant encore repris par certains éditeurs modernes...)}\\
ἔξεστιν. οἱ γὰρ ἄνδρες ἤδη κλῄδια\testimonia{421-8}{(οἱ γὰρ ἄνδρες -- ἐξαψαμένους) Suidas λ 54 (αὐτοὶ, τρεῖς et ᾡκότριψ Εὐριπίδης omissis, Λακωνίκα pro Λακωνίκ᾽, θριπηδέστατ᾽ pro θριπήδεστ᾽ et σφραγῖδα pro σφραγίδια)}\testimonia{421-7}{(οἱ γὰρ ἄνδρες -- θριπήδεστ᾽) Suidas κ 159 (αὐτοὶ et ᾡκότριψ Εὐριπίδης omissis, Λακωνίκα pro Λακωνίκ᾽ et θριπηδέστατ᾽ pro θριπήδεστ᾽)}\testimonia{421-7}{(οἱ γὰρ ἄνδρες -- θριπήδεστ᾽)}\\
αὐτοὶ φοροῦσι κρυπτὰ κακοηθέστατα\testimonia{422}{(κακοηθέστατα) \Eustathe 974.37, 1108.25, 1603.52}\\
Λακωνίκ᾽ ἄττα, τρεῖς ἔχοντα γομφίους.\testimonia{423}{(γομφίους) \Phryn 61.1}\\
πρὸ τοῦ μὲν οὖν ἦν ἀλλ᾽ ὑποῖξαι τὴν θύραν
\commentaire{424}{ὑποῖξαι τὴν θύραν.}{Il s'agit de la porte de l'appartement (cf. v. 414) et non de celle du cellier.}\\
ποιησαμέναισι δακτύλιον τριωβόλου,\\
νῦν δ᾽ οὗτος αὐτοὺς ᾡκότριψ Εὐριπίδης\testimonia{426}{(οἰκότριψ $=$ οἰκοτραφής) Moeris ο 25 ; (ᾡκότριψ) Hesychius ο 265-266 ; \AelDion ο 8-9}\apparat{426}{ὡκότριψ R : \lecon{ᾡκότριψ} Bentley : οἱκότριψ Kühner-Blass \ligneApp{427} ἐδίδαξεν R : \lecon{ἐδίδαξε} Suid. λ 64}%
\commentaire{426}{ᾡκότριψ Εὐριπίδης.}{Ce terme, \italique{tâte-poule}, dont un Grec aurait saisi le sens métaphorique (cf. les verbes βλιμάζειν, σκιμαλίζειν), rend bien le mot ᾡκότριψ. Cf. ἀστύτριψ, \italique{pilier de cité}, Θησειότριψ, \italique{pilier du Théseion} (fragm. du Polyides), πορνότριψ, etc. (Note de Willems.)}\\
ἐδίδαξε θριπήδεστ᾽ ἔχειν σφραγίδια\\
ἐξαψαμένους. νῦν οὖν ἐμοὶ τούτῳ δοκεῖ\testimonia{428-9}{(ἐμοὶ -- ἁμωσγέπως) Suidas κ 2771 (τινα pro τιν᾽ ἡμᾶς)}\\
ὄλεθρόν τιν᾽ ἡμᾶς κυρκανᾶν ἁμωσγέπως,\testimonia{429}{(κυρκανᾶν) Photius κ 1254 κυρκανᾶν $=$ Hesychius κ 4691 κυρκανᾷ ; \Eustathe 1213.50 \espTest (ἁμωσγέπως) Hesychius α 4182 ἁμωσγέπως ; \Syn α 449 $=$ Suidas α 1645}\apparat{429}{ἀμωσγέπως R, Suid. : \lecon{ἁμωσγέπως} Blaydes}%
\commentaire{429}{ὄλεθρόν τιν᾽ ἡμᾶς κυρκανᾶν ἁμωσγέπως.}{Κυρκανᾶν est un synonyme très rare de κυκᾶν et se dit proprement d'un cuisinier préparant une mixture\marqueur{1}. Il est pris ici au sens de \italique{comploter la perte} de quelqu'un. Il est clair que l'expression ὄλεθρόν κυρκανᾶν est faite sur le modèle de κυρκανᾶν φάρμακα (cf. \italique{Ploutos} 302, 309 : τὰ φάρμακ᾽ ἀνακυκῶσαν) [...] Pour le ton, le verbe κυρκανᾶν paraît correspondre à nos métaphores familières \italique{cuisiner}, \italique{mitonner une affaire}. (Note de Taillardat, §421.)}%
\debutNote{429}%
\note{1}{Cf. Épinicus 1,8 : συγκυρκανᾶν (préparation d'une boisson faite de farine, de miel et de vin mêlés).}\\
ἢ φαρμάκοισιν ἢ μιᾷ γέ τῳ τέχνῃ,\\
\lednopb%
ὅπως ἀπολεῖται. ταῦτ᾽ ἐγὼ φανερῶς λέγω,\\
\lednopb%
τὰ δ᾽ ἄλλα μετὰ τῆς γραμματέως συγγράψομαι.\testimonia{432}{(μετὰ τῆς γραμματέως συγγράψομαι) \Eustathe 1484.53 \espTest (τῆς γραμματέως) Pollux IV 19}%
\commentaire{432}{τῆς γραμματέως.}{Aristophane faisait-il sourire sa mâle assemblée lorsqu'il parlait de \italique{la} greffièr\italique{e} ? c'est ce que suggère Pollux (IV, 19) :%
\citComm{Ἀριστοφάνης δὲ παίζει ἐν Θεσμοφοριαζούσαις, λέγων ἡ γραμματεύς.}%
}\\\passeLigne}%
\end{groupe}%
%% v. 433-442
\begin{groupe}%
\ledpb%
\personnageGrec{Χορός\strophe{στροφή}}\typeVers{par.}\chantChoeur{οὔπω ταύτης ἤκουσα\apparat{433}{οὔπω τε \Rac : οὐπώποτε \Rpc : οὔπω Dindorf}%
\commentaire{433-42}{}{Ce chant du chœur trouvera son écho aux vers 520-530 : semblable réplique au point de vue métrique\marqueur{1} qui permettait aux deux passages d'être chantés sur la même mélodie. Là comme ici, le chœur s'exprime à la première personne.}%
\debutNote{433-42}%
\note{1}{À une exception près : au vers 437b où le texte est probablement corrompu.}}%
\end{groupe}%
%% v. 434-5
\begin{groupe}%
\decoupeVers{5}\typeVers{$2$ troch.}\chantChoeur{πολυπλοκωτέρας γυναικὸς}%
\end{groupe}%
%% v. 436
\begin{groupe}%
\ligne{436}\typeVers{$2$ troch.}\chantChoeur{οὐδὲ δεινότερον λεγούσης.\apparat{435}{\lecon{δεινότερον} R : δεινότερα Hermann}}%
\end{groupe}%
%% v. 437a
\begin{groupe}%
\decoupeVers{a}\typeVers{$2$ troch.}\chantChoeur{πάντα γὰρ λέγει δίκαια,}%
\end{groupe}%
%% v. 437b
\begin{groupe}%
\ligne{437}\decoupeVers{b}\typeVers{$2$ iamb.}\chantChoeur{†πάσας δ᾽ ἰδέας ἐξήτασεν†,\testimonia{437-9}{(πάσας -- ἀνηῦρεν) Suidas ε 22 πάσας δ' ἰδέας ἐξήτασε πάντα δ' ἐβάστασε φρενί, πυκνῶς τε ποικίλως λόγους ἀνεῦρεν}\apparat{437}{\lecon{πάσας} εἰδέας \lecon{ἐξήτασεν} R, Suid. ε 22 (sed \lecon{ἰδέας} ἐξήτασε Suid.) : πᾶσαν ἰδέαν ἐξετάζει Hermann \ligneApp{438} ἐβάστασεν R : \lecon{ἐβάστασε} Suid. \ligneApp{439} \lecon{ποικίλους} R : ποικίλως Suid. \espApp ἀνεῦρεν R, Suid. : \lecon{ἀνηῦρεν} Meineke}}% Austin comme Fritzsche lisent ἀνεῦρεν dans R, mais semblent se tromper : Hall, Coulon & Prato contre eux
\end{groupe}%
%% v. 438
\begin{groupe}%
\typeVers{$2$ troch.}\chantChoeur{πάντα δ᾽ ἐβάστασεν φρενὶ πυκνῶς τε\testimonia{438}{(πυκνῶς) Hesychius π 4335 πυκνῶς}}%
\end{groupe}%
%% v. 439a
\begin{groupe}%
\decoupeVers{a}\typeVers{$2$ troch.}\chantChoeur{ποικίλους λόγους ἀνηῦρεν}%
\end{groupe}%
%% v. 439b
\begin{groupe}%
\ligne{439}\decoupeVers{b}\typeVers{$2$ troch\cat}\chantChoeur{εὖ διεζητημένους·}%
\end{groupe}%
%% v. 440
\begin{groupe}%
\typeVers{$2$ troch.}\chantChoeur{ὥστ᾽ ἂν εἰ λέγοι παρ᾽ αὐτῆς\apparat{440}{\lecon{αὐτῆς} R : αὐτὴν Zanetti}}%
\end{groupe}%
%% v. 441a
\begin{groupe}%
\decoupeVers{a}\typeVers{$2$ troch.}\chantChoeur{Ξενοκλέης ὁ Καρκίνου, δο-\apparat{441}{Ξενοκλῆς R : \lecon{Ξενοκλέης} Bentley}%
\commentaire{441}{Ξενοκλῆς ὁ Καρκίνου.}{Xénoclès, le plus célèbre des trois carcinites était un poète tragique non dépourvu de talent puisqu'en 415 il l'emporta sur Euripide\marqueur{1}. Mais on lui reprochait de chercher à captiver son auditoire par des moyens peu relevés. Platon le comique l'appelle δωδεκαμήχανος [l'homme] \italique{aux douze postures}, surnom qu'on donnait à la courtisane Cyrène\marqueur{2}. Aristophane le traite de μηχανοδίφης, inventeur de \italique{trucs}. Il résulte d'un passage d'Athénée (IV, 134d), mis en rapport avec le nôtre, que Xénoclès cumulait avec la profession d'acteur et de poète tragique, celle d'orateur. Ce fut aussi le cas d'Eschine. Cela n'a rien d'étonnant, le métier d'acteur étant une excellente préparation à celui d'orateur. Il n'y a qu'à se souvenir de Démosthène et de ses cailloux. (Note de Willems ; voir aussi la note (t. I pp. 553-) qu'il a consacrée au banquet de Philocléon conté par Xanthias (\italique{Guêpes}, vv. 1299-1323) où il est question des carcinites.)}%
\debutNote{441}%
\note{1}{Aristophane se montre toutefois sévère avec Xénoclès : il est méchant comme ses vers (v. 169) et son art comme sa personne ne méritent pas de survivre (\italique{Grenouilles}, v. 86).}}%
\note{2}{Voir la note au vers 235.}%
\end{groupe}%
%% v. 441b
\begin{groupe}%
\ligne{441}\decoupeVers{b}\typeVers{$2$ troch.}\chantChoeur{κεῖν ἂν αὐτόν, ὡς ἐγᾦμαι,}%
\end{groupe}%
%% v. 442a
\begin{groupe}%
\decoupeVers{a}\typeVers{troch.}\chantChoeur{πᾶσιν ὑμῖν\commentaire{442}{πᾶσιν ὑμῖν.}{Les spectateurs.}}%
\end{groupe}%
%% v. 442b
\begin{groupe}%
\ligne{442}\decoupeVers{b}\typeVers{$2$ troch\cat}\chantChoeur{ἄντικρυς μηδὲν λέγειν.\testimonia{442}{(ἄντικρυς) Photius α 2210 ἄντικρυς ; \AelDion α 107}\\\passeLigne}%
\end{groupe}%
%% v. 443-458
\begin{groupe}%
\ledpb%
\personnageGrec{Στεφανό-\\πωλις}%
\typeVers{trim. iamb.}\trimIamb{\ligne{443}ὀλίγων ἕνεκα καὐτὴ παρῆλθον ῥημάτων.\apparat{443}{γυνὴ ἑτέρα \Rpc : Γυ. Β. Brunck : Κρίτυλλα Sommerstein : \lecon{Στεφανόπωλις} Rossella \espApp ἕνεκ᾽ αὐτὴ R : \lecon{ἕνεκα καὐτὴ} Porson}\\
τὰ μὲν γὰρ ἄλλ᾽ αὕτη κατηγόρηκεν εὖ·\\
ἃ δ᾽ ἐγὼ πέπονθα, ταῦτα λέξαι βούλομαι.\\
ἐμοὶ γὰρ ἁνὴρ ἀπέθανεν μὲν ἐγκύμων\apparat{446}{ἀνὴρ R : \lecon{ἁνὴρ} Brunck \espApp ἐν Κύπρῳ R : \lecon{ἐγκύμων} Bibi}
\commentaire{446}{ἐμοὶ γὰρ ἁνὴρ ἀπέθανεν ἐν Κύπρῳ.}{La précision du lieu de décès du mari (ἁνὴρ ἀπέθανεν ἐν Κύπρῳ) n'a pas manqué d'intriguer les modernes scholiastes.
Leur inclination naturelle étant portée au scabreux, ils ont pour une bonne partie voulu y voir un tour signifiant que la matrone athénienne, mère de cinq enfants, avait tué son mari au déduit avec ses appétits démesurés ; «~mourir de sa belle mort~», comme l'heureux \italique{Félix Faurenicator}. 
L'hypothèse est plaisante mais gratuite et donc à rejeter ! tous les lexiques consultés ne soufflent mot d'une telle expression qu'on aurait nécessairement relevée. 
D'autres se sont essayés à prendre ce vers littéralement. Là encore, ils se sont heurtés à un mur : les dernière expéditions sur l'île de Chypre feraient de notre procurette une septuagénaire au moins !.. et comme il ne semble pas y avoir de jeux de mots, il ne reste qu'une seule voie : une erreur de transcription du texte. 
(C'est d'autant plus vraisemblable que l'expression « ἀπέθανεν ἐν Κύπρῳ », si elle s'était véritablement trouvée dans le texte original, n'aurait pas manqué d'intriguer les commentateurs anciens qui, même sans y rien connaître, auraient ouvert les vannes de leur verve intarissable... au contraire : ce vers n'a étrangement pas laissé la moindre scholie ancienne derrière lui.)\\
L'intervention de la tresseuse de couronnes n'a donc rien de burlesque : elle essaie au contraire d'introduire du pathétique. Or, qu'y a-t-il de plus banal pour apitoyer son auditoire à peu de frais que de sortir le lieu-commun de la femme qui se retrouve veuve alors qu'elle est enceinte ? 
\citComm{ἐμοὶ γὰρ ἁνὴρ ἀπέθανεν μὲν \textbf{ἐγκύμων}\\
Mon mari mourut alors que j'étais enceinte \\
Me laissant cinq enfants à charge etc. etc.}}\\
παιδάρια πέντε καταλιπών, ἁγὼ μόλις\testimonia{447-9}{(μόλις -- τέως μὲν οὖν ἀλλ᾽) scholium\expo{T} Plato \italique{Hipparch.} 229d (μυρίναις pro μυρρίναις et μέν οὐκ ἄν pro μὲν οὖν ἀλλ᾽)}\\
στεφανηπλοκοῦσ᾽ ἔβοσκον ἐν ταῖς μυρρίναις.\testimonia{448}{(μυρρίνη $=$ μυρσίνη) Moeris μ 23}\\
\lednopb%
τέως μὲν οὖν ἀλλ᾽ ἡμικάκως ἐβοσκόμην·\testimonia{449}{Suidas τ 321 \espTest (ἡμικάκως) Pollux VI 161}\\
νῦν δ᾽ οὗτος ἐν ταῖσιν τραγῳδίαις ποιῶν\apparat{450}{ταῖσι R : \lecon{ταῖσιν} Brunck}\\
\lednopb%
τοὺς ἄνδρας ἀναπέπεικεν οὐκ εἶναι θεούς·
\commentaire{450-1}{ἐν ταῖσιν τραγῳδίαις ποιῶν / τοὺς ἄνδρας ἀναπέπεικεν οὐκ εἶναι θεούς·}{Petit florilège des \italique{\grec{γνῶμαι}} théoclastes d'Euripide :\\
\italique{Hécube}, v. 489-492 :%
\citComm{Ταλθύβιος — ὦ Ζεῦ, τί λέξω; πότερά σ᾽ ἀνθρώπους ὁρᾶν;\\
ἢ δόξαν ἄλλως τήνδε κεκτῆσθαι μάτην,\\
ψευδῆ, δοκοῦντας δαιμόνων εἶναι γένος\\
τύχην δὲ πάντα τἀν βροτοῖς ἐπισκοπεῖν;}%
\italique{Béllérophon}\marqueur{1} :%
\citComm{Βελλεροφόντης — φῆσίν τις εἶναι δῆτ᾽ ἐν οὐρανῷ θεούς;\\
οὐκ εἰσίν, οὐκ εἴσ᾽,}%
Prologue de \italique{Mélanippe la sage}\marqueur{2} (douteux) :%
\citComm{Ζεὺς ὅςτις ὁ Ζεύς, οὐ γὰρ οἶδα πλὴν λόγῳ\\
La véritable version euripidéenne étant peut-être\marqueur{3} :\\
Μελάνιππη — Ζεύς, ὡς λέλεκται τῆς ἀληθείας ὕπο,\\
Ἕλλην᾽ ἔτιχθ᾽, ὃς ἐξέφυσεν Αἴολον·}%
}%
\debutNote{450-1}%
\note{1}{Deux premières lignes du \nieme{10} fragment de l'édition Jouan-van Looy aux Belles Lettres ; fr. 286 de l'édition Kannicht ; fr. 385 Mette.}%
\note{2}{Fragment douteux car on recense différentes versions de ce début de prologue. Citation apocryphe ? retouche d'Euripide d'un vers ayant fait scandale ? quoi qu'il en soit, Jouan \& van Looy le mettent en tête de leur \italique{Mélanippe} sans le comptabiliser ; fr. 480 de l'édition Kannicht.}%
\note{3}{Premier fragment de l'édition Jouan-van Looy ; 481 chez Kannicht et 665a-c chez Mette.}%
\\\passeLigne}
\end{groupe}%
%% v. 443-458
\begin{groupe}%
\ledpb%
\trimIamb{ὥστ᾽ οὐκέτ᾽ ἐμπολῶμεν οὐδ᾽ εἰς ἥμισυ.\\
%\commentaire{452}{ὥστ᾽ οὐκέτ᾽ ἐμπολῶμεν οὐδ᾽ εἰς ἥμισυ.}{Le reproche apparent de l'oratrice est que l'impiété d'Euripide a entraîné... la chute des ventes de couronnes ; accusation on-ne-peut-plus frivole ! et qui déchaîne pourtant l'enthousiasme du chœur (\grec{κομψότερον ἔτ᾽ ἢ τὸ πρότερον ... οἷα κατεστωμύλατο ... φρένας ἔχουσα \&c.}). En outrec cette crise couronnomique est en pleine contradiction avec la fin de son discours~: quand elle dit s'en aller car elle a quantité de commandes à honorer. Il y a donc certainement jeu de mots et double discours ; je mettrais ma main à couper qu'il se trouve au vers 452, à rapprocher du vers 456. Je laisse à d'autres le soin d'expliquer que la vendeuse de couronnes dit à mi-mots se prostituer~; ce sont d'ailleurs les mêmes dignes professeurs qui, après avoir estimé son âge à une soixantaine d'année, nous lancent avec le même ton que cette mère de cinq enfants ferait le tapin dans le marché aux myrthes (hypothèses qui, conjuguées, ne méritent qu'un grand éclat de rire, le même que celui du film \italique{La traversée de Paris}). Restons sérieux~: s'il y a quelques sous-entendus salaces, ils ne sont certainement pas là pour que l'oratrice se décrédibilise elle-même~; faut-il le rappeler~: la cible est Euripide !
%Le poète est ridiculisé ailleurs (v. 1044-1048 des \italique{Grenouilles}) comme cocu notoire. Qu'il soit ici aussi joué sur ce plan et sur sa consommation excessive de certains légumes réputés par les Grecs ne pas rendre très vaillants~; ces déboires entraînant chez le sujet l'agraire aigreur de ses euripidéennes tragédies... voici qui, à défaut d'une démonstration, a au moins le mérite de la vraisemblance et de justifier les éloges du chœur. Et si, comme je le crois, il n'y a pas là de sous-entendu obscène, alors il s'agirait simplement d'une explication de l'impiété d'Euripide du fait que ce ne soit pas un homme du ciel, mais de la terre (le verbe \grec{πολεῖν} signifie \italique{retourner la terre} ; attaque identique au vers 840 des \italique{Grenouilles})}%
%
νῦν οὖν ἁπάσαισιν παραινῶ καὶ λέγω\testimonia{453-6}{Aulus Gellius XV 20.7}\apparat{453}{ἁπάσαισι R : \lecon{ἁπάσαισιν} Gellius \italique{Ν.Α.} 15.20}\\
τοῦτον κολάσαι τὸν ἄνδρα πολλῶν οὕνεκα·\\
ἄγρια γὰρ ἡμᾶς ὦ γυναῖκες δρᾷ κακά,\testimonia{455-6}{Plut. \italique{Moral.}, \italique{Compar. Arist. et Menandr.} (853c)}\\
ἅτ᾽ ἐν ἀγρίοισι τοῖς λαχάνοις αὐτὸς τραφείς.\apparat{456}{\lecon{τοῖς} Plut., \italique{Mor.}, 853C, Gellius \italique{N.Α.}15.20 : om. R}
\commentaire{455-6}{ἄγρια γὰρ ἡμᾶς ὦ γυναῖκες δρᾷ κακά / ἅτ᾽ ἐν ἀγρίοισι τοῖς λαχάνοις αὐτὸς τραφείς.}{Le procès d'Euripide a été ouvert et la parole est à l'accusation. Après un premier réquisitoire-fleuve, une seconde oratrice s'avance et prend la parole (vers 443-458).
Le reproche apparent de l'oratrice est que l'impiété d'Euripide a entraîné… la chute des ventes de couronnes ; accusation on-ne-peut-plus frivole ! et qui déchaînera pourtant l'enthousiasme du chœur quelques vers plus loin (κομψότερον ἔτ᾽ ἢ τὸ πρότερον, « plus intelligent encore que le premier [discours] » […] οἷα κατεστωμύλατο, « comme elle a bien parlé » […] φρένας ἔχουσα, « elle a de l'esprit » \&c.). 
Cette crise coronéconomique, terrible pour elle – comme nous pouvons nous en douter – n'est cependant peut-être pas à prendre au pied de la lettre. En outre, la fin de sa tirade semble en contradiction avec sa plainte : elle dit devoir s'en aller car elle a quantité de commandes à honorer. 
Les allusions, doubles sens et sous-entendus sont donc certainement au rendez-vous et nous allons tenter de défricher un peu tout cela !\\
Je laisse à d'autres commentateurs le soin d'expliquer que la vendeuse de couronnes avoue à mi-mot être une petite vertu rentrant de travailler ; songez donc ! évoquer à la fois Chypre, le marché aux myrtes, la tresse de couronnes ! c'en était bien assez pour que nos fouines dégainent mécaniquement leur lecture \italique{sensu obsceno} ! On repassera par contre pour obtenir de leur part un embryon d'explication : \ \ \	1° en quoi l'impiété d'Euripide nuirait-elle effectivement à la prostitution ?	\ \ \	2°	 comment une mère de cinq enfants pourrait-elle se permettre de faire le tapin sans provoquer l'hilarité générale\marqueur{1} ?
S'il y a quelques sous-entendus salaces, ils ne sont certainement pas là pour que l'oratrice se décrédibilise elle-même ; faut-il le rappeler : la cible est Euripide !\\
Le point culminant de l'attaque se trouve dans notre distique, aux vers 455-456.
Aristophane joue sur la répétition d'ἄγριος, qualifiant les mauvais comportements qu'Euripide impute à l'engeance féminine dans ses pièces\marqueur{2} et, au vers suivant, les λάχανα, les \italique{légumes} qui, lorsqu'ils sont ἄγρια, « sauvages », désignent des \italique{herbes sauvages}, de celles qui poussent sans qu'on les cultive (par opposition à ἥμερα). 
L'allusion porte sur le métier (supposé) de la mère d'Euripide, vendeuse d'herbes\marqueur{3} (ou de légumes) dont les poètes comiques firent leurs choux gras.\\
Mais ce jeu de mots à lui tout seul ne justifie pas – loin s'en faut – les éloges du chœur qui s'ensuivent\marqueur{4}. Il faut chercher davantage !\\
En consultant différents recueils de proverbes\marqueur{5}, je suis tombé sur cette citation du poète Dikaïogène\marqueur{6}, que Stobée recueillit dans son \italique{Anthologie}\marqueur{7} pour son chapitre consacré au respect filial.
\citComm{Αὐτὸς τραφεὶς δὲ τῶν φυτευσάντων ὕπο\\
καλῶς τὸν αὐτὸν ἔρανον αὐτοῖσιν νέμεις\\
Toi-même élevé [litt. : nourri] par tes parents [litt. : ceux qui t'ont semé],\\
tu leur partages comme de juste ton propre repas.}
Dikaïogène est postérieur de peu à Aristophane, mais le tour paraît effectivement proverbial. 
Si tel est le cas – au moins pour la première partie : αὐτὸς τραφεὶς ὑπὸ τῶν φυτευσάντων –, alors nul doute que le spectateur pouvait instinctivement compléter la fin du discours de la vendeuse de couronne, où αὐτὸς τραφεὶς avait sciemment été rejeté à la fin : ἅτ᾽ ἐν ἀγρίοισι τοῖς λαχάνοις αὐτὸς τραφείς … ὑπὸ τῶν φυτευσάντων… et comprendre que les parents d'Euripide ayant semé de la mauvaise graine – un comble pour sa mère, la λαχανοπωλήτρια ! –, il n'est guère étonnant, d'après le proverbe, que le rejeton \italique{sème la zizanie en retour}\marqueur{8}. 
L'arbre se jugeant à ses fruits, et les fruits d'Euripide étant infestés de petits vers féminassassins, il est naturel que les femmes les aient en travers de la gorge et cherchent \italique{en retour} à lui faire rendre gorge !\\
Je ne prétends pas avoir littéralement rendu les images suggérées par l'oratrice, ni que ces jeux sur les expressions françaises soient aussi bons que l'original. Néanmoins, j'espère simplement que par cette méthode le lecteur appréhendera mieux le texte d'Aristophane qu'en recourant à une longue analyse des divers sens et interprétations possibles des mots `ἄγρια', `αὐτὸς τραφεὶς' (\italique{nourri} ou \italique{éduqué}), `φυτευσάντων' (\italique{semeurs} ou \italique{géniteurs}) et `τὸν αὐτὸν ἔρανον' (\italique{le même repas}, \italique{écot} ou \italique{retour}).\\
Vous trouvez cela un peu léger ? que j'extrapole ? pourtant, c'est exactement la même figure de rhétorique que Dion Cassius (\italique{Histoire romaine}, XLVI, 4) met dans la bouche d'un orateur qui répond dans le même ton à Cicéron et ses violentes \italique{Philippiques} contre Marc Antoine. Cicéron... qui est-ce, au juste ?
\citComm{ὁ πατὴρ αὐτῷ ὁ κναφεύς [ ... ] ἄνθρωπος ἀγαπητῶς ἔκ τε τούτων καὶ ἐκ τῶν πλυνῶν διατρεφόμενος [ ... ] Ἐν οἷς \textbf{αὐτὸς τραφεὶς} οὐκ ἀπεικότως τοὺς κρείσσονας αὐτοῦ καὶ πατεῖ καὶ πλύνει.\\
\francais{Son père était foulon, lui qui avait grand-peine à vivre de ce métier et de son lavoir. À la suite d'une telle éducation, non sans invraisemblance, ce personnage foule aux pieds ceux qui valent mieux que lui et leur lave la tête.}}
Profitons à présent de cette note sur les λάχανα pour élargir notre champ botanique de recherche à percer le véritable sens des différentes attaques d'Aristophane lancées contre Euripide, ce fils de la λαχανόπωλις. Comprendre ce qu'il lui reproche exactement à travers ces attaques nous prémunira du même coup de bien des hypothèses fantaisistes\marqueur{9}.\\
Nous allons voir que les multiples attaques contre Euripide le fils de la vendeuse d'herbes (et notamment de σκάνδιξ, identifié à notre \italique{scandix peigne de Vénus}\marqueur{10}) portent contre son style, tout en finasseries, qui essaye vainement d'enrober de jolies phrases ses billevesées.\\
Rappelons pour mémoire qu'au vers 387, la précédente oratrice, plus prolixe, avait commencé son discours en désignant Euripide du sobriquet d'Εὐριπίδου τοῦ τῆς λαχανοπωλητρίας, \italique{fils de la marchande de légumes}.\\
Au vers 478 des \italique{Acharniens}, alors que Dikaïopolis emprunte à Euripide tout un fatras de matériel puisé dans ses propres drames, voilà qu'il déclenche le courroux d'Euripide en lui demandant :
\citComm{Δικαιόπολις – σκάνδικά μοι δὸς μητρόθεν δεδεγμένος.\\
Εὐριπίδης – ἁνὴρ ὑβρίζει· κλῇε πηκτὰ δωμάτων.\\
\francais{Dikaïopolis – Donne-moi de ce scandix que tu tiens de ta mère.\\
Euripide – Il m'insulte. « Ferme sur lui les ais solides de la porte. »}}
Qu'est-ce que Dikaïopolis veut sous-entendre avec ce \italique{scandix}, qu'Euripide posséderait par héritage maternel ? Le début des \italique{Cavaliers} aidera à éclaircir cette affaire : les deux personnages présents sur scène sont les généraux Nicias et Démosthène, généraux qui avaient préparé la victoire de Sphactérie (victoire éclatante par la capture de la garnison spartiate). 
Mais ces deux généraux s'étaient fait voler la vedette par le braillard Cléon débarqué avec un contingent sur place au dernier moment et qui – coup du sort – l'avait emporté\marqueur{11}. Ils ressassent leur amertume au début de la pièce en se renvoyant la balle pour trouver une idée de salut : que faire ? — Dis-le ! — Non, toi dis-le\marqueur{12}… Nicias s'interroge : comment tourner cela précieusement à la manière d'Euripide (κομψευριπικῶς). Démosthène rétorque alors (vers 19-20) :
\citComm{Δημοσθένης – μὴ 'μοί γε, μὴ 'μοί, μὴ διασκανδικίσῃς·\\
ἀλλ᾽ εὑρέ τιν᾽ †ἀπόκινον† ἀπὸ τοῦ δεσπότου.\\
Démosthène – Non pas, non pas, ne diascandise pas !\\
mais imagine quelque branle de sortie qui le décroche du maître\marqueur{13}.}
Les dictionnaires traduisent le verbe διασκανδικίζω par « mâcher du cerfeuil »\marqueur{14} avec un sens imagé de « parler comme Euripide », sens directement tiré du vers 19 des \italique{Cavaliers}… nous tournons en rond et n'en sommes guère avancés !\\
Le verbe διασκανδικίζω n'est pas un hapax d'Aristophane ; il se trouve aussi dans un fragment de Téléclide, un autre poète comique contemporain d'Euripide. Ce fragment\marqueur{15} nous a été conservé par Athénée\marqueur{16} :
\citComm{Ξυγγενέσθαι διὰ χρόνου †λιπαρείτω† με\\
δρυπεπέσι μάζαις καὶ διασκανδικίσαι.}
Tout ce que l'on peut inférer du contexte chez Athénée, c'est que l'adjectif δρυπεπής, \italique{bien mûr}, lorsqu'il est substantivé comme ici, désigne une \italique{olive}. 
Il me semble bien avoir trouvé un sens possible à διασκανδικίζω qui rendrait compte des deux passages : non pas comme il est indiqué dans les dictionnaires « se nourrir de σκάνδιξ » mais plutôt : « assaisonner à l'aide de σκάνδιξ » et, dans un sens imagé, « agrémenter », « enjoliver » ; le σκάνδιξ jouant un rôle analogue à notre \italique{piment}, et les multiples sens du verbe « pimenter » (penser aussi au verbe \italique{épicer}).\\
Je traduis donc : « Il m'importune (λιπαρεῖ ?) depuis longtemps pour ajouter des olives aux galettes ainsi que [les] agrémenter de scandix ».\\
Ce sens pour διασκανδικίζω d' « agrémenter », « épicer [un discours de circonlocutions] » avec une nuance péjorative comme cela paraît être le cas chez Téléclide convient aussi au passage des \italique{Cavaliers}.\\
On peut ainsi proposer la traduction pour le vers 19 des \italique{Cavaliers} :
\citComm{Δημοσθένης – μὴ 'μοί γε, μὴ 'μοί, μὴ διασκανδικίσῃς·\\
Démosthène – Arrête-moi ça : ne me saupoudre pas ça de scandix !}
Le vers suivant : « ἀλλ᾽ εὑρέ τιν᾽ †ἀπόκινον† ἀπὸ τοῦ δεσπότου » laisse cependant circonspect : on ne comprend pas bien le rapport entre le scandix et l'ἀπόκινος, défini par le scholiaste comme une danse licencieuse et qui fait entendre par son étymologie le verbe ἀπο-κινέω, \italique{se retirer}\marqueur{17}. Certes tout ce passage des \italique{Cavaliers} est porté à la licence… Certes Cléon est raillé ailleurs pour ses mœurs infâmes ! 
Mais l'on ne perçoit pas pour autant la transition entre les vers 19 et 20, pourtant soulignée par le « ἀλλ᾽ » au début du vers 20\marqueur{18}. 
Comme toujours devant les difficultés, nous savons que nous pouvons compter sur notre bon vieil et fidèle ami : Hésychius (α 6416) ! Voici ce que grand érudit a répondu à nos questions :%
\citComm{ἀπόκυνον: μάζα μεμαγμένη φαρμάκῳ, πρὸς ἀναίρεσιν κυνῶν. ἢ εἶδος βοτάνης\\
apokynon : galette préparée avec du poison, pour la destruction des chiens. \italique{Item} : sorte d'herbe.}%
Et la lumière fut : l'ἀπόκυνον, comme le σκάνδιξ, est une \italique{herbe sauvage}… à cette différence que celui-ci sert de condiment aux pauvres\marqueur{19} quand celui-là nous débarrasse des nuisibles : chiens selon Hésychius, Cléon pour Aristophane. 
Ce poison n'est d'ailleurs pas choisi par hasard : il fait la transition avec le scandix, l'autre herbe et renvoie Cléon à son canin statut de \italique{Chien de Dêmos}, animal flatteur et cajoleur (tout au long des \italique{Cavaliers}, Cléon est moqué de la sorte – sûrement par reprise ironique de quelque propos du véritable Cléon).\\
Je propose donc la lecture suivante pour les vers 19-20 des \italique{Cavaliers} :%
%\eff{Μὴ ἀποσκορακίσῃς με Souda + Hésychius mu1136}
\citComm{Δημοσθένης – μὴ 'μοί γε, μὴ 'μοί, μὴ διασκανδικίσῃς·\\
ἀλλ᾽ εὑρέ τιν᾽ ἀπόκυνον ἀπὸ τοῦ δεσπότου. \decalage{aaaaaa} ἀπόκινον codices omnes : \lecon{ἀπόκυνον} Bibi ap. Hesych.\\
Démosthène – Arrête-moi ça : cesse de condimentir !\\
mais trouve-moi plutôt quelque mort-aux-chiens qui en débarrasse son maître.}%
En conclusion, il semblerait que les lazzis visant Euripide, le fils de la vendeuse d'herbe, soient davantage des métaphores stylistiques que basés sur des faits réels\marqueur{20} ; la vendeuse d'herbe ayant sûrement la réputation de parler de façon détournée, d'être une embobineuse, comparable en cela à la réputation chez nous des harengères dont le franc-parler est devenu proverbial.
De la même façon, certaines \italique{Vies} d'Euripide en font le fils d'un κάπηλος\marqueur{21}, d'un \italique{boutiquier} dirait-on en français ; là encore, le détail est certainement à prendre de façon imagée\marqueur{22}.\\
Intéressons-nous maintenant plus particulièrement au σκάνδιξ : parmi tout l'étal que pouvait proposer ces vendeuses de verdure, seul le σκάνδιξ est raillé par les Comiques\marqueur{23} (je tiens pour rien la pseudo-reconnaissance de Ménélas par Hélène au vers 910 `ἐκ τῶν ἰφύων', d'après les \italique{herbes} qu'il arborerait ; lecture inepte et que nous balaierons le moment venu).
Quelle est donc la particularité de ce scandix ?
Dans un article fort intéressant\marqueur{24}, Suzanne Saïd, à la suite de P. Fornaro et G. Arrighetti, opte elle aussi pour une interprétation stylistique de la généalogie d'Euripide : « Un auteur qui a un style «~vulgaire~»\marqueur{25} et représente le \italique{genus tenue} (ou \italique{subtile}) ou, pour parler grec, le Χαρακτὴρ ἰσχνός\marqueur{26}, ne peut qu'être le fils d'un κάπηλος et d'une marchande de σκάνδιξ, surtout quand on sait, par Dioscoride, que cette plante avait des effets laxatifs (εὐκοίλιον). 
De fait, dans les \italique{Acharniens}, cette fameuse σκάνδιξ qu'Euripide tient de sa mère est étroitement associée aux «~plantes amaigrissantes~» (ἰσχνὰ φυλλεῖα)\marqueur{27}. 
Dans les \italique{Grenouilles}, le poète se vante d'ailleurs d'avoir joué les diététiciens et d'avoir «~fait maigrir~» (ἰσχναίνειν) «~au moyen de versiculets, de marches et de bettes blanches~»\marqueur{28} (c'est-à-dire d'une plante qui a exactement les mêmes propriétés que la σκάνδιξ)\marqueur{29} une tragédie qu'Eschyle avait rendue obèse. 
Aristophane fait peut-être même ici d'une pierre deux coups : grâce à cette généalogie qui fait de lui le fils d'une «~marchande~» (-πωλις), il associe Euripide, le «~nouveau~» poète tragique, aux nouveaux politiciens qui appartiennent tous à la catégorie des «~marchands~» (-πώλης). »\\
L'hypothèse d'un Euripide raillé pour son σκάνδιξ en raison des vertus amaigrissantes de cette plante méritait qu'on s'y arrêtât longuement et justifiait la copieuse citation. Elle nous semble néanmoins insuffisante : pourquoi l'Euripide des \italique{Acharniens} prendrait-il la mouche lorsqu'on lui parle de son scandix quand l'Euripide des \italique{Grenouilles} se vante de recourir lui-même aux plantes amaigrissantes ? De plus, cela ne nous explique pas davantage la redondance du thème du seul scandix alors qu'en ne consultant que Théophraste, Pline ou Dioscoride, on pourrait établir une liste à n'en plus finir des plantes réputées amaigrissantes ou laxatives.
L'attaque ne semble pas non plus porter le fait qu'Euripide soit un \italique{fils de la terre}, \latin{terrae filius}, c'est-à-dire dont les origines sont obscures et l'arbre généalogique difficile à établir\marqueur{30}, mais bien sur le style d'Euripide.\\
Je crois avoir établi dans cette note – note qui tend à l'obésité malgré le recours massif au σκάνδιξ – que si Aristophane parle si souvent du σκάνδιξ d'Euripide, c'est parce que le Grec usait lui-même d'un verbe correspondant peu ou prou à notre « enfariner » : διασκανδικίζειν.\\
Cette trop longue note touchant enfin à sa fin, nous nous permettrons quelque aventureuse conjecture sur l'origine de cette utilisation imagée du σκάνδιξ : il nous semble que si le grec a choisi cette plante, c'est non pas pour ses vertus (réelles ou supposées), mais bien plutôt pour des raisons de sonorité. Cette intuition m'est venue de la consultation de l'\latin{Etymologicum Genuium}\marqueur{31} à l'entrée σκάνδαλον, la \italique{chausse-trappe}, pris ici en un sens imagé : σκάνδαλον: τὸ σκάζειν περὶ τὴν ἀλήθειαν. 
Image saisissante et évocatrice s'il en est ! « boîter autour de la vérité », « tourner autour du pot-aux-roses ». 
Elle s'appuie sur un rapprochement pseudo-étymologique\marqueur{32} entre le mot σκάνδ-αλον et le verbe σκάζειν, boîter… le σκάνδ-ιξ ne pouvait-il pas suggérer la même image : celle de \italique{tordre}, \italique{rendre boiteuse} ou de \italique{camoufler} la vérité en la saupoudrant d'une herbe de vil prix ?}%
\debutNote{455-6}%
\note{1}{Pour ce genre de situation, le rire du soldat allemand dans la \italique{Traversée de Paris} est, ce me semble, la réaction la plus saine qui soit.}%
\note{2}{Aristophane, au vers 349 des \italique{Nuées}, qualifie ainsi d' ἄγριος le chevelu Hiéronyme, un poète dithyrambique, fils de Xénophane qui passait pour pédéraste ; et le scholiaste de noter qu'on appelait ἀγρίους (et débauchés) les pédérastes : ἀγρίους (καὶ κολλοποδιώκτας) ἐκάλουν τοὺς παιδεραστάς. De même, à l'entrée Ἀγρίους de la \italique{Souda}, on peut lire : Αἰσχίνης ἐν τῷ κατὰ Τιμάρχου τοὺς σφόδρα ἐπτοημένους περὶ τὰ παιδικὰ οὕτως ὠνόμασε. καὶ Μένανδρος δὲ ἄγριον εἶπε κυβευτὴν τὸν λίαν περὶ τὸ κυβεύειν ἐσπουδακότα, « Eschine dans son \italique{Contre Timarque} nomme ainsi les hommes en proie à une excessive attirance pour les enfants. Et Ménandre dit qu'un joueur est ἄγριον, celui qui est trop porté au jeu ». L'on voit que le terme ἄγριος, au sens moral, désigne un \italique{intempérant}.}%
\note{3}{L'oratrice précédente, au vers 387, avait déjà traité Euripide de fils de la λαχανοπωλήτρια.}%
\note{4}{En rester là donnerait raison au grave Plutarque qui citait ces deux vers comme exemple d'un jeu de mots qu'il goûtait fort peu (voir ses \italique{Œuvres morales}, 853C, \italique{Abrégé de la comparaison d'Aristophane et de Ménandre}). Je ne suis pas sûr que cet inénarrable raseur, le Cicéron des lettres grecques, ait bien saisi intégralement le détournement d'Aristophane que nous allons voir.}%
\note{5}{\italique{Corpus paroemiographorum graecorum}, II, ed. Leutsch, Apostolii centuria V 30a, p. 341.}%
\note{6}{L'on ne possède pas grand-chose sur ce poète grec, auteur de tragédies et dithyrambes, qui vécut probablement au \nieme{IV} siècle av. J.-C. Tout ce que l'on sait de lui est regroupé dans les \italique{Tragicorum Graecorum Fragmenta}, vol. I, ed. Snell \& Kannicht. Il y est classé en tant que « poète mineur » n° 52.}%
\note{7}{Stobée, livre IV, chapitre 25, 6.}%
\note{8}{L'on trouve chez Démosthène dans son \italique{Contre Neæra}, 8, l'expression ironique τὸν αὐτὸν ἔρανον ἀποδοῦναι, \italique{payer quelqu'un de la même monnaie}.}%
\note{9}{Voir la note sur le vers 910 et ces étranges plantes, ἰφύων, des \italique{nielles des blés} (Agrostemma githago), qui ont envahi bien mal à propos jusqu'aux sources les plus fiables.}%
\note{10}{Pour un article passionnant sur quelques plantes apparaissant chez Aristophane, passionnant car il mêle avec un rare bonheur de savantes études de textes et des connaissances botaniques fort appréciables, le lecteur se reportera à l'article de Suzanne Amigues : « Quelques légumes de disette chez Aristophane et Plutarque » paru en 1988 dans le \italique{Journal des Savants} (pp. 157-171). Passionnant quand bien-même je suis obligé de m'inscrire en faux contre l'interprétation qu'elle donne du vers 910 des \italique{Thesmophories} (nous en reparlerons).}%
\note{11}{Pour plus de détails, le lecteur pourra utilement se reporter au quart livre de Thucydide, chapitres 28 et suivants ainsi qu'à la longue tirade du Démosthène des \italique{Cavaliers}, à partir du vers 40.}%
\note{12}{Aristophane fait-il allusion à un complot qui aurait été tramé contre Cléon ? contre ce va-t-en-guerre de Cléon qui empêchait toute négociation de paix au moment-même où Athènes était en position de force pour se sortir favorablement de cette guerre fratricide qui n'avait que trop duré.}%
\note{13}{Le \italique{maître}, c'est Δῆμος Πυκνίτης, le Populo-citoyen-roi-de-l'Assemblée, mené par le bout du nez par la pire espèce rustique d'embabouineur qui soit : le Marchand-de-cuir Cléon.}%
\note{14}{Bailly. Idem pour le Liddell-Scott : « feed on chervil (σκάνδιξ) ». Quant au Diccionario Griego-Español, il donne : « dar de comer perifollo ».}%
\note{15}{Fragment 40(38) dans l'édition de Kassel-Austin, \italique{Poetae Comici Graeci}, vol. VII.}%
\note{16}{\italique{Deipnosophistes}, VI, 56d.}%
\note{17}{Je dois au grand Alponse Willems cette jolie trouvaille qui évoque tout en un : le « branle de sortie »… trouvaille qui porte malheureusement sur une transmission fautive du texte comme nous l'allons voir tout à l'heure.}%
\note{18}{Et je n'ai encore jamais trouvé chez Aristophane d'analogue à notre éculée « sans transition ».}%
\note{19}{Le scandix semble être une épice peu appréciée : Andocide parle ainsi des privations des premières années de la guerre du Péloponnèse : μηδὲ ἄγρια λάχανα καὶ σκάνδικας ἔτι φάγοιμεν, « puissions-nous ne plus manger des herbes sauvages ni de scandix ».}%
\note{20}{J'accorderais ainsi plus de crédit aux biographies d'Euripide qui le présentent comme ayant eu des parents riches que celles qui en font – peut-être d'après des interprétations trop littérales des Comiques – un \italique{self-made-man} né dans l'indigence.}%
\note{21}{Voir les \italique{Testimonia} des \italique{Vitae Euripidis} réunis par Kannicht au cinquième volume, consacré à Euripide, des \italique{Tragicorum Graecorum Fragmenta}.}%
\note{22}{Pour le sens imagé de κάπηλος, voir par exemple ce qu'en dit Hésychius : κίβδης· κακοῦργος, κάπηλος, « frelateur : malfaiteur, \italique{boutiquier} » ; on trouve encore ceci : Φοινικελίκτην καὶ λόγων ἀλαζόνα: ἀπατηλὸν καὶ κάπηλον, « Phénichien et fort en gueule : un fourbe et un \italique{boutiquier} » ; les Phéniciens étant un peuple qui se livrait au trafic, la langue populaire leur rendait justice. Ainsi pour le français ; en ouvrant par exemple \italique{Le livre des proverbes français}, proverbes compulsés par Le Roux de Lincy, et allant visiter la ville chère à Rousseau \& Calvin, croquée comme il se doit par le blason populaire, on trouve : « Trois Juifs font un Balois, trois Balois font un Genevois » ou encore : « Les Genevois ont vertu de cent lieues de loing ». Chez Aristophane, ce sont plus particulièrement les débitants de vin qui sont accusés de gruger les mesures légales.}%
\note{23}{Nous n'avons pas encore parlé d'une citation non identifiée à l'entrée Σκάνδιξ de la \italique{Souda} : καὶ μή μοι τὰ τοῦ σκανδικοπώλου παρατραγῴδει, σεμνὸν τὸ τῆς Πειρήνης ὕδωρ ἀποκαλοῦντος, « et ne m'use pas de l'emphase du vendeur de scandix qui appelle `sacrée' la fontaine de Pirène ». La cible est clairement Euripide : la citation peut faire allusion au vers 69 de sa \italique{Médée} ou aux vers 205-206 des \italique{Troyennes}. Il est toujours délicat de proposer des hypothèses sur des fragments aussi brefs… proposons tout-de-même notre hypothèse qui vaut ce qu'elle vaut – c'est-à-dire pas grand-chose : cette fontaine située sur l'Acrocorinthe délivrait effectivement une eau pure et abondante (cf. ce qu'en disent Pausanias, Strabon et Suidas à l'entrée Πειρήνη ainsi que cette citation d'Antipater prise dans l'\italique{Anthologie grecque} : Λαί̈δ' ἔχω πολιῆτιν ἐϋζώνοιο Κορίνθου, Πειρήνης λευκῶν φαιδροτέρην λιβάδων, « Je tiens Laïs, citoyenne de Corinthe, bien ceinturée, plus brillante que les gouttes blanches de Pirène ») ; ce n'est donc pas sur la qualité de l'eau qu'Euripide enrobe de scandix la vérité… peut-être sur ce qu'on y pratiquait ? La fontaine de Pirène se trouvait juste en contrebas du temple d'Aphrodite dont certaines de ses préposées se livraient à une prostitution « sacrée » (la réputation des Corinthiennes était telle que les Grecs en ont inventé le verbe κορινθιάζειν).}%
\note{24}{« Entre l'artisan et le politique », in : \italique{Dieux, héros et médecins grecs} (2001), pp. 65-96.}%
\note{25}{Tzetzes (dans une scholie citée par Fornaro) parle, à propos d'Euripide, de κάπηλα γράμματα et oppose ce poète κάπηλον et ἀγοραῖον à Eschyle. De même Lucien (\italique{De Historia conscribenda}, 44) emploie ces deux adjectifs pour caractériser un vocabulaire vulgaire (ὀνόματα ἀγοραῖα καὶ καπηλικά).}%
\note{26}{Quintilien \italique{Inst. orat.} XII, 10, 58.}%
\note{27}{\italique{Acharniens}, vers 469.}%
\note{28}{\italique{Grenouilles}, vers 941-942.}%
\note{29}{Voir Dioscoride, \italique{De Materia medica Libri V}, ed. M. Wellmann, Berlin 1907, vol. I, livre II 123, 138.}%
\note{30}{Cf. lorsque Ion demande à Xouthos (\italique{Ion}, v. 542) :%
\citComm{Ἴων -- γῆς ἄρ᾽ ἐκπέφυκα μητρός;\\
J'ai donc eu la terre pour mère ?}%
L'adage fut recensé par Érasme : (n° 786, \latin{terrae filius}).
Est-ce là qu'Eschyle veut en venir dans les \italique{Grenouilles} (v. 840), lorsqu'il lance à Euripide la cinglante réponse : ἄληθες ὦ παῖ τῆς ἀρουραίας θεοῦ.}%
\note{31}{Lexique du \nieme{IX} siècle souvent associé à d'autres, postérieurs mais plus complets : l'\latin{Etymologicum magnum} et l'\latin{Etymologicum Symeonis}.}%
\note{32}{Je me vois forcé de préciser \italique{pseudo}-étymologique, car mes velléités de rapprocher étymologiquement σκάνδαλον de σκάζειν furent foudroyées par Alexandre Perret, dont j'eus la chance d'assister à des cours de linguistique, cours qui ressemblaient davantage à d'enivrants voyages à travers les contrées et les parlers d'Europe (et ce jusqu'aux confins de l'Asie). Voici ce que m'a très aimablement répondu M. Perret : « Après recherches fructueuses, je pense avoir trouvé une réponse dans le \italique{Lexikon der Indogermanischen Verben} (LIV) mais il ne semble pas que les deux racines soient liées. Σκανδάλον semble en effet rattaché à un thème *skend- qui signifierait « bondir » (on a skándati en sanskrit pour « sauter »), mais le suffixe -αλον me dérange dans la mesure où ça semble pré-grec. Aucune étymologie n'est assez bien solide pour ce mot qui peut être un emprunt. Quelle motivation pour passer de « bondir » à « faire tomber » ? À la limite, un autre thème *sk'end- (avec un k- palatal donc invisible en grec mais pas en sanskrit) avec le sens parfois de « faire tomber » (skr. Chándayati: « tomber »), est en soit beaucoup plus plausible !! Mais pour σκάζω, on rapproche ça d'une racine *skeng- avec effectivement le sens de « boiteux, tors » (All. \italique{Hinken}). Ce qui m'étonne est qu'Hesychius recense à la fois σκάνδαλον: \italique{piège à souris} et σκάνδαλος: \italique{obstacle}, mais est-ce que les pièges à souris étaient des trappes où tomber, ou bien des claquoirs comme à notre époque ? Ça mérite investigation. En soit passer de *skend- à σκάζω avec un suffixe -yo de présent n'est pas choquant. »}\\
ἀλλ᾽ εἰς ἀγορὰν ἄπειμι· δεῖ γὰρ ἀνδράσιν\testimonia{457-8}{(δεῖ γὰρ -- εἴκοσιν) Suidas σ 1591 (sed ἀνδράσι et εἴκοσι pro ἀνδράσιν et εἴκοσιν)}\\
πλέξαι στεφάνους συνθηματιαίους εἴκοσιν.\testimonia{458}{Athenaeus XV 680c \espTest (στεφάνους συνθηματιαίους) Pollux VII 200 \espTest (συνθηματιαίους) Hesychius σ 2624 συνθηματιαῖοι}\commentaire{458}{στεφάνους συνθηματιαίους.}{Athénée (XV, 680c), dans le long chapitre consacré aux couronnes de ses \italique{Deipnosophistes}, parle de ces couronnes \italique{synthématiées} :%
\citComm{συνθηματιαῖοι στέφανοι, ἠργολαβημένοι καὶ ἐκδόσιμοι. Ἀριστοφάνης ἐν Θεσμοφοριαζούσαις.}%
Glose identique chez Hésychius (σ 2624) :%
\citComm{συνθηματιαῖοι: οἱ ἠργολαβημένοι καὶ ἐκδεδομένοι, ἐκδόσιμοι.}}\\\passeLigne}%
\end{groupe}%
%% v. 459-465
\begin{groupe}%
\ledpb%
\personnageGrec{Χορός}%
\typeVers{$2$ toch.}\chantChoeur{ἕτερον αὖ τι λῆμα τοῦτο\testimonia{459-60}{(ἕτερον -- ἀναπέφηνεν) Suidas λ 441 (ἀναπέφηνε pro ἀναπέφηνεν)}}%
\end{groupe}%
%% v. 460a
\begin{groupe}%
\decoupeVers{a}\typeVers{$2$ troch.}\chantChoeur{κομψότερον ἔτ᾽ ἢ τὸ πρότερον}%
\end{groupe}%
%% v. 460b
\begin{groupe}%
\ligne{460}\decoupeVers{b}\typeVers{troch.}\chantChoeur{ἀναπέφηνεν.}%
\end{groupe}%
%% v. 461
\begin{groupe}%
\typeVers{choria. $+$ ia.}\chantChoeur{οἷα κατεστωμύλατο\apparat{461}{\lecon{κατεστωμύλατο} R : κἀστωμύλατο Dobree}}%
\end{groupe}%
%% v. 462
\begin{groupe}%
\typeVers{$2$ troch.}\chantChoeur{οὐκ ἄκαιρα, φρένας ἔχουσα\apparat{462}{ἄκερα R : \lecon{ἄκαιρα} Zanetti}}%
\end{groupe}%
%% v. 463
\begin{groupe}%
\typeVers{$2$ troch.}\chantChoeur{καὶ πολύπλοκον νόημ᾽, οὐδ᾽\apparat{463}{\lecon{πολύπλοκον} R : τι πολύπλοκον Dobree : πολύστροφον Hermann}%
\commentaire{463}{πολύπλοκον.}{Les diverses corrections qui ont été proposées pour faire rentrer de force ce vers dans le mètre trochaïque ne me paraissent pas indispensables : on peut simplement considérer qu'il n'y a pas de \italique{correptio attica} et scander :%
\citComm{καὶ πολύπ|λοκον νόημ᾽, οὐδ᾽\\
$ - \cup - \cup - \cup - - $}%
}%
}%
\end{groupe}%
%% v. 464
\begin{groupe}%
\typeVers{$2$ troch.}\chantChoeur{ἀσύνετ᾽ ἀλλὰ πιθανὰ πάντα.}%
\end{groupe}%
%% v. 465a
\begin{groupe}%
\decoupeVers{a}\typeVers{$2$ troch.}\chantChoeur{δεῖ δὲ ταύτης τῆς ὕβρεως ἡ-}%
\end{groupe}%
%% v. 465b
\begin{groupe}%
\ligne{465}\decoupeVers{b}\typeVers{troch.}\chantChoeur{μῖν τὸν ἄνδρα}%
\end{groupe}%
%% v. 465c
\begin{groupe}%
\ligne{465}\decoupeVers{c}\typeVers{$2$ troch\cat}\chantChoeur{περιφανῶς δοῦναι δίκην.\\\passeLigne}%
\end{groupe}%
%% v. 466-519
\begin{groupe}%
\ledpb%
\personnageGrec{Μνησίλοχος}%
\trimIamb{τὸ μὲν ὦ γυναῖκες ὀξυθυμεῖσθαι σφόδρα\testimonia{466}{(ὀξυθυμεῖσθαι) Photius ο 380 ὀξυθυμεῖσθαι}%
\commentaire{466 sqq.}{}{Scholiaste : \grec{ἐπιτηέδες ὀ κηδεστὴς ὑπὲρ τὰς ἄλλας γυναῖκας γυναικίζεται, ἵνα διὰ τούτου λαθῃ}... prudence qui va bientôt le quitter.}\\
Εὐριπίδῃ, τοιαῦτ᾽ ἀκουούσας κακά,\apparat{467}{ἀκουούσαις R : \lecon{ἀκουούσας} Zanetti (\italique{audientes} Divus)}\\
οὐ θαυμάσιόν ἐστ᾽, οὐδ᾽ ἐπιζεῖν τὴν χολήν.\testimonia{468}{(θαυμάσιον $=$ θαυμαστόν) Moeris θ 16 $\simeq$ Hesychius θ 146 θαυμάσιος}\\
καὐτὴ γὰρ ἔγωγ᾽, οὕτως ὀναίμην τῶν τέκνων,\testimonia{469}{Suidas ο 330 (sed ἐγώ pro ἔγωγ᾽) ; cf. (ὀναίμην) Hesychius ο 852 ὀναίμην ; \Syn ο 162 ὀναίμην $=$ Photius ο 332}\apparat{469}{ὀνοίμην R : \lecon{ὀναίμην} Suid. ο 330}\\
μισῶ τὸν ἄνδρ᾽ ἐκεῖνον, εἰ μὴ μαίνομαι.\testimonia{470}{Suidas μ 328 (τἄνδρ᾽ pro τὸν ἄνδρ᾽)}\\
ὅμως δ᾽ ἐν ἀλλήλαισι χρὴ δοῦναι λόγον·\apparat{471}{ἀλλήλοισι R : \lecon{ἀλλήλαισι} Brunck}%
\commentaire{471}{ἐν ἀλλήλοισι χρὴ δοῦναι λόγον.}{Même expression dans Hérodote (I, 97) : « καὶ ἐδίδοσαν σφίσι λόγον ».}\\
αὐταὶ γάρ ἐσμεν, κοὐδεμί᾽ ἔκφορος λόγου.\\%\apparat{472}{\lecon{ἔκφορος} R : ἐκφορὰ Valckenaer ad Eur. \italique{Hipp.} 294 \ligneApp{473} κεῖνον R, Suid τ 683 : \lecon{᾽κεῖνον} Meineke}\\
τί ταῦτ᾽ ἔχουσαι ᾽κεῖνον αἰτιώμεθα\testimonia{473-4}{(τί ταῦτ᾽ -- φέρομεν) Suidas τ 683}\apparat{473}{κεῖνον R, Suid τ 683 : \lecon{᾽κεῖνον} Meineke}%
\commentaire{473}{τί ταῦτ᾽ ἔχουσαι κεῖνον αἰτιώμεθα.}{Ce vers rappelle celui des \italique{Acharniens} (vers 514) :
\citComm{τί ταῦτα τοὺς Λάκωνας αἰτιώμεθα;}
Ils pourraient donc bien être tous deux une reprise parodique du \italique{Télèphe} d'Euripide, pièce maintes fois détournée ici et là.}\\
βαρέως τε φέρομεν, εἰ δύ᾽ ἡμῶν ἢ τρία\apparat{474}{ἢ δύ᾽ R : \lecon{εἰ δύ᾽} Grynaeus (\italique{si} Divus)}\\
κακὰ ξυνειδὼς εἶπε δρώσας μυρία;\\
ἐγὼ γὰρ αὐτὴ πρῶτον, ἵνα μἄλλην λέγω,\apparat{476}{μὴ ἄλλην R : \lecon{μἄλλην} Dindorf (1839)}\\
ξύνοιδ᾽ ἐμαυτῇ πολλά γ᾽· ἀλλ᾽ ἐκεῖνο δ᾽ οὖν\apparat{477}{πολλ᾽ R : δεινὰ πολλ᾽ Bentley : \lecon{πολλά γ᾽· ἀλλ᾽} Bibi \ligneApp{478} ἦν R : \lecon{ἦ} Dindorf (1825)}%
\commentaire{477}{ξύνοιδ᾽ ἐμαυτῇ πολλ᾽ ἐκεῖνο δ᾽ οὖν.}{Il manque un pied au vers 477 ; Bentley y suppléa en ajoutant <δεινὰ> πολλ᾽ pour introduire le superlatif δεινότατον du vers suivant. C'est une hypothèse (Dawes préféra πολλὰ δείν᾽). J'en propose une autre qui a l'avantage d'expliquer paléographiquement la lacune par un \italique{saut du même au }quasi-\italique{même} (il n'y a qu'une patte de mouche de différence entre un π et un γ) : \textbf{ΠΟΛΛ'} à la place de ΠΟΛΛΑ\textbf{Γ'ΑΛΛ'}.}\\
δεινότατον, ὅτε νύμφη μὲν ἦ τρεῖς ἡμέρας,\\\passeLigne}
\end{groupe}%
%% v. 466-519
\begin{groupe}%
\ledpb%
\trimIamb{ὁ δ᾽ ἀνὴρ παρ᾽ ἐμοὶ καθηῦδεν· ἦν δέ μοι φίλος,\apparat{479}{καθεύδειν R : \lecon{καθηῦδεν} Scaliger : καθεῦδεν Bentley, Brunck \ligneApp{480} διεκόρευσεν R, Pollux III.42 codd.FS : \lecon{διεκόρησεν} Pollux codd.A, Fritzsche (cf. \italique{Pax} 59) \espApp ἑπτέτιν οὖσαν R : \lecon{οὖσαν ἑπτέτιν} Pollux \ligneApp{481} μ᾽ οὔκνυεν R : \lecon{μου ᾽κνυεν} Suid. ε 520 \ligneApp{482} κᾶτ᾽ R : \lecon{κᾆτ᾽} Küster : κἀγὼ Lenting}\\
ὅσπερ με διεκόρησεν οὖσαν ἑπτέτιν.\testimonia{480}{Pollux III 42 (οὗτος pro ὅσπερ et διεκόρησεν pro διεκόρευσεν) \espTest (διεκόρησεν) Hesychius δ 1558 διεκόρησε}%
\commentaire{480}{με διεκόρευσεν ἑπτέτιν οὖσαν.}{Mnésiloque n'aura pas mis longtemps à perdre toute retenue : il s'invente bien imprudemment une histoire repoussante qui n'a d'égale que celle de la Quartilla de Pétrone (\italique{Satyrico}, XXV) ou de la marquise de Brinvilliers...}\\
οὖτος πόθῳ μου ᾽κνυεν ἐλθὼν τὴν θύραν·\testimonia{481}{Suidas ε 520 ; Zonaras 669 ἔκνυεν; \Eustathe 1746.6 (μοὔκνυεν ἕλκων pro μου ᾽κνυεν ἐλθὼν)}\\
κᾆτ᾽ εὐθὺς ἔγνων· εἶτα καταβαίνω λάθρᾳ.\\
ὁ δ᾽ ἀνὴρ ἐρωτᾷ `ποῖ σὺ καταβαίνεις;' `ὅποι;\\
στρόφος μ᾽ ἔχει τὴν γαστέρ᾽ ὦνερ κὠδύνη·\testimonia{484}{\Eustathe 1746.45 \espTest (στρόφος) \Eustathe 503.20}%
\commentaire{484}{στρόφος μ᾽ ἔχει τὴν γαστέρ᾽ ὦνερ κὠδύνη·}{\grec{στρόφος} est le terme «~propre~» pour désigner la \italique{colique} (cf. Hippocrate, \italique{Des maladies aiguës}, XXXVII ; \italique{Aphorismes}, V, 41).}\\
ἐς τὸν κοπρῶν᾽ οὖν ἔρχομαι.' `βάδιζέ νυν.'\\
κᾆθ᾽ ὁ μὲν ἔτριβεν κεδρίδας ἄννηθον σφάκον·\testimonia{486}{Suidas σ 1712 \espTest (κεδρίδας) Hesychius κ 1980 κεδρίδες ; Photius κ 1980 κ 517}\apparat{486}{\lecon{ἔτριβεν} R : ἔτριβε Suid. σ 1712 \espApp ἄνηθον R, Suid. : \lecon{ἄννηθον} Dindorf (1830)}%
\commentaire{486}{κᾆθ᾽ ὁ μὲν ἔτριβεν κεδρίδας ἄννηθον σφάκον·}{Le niais mari se jette sur les graines de genévrier, d'aneth et de sauge car elles sont réputées préserver du mal dont sa femme dit souffrir. À toutes fins utiles pour le lecteur – qui pourrait bien trouver un réel secours grâce à ces connaissances qui ne sont pas toutes théoriques –, nous reproduisons ce que Pline l'Ancien nous apprend des vertus de ces plantes :
\citComm{Pour le genévrier (ἡ κεδρίς)\marqueur{1} :\\
\latin{Juniperus vel ante caetera omnia excalfacit, extenuat, cedro assimilis. Et ejus duo genera : altera major, altera minor. Utraque accensa serpentes fugat. Semen \textbf{stomachi}, pectoris, lateris \textbf{doloribus utile}.}\\
\francais{Le genévrier est plus que toute autre plante échauffant et atténuant, semblable en cela au cèdre. Il en est de deux espèces : l'un, plus grand que l'autre. Tous deux, brûlés, font fuir les serpents. La graine est bonne dans les \italique{douleurs d'estomac}, de poitrine et de côté.}\\
Pour l'aneth (τό ἄννηθον)\marqueur{2} :\\
\latin{Anethum quoque ructus movet , et tormina sedat : alvum sistit : epiphoris radices illinuntur ex aqua vel vino.}\\
\francais{L'aneth (anethum graveolens, L.) aussi est carminatif, et calme les tranchées ; Il arrête les flux de ventre. On applique la racine dans de l'eau ou dans du vin sur les épiphoras.}\\
Et enfin la sauge (ὁ σφάκος)\marqueur{3} :\\
\latin{Est et silvestris elelisphacos} \grec{[le scholiaste du \italique{Ravennas} glose : σφάκον· ἐλελίσφακον. ὡς ἐπιτήδεια παῦσαι στρόφον]} \latin{dicta a Graecis, ab aliis \textbf{sphacos}. Ea est sativa lente levior, et folio minore, atque sicciore, et odoratiore. [...] \textbf{Bibitur cum absinthio ad dysenteriam}. Cum vino eadem commorantes menses trahit : abundantes sistit decocto ejus poto. Per se imposita herba vulnerum sanguinem cohibet.}\\
\francais{La plante appelée par les Grecs élélisphacos ou \italique{sphacos} (salvia pomifera, L.) est une espèce de lentille sauvage plus légère que la lentille cultivée, à feuilles plus petites, plus sèches et plus odorantes. [...] On la fait boire avec l'absinthe pour la dysenterie. Avec le vin, elle fait venir les règles en retard; en décoction, elle arrête les règles trop abondantes ; appliquée seule, elle étanche le sang des plaies.}}}%
\debutNote{486}%
\note{1}{Pline l'Ancien, \italique{Histoire naturelle}, XXIV, 36. Je reprends l'édition et la traduction d'Émile Littré, \italique{Histoire naturelle de Pline}, tome second (1877).}%
\note{2}{Pline l'Ancien, \italique{Histoire naturelle}, XX, 74.}%
\note{3}{Pline l'Ancien, \italique{Histoire naturelle}, XXII, 71.}\\
ἐγὼ δὲ καταχέασα τοῦ στροφέως ὕδωρ
\commentaire{487}{ἐγὼ δὲ καταχέασα τοῦ στροφέως ὕδωρ.}{La perfide est rouée (στρόφις) : elle asperge d'eau les gonds de la porte pour ne pas qu'ils fassent de bruit. Le scholiaste l'avait compris ainsi : ἵνα ὀλισθηρὸς γενόμενος μὴ ψοφῇ. On a le même tour chez Plaute (\italique{Curculio}, I, 3, 1) :%
\citComm{\latin{Anus — Placide egredere, et sonitum prohibe forium, et \textbf{crepitum cardinum},\\
ne quod hic agimus, herus percipiat fieri, mea Planesium.\\
mane, \textbf{subfundam aquolam}.}%
}%
Je me permets d'attirer l'attention du lecteur sur le jeu de mots entre καταχέω / καταχέζω d'une part et στροφεύς / στρόφος d'autre part. Je n'en dis pas davantage car en voilà-z-assez !}\\
ἐξῆλθον ὡς τὸν μοιχόν· εἶτ᾽ ἠρειδόμην\apparat{488}{ἐρειδόμην R : \lecon{ἠρειδόμην} Bentley, Küster}\\
παρὰ τὸν Ἀγυιᾶ κύβδ᾽ ἐχομένη τῆς δάφνης.\testimonia{489}{(Ἀγυιᾶ) Stephanus Byzantinus 22.15 ; cf. Moeris α 13}\apparat{489}{κυνδ᾽ \Rac : \lecon{κύβδ᾽} \Rpc, \Mon}%
\commentaire{489}{παρὰ τὸν Ἀγυιᾶ κύβδ᾽ ἐχομένη τῆς δάφνης.}{L'Agyieus, surnom tiré d'ἀγυιά, la \italique{rue}, donné à Apollon, considéré comme protecteur des rues, devint par extension le nom des images et autels élevés en l'honneur du dieu, à Athènes, devant les portes des maisons (ἐν προθύροις). (Note d'Émile Saglio\marqueur{1}.) Mnésiloque parle bien d'un autel, ceint selon l'usage de laurier ; Hésychius (α 856) :
\citComm{ἀγυιεύς· ὁ πρὸ τῶν θυρῶν ἑστὼς βωμὸς ἐν σχήματι κίονος\\
\francais{Agyieus : l'autel placé devant les portes, en forme de pilier.}}%
Sachant ce qu'il y fait, Mnésiloque ajoute donc l'insulte au dieu à l'outrage de l'adultère ! Au vu de la forme de l'ἀγυιεύς, on peut imaginer comment il prétend s'être tenu `παρὰ τὸν Ἀγυιᾶ κύβδα', mais pourquoi `ἐχομένη τῆς δάφνης', \italique{s'agrippant au laurier} ?! Quoique porté aux \italique{Realia} par ma formation à l'école de Willems, je me suis dispensé de reconstituer la scène ; cela ne m'empêchera pas d'affirmer : voilà une posture fort incommode !
Pourtant, cette image érotique d'un \italique{Penseur} de Rodin (κύβδα) se faisant levrer alors qu'il se tient à un arbre se retrouve ailleurs\marqueur{2} ; il y a donc un lièvre à lever !
Une fois encore, ce sont les parémiographes qui nous furent d'une aide précieuse : ils ont recensé une des \italique{formules d'Éphèse} (Ἐφέσια γράμματα), incantations magiques pour conjurer le mauvais sort\marqueur{3} :%
\citComm{Ἐφέσια γράμματα: ἐπῳδαί τινες ἦσαν, ἅσπερ οἱ φωνοῦντες ἐνίκων ἐν παντί. ταὐτὸ τῇ [lire : ταὐτότητι ?]·\\ 
Δαφνίνην φορῶ βακτηρίαν.\\
Lettres éphésiennes : paroles magiques que ceux qui les prononçaient obtenaient la victoire en tout. En répétant :\\
Je porte la branche de laurier.}%
Ainsi, les Grecs \italique{touchaient du bois} (de laurier) pour éviter les malheurs ; cela explique l'attitude de notre libidineuse : s'accrochant aux branches et à de superstitieuses croyances, pensant ainsi échapper au mauvais œil d'un flagrant délit d'adultère.}%
\debutNote{489} %
\note{1}{Début de l'article \latin{Agyieus} de l'indispensable \italique{Dictionnaire des antiquités grecques et romaines} de Daremberg \& Saglio ; l'on peut y voir des représentations d'Agyeus.}%
\note{2}{Lire Théocrite, \italique{Idylles}, V, 115-119 et voir, par exemple, – ce qui n'est d'ailleurs pas indispensable ! – \italique{Antike Denkmäler}, Band II, 1908, Tafel 41, Figure 3.}%
\note{3}{\italique{Corpus paroemiographorum graecorum}, I, ed. E. L. a Leutsch, Diogeniani centuria IV 78, p. 244. On pourra consulter la \italique{Souda} ou Photius à l'entrée Ἀλεξιφάρμακα : ces charmes pouvaient concerner des époux. Noter que ce fameux bâton de laurier que pouvaient prendre les convives lors de symposies avant d'entonner le chant était appelé αἴσακος, qui est aussi le nom d'un centaure.}\\
ταῦτ᾽ οὐδεπώποτ᾽ εἶφ᾽, ὁρᾶτ᾽, Εὐριπίδης·\\
οὐδ᾽ ὡς ὑπὸ τῶν δούλων τε κὠρεωκόμων\\
\lednopb%
σποδούμεθ᾽, ἢν μὴ 'χωμεν ἕτερον, οὐ λέγει·\\\passeLigne}
\end{groupe}%
%% v. 466-519
\begin{groupe}%
\ledpb%
\trimIamb{οὐδ᾽ ὡς ὅταν μάλισθ᾽ ὑπό του ληκώμεθα\testimonia{493-4}{Suidas λ 434 Ληκώμεθα: κινώμεθα. Ἀριστοφάνης: μάλισθ' ὅταν ὑπό του ληκώμεθα τὴν νύχθ' ὅλην}\apparat{493}{\lecon{ληκώμεθα} \Rac (?), Suid. λ 434 : κινώμεθα \Rpc, \Mon \ligneApp{494} post νύχθ' add. ὅλην Suid. (reminiscentia Pherecratis ? cf. Photius λ 254) \espApp σκορόδια R : \lecon{σκόροδα δια} Bentley \ligneApp{495} ὀσφραινόμενος R : \lecon{ὀσφρόμενος} Bentley, Küster \espApp ἀνὴρ R : \lecon{ἁνὴρ} Dawes \espApp ἀπὸ τοῦ R : \lecon{ἀπὸ} Biset}%
\commentaire{493}{μάλισθ᾽ ὑπό του ληκώμεθα.}{Le scribe ici encore a hésité entre \grec{ληκώμεθα}, qu'il écrivit probablement en premier, avant de gratter les trois premières lettres et de réécrire par dessus \grec{κινώμεθα}. N'étant pas un \italique{aficionado} de la \italique{lectio difficilior}, il aura probablement préféré le verbe \grec{κινεῖν}, bien connu de lui, au fort rare \italique{ληκᾶν} (ou \grec{ληκᾶσθαι} dans la bouche d'une femme). Suidas (λ 434) ne s'y est pas trompé qui glose ληκώμεθα par κινώμεθα en citant ce passage d'Aristophane.}\\
τὴν νύχθ᾽, ἕωθεν σκόροδα διαμασώμεθα,%\apparat{494}{post νύχθ' add. ὅλην Suid. (reminiscentia Pherecratis ? cf. Photius λ 254) \espApp σκορόδια R : \lecon{σκόροδα δια} Bentley \ligneApp{495} ὀσφραινόμενος R : \lecon{ὀσφρόμενος} Bentley, Küster \espApp ἀνὴρ R : \lecon{ἁνὴρ} Dawes \espApp ἀπὸ τοῦ R : \lecon{ἀπὸ} Biset}%
\commentaire{494}{ἕωθεν σκόροδα διαμασώμεθα.}{(Pour ne pas que le mari cocu fleure le sacrifice à Vénus de sa femme.) 
À quand une thesmophorienne ecclésie pour châtier Aristophane qui enseigne aux Athéniens toutes les fourberies de leur femme ?.. à moins que le tour ne fût déjà connu :%
\citComm{Διὰ ταῦθ' ὁ πόρνος οὗτος οὐδὲ τῶν πράσων\\
ἑκάστοτ' ἐπεδείπνει μεθ' ῆμῶν· τοῦτο δ' ἦν,\\
ἵνα μή τι λυπήσειε τὸν ἐραστὴν φιλῶν.\marqueur{1}}%
et encore avec de l'oignon :%
\citComm{Καὶ ὁ Χαρμίδης εἶπεν· `Ὦ ἄνδρες, ὁ Νικήρατος κρομμύων ὄζων ἐπιθυμεῖ οἴκαδε ἐλθεῖν, ἵν᾽ ἡ γυνὴ αὐτοῦ πιστεύῃ μηδὲ διανοηθῆναι μηδένα ἂν φιλῆσαι αὐτόν'\marqueur{2}.}%
}%
\debutNote{494}%
\note{1}{Athénée, XIII, 572c.}%
\note{2}{Xénophon, \italique{Symposium}, IV, 8.}\\
ἵν᾽ ὀσφρόμενος ἁνὴρ ἀπὸ τείχους εἰσιὼν%
\commentaire{495}{ἀπὸ τοῦ τείχους εἰσιὼν.}{Thucydide (VIII, 69) :%
\citComm{ἦσαν [δ᾽] Ἀθηναῖοι πάντες αἰεὶ οἱ μὲν ἐπὶ τείχει, οἱ δ᾽ ἐν τάξει, τῶν ἐν Δεκελείᾳ πολεμίων ἕνεκα ἐφ᾽ ὅπλοις.\\
\francais{Tous les Athéniens restaient continuellement, les uns sur les murailles, les autres à leurs postes, à cause des ennemis qui occupaient Décélie.}%
}%
}\\
μηδὲν κακὸν δρᾶν ὑποτοπῆται. ταῦθ᾽, ὁρᾷς,%
%\commentaire{496}{ὁρᾷς}{\eff{Pourquoi Mnésiloque passe-t-il du vouvoiement au tutoiement ?}}
\\
οὐπώποτ᾽ εἶπεν. εἰ δὲ Φαίδραν λοιδορεῖ,\\
\ligne{498}ἡμῖν τί τοῦτ᾽ ἔστ᾽; οὐδ᾽ ἐκεῖν᾽ εἴρηκέ πω,\\
ὡς ἡ γυνὴ δεικνῦσα τἀνδρὶ τοὔγκυκλον%
\commentaire{499}{τοὔγκυκλον.}{L'ἔγκυκλον était une sorte de jupe courte, portée par les femmes à l'intérieur. Cf. la note au vers 643.}%
\apparat{v. 500}{ὑπ᾽ αὐγὰς οἷόν ἐστιν ἐγκεκαλυμμένον R : ὑπαυγάσ᾽ οἷόν κτλ. Hermann (1838) : \lecon{del} van Herwerden (\italique{Mnemosyne N. S.}, 24, 1896, p.293), Willems \ligneApp{501} μυχὸν R : \lecon{μοιχὸν} Gryaneus}%
\commentaire{500}{ὑπ᾽ αὐγὰς οἷόν ἐστιν ἐγκεκαλυμμένον.}{Le manuscrit de Ravenne comporte :
\citComm{ὡς ἡ γυνὴ δεικνῦσα τἀνδρὶ τοὔγκυκλον\\
ὑπ᾽ αὐγὰς οἷόν ἐστιν ἐγκεκαλυμμένον\\
τὸν μυχὸν ἐξέπεμψεν, οὐκ εἴρηκέ πω.}
Gardons-nous, dans un premier temps, de biffer promptement le \nieme{500} vers, à l'instar de Willems\marqueur{1}, en arguant de sa faute de mètre et de sa construction des plus mal embouchées. Si interpolation il y a -- ce que nous allons essayer de démontrer --, celle-ci ne peut être que très ancienne ; en effet, une scholie érudite notait déjà l'abrègement violent (συστολή βίαιος) que devait subir ὑπ᾽ αὐγὰς pour rentrer dans le mètre iambique. \\
J'ai longtemps voulu conserver ὑπ᾽ αὐγὰς \italique{contra metrum} (quitte à le remplacer par des formes proches comme ὑπαυγάσ᾽ ou ὑπαυγάσας) car il me semblait donner des indications sur le stratagème employé par la pécheresse.
Le `οἷόν ἐστιν', lui, me laissait perplexe !.. (Willems parlait d'une «~pure cheville~» ; le cas me semble plus grave : l'on doit faire violence au grec pour en proposer un sens acceptable.) Intéressons-nous y donc !\\
Après avoir retourné l'affaire dans tous les sens, je ne m'explique cet οἷόν ἐστιν que par la présence originelle de : ᾦον ἐστιν\marqueur{2}. Athénée (II, 57e) nous précise que ᾦον pouvait être utilisé pour ὑπερῷον, l'\italique{étage supérieur d'une maison}. Le scholiaste T de l'\italique{Iliade} (XVI, 184) emploie d'ailleurs le mot en ce sens -- était-ce le même que le scribe qui écrivit notre scholie devenue le vers 500\marqueur{3} ? Quoi qu'il en soit, l'adjectif ὑπερῷος signifie \italique{situé à l'étage supérieur}.
Si l'hypothèse est exacte, le sens se laisse alors clairement dégager : on comprend que l'amant (τὸν μοιχόν) regarde vers le bas (αὐγάζειν, au sens d'\italique{observer}, est préfixé de ὑπ- par décalque sur ὑφ-ορᾶν) la femme embobiner son cocu de mari, bien caché (ἐγκεκαλυμμένον : il est à l'étage supérieur, ᾦον ἐστιν).
Une fois cette correction minime apportée, le \nieme{500} vers se peut lire ainsi : 
\citComm{ὑπαυγάσας· [τὸν μοιχόν] ᾦον ἐστιν ἐγκεκαλυμμένον}
Ainsi compris, il devient tout-à-fait clair que ce «~vers~» est une glose à propos de τὸν μοιχόν du vers 501. Le malheur voulut qu'elle fît peu ou prou un trimètre iambique et soit alors très tôt intégrée au texte même d'Aristophane ; l'inintelligibilité du nouveau texte aidant à sa corruption, ΥΠΑΥΓΑΣΑΣ devint ΥΠΑΥΓΑΣ par haplographie et ΩΙΟΝ ΟΙΟΝ.\\
Lisons donc avec Willems : %
\citComm{ὡς ἡ γυνὴ δεικνῦσα τἀνδρὶ τοὔγκυκλον\\
τὸν μοιχὸν ἐξέπεμψεν, οὐκ εἴρηκέ πω.\\
\francais{Comment la femme qui exhibant l'encycle à son mari\\
fit échapper son amant, il n'en a pas soufflé mot.}}%
Le passage, ici comme ailleurs, n'est qu'allusif ; aussi fait-il probablement référence à une histoire bien connue des spectateurs. Victor Coulon\marqueur{4} l'a rapprochée de contes au motif semblable, couchés par écrit au \nieme{XI} siècle, et de source probablement arabe (parfois, la femme est aidée de sa mère dans son entreprise de mystification du cocu)~; les conteurs orientaux ayant eux-mêmes été familiarisés aux contes grecs. Ainsi, il nous apparaît plus que plausible qu'Aristophane rapporte ici non pas un fait-divers de la chronique scandaleuse du temps mais une historiette populaire ; ces «~révélations~» sur les Athéniennes ne seraient pas alors aussi compromettantes pour elles que ne l'étaient celles d'Euripide.\\
Voir le commentaire au vers 643 où j'y montre que l'encycle (ἔγκυκλον) est une courte jupe. Cette donnée finit d'anéantir les hypothèses loufoques proposées par de savants universitaires selon lesquelles la femme tendrait son encycle pour cacher l'amant aux yeux du mari.
%Ce vers pèche contre le mètre. Willems, dans sa \nieme{13} note reproduite en annexe argue de cela et de sa construction des plus mal embouchées pour le biffer promptement. Pourtant, la présence de «~\grec{ὑπ᾽ αὐγὰς}~» paraît indispensable pour la révélation du stratagème : la femme emmène son mari à l'extérieur, \italique{au soleil} pour lui montrer son linge bien tissé\marqueur{1} et ce afin de laisser à son amant la possibilité de filer à l'anglaise\marqueur{2}. Je reprends la première conjecture de Hermann qui a l'avantage de ne changer aucune lettre : ὑπαυγάσ(αι)᾽ οἷόν ἐστιν. Le verbe \grec{ὑπαυγάζειν} n'est que peu attesté~: transitivement, Aristote\marqueur{3} l'emploie avec le sens de \italique{faire briller}, cependant Sophocle emploie bien le verbe αὐγάζειν avec le sens de \italique{voir} (au vers 217 de son \italique{Philoctète} ; de toute façon s'il fallait justifier par des exemples de chaque emploi d'un mot par Aristophane, nos traductions resteraient bien lacunaires !..) Ici, il est clair que le sens est : \italique{montrer \grec{ὑπ᾽ αὐγὰς}}~: \italique{montrer en plein jour}. Victor Coulon\marqueur{4} a rapproché cette histoire de contes au motif semblable, couchés par écrit au \nieme{XI} siècle, et de source probablement arabe (parfois, la femme est aidée de sa mère dans son entreprise de mystification du cocu)~: les conteurs orientaux ayant eux-mêmes été familiarisés aux contes grecs. Ainsi, il nous apparaît plus que plausible qu'Aristophane rapporte ici non pas un fait-divers de la chronique scandaleuse du temps mais une historiette populaire ; aussi ces «~révélations~» sur les Athéniennes ne sont peut-être pas aussi compromettantes pour elles que ne l'étaient celles d'Euripide. Passons à ἐγκεκαλυμμένον, littéralement \italique{bien affublé} d'où le sens dérivé de \italique{caché}. Ce terme, qui se rapporte au galant (τὸν μοιχὸν), est de prime abord étonnant : on comprend que l'amant veuille se cacher pour ne pas être reconnu en sortant, mais comment alors qu'il n'a que le bref laps de temps où la femme attire son mari à l'extérieur, comment ainsi pris de court trouverait-il le temps de bien se vêtir ? On l'imaginerait plutôt sortir précipitamment dans la même tenue – celle d'Adam – que lorsqu'il fut interrompu par l'arrivée inopportune du cocu. À la réflexion, c'est évident : le manteau (τοὔγκυκλον) prétendûment nouvellement tissé ne peut être, pour tromper le mari, que celui de l'amant ; de plus, il ne peut dans sa fuite dérober à la volée des habits sous peine d'éveiller les soupçons du cocu. Aussi, je mettrais ma main au feu que la version populaire – c'est-à-dire connue de tous les spectateurs – du conte faisait sortir l'amant précipitamment dans son plus simple appareil. Alors pourquoi cet ἐγκεκαλυμμένον ? Comme toujours : Hésychius à la rescousse ! Il glose comme exprès pour nous (ε 186) :
%\citComm{ἐγκαλύπτεσθαι· διατρέπεσθαι, αἰσχύνεσθαι.}%
%C'est bien cela : le cocufieur sort \italique{troublé}, \italique{honteux} comme un renard qu'une poule aurait pris... et nu comme un ver ! Aristophane use plaisamment à dessein de l'adjectif ἐγκεκαλυμμένον pour décrire cette situation ; il me semble que l'équivoque est plutôt bien rendue par l'expression : « partir la queue entre les jambes ».
}%
\debutNote{500}%
%\note{1}{De tout temps, la femme a pris plaisir à admirer de belles robes ; \cf Eur. \italique{Héc.} 1154 où les captives troyennes louent le tissage des femmes thraces \grec{ὑπ᾽ αὐγὰς τούσδε λεύσσουσαι πέπλους}. (Note de Victor Coulon, référence ci-dessous.)}%
%\note{2}{Et non pas l'invraisemblable interprétation de ces vers qu'on peut lire çà et là : la femme déploierait son \italique{encycle} pour camoufler au mari la fuite de l'amant caché derrière (sic) avant que ce dernier ne trouve encore le moyen de s'enfuir (re-sic). Il m'a été donné de lire un papier traitant de l'identification de ce vêtement qui donnait du crédit à cette interprétation. Et voilà comment une exégèse imbécile peut conduire à voir en l'ἔγκυκλον non une sorte de jupe, ce qu'il était (comme nous l'allons montrer à la note du vers 643), mais un \italien{mantello ampio, pesante}...}%
%\note{3}{Aristote, \italique{Ars Rhetorica}, 3, 1378}%
\note{1}{Le lecteur pourra néanmoins utilement se reporter à \nieme{13} note ainsi que son introduction aux \italique{Notes sur la} Paix (t. II, pp. 93-99), toutes deux reproduites en annexe.}%
\note{2}{Rappelons que les iotas souscrits n'apparaissent qu'au cours du \nieme{XIII} siècle ; le manuscrit de Ravenne, du \nieme{IX} ou \nieme{X} siècle est adscrit tout comme les manuscrits qu'il transcrivait. La faute -- si faute il y a -- ne consiste par conséquent qu'en la confusion entre un Ω et un Ο.}%
\note{3}{Ce même scholiaste T connaissait bien notre comédie : il en cite des passages comme les vers 536-537 (voir les \latin{testimonia} \italique{ad locum}.).}%
\note{4}{«~Interprétation d'un passage d'Aristophane imité par les conteurs du Moyen-Âge~», in \italique{Annuaire de l'institut de philologie et d'histoire orientales}, tome II(1934) : \italique{Mélanges Bidez}, premier volume, p. 121-34.}\ligne{500}\\
τὸν μοιχὸν ἐξέπεμψεν, οὐκ εἴρηκέ πω.\\
\ligne{502}ἑτέραν δ᾽ ἐγᾦδ᾽ ἣ 'φασκεν ὠδίνειν γυνὴ\\%\apparat{502}{ήφασκεν R : \lecon{ἣ 'φασκεν} Zanetti}\\
δέχ᾽ ἡμέρας, ἔως ἐπρίατο παιδίον·\\\passeLigne}
\end{groupe}%
%% v. 520-530
\begin{groupe}%
\ledpb%
\trimIamb{ὁ δ᾽ ἀνὴρ περιήρχετ᾽ ὠκυτόκι᾽ ὠνούμενος·\testimonia{504}{(ὠκυτόκι᾽) Pollux II 7 ; \Phryn 129.3}\apparat{504}{\lecon{περιήρχετ᾽} R : περιέτρεχ᾽ Bergk \espApp ὠκυτόκειa R : \lecon{ὠκυτόκι᾽} Pollux II.7 \ligneApp{505} \lecon{τὸ δ᾽} R : τότ᾽ Biset : τῇδ᾽ Daubuz \ligneApp{506} βοῴη suppl. \Rpc}\\
τὸ δ᾽ εἰσέφερε γραῦς ἐν χύτρᾳ τὸ παιδίον,\\
ἵνα μὴ βοῴη, κηρίῳ βεβυσμένον·\testimonia{506}{(κηρίῳ βεβυσμένον) Hesychius κ 2549 κηρίῳ βύσασα ; Photius κ 675 κηρίῳ βῦσαι}\\
εἶθ᾽ ὡς ἔνευσεν ἡ φέρουσ᾽, εὐθὺς βοᾷ,\\
`ἄπελθ᾽ ἄπελθ᾽, ἤδη γὰρ ὦνέρ μοι δοκῶ\\
τέξειν.' τὸ γὰρ ἦτρον τῆς χύτρας ἐλάκτισεν·\testimonia{509}{(τὸ γὰρ ἦτρον -- ἐλάκτισεν) Suidas η 637 ; Zonaras 1012 ἦτρον \espTest (ἦτρον) Moeris η 14}%
\commentaire{509}{τέξειν.}{Cf. sur le mot τέξειν la note 14 de Willems à lire dans les annexes.}
\commentaire{509}{τὸ γὰρ ἦτρον τῆς χύτρας ἐλάκτισεν·}{Il y a une plaisanterie \italique{contre l'attente} (παρὰ προσδοκίαν) : τῆς χύτρας est mis pour τῆς μητρός.}\\
χὠ μὲν γεγηθὼς ἀπέτρεχ᾽, ἡ δ᾽ ἐξέσπασεν\apparat{510}{ἔτρεχεν R : \lecon{ἀπέτρεχ᾽} van Herwerden}\\
ἐκ τοῦ στόματος τοῦ παιδίου, τὸ δ᾽ ἀνέκραγεν.\apparat{511}{\lecon{τοῦ παιδίου} R : τὸ κηρίον Hirschig}\\
εἶθ᾽ ἡ μιαρὰ γραῦς, ἣ ᾽φερεν τὸ παιδίον,\apparat{512}{εἰσέφερε \Rac : ἥ ἔφερε \Rpc (cum η supra ει scripto et σ eraso) : \lecon{ἥ ᾽φερεν} Bentley}\\
θεῖ μειδιῶσα πρὸς τὸν ἄνδρα καὶ λέγει,\\
`λέων λέων σοι γέγονεν, αὐτέκμαγμα σόν,\testimonia{514-6}{Suid. λ 271 $\simeq$ Suid. π 2114}\testimonia{514}{Pollux IX 131 (λέων\expo{2} omisso et αὐτὸ ἔκμαγμα pro αὐτέκμαγμα) ; \Eustathe 1857.10, 462.40 ; cf. \AelDion ε 26}\apparat{514}{αὖτ᾽ ἔκμαγμα (ex έγμαγμα) \Rpc, Pollux IX.131 : ἔκμαγμα Suid. λ 271, π 2114 : \lecon{αὐτέκμαγμα} Reiske \ligneApp{516} \lecon{κύτταρον} R, Suid. κ 2787 : κύτταρος Fritzsche}\\
τά τ᾽ ἄλλ᾽ ἁπαξάπαντα καὶ τὸ πόσθιον\testimonia{515-6}{(πόσθιον -- κύτταρον) Suidas κ 2787}%
\commentaire{515}{ἁπαξάπαντα.}{La Fontaine, \italique{La chatte métamorphosée en femme} (II, 18) :%
\citComm{\francais{La croit femme en tout et partout.}}%
}\\
τῷ σῷ προσόμοιον, στρεβλὸν ὥσπερ κύτταρον.'\testimonia{516}{(στρεβλὸν) Hesychius σ 1977 στρεβλὸν \espTest{(κύτταρον)} Hesychius κ 4747 ; cf. Photius κ 1270}%\apparat{516}{\lecon{κύτταρον} R, Suid. κ 2787 : κύτταρος Fritzsche}%
\commentaire{516}{κύτταρον,}{synonyme de κῶνος, \italique{pomme de pin}, que les Grecs appelaient aussi στρόβιλος (στρόμβος dans Nicandre. Hehn., p.292). (Note de Willems.)}\\
ταῦτ᾽ οὐ ποιοῦμεν τὰ κακά; νὴ τὴν Ἄρτεμιν\\
ἡμεῖς γε. κᾆτ᾽ Εὐριπίδῃ θυμούμεθα,\\
οὐδὲν παθοῦσαι μεῖζον ἢ δεδράκαμεν;
\commentaire{518-9}{θυμούμεθα, / οὐδὲν παθοῦσαι μεῖζον ἢ δεδράκαμεν;}{On sait grâce au scholiaste que ces vers sont tirés du \italique{Télèphe} : 
\citComm{\grec{παρὰ τὰ ἐκ Τηλέφου Εὐριπίδου· `εἰ δὴ θυμούμεθα παθόντες οὐδὲν μᾶλλον ἢ δεδρακότες'\marqueur{1}.}}}%
\debutNote{518-9}%
\note{1}{fragment 20 du \italique{Télèphe} de l'édition Jouan-van Looy aux Belles Lettres, tome VIII, \nieme{3} partie. Correspond aux fragments 711 de l'édition Kannicht et 965 de celle de Mette.}\\\passeLigne}%
\end{groupe}%
%% v. 520-530
\begin{groupe}%
\ledpb%
\personnageGrec{Χορός\strophe{ἀντιστρ.}}%
\typeVers{par.}\chantChoeur{τουτὶ μέντοι θαυμαστόν,\commentaire{520-30}{}{Antistrophe du chant du chœur des vers 433-442.}}%
\end{groupe}%
%% v. 520-530
\begin{groupe}%
\typeVers{$2$ troch.}\chantChoeur{ὁπόθεν ηὑρέθη τὸ χρῆμα,\apparat{521}{εὑρέθη R : \lecon{ηὑρέθη} Meineke}}%
\end{groupe}%
%% v. 520-530
\begin{groupe}%
\typeVers{$2$ troch.}\chantChoeur{χἤτις ἐξέθρεψε χώρα}%
\end{groupe}%
%% v. 520-530
\begin{groupe}%
\typeVers{$2$ troch.}\chantChoeur{τήνδε τὴν θρασεῖαν οὕτω.}%
\end{groupe}%
%% v. 520-530
\begin{groupe}%
\typeVers{$2$ troch.}\chantChoeur{τάδε γὰρ εἰπεῖν τὴν πανοῦργον}%
\end{groupe}%
%% v. 520-530
\begin{groupe}%
\typeVers{$2$ troch.}\chantChoeur{κατὰ τὸ φανερὸν ὧδ᾽ ἀναιδῶς}%
\end{groupe}%
%% v. 526a
\begin{groupe}%
\decoupeVers{a}\typeVers{$2$ troch.}\chantChoeur{οὐκ ἂν ᾠόμην ἐν ἡμῖν\apparat{526}{ὑμῖν R : \lecon{ἡμῖν} Bonini}}%
\end{groupe}%
%% v. 526b
\begin{groupe}%
\ligne{526}\decoupeVers{b}\typeVers{$2$ troch\cat}\chantChoeur{οὐδὲ τολμῆσαί ποτ᾽ ἄν.}%
\end{groupe}%
%% v. 527
\begin{groupe}%
\typeVers{$2$ troch.}\chantChoeur{ἀλλὰ πᾶν γένοιτ᾽ ἂν ἤδη·\testimonia{527-30}{Suidas α 1073}\apparat{527}{ἀλλ᾽ ἅπαν R : \lecon{ἀλλὰ πᾶν} Cobet}%
\commentaire{527}{ἀλλ᾽ ἅπαν γένοιτ᾽ ἂν ἤδη·}{Le lieu-commun du \latin{omnia mutantur}\marqueur{1} avait bien-sûr cours aussi chez les Grecs. 
Brève recension : \\
Hérodote (V, 9) :%
\citComm{γένοιτο δ̓ ἂν πᾶν ἐν τῷ μακρῷ χρόνῳ.}%
Sophocle (\italique{Philoctète}, vers 305-306) : %
\citComm{πολλὰ γὰρ τάδε / ἐν τῷ μακρῷ γένοιτ᾽ ἂν ἀνθρώπων χρόνῳ.}%
Eupolis\marqueur{2} :%
\citComm{ἦ πολλά γ' ἐν μακρῷ χρόνῳ γίγνεται μεταλλαγῇ / (τῶν) πραγμάτων· μένει δὲ χρῆμ' οὐδὲν ἐν ταὐτῷ ῥυθμῷ.}%
Ménandre (fragment 685(466)\marqueur{3}) : 
\citComm{οὐκ ἔστ' ἄπιστον οὐδὲν ἐν θνητῷ βίῳ / οὐδ' ἂν γένοιτο.}
Et dans la pièce d'Aristophane d'antan, \italique{Lysistrata} (vers 258) : 
\citComm{Χορὸς γερόντων --- ἦ πόλλ᾽ ἄελπτ᾽ ἔνεστιν ἐν τῷ μακρῷ βίῳ φεῦ.}
Remarquons que dans tous ces extraits, le grec insiste sur la longueur de temps que prennent les changements pour s'opérer (ἐν τῷ μακρῷ χρόνῳ / ἐν θνητῷ βίῳ / ἐν τῷ μακρῷ βίῳ)... or c'est tout le contraire ici : ἤδη, \italique{aujourd'hui}, \italique{actuellement}. Ce contraste est certainement une allusion d'Aristophane aux temps mouvementés politiquement de la représentation des \italique{Thesmophories} -- où tout semble s'accélérer.}%
\debutNote{527}%
\note{1}{La citation est d'Ovide (\italique{Métamorphoses}, XV, 165). Érasme (\italique{Adages}, 663) recense la formule équivalente de Térence : \latin{Omnium rerum vicissitudo est}.}%
\note{2}{Fragment 391(356) des \italique{Poetae Comici Graeci}, ed. Kassel \& Austin, V (1986).}%
\note{3}{\italique{Poetae Comici Graeci}, VI 2 (1998). (Le fragment est extrait de Stobée I, 8, 31.)}}%
\end{groupe}%
%% v. 520-5308
\begin{groupe}%
\typeVers{$2$ troch.}\chantChoeur{τὴν παροιμίαν δ᾽ ἐπαινῶ\testimonia{528-30}{\Stobee III 13.29}}%
\end{groupe}%
%% v. 529
\begin{groupe}%
\typeVers{$2$ troch.}\chantChoeur{τὴν παλαιάν· ὑπὸ λίθῳ γὰρ}%
\end{groupe}%
%% v. 530a
\begin{groupe}%
\decoupeVers{a}\typeVers{troch.}\chantChoeur{παντί που χρὴ}%
\end{groupe}%
%% v. 530b
\begin{groupe}%
\ligne{530}\decoupeVers{b}\typeVers{$2$ troch\cat}\chantChoeur{μὴ δάκῃ ῥητωρ ἀθρεῖν.
\commentaire{528-30}{τὴν παροιμίαν δ᾽ ἐπαινῶ / τὴν παλαιάν· ὑπὸ λίθῳ γὰρ / παντί που χρὴ / μὴ δάκῃ ῥητωρ ἀθρεῖν.}{Le proverbe – avec un scorpion à la place de l'orateur – était effectivement populaire puisqu'il nous est diversement parvenu dans les recueils de Proverbes. Ainsi des parémiographes Zénobie\marqueur{1} ou Diogénianus\marqueur{2} et d'Hésychius\marqueur{3} qui rapportent la version la plus courante : 
\citComm{ὑπὸ παντὶ λίθῳ σκόρπιος εὕδει: ἐπὶ τῶν κακοήθων λέγεται.}%
Macaire\marqueur{4} en précise le sens : \grec{ἐπὶ πολυπληθίᾳ τῶν πανταχοῦ ὄντων πονηρῶν ἀνθρώπων}, « rapport au grand nombre d'hommes qui sont méchants sous tout rapport ».
Le proverbe est attesté aussi par Praxilla\marqueur{5} et Sophocle\marqueur{6}.\\
Ici, la \italique{pierre} (λίθος) dont parle la coryphée est bien-sûr le piédestal (βῆμα) de la Pnyx sur lequel les orateurs -- Mnésiloque, en l'occurrence -- parlaient. (Aristophane use du mot dans le même sens aux vers 683 des \italique{Acharniens} ; 680 de la \italique{Paix} et 87 de l'\italique{Assemblée des femmes}.)
Hésychius (λ 999) :%
\citComm{λίθος· βῶλος. βωμός. καὶ βάσις. τὸ ἐν τῇ Ἀθηναίων ἐκκλησίᾳ βῆμα.}%
Enfin, la substitution au scorpion de l'orateur n'est pas gratuite : on parlait couramment de la \italique{morsure de la langue}\marqueur{7}.
Ainsi de la \italique{morsure du sycophante}, δῆγμα συκοφάντου, qui paraît proverbiale\marqueur{8}. Diogène le cynique eut cette réplique à l'ironie \italique{mordante} (rapportée par son homonyme laërtien) :
\citComm{ἐρωτηθεὶς τί τῶν θηρίων κάκιστα δάκνει ἔφη: τῶν μὲν ἀγρίων συκοφάντης, τῶν δὲ ἡμέρων κόλαξ.\\
Interrogé sur la bête dont la morsure est la plus mauvaise, [il répondit] que des bêtes sauvages c'était le sycophante, et des animaux domestiques, c'était le flatteur.}}%
\debutNote{528-30}%
\note{1}{\latin{Zenobii centuria} VI, 20 (tomus I). Je cite les proverbes d'après le recensement d'E. Leutsch et F. G. Schneidewin, \italique{Corpus paroemiographorum Graecorum}, en deux tomes. Signalons qu'une nouvelle édition critique des recensions de Zénobius et Diogénianus avec traduction en italien a paru en 2019 sous le titre \italique{I proverbi greci}, dir. Emanuelle Lelli.}%
\note{2}{\latin{Diogeniani centuria} VIII, 59 (tomus I).}%
\note{3}{Hésychius (υ 717). Photius et Suidas rapportent aussi ce proverbe : il est partout !}%
\note{4}{\latin{Macarii centuria} VIII, 69 (tomus II).}%
\note{5}{Le scholiaste attribue même le proverbe à Praxilla. Athénée (XV, 695d) donne un exemple de scolie (σκόλιον), sorte de chanson à boire, écrite sur ce proverbe ; comme Praxilla excella dans ce genre, il n'est pas impossible que la scolie fût d'elle, et que le scholiaste l'ait en tête lorsqu'il attribue un peu hâtivement la maternité du proverbe à Praxilla. Cette scolie correspond au fragment 903 des \italique{Poetae Melici Graeci}.}%
\note{6}{Fragment 37 de l'édition Radt.}%
\note{7}{Voir les sens figurés du verbe δάκνειν.}%
\note{8}{Voir le \italique{Ploutos}, vers 885, et le commentaire du scholiaste \italique{ad. loc.}}\\\passeLigne}%
\end{groupe}%
%% v. 531-532
\begin{groupe}%
\personnageGrec{Κορυφαῖος}%
\typeVers{tetr. iamb\cat}\tetramIambCat{\ligne{531}ἀλλ᾽ οὐ γάρ ἐστι τῶν ἀναισχύτων φύσει γυναικῶν\testimonia{531-2}{Suidas α 1349 (sed ἄλλαι pro ἄρ᾽ εἰ)}\apparat{531}{(Κορ.) add. von Velsen}%
\commentaire{531-73}{}{Tout ce passage est composé en tétramètres iambiques catalectiques ; mètre du conflit et des débats animés. Cf. aussi les \italique{Cavaliers} v. 335-366 ; \italique{Lysistrata} v. 350-381, 467-475.}\\
\lednopb%
οὐδὲν κάκιον εἰς ἅπαντα πλὴν ἄρ᾽ εἰ γυναῖκες.\apparat{532}{\lecon{ἄρ᾽ εἰ} R : ἄλλαι Suid. α 1349}%
\commentaire{531-2}{ἀλλ᾽ οὐ γάρ ἐστι τῶν ἀναισχύτων φύσει γυναικῶν / οὐδὲν κάκιον εἰς ἅπαντα πλὴν ἄρ᾽ εἰ γυναῖκες.}{Oct. Feuillet, \italique{Dalila}, partie IV :
\citComm{\francais{Rien de plus semblable sous le soleil qu'une femme – et une femme.}}%
C'est la traduction presque textuelle de Ménandre (\italique{Sentences monostiques}, Γνῶμαι μονόστιχοι 109) :%
\citComm{Γυνὴ γυναικὸς πώποτ᾽ οὐδὲν διαφέρει.}%
(Note de Willems ; pour voir ce que Willems pense de la « manie des conjectures » visant à retoucher le texte d'Aristophane, on lira avec un plaisir non dissimulé la note entière n° 15, reproduite à la fin.)\\
Il serait piquant qu'Aristophane mette dans la bouche de la coryphée une citation d'Euripide lui-même pour dénigrer le comportement de certaines femmes ; le scholiaste n'en a malheureusement pas soufflé mot... on dispose néanmoins de citations d'Euripide très approchantes (fr. 494.27 Kannicht $=$ fr. 14.27 \italique{Mélanippe} Jouan-van Looy) : %
\citComm{τῆς μὲν κακῆς κάκι[ο]ν οὐδὲν γίγνεται / γυναικός.}%
}%
\\\passeLigne}%
\end{groupe}%
%% v. 533-539
\begin{groupe}%
\ledpb%
\personnageGrec{Μίκα}%
\tetramIambCat{οὔτοι μὰ τὴν Ἄγλαυρον ὦ γυναῖκες εὖ φρονεῖτε,\apparat{533}{(γυνὴ τις) \Rpc : \lecon{(Μίκα)} Rogers \espApp Ἄγραυλον R : \lecon{Ἄγλαυρον} Daubuz}%
\commentaire{533}{μὰ τὴν Ἄγραυλον.}{Scholie d'un érudit :%
\citComm{κατὰ τῆς ᾿Αγραύλου ὦμνυον, κατὰ δὲ τῆς Πανδρόσου σπανιώτερον· κατὰ δὲ τῆς Ἕρσης οὐχ εὑρήχαμεν.\\
\francais{L'on jurait sur Agraule, plus rarement sur Pandrosos ; nous n'avons pas trouvé d'exemple sur Ersès.}} %
Aglaure (correction d'Agraule confirmée par l'épigraphie), Pandrose et Hersé étaient les trois filles de Kécrops, mythique fondateur et premier roi d'Athènes.}\\
ἀλλ᾽ ἢ πεφάρμαχθ᾽ ἢ κακόν τι μέγα πεπόνθατ᾽ ἄλλο,\\
ταύτην ἐῶσαι τὴν φθόρον τοιαῦτα περιυβρίζειν\\
ἡμᾶς ἁπάσας. εἰ μὲν οὖν τις ἔστιν· εἰ δὲ μή, ἡμεῖς\testimonia{536-7}{(εἰ μὲν -- λαβοῦσαι) scholium\expo{bT} Homerus \italique{Il.} I 135-7 (ἔστι τις pro τις ἔστιν et αὐταί omisso)}\apparat{536}{\lecon{μή ἡμεῖς} R : μἠμεῖς Bentley}\\
αὐταί γε καὶ τὰ δουλάρια τέφραν ποθὲν λαβοῦσαι\apparat{537}{αὐτοί R : \lecon{αὐταί} Bonini \espApp \lecon{γε} R : τε schol. T ad Hom.\italique{Il}.I.135, Reiske \espApp πόθεν R : \lecon{ποθὲν} Küster}\\
ταύτης ἀποψιλώσομεν τὸν χοῖρον, ἵνα διδαχθῇ%
%\commentaire{537-8}{τέφραν ποθὲν λαβοῦσαι / ταύτης ἀποψιλώσομεν τὸν χοῖρον}{.}
\\
γυνὴ γυναῖκας οὖσα μὴ κακῶς λέγειν τὸ λοιπόν.\\\passeLigne}%
\end{groupe}%
%% v. 540-543
\begin{groupe}%
\personnageGrec{Μνησίλοχος}%
\tetramIambCat{μὴ δῆτα τόν γε χοῖρον ὦ γυναῖκες. εἰ γὰρ οὔσης\\
παρρησίας κἀξὸν λέγειν ὅσαι πάρεσμεν ἀσταί,\\
εἶτ᾽ εἶπον ἅ ᾽γίγνωσκον ὑπὲρ Εὐριπίδου δίκαια,\apparat{542}{ἅ γιγνωσκον R : \lecon{ἅ ᾽γίγνωσκον} Bentley}\\
διὰ τοῦτο τιλλομένην με δεῖ δοῦναι δίκην ὑφ᾽ ὑμῶν;\\\passeLigne}%
\end{groupe}%
%% v. 544-548
\begin{groupe}%
\personnageGrec{Μίκα}%
\tetramIambCat{οὐ γάρ σε δεῖ δοῦναι δίκην; ἥτις μόνη τέτληκας\\
ὑπὲρ ἀνδρὸς ἀντειπεῖν, ὃς ἡμᾶς πολλὰ κακὰ δέδρακεν\apparat{545}{ὡς R : \lecon{ὃς} Grynaeus (\italique{qui} Divus)}\\
ἐπίτηδες εὑρίσκων λόγους, ὅπου γυνὴ πονηρὰ\apparat{546}{ἐξεπίτηδες R : \lecon{ἐπίτηδες} Scaliger}\\
ἐγένετο, Μελανίππας ποιῶν Φαίδρας τε· Πηνελόπην δὲ\\
οὐπώποτ᾽ ἐποίησ᾽, ὅτι γυνὴ σώφρων ἔδοξεν εἶναι.\apparat{548}{ἐποίησεν R : \lecon{ἐποίησ᾽} Scaliger}\\\passeLigne}%
\end{groupe}%
%% v. 549-550
\begin{groupe}%
\personnageGrec{Μνησίλοχος}%
\tetramIambCat{\ligne{549}ἐγὼ γὰρ οἶδα ταἴτιον. Μίαν γὰρ οὐκ ἂν εἴποις\apparat{549}{ἂν ex ἂρ corr. \Rpc}%
\commentaire{549-50}{Μίαν γὰρ οὐκ ἂν εἴποις / τῶν νῦν γυναικῶν Πηνελόπην, Φαίδρας ἁπαξαπάσας.}{Le comique Euboulos paraît s'être inspiré de ce passage :%
\citComm{ὦ Ζεῦ πολυτίμητ', εἶτ' ἐγὼ κακῶς ποτε\\
ἐρῶ γυναῖκας ; νὴ Δί' ἀπολοίμην ἄρα\\
πάντων ἄριστον κτημάτων. Εἰ δ' ἐγένετο\\
κακὴ γυνὴ Μήδεια, Πηνελόπη δέ <γε>\\
μέγα πρᾶγμ'. ᾿Ερεῖ τις ὡς Κλυταιμνήστρα κακή·\\
Ἄλκηστιν ἀντέθηκα χρηστήν. ᾿Αλλ' ἴσως\\
Φαίδραν ἐρεῖ κακῶς τις · ἀλλὰ νὴ Δία\\
χρηστή τις ἦν μέντοι - τίς ; οἴμοι δείλαιος,\\
ταχέως γέ μ' αἱ χρησταὶ γυναῖκες ἐπέλιπον,\\
τῶν δ' αὖ πονηρῶν ἔτι λέγειν πολλὰς ἔχω.\\
\francais{Zeus très honoré, moi, dire du mal des femmes ? Ah! plutôt mourir ! Des femmes, le plus précieux de tous les biens ! Car si Médée s'est montrée méchante femme, Pénélope, elle, fut un vrai trésor. On dira que Clytemnestre fut mauvaise ? Je mets en regard la vertueuse Alceste. Mais peut-être médira-t-on de Phèdre ? Eh, par Zeus, vertueuse sans contredit fut.... qui donc ? Ah! malheureux ! les femmes vertueuses m'ont de suite fait faute, quand des mauvaises j'en ai encore un tas à citer.}}%
(Meineke, \italique{Fragm. Com. Gr}., t. III, p. 260. [$=$ fragment 115 (116. 117) de l'édition Kassel-Austin, \italique{Poet. Com. Gr.}]) (Note de Willems.)}
\\
τῶν νῦν γυναικῶν Πηνελόπην, Φαίδρας δ᾽ ἁπαξαπάσας.\apparat{550}{ἁπαξαπάσας R, M\expo{pc} : \lecon{δ᾽ ἁπαξαπάσας} M\expo{ac}}%
\commentaire{550}{Πηνελόπην, Φαίδρας.}{Si Phèdre est souvent citée par Aristophane comme (anti-)modèle d'impudicité (voir le vers 153 ; les \italique{Grenouilles} v. 1043 ; fr. 469(453) K.-A.(Kock)), Mélanippe apparaît seulement ici dans les œuvres conservées d'Aristophane.%
%\eff{développer sur Mélanippe ?}
}%
\\\passeLigne}%
\end{groupe}%
%% v. 551-552a
\begin{groupe}%
\ledpb%
\personnageGrec{Μίκα}%
\tetramIambCat{ἀκούετ᾽ ὦ γυναῖκες οἷ᾽ εἴρηκεν ἡ πανοῦργος\\
ἡμᾶς ἁπάσας αὖθις αὖ.\\\passeLigne}%
\end{groupe}%
%% v. 552b-553
\begin{groupe}%
\personnageGrec{Μνησίλοχος}%
\tetramIambCat{\ligne{552}\decalage{ἡμᾶς ἁπάσας αὖθις αὖ.}καὶ νὴ Δί᾽ οὐδέπω γε\\
εἴρηχ᾽ ὅσα ξύνοιδ᾽· ἐπεὶ βούλεσθε πλείον᾽ εἴπω;\apparat{553}{πλεῖον R : \lecon{πλείον᾽} Bentley}\\\passeLigne}%
\end{groupe}%
%% v. 554
\begin{groupe}%
\personnageGrec{Μίκα}%
\tetramIambCat{ἀλλ᾽ οὐκ ἂν ἔτ᾽ ἔχοις· ὅσα γὰρ ᾔδησθ᾽ ἐξέχεας ἅπαντα.\testimonia{554}{(ᾔδησθ᾽) \EtymGen (A) ᾔδεισθα}\apparat{554}{ᾔδεις R : \lecon{ᾔδησθ᾽} Hirschig (Ar. \italique{Vesp}.146)}\\\passeLigne}%
\end{groupe}%
%% v. 555-557a
\begin{groupe}%
\personnageGrec{Μνησίλοχος}%
\tetramIambCat{μὰ Δί᾽ οὐδὲ τήν γε μυριοστὴν μοῖραν ὧν ποιοῦμεν.\apparat{555}{οὐδὲ τὴν R : οὐδέπω τὴν Dobree : \lecon{οὐδὲ τήν γε} Faber \ligneApp{556} ἔπειτά γ᾽ R : \lecon{ἐπεὶ τάδ᾽} Dindorf (1825) \espApp οὐχ R : \lecon{οὐκ} Μυ2 \espApp τλεγγίδας R : \lecon{στλεγγίδας} \Mon \ligneApp{557} τὸν σῖτον R, Schol., Suid. σ 512 : \lecon{τὸν κύσθον} Bibi, cf. v.1114}\\
ἐπεὶ τάδ᾽ οὐκ εἴρηχ᾽, ὁρᾷς, ὡς στλεγγίδας λαβοῦσαι
\commentaire{556-7}{ὁρᾷς ὡς τλεγγίδας λαβοῦσαι / ἔπειτα σιφωνίζομεν τὸν σῖτον.}{Voilà un trait sur les mœurs féminines qui n'a pas été compris des différents exégètes ; c'est donc à nous que revient la tâche – ô combien délicate ! – d'expliquer ce qui répugne aux oreilles. Muses, venez m'aider... mais vous êtes pucelles ; Muses, ne bougez donc. Ah ! que l'on se sent seul ; je ne peux même pas lancer en avant l'honorable professeur Jeffrey Henderson, le plus fin limier des obscénités d'Aristophane et dont le groin fureteur parvient même à les débusquer là où elles ne sont pas (celle-ci, du moins, lui aura échappé.) Mais la recherche est un plaisir solitaire. Attelons-nous y donc ! Pris littéralement, le texte a de quoi surprendre : « tu vois comme nous saisissons les strigiles, et qu'ensuite nous siphonnons le pain. (sic) » Si vous êtes pris de l'envie de corriger, avec Küster, τὸν σῖτον en τὸν οἶνον (l'on siphonnera plutôt du vin que du pain), renseignez-vous sur ce qu'est un strigile (une sorte de spatule ou plaque rectangulaire tout en longueur, légèrement recourbée pour servir de racloir lors de la toilette) ; j'offre l'édition d'Aristophane de Küster à qui parviendra à se saouler en buvant avec ce genre de coupe. (Il faudra d'ailleurs un jour que l'on m'explique l'intérêt d'introduire des émendations si c'est pour laisser un texte tout aussi obscur – et retouché en prime.) Qu'à cela ne tienne, ce n'est pas ce genre de peccadille qui doit arrêter Austin \& Olson : Aristote a dit un jour (\italique{Tropiques},145a24) qu'on pouvait se servir d'un strigile pour puiser de l'eau (hé ! ce sont des choses qui arrivent lorsqu'on est au bain ; A.\&O. auraient pu citer la phrase suivante qui leur était adressée~: Ἀλλ᾽ ὅμως εἴ τις ὁρίσαιτο τὴν στλεγγίδα ὄργανον πρὸς τὸ ἀρύειν, ἡμάρτηκεν, « Mais cependant, celui qui définirait le strigile comme instrument pour puiser se trompe »). Non, l'hypothèse d'ivrognesses s'enivrant à grandes rasades de strigiles est à écarter, sauf à m'expliquer : 1° comment boire au strigile \ \ \ 2° pourquoi aller chercher dans la salle de bain un instrument qui n'est pas censé en bouger \ \ \ 3° et quand bien même, qu'y aurait-il de déshonorant dans cette pratique ? \ \ \ Il faut chercher ailleurs, car les tentations des bains ne sont pas d'ordre biberonnesque. Aristophane vitupère ailleurs\marqueur{1} contre cette propension des Athéniens décadents à se délicater : les bains chauds sont un exemple parmi d'autres de ces mauvaises habitudes qui amollissent les cœurs et durcissent les vices. Pour comprendre exactement où il veut en venir, c'est Euripide lui-même qui nous donnera la clef de la salle de bain : dans son drame satyrique du \italique{Cyclope}, le héros aux-mille-tours-lubriques, se désole sur « son cher siphon » (v. 339, τὸν σίφωνα τὸν φίλον) qui depuis trop longtemps vit dans la solitude – solitude vis-à-vis du beau sexe ; le verbe χηρεύω servant aussi à désigner l'état de veuvage. Il parle évidemment de son πέος~; une fois attesté ce sens pour σίφων, celui du verbe σιφωνίζω s'éclaire d'une lumière nouvelle (le verbe est d'autant plus amusant qu'il joue sur le rapprochement, cher à Aristophane, entre βινεῖν et πίνειν).\\
La présence étonnante de τὸν σῖτον s'explique alors aisément : la leçon originale était τὸν κύσθον $=$ \latin{pudenda muliebria}, très proche phonétiquement de τὸν σῖτον. 
Iotacisme aidant, il s'agit d'une simple \italique{métathèse} du `s' (avec chute -- sans doute postérieure -- du `k').
Il est remarquable de noter que la même métathèse se produit au vers 1114 : le rustre archer Scythe, dans l'optique de prouver que Mnésiloque n'est pas une femme, conseille d'aller y voir de plus près :%
\citComm{σκέψαι †τὸ σκύτο†.}%
Scaliger \italique{à la grâce du mètre} -- et surtout à la grâce du sens -- avait rétabli :%
\citComm{σκέψαι τὸ κύστο ($=$ τὸν κύσθον dans le parler scythe).}%
Il est à croire qu'un scribe des \italique{Thesmophories} confondait avec un autre le mot grec pour désigner le « plus bel apanage » du corps féminin, qui tentait tendrement tant Brassens \& Villon : il s'y sera trompé par deux fois dans notre pièce (\latin{bis repetita...}).
J'arrête mes explications ici, de peur qu'un Calliclès de passage ne m'assène :%
\citComm{Καλλίκλης-Φιλόλογος --- οὐκ αἰσχύνῃ εἰς τοιαῦτα ἄγων, ὦ Πέπωκες, τοὺς λόγους;\marqueur{2}}%
%La présence étonnante de τὸν σῖτον s'explique alors aisément : la leçon originale était τὴν μᾶζαν, qu'on traduit faute de mieux par \italique{la galette}, et glosée τὸν σῖτον, \italique{le pain} par un érudit n'ayant pas tout saisi... car la μᾶζα peut désigner aussi le « plus bel apanage » du corps féminin, qui tentait tendrement tant Brassens \& Villon (cf. l'ingénieuse – et crue ! – explication du vers 835 des \italique{Acharniens} par Jeffrey Henderson\marqueur{2} – oreilles chastes s'abstenir ! – et, si la démonstration ne suffit pas, on pourra finir de se convaincre avec l'ordurier Hipponax\marqueur{3}). 
%Un scribe aura alors pris le Pirée pour un homme, l'attribut de la femme pour quelque denrée peu appétissante\marqueur{4} et une glose intempestive pour la leçon originale\marqueur{5}. 
(Moi qui m'étais lancé dans ce commentaire avec la ferme intention de disperser la horde libidineuse des commentateurs qui collent aux basques d'Aristophane et maculent\marqueur{3} notre poète d'une horrible réputation de pornographe, il me déplaît de leur donner du grain à moudre. Las !.. Philologus semble moins à l'affût -- ou moins au fait -- des travers féminins. Aussi l'avons-nous suppléé en ce passage -- une fois n'est pas coutume !)\\% -- pour un autre passage d'Aristophane
%Plutus 1036 Γραῦς -- διὰ δακτυλίου μὲν οὖν ἔμεγ᾽ ἂν διελκύσαις.
%Χρεμύλος -- εἰ τυγχάνοι γ᾽ ὁ δακτύλιος ὢν τηλία.
%Διὰ δακτυλίου δεῖ σε ἑλκυσθῆναι Shakespeare (Souda online ; français Willems ; italien i proverbi greci ; dia koskinou la véritable expression ?)
Profitons-en pour remarquer à quel point les griefs des femmes envers Euripide ne sont pas sérieux dans la bouche d'Aristophane (si Euripide est misogyne, c'est bien là le moindre de ses défauts) ; car si les maris athéniens, s'en revenant d'une pièce d'Euripide, regardaient d'un œil soupçonneux leur femme tenant une marmite, avec quels yeux vont-ils dorénavant observer leur femme se saisir d'un strigile ou de tout instrument de forme analogue ?! Quoi qu'il en soit, comme pour l'histoire de l'amant qui s'échappe (v. 501), Mnésiloque n'invente rien mais reprend des thèmes certainement populaires du Landerneau athénien : ainsi, le détournement à fin lubrique du strigile se retrouve illustrée, ce me semble, sur une coupe d'Onésimos\marqueur{4}, antérieure à la pièce d'Aristophane. Pour finir – était-il besoin de préciser ?.. – précisons tout-de-même : Mnésiloque mime en lançant son « ὁρᾷς » à Mika : « Ainsi font, font, font, les petites midinettes. Un siphon, fond, fond, on se siphonne toutes à fond.»}%
\debutNote{556-7}%
\note{1}{C'est un point de la querelle du Juste et de l'Injuste des \italique{Nuées} (v. 1044-1046) :
\citComm{Ἄδικος Λόγος — σε θερμῷ φησι λοῦσθαι πρῶτον οὐκ ἐάσειν.\\
καίτοι τίνα γνώμην ἔχων ψέγεις τὰ θερμὰ λουτρά;\\
Δίκαιος Λόγος — ὁτιὴ κάκιστόν ἐστι καὶ δειλὸν ποιεῖ τὸν ἄνδρα.\\
\francais{Raisonnement Injuste — il ne te permettra pas, dit-il, de te baigner dans l'eau chaude.\\
(Au Juste :) Mais d'après quel principe improuves-tu les bains chauds? \\
Raisonnement Juste — Parce qu'ils sont très pernicieux et qu'ils rendent l'homme lâche.}}%
Platon, d'accord sur ce point avec Aristophane comme sur beaucoup d'autres, nomme les bains chauds γεροντικὰ λουτρά, « des bains de vieillards » (\italique{Lois}, VI, 761c). Érasme en a fait un adage (2061): ἁπαλοὶ θερμολουσίαις, \italique{amollis par les bains chauds}.}%
%\note{2}{p. 113 de sa \italique{Maculate Muse}.}%
%\note{3}{fragment 28 de l'édition Kassel-Austin ; se reporter aux notes du \nieme{6} fragment pour l'éclaircissement.\eff{Et chez Masson ?}}%
%\note{4}{La μᾶζα étant un mets tout-à-fait commun et fait à la va-vite quand le σῖτος s'achetait chez un boulanger. Pour en savoir davantage sur la μᾶζα, l'on se reportera avec intérêt à la note essentielle de Willems (t. I, p. 96-99), reproduite à la fin.}%
%\note{5}{Confusion facilitée par la scansion identique.}%
%\note{2}{Iotacisme aidant, il s'agit d'une simple \italique{métathèse}.}%
\note{2}{Platon, \italique{Gorgias} 494e.}%
\note{3}{Le chef de cette meute, le sus-nommé Jeffrey Henderson, ne s'en cache même pas : il affiche clairement ses `ambitions' jusque dans le titre de son exégèse porcine d'Aristophane : \italique{Maculate Muse}. Le livre fait autorité et cependant une bonne moitié -- sans exagération aucune -- n'est que pure élucubrication ! J'en viens à me demander si Hipponax ou Archiloque, qui traînent derrière eux une réputation plus ordurière encore que celle d'Aristophane, ne seraient pas eux aussi victimes du même genre de calomniateurs -- l'affreux Philologus ! ὦ βορβοροτάραξι Φιλόλογε !}%
\note{4}{Bâle, Antikenmuseum BS 440. ARV 326, 86 \italique{bis}. \italique{Para} 359. Elle est reproduite par John Boardman, fig. 233 in \italique{Les vases athéniens à figures rouges}, \italique{La période archaïque} ; première édition sous le titre \italique{Athenian Red Figure Vases}, \italique{The Archaic Period}, 1975. (Le lecteur pourra la retrouver reproduite aux annexes de ce livre.) En parcourant le livre de Boardmann, cette coupe a retenu mon attention : elle me paraît parfaitement illustrer l'anecdote de Mnésiloque. (Reconnaissons tout d'abord que l'interprétation est sujette à caution : j'ai montré cette coupe à différents amis et ai eu droit à autant d'interprétations qu'il y avait d'amis... interprétations qui en disaient d'ailleurs plus long sur leur propre psychologie que sur les mœurs des Athéniens ; profitons-en donc pour suggérer à l'éventuel psychanalyste qui se serait égaré à lire ces lignes d'abandonner fissa ses infantiles tests de Rorschach et de passer aux coupes athéniennes du \nieme{VI} siècle : dix vases athéniens d'Onésime seraient aussi révélateurs des obsessions ou perversions des supposés psychotiques que les taches d'encre du regretté Hermann, arraché trop tôt à l'affection des siens... Sans parler du gain en valeur artistique : incommensurable ! Refermons cette parenthèses et retournons à notre coupe :) qu'y ai-je donc vu ? À droite, un satyre (sa calvitie naissante est caractéristique des représentations de satyres de cette époque ; c'est entre autre sur ce détail de la laideur physique de Socrate, visible sur ses bustes, qu'Alcibiade rapproche selon moi Socrate du satyre Silène), un satyre donc, en train de se strigiler le fondement, pointe (du doigt) le Mont de Vénus d'une femme nue, lui désignant par là l'usage impropre du strigile auquel il lui arrive de céder. La clef qui me manque pour être tout-à-fait sûr de la bonne compréhension de cette scène est l'identification de l'objet dans la main de la femme : je crois y reconnaître un strigile enroulé dans du tissu. (plus recommandé pour atteindre le plaisir que la lame de bronze brute.) D'aucun y a vu le pagne du jeune homme à gauche. J'interprète cette scène comme un résumé des vices féminins de l'imaginaire grec : la femme est incapable de se contenir, elle cède à tous les plaisirs : le plaisir de la vigne (la femme est adossée à une outre) ; le plaisir de la chair (l'adultère est suggéré avec l'amant à gauche qui se rhabille) et -- toujours insatiable ! -- le plaisir des dieux ! (nue, la voilà habillée pour l'hiver !..)\\
Une autre hypothèse d'un mien ami quant à cette scène vaut d'être rapportée : il l'avait interprétée quelque peu différemment : tout en restant dans le domaine d'Onan, il proposait de garder τὸν σῖτον et de comprendre : \italique{le pain} [d'argile] servant à récupérer les peaux mortes \& autres hormones raclées avec l'huile à l'aide du strigile (les στλεγγίσματα et autre γλοιός), raclures qui étaient ensuite parfois monnayées (notamment lorsqu'elle provenaient de virils athlètes)... et pourquoi pas par quelque vicelarde pour s'en foutre plein les zones érogènes ? Son explication avait l'avantage indéniable de conserver la leçon du manuscrit... Mais je crains que son imagination ne se soit quelque peu emballée : je n'ai pas trouvé d'exemple d'acception du mot σῖτος pour désigner ce «~pain d'argile~». Et si Pline nous confirme un usage médical des στλεγγίσματα, l'usage détourné, qui lui a été suggéré par certains papiers de « recherche », est un pur fantasme des modernes scholiastes.}\\
ἔπειτα σιφωνίζομεν τὸν κύσθον.\testimonia{557}{(σιφωνίζομεν τὸν σῖτον) Suidas σ 512 \espTest (σιφωνίζομεν) Pollux VI 19}\\\passeLigne}%
\end{groupe}%
%% v. 557b
\begin{groupe}%
\personnageGrec{Μίκα}%
\tetramIambCat{\ligne{557}\decalage{ἔπειτα σιφωνίζομεν τὸν κύσθον.}ἐπιτριβείης.\\\passeLigne}%
\end{groupe}%
%% v. 558-559a
\begin{groupe}%
\ledpb%
\personnageGrec{Μνησίλοχος}%
\tetramIambCat{ὥς τ᾽ αὖ τὰ κρέ᾽ ἐξ Ἀπατουρίων ταῖς μαστροποῖς διδοῦσαι\apparat{558}{ὥς τ᾽ αὐτὰ R : ὥς ταῦτα \Mon : \lecon{ὥς τ᾽ αὖ τὰ} Zanetti}\\
ἔπειτα τὴν γαλῆν φαμεν–
\commentaire{558-9}{ὥς τ᾽ αὐτὰ κρέ᾽ ἐξ Ἀπατουρίων ταῖς μαστροποῖς διδοῦσαι / ἔπειτα τὴν γαλῆν φαμεν.}{Si Mnésiloque avait commencé un discours tout en nuance, il a bien vite perdu de sa prudence et se met à surpasser le maître Euripide dans la diffamation du genre féminin. Le ton monte et Mnésiloque en est maintenant à dénoncer les pratiques les plus honteuses des Athéniennes les plus dissolues.\\
Nous ne nous attarderons pas sur ce que pouvaient bien procurer les procureuses en échange de ces détournements de viandes\marqueur{1}.
Plus intéressante est l'excuse alléguée pour expliquer cette soudaine disparition carnée. La \italique{team 1\expo{er}} degré des commentateurs d'Aristophane n'a pas manqué de faire remarquer que la belette était un animal chapardeur\marqueur{2}. Certes ! mais ce pauvre animal a surtout \italique{bon dos} : on l'accuse bien facilement ; ainsi aux vers 1150-1152 de la \italique{Paix} :%
\citComm{ἦν δὲ καὶ πυός τις ἔνδον καὶ λαγῷα τέτταρα,\\
εἴ τι μὴ 'ξήνεγκεν αὐτῶν ἡ γαλῆ τῆς ἑσπέρας·\\
ἐψόφει γοῦν ἔνδον οὐκ οἶδ᾽ ἅττα κἀκυδοιδόπα\\
Il y avait aussi au logis de l'amouille et quatre morceaux de lièvre ;\\
à moins que la belette n'en ait volé hier soir ;\\
car il se passait là-dedans je ne sais quel bruit et quel remuement.}%
L'allusion au remue-ménage n'est pas limpide mais il est clair que le sens ne peut être pris littéralement : qui, entendant une belette le dévaliser, ne s'armerait aussitôt d'un bâton pour châtier la voleuse ? (Noter d'ailleurs ici aussi la présence de l'article : ἡ γαλῆ.) 
Aristophane joue certainement sur un détournement d'une locution proverbiale : la belette attique étant notre \italique{bouc-émissaire} à qui l'on fait \italique{porter le chapeau}. 
Conformité du grec et du français : la belette joue le même rôle que notre chat, animal qui n'a pas encore été introduit dans l'Athènes du \nieme{V} siècle : il chasse les rongeurs. Or, que lit-on chez La Fontaine ?%
\citComm{Le sommeil quitta son logis,\\
Il eut pour hôte les soucis,\\
Les soupçons, les alarmes vaines.\\
Tout le jour il avait l'œil au guet ; \textbf{et la nuit,\\
Si quelque chat faisait du bruit,\\
Le chat prenait l'argent.}\marqueur{3}}%
Ce n'est pas moi qui m'abuse : remplacez l'argent par la nourriture et les deux passages disent exactement la même chose.
Voici ce que l'on trouve dans le \italique{Dictionnaire de la langue verte} d'Alfred Delvau :%
\citComm{« C'est le chat ! » : Expression de l'argot du peuple, qui souligne ironiquement un doute, une dénégation. Ainsi, quelqu'un disant : Ce n'est pas moi qui ai fait cela. – Non ! c'est le chat, lui répondra-t-on.}%
Un autre passage de la \italique{Paix} montre encore mieux ce sens imagé : aux vers 793-795 :%
\citComm{καὶ γὰρ ἔφασχ᾽ ὁ πατὴρ ὃ παρ᾽ ἐλπίδας\\
εἶχε τὸ δρᾶμα γαλῆν τῆς\\
ἑσπέρας ἀπάγξαι.\\
Aussi bien leur père prétend-il qu'un drame qu'il pensait tenir contre toute attente fut étranglé le soir par une belette.}%
Juste avant ce passage, Aristophane faisait un honnête résumé des bienfaits qu'il a apportés à la Comédie, il en vient ensuite à étriller les nullités de ses poétereaux de contemporains. Nous en sommes à Carcinos \& Fils qui font des pieds et des mains pour pouvoir danser – c'est-à-dire obtenir la subvention d'un chœur.
Mais une fois le chœur accordé – à rebours de ce qui était attendu ! (παρ᾽ ἐλπίδας) [du bon goût ?] – Carcinos-père fait volte-face et prétend que son drame « fut étranglé le soir par une belette » (c'est encore le soir que la belette commet son forfait). 
Le scholiaste nous est ici précieux puisqu'il précise que Carcinos voulait représenter une pièce intitulée les \italique{Souris} (Μῦες\marqueur{4}) ; encore que la chronique du temps ne nous ait pas conservé l'affaire, il paraît aisé de reconstituer peu ou prou ce qui s'est passé. Les poètes devaient présenter à l'archonte éponyme leur projet. 
Pour les Dionysies, l'archonte sélectionnait trois tétralogies qui recevaient chacune le soutien financier d'un riche Athénien : le chorège. 
Ce n'est qu'alors que le poète se mettait à écrire. Et c'est à cette étape que les forces auront probablement manqué à notre Carcinos… ce dernier croyant s'en tirer avec un trait d'esprit : s'il ne peut pas présenter ses \italique{Souris}, c'est que la belette les a étranglées. 
Le mot ne peut être \italique{bon} que s'il existait effectivement une expression pour faire endosser la responsabilité d'un méfait à cette belette proverbiale.}%
\debutNote{558-9}%
\note{1}{Contentons-nous de renvoyer à la magistrale note d'Alphonse Willems sur la fête des Apaturies, t. II, pp. 118-120 ; le lecteur y découvrira l'importance démocratique de la κρεανομία, l'institution de la distribution de viande. Sur les débordements que de telles orgies pouvaient entraîner, l'on pourra – un exemple parmi d'autres – essayer de comprendre le jeu de mots d'Aristophane sur l'ἀνάρρυσις (au vers 890 de la \italique{Paix}), le second jour des Apaturies.}%
%\note{1}{Voir par exemples les vers 363-364 des \italique{Guêpes}. Lucien cite aussi l'adage ἁρπακτικώτεροι τῶν γαλῶν (\italique{Le pêcheur}, 34).}%
\note{2}{Voir par exemples les vers 363-364 des \italique{Guêpes}. Lucien, \italique{Piscator}, 34 : ἁρπακτικώτεροι τῶν γαλῶν.}%
\note{3}{Il s'agit – vous l'aurez bien sûr reconnue ; sinon, cessez immédiatement votre lecture et ne revenez que lorsque vous la connaîtrez par cœur ! – de la seconde fable du huitième livre : \italique{Le Savetier et le Financier}.}%
\note{4}{Sans doute aucun le drame satyrique. Μῦς étant probablement un nom propre, comme chez Hérodote (VIII, 133), dans la \italique{Souda} (ο 680) ou chez Athénée (XI, 782b), les Μῦες désignant les enfants de ce Μῦς. Je m'inscris en faux contre l'interprétation de Paul Mazon (\italique{La} Paix\italique{, texte grec publié avec une introduction, des notes critiques et explicatives} (1904), v. 793) reprise par Taillardat (§792, p.463 n.5) qui pensent que Μῦες était le surnom de ces débauchés de Carcinites. L'on s'expliquerait difficilement la confusion du scholiaste (présent sur tous les bons manuscrits et repris par Suidas κ 396). Objection plus forte : cela n'expliquerait pas le texte de la \italique{Paix} : la belette (γαλῆ) est accusé d'avoir étranglé nuitamment non les choreutes mais la pièce elle-même (τὸ δρᾶμα).}\\\passeLigne}%
\end{groupe}%
%% v. 559b
\begin{groupe}%
\personnageGrec{Μίκα}%
\tetramIambCat{\ligne{559}\decalage{ἔπειτα τὴν γαλῆν φαμεν–}τάλαιν᾽ ἐγώ· φλυαρεῖς.\\\passeLigne}%
\end{groupe}%
%% v. 560-562a
\begin{groupe}%
\personnageGrec{Μνησίλοχος}%
\tetramIambCat{οὐδ᾽ ὡς τὸν ἄνδρα τῷ πελέκει γυνὴ κατεσπόδησεν,\testimonia{560}{Suidas κ 996 (ἑτέρα pro γυνὴ et κατεσπόδησε vel κατεσπόδωσεν)}\apparat{560}{\lecon{τὸν ἄνδρα τῷ πελέκει γυνὴ} R : ἑτέρα τὸν ἄνδρα τῷ πελέκει Suid. κ 996}\\
οὐκ εἶπον· οὐδ᾽ ὡς φαρμάκοις ἑτέρα τὸν ἄνδρ᾽ ἔμηνεν,\\
οὐδ᾽ ὡς ὑπὸ τῇ πυέλῳ κατώρυξέν ποτε—\apparat{562}{\lecon{ποτε} R : ποτ᾽ ed. pr.}\\\passeLigne}%
\end{groupe}%
%% v. 562b
\begin{groupe}%
\personnageGrec{Μίκα}%
\tetramIambCat{\ligne{562}\decalage{οὐδ᾽ ὡς ὑπὸ τῇ πυέλῳ κατώρυξέν ποτε—}ἐξόλοιο.\\\passeLigne}%
\end{groupe}%
%% v. 563a
\begin{groupe}%
\personnageGrec{Μνησίλοχος}%
\tetramIambCat{Ἀχαρνικὴ τὸν πατέρα.\apparat{563}{\lecon{ἀχαρνικὴ} R : Ἁχαρνικὴ Dobree}%
%\commentaire{563}{Ἀχαρνικὴ}{\eff{Une charbonnière ? qui aura charbonné son mari père ? Ou alors : après la hache et le poison, reste la corde (kat)ἀρτῶ à la la place de katwrucen Τὰ τρία τῶν εἰς θάνατον .}Ἀχαρνικῇ πύλῃ (épigraphie)
%κατώρης· κάτω ῥέπων <- (kat -ὄρνυμι) l'aura-t-elle précipité au pied de la porte d'acharne ? (Pris littéralement, le texte indique seulement la façon dont la femme a tenté de se débarrasser du corps et non comment elle l'a tuée -- ce qui constitue le réquisitoire de Mnésiloque ; de plus, si l'affaire était connue de tous, pourquoi préciser ici que c'est une acharnienne qui a fait e coup ? Cédons donc donc à la philologienne manie de vouloir à tout prix retoucher le texte -- au moins, contentons- nous de circonscrire cette retouche au commentaire et non au texte présenté au lecteur) "Je nai pas dit : ni comment, par un poison, une autre femme affolit son mari, ni comment au pied de la porte}
\\\passeLigne}%
\end{groupe}%
%% v. 563b
\begin{groupe}%
\personnageGrec{Μίκα}%
\tetramIambCat{\ligne{563}\decalage{Ἀχαρνικὴ τὸν πατέρα.}ταυτὶ δῆτ᾽ ἀνέκτ᾽ ἀκούειν;\\\passeLigne}%
\end{groupe}%
%% v. 564-565
\begin{groupe}%
\personnageGrec{Μνησίλοχος}%
\tetramIambCat{οὐδ᾽ ὡς σὺ τῆς δούλης τεκούσης ἄρρεν εἶτα σαυτῇ\\
τοῦθ᾽ ὑπεβάλου, τὸ σὸν δὲ θυγάτριον παρῆκας αὐτῇ.\\\passeLigne}%
\end{groupe}%
%% v. 566-567a
\begin{groupe}%
\personnageGrec{Μίκα}%
\tetramIambCat{οὔ τοι μὰ τὼ θεὼ σὺ καταπροίξει λέγουσα ταυτί,\testimonia{566}{Suidas κ 730 (sed καταπροίξῃ pro καταπροίξει)}\\
ἀλλ᾽ ἐκποκιῶ σου τὰς ποκάδας.\testimonia{567}{(ἐκποκιῶ σου τὰς ποκάδας) Suidas ε 588 (ἐκπιῶ cod. A) ; Zonaras 672 ἐκποκιῶ \espTest (ἐκποκιῶ) Heschius ε 1637 ἐκποκιῶ ; \EtymMagnum 323.25 \espTest (ποκάδας) Hesychius π 2750 ποκάδες $=$ Photius π 1010 π\{λ\}οκάδες}\apparat{567}{\lecon{ποκάδας} R, Suid. ε 588 : πλοκάδας Schol. \espApp δὲ R : \lecon{δὴ} Bothe}\\\passeLigne}%
\end{groupe}%
%% v. 567b
\begin{groupe}%
\personnageGrec{Μνησίλοχος}%
\tetramIambCat{\ligne{567}\decalage{ἀλλ᾽ ἐκποκιῶ σου τὰς ποκάδας.}οὐ δὴ μὰ Δία σύ γ᾽ ἅψει.\\\passeLigne}%
\end{groupe}%
%% v. 568a
\begin{groupe}%
\personnageGrec{Μίκα}%
\tetramIambCat{καὶ μὴν ἰδού.\\\passeLigne}%
\end{groupe}%
%% v. 568b
\begin{groupe}%
\personnageGrec{Μνησίλοχος}%
\tetramIambCat{\ligne{568}\decalage{καὶ μὴν ἰδού.}καὶ μὴν ἰδού.\\\passeLigne}%
\end{groupe}%
%% v. 568c
\begin{groupe}%
\personnageGrec{Μίκα}%
\tetramIambCat{\ligne{568}\decalage{καὶ μὴν ἰδού.καὶ μὴν ἰδού.}λαβὲ θοἰμάτιον Φιλίστη.
\commentaire{568}{λαβὲ θοἰμάτιον Φιλίστη.}{Mettre bas les manteaux était de mise avant d'en venir aux mains. Ainsi au vers 408 des \italique{Guêpes} ; autre exemple : le fragment d'Hipponax\marqueur{1} contre Boupalos\marqueur{2} : Boupalos, sculpteur de Chios, avait fait d'Hipponax un portrait où il exagérait encore la laideur du modèle. Pour se venger, le poète écrivit contre lui des vers mordants et pleins de menaces. Suidas (β 452) nous a conservé celui-ci, qui paraît avoir passé en proverbe : %
\citComm{λάβετε μευ ταἰμάτια [θοιμάτιον codd.], κόψω Βουπάλου τον οφθαλμόν.\\
\francais{Prenez mon manteau, que je fasse sauter un œil à Boupalos.}}%
}%
\debutNote{568}%
\note{1}{Fr. 121 de l'édition Degani : \italique{Hipponactis testimonia et fragmenta}, Teubner, 1991 ; fr. 120 chez Olivier Masson, voir aussi son introduction p. 13.}%
\note{2}{Il est question de ce Boupalos au vers 361 de \italique{Lysistrata} -- le fragment sus-cité d'Hipponax nous a d'ailleurs été transmis par le scholiaste à ce passage ; je reprends la note de Willems de ce passage, p. 351 du second tome.}\\\passeLigne}%
\end{groupe}%
%% v. 569a
\begin{groupe}%
\personnageGrec{Μνησίλοχος}%
\tetramIambCat{πρόσθιγε μόνον, κἀγώ σε νὴ τὴν Ἄρτεμιν—\apparat{569}{πρόσθες R : πρόσιθι Cobet : \lecon{πρόσθιγε} Willems}%
\commentaire{569}{πρόσθες μόνον.}{Willems substitue πρόσθιγε (du verbe προσθιγγάνω, \italique{toucher}) à πρόσθες qu'il soupçonne de ne pas être le mot propre. Il s'en explique succinctement dans la \nieme{16} note reproduite à la fin. D'aucuns ont essayé de justifier la leçon du texte en avançant une tournure elliptique : πρόσθες [τὴν χεῖρα] (expression qui se trouve effectivement au vers 1669 des \italique{Phéniciennes} d'Euripide\marqueur{1}) mais tous les exemples proposés comportent χεῖρα, qui semble indispensable au sens. Cobet avait donc proposé de lire πρόσιθι μόνον, \italique{approche seulement} ; ce que l'on dirait effectivement dans nos langues modernes mais le grec use d'une image différente : il dit \italique{jette(-moi) seulement ta main} (cf. les exemples déjà cités) ou bien \italique{touche(-moi) seulement} (cf. les exemples proposés par Willems auxquels on peut ajouter \italique{Lysistrata} v. 365, qui est l'exacte correspondance de son émendation : ἅψαι μόνον).}%
\debutNote{569}%
\note{1}{L'on cite d'autres exemples équivalents : le vers 933 des \italique{Nuées}: κλαύσει, τὴν χεῖρ᾽ ἢν ἐπιβάλλῃς ; le fragment 309(277) (éd. K.-A. = \italique{Zonarae lexikon} p.1168) : τὴν χεῖρα μὴ ᾿πίβαλλε, μὴ κλάων καθῇ ; \italique{Lysistrata}, vers 435-436 \& 439-440 : τὴν χεῖρά μοι ... προσοίσει ... κλαύσεται \& εἴ ... μόνον τὴν χεῖρ᾽ ἐπιβαλεῖς, ἐπιχεσεῖ πατούμενος.}%
\\\passeLigne}%
\end{groupe}%
%% v. 569b
\begin{groupe}%
\personnageGrec{Μίκα}%
\tetramIambCat{\ligne{569}\decalage{πρόσθιγε μόνον, κἀγώ σε νὴ τὴν Ἄρτεμιν—}τί δράσεις;\\\passeLigne}%
\end{groupe}%
%% v. 570
\begin{groupe}%
\personnageGrec{Μνησίλοχος}%
\tetramIambCat{τὸν σησαμοῦνθ᾽ ὃν κατέφαγες, τοῦτον χεσεῖν ποιήσω.\testimonia{570}{(χεσεῖν) \EpimAlphab ε 201}%
\commentaire{570}{τὸν σησαμοῦνθ᾽ ὃν κατέφαγες, τοῦτον χεσεῖν ποιήσω.}{«~Le gâteau de sésame que tu as engoulé, je vais te le faire chier.~» Mnésiloque aurait pu ajouter : que tu viens d'engouler \italique{en cachette} ; car ce \grec{σησαμοῦς}\marqueur{1} a été avalé en plein milieu des Thesmophories, le jour du \italique{jeûne}, \grec{νηστεία}.
Athénée (XIV, 647a) nous rapporte que les Syracusaines préparaient des gâteaux semblables de miel et de sésame, les \italique{mylles} pour leurs Thesmophories~:%
\citComm{ΜΥΛΛΟΙ. Ἡρακλείδης ὁ Συρακόσιος ἐν τῷ περὶ Θεσμῶν ἐν Συρακούσαις φησὶ τοῖς Παντελείοις τῶν Θεσμοφορίων ἐκ σησάμου καὶ μέλιτος κατασκευάζεσθαι ἐφήβαια γυναικεῖα, ἃ καλεῖσθαι κατὰ πᾶσαν Σικελίαν μυλλοὺς καὶ περιφέρεσθαι ταῖς θεαῖς.\\
\francais{Mylloi. Héraclide de Syracuse rapporte ce qui suit, dans son \italique{Traité des Rites religieux} : « Le jour des grandes Thesmophories, on fait à Syracuse avec de la farine, du miel et du sésame, la figure de la partie sexuelle d'une femme, que l'on porte en pompe, pour l'offrir aux déesses. On l'appelle \italique{mylle} par toute la Sicile. »}}%
C'est de peur que Mica risque de relâcher son gâteau. En effet, chez Aristophane la peur se manifeste presque toujours physiquement par des troubles du bas-ventre\marqueur{2}. On dira en français \italique{avoir les pétoches} -- littéralement.
\latin{Pulchre, bene, recte !} Mais l'essentiel nous échappe encore : Mnésiloque ne promet pas à Mica un petit coup en cas d'agression de sa part, mais une riposte non graduée : une pluie de horions à lui assainir le côlon ! Car il faut comprendre : «~tu vas en chier [de peur] \italique{malgré} le gâteau de sésame par toi engueulé~».
Ce genre de gâteaux avait la vertu d'être astringent\marqueur{3}. C'est ce qu'il ressort d'un passage des \italique{Nuées} (vers 506-508) : Strepsiade hésite à rentrer dans l'antre de Socrate, qu'il compare aimablement à celui de Trophonius : %
\citComm{\ \ \ \ \ \ \ \ \ \ \ \ \ \ \ \ \ \ \ \ \ \ \ \ \ \ \ \ \ \ εἰς τὼ χεῖρέ νυν\\
δός μοι μελιτοῦτταν πρότερον, ὡς δέδοικ᾽ ἐγὼ\\
εἴσω καταβαίνων ὥσπερ εἰς Τροφωνίου.\\
Mets-moi d'abord en main un gâteau au miel ; car j'ai peur de descendre, comme si je pénétrais dans l'antre de Trophonios.}%
Il faut bien-sûr prononcer d'un trait `εἰς Τροφωνίου' pour faire entendre στρόφος, la \italique{colique}\marqueur{4}.\\
%()\\
C'est aussi ce qui explique le bon mot de table qu'Athénée (I, 4d) rapporte :%
\citComm{ἀμίδα δότω τις ἢ πλακοῦντά τις δότω\\
Qu'on m'apporte un pot de chambre ou bien un gâteau !}%
Au contraire de tous ces gâteaux de miel, la συρμαία (sorte de rave noire que prisaient les Égyptiens pour leur régime ; voir le commentaire \italique{infra} aux vers 855-857) possède des vertus purgatives. 
Cette opposition est probablement la raison de la vraisemblable confrontation dans la pièce perdue des \italique{Danaïdes}\marqueur{5} entre deux vendeurs : celui de syrméa et celui d'ἐγκρίδες\marqueur{6} :%
\citComm{ἐν δ᾿ Ἀριστοφάνης Δαναΐσι συρμαιοπῶλαι καὶ ἐγκριδοπῶλαι.}%
%Profitons-en pour nous pencher sur un problème d'établissement de texte : un passage des \italique{Deipnosophistes} d'Athénée (XIV, 645e), fragment n° 269(256) d'Aristophane dans l'édition de Kassel \& Austin des \italique{Poetae Comici Graeci}. Athénée y définit ce que sont que les ἐγκρίδες :%
%\citComm{ΕΓΚΡΙΔΕΣ. πεμμάτιον ἑψόμενον ἐν ἐλαίῳ καὶ μετὰ τοῦτο μελιτούμενον.\\
%μνημονεύει αὐτῶν Στησίχορος διὰ τούτων· χόνδρον τε καὶ ἐγκρίδας ἄλλα τε πέμματα καὶ μέλι χλωρόν.\\
%μνημονεύει αὐτῶν καὶ Ἐπίχαρμος καὶ ἐν τοῖς Ἐγχειρογάστορσι Νικοφῶν.\\
%Ἀριστοφάνης δʼ ἐν Δαναίσιν καὶ ποιητήν [sic] φησιν αὐτῶν εἶναι ἐν τούτοις· μήτʼ αρμα [re-sic] εἶναι ἐγκριδοπώλην.\\
%ENCRIDES. C'était un petit gâteau qu'on faisait frire dans l'huile, après quoi on le frottait de miel.\\
%Stésichore en parle ainsi : « Du gruau, des \italique{encrides}, et autres pâtisseries, et du miel tout récent. »\\
%Il en est aussi fait mention dans Épicharme et dans les \italique{Mercenaires} de Nicophon.\\
%Aristophane dit dans les \italique{Danaïdes} $[$ ... $]$}%
%Faut-il voir dans le titre de la pièce de Nicophon, les \italique{Mercenaires}, ἐγχειρογάστορες, un jeu de mots qui ferait entendre des problèmes gastriques, raison de la présence desdits ἐγκρίδες, remèdes supposés ? C'est possible, mais peu nous importe : concentrons-nous sur Aristophane !\\
%Ce fragment des \italique{Danaïdes} est à rapprocher de Pollux (VII 199 $=$ fr. 276(265) K.-A.) : %
%\citComm{ἐν δ᾿ Ἀριστοφάνης Δαναΐσι συρμαιοπῶλαι καὶ ἐγκριδοπῶλαι.\\
%Aristophane dans ses \italique{Danaïdes} [emploie les mots de] "vendeurs de syrméa" (syrméa $=$ purgatif, voir le commentaire au vers 857) et de "vendeurs d'encrides".}%
%\eff{Une fois tous ces éléments recoupés, il devient aisé de saisir l'idée du premier fragment : le "μήτʼ αρμα" inintelligible est en fait un κάθαρμα, un \italique{purgatif}. Vous me direz qu'il y a loin de μήτʼ αρμα à κάθαρμα ? détrompez-vous : on trouve de nombreuses abréviations dans les manuscrits ; les prépositions sont bien-sûr parmi les plus abrégées et κατά $=$ κ`  quand μετά $=$ μ`. Il n'y a donc, à la base, qu'un kappa mélu en mu\marqueur{5}.
%Supposer\marqueur{6} que dans le manuscrit copié les kappa étaient semblables aux mu permet de corriger l'autre aberration du fragment : καὶ ποιητήν est une mauvaise transcription de μαιοπώλην du mot συρμαιοπώλην (le συρ ayant sauté par haplographie, comme on dit, précédé qu'il est du σιν de Δαναίσιν).\\
%On obtient :%
%\citComm{Ἀριστοφάνης δʼ ἐν Δαναίσιν· συρμαιοπώλην φησιν αὐτῶν εἶναι ἐν τούτοις κάθαρμα εἶναι (ἐγκριδοπώλην).}%
%À la lettre, ce n'est toujours pas grec. Mais le sens général se laisse dégager : il s'agit d'une querelle entre les différents vendeurs ; celui d'encrides (ἐγκριδοπώλης) -- à moins que ce ne soit le vendeur de syrméa (συρμαιοπώλης) qui plastronne -- traite le vendeur de syrméa de purgatif (κάθαρμα). (Avec sans doute des jeu de mots sur la polysémie.)}%
}%
\debutNote{570}%
\note{1}{Suidas (π 3192) nous apprend que c'est une sorte de gâteau fait à partir de miel bouilli et de sésame : Πυραμοῦς: εἶδος πλακοῦντος ἐκ μέλιτος ἑφθοῦ καὶ πυροῦ πεφρυγμένου, ὡς σησαμοῦς ὁ διὰ σησαμῆς. ταῦτα δὲ ἐτίθεσαν ἆθλα τοῖς διαγρυπνηταῖς ἐν τοῖς συμποσίοις, καὶ ὁ διαγρυπνήσας μέχρι τῆς ἔω ταῦτα ἐλάμβανεν ἔπαθλον. ἡμέτερος ὁ πυραμοῦς.}%
\note{2}{On en ramasse à la pelle de ces exemples de déjections consécutives à une frayeur ; Taillardat (§309) en cite bon nombre : \italique{Assemblée des femmes}, v. 1060 ; \italique{Acharniens}, v. 350 ; et d'autres, mais la liste n'est toujours pas exhaustive : on pourrait encore la compléter avec les \italique{Guêpes}, v. 627 ; \italique{Oiseaux}, v. 66 ; etc. etc. Signalons la même promesse de relâchement des intestins après correction dans \italique{Lysistrata}, v. 440 : \grec{ἐπιχεσεῖ πατούμενος}.}%
\note{3}{Hypothèse, certes corroborée par les textes d'Aristophane que nous allons voir mais point testée sur mon propre appareil digestif. Est-ce digne d'un disciple de Willems, ce grand amateur de \italique{Realia} ?}%
\note{4}{En passant, Euripide use du même jeu de mots scatologique : au vers 300 de son \italique{Ion} : σὺν ἀνδρί $[$Créuse dit à Ion qu'elle est venue avec son mari, mais... ] σηκοὺς δ᾽ ἐκστρέφει Τροφωνίου... que le mari, Xuthus, était d'ailleurs véritablement allé consulter. Le tour était proverbial et Érasme l'a recensé (\italique{Adages} n° 677)}%
\note{5}{Pollux (VII 199 $=$ fr. 276(265) Kassel-Austin). Cf. aussi le passage -- corrompu -- d'Athénée (XIV, 645e) où il est question de l'ἐγκριδοπώλης dans les \italique{Danaïdes}.}%
\note{6}{Athénée (XIV, 645e) : ΕΓΚΡΙΔΕΣ. πεμμάτιον ἑψόμενον ἐν ἐλαίῳ καὶ μετὰ τοῦτο μελιτούμενον.}%
%\note{5}{Pour tout ce qui touche à la paléographie, je vais chercher les informations principalement dans le livre d'Édouard Thomson, \italique{Handbook of Greek and Latin Palaeography} (1906). Les abréviations et autres contractions en sont le septième chapitre (pp. 86 sqq.). On trouvera un utile tableau de l'évolution des cursives grecques à la page 149. Les abréviations comme κατά $=$ κ` ou μετά $=$ μ` ainsi que des kappa très semblables aux mu laissent à penser que le manuscrit mal interprété était contemporain d'Athénée : \nieme{II} siècle ap. J.-C.}%
%\note{6}{Supposition gratuite. Il faudrait qu'un fin connaisseur des \italique{Deipnosophistes} confirme la présence d'erreurs de ce genre ou bien infirme l'hypothèse. Pour une introduction à la tradition manuscrite d'Athénée, l'on se reportera avec intérêt à l'article de l'un de ces fins connaisseurs : Benoît Louyest, "L'épitomé du \italique{Banquet des sophistes} d'Athénée", in \italique{Rursus} 8 (2012).}%
\\\passeLigne}%
\end{groupe}%
%% v. 571-573
\begin{groupe}%
\personnageGrec{Κορυφαῖος}%
\tetramIambCat{παύσασθε λοιδορούμεναι· καὶ γὰρ γυνή τις ἡμῖν\apparat{571}{(Χορ.) R : \lecon{(Κορ.)} von Velsen \espApp παύσασθαι R : \lecon{παύσασθε} Μυ2}\\
ἐσπουδακυῖα προστρέχει. πρὶν οὖν ὁμοῦ γενέσθαι,\testimonia{571-2}{(πρὶν οὖν -- πυθώμεθ᾽) Suidas ο 292 (κοσμίως omisso)}\\
σιγᾶθ᾽, ἵν᾽ αὐτῆς κοσμίως πυθώμεθ᾽ ἅττα λέξει.\\\passeLigne}%
\end{groupe}%
%% v. 574-581
\begin{groupe}%
\personnageGrec{Κλεισθένης}%
\typeVers{trim. iamb.}\trimIamb{φίλαι γυναῖκες ξυγγενεῖς τοὐμοῦ τρόπου·\testimonia{574}{(ξυγγενεῖς τοὐμοῦ τρόπου) Suidas ξ 83}\\
ὅτι μὲν φίλος εἴμ᾽ ὑμῖν, ἐπίδηλος ταῖς γνάθοις·\\
γυναικομανῶ γὰρ προξενῶ θ᾽ ὑμῶν ἀεί.\testimonia{576}{Suidas  π 2541 προξενῶ θ' ὑμᾶς ἀεί}\apparat{576}{\lecon{ὑμῶν} R : ὑμᾶς Suid. π 2541}%
\commentaire{575-6}{}{Sur la conservation de ces deux vers que d'aucuns auraient voulu biffer, voir la note de Willems n° 17.% ; cette note est surtout instructive pour ses remarques de bas de page à propos de mesquins vols d'émendations que certains parmi les plus connus hellénistes se permettaient, pillant sans vergogne leurs confrères ne jouissant pas de la même notoriété. Ô le joli monde de ces esprits supérieurs versés dans la recherche.
}\\
καὶ νῦν ἀκούσας πρᾶγμα περὶ ὑμῶν μέγα\\
ὀλίγῳ τι πρότερον κατ᾽ ἀγορὰν λαλούμενον,
\commentaire{578}{ὀλίγῳ τι πρότερον κατ᾽ ἀγορὰν λαλούμενον,}{Clisthène rapporte les cancans de l'Agora, qui font grand cas de ce parent d'Euripide introduit chez les femmes pou plaider sa cause. Ce n'est certainement pas Euripide qui est à l'origine de la fuite et ce pour des raisons évidentes : \ \ \ 1° Si Aristophane, l'avait imaginé ainsi, il nous l'aurait fait comprendre directement ou par quelque allusion \ \ \ 2° Euripide n'a aucun intérêt à voir son complice être démasqué, qui ne pourra donc plus plaider sa cause \ \ \ 3° Euripide a beaucoup de défauts, mais pas celui d'être ingrat au point d'envoyer à une mort certaine son dévoué parent. \ \ \ Savoir quel domestique d'Agathon, quel voisin, ayant aperçu Mnésiloque attifé en femme, sera allé colporter la nouvelle sur l'Agora... voilà bien des détails sans importance sur lesquels Aristophane n'a pas daigné s'arrêter ; nous aurions fait de même si les professeurs parmi les plus admis ne s'étaient lancés dans ce genre d'extravagantes et vaines conjectures.}\\
ἥκω φράσων τοῦτ᾽ ἀγγελῶν θ᾽ ὑμῖν, ἵνα
\commentaire{579-87}{ἥκω φράσων...}{Passées les retrouvailles, Clisthène prend un ton paratragique -- l'heure est grave ! Noter à ce propos les redoublements\marqueur{1} : φράσων/ἀγγελῶν, σκοπῆτε/τηρῆτε, δεινὸν/μέγα, βουλεύοισθε/μέλλοιτε.}%
\debutNote{579-87}%
\note{1}{Redoublements propres au style eschyléen et raillés par l'Euripide des \italique{Grenouilles} (vv. 1153 sqq.).}\\
σκοπῆτε καὶ τηρῆτε μὴ καὶ προσπέσῃ\apparat{580}{σκοπεἴτε καὶ τηρεἴτε R : \lecon{σκοπῆτε καὶ τηρῆτε} Biset \espApp καὶ μὴ R : \lecon{μὴ καὶ} Porson (\italique{Aristophanica} p. 216) \espApp \textbf{581} ἡμῖν R : \lecon{ὑμῖν} \Mon \espApp ἀφράκτοις R : \lecon{ἀφάρκτοις} Dindorf (1830)}\\
ὑμῖν ἀφάρκτοις πρᾶγμα δεινὸν καὶ μέγα.\\\passeLigne}%
\end{groupe}%
%% v. 582-583
\begin{groupe}%
\personnageGrec{Κορυφαῖος}%
\trimIamb{τί δ᾽ ἔστιν ὦ παῖ; παῖδα γάρ σ᾽ εἰκὸς καλεῖν,\\
ἕως ἂν οὕτως τὰς γνάθους ψιλὰς ἔχῃς.%
\commentaire{582-3}{τί δ᾽ ἔστιν ὦ παῖ; παῖδα γάρ σ᾽ εἰκὸς καλεῖν, /	ἕως ἂν οὕτως τὰς γνάθους ψιλὰς ἔχῃς.}{Pique de la coryphée analogue à celle du chœur des \italique{Guêpes}, aux vers 1297-1298 : %
\citComm{Χορός – τί δ᾽ ἔστιν ὦ παῖ; παῖδα γάρ, κἂν ᾖ γέρων,\\
καλεῖν δίκαιον ὅστις ἂν πληγὰς λάβῃ.}%
}%
\\\passeLigne}%
\end{groupe}%
%% v. 584-585
\begin{groupe}%
\personnageGrec{Κλεισθένης}%
\trimIamb{Εὐριπίδην φάσ᾽ ἄνδρα κηδεστήν τινα\apparat{584}{ἔφας᾽ R : \lecon{φάσ᾽} Zanetti}\\
αὑτοῦ γέροντα δεῦρ᾽ ἀναπέμψαι τήμερον.\apparat{585}{αὐτου R : \lecon{αὑτοῦ} Biset}%
\commentaire{585}{ἀναπέμψαι.}{Le suffixe `ἀνα' s'explique par le fait que le Thesmophorion se trouvait probablement sur la colline de la Pnyx. (Cf. le vers 658.)}\\\passeLigne}%
\end{groupe}%
%% v. 586
\begin{groupe}%
\personnageGrec{Κορυφαῖος}%
\trimIamb{πρὸς ποῖον ἔργον ἢ τίνος γνώμης χάριν;%\commentaire{586}{πρὸς ποῖον ἔργον ἢ τίνος γνώμης χάριν;}{Vers paratragique.}
\\\passeLigne}%
\end{groupe}%
%% v. 587-588
\begin{groupe}%
\personnageGrec{Κλεισθένης}%
\trimIamb{ἵν᾽ ἅττα βουλεύοισθε καὶ μέλλοιτε δρᾶν,\\
\lednopb%
ἐκεῖνος εἴη τῶν λόγων κατάσκοπος.\\\passeLigne}%
\end{groupe}%
%% v. 589
\begin{groupe}%
\personnageGrec{Κορυφαῖος}%
\trimIamb{καὶ πῶς λέληθεν ἐν γυναιξὶν ὢν ἀνήρ;\\\passeLigne}%
\end{groupe}%
%% v. 590-591
\begin{groupe}%
\personnageGrec{Κλεισθένης}%
\trimIamb{ἄφευσεν αὐτὸν κἀπέτιλ᾽ Εὐριπίδης\apparat{590}{\lecon{ἄφευσεν} R : ἀφηῦσεν Meineke \espApp καπέτιλλ᾽ R : \lecon{κἀπέτιλ᾽} Bekker}\\
καὶ τἄλλ᾽ ἅπανθ᾽ ὥσπερ γυναῖκ᾽ ἐσκεύασεν.\\\passeLigne}%
\end{groupe}%
%% v. 592-594
\begin{groupe}%
\personnageGrec{Μνησίλοχος}%
\trimIamb{πείθεσθε τούτῳ ταῦτα; τίς δ᾽ οὕτως ἀνὴρ\\
ἠλίθιος ὅστις τιλλόμενος ἠνείχετο;
\commentaire{593}{ἠλίθιος ὅστις τιλλόμενος ἠνείχετο;}{Willems, chevalier servant du \italique{Ravennas}, défend encore une fois la leçon du manuscrit dans sa \nieme{18} note, à retrouver à la fin.}\\
οὐκ οἴομαι 'γωγ᾽ ὦ πολυτιμήτω θεώ.\apparat{594}{οἴομ᾽ ἔγωγ᾽ R : \lecon{οἴομαι 'γωγ᾽} anonymus Parisinus}\\\passeLigne}%
\end{groupe}%
%% v. 595-596
\begin{groupe}%
\personnageGrec{Κλεισθένης}%
\trimIamb{ληρεῖς· ἐγὼ γὰρ οὐκ ἂν ἦλθον ἀγγελῶν,\\
εἰ μὴ ᾽πεπύσμην ταῦτα τῶν σάφ᾽ εἰδότων.\apparat{596}{πεπύσμην R, \PsiUn : \lecon{᾽πεπύσμην} Faber \espApp ταυτὶ R : \lecon{ταῦτα} Bentley}\\\passeLigne}%
\end{groupe}%
%% v. 597-602
\begin{groupe}%
\personnageGrec{Κορυφαῖος}%
\trimIamb{τὸ πρᾶγμα τουτὶ δεινὸν εἰσαγγέλλεται.\\
ἀλλ᾽ ὦ γυναῖκες οὐκ ἐλινύειν ἐχρῆν,\testimonia{598}{(ἐλινύειν) Hesychius ε 2095 ; cf. Moeris ε 61 (ἐλινύων $=$ ἀναπαυόμενος) $\sim$ Hesychius ε 2097}\apparat{598}{ἐλιννύειν R : \lecon{ἐλινύειν} Bourdin, cf. Moeris \italique{Gram.} ε 61, Hesychius ε 2097}\\
ἀλλὰ σκοπεῖν τὸν ἄνδρα καὶ ζητεῖν ὅπου\\
λέληθεν ἡμᾶς κρυπτὸς ἐγκαθήμενος.\apparat{600}{ὑμᾶς R : \lecon{ἡμᾶς} \Mon}\\
καὶ σὺ ξυνέξευρ᾽ αὐτόν, ὡς ἂν τὴν χάριν\testimonia{601-2}{Suidas π 2540}\apparat{601}{συνέξευρ᾽ R, Suid. π 2540 codd. AG : \lecon{ξυνέξευρ᾽} Suid. codd. FVM \ligneApp{602} \lecon{ταύτην} R : αὐτὴν Suid.}\\
ταύτην τε κἀκείνην ἔχῃς ὦ πρόξενε.%
\commentaire{602}{ὦ πρόξενε.}{Les proxènes ou hôtes publics, agents officiels dont les fonctions consistaient principalement à recevoir et à héberger les délégués de l'État dont ils exerçaient la proxénie. (Note de Willems.)}\\\passeLigne}%
\end{groupe}%
%% v. 603a
\begin{groupe}%
\personnageGrec{Κλεισθένης}%
\trimIamb{φέρ᾽ ἴδω· τίς ἡ πρώτη σύ;\\\passeLigne}%
\end{groupe}%
%% v. 603b
\begin{groupe}%
\personnageGrec{Μνησίλοχος}%
\trimIamb{\ligne{603}\decalage{φέρ᾽ ἴδω· τίς ἡ πρώτη σύ;}ποῖ τις τρέψεται;\\\passeLigne}%
\end{groupe}%
%% v. 604a
\begin{groupe}%
\ledpb%
\personnageGrec{Κλεισθένης}%
\trimIamb{ζητητέαι γάρ ἐστε.%\apparat{604}{\lecon{ζητητέαι γάρ ἐστε} R : ζητητέον τἄρ᾽ ἐστί Maas \espApp 604 ante 603 transp. Maas, post 607 collocare malui}
\\\passeLigne}%
\end{groupe}%
%% v. 604b
\begin{groupe}%
\personnageGrec{Μνησίλοχος}%
\trimIamb{\ligne{604}\decalage{ζητητέαι γάρ ἐστε.}κακοδαίμων ἐγώ.\\\passeLigne}%
\end{groupe}%
%% v. 605
\begin{groupe}%
\personnageGrec{Μίκα}%
\trimIamb{ἔμ᾽ ἥτις εἴμ᾽ ἤρου; Κλεωνύμου γυνή.\apparat{605}{εἰτις ἤρου; R : \lecon{ἥτις} (Zanetti) \lecon{εἴμ᾽ ἤρου;} Bentley : ἥτις ἤρου; τοῦ Biset : alii alia}%
\commentaire{605}{Κλεωνύμου γυνή.}{L'évocation de Cléonyme est pour nous l'occasion de faire un bref résumé des railleries à l'encontre de ce personnage parmi les plus plaisantés de l'Ancienne comédie. Cet Athénien nous est autrement connu par trois décrets qu'il proposa à la \grec{βουλή}\marqueur{1} l'an 426/5 ainsi qu'une proposition de fixer à 1000 drachmes la récompense contre des informations touchant à l'affaire de la parodie des Mystères d'Éleusis\marqueur{2} (contemporaine de celle des Hermocopides, l'an 415). Cléonyme est victime d'un tir de barrage d'Aristophane portant à la fois :\\
— au moral : sa lâcheté (il lui était arrivé un jour de jeter son bouclier pour mieux fuir. C'est surtout là-dessus que le poète ne tarit point : \italique{Cavaliers} v. 1372, \italique{Nuées} v. 353, \italique{Guêpes} v. 19, 592, \italique{Paix} v. 446, 672, 1186 (non nommé), 1275, 1295, \italique{Oiseaux} v. 290, 1475, 1481) ; sa duplicité (\italique{Acharniens} v. 88, il est traité de parjure v. 399 des \italique{Nuées}, de flatteur au v. 592 des \italique{Guêpes} (il est appelé Κολακώνυμος), de sycophante au v. 1479 des \italique{Oiseaux}) ; ses manières brusques de grand dadais (\italique{Acharniens} v. 844) ; ses mœurs peu ragoûtantes (accusé d'être un \grec{καταπύγων} au v. 675 des \italique{Nuées}) ;\\
— au physique : une laideur repoussante (v. 822 des \italique{Guêpes}, ) ; son androgynie (Socrate, dans une intention de parler exactement, emploie le féminin \grec{Κλεωνύμη} au v. 680 des \italique{Nuées}) ; une taille excessive (\italique{Acharniens} v. 88, \italique{Oiseaux}, v. 1476) ; sa gloutonnerie (\italique{Cavaliers} v. 1290, \italique{Oiseaux} v. 289).\\
Jouant sur les contrastes, Aristophane nous présente une femme de Cléonyme \italique{petite} de nom (c'est le sens de \grec{Μίκα}), virile et hardie dans les débats.\\
%Cette note sur l'occasion de tordre le cou à l'une des traditions philologiques les plus tenaces et qui laissait déjà en son temps Willems devant des abîmes de perplexité : pourquoi sont-ce les mêmes éminents professeurs qui se lancent dans des conjectures les plus extravagantes lorsqu'ils sont confrontés à d'obscurs passages allusifs et qui, face aux passages les plus clairs et dont le sens 
Revenons sur le grief le plus reproché à Cléonyme : avoir jeté son bouclier pour mieux prendre ses jambes à son cou, sa \grec{ῥίψασπις}. Il n'y a rien à y redire : le crime était théoriquement puni de mort\marqueur{3} et l'on comprend aisément que le poète ait exposé à l'opprobre publique le couard qui avait encore l'outrecuidance de venir palabrer après son forfait. Exégèse par trop simple pour nos éminents professeurs – dont le délassement consiste à extravaguer sur les passages les plus obscurs et remettre systématiquement en doute les passages les plus clairs\marqueur{4}. On ne la fait pas à ces \italique{viri docti}, comme disait Fritzsche : le dernier en date\marqueur{5} argue de ce que Cléonyme serait mort – condamné à mort ! – bien avant les dernière traces archéologiques de sa vie démocratique s'il avait effectivement commis ce qu'Aristophane lui reproche... avec le même raisonnement, les historiens des prochains millénaires déduiront que l'\italique{Élysée}, \italique{Matignon} et le \italique{palais Bourbon} ne pouvaient désigner que des prisons au vu des canailles hors-la-loi qui les fréquentaient ; prisons flanquées, semblera-t-il, d'une maison de santé pour ceux dont la panse menacait de crever à tout instant : le \italique{palais du Luxembourg}\marqueur{6}. Mais il n'y a pas besoin d'extrapoler : le même raisonnement intenable entraînait déjà de refuser le vers 1166 des \italique{Acharniens}\marqueur{7} : comment un bandit (un \italique{Oreste}, nom générique ; un \italique{Cartouche} ou \italique{Mandrin} chez nous) aurait-il pu sévir la nuit à Athènes alors que le banditisme était aussi puni de mort ?..
Certes, « \anglais{in interpreting the personal jokes in Aristophanes one must always be wary of the too serious or too literal explanation of the text} »\marqueur{8}, mais l'\italique{on ne détruit que ce que l'on remplace}, et pour remplacer une explication qui coule de source comme celle de nos passages, il vaut mieux présenter de solides arguments. Or M. Storey ne me fera pas accroire sans preuve\marqueur{9} que le motif pour lequel Aristophane poursuivit inlassablement de sa verve pendant des années Cléonyme fût, non pas une désertion devant l'ennemi doublée d'un abandon d'armes, mais une banale brigue pour échapper au service militaire\marqueur{10}.
}%
\debutNote{605}%
\note{1}{\italique{Inscriptiones Graecae}, I\expo{3}, 61.34, 68.5, 69.3-4.}%
\note{2}{Andocide, I-\italique{\grec{Περὶ τῶν μυστηρίων}}, 27. La démagogie de Cléonyme se voit elle-même débordée par celle de Pisandre qui renchérit à 10000 drachmes !}%
\note{3}{Andocide, I-\italique{\grec{Περὶ τῶν μυστηρίων}}, 74.}%
\note{4}{Willems, tome I, p. 108, s'en plaignait déjà en son temps.}%
\note{5}{I. C. Storey, «~The `Blameless Shield' of Kleonymus~», in \italique{Rheinisches Museum für Philologie}, n°132 (1989), p. 247-261.}%
\note{6}{«~Drôles de noms pour des prisons~!~» objecteront les plus prudents ; ce à quoi les plus érudits répondront que l'\italique{Élysée} fait référence ironiquement au lieu que devait rejoindre son unique résident après avoir reçu les douze balles qui devaient immanquablement être promises à ce ``héros jupitérien'', multirécidiviste de la haute-trahison ! « Mate ou gnons ! », tel devait être le mot d'ordre à ``Matignon'' %-- noter la eff{demander à Alexandre + exemples} -- 
qui n'accueillit que de serviles et soumises serpillières. Les dénominations de ``palais'' sont à prendre dans leur probable acception technique judiciaire et n'étaient que la suite logique des \italique{palais} de justice ; ``Bourbeux'' devait être le nom d'origine du premier quand on y constate l'entassement de la valetaille parlementaire et l'ambiance de basse-cour qui y régnait~; le second porte évidemment le nom de son plus célèbre prisonnier~: Rosa Luxembourg (et les jardins du même nom – attestés parfois au singulier – étaient probablement la cour dans laquelle les prisonniers allaient se délassant quand la folie régnant à l'intérieur se faisait trop oppressante).}%
\note{7}{voir Willems, t. I p. 106 sq.}%
\note{8}{«~The `Blameless Shield' of Kleonymus~», p. 248.}%
\note{9}{Car je tiens pour nulle et non avenue sa citation du passage des \italique{Cavaliers} (v. 1369-1372) qui \italique{prouve} simplement que M. Storey n'a pas saisi l'image et le jeu de mots d'Aristophane sur \grec{κατὰ σπουδὰς} (mis à dessein pour faire entendre κατ᾿ ἀσπίδας en lieu et place de l'expression figée `κατὰ σπουδήν' qu'on trouve chez Thucydide, Xénophon, Élien, Plutarque, etc.). Les autres passages sont encore moins probants : \italique{Paix} v. 1172-1187 et \italique{Guêpes} v. 1112-1121. Au reste, si M. Storey convient -- comme il ne fait pas de doute et comme le scholiaste l'avait déjà remarqué -- que les vers 827-829 des \italique{Thesmophories} visent Cléonyme, je lui saurais gré d'expliquer le tour `ἀπὸ τῶν ὤμων' à l'aune de sa brillante exégèse.}%
\note{10}{Si banale qu'on n'en verrait pas le bout de la liste si l'on s'avisait de recenser ceux qui y ont recouru parmi les modernes brigands qui sont au pouvoir. Citons en seulement l'exemple le plus caricatural : M. Ciotti le-bien-nommé, alors qu'il n'était encore que vil attaché-parlementaire fit jouer ses relations pour échapper au service militaire... maintenant que son derrière est confortablement installé au palais bourbeux et ne risque plus de se salir dans la fange de quelque camp d'entraînement, il n'a de cesse de haranguer le veau français pour un rétablissement dudit service militaire.% \italique{Ubi quae fuerant vitia, mores sunt.}%
}%
\\\passeLigne}%
\end{groupe}%
%% v. 606
\begin{groupe}%
\personnageGrec{Κλεισθένης}%
\trimIamb{γιγνώσκεθ᾽ ὑμεῖς ἥτις ἔσθ᾽ ἡδὶ γυνή;\apparat{606}{ἥδε R : ἥδ᾽ ἡ \Mon : \lecon{ἡδὶ} Grynaeus}\\\passeLigne}%
\end{groupe}%
%% v. 607
\begin{groupe}%
\personnageGrec{Κορυφαῖος}%
\trimIamb{γιγνώσκομεν δῆτ᾽. ἀλλὰ τὰς ἄλλας ἄθρει.\\\passeLigne}%
\end{groupe}%
%% v. 608-609a
\begin{groupe}%
\personnageGrec{Κλεισθένης}%
\trimIamb{ἡδὶ δὲ δὴ τίς ἐστιν ἡ τὸ παιδίον\\
ἔχουσα;\\\passeLigne}%
\end{groupe}%
%% v. 609b
\begin{groupe}%
\personnageGrec{Μίκα}%
\trimIamb{\ligne{609}\decalage{ἔχουσα;}τίτθη νὴ Δί᾽ ἐμή.\apparat{609}{τίτθην η δί᾽ R : \lecon{τίτθη νὴ Δί᾽} Ellebodius}%
\commentaire{609}{τίτθη.}{Remarquons en passant qu'une nourrice à Athènes n'était pas nécessairement une esclave. Cf. Démosth., LVII, 35. (Note de Willems.)}\\\passeLigne}%
\end{groupe}%
%% v. 609c
\begin{groupe}%
\personnageGrec{Μνησίλοχος}%
\trimIamb{\ligne{609}\decalage{ἔχουσα;τίτθη νὴ Δί᾽ ἐμή.}διοίχομαι.\\\passeLigne}%
\end{groupe}%
%% v. 610
\begin{groupe}%
\personnageGrec{Κλεισθένης}%
\trimIamb{αὕτη σὺ ποῖ στρέφει; μέν᾽ αὐτοῦ. τί τὸ κακόν;%\apparat{610}{τί τὸ κακόν; continuat Clistheni R : Parenti tribuit (cum versu sequenti) \Rpc}
\\\passeLigne}%
\end{groupe}%
%% v. 611a
\begin{groupe}%
\personnageGrec{Μνησίλοχος}%
\trimIamb{ἔασον οὐρῆσαί μ᾽· ἀναίσχυντός τις εἶ.\apparat{611}{εἶ R : \lecon{τις εἶ} Küster}\\\passeLigne}%
\end{groupe}%
%% v. 611b-612
\begin{groupe}%
\personnageGrec{Κλεισθένης}%
\trimIamb{%\ligne{611}\decalage{ἔασον οὐρῆσαί μ᾽.}\\
%\commentaire{611}{ἀναίσχυντός εἶ.}{J'attribue à Clisthène la réplique (qui traitera encore Mnésiloque au vers 638 d'\grec{ἀναίσχυντός}).}%
σὺ δ᾽ οὖν ποίει τοῦτ᾽. ἀναμενῶ γὰρ ἐνθάδε.\apparat{612}{ἀναμένω R : \lecon{ἀναμενῶ} Grynaeus}\\\passeLigne}%
\end{groupe}%
%% v. 613-614
\begin{groupe}%
\personnageGrec{Κορυφαῖος}%
\trimIamb{ἀνάμενε δῆτα καὶ σκόπει γ᾽ αὐτὴν σφόδρα·\\
μόνην γὰρ αὐτὴν ὦνερ οὐ γιγνώσκομεν.\\\passeLigne}%
\end{groupe}%
%% v. 615a
\begin{groupe}%
\personnageGrec{Κλεισθένης}%
\trimIamb{πολύν γε χρόνον οὐρεῖς σύ.\apparat{615}{πολύ R : \lecon{πολύν} Gelenius}\\\passeLigne}%
\end{groupe}%
%% v. 615b-616
\begin{groupe}%
\personnageGrec{Μνησίλοχος}%
\trimIamb{\ligne{615}\decalage{πολύν γε χρόνον οὐρεῖς σύ.}νὴ Δί᾽ ὦ μέλε·\\
στραγγουριῶ γάρ· ἐχθὲς ἔφαγον κάρδαμα.\testimonia{616}{Suidas σ 1163 (χθὲς γὰρ pro γάρ ἐχθὲς)}%
\commentaire{616}{κάρδαμον}{%Le κάρδαμον est identifié au \italique{cresson alénois} par les différents dictionnaires de grec\marqueur{1}. Alphonse Willems, porté aux \italique{realia}, avait déjà formulé en son temps une objection à cette identification : Mnésiloque justifie le long temps qu'il met à uriner au motif d'avoir mangé du cardame la veille. Hippocrate (\italique{Du Régime}, II, 54) confirme que le cardame provoque la dysurie. Or, de ce que j'ai pu savoir grâce à l'internet -- le temple de la culture pour un néophyte en herbes comme moi --, le cresson, loin d'être dysurique, est au contraire diurétique !
%Une anecdote de Xénophon confirme les dires d'Hippocrate : chez les Perses, où il passait pour indécent de cracher, se moucher ou s'écarter pour faire de l'eau (\italique{Cyropédie}, I, 2, 16), on ne donnait aux enfants à leur repas du midi que du pain, avec du cardame pour tout assaisonnement (\italique{Cyrop}., I, 2, 8 et 11). La consommation de cresson paraît donc impropre aux mœurs perses. 
%Il me semble que le cardame grec n'était peut-être pas un cresson mais plutôt une crucifère cousine : une variété de moutarde ? Quelques indices me le font penser. 
%D'abord, le cardame est très souvent rapproché d'autres noms de moutarde : ainsi dans le passage déjà cité d'Hippocrate, le κάρδαμον est qualifié de θερμόν, tout comme la plante suivante : νᾶπυ. 
%Autre exemple : Athénée (IX, 366d) cite ce vers de Nicandre : Κάρδαμον, ἄρρινόν τε, μελάμφυλλόν τε σίνηπυ. Plus importante encore est cette glose d'Hésychius : \citComm{ἐρύσιμον· κάρδαμον ἄγριον}
%L'érysimon est identifié à notre \italique{sisymbre à cornes} (\latin{Sisymbrium polyceratium}) ; ce qui se recoupe avec mes cyber-recherches qui m'ont appris que le sisymbre était un succédané sauvage de la moutarde.
%
%\eff{Érasme 754}
Les botanistes s'accordent, je crois, à reconnaître dans le κάρδαμον, non le \italique{cresson}, comme on l'a cru jusqu'en ces derniers temps, mais une autre crucifère, connue des savants sous le nom de \latin{erucaria aleppica}. Suivant Fraas, \italique{Synopsis florae classicae}, les Grecs lui donnent encore aujourd'hui le même nom. Hippocrate dit positivement que le \grec{κάρδαμον} provoque la dysurie (\italique{Du Régime}, II, 54). On cite à ce sujet une singulière anecdote, recueillie par Xénophon. Chez les Perses, où il passait pour indécent de cracher, se moucher ou s'écarter pour faire de l'eau, on ne donnait aux enfants à leur repas du midi que du pain, avec du cardame pour tout assaisonnement. (\italique{Cyrop.}, I, 2, 8, 11 et 16) Mettez cela d'accord avec les traités de thérapeutique où le cresson est donné pour un sialagogue. (Note de Willems\marqueur{1}.)}%
\debutNote{616}%
\note{1}{Les dictionnaires  \LSJ , Bailly et Pape désignent le \grec{κάρδαμον} respectivement comme « \anglais{nose-smart}, \latin{Lepidium sativum} » (\italique{cresson alénois}), « cresson » et « \allemand{eine Art Kresse mit bitterem Kraut} ».%
Pour aller dans le sens de Willems, ajoutons que le cresson, outre ses vertus sialagogiques, est loin d'être dysurique : il a aussi des vertus au contraire diurétiques. Cf. le \italique{Dictionnaire des plantes médicinales et vénéneuses de France} de Paul-Victor Fournier, p. 731. Mais la raison de sa consommation est peut-être à chercher ailleurs : cette plante était réputée donner du courage.}%
\\\passeLigne}%
\end{groupe}%
%% v. 617
\begin{groupe}%
\personnageGrec{Κλεισθένης}%
\trimIamb{τί καρδαμίζεις; οὐ βαδιεῖ δεῦρ᾽ ὡς ἐμέ;\\\passeLigne}%
\end{groupe}%
%% v. 618a
\begin{groupe}%
\personnageGrec{Μνησίλοχος}%
\trimIamb{τί δῆτά μ᾽ ἕλκεις ἀσθενοῦσαν;\\\passeLigne}%
\end{groupe}%
%% v. 618b-619a
\begin{groupe}%
\personnageGrec{Κλεισθένης}%
\trimIamb{\ligne{618}\decalage{τί δῆτά μ᾽ ἕλκεις ἀσθενοῦσαν;}εἰπέ μοι,\\
τίς ἔστ᾽ ἀνήρ σοι;\apparat{619}{ἔσθ᾽ R : \lecon{ἔστ᾽} \Mon}\\\passeLigne}%
\end{groupe}%
%% v. 619b-620
\begin{groupe}%
\personnageGrec{Μνησίλοχος}%
\trimIamb{\ligne{619}\decalage{τίς ἔστ᾽ ἀνήρ σοι;}τὸν ἐμὸν ἄνδρα πυνθάνει;\\
τὸν δεῖνα γιγνώσκεις, τὸν ἐκ Κοθωκιδῶν.\testimonia{620}{(Κοθωκιδῶν) Photius κ 859 Κοθωκίδαι (ap Scholium) $\sim$ Hesychius  κ 3218}%\apparat{620}{add. signum interrogationis Canini}%
\commentaire{619}{τὸν ἐκ Κοθωκιδῶν;}{« Du $[$dème$]$ des Cothôcides ». Ce dème faisait partie de la tribu des \grec{Οἰνηίς} au nord-ouest de la ville\marqueur{1}. Y a-t-il une raison particulière pour que Mnésiloque évoque ce dème ? certainement... mais cette raison nous reste obscure et aucune explication satisfaisante n'a jusqu'ici été avancée.}%
\debutNote{619}%
\note{1}{L'on pourra consulter le travail de John S. Traill, \italique{The Political Organization of Attica}, \italique{Hesperia} supplement XIV (1975). Voir notamment les pages 49, 62 et 68 ainsi que les cartes d'Athènes à la fin.}\\\passeLigne}%
\end{groupe}%
%% v. 621a
\begin{groupe}%
\ledpb%
\personnageGrec{Κλεισθένης}%
\trimIamb{τὸν δεῖνα; ποῖον;\\\passeLigne}%
\end{groupe}%
%% v. 621b-622a
\begin{groupe}%
\personnageGrec{Μνησίλοχος}%
\trimIamb{\ligne{621}\decalage{τὸν δεῖνα; ποῖον;}ἔσθ᾽ ὁ δεῖν᾽, ὃς καί ποτε\apparat{621}{(Κλεισθ.) continuat R : \lecon{(Μνη.)} \Mon}\\
τὸν δεῖνα τὸν τοῦ δεῖνα—\\\passeLigne}%
\end{groupe}%
%% v. 622b-623a
\begin{groupe}%
\personnageGrec{Κλεισθένης}%
\trimIamb{\ligne{622}\decalage{τὸν δεῖνα τὸν τοῦ δεῖνα—}ληρεῖν μοι δοκεῖς.\\
ἀνῆλθες ἤδη δεῦρο πρότερον;\\\passeLigne}%
\end{groupe}%
%% v. 623b-624a
\begin{groupe}%
\personnageGrec{Μνησίλοχος}%
\trimIamb{\ligne{623}\decalage{ἀνῆλθες ἤδη δεῦρο πρότερον;}νὴ Δία\\
ὁσέτη γε.\apparat{624}{ὅς᾽ ἔτη R : \lecon{ὁσέτη} G. H. Schaefer ad Gregor. Cor. p.65 \espApp σοὐστὶν R : \lecon{σοὐστὶ} Dindorf (σοῦ ᾽στὶ iam Biset)}\\\passeLigne}%
\end{groupe}%
%% v. 624b
\begin{groupe}%
\personnageGrec{Κλεισθένης}%
\trimIamb{\ligne{624}\decalage{ὁσέτη γε.}καὶ τίς σοὐστὶ συσκηνήτρια;%
\commentaire{624}{συσκηνήτρια.}{Il faut croire d'après ce passage et le vers 658 que lors des Thesmophories, des tentes pour deux femmes étaient installées à proximité du Pnyx.}\\\passeLigne}%
\end{groupe}%
%% v. 625a
\begin{groupe}%
\personnageGrec{Μνησίλοχος}%
\trimIamb{ἡ δεῖν᾽ ἔμοιγ᾽.\apparat{625}{ἔμοιγε R : \lecon{ἔμοιγ᾽} edd. pr.}\\\passeLigne}%
\end{groupe}%
%% v. 625b
\begin{groupe}%
\personnageGrec{Κλεισθένης}%
\trimIamb{\ligne{625}\decalage{ἡ δεῖν᾽ ἔμοιγ᾽.}οἴμοι τάλας. οὐδὲν λέγεις.\\\passeLigne}%
\end{groupe}%
%% v. 626-629
\begin{groupe}%
\personnageGrec{Μίκα}%
\trimIamb{ἄπελθ᾽. ἐγὼ γὰρ βασανιῶ ταύτην καλῶς\\
ἐκ τῶν ἱερῶν τῶν πέρυσι· σὺ δ᾽ ἀπόστηθί μοι,\\
ἵνα μὴ 'πακούσῃς ὢν ἀνήρ. σὺ δ᾽ εἰπέ μοι\\
ὅ τι πρῶτον ἡμῖν τῶν ἱερῶν ἐδείκνυτο.\\\passeLigne}%
\end{groupe}%
%% v. 630
\begin{groupe}%
\ledpb%
\personnageGrec{Μνησίλοχος}%
\trimIamb{φέρ᾽ ἴδω, τί μέντοι πρῶτον ἦν; ἐπίνομεν.\testimonia{630-1}{Suidas π 2563 (sed προπίνομεν pro προὐ- cod. V)}\apparat{630}{\lecon{μέντοι} \Rpc in mg., Suid. π 2563 : om. \Rac, \Mon}\\\passeLigne}%
\end{groupe}%
%% v. 631a
\begin{groupe}%
\personnageGrec{Μίκα}%
\trimIamb{τί δὲ μετὰ τοῦτο δεύτερον;\apparat{631}{\lecon{δὲ} R, Suid. π 2563 : δαὶ Bentley \espApp με R : \lecon{μετὰ} Suid.}\\\passeLigne}%
\end{groupe}%
%% v. 631b
\begin{groupe}%
\personnageGrec{Μνησίλοχος}%
\trimIamb{\ligne{631}\decalage{τί δὲ μετὰ τοῦτο δεύτερον;}προὐπίνομεν.\\\passeLigne}%
\end{groupe}%
%% v. 632
\begin{groupe}%
\personnageGrec{Μίκα}%
\trimIamb{ταυτὶ μὲν ἤκουσάς τινος· τί δ᾽ ἦν τρίτον;\apparat{632}{τί δὲ τρίτον R : τί δ᾽ αὖ Bentley : τί δαὶ τρίτον Elmsley : \lecon{τί δ᾽ ἦν τρίτον} Teubner}\\\passeLigne}%
\end{groupe}%
%% v. 633
\begin{groupe}%
\personnageGrec{Μνησίλοχος}%
\trimIamb{σκάφιον Ξένυλλ᾽ ᾔτησεν· οὐ γὰρ ἦν ἀμίς.\testimonia{633}{Pollux X 45 ἐν Πολυίδῳ (sic) Ἀριετοφάνης· σκάφιον Ξένυλλ᾽ ἤτησεν· οὐ γὰρ ἦν ἁμίς \espTest Photius σ 286 σκάφιον}%
\commentaire{633}{σκάφιον Ξένυλλ᾽ ᾔτησεν· οὐ γὰρ ἦν ἀμίς.}{Acculé, c'est à ce vers que Mnésiloque fait son hara-kiri : il lance une dernière boutade en guise de baroud d'honneur ! L'explication communément admise laisse songeur quant à l'humour des commentateurs d'Aristophane : Mnésiloque emploierait selon eux un terme impropre (l'\italique{amis} serait une sorte de pot de chambre à usage exclusivement masculin en raison d'un col étroit... mais sa forme varie et le terme pourrait très bien aussi désigner ce genre d'ustensile pour les femmes\marqueur{1}). Quoi qu'il en soit, c'est un peu \italique{gros} : tout rustre qu'il soit, Mnésiloque ne peut pas commettre une telle bévue ; je dirais même plus~: dans ce domaine, il est certainement tout aussi savant que l'\italique{Athénienne de la rue}. 
Voici une explication bien meilleure : le \grec{σκάφιον} désignait à l'origine une \italique{coupe} (en forme de \italique{barque} selon l'étymologie). Il n'acquiert le sens quasi-exclusif d'\italique{urinoir} qu'en passant au latin \latin{scaphium}. À l'époque d'Aristophane, les deux sens devaient encore co-exister. La première partie du vers « σκάφιον Ξένυλλ᾽ ᾔτησεν » était donc interprétée par Mica favorablement~: Xénylla a demandé une coupe $[$pour faire une libation$]$... la chute est rude et Mnésiloque dans la seconde partie accuse les femmes d'un usage détourné des coupes consacrées qui passe la mesure : cet individu n'a rien des manières d'une femme : après avoir uriné publiquement, voilà qu'il prononce maintenant des \italique{insanités} (οὐδὲν ὑγιὲς v. 636).\\
%Ce n'est d'ailleurs pas autrement que le comprenait le scholiaste :
%\citComm{μὴ εὑρίσχων τί εἴπῃ ὁ κηδεστής φησιν ὅτι ἀντὶ ἀμίδος σκάφιον ᾔτησεν ἣ Ξένυλλα. ἀμὶς δέ ἐστιν οὐροδόχον ἀγγεῖον.\\
%\francais{Ne trouvant pas quoi répliquer, le Parent dit qu'en place d'un amis, c'est un scaphion qu'a demandé Xénylla. Un amis est un récipient pour recueillir l'urine.}}
Voici pour le sens propre -- si l'on peut dire. Mais, comme de coutume, Aristophane joue à différents niveaux : la réplique est aussi un clin d'œil au proverbe grec bien connu\marqueur{2} :%
\citComm{Ἄμας ἀπῄτουν, οἱ δ' ἀπηρνοῦντο σκάφας\\
Je demandais des bêches, ils me refusent des seaux.}%
ἄμη et σκάφη pouvaient en effet être tous deux synonymes de \italique{seau} -- comme dans le proverbe pour évoquer abondance de biens\marqueur{3} : ἄμαις καὶ σκάφαις ἀρύσασθαι --, mais c'était des \italique{houes} (autre sens de ἄμη) que demandait le paysan du proverbe...\\
Apprécions à sa juste valeur l'à-propos de l'allusion au proverbe : ce dernier repose sur la polysémie d'ἄμη comme est polysémique le σκάφιον du début de la réplique de Mnésiloque ; de plus, le proverbe évoque ceux qui fournissent des réponses sans rapport avec la question posée -- ce qui est aussi le cas de notre pauvre Mnésiloque, bien embarrassé.}%
\debutNote{633}%
\note{1}{Pour toutes les remarques de cette note, je renvoie aux articles \latin{amis} et \latin{scaphium} du \italique{Dictionnaire des antiquités grecques et romaines} de Daremberg \& Saglio.}%
\note{2}{Il a été relevé par Suidas en α 1504 et α 1574. Érasme y consacre une entrée dans ses \italique{Adages} (au n°1149) et l'évoque aussi au n°1880. Maintes fois recensé par les parémiographes (par exemple Zénobius I, 83), il était connu de Plutarque (\italique{Du bavardage}, OM, 6, 35, 512e) et le scholiaste de la \italique{Paix} d'Aristophane le cite encore à propos de la chanson de Datis le Lydien aux vers 289-291 et de l'équivoque en grec que suggère la forme barbaristique χαίρομαι -- au passage : Aristophane joue sur le stéréotype du débauché « Lydien ($=$ Datis) à midi », Λυδὸς ἐν μεσημβρίᾳ (rapporté par Athénée, XII, 515d-517a et recensé par Érasme dans ses \italique{Adages}, n° 1594). Enfin, le scholiaste de Théocrite (\italique{Idylles}, IV, 10) cite aussi de ce proverbe. Pour finir, Aristophane y fait de nouveau allusion au fr. 927 Kassel-Austin : ἄγροικός εἰμι· τὴν σκάφην σκάφην λέγω.}%
\note{3}{Abondance de biens nuit -- soit dit en passant -- à la santé du corps comme celle de l'esprit ! (contrairement à un lieu commun par trop répandu dont Bloy n'a, à ma connaissance, étrangement pas fait l'exégèse -- Flaubert aurait parlé d'une idée reçue (c'est-à-dire bourgeoise). Les latins et leur solide bon sens auraient rétorqué que \latin{Copia parit fastidium} et que \latin{Neque laus in copia, neque culpa in penuria constitit}.)}\\\passeLigne}%
\end{groupe}%
%% v. 634-635a
\begin{groupe}%
\personnageGrec{Μίκα}%
\trimIamb{οὐδὲν λέγεις. δεῦρ᾽ ἐλθὲ δεῦρ᾽ ὦ Κλείσθενες·\apparat{634}{δεῦρο \Rac, \Mon : \lecon{δευρ ὦ} \Rdeux \espApp Κλεισόσθενες R : \lecon{Κλείσθενες} Bonini \ligneApp{635} ἀνὴρ R : \lecon{ἁνὴρ} Bentley}\\
ὅδ᾽ ἐστὶν ἁνὴρ ὃν λέγεις.%\apparat{635}{ἀνὴρ R : \lecon{ἁνὴρ} Bentley}
\\\passeLigne}%
\end{groupe}%
%% v. 635b
\begin{groupe}%
\personnageGrec{Κλεισθένης}%
\trimIamb{\ligne{635}\decalage{ὅδ᾽ ἐστὶν ἁνὴρ ὃν λέγεις.}τί οὖν ποιῶ;\\\passeLigne}%
\end{groupe}%
%% v. 636
\begin{groupe}%
\personnageGrec{Μίκα}%
\trimIamb{ἀπόδυσον αὐτόν· οὐδὲν ὑγιὲς γὰρ λέγει.\\\passeLigne}%
\end{groupe}%
%% v. 637-8
\begin{groupe}%
\personnageGrec{Μνησίλοχος}%
\trimIamb{κἄπειτ᾽ ἀποδύσετ᾽ ἐννέα παίδων μητέρα;\apparat{637-40}{(Μνη.) R : (Μνη.) (637), (Κλεισθ.) (638) et (Γυν.) (639-40) Brunck, edd. post. : \lecon{(Μνη.) (637-8) et (Μίκα) (639-40)} Bibi \ligneApp{638} χάλᾳ R : χάλα Grynaeus : \lecon{χαλᾷ} Bibi%\espApp χάλα.... στρόφιον Micae, ἀναίσχυντε σύ Affini tribuit Degani
}%
\\
χαλᾷ ταχέως τὸ στρόφιον ὦναίσχυντε σύ.%
\commentaire{637-40}{κἄπειτ᾽ ἀποδύσετ᾽ ἐννέα παίδων μητέρα; ... καὶ νὴ Δία τίτθος γ᾽ ὥσπερ ἡμεῖς οὐκ ἔχει.}{%Les tripatouillomanes ont encore sévi à ce vers !.. ὦναίσχυντοι, \italique{quousque tandem...} ô Willems, combien la tâche que tu t'étais confié est loin d'être terminée !\\
%Le vers coulait pourtant d'une source claire : Clisthène est en train de déshabiller Mnésiloque ; ce dernier feint donc l'indignation (χαλᾷ est un subjonctif dit \italique{délibératif}) devant l'empressement (ταχέως) que cet homme -- car Clisthène reste un homme -- met à enlever le \italique{soutien gorge} d'une femme -- car Mnésiloque continue de jouer une femme (vertueuse même, ici, mais un peu tard). Qu'y a-t-il de plus limpide ? l'apostrophe ὦναίσχυντε σύ conclue logiquement la tirade de mère effarouchée qu'affecte le peu crédible Mnésiloque (et qui fait d'ailleurs écho au précédent ἀναίσχυντός τις εἶ lancé par le même Mnésiloque au vers 611\marqueur{1}).\\
%Qu'ont eu besoin les éditeurs, tous, de suivre Brunck\marqueur{2} ? et d'inventer un changement d'interlocuteurs là où il n'y avait pas lieu d'être -- alors que le \italique{Ravennas} est on ne peut plus fiable sur ce point -- et, \italique{mutatis mutandis}, de surajouter une autre retouche devenue nécessaire : χάλα à la place de χαλᾷ.
Le manuscrit de Ravenne porte en marge du vers 638 μν\auDessus{η}{σ}, indiquant ainsi que Mnésiloque prend la parole ; le manuscrit ne présente alors plus aucun signe de changement de locuteur jusqu'au vers 641\marqueur{1}. Or s'il est sûr que Mnésiloque prononce le vers 637, ἀποδύσετ᾽ ἐννέα παίδων μητέρα, il ne peut aucunement prononcer le vers 639, ὡς καὶ στιβαρά τις φαίνεται καὶ καρτερά·
Il y a donc un \italique{paragraphos} à rétablir !
Je dis \italique{un} car les éditeurs suivent quasi-unanimement Brunck en `rétablissent' deux : au début du vers 638 (Clistène) et au début du vers 639 (la femme -- Mika -- reprend la parole).
Une autre solution me semble plus légère : laisser à Mnésiloque les deux vers 637-638 et à Mika les vers 639-640.
Cette retouche comporte surtout l'avantage d'être bien meilleure quant à l'action et au jeu de scène.\\
Presque tous les éditeurs depuis Brunck, disais-je, retirent le vers 638 à Mnésiloque pour l'accorder à Clisthène. Double erreur ! \ \ \ \italique{Primo}, parce qu'ils sont alors contraints de surajouter une seconde émendation : transformer le χάλᾳ du manuscrit en un χάλα (impératif lancé par Clisthène). \ \ \ \italique{Secundo}, parce que Mika, au vers 636, avait déjà ordonné à Clisthène de déshabiller Mnésiloque (ἀπόδυσον αὐτόν) ; il n'y a aucune de raison de douter qu'il ne soit en train de lui obéir ; un peu d'action que diable ! une pièce d'Aristophane ne doit pas avoir la même pesanteur que l'administration française (Mika ordonnerait, selon nos éditeurs, à Clisthène de déshabiller Mnésiloque, qui transmettrait alors l'ordre à l'intéressé... qui finirait par ne rien faire ?..)\\
Le vers était pourtant excellent dans la bouche de Mnésiloque : Clisthène est en train de le déshabiller ; il feint donc l'indignation (χαλᾷ est un subjonctif dit \italique{délibératif}) devant l'empressement (ταχέως) que cet homme -- car Clisthène reste un homme -- met à enlever le \italique{soutien gorge} d'une femme -- car Mnésiloque continue de jouer une femme (vertueuse même -- mais un peu tard). Qu'y a-t-il de plus limpide ? l'apostrophe ὦναίσχυντε σύ conclue logiquement la tirade de mère effarouchée qu'affecte le peu crédible Mnésiloque (et qui fait d'ailleurs écho au précédent ἀναίσχυντός τις εἶ (v. 611) lancé par le même Mnésiloque au même Clisthène\marqueur{2}).
}%
\debutNote{637-40}%
%\note{1}{D'aucuns avaient même prétexté de ce double ἀναίσχυντος pour attribuer ἀναίσχυντός τις εἶ à Clisthène, histoire que leur texte présente une homogénéité certaine : foutaises partout, justesse nulle part.}%
%\note{2}{Quelques uns, croyant faire preuve d'originalité, attribuèrent la réplique à Mika plutôt qu'à Clisthène. Mais tous tripatouillent !.. tous !!! La pire espèce étant celle qui tente de camoufler son forfait en le faisant disparaître de l'apparat critique.}%
\note{1}{Il est d'usage dans notre manuscrit que le \italique{signe} soit un παράγραφος marginal ` --- ' (c'est-à-dire un tiret cadratin) lorsque le changement de locuteur se fait en début de vers ou un \latin{punctum duplex} ` : ' lorsqu'il se produit en milieu de vers. Par abus de langage, le terme de \italique{paragraphos} est souvent utilisé pour désigner indifféremment l'un de ces signes.}%
\note{2}{Quelques éditeurs avaient même prétexté de ce double ἀναίσχυντος pour attribuer l'ἀναίσχυντός τις εἶ du vers 611 à Clisthène, histoire que leurs multiples tripatouillages présentent une apparence d'homogénéité...}%
\\\passeLigne}%
\end{groupe}%
%% v. 639-640
\begin{groupe}%
\personnageGrec{Μίκα}%
\trimIamb{ὡς καὶ στιβαρά τις φαίνεται καὶ καρτερά·\\%\apparat{639-40}{(Μνη.) 639 et (Γυνή) 640 R : \lecon{(Μίκα)} Bibi}\\
καὶ νὴ Δία τιτθούς γ᾽ ὥσπερ ἡμεῖς οὐκ ἔχει.\apparat{640}{τίτθος R : \lecon{τιτθούς} ed. Iuntina}\\\passeLigne}%
\end{groupe}%
%% v. 641
\begin{groupe}%
\personnageGrec{Μνησίλοχος}%
\trimIamb{στερίφη γάρ εἰμι κοὐκ ἐκύησα πώποτε.\testimonia{641}{(στερίφη) Hesychius σ 1769 στερίφη $=$ Photius σ 532 ; cf. Moeris σ 46 (στερίφη $=$ στερρά)}\apparat{641}{στεριφῇ R : \lecon{στερίφη} (sed στεριφὴ) \Mon}%
\commentaire{641}{στερίφη γάρ εἰμι.}{Jeu de mots d'Aristophane : comment Mnésiloque explique-t-il la platitude de sa poitrine ? tout simplement parce qu'il a les chairs \italique{fermes} (sens premier de στέριφος, le Scythe l'emploie avec cette acception au vers 1185). Mais στέριφος peut aussi avoir le sens dérivé de \italique{stérile} ; sens que le spectateur ne privilégiera qu'en entendant la seconde partie de la réplique de Mnésiloque (κοὐκ ἐκύησα). C'est la raison pour laquelle Aristophane recourt à un mot rare, στέριφος (le scholiaste de noter : ἀντὶ τοῦ στείρα).\\
%Y a-t-il ici aussi détournement d'un tour proverbial ? toujours est-il qu'on retrouve un tour identique chez Hipponax (fr. 122 ed. Degani) :%
%\citComm{ἀμφιδέξιος γάρ εἰμι κοὐκ ἁμαρτάνω κόπτων.}%
}%
\\\passeLigne}%
\end{groupe}%
%% v. 642
\begin{groupe}%
\personnageGrec{Κλεισθένης}%
\trimIamb{νῦν· τότε δὲ μήτηρ ἦσθα παίδων ἐννέα.\apparat{642-3}{nulla nota mutatae personae R : (Γυνή) (642) et (Κλεισθ.) (643) Ellebodius : (Κλεισθ.) (642-3) \ligneApp{642} νῦν τότε δὴ R : \lecon{νῦν τότε δὲ} \Mon}%
\commentaire{642-3}{νῦν· τότε δὲ μήτηρ ἦσθα παίδων ἐννέα. / ἀνίστασ᾽ ὀρθός. ποῖ τὸ πέος ὠθεῖς κάτω;}{Il n'y a aucun \italique{paragraphos} aux vers 642 et 643 alors qu'assurément ces vers ne peuvent être dits par Mnésiloque. Les éditeurs ajoutent deux changement d'interlocuteurs. Par simplicité, je préfère n'en ajouter qu'un et attribue donc le distique entier à Clisthène.}\\
ἀνίστασ᾽ ὀρθός. ποῖ τὸ πέος ὠθεῖς κάτω;%
\commentaire{643}{ἀνίστασ᾽ ὀρθός.}{«~Redresse-toi !~» L'injonction adressée à Mnésiloque est bien-sûr à double entente, et Clisthène se fera doublement obéir. Mais interrogeons-nous : pourquoi Mnésiloque s'était-il assis (ou plutôt recroquevillé) ? réponse simple : pour cacher sa virilité (nous y reviendrons à la fin de cette note). 
Ils se redressent ; Mnésiloque tente une dernière fois de cacher ce qui doit être caché et Clisthène : ποῖ τὸ πέος ὠθεῖς κάτω; «~Où repousses-tu ce \italique{péos} ?!~» Le \italique{péos} passe alors au verso de Mnésiloque, d'où peut le contempler Mika, avant que ne s'ensuive la scène burlesque.\\
Revenons à ce qui nous intéresse pour la bonne compréhension de l'action présente : Clisthène vient de retirer le στρόφιον de Mnésiloque, ce qui entraîna quelques remarques sur sa poitrine. Or, lors de la scène d'habillement (vers 253-261), Mnésiloque avait successivement enfilé le κροκωτός, le στρόφιον \italique{puis à la fin} l'ἔγκυκλον. Il résulte donc que l'ἔγκυκλον \italique{est un vêtement qui ne peut pas recouvrir la poitrine}. (Sinon Clisthène aurait dû commencer par l'encycle.) 
Et puisque, jusqu'ici, seules d'ineptes sottises ont été écrites à propos de l'ἔγκυκλον, il nous a paru profitable d'allonger cette note afin que soit irréfragable l'identification de l'ἔγκυκλον avec une sorte de jupe courte\marqueur{1}.\\
Ce vêtement qui se porte plutôt à l'intérieur aux dires d'Eustathe (voir ci-dessous).
Ce vêtement que l'étymologie, ἔγ-κυκλον, désigne comme ce qui \italique{en-toure}, ce qui \italique{en-ceint}.
Ce vêtement qui sitôt retiré laisse tout émoustillé le mari de Praxagora ; voyez plutôt (\italique{Assemblée des femmes}, vers 535-538) :%
\citComm{Βλέπυρος -- εἶτ' οὐ τὸ σαυτῆς ἱμάτιον ἐχρῆν σ' ἔχειν;\\
ἀλλ' ἔμ' ἀποδύσασ' ἐπιβαλοῦσα \textbf{τοὔγκυκλον}\\
ᾤχου καταλιποῦσ' ὡσπερεὶ προκείμενον,\\
μόνον οὐ στεφανώσασ' οὐδ' ἐπιθεῖσα λήκυθον.\\
Alors que n'as-tu pris ton propre manteau ? Mais tu me dévêts et jettes sur moi \italique{ton encycle}, puis tu t'en vas, me laissant là comme un cadavre exposé, hors que tu as omis de me couronner et de placer près de moi un lécythe.}%
Le champ lexical est celui d'une exposition funèbre\marqueur{2} comme la traduction de Willems le laisse entendre, mais le λήκυθον du vers 538 doit être pris \italique{sensu obsceno} selon la matrice métaphorique : λήκυθος $=$ ἀλάβαστος\marqueur{3} $=$ \latin{phallus}.
Blépyros reproche à sa femme de ne pas s'être acquitté du \italique{nome de mâtines} (ὄρθριος νόμος\marqueur{4}).\\
Ce vêtement qui permet à Aristophane un jeu de mot avec l'\italique{enceinte} de la ville fortifiée de Pylos\marqueur{5} -- oui ! --, la Πύλος, cette ville qui déclenche intarissablement la verve éroticomique de notre poète, jonglant avec l'autre matrice métaphorique (féminine cette fois) : πύλος $=$ πύλη $=$ θύρα $=$ \latin{pudenda muliebria} (raison pour laquelle l'Athénien ne veut pas la céder). 
%Excusez-moi des longueurs ; je me vois contraint d'expliciter quelque peu Aristophane à ses propres commentateurs devant toutes les conneries qu'ils débitent -- littéralement : ils pensent que l'\italique{encycle} désigne le \italique{con}. 
%Cela dit, revenons à \italique{Lysistrata} : 
Il y avait certainement à l'époque de la représentation de \italique{Lysistrata} -- qui, rappelons-le, est de quelques mois antérieure à celle des \italique{Thesmophories} -- des tractations pour que la paix fût conclue en échange de la remise de Pylos ou, par synecdoque, de son \italique{encycle}. 
Ce n'est pas autrement qu'il faut comprendre le passage au début de la même pièce (vers 111-114) : Lysistrata demande : %
\citComm{ἐθέλοιτ᾽ ἂν οὖν, εἰ μηχανὴν εὕροιμ᾽ ἐγώ,\\
μετ᾽ ἐμοῦ καταλῦσαι τὸν πόλεμον;\\
Consentirez-vous donc, si j'en trouve le moyen, à mettre fin avec moi à la guerre ?}%
et Aristophane de répondre, par la voix d'une femme :%
\citComm{\decalage{μετ᾽ ἐμοῦ καταλῦσαι τὸν πόλεμον;}νὴ τὼ θεώ\\
ἔγωγ᾽ ἂν οὖν κἂν εἴ με χρείη \textbf{τοὔγκυκλον}\\
τουτὶ καταθεῖσαν ἐκπιεῖν αὐθημερόν.\\
Oui, par les deux déesses. J'en suis, quand il me faudrait déposer l'\italique{encycle} que voici en gage pour boire [à la paix : les traités, σπονδαί, sont en grec, littéralement, des \italique{libations}] dès aujourd'hui.}%
Il faut donc bien-sûr aussi attribuer à cette même femme -- et comme ne le font pas les éditeurs d'Aristophane -- le vers 104 :%
\citComm{ὁ δ᾽ ἐμός γε τελέους ἑπτὰ μῆνας ἐν Πύλῳ.}%
(Elle se dit prête à faire tomber l'encycle dès le retour de son mari en garnison depuis sept longs mois à Pylos.)\\
Aristophane, par deux fois dans la même pièce, exhorte ses contemporains à renoncer à Pylos, c'est dire l'importance qu'il accordait à l'éventualité de la paix ! On connaît la suite...\\
Ce vêtement, donc, qui enclot le sadinet, puisqu'il faut parler net... Alors ? que peut-il bien être ce vêtement ? S'il vous plaît, ne me répondez pas de chœur avec tous les fonctionnaires de la recherche : «~un lourd et pesant manteau~».\\
Mais délaissons un peu nos primes amours, les comédies d'Aristophane, pour aborder un terrain plus aride et cependant qui n'est certainement pas à négliger, je veux parler de la lexicologie.
Ce sens de \italique{jupe} pour l'ἔγκυκλον est confirmé, ce me semble, par le commentaire le plus précis : celui du scholiaste de \italique{Lysistrata} (vers 113), qui glose ἔγκυκλον ainsi :%
\citComm{τὸ περιβόλαιον τὸ ζωνάριον.}%
J'interprète ζωνάριον comme un adjectif se rapportant à περιβόλαιον\marqueur{6} : \italique{le vêtement / la couverture}, celle qui enceint, \italique{au niveau de la taille} ; une jupe donc !\\
J'en finirais avec l'\italique{encycle} en citant ce passage du copieux commentaire à l'\italique{Iliade} (976.13) d'Eustathe de Thessalonique (il commente \italique{Iliade} XIV, v. 178)\marqueur{7} :%
\citComm{ἰστέον δὲ ὅτι, καθάπερ ὁ ἑανὸς οὕτω καὶ τὸ ἔκκυκλον ἢ ἔγκυκλον, οὗ καὶ ὁ κωμικὸς μέμνηται, γυναικεῖον ἦν ποτε φόρημα, φησὶ γοῦν Παυσανίας ὅτι ἔγκυκλον περιπόρφυρον\marqueur{8} ἱμάτιον καὶ χιτὼν γυναικεῖος, ὃν ἔνδοθεν ἐνδύονται γυναῖκες, εἶτα τὸ ἔνδυμα. λέγεται δὲ φησίν [sic, je lis : σφῇσιν] ὑπ' ἐνίων ἐπιγονατίς, διότι μέχρι γονάτων διήκει.\\
Il faut savoir que, de même que l'ἑανὸς [= riche robe de femme dans le passage d'Homère qu'Eustathe commente], il y a aussi l'ἔκκυκλον (ou ἔγκυκλον), dont s'est aussi souvenu le Comique [= Aristophane]. C'était un vêtement (φόρημα) féminin porté jadis. Ce qui est sûr, c'est que Pausanias indique que l'ἔγκυκλον était un vêtement (ἱμάτιον) aux franges pourpres (περιπόρφυρον) et une tunique (χιτὼν) féminine, que les femmes revêtaient à l'intérieur, et ensuite un habit (ἔνδυμα). Il est qualifié par elles-mêmes, d'après certains, d'ἐπιγονατίς, parce qu'il descend jusqu'aux genoux.}%
Fermons cette longue parenthèse sur l'\italique{encycle} qui, assurément, est une \italique{jupe}.\\
Nous disions que Mnésiloque s'était habillé dans des cet ordre : κροκωτός, στρόφιον et ἔγκυκλον. Le στρόφιον et l'ἔγκυκλον viennent d'être retirés, mais pas le κροκωτός (au vers 941, Mnésiloque le porte toujours) ; comment se fait-il que son olisbos soit apparent malgré la présence de cette robe safranée ? parce que la crocote est composée d'un tissu très fin qui la rend transparente\marqueur{9}.%
}%
\debutNote{643}%
\note{1}{Et certainement pas un \italien{mantello ampio, pesante,} comme le conclue Franca Perusino dans son article «~\italien{L'ἔγκυκλον, un mantello femminile nelle commedie di Aristofane}~», \italique{Quaderni Urbinati di Cultura Classica}, 72.3 (2002), pp. 129-133. Je ne cite cet article que par pur scrupule bibliographique : il présente tous les travers philologiques rencontrés chez les commentateurs d'Aristophane. Ses conclusions erronées se sont reproduites dans la plupart des travaux portant sur le costume de l'Ancienne comédie... espérons que cette note rectificatrice mettra enfin un terme à la contagion ! Ne soyons pas trop injuste et signalons la belle trouvaille de cet article : une probable identification d'un ἔγκυκλον sur un vase à figures rouges de la première moitié du \nieme{V} siècle. Ce vase, reproduit en annexe, montre une femme nue, un χιτών à la main, et, à côté d'elle, des bottes ainsi qu'un tissu à franges (περιπόρφυρον, cf. le commentaire d'Eustathe) posé sur un tabouret... notre ἔγκυκλον ? C'est bien ce qu'il m'en semble ! Chacun prêchant pour sa propre paroisse, j'y vois effectivement une \italique{jupe} et non un \italien{mantello ampio e pesante}. Je laisse le lecteur juge de l'iconographie ; pour ce qui est du sens d'ἔγκυκλον chez Aristophane, le doute n'est pas possible. Ce vase est le n° 204020 du \italique{Corpus Vasorum Antiquorum}, Beazley l'a recensé (378.122) dans son monumental \italique{Attic Red-Figure Vase-Painters}\expo{2} (1963).}%
\note{2}{Cf. par exemple \italique{Lysistrata}, vv. 599 sqq.}%
\note{3}{Cf. la saillie du chaud-lapin Kinésias soumis à la torture par Myrrhine (\italique{Lysistrata}, v. 947) :%
\citComm{\textbf{Μυρρίνη} -- λαβὲ τόνδε τὸν ἀλάβαστον. --- \textbf{Κινησίας} -- ἀλλ᾽ ἕτερον ἔχω.}%
}%
\note{4}{\italique{Assemblée des femmes}, vers 741. Par jeu de mots avec l'ὄρθιος νόμος, le \italique{mode orthien}, qu'Hésychius (ο 1188) n'a pas manqué de nous définir :%
\citComm{ὄρθιον νόμον· οὕτως ὠνομάσθαι φασὶν αὐτὸν ἀπὸ τοῦ ῥυθμοῦ, καθάπερ καὶ τὸν τροχαῖον.}%
(Cf. aussi sur le mode orthien \italique{Cavaliers} 1280 ; \italique{Acharniens} 16.)\\
Les allusions d'Aristophane à cette tradition matinale à laquelle sacrifient ses personnages sont légions ; ne citons que \italique{Lysistrata} v. 966 puisque le passage est encore sous nos yeux.}%
\note{5}{\italique{Lysistrata}, vers 1162-1163. Pour conclure la paix, l'ambassadeur laconien demande à récupérer l'\italique{encyle} :%
\citComm{Λάκων -- ἁμές γε λῶμες, αἴ τις ἁμὶν τὤγκυκλον\\
λῇ τοῦτ᾽ ἀποδόμεν. --- Λυσιστράτη -- ποῖον ὦ τᾶν; --- Λάκων -- τὰν Πύλον,\\
ἇσπερ πάλαι δεόμεθα καὶ βλιμάττομες.}%
}%
\note{6}{Chantraine, page 74 de sa \italique{Formation des noms en grec ancien}, précise : «~la finale -άριον qui a fourni essentiellement des diminutifs s'observe parfois dans des adjectifs, cf. σωμα παιδαριον dans des inscription de Delphes~».}%
\note{7}{La faussaire précitée avait mutilé ce texte, et ce afin de le tordre à ses conclusions -- selon la pente bien connue de ces gens-là. On constate qu'il n'était nul besoin de le charcuter : une minuscule émendation (ajouter un sigma à φησίν et lire σφῇσιν) suffit à le rendre intelligible.}%
\note{8}{On lit en épigraphie : ἔνκυκλον περιποίκιλο, ἔγκυκλον ποικίλον. À l'inverse, on croise aussi un ἔνκυκλον λευκ[όν].}%
\note{9}{Cf. le cratère à figures rouges (ca. -380) représentant une scène de comédie : un vieillard entouré de deux gardes et vêtu de la crocote... qui laisse tout transparaître ; il est reproduit en annexe. Cratère «~St. Agata Antigone~», St. Agata dei Goti I, collection Rainone-Mustilli. Inventorié par Arthur D. Trendall , \italique{Phlyax Vases}, «~BICS~» suppl. LXIX (1967).}%
\\\passeLigne}%
\end{groupe}%
%% v. 644
\begin{groupe}%
\personnageGrec{Μίκα}%
\trimIamb{τοδὶ διέκυψε καὶ μάλ᾽ εὔχρων ὦ τάλαν.\apparat{644}{τὸ δὶ δὴ ἔκυψε R : \lecon{τοδὶ διέκυψε} Dobree, Dindorf}%
\commentaire{644}{τάλαν.}{Ce mot, le vocatif de τάλας, tant masculin que féminin, n'est mis par Aristophane que dans la bouche de femmes\marqueur{1}.}%
\debutNote{644}%
\note{1}{\italique{Lysistrata}, vers 102, 910 et 914 ; \italique{Grenouilles}, vers 559 ; \italique{Assemblée des femmes}, vers 124, 658 et 1005 ; \italique{Ploutos}, vers 706 et 1055.}\\\passeLigne}%
\end{groupe}%
%% v. 645a
\begin{groupe}%
\personnageGrec{Κλεισθένης}%
\trimIamb{καὶ ποῦ 'στιν;\\\passeLigne}%
\end{groupe}%
%% v. 645b
\begin{groupe}%
\personnageGrec{Μίκα}%
\trimIamb{\ligne{645}\decalage{καὶ ποῦ 'στιν;}αὖθις εἰς τὸ πρόσθεν οἴχεται.\\\passeLigne}%
\end{groupe}%
%% v. 646a
\begin{groupe}%
\personnageGrec{Κλεισθένης}%
\trimIamb{οὐκ ἐνγεταυθί.\apparat{646}{\lecon{ἐνγεταυθί} R : ἐγγεταυθί Thiersch}\\\passeLigne}%
\end{groupe}%
%% v. 646b
\begin{groupe}%
\personnageGrec{Μίκα}%
\trimIamb{\ligne{646}\decalage{οὐκ ἐνγεταυθί.}ἀλλὰ δεῦρ᾽ ἥκει πάλιν.\apparat{646}{\lecon{ἀλλὰ} R : μἀλλὰ Bentley}\\\passeLigne}%
\end{groupe}%
%% v. 647-648
\begin{groupe}%
\personnageGrec{Κλεισθένης}%
\trimIamb{ἰσθμόν τιν᾽ ἔχεις ὤνθρωπ᾽· ἄνω τε καὶ κάτω\testimonia{647-8}{Suidas ι 640}\apparat{647}{ῶνθρωπε R : ἄνθρωπε \Mon : ἄνθρωπ᾽ Suid. ι 640 : \lecon{ὤνθρωπ᾽} Ellebodius}%
\commentaire{647-8}{ἰσθμόν τιν᾽ ἔχεις ὦνθρωπ᾽· ἄνω τε καὶ κάτω / τὸ πέος διέλκεις πυκνότερον Κορινθίων.}{Clisthène fait référence au \grec{δίολκος} : un chemin dallé de 8 kilomètres qui permettait de faire traverser aux bateaux l'isthme de Corinthe grâce à une sorte de rail rudimentaire et d'éviter ainsi le long et périlleux contournement du Péloponnèse\marqueur{1}. Pendant la scène burlesque de déshabillage qui vient d'avoir lieu, Mnésiloque, pour cacher sa virilité – et certainement aussi émoustillé d'être effeuillé de la sorte –, a actionné à plusieurs reprises le mécanisme permettant de lever et baisser son \italique{phallus}. Admiratif devant de telles performances, Clisthène le flatte d'être encore mieux outillé que les machines qui permettaient le transbordement des bateaux.\\
Voilà pour l'image `τὸ πέος διέλκεις πυκνότερον Κορινθίων' ! il reste à expliquer la première partie de la réplique de Clisthène : ἰσθμόν τιν᾽ ἔχεις, «~tu possèdes un [drôle de] genre d'\italique{isthme}~».
ἰσθμός désigne proprement un \italique{passage étroit} et, par analogie, le \italique{cou}. Le spectateur athénien saisissait instamment l'image : par \italique{isthme} il fallait entendre \italique{pénis} grâce à la matrice métaphorique suivante : ἰσθμός = αὐχήν\marqueur{2} = τράχηλος\marqueur{3} = πέος.}%
\debutNote{647-8}%
\note{1}{Voir à ce sujet l'article de M. J. T. Lewis, « Railways in the Greek and Roman world », in \italique{Early Railways : a Selection of Papers from the First International Early Railways Conference}, 2001, p. 8-19. Thucydide (III, 15 ; VIII, 7) évoque le δίολκος dans sa \italique{Guerre du Péloponnèse}.}%
\note{2}{αὐχήν a le sens de πέος par exemple au vers 681 de \italique{Lysistrata} : le chœur des vieillards a maille à partir avec l'engeance féminine ; après avoir rappelé par une métaphore hippique combien la femme est portée à la luxure, il en vient à la conclusion que le meilleur moyen pour dominer ces fougueuses juments est de les chevaucher -- la fameuse \italique{phallocratie} ? Il formule cette idée à l'aide d'une image carcérale : il faut assujettir leur \italique{bois percé} (= \italique{carcan}) au \italique{cou} (αὐχήν) :
\citComm{ἀλλὰ τούτων χρῆν ἁπασῶν ἐς τετρημένον ξύλον\\
ἐγκαθαρμόσαι λαβόντας \textbf{τουτονὶ τὸν αὐχένα}.\\
Eh bien, il importe de s'assurer d'elles toutes, et d'assujettir dans un carcan \italique{ce cou-là}.}
On peut avoir un autre exemple d'αὐχήν avec ce sens-là au vers 184 du \italique{Cyclope} d'Euripide.}%
\note{3}{τράχηλος est phallossémique au vers 20 des \italique{Grenouilles}. Aristophane commence la pièce avec un rapide survol des blagues lourdes et vulgaires de ses concurrents aux concours comiques ; ce n'est pas sa muse qui s'abaisserait à de si grossières manœuvres ! non, \italique{on n'en soufflera mot}, οὐκ ἐρεῖ ! οὐκ ἐρεῖ !.. et puis au vers 19 :
\citComm{ὢ τρισκακοδαίμων ἄρ᾽ \textbf{ὁ τράχηλος οὑτοσί},\\
ὅτι θλίβεται μέν, τὸ δὲ γέλοιον οὐκ ἐρεῖ.\\
Ô trois fois malencontreuse \italique{nuque}, qui est moulue, et n'a pas le mot pour rire !}
Il faut bien-sûr lire -- je dis «~bien-sûr~», mais personne ne l'avait proposé jusqu'ici... --, il faut donc lire :
\citComm{ἄρ᾽ ὁ τράχηλος οὑτοσί,\\
ὅτι θλίβεται μέν, τὸ δὲ γέλοιον \textbf{οὐ κε ῥεῖ;}\\
Est-ce que le τράχηλος que voici, puisqu'il est pressé, pour la farce ne devrait pas couler ?}}\\
τὸ πέος διέλκεις πυκνότερον Κορινθίων.\\\passeLigne}%
\end{groupe}%
%% v. 649-650a
\begin{groupe}%
\personnageGrec{Μίκα}%
\trimIamb{ὤ μιαρὸς οὗτος· ταῦτ᾽ ἄρ᾽ ὑπὲρ Εὐριπίδου\\
ἡμῖν ἐλοιδορεῖτο.\\\passeLigne}%
\end{groupe}%
%% v. 650b-651
\begin{groupe}%
\personnageGrec{Μνησίλοχος}%
\trimIamb{\ligne{650}\decalage{ἡμῖν ἐλοιδορεῖτο.}κακοδαίμων ἐγώ,\\
εἰς οἷ᾽ ἐμαυτὸν εἰσεκύλισα πράγματα.\apparat{651}{εἰσεκύλησα R : \lecon{εἰσεκύλισα} Scaliger}%
\commentaire{651}{εἰσεκύλισα.}{Homère emploie le verbe κυλίνδειν en parlant de malheurs que le ciel \italique{fait rouler} sur quelqu'un : πῆμα θεὸς Δαναοῖσι κυλίνδει (\italique{Il}. 17, 688)\marqueur{1}. La même image se trouve chez Aristophane : νὴ τὸν Διόνυσον ἄπορά γ᾽ ἡμῖν πράγματα / δαίμων τις ἐσκεκύκληκεν ἐς τὴν οἰκίαν (\italique{Guêp.} 1474 s.), « ah! les inextricables tracas qu'un dieu a fait dévaler sur notre maison. » Le même verbe est employé d'une manière un peu différente, \italique{Thesm}. 767, mais toujours à propos de malheurs : ὁ μὲν γὰρ αἴτιος / κἄμ᾽ ἐσκυλίσας εἰς τοιαυτὶ πράγματα / οὐ φαίνεταί πω, « le coupable, celui qui m'a fait rouler dans de tels malheurs, n'apparaît toujours pas. » Cf. \italique{Thesm}. 650 s. : κακοδαίμων ἐγώ, / εἰς οἷ᾽ ἐμαυτὸν εἰσεκύλισα πράγματα. Il est clair que, dans ce cas, ἐμβάλλειν τινὰ εἰς συμφοράς n'est qu'un synonyme expressif\marqueur{2} de ἐμβάλλειν (p. ex. ἐμβάλλειν τινὰ εἰς συμφοράς, Antiphon 3, 4, 10 ; εἰς ἀτυχίας, Eschin. 3, 79 ; εἰς κακόν, Ar. \italique{Nuées} 1460). (Note de Taillardat, §558.)}%
\debutNote{651}%
\note{1}{Le moyen κυλίνδεσθαι se dit de malheur lui-même «~roulant~» sur quelqu'un : νῶιν δὴ τόδε πῆμα κυλίνδεσθαι, \italique{Il.}, 11, 347 ; cf. \italique{Od}. 2, 163 ; 8, 81.}%
\note{2}{Connu d'autres auteurs ; Phérécr. 146 [K = 156 K.-A.] : εἰκῆ μ᾿ ἐπῆρας ὄντα τηλικουτονὶ / πολλοῖς ἐμαυτὸν ἐγκυλῖσαι πράγμασιν ; Cf. Xén. \italique{Mém.} I, 2, 22 : τοὺς εἰς ἔρωτας ἐγκυλισθέντας ; Ἀπολλόδωρος ὁ Καρύστιος 5, 7 [ $=$ Athénée 280e] : εἰκῆ τέ πως / ἡμᾶς κυλίνδους᾽ ὅντιν᾽ ἂν τύχῃ τρόπον (sc. τύχη ἄγροικος).}\\\passeLigne}%
\end{groupe}%
%% v. 652a
\begin{groupe}%
\personnageGrec{Μίκα}%
\trimIamb{ἄγε δὴ τί δρῶμεν;\\\passeLigne}%
\end{groupe}%
%% v. 652b-654
\begin{groupe}%
\personnageGrec{Κλεισθένης}%
\trimIamb{\ligne{652}\decalage{ἄγε δὴ τί δρῶμεν;}τουτονὶ φυλάττετε\\
καλῶς, ὅπως μὴ διαφυγὼν οἰχήσεται·\apparat{653}{οἴχεται R : \lecon{οἰχήσεται} Faber (p.788)}\\
ἐγὼ δὲ ταῦτα τοῖς πρυτάνεσιν ἀγγελῶ.\apparat{654}{πρυτάνεσσιν R : \lecon{πρυτάνεσιν} \Mon}\\\passeLigne}%
\end{groupe}%
%% v. 655-662
\begin{groupe}%
\personnageGrec{Κορυφαῖος}%
\typeVers{\decalageN{4}$4$ an\cat}\tetramAnap{ἡμᾶς τοίνυν μετὰ τοῦτ᾽ ἤδη τὰς λαμπάδας ἁψαμένας χρὴ%
\commentaire{655-86}{Le chœur compris entre les vv. 655-686.}{La \nieme{19} note d'Alphonse Willems est d'une telle richesse d'explication que je ne peux que la reproduire intégralement ci-dessous. Pour ce qui touche à la colométrie, je reprends celle de Lætitia\marqueur{1}. (Le passage des vers 667-686 a de fortes correspondances métriques avec les vers 707-725).\\
Ce petit chœur n'a pas été examiné avec assez d'attention. On n'a pas remarqué qu'il s'accompagne d'un jeu de scène unique dans l'œuvre d'Aristophane. Trois morceaux le composent : \\
\ \ \ 1° Les vv. 655 à 658 appartiennent au coryphée. Mnésiloque n'est peut-être pas le seul homme qui ait pénétré parmi les femmes. Il importe de s'en assurer, et, après avoir déposé les manteaux et allumé les torches, de « tourner et virer tout à l'entour du Pnyx, et d'inspecter avec soin les tentes et les passages ».\\
\ \ \ 2° Les quatre suivants (659-662) sont des trochées, et le changement de mètre indique un changement d'interlocuteur. C'est, en effet, la réponse du \italique{parastate} ou chef du demi-chœur [le \italique{Ravennas} porte sur sa marge, η\auDessus{μ}{χ}ι γ\auDessus{υ}{ν}, écrite d'une seconde main, mais dont l'indication se retrouve sur le manuscrit de Münich]. En disant qu'il y avait lieu d'inspecter le Pnyx, les tentes et les passages, la coryphée indiquait clairement qu'il fallait quitter l'orchestre, c'est-à-dire l'enceinte sacrée adjacente au temple\marqueur{2}. La parastate annonce donc que le chœur va « s'élancer dehors (ἐξορμᾶν) d'un pied léger\marqueur{3} ». Puis elle engage celle qu'elle appelle τὴν πρώτην à commencer tout de suite la reconnaissance. Comment faut-il interpréter τὴν πρώτην ? D'y voir, comme on le fait généralement, une locution adverbiale (\latin{quam primum} ou \latin{prima quaque via}) me semble inadmissible. Avec M. Mazon, je crois que ce terme désigne une personne déterminée\marqueur{4} ; suivant moi, ce serait τὴν πρώτην ἱέρειαν, appelée aussi τὴν προστάτιν, savoir la première, il y en avait beaucoup\marqueur{5}, des prêtresses temporaires, ἀρήτειραι, de Déméter, celle qui, tout à l'heure, avait prononcé les prières d'usage (vv. 295 et suiv.), en un mot, la coryphée. 
Callimaque, parlant d'une de ces prêtresses, ἱερέη Δήμητρος, se sert d'une désignation analogue : πολλῶν προστασίη νέων γυναικῶν\marqueur{6}. \\
\ \ \ 3° Le morceau qui suit (vv. 663-684) est un \italique{hyporchême}. Une partie du chœur reste en place et chante, tandis que l'autre, en quête de quelque complice de Mnésiloque\marqueur{7}, évolue silencieusement, σιωπῇ (v. 660), va et vient, entre et sort, cherchant partout, dans l'orchestre et dehors, καὶ τὰ τῇδε καὶ τὰ δεῦρο (v. 666), jusqu'à ce qu'il déclare n'avoir rien trouvé (687-688).\\
%\eff{Après cette fine analyse de Willems, s'agit-il de roues (trokhos) cf. Daremberg Salvatio 1027a ? τὴν πρώτην τρέχειν χρῆν ὡς τάχιστ᾽ ἤδη κύκλῳ}
}%
\debutNote{655-86}%
\note{1}{\italique{The songs of Aristophanes}, pp. 418-422.}%
\note{2}{On peut conclure de là que le Thesmophorion, qu'on savait situé sur une hauteur, l'était sur la colline du Pnyx et voisin du lieu d'assemblée.}%
\note{3}{Le simple, ὀρμᾶν, qui se lit plus loin, v. 953, et \italique{Gren}., 478, signifierait « prendre sa course ».}%
\note{4}{\italique{Essai sur la composition des comédies d'Aristophane}, Paris (1904), p. 131, n. 4.}%
\note{5}{Il y en avait deux par dème (Isée, \italique{Succession de Kiron}, 19-20), donc plus de deux cents.}%
\note{6}{\italique{Anthol. Palat}., VII, 728.}%
\note{7}{εἴ τις ἐν τόποις ἑδραῖος. Qu'avons-nous affaire d'ajouter, avec MM. Blaydes et van Leeuwen, τούτοισι τοῖς devant τόποις ? La même locution adverbiale se lit dans un fragment de Nausicratès : [τρίγκλαι] ἃς κλύδων Αἰξωνικὸς ἐν τόποις παιδεύεται (\italique{Fragm. Com. Gr}., t. IV, p. 575 [= Athénée 325e = Nausicratès, fr. 1.6-11 Kassel-Austin ; le texte n'est pas établi avec certitude : on trouve aussi : ἐντόπους à la place de ἐν τόποις. Même si l'exemple n'est pas probant, Willems n'en a pas moins raison ; Lætitia remarque que l'adjonction de Blaydes accroît la disparité entre les deux stances des vers 663-666 -- là où l'on chercherait à rétablir une symétrie.]) et la locution inverse se lit dans Sophocle : ἥτις μ᾽ ἀποικίσειεν ἐκ τόπων (\italique{Trach}., 955), où jamais personne n'a trouvé à redire (ἐκ τούτων τῶν τόπων.). Je ne vois pas que ces tournures soient plus bizarres que ἐν ποσί et ἐκ ποδῶν, ou bien que ἐν χώρᾳ et ἐκ χώρας, πρὸ ou παρὰ χειρῶν, etc.}%
\\
ξυζωσαμένας εὖ κἀνδρείως τῶν θ᾽ ἱματίων ἀποδύσας%
\commentaire{656}{τῶν θ᾽ ἱματίων ἀποδύσας.}{Willems, fidèle à son poste de \italique{gardien de la tradition}, justifie contre les retoucheurs de texte l'expression ἀποδῦναι $+$ gén. (C'est sa \nieme{20} note aux \italique{Thesmophories}.)}\\
ζητεῖν, εἴ που κἄλλος τις ἀνὴρ ἐπελήλυθ, καὶ περιθρέξαι\apparat{657}{εἰσελήλυθεν R : ἀνελήλυθε Fritzsche : \lecon{ἐπελήλυθε} Handley \espApp \textbf{658} ἀθρῆσαι R : \lecon{διαθρῆσαι} Daubuz}%
\commentaire{657}{εἰσελήλυθεν.}{La leçon du manuscrit de Ravenne est \italique{contra metrum}. Tannegui Le Fèvre (Faber dans l'apparat), le père de Mme Dacier, avait proposé la simple correction ἐσελήλυθε. Austin\marqueur{1} la réfuta : ἐς ne précède jamais une voyelle chez Aristophane. Restent les conjectures de Fritzsche et de Handley, respectivement : ἀνελήλυθε, s'expliquant par une haplographie (ἀνὴρ) ἐλήλυθε suivie d'une complétion postérieure εἰσ- (on reste circonspect devant de telles explications) ; ou bien -- hypothèse beaucoup plus simple et, partant, plus satisfaisante --, ἐπελήλυθε (la confusion entre ΙC et Π en onciales étant une erreur courante).}%
\debutNote{657}%
\note{1}{Voir ses notes : « Textual problems in Ar. \italique{Thesm}. » (1987) ; je reprends la trame de celle qui nous intéresse présentement.}\\
τὴν Πύκνα πᾶσαν καὶ τὰς σκηνὰς καὶ τὰς διόδους διαθρῆσαι.\\\passeLigne}%
\end{groupe}%
\begin{groupe}%
\personnageGrec{Κορυφ. β}%
\typeVers{$4$ troch.\cat}\tetramTroch{εἶα δὴ πρώτιστα μὲν χρὴ κοῦφον ἐξορμᾶν πόδα\testimonia{659}{(εἶα δὴ) Hesychius ε 8 ἔα δή· ἄγε δή $=$ Photius ε  1 $=$ \Syn ε 1 $=$ Suidas ε 3}\apparat{659}{(Ἡμιχ. γυν.) in mg. \Rpc, \Mon}\\
καὶ διασκοπεῖν σιωπῇ πανταχῇ· μόνον δὲ χρὴ\apparat{660}{πανταχῇ post χρὴ R : \lecon{ante μόνον} Küster}\\
μὴ βραδύνειν, ὡς ὁ καιρός ἐστι μὴ μέλλειν ἔτι,\\
ἀλλὰ τὴν πρώτην τρέχειν χρῆν ὡς τάχιστ᾽ ἤδη κύκλῳ.\apparat{662}{χρὴ R : \lecon{χρῆν} Bentley : χρῆ μ᾽ Austin}%
%\commentaire{662}{τὴν πρώτην}{La plupart des traductions et des commentaires donnent une valeur adverbiale à \grec{τὴν πρώτην}. L'interprétation de Willems\marqueur{1} me semble bien préférable ...}%
%\debutNote{662}%
%\note{1}{Cf. sa note 19 reproduite en annexe.}
\\\passeLigne}%
\end{groupe}%
%% v. 663-666
\begin{groupe}%
\personnageGrec{Χορός}%
\typeVers{$5$ troch.}\chantChoeur{\ligne{663}εἶά νυν ἴχνευε καὶ μά-}%
\end{groupe}%
%% v. 663-666
\begin{groupe}%
\decoupeVers{4a}\typeVers{}\chantChoeur{\ligne{663}\alinea τευε ταχὺ πάντ᾽, εἴ τις ἐν τόποις ἑδραῖος}%
\end{groupe}%
%% v. 663-666
\begin{groupe}%
\decoupeVers{b}\typeVers{lecyth.}\chantChoeur{ἄλλος αὖ λέληθεν ὤν.}%
\end{groupe}%
%% v. 663-666
\begin{groupe}%
\typeVers{$2$ troch.}\chantChoeur{πανταχῇ δὲ ῥῖψον ὄμμα,\apparat{665}{διάρριψον R : \lecon{δὲ ῥῖψον} Hermann}}%
\end{groupe}%
%% v. 663-666
\begin{groupe}%
\decoupeVers{a}\typeVers{$2$ troch.}\chantChoeur{καὶ τὰ τῇδε καὶ τὰ δεῦρο\apparat{666}{Post καὶ τὰ τῇδε add. καὶ τὰ κεῖσε Kaibel}}%
\end{groupe}%
%% v. 663-666
\begin{groupe}%
\decoupeVers{b}\typeVers{lecyth.}\chantChoeur{\ligne{666}πάντ᾽ ἀνασκόπει καλῶς.\\\passeLigne}%
\end{groupe}%
%% v. 667-686
\begin{groupe}%
\typeVers{$2$ anap.}\chantChoeur{\ligne{667}ἢν γάρ ληφθῇ δράσας ἀνόσια,\apparat{667}{μὴ λάθῃ R : \lecon{ληφθῇ} Reisig : με λάθῃ Bergk}}%
\end{groupe}%
%% v. 663-666
\begin{groupe}%
\typeVers{$2$ anap.}\chantChoeur{δώσει τε δίκην καὶ πρὸς τούτῳ}%
\end{groupe}%
%% v. 663-666
\begin{groupe}%
\typeVers{$2$ anap\cat}\chantChoeur{τοῖς ἄλλοις ἀνδράσιν ἔσται\apparat{669}{ἅπασιν \inRas{ἅπασιν}{} ἔσται \Rpc : ἅπασιν ἔσται \Mon : \lecon{ἀνδράσιν ἔσται} Beer ap. Hermann (1845) : ἔσται ἅπασιν Bergk : alii alia}}%
\end{groupe}%
%% v. 663-666
\begin{groupe}%
\typeVers{anap.}\chantChoeur{παράδειγμ᾽ ὕβρεως}%
\end{groupe}%
%% v. 663-666
\begin{groupe}%
\typeVers{$2$ anap.}\chantChoeur{ἀδίκων τ᾽ ἔργων ἀθέων τε τρόπων·}%
\end{groupe}%
%% v. 663-666
\begin{groupe}%
\typeVers{$2$ anap.}\chantChoeur{φήσει δ᾽ εἶναί τε θεοὺς φανερῶς,}%
\end{groupe}%
%% v. 663-666
\begin{groupe}%
\typeVers{anap.}\chantChoeur{δείξει τ᾽ ἤδη}%
\end{groupe}%
%% v. 663-666
\begin{groupe}%
\typeVers{$3$ troch\cat}\chantChoeur{πᾶσιν ἀνθρώποις σεβίζειν δαίμονας}%
\end{groupe}%
%% v. 663-666
\begin{groupe}%
\decoupeVers{6}\typeVers{$2$ dochm. ?}\chantChoeur{†δικαίως τ᾽ ἐφέποντας † ὅσια καὶ νόμιμα\apparat{675}{\lecon{ἐφέποντας} R : ἐφέπειν Hermann}}%
\end{groupe}%
%% v. 663-666
\begin{groupe}%
\typeVers{$2$ dochm.}\chantChoeur{\ligne{677}μηδομένους ποιεῖν ὅ τι καλῶς ἔχει.}%
\end{groupe}%
%% v. 663-666
\begin{groupe}%
\typeVers{$3$ iamb\cat}\chantChoeur{κἂν μὴ ποιῶσι ταῦτα τοιάδ᾽ ἔσται·}%
\end{groupe}%
%% v. 663-666
\begin{groupe}%
\typeVers{$2$ iamb...}\chantChoeur{αὐτῶν ὅταν ληφθῇ τις †ὅσια δρῶν†,\apparat{679}{\lecon{αὐτῶν} R : τούτων Richards \espApp \lecon{ὅσια} R : ἀνόσια Zanetti : ὅσια μὴ Burges : ἀνόσιόν τι Hermann}}%
\end{groupe}%
%% v. 680
\begin{groupe}%
\typeVers{$5$ iamb.}\chantChoeur{μανίαις φλέγων λύσσῃ παράκο-
\commentaire{680}{μανίαις φλέγων λύσσῃ παράκοπος,}{Παράκοπος, \italique{délirant}, est une métaphore usée\marqueur{1} et d'origine incertaine. Elle paraît sans rapport avec la frappe de la monnaie\marqueur{2} ; mais elle en a peut-être un avec la notion d'« ébranler le cerveau »\marqueur{3}, cf. Euripide \italique{Hippolyte}, v. 238, etc... παρακόπτειν φρένας. (Note de Taillardat §476.)}%
\debutNote{680}%
\note{1}{Eschyle \italique{Prom.} 581 ; Eur. \italique{Bacch.} 33, 1000, etc.}%
\note{2}{Malgré παρακόπτειν (\italique{Ach.} 517), « battre de la fausse monnaie ».}%
\note{3}{Cf. §474 du même livre : \italique{Les images d'Aristophane} de Jean Taillardat.}}%
\end{groupe}%
%% v. 681-2
\begin{groupe}%
\decoupeVers{2}\typeVers{}\chantChoeur{\alinea πος, εἴ τι δρῴη πᾶσιν ἐμφανὴς ὁρᾶν}%
\end{groupe}%
%% v. 683
\begin{groupe}%
\typeVers{$2$ iamb.}\chantChoeur{\ligne{683}ἔσται γυναιξὶ καὶ βροτοῖς,\apparat{683}{ἐστὶν R : \lecon{ἔσται} Bothe (\italique{erit} Divus) \espApp γυναιξὶν : \lecon{γυναιξὶ} Brunck}%
\commentaire{683}{γυναιξὶ καὶ βροτοῖς,}{L'on s'attendait à la formule consacrée : θεοῖσι καὶ βροτοῖς\marqueur{1}.}%
\debutNote{683}%
\note{1}{Homère, \italique{Iliade} XXIV, 553 ; Pindare, \italique{Pythiques} 3.30 ; Sophocle, \italique{Électre} v. 199.}}%
\end{groupe}%
%% v. 684-5
\begin{groupe}%
\decoupeVers{5}\typeVers{$2$ iamb.}\chantChoeur{ὅτι τά παράνομα τά τ᾽ ἀνόσια παρὼν}%
\end{groupe}%
%% v. 686
\begin{groupe}%
\typeVers{dochm...}\chantChoeur{\ligne{686}θεὸς παραχρῆμ᾽ ἀποτίνεται.\apparat{686}{ἀποτίνεται· / παραχρῆμά τε τίνεται R : \lecon{παραχρῆμ᾽ ἀποτίνεται} Hermann}\\\passeLigne}%
\end{groupe}%
%% v. 687-688
\begin{groupe}%
\personnageGrec{Κορυφαῖος}\typeVers{tetr. tr\cat}\tetramTroch{\ligne{687}ἀλλ᾽ ἔοιχ᾽ ἡμῖν ἅπαντά πως διεσκέφθαι καλῶς.\apparat{687}{Κορ. add. Velsen \espApp \lecon{πως} R : που Blaydes}\\
οὐχ ὁρῶμεν γοῦν ἔτ᾽ ἄλλον οὐδέν᾽ ἐγκαθήμενον\\}%
\end{groupe}%
%% v. 689-691
\begin{groupe}%
\ledpb%
\personnageGrec{Μίκα}%
\typeVers{extr. metr.}\trimIamb{\ligne{689}ἆ ἆ}%
\end{groupe}%
%% v. 689-691
\begin{groupe}%
%\personnageGrec{Μίκα}%
\typeVers{trim. iamb.}\trimIamb{\ligne{689}ποῖ ποῖ σὺ φεύγεις; οὗτος οὗτος οὐ μενεῖς;\apparat{689}{ἀά supra lin. R : ἔα Farraeus : ἆ ab init. v. 689 Bergler : \lecon{ἆ ἆ} Thiersch \espApp ποῖ R : \lecon{ποῖ ποῖ} Biset}%
\commentaire{689-761}{Parodie du \italique{Télèphe} d'Euripide.}{Comme dans ses \italique{Acharniens} -- représentés quatorze ans plus tôt et qui comportent nombre de similarités avec les \italique{Thesmophories}--, Aristophane détourne la pièce d'Euripide. On ne sait pas grand chose de cette dernière, la majorité des fragments conservés étant eux-mêmes tirés d'Aristophane et de ses anciens commentateurs. Je renvoie à l'édition de François Jouan \& Herman van Looy\marqueur{1} pour une présentation détaillée du \italique{Télèphe}.
Télèphe était un héros populaire puisque les trois grands tragiques lui ont consacré pas moins de sept pièces ; il fut blessé par un coup de lance d'Achille.
Comme sa blessure ne guérissait pas, il finit par consulter l'oracle d'Apollon, réponse de ce dernier : ὁ τρώσας ἰάσεται, « Qui t'a blessé te guérira ». Il se rend donc, déguisé en mendiant, dans le camp de l'armée achéenne préparant son expédition sur Troie. Allait-il d'abord trouver Clytemnestre qui lui confiait son petit Oreste pour appuyer son plaidoyer ? toujours est-il que Télèphe se retrouva la proie exposée à la colère des chefs grecs, il se réfugiait alors en suppliant sur l'autel en prenant en otage Oreste. Le sort de l'enfant était plus heureux que dans sa thesmophorienne parodie : grâce à l'entremise d'Agamemnon, Achille se laissait convaincre de soigner Télèphe et tout était bien qui finissait bien. Outre cette scène d'enlèvement du petit Oreste, de nombreux autres passages des \italique{Thesmophories} sont sans doute aussi inspirés du \italique{Télèphe} (v. 39-70 : où l'altercation entre le héros et un gardien attirait Clytemnestre à l'extérieur et permettait sa rencontre avec Télèphe ; v. 466-470 : lorsqu'au début de son plaidoyer, le héros -- sous une fausse identité -- usait de précautions oratoires ; v. 598-600, 659-667 avec l'annonce de la présence d'un espion dans le camp grec.)}%
\debutNote{689-761}%
\note{1}{Aux Belles Lettres, tome VIII, \nieme{3} partie.}\\
τάλαιν᾽ ἐγὼ τάλαινα, καὶ τὸ παιδίον\\
ἐξαρπάσας μοι φροῦδος ἀπὸ τοῦ τιτθίου.\\\passeLigne}%
\end{groupe}%
%% v. 692-695a
\begin{groupe}%
\personnageGrec{Μνησίλοχος}\trimIamb{κέκραχθι· τοῦτο δ᾽ οὐδέποτε σὺ ψωμιεῖς,\\
ἢν μή μ᾽ ἀφῆτ᾽· ἀλλ᾽ ἐνθάδ᾽ ἐπὶ τῶν μηρίων%
\apparat{693}{ἀφε\auDessus{ῖ}{ῆ}τ᾽ \Rdeux : ἀφεῖτ᾽ \Mon}%
\commentaire{693-5}{ἐπὶ τῶν μηρίων ... βωμόν.}{Vers tragiques, très sûrement empruntés à Euripide. (Note de Willems.)}\\
πληγὲν μαχαίρᾳ τῇδε φοινίας φλέβας\\
καθαιματώσει βωμόν.\\\passeLigne}%
\end{groupe}%
%% v. 695b-698
\begin{groupe}%
\personnageGrec{Μίκα}%
\trimIamb{\ligne{695}\decalage{καθαιματώσει βωμόν.}ὢ τάλαιν᾽ ἐγώ.\\
γυναῖκες, οὐκ ἀρήξετ᾽; οὐ πολλὴν βοὴν\testimonia{696-7}{(γυναῖκες -- τροπαῖον) \GregCor pp. 20-22 ; cf. scholium\expo{Ald} \italique{Plutus} 453c}\commentaire{696-7}{οὐ πολλὴν βοὴν / στήσεσθε καὶ τροπαῖον.}{Dans le sens de μήστωρ φόβοιο d'Homère. « Un vacarme à faire sauver ». (Musset, \italique{Confess}., p. 111 ; note de Willems.)}\\
στήσεσθε καὶ τροπαῖον, ἀλλὰ τοῦ μόνου\apparat{697}{καὶ om. R : καὶ suppl. Schol. ad \italique{Pl.} v.453 et Greg. Cor. \italique{de dialectis} p.22}\\
τέκνου με περιόψεσθ᾽ ἀποστερουμένην;\\\passeLigne}%
\end{groupe}%
%% v. 699-701
\begin{groupe}%
\ledpb%
\personnageGrec{Χορός}\typeVers{προαναφώνημα}\chantChoeur{\decalage{aa}ἔα ἔα.%\commentaire{699-725}{Chant amébée.}{\eff{Retrouver dans Willems.}}
}%
\end{groupe}%
%% v. 699-701
\begin{groupe}%
\typeVers{2 dochm.}\chantChoeur{ὦ πότνιαι Μοῖραι τί τόδε δέρκομαι\apparat{700}{δὲ R : \lecon{τόδε} Dobree : δὴ Reisig}}%
\end{groupe}%
%% v. 699-701
\begin{groupe}%
\typeVers{dochm.}\chantChoeur{νεοχμὸν αὖ τέρας;\\\passeLigne}%
\end{groupe}%
%% v. 702-703
\begin{groupe}%
\typeVers{$4$ troch\cat}\testimonia{702}{(ἅπαντ' -- κἀναισχυντίας) Suidas α 2902 (sed ἅπαντα pro ἅπαν γὰρ ἐστὶν et πλέα pro ἔργα) ; \italique{Corpus paroemiographorum Graecorum}, Appendices, I, 35 (sed ἅπαντα pro ἅπαν γὰρ ἐστὶν et ἔργα omisso) ; \italique{Laurentiamus} LVIII 24 p. 40 (ed. Cohn) (sed ἅπαντα pro ἅπαν γὰρ ἐστὶν et δοῦλα pro ἔργα)}\tetramTroch{ὡς ἅπαντ᾽ ἄρ᾽ ἐστὶ τόλμης διπλὰ κἀναισχυντίας.\apparat{702}{ἅπαν γὰρ ἐστὶν R : ἅπαντα Suid. α 2902, cod. Laurent. LVIII 24 : \lecon{ἅπαντ᾽ ἄρ᾽ ἐστὶ} Bentley \espApp ἔργα R : πλέα Suid. : δοῦλα cod. Laur. : \lecon{διπλὰ} Bibi}%
\commentaire{702}{ὡς ἅπαν γὰρ ἐστὶν τόλμης ἔργα κἀναισχυντίας.}{Le scholiaste met en avant le tour proverbial de la formule en commentant d'un sobre «~παροιμία~» ; à nous donc de mener une enquête plus approfondie !\\
Pour ce faire, commençons par recenser quelques autres attestations de l'expression (ou de ses proches cousines). Et tout d'abord, remarquons que les termes τόλμη et ἀναισχυντία sont souvent mis en collocation : outre le présent passage, Platon\marqueur{1} et Démosthène\marqueur{2} associent les deux mots avec un sens analogue. Isocrate nous apporte quelque précision, l'on comprend en le lisant\marqueur{3} que l'audace (τόλμα) pouvait être interprétée comme un signe (σημεῖον) d'impudence (ἀναισχυντία). Les parœmiographes n'ont pas manqué de la relever, avec de légères variations. Les voici :
Suidas (α 2902) : 
\citComm{ἅπαντα τόλμης πλέα κἀναισχυντίας: ἐπὶ τῶν ἰταμῶν καὶ ἀναιδῶν}
Dans le \italique{Corpus paroemiographorum Graecorum}, on trouve\marqueur{4} :
\citComm{ἅπαντα τόλμης καὶ ἀναισχυντίας: ὅτι ἡ ἀναισχυντία νικᾷ τὰ πολλά.}
Si l'on consulte le codex \italique{Laurentiamus} LVIII 24\marqueur{5}, l'on tombe sur la formulation suivante :%
\citComm{ἅπαντα τόλμης δοῦλα καὶ ἀναισχυντίας.}%
Ces trois versions pourraient-elles découler d'une même source ? (savoir : l'\allemand{Urtext} des \italique{Thesmophories}.)
Il me semble que oui... \latin{sed quod gratis asseritur gratis negatur}. Aussi nous faut-il présenter un embryon de `démonstration' ; ce qui devrait infailliblement nous conduire à formuler d'aventureuses conjectures. Partons donc à l'aventure ! \\%(Je prie le lecteur rigoriste de ne pas trop m'en tenir rigueur.)\\
La leçon du \italique{Ravennas}, ὡς ἅπαν γὰρ ... ἔργα κτλ., est fautive. La correction la plus simple à apporter est de lire : ὡς ἅπαντ᾽ ἄρ᾽ κτλ.\marqueur{6} Il n'y aurait alors qu'une simple confusion d'un Τ lu en Γ : ΑΠΑΝΤΑΡ $ \to $ ΑΠΑΝΓΑΡ.
Cette correction fait concorder le texte d'Aristophane avec les parémiographes sur la présence d'ἅπαντα. Les variantes ne différent plus que sur un seul mot : ἔργα (dans le \italique{Ravennas}) ; πλέα (\italique{Souda}) ; δοῦλα (Laurent.) ; $\emptyset$ (\italique{Corpus paroemiographorum}).
À présent, si l'on fait l'hypothèse -- avec des `si', l'on scierait le texte entier !.. -- si donc l'on fait l'hypothèse que le texte original d'Aristophane contenait à cette place : διπλά, alors les différentes leçons pourraient s'expliquer.
Très tôt, un manuscrit aura porté un δ surmonté de deux points pour abréger le δι\marqueur{7} de διπλά. L'omission des deux points donne δπλα ; la chute du δ pouvait alors entraîner la lecture πλέα tandis qu'un π cursif ressemblant au digramme ου faisait écrire à un scribe δοῦλα. 
%Si maintenant, sur une autre branche, un autre archétype voyait sa leçon ΔΙΠΛΑ réduite en ΔΙΠΑ (chute classique lorsque l'on se trouve en présence de deux lettres consécutives identiques ou ressemblantes), le Δ lu en Ε, le Ι en Ρ et le Π en Γ (toutes erreurs fréquentes prises individuellement, peut-être induites en l'occurrence par l'ἔργον du vers suivant), cette tradition aurait bien pu donner notre leçon ἔργα. 
\citComm{ὡς ἅπαντ᾽ ἄρ᾽ ἐστὶ τόλμης διπλὰ κἀναισχυντίας.\\
Ô combien toutes [les actions] sont doublées d'audace et d'impudence.}%
L'on s'aperçoit à la traduction qu'on ressent l'envie de gloser ἅπαντα par ἔργα. Mettons-nous alors à la place d'un scribe ayant devant lui la leçon déjà corrompue de διπλὰ\marqueur{8}. Devant donc cette difficulté surmontée de la glose ἔργα, il aura choisi la \latin{lectio facilior} ; c'est-à-dire recopier ἔργα. C'est ainsi que je m'explique sa présence dans le \italique{Ravennas}.\\
La prudence m'avait dans un premier temps incité à circonscrire cette conjecture à l'apparat critique. Las ! l'esprit d'imagination se laisse si facilement séduire qu'à la relecture διπλά rejoignit l'\allemand{Urtext}.
%\eff{Il ne me semble pas impossible que ces trois versions du proverbe -- avec ἔργα (notre texte, le \italique{Ravennas}), πλέα (\italique{Souda}) et δοῦλα (cod. Laur.) -- découlent d'une même source. Qu'il me soit permis de céder à la tentation des aventureuses conjectures, en espérant qu'une fois ne soit pas coutume. 
%Supposez, pour la démonstration, que le texte original d'Aristophane contenait διπλά. Très tôt, un manuscrit aurait porté un δ surmonté de deux points pour abréger ce δι (cf. le vers 122 et la correction proposée : δινεύματα). L'omission des deux points donne δπλα ; la chute du δ pouvait alors entraîner la lecture πλέα tandis qu'un π cursif ressemblant à ου faisait écrire à un scribe δοῦλα. Si maintenant, sur une autre branche, un autre archétype voyait sa leçon ΔΙΠΛΑ réduite en ΔΙΠΑ (chute classique lorsque l'on se trouve en présence de deux lettres consécutives identiques ou ressemblantes), le Δ lu en Ε, le Ι en Ρ et le Π en Γ (toutes erreurs fréquentes prises individuellement, peut-être induites en l'occurrence par l'ἔργον du vers suivant), cette tradition aurait bien pu donner notre leçon ἔργα. 
%\citComm{ὡς ἅπαντ᾽ ἄρ᾽ ἐστὶ τόλμης διπλὰ κἀναισχυντίας.\\
%Ô combien vont de pair audace et impudence.}
%Restons néanmoins prudent : si séduisante soit-elle pour l'esprit d'imagination, cette conjecture n'en restera pas moins circonscrite à l'apparat critique.\\
Concluons cette note en digressant légèrement sur un détail amusant et plein de sens : les Grecs avaient fait d'\italique{Implacabilité} un dieu : θεὸς ἡ Ἀναίδεια\marqueur{9}.%
%P.-S. : faisons aussi remarquer au lecteur que l'émendation de Bentley -- ἅπαντ᾽ ἄρ᾽ ἐστὶ à la place de ἅπαν γὰρ ἐστὶ ; un simple Τ lu en Γ -- ne se justifie que si le scribe recopiait un manuscrit déjà élidé.}
%\eff{C'est le verbe apantaw ?
%\italique{Ion} (940) :
%\citComm{Πρεσβύτης -- Τίν'; Ὡς ἀπαντᾷ δάκρυά μοι τοῖς σοῖς λόγοις.}
%Ménandre 706(483) K.-A.
%Loci communes (Capita theologica id est scite dicta atque electa, ex diversis, tum Christianorum, tum profanorum libris (PG 91, 721–1018).)
%κεφάλαια θεολογικὰ ἤτοι ἐκλογαὶ
%ΛΟΓΟΣ Β´.
%Περί φρονήσεως καί βουλῆς. 
%Μενάνδρου. Ἅπαντα δοῦλα τοῦ φρονεῖν καθίσταται.
%Βουλευόμενοι παραδείγματα ποιοῦνται παρεληλυθότα τῶν μελλόντων. Τό γάρ ἀφανές, ἐκ τοῦ φανεροῦ ταχίστην ἔχει τήν διάγνωσιν. 
%"Loci communes" du Pseudo-Maxime le Confesseur. Éditions :
%Florilège sacro-profane du Pseudo-Maxime (2001), Sargologos, Étienne
%IHM, Sibylle : "Ps.-Maximus Confessor. Erste kritische Edition
%einer Redaktion des sacro-profanen Florilegiums Loci Communes" (2001)}
}%
\debutNote{702}%
\note{1}{Platon, \italique{Apologie de Socrate}, 38d. Voir aussi son \italique{Théétète} (196d) où le tour proverbial est encore plus sensible : ὅμως δέ — πάντα γὰρ τολμητέον — τί εἰ ἐπιχειρήσαιμεν ἀναισχυντεῖν; «~De l'audace ! encore de l'audace ! toujours de l'audace !~»}%
\note{2}{Παραγραφὴ πρὸς Ναυσίμακον καὶ Ξενοπείθην, 5 (Νόμος).}%
\note{3}{Περὶ ἀντιδόσεως, 190 : ἔτι δὲ τὴν τόλμαν μὴ τὴν ἀναισχυντίας σημεῖον γιγνομένην, ἀλλὰ τὴν μετὰ σωφροσύνης.}%
\note{4}{\italique{Corpus paroemiographorum Graecorum}, vol. I, Appendices, I, 35, p. 384 ($=$ Vat. K II 36).}%
\note{5}{Il s'agit d'un manuscrit conservé à la Biblioteca Medicea Laurenziana (BML) de Florence, au dépôt Plutarque, sous la cote 58. 24. (J'en profite pour signaler au lecteur un outil sur l'internet fort pratique pour tout ce qui touche aux manuscrits : Πίνακες. Il y trouvera, par exemples, l'historique de notre manuscrit ainsi que sa numérisation.) Je me reporte à l'édition qu'en fit Leopold Cohn en 1887 : \italique{Zu den Paroemoigraphen} (notre proverbe y figure à la page 40).}%
\note{6}{Cette émendation, proposée par Bentley et unanimement adoptée se justifierait d'autant mieux si le scribe du \italique{Ravennas} recopiait un manuscrit déjà élidé -- du moins partiellement, le modèle restant majoritairement à \latin{scriptio plena}.}%
\note{7}{Cf. l'article de Victor Coulon, «~Notes critiques et exégétiques sur Aristote, \italique{Poétique} , 1450a12, 1450a29, sur Thucydide I, 1, 2 et sur cinq passages d'Aristophane~», in \italique{Revue des Études Grecques} (1937), t. 50 fasc. 238, p. 458. (À propos du vers 122 et de la correction de Bentley δινεύματα à la place de διανεύματα R.)}%
\note{8}{Et donc incompréhensible ; cf. la version courte ἅπαντα τόλμης καὶ ἀναισχυντίας citée plus haut : le scribe du \italique{Corpus} a tout simplement décidé de ne pas recopier le mot qu'il n'arrivait pas à lire.}%
\note{9}{\italique{Implacabilité} et non \italique{Impudence}. Cf. le \italique{Dictionnaire des antiquités grecques et romaines} (Daremberg \& Saglio), t. I, p. 398b (\latin{Areopagus}). Les Athéniens avaient dressé des autels à Ὕβρις et Ἀναίδεια dans l'Aréopage, contre lesquels s'était insurgé ce rabat-joie de Cicéron -- non sans faire un contre-sens comme le souligne l'auteur de l'article. Dans sa querelle avec Agamemnon, l'attitude d'Achille -- qui par bien des aspects est celle d'un adolescent attardé -- nous révèle en fait son côté divin : il est \italique{implacable}, tout autant que le destin qu'il connaît et a accepté. Pour revenir au proverbe `θεὸς ἡ Ἀναίδεια', certainement inspiré par ces autels de l'Aréopage, voir Zénobius (IV, 36). Érasme l'a aussi recensé dans ses \italique{Adages} (n° 2063).}\\
οἷον αὖ δέδρακεν ἔργον, οἷον αὖ φίλαι τόδε.\\\passeLigne}%
\end{groupe}%
%% v. 704
\begin{groupe}%
\personnageGrec{Μνησίλοχος}%
\tetramTroch{οἷον; ὑμῶν ἐξαρήξω τὴν ἄγαν αὐθαδίαν.\apparat{704}{οἷον R : \lecon{οἷον;} Ellebodius \espApp ἐξάρξω R : ἐξαράξει Ellebodius : ἐξαράξω Bentley : \lecon{ἐξαρήξω} Bibi}%
\commentaire{704}{οἷον ὑμῶν ἐξάρξω τὴν ἄγαν αὐθαδίαν.}{Il manque une syllabe à ἐξάρξω, la leçon du manuscrit de Ravenne. Les différents éditeurs y suppléent avec le verbe ἐξαράττειν (ἐξαράξω au futur ou, en s'éloignant davantage de la leçon -- quitte à corriger le texte... --, ἐξαράξει, pour faire de οἷον le sujet). La correction est légère et pourtant elle ne manque pas de laisser dubitatif : si l'image de \italique{briser} (ἐξαράξω) l'\italique{audace} (αὐθαδίαν) de \italique{quelqu'un} (ὑμῶν) ne paraît aucunement étrangère à une imagination française, il n'en allait peut-être pas de même pour les Grecs. L'incontournable Taillardat\marqueur{1} l'avait déjà noté : considérer l'arrogance «~comme un objet matériel qu'on peut briser [...] est une image unique~». En revanche, en grec comme en français, l'image de \italique{chasser}, \italique{pousser} hors (ἐξ-), \italique{repousser} un sentiment ou un trait de caractère est courante. Ainsi, l'on chassera l'audace par l'audace\marqueur{2} :%
\citComm{Ἀντιφάνης δέ· οἴνῳ .... τὸν οἶνον ἐξελαύνειν, [...] τριωβόλῳ δὲ πόρνην, αὐθαδίαν αὐθαδίᾳ.}%
L'on se débarrassera de son chagrin\marqueur{3} :%
\citComm{ἐκκρούει τὴν λύπην}%
comme les vieillards cagneux secouent leurs ennuis chez Aristophane\marqueur{4} :%
\citComm{ἀποσείονται δὲ λύπας.}%
Euripide met dans la bouche du fourbe Polydectès qui veut s'accaparer Danaé que\marqueur{5} :%
\citComm{πατέρα τε παισὶν ἡδέως συνεκφέρειν\\
φίλους ἔρωτας \textbf{ἐκβαλόντ' αὐθαδίαν}.\\
Un père doit supporter avec indulgence les affaires amoureuses de son fils et \italique{repousser l'entêtement}.}
L'occasion fait le larron : profitons-en pour placer cet aphorisme rapporté par Athénée\marqueur{6} illustrant le sens de la mesure grec, même quand il s'agit de boire...%
\citComm{Πάσας δ᾽ ἐκ κραδίας ἀνίας ἀνδρῶν ἀλαπάζει\\
πινόμενος κατὰ μέτρον· ὑπὲρ μέτρον δὲ χερείων.\\
\francais{Le vin est souverain pour chasser les soucis du cœur de qui n'en boit pas trop, mais l'excès nuit.}}%
(D'aucun franc-buveur espiègle répondrait que la mesure du boire est de boire sans mesure\marqueur{7} !..)\\
Le verbe ἀρήγω (futur : ἀρήξω) signifie \italique{secourir} et, par suite, \italique{écarter}, \italique{repousser}. Il conviendrait tant pour le sens que paléographiquement s'il était préfixé de ἐξ-.
Comme pour nous conforter, Hésychius (ε 3599) nous atteste justement l'existence du composé :%
\citComm{ἐξαρῆξαι· ἐκφορῆσαι.}%
Lisons donc :%
\citComm{ὑμῶν ἐξαρήξω τὴν ἄγαν αὐθαδίαν\\
Je vous extirperai votre trop-plein d'audace.}
P.-S. : l'idée de recourir aux matrices métaphoriques pour le bon établissement d'un texte revient à l'immense Jean Taillardat. Je ne peux que renvoyer à son coruscant article « Images et matrices métaphoriques »\marqueur{8}. Il le termine par ce vœu :%
\citComm{En guise de conclusion, je formulerai un souhait, celui qu'on puisse, un jour, consulter un dictionnaire des matrices métaphoriques pour les langues indo-européennes vivantes (langues populaires et argots compris) et mortes. On l'imaginerait volontiers comme une sorte de dictionnaire analogique, mais réduit aux images, les notions à exprimer (par exemple « colère », « inconsistance », « jeunesse », etc.) étant classées en ordre alphabétique ; quel merveilleux outil serait ainsi offert à qui s'intéresse à la recherche étymologique et combien de tâtonnements, d'erreurs même, seraient alors évités ! Mais quelle équipe aura assez d'abnégation et de courage pour se lancer jamais dans cette énorme entreprise ?}%
Puisse ce grand et beau dictionnaire\marqueur{9} un jour voir le jour !%
}%
\debutNote{704}%
\note{1}{§628 de ses \italique{Images d'Aristophane}, ouvrage dont on ne soulignera jamais assez l'importance pour toute personne désireuse d'avoir un avant-goût du génie de la langue aristophanesque.}%
\note{2}{Athénée, \italique{Deipnosophistes}, II, 44a.}%
\note{3}{Aristote, \italique{Éthique à Nicomaque, 1154a27}.}%
\note{4}{\italique{Grenouilles}, v. 346.}%
\note{5}{Fragment du \italique{Dictys}. Stobée, IV, 26, 16 $=$ fr. 7 \italique{Dictys} Jouan \& van Looy $=$ fr. 339 Kannicht $=$ fr. 433 Mette.}
\note{6}{Athénée, II, 37b.}%
\note{7}{Aristophane, \italique{Guêpes}, vers 79-80 :%
\citComm{εἶναι φιλοπότην αὐτόν. [ ... ]\\
αὕτη γε χρηστῶν ἐστιν ἀνδρῶν ἡ νόσος.}.}%
\note{8}{\italique{Bulletin de l'Association Guillaume Budé} (décembre 1977), Lettre d'humanité n° 36, pp. 344-354.}%
\note{9}{Dictionnaire dont le grand'œuvre de Taillardat, ses \italique{Images d'Aristophane}, en est la première pierre... et pourrait en être la pierre d'angle !}%
\\\passeLigne}%
\end{groupe}%
%% v. 705
\begin{groupe}%
\personnageGrec{Χορός}%
\tetramTroch{ταῦτα δῆτ᾽ οὐ δεινὰ πράγματ᾽ ἐστὶ καὶ περαιτέρω;\\\passeLigne}%
\end{groupe}%
%% v. 706
\begin{groupe}%
\personnageGrec{Μίκα}%
\tetramTroch{δεινὰ δῆθ᾽, ὁτιή γ᾽ ἔχει μου 'ξαρπάσας τὸ παιδίον.\apparat{706}{ὅτι R : \lecon{ὁτιή} Faber : ὅστις Porson}\\\passeLigne}%
\end{groupe}%
%% v. 707-708
\begin{groupe}%
\personnageGrec{Χορός}%
\typeVers{dim. anap.}\dimAnap{τί ἂν οὖν εἴποι πρὸς ταῦτά τις, ὅτε\\
τοιαῦτα ποιῶν ὅδ᾽ ἀναισχυντεῖ;%
\commentaire{707-25}{}{Ce passage, sans correspondre parfaitement du point de vue métrique à celui des vers 667-686, présente néanmoins de troublantes ressemblances avec ce dernier. Se reporter à Lætitia\marqueur{1} pour les parallélismes qu'on peut faire avec le texte tel qu'il nous a été conservé.}%
\debutNote{707-25}%
\note{1}{\italique{The songs of Aristophanes}, pp. 421 sqq. Je reprends pour une bonne part sa colométrie.}\\\passeLigne}%
\end{groupe}%
%% v. 709
\begin{groupe}%
\personnageGrec{Μνησίλοχος}%
\typeVers{paroem.}\dimAnap{κοὔπω μέντοι γε πέπαυμαι.\\\passeLigne}%
\end{groupe}%
%% v. 710
\begin{groupe}%
\personnageGrec{Χορός}%
\typeVers{paroem.}\chantChoeur{ἀλλ᾽ οὖν ἥκειν γ᾽ ὅθεν ἥκεις.\testimonia{710-3}{Suidas η 170 (sed λήψῃ pro λήψει)}\testimonia{710}{Photius η 88 ἥκειν· ἀντὶ τοῦ ἥκων}\apparat{710}{ἥκεις τ᾿ R : ἥκεις γ᾽ Suid. η 170 : ἥκων γ᾽ Willems : \lecon{ἥκειν γ᾽} Bibi ap. Phot. η 88 \ligneApp{711} τ᾿ R, Suid : \lecon{δ᾿} Bentley \espApp οὐ λέξεις R, Suid. : \lecon{οὐκ ἔξει.} Bibi ap. Schol.}%
\commentaire{710-2}{ἀλλ᾽ οὖν ἥκεις τ᾿ ὅθεν ἥκεις.}{Je reproduis ci-dessous la \nieme{21} note d'Alphonse Willems. La correction qu'il y propose fut sauvée du néant de l'oubli et adoptée par la plupart des éditeurs modernes grâce à la diligente entremise de Coulon :

«~Ce passage a été refait de cinq ou six façons, plus bizarres les unes que les autres. 
La correction me paraît élémentaire. Les manuscrits concordent, sauf sur le premier ἥκεις (ἥκεις γ' dans Suidas, ἥκεις τ' dans le \italique{Ravennas}, ἥκισθ' dans le manuscrit de Munich [qui n'est qu'une copie du manuscrit de Ravenne, comportant quelques émendations parfois bien senties]). Lisez ἥκων γ', le sens devient parfait :
\citComm{ἀλλ᾽ οὖν ἥκων γ᾽ ὅθεν ἥκεις}
Il est à peine besoin d'ajouter que ἥκων est pris dans le sens de ἀνήκων, \italique{redux}, ce qui est des plus fréquents\marqueur{1}. Et quant à ἀλλ᾽ οὖν..γε, chacun sait qu'il signifie \italique{toujours}, \italique{tout au moins}, \italique{en tous cas}\marqueur{2}.\\
%\eff{Traduisez : 
%\citComm{\francais{Toujours tu n'iras pas, revenu au lieu dont tu es parti, après une fuite facile, raconter ton crime et ton évasion...}}}
% [...]\marqueur{3}\\
C'est le mot d'Épicharme, fr. 126\marqueur{3} : κἀπῆνθεν ὅθεν ἦνθεν πάλιν ; ou de Ménandre (t. IV, p. 211 [$=$ fr. 373 K.-A.]) : ἀπῆλθεν ὅθεν ἦλθεν ταχύ. Cf. aussi Soph., \italique{Œd. Col.}, 1227 : βῆναι κεῖσ' ὁπόθεν περ ἥκει ; Eur., \italique{Her. F.}, 1247 : θανών, ὅθενπερ ἦλθον, εἶμι γῆς ὕπο, etc.~»

Le chœur enjoint Mnésiloque de cesser cette escalade de l'impudence : qu'il quitte au moins l'autel et les rejoigne\marqueur{4}.\\
À la lumière de la démonstration de Willems, je me permets d'ajouter la glose suivante du \latin{Lexicon} de Photius (η 88) :%
\citComm{ἥκειν· ἀντὶ τοῦ ἥκων.}%
Supposer que la glose de Photius se rapporte à notre vers offre l'avantage de rester plus proche du texte des manuscrits : il y aurait eu simple permutation d'un ν par un ς.
L'on peut donc lire après Willems :%
\citComm{ἀλλ᾽ οὖν ἥκειν γ᾽ ὅθεν ἥκεις.\\
\francais{Au moins, reviens d'où tu es parti.}}%
}%
\debutNote{710-2}%
\note{1}{\italique{Paix}, 265, 275; \italique{Lysistr}., 595, 731, 924 ; \italique{Gren}., 503 ; \italique{Ploutos}, 521, 1179, etc.}%
\note{2}{\italique{Nuées}, 985, 1002 ; \italique{Guêpes}, 1129, 1190 ; Esch., \italique{Prom}., 1058 ; Soph., \italique{El}., 233 ; \italique{Ant}., 84, etc. M'éloignant de Willems quant au sens à donner à ce vers, je ne retranscris pas ici la traduction qu'il en propose. Le lecteur la trouvera en annexe où toutes les notes d'Alphonse Willems sur les \italique{Thesmophories} sont intégralement reproduites.}%
\note{3}{$[=$ Plutarque, \italique{Moralia}, Παραμυθητικὸς πρὸς Ἀπολλώνιον, 110a $=$ fr. 213(245) Epich. Kassel-Austin, \italique{Poet. Com. Gr.}, I, \italique{Comoedia Dorica}.$]$}%
\note{4}{C'est l'explication de la scholie aussi conservée par Suidas (η 170) : ἥκεις φησίν εἰς τὰς χεῖρας ἡμῶν, ὅθεν οὐ φεύξῃ.}%
}%
\end{groupe}%
%% v. 711-713
\begin{groupe}%
\typeVers{paroem.}\chantChoeur{φαύλως δ᾿ ἀποδρὰς οὐκ ἔξει.%
\commentaire{711-2}{φαύλως τ᾽ ἀποδρὰς οὐ λέξεις / οἷον δράσας διέδυς ἔργον.}{%. ARISTOPHANES Comic. Acharnenses {0019.001} Line 215
%εἴ τις οἶδ’ ὅποι τέτραπται γῆς ὁ τὰς σπονδὰς φέρων.
%Ἐκπέφευγ’, οἴχεται φροῦδος. Οἴ-    (208-209)
%μοι τάλας τῶν ἐτῶν τῶν ἐμῶν·   (210)
%οὐκ ἂν ἐπ’ ἐμῆς γε νεότητος, ὅτ’ ἐ-   (211-212)
%γὼ φέρων ἀνθράκων φορτίον    (213-214)
%ἠκολούθουν Φαΰλλῳ τρέχων, ὧδε φαύλως ἂν ὁ    (215)
%σπονδοφόρος οὗτος ὑπ’ ἐμοῦ τότε διωκόμενος    (216-217)
%ἐξέφυγεν οὐδ’ ἂν ἐλαφρῶς ἂν ἀπεπλίξατο.    (218)
%
%νῦν δ᾽ ἐπειδὴ στερρὸν ἤδη τοὐμὸν ἀντικνήμιον,
%220καὶ παλαιῷ Λακρατείδῃ τὸ σκέλος βαρύνεται,
%οἴχεται. διωκτέος δέ: μὴ γὰρ ἐγχάνῃ ποτὲ
%μηδέ περ γέροντας ὄντας ἐκφυγὼν Ἀχαρνέας.
%
%ὅστις ὦ Ζεῦ πάτερ καὶ θεοὶ τοῖσιν ἐχθροῖσιν ἐσπείσατο,
%οἷσι παρ᾽ ἐμοῦ πόλεμος ἐχθοδοπὸς αὔξεται τῶν ἐμῶν χωρίων:
%230κοὐκ ἀνήσω πρὶν ἂν σχοῖνος αὐτοῖσιν ἀντεμπαγῶ
%ὀξὺς ὀδυνηρὸς ... ἐπίκωπος, ἵνα
%μήποτε πατῶσιν ἔτι τὰς ἐμὰς ἀμπέλους.
%
%ἀλλὰ δεῖ ζητεῖν τὸν ἄνδρα καὶ βλέπειν βαλληνάδε
%235καὶ διώκειν γῆν πρὸ γῆς, ἕως ἂν εὑρεθῇ ποτέ:
%ὡς ἐγὼ βάλλων ἐκεῖνον οὐκ ἂν ἐμπλῄμην λίθοις.
%
%
%
%
%
%Βδ.        ἀπολεῖς με. ποίας χάρακας; ἀλλ’ ὡς ἢ κάπρον
%ἐδιώκαθές ποτ’ ἢ λαγών, ἢ λαμπάδα
%ἔδραμες, ἀνευρὼν ὅ τι νεανικώτατον.
%Φι.        ἐγᾦδα τοίνυν τό γε νεανικώτατον· @1   (1205)
%ὅτε τὸν δρομέα Φάυλλον ὢν βούπαις ἔτι
%εἷλον διώκων λοιδορίας ψήφοιν δυοῖν.
%Βδ.        παῦ’· ἀλλὰ δευρὶ κατακλινεὶς προσμάνθανε
%ξυμποτικὸς εἶναι καὶ ξυνουσιαστικός.
%Φι.        πῶς οὖν κατακλινῶ; φράζ’ ἁνύσας.    (1210)
%Βδ.                                        εὐσχημόνως.   (1210)
%
%
%
%faulws logizw : vers 656
% Φι.                                        παῦσαι καὶ μὴ πατέριζε.   (652)
%εἰ μὴ γάρ, ὅπως δουλεύω ’γώ, τουτὶ ταχέως με διδάξεις,
%οὐκ ἔστιν ὅπως οὐχὶ τεθνήσει, κἂν χρῇ σπλάγχνων μ’
%ἀπέχεσθαι.    (654)
%Βδ.        ἀκρόασαί νυν, ὦ παπίδιον, χαλάσας ὀλίγον τὸ μέτωπον. @1   (655)
%καὶ πρῶτον μὲν λόγισαι φαύλως, μὴ ψήφοις ἀλλ’ ἀπὸ
%χειρός,    (656)
%τὸν φόρον ἡμῖν ἀπὸ τῶν πόλεων συλλήβδην τὸν προσιόντα,
%κἄξω τούτου τὰ τέλη χωρὶς καὶ τὰς πολλὰς ἑκατοστάς,
%πρυτανεῖα, μέταλλ’, ἀγοράς, λιμένας, μισθώσεις, δημιό-
%πρατα·    (659)
%
%
%REFAIRE NOTE !!!!! Il n'y a pas lieu de corriger le texte ! (même si les précédents éditeurs disaient effectivement n'importe quoi. Il s'agit de Phaüllos ! un célèbre coureur cf Acharniens 215 , Guêpes 1206 (voir Dargemberg et Saglio, Realencyclopädie). Traduire : "Tu ne diras pas avoir décampé à la manière de Phaüllos" ?)
% Voir Nuées 167-168
Où se trouve le complément réclamé par λέγειν ? Je ne suis pas sûr qu'il soit de bonne grécité de construire la phrase ainsi que tous les éditeurs : λέξεις οἷον δράσας, «~tu raconteras tout ce que tu as fait~». (Sauf erreur de ma part, le grec aurait recouru à un verbe conjugué et non au participe de δράω.)
%Une retouche du texte -- très légère, accordez-moi au moins cela avant de crier au pourfendeur des éditeurs-tripatouilleurs pris lui-même en flagrant délit de bidouillage --, 
Une légère retouche du texte permettrait une bien meilleure lecture, tant pour la syntaxe grecque qu'au niveau du sens.
Lisez : 
\citComm{φαύλως γ᾽ ἀποδρὰς οὐκ ἔξει.\\
\francais{Ce n'est pas avec une fuite facile que tu t'en iras !}}
(Confusion entre un Κ et un Λ.)\\
Le vers 712, οἷον δράσας διέδυς ἔργον, est en fait apposé au vers 713 et en explique la sentence : λήψει δὲ κακόν.\\
Mais es raisons exposées jusqu'ici sont trop maigres pour justifier une émendation du texte\marqueur{1}. 
En voici une autre. Je m'appuierai sur un débris de scholies qui se trouve à l'intérieur du texte même du \italique{Ravennas}, et que j'interprète ainsi :
\citComm{ἥκεις, φησίν, εἰς τὰς χεῖρας ἡμῶν, ὅθεν οὐ φεύξῃ.\\
Tu [re]viens, dit-elle, dans nos bras, et d'où tu ne fuiras pas.}
Comprise ainsi, cette scholie est clairement une paraphrase du vers 710 pour la première partie (ὅθεν ἥκεις εἰς τὰς χεῖρας ἡμῶν $=$ ἥκειν γ᾽ ὅθεν ἥκεις). C'est la seconde partie de la scholie qui nous intéresse : ὅθεν οὐ φεύξῃ, «~[et] d'où tu ne fuiras pas~». Peut-on espérer meilleure glose pour le `οὐκ ἔξει' que nous rétablissons ?\\
% et du vers 711 pour la seconde (ὅθεν οὐ φεύξῃ $=$ οὐκ ἔξει), du moins tels que je les lis.
%L'autre scholie, encore plus abîmée, est extérieure ; elle se trouve dans les marges du manuscrit :
%\citComm{ἐκ τῶν ἀνδρῶν φησίν [bis repetita : je lis σφῇσιν] ὅθεν οὐκ ἔξεστιν}
%L'on comprend que ce scholiaste interprétait différemment le vers 710 -- d'accord en cela avec Willems : ὅθεν ἥκεις, "d'où tu viens", sous-entendu ἐκ τῶν ἀνδρῶν, "de chez les hommes", jusqu'à nous (σφῇσιν). La fin de la scholie, ὅθεν οὐκ ἔξεστιν, n'offre aucun sens acceptable. Il faut bien-sûr lire : 
%\citComm{ὅθεν οὐκ ἔξεσθα\\
%Et d'où tu ne sortiras pas...}
%Je plie les scholies pour les accorder à mes vues sur le texte, me direz-vous. Peut-être, mais convenez que les corrections sont minimes par rapport au gain en intelligibilité.%
}%
\debutNote{711-2}%
\note{1}{D'autant que le manuscrit de Ravenne concorde parfaitement à ce vers avec la \italique{Souda} qui le cite.}%
%\note{2}{Il y a nombre de σφῇσιν à rétablir dans les commentaires là où des scribes peu attentifs auront laissé le trop commun φησίν. (Nous en avons déjà croisé un sur notre route chez Eustathe de Thessalonique, cf. la note au vers 643.) L'erreur s'explique d'autant plus aisément lorsque -- comme c'est le cas ici -- le mot précédent se termine par un σ.}%
%\commentaire{711}{φαύλως,}{\eff{C'est bien en mauvaise part qu'il faut prendre ici φαύλως : \italique{injurieusement}, \italique{sans façon}, \italique{avec mépris}, \italique{grossièrement}.
%Aristophane l'emploie une fois ailleurs avec un verbe qui se rapporte à la course (\italique{Acharniens}, vers 215) :
%\citComm{οὐκ ἂν ἐπ᾽ ἐμῆς γε νεότητος, ὅτ᾽ ἐ-\\
%γὼ φέρων ἀνθράκων φορτίον\\
%ἠκολούθουν Φαΰλλῳ τρέχων, ὧδε \textbf{φαύλως} ἂν ὁ\\
%σπονδοφόρος οὗτος ὑπ᾽ ἐμοῦ τότε διωκόμενος\\
%ἐξέφυγεν οὐδ᾽ ἂν ἐλαφρῶς ἂν ἀπεπλίξατο.\\
%Ce n'est pas dans mon jeune âge, quand avec une charge de charbons je luttais de vitesse avec Phaÿllos, que ce portepaix aurait si φαύλως échappé à ma poursuite, qu'il se serait si prestement dérobé.}}
%Cf. les Guêpes et Phaÿllos : pourquoi est-il accusé d'outrage ? est-ce à cause du φαύλως ? ou bien parce que Phaÿllos était proverbialement réputé parler gauloisement ? (cf. le proverbe "passer les bornes", recensé par Suidas et Érasme) ou allusion à un Phaÿllos contemporains ?
%}%
}%
\end{groupe}%
%% v. 710
\begin{groupe}%
\typeVers{2 anap.}\chantChoeur{οἷον δράσας διέδυς ἔργον·}%
\end{groupe}%
%% v. 710
\begin{groupe}%
\typeVers{anap.}\chantChoeur{λήψει δὲ κακόν.\\\passeLigne}%
\end{groupe}%
%% v. 714
\begin{groupe}%
\personnageGrec{Μνησίλοχος}%
\typeVers{$4$ troch\cat}\tetramTroch{τοῦτο μέντοι μὴ γένοιτο μηδαμῶς, ἀπεύχομαι.\\\passeLigne}%
\end{groupe}%
%% v. 715-716
\begin{groupe}%
\personnageGrec{Χορός}%
\typeVers{2 dochm.}\diDoch{τίς οὖν σοι, τίς ἂν σύμμαχος ἐκ θεῶν\apparat{715}{\lecon{οὖν σοι} R : ἂν σοι Hamaker}\\
ἀθανάτων ἔλθοι ξὺν ἀδίκοις ἔργοις;\\\passeLigne}%
\end{groupe}%
%% v. 717
\begin{groupe}%
\personnageGrec{Μνησίλοχος}%
\typeVers{trim. ia\cat}\trimIambCat{μάτην λαλεῖτε· τὴνδ᾽ ἐγὼ οὐκ ἀφήσω.\apparat{717}{τὴν δ᾽ R : \lecon{τήνδ᾽} Biset}\\\passeLigne}%
\end{groupe}%
%% v. 718-725
\begin{groupe}%
\ledpb%
\personnageGrec{Χορός}%
\typeVers{2 iamb.}\chantChoeur{ἀλλ᾽ οὐ μὰ τὼ θεὼ τάχ᾽ οὐ\\
χαίρων ἴσως ἐνυβριεῖς\apparat{719}{ἐνυβρίσεις R : \lecon{ἐνυβριεῖς} anonymus Romanus}\\
λόγους τε λέξεις ἀνοσίους}
\end{groupe}%
%% v. 718-725
\begin{groupe}%
\typeVers{dochm.}\chantChoeur{ξὺν ἀθέοις ἔργοις\apparat{721}{ἀθέοις ἔργοις γὰρ R : ἀθέοις γὰρ ἔργοις Burges : ἐπ' ἀθέοις ἔργοις γὰρ Enger : \lecon{ξὺν ἀθέοις ἔργοις γὰρ} Bibi ap. Taillardat, cf. v.716 \espApp γὰρ R : \lecon{καὶ γὰρ} Hermann \ligneApp{724} τάχα δέ σε R : \lecon{τάχα δὲ} Fritzsche \ligneApp{725} τις τύχη R : \lecon{τύχη} Bergk grat. metr. \ligneApp{726} \lecon{\expo{χ}ρῆν σ᾽} \Rdeux : χρὴ μ᾽ \Mon}%
\commentaire{721-5}{ἀθέοις ἔργοις γὰρ ἀνταμει- / ψόμεσθά σ᾽ ὥσπερ εἰκὸς ἀντὶ τῶνδε / τάχα δὲ \{σε\} μεταβαλοῦσ᾽ ἐπὶ κακὸν ἑτερότροπον / ἐπέχει \{τις\} τύχη.}{Le grand Jean Taillardat avait donné un exposé sur son interprétation des vers 721-725 du manuscrit de Ravenne. Il nous offre une explication limpide de ce passage fort débattu ; je transcris pieusement ci-dessous le résumé de son intervention, tout en reprenant sa traduction\marqueur{1} : \\
« Car toi, avec tes impiétés, nous te châtierons de ta conduite, comme de juste. Et puis la Fortune, changeant brusquement d'humeur, aura vite fait de poursuivre un méchant homme qui a pris le chemin du crime ».\\
Cette traduction appelle les justifications suivantes : \\
— Ἀθέοις ἔργοις : selon Jean Taillardat, il est impossible d'interpréter ce datif comme le complément de moyen ἀνταμειψόμεσθα et de traduire « \latin{\italique{Impiis factis} enim remunerabimus te sicut par est pro eis} » (Andrea Divo) ou « \latin{\italique{Nefariis} enim \italique{factis} vicissim, ut par est, haec tua ulciscemur} » (Brunck) ou « Nous y répondrons, comme de raison, \italique{par des actes sacrilèges} » (Willems)\marqueur{2}. En effet, on ne pouvait pas considérer comme sacrilège « \italique{la coutume} (qui) était d'apporter du feu près des suppliants réfugiés sur l'autel » afin de les contraindre à le lâcher (schol. Eur. \italique{Androm}. 256 : ἔθος ἦν πῦρ προσφέρειν τοῖς εἰς βωμὸν καταφεύγουσιν ; cf. Eur. \italique{Herc. fur}. 243, 285). Ἀθέοις ἔργοις est donc : \ \ \ 1) \ soit un datif-instrumental d'accompagnement lié à σ(ε), cf. Kühner-Gerth, I, p. 435; ἀθέοις ἔργοις a, dans ce cas, la même fonction syntaxique que ξὺν ἀδίκοις ἔργοις\marqueur{3} (v. 716) lié à σοι (v. 715) ; \ \ \ 2) \ soit un datif-instrumental de cause lié à ἀνταμειψόμεσθα. Jean Taillardat penche pour la première solution.\\
— Ἀντὶ τῶνδε (v. 723) $=$ ἀντὶ τῶνδε τῶν ἀθέων ἔργων « nous te payerons en retour de tes impiétés».\\
— Μεταβαλοῦσ(α), v. 724, est employé absolument, ce qui n'est pas rare, cf. L.S.J. \italique{s. v.} μεταβάλλω, A. III. 2.\\
On reconnaîtra d'autre part la locution ἐπέχειν ἐπί τινα (Hérodote 9.59) à peu près synonyme de διώκειν τινά (Id. \italique{ibid}.), car τάχα ... ἐπὶ κακὸν ἑτερότροπον ἐπέχει (présent de prophétie) τύχη équivaut à τάχα ... ἐπὶ κακὸν ἄνδρ᾽ ὃς τὴν ἑτέραν ὁδὸν ἐτράπετ' ἐπέχει τύχη.\\
Pour κακός $=$ κακός τις ἀνήρ, cf. Théognis 189, 190, 577, 972, 1112 (de même κακὴ $=$ κακή τις γυνή, id. 185). Le composé ἑτερότροπος « qui a pris la mauvaise route » (avec ἕτερος euphémique pour κακός, cf. L.S.J. III. 2) est une allusion au mythe bien connu des deux chemins, celui du bien et celui du mal (Hésiode \italique{Op}. 287-292 ; Prodicos in Xénophon \italique{Mém}. 2. 1. 21-27). \\
Quant à σε de τάχα δέ σε μεταβαλοῦσα, κτλ. (cod. R), Jean Taillardat le considère comme une glose qui expliquait κακόν ἑτερότροπον et qui s'est glissée dans le texte même. \\
Quoi qu'on en ait dit, le sens de ce passage est très clair : le chœur fait entendre au Parent d'Euripide qu'à défaut d'une punition donnée par les femmes, il recevra un châtiment infligé par le Ciel.}
\debutNote{721-5}
\note{1}{À retrouver dans la \italique{Revue des Études Grecques}, tome 101, fascicule 482-484 (juillet-décembre 1988). Actes de l'association, p. XVIII-XIX.}
\note{2}{Ou « \anglais{Since your crimes are godless, it is fitting that our punishment of you be \italique{godless} as well} » (Austin \& Olson.)}%
\note{3}{Fort de cette remarque on-ne-peut-plus pertinente, je [Bibi] n'ai pas hésité à introduire ce ξύν dans le texte sachant qu'il satisfait pleinement en outre la métrique. Lire au vers 721 : [ξὺν] ἀθέοις ἔργοις. C'est aussi \italique{gratia metri} qu'Enger proposait en son temps l'ajout d' [ἐπ'] ἀθέοις ἔργοις, conjecture acceptée par Lætitia (\italique{The Songs of Aristophanes}, p. 427). N.D.É.}}%
\end{groupe}%
%% v. 718-725
\begin{groupe}%
\typeVers{tetram. troch.}\chantChoeur{καὶ γὰρ ἀνταμειψόμεσθά σ᾽ ὥσπερ εἰκὸς ἀντὶ τῶνδε.}%
\end{groupe}%
%% v. 718-725
\begin{groupe}%
\typeVers{2 dochm.}\chantChoeur{\ligne{724}τάχα δὲ μεταβαλοῦσ᾽ ἐπὶ κακὸν ἑτερότρο-%
\commentaire{724-5}{τάχα δὲ \{σε\} μεταβαλοῦσ᾽ ἐπὶ κακὸν ἑτερότροπον ἐπέχει \{τις\} τύχη.}{Il suffit que le chœur dise `τάχα δὲ' ... et le rythme même du vers s'accélère : les trois dochmiaques des vers 724-725 contiennent force résolutions.}}%
\end{groupe}%
%% v. 718-725
\begin{groupe}%
\typeVers{dochm.}\chantChoeur{\alinea πον ἐπέχει τύχη.\\\passeLigne}%
\end{groupe}%
%% v. 726-727
\begin{groupe}%
\personnageGrec{Κορυφαῖος}%
\typeVers{$4$ troch\cat}\tetramTroch{\ligne{726}ἀλλὰ τάσδε μὲν λαβεῖν χρῆν σ᾽ ἐκφέρειν τε τῶν ξύλων,%
\commentaire{726-7}{ἀλλὰ τάσδε μὲν λαβεῖν χρῆν σ᾽ ἐκφέρειν τε τῶν ξύλων, / καὶ καταίθειν τὸν πανοῦργον πυρπολεῖν θ᾽ ὅσον τάχος.}{C'était, en effet, l'habitude de faire périr par le feu ceux qui avaient cherché refuge sur un autel. Cf. Euripide, \italique{Androm}., v. 256 ; \italique{Herc. Fur}., v. 243 et 285 et les interprètes. (Note de Willems.)\\
Cette \italique{exhortation} -- κατακελευσμός selon la dénomination consacrée de Zieliński\marqueur{1} -- de la Coryphée pour appeler à la reprise de l'action est selon l'usage un distique de tétramètres trochaïques. (Et, comme au vers 531-532, débute par un ἀλλὰ.)}%
\debutNote{726-7}%
\note{1}{Voir \italique{The Verse of Greek Comedy} (1912), p. 316, de John W. White. Le lecteur pourra utilement consulter ce livre de référence, entre autre s'il s'intéresse à la structure des comédies d'Aristophane.}%
\commentaire{726}{Τάσδε}{renvoie à des compagnes de Mica et non pas à des torches, selon l'hypothèse invraisemblable de Fritzsche -- qui a malheureusement fait florès. Logique pratique élémentaire%\marqueur{1} 
: avant d'aller recueillir du bois, il faut d'abord se \italique{débarrasser} de ce qui encombre -- en l'occurrence : les torches -- et non s'en embarrasser. Ce n'est qu'une fois le bois entassé que Mica pourra effectivement se saisir d'une torche et bouter le feu au bûcher.}%
%\debutNote{726}%
%\note{1}{Il m'est utile d'avoir été professeur pour ce commentaire des \italique{Thesmophories} : on apprend à tout préciser, même des évidences telles qu'on aurait d'abord cru insulter l'intelligence de l'auditoire en s'abaissant à les souligner. Un éditeur d'Aristophane, tel le jeune élève dont l'esprit est en cours de formation, a parfois besoin d'être pris par la main.}%
%\commentaire{726}{ἐκφέρειν τε τῶν ξύλων,}{\eff{parler de la porte du fond de scène d'où serait tiré ce bois ?}}
\\
καὶ καταίθειν τὸν πανοῦργον πυρπολεῖν θ᾽ ὅσον τάχος.\\\passeLigne}%
\end{groupe}%
%% v. 728-729
\begin{groupe}%
\personnageGrec{Μίκα}%
\typeVers{trim. iamb.}\trimIamb{ἴωμεν ἐπὶ τὰς κληματίδας ὦ Μανία.\commentaire{728}{ἴωμεν ἐπὶ τὰς κληματίδας.}{Blaydes, jugeant l'article de trop, propose : ἴωμεν οὖν ἐπὶ κληματίδας. Nous répondrons par le texte suivant de Plutarque, où il s'agit, comme ici, d'un sanctuaire où des femmes se trouvent réunies pour la célébration de leurs mystères :%
\citComm{$[$ ἡ Δαμοκρίτα $]$ ξύλα ταῖς θύραις πολλὰ προσνήσασα ταῦτα δ᾽ ἦν \textbf{εἰς τὴν τῆς ἑορτῆς θυσίαν ὑπ᾽ ἐκείνων παρεσκευασμένα}, πῦρ ἐνῆκε.\\
Damocrita ayant entassé contre les portes une quantité de bois, \italique{que les femmes avaient préparé pour le sacrifice de la fête}, elle y mit le feu.}%
(Plutarque, \italique{Narrat. amour}., 5, p. 775e ; Note de Willems.)}\\
κἀγώ σ᾽ ἀποδείξω θυμάλωπα τήμερον.\testimonia{729}{(θυμάλωπα) Hesychius θ 862 θυμάλωψ ; cf. Photius θ 256 Θυμάλωψ, Hesychius θ 861}\apparat{729}{\lecon{κἀγώ} R : ἦ ᾽γώ Sommerstein}\\\passeLigne}%
\end{groupe}%
%% v. 730-732
\begin{groupe}%
\ledpb%
\personnageGrec{Μνησίλοχος}%
\trimIamb{ὕφαπτε καὶ κάταιθε· σὺ δὲ τὸ Κρητικὸν\testimonia{730-1}{(σὺ δὲ -- ταχέως) Suidas κ 2408}\apparat{730}{σὺ τόδε τὸ R : \lecon{σὺ δὲ τὸ} Suid. κ 2408}%
\commentaire{730}{τὸ Κρητικὸν.}{Hésychius (κ 4087) : %
\citComm{Κρητικόν· ἱματίδιον λεπτὸν καὶ βραχύ.\\
\francais{Crétique : petit habit fin et court.}}%
}\\%
ἀπόδυθι ταχέως· τοῦ θανάτου δ᾽ ὦ παιδίον\\
μόνην γυναικῶν αἰτιῶ τὴν μητέρα.}%
\end{groupe}%
%% v. 733-735
\begin{groupe}%
\trimIamb{τουτὶ τί ἔστιν; ἀσκὸς ἐγένεθ᾽ ἡ κόρη\apparat{733}{ἐγένηθ᾽ R : \lecon{ἐγένεθ᾽} Gryaneaus}%
\commentaire{733}{ἀσκὸς ἐγένεθ᾽ ἡ κόρη.}{Il était probablement interdit d'apporter du vin lors des Thesmophories, fêtes consacrées à Déméter et sa fille. (Cf. par exemple l'adage latin \latin{Cereri sacrificant}\marqueur{1}, \italique{sacrifier à Cérès}, c'est-à-dire faire abstinence de vin.) Cette interdiction devait souffrir quelques écarts, si l'on en croit le compte-rendu de Mnésiloque (vers 626-633).}%
\debutNote{733}%
\note{1}{Recensé par Érasme, \italique{Adages} n° 3071. Voir aussi les adage n° 1895, 1896.}\\
οἴνου πλέως καὶ ταῦτα Περσικὰς ἔχων.\testimonia{734}{(Περσικὰς) \Eustathe 1522.12}%
\commentaire{734}{Περσικὰς ἔχων.}{On comprend d'après d'autres passages d'Aristophane que les \italique{persiques} désignent des chaussures légères et recouvrant le pied\marqueur{1}, peut-être des \italique{chaussons}\marqueur{2} ; en tout cas, ces \italique{persiques} sont portées par les femmes\marqueur{3}, ou comme ici les enfants.\\
Les «~persiques~» dont parle Mnésiloque sont visibles sur le fameux cratère de Würzbourg\marqueur{4}. Il faut savoir que les outres à l'époque étaient faites de la peau d'un animal évidé\marqueur{5}, bouc ou chèvre la plupart du temps. Ces «~persiques~» sont en réalité les pattes arrière de l'animal outré, fermées par un lacet.}%
\debutNote{734}%
\note{1}{La puce des socratiques expériences des \italique{Nuées}, v. 151, se retrouve chaussée de `persiques' après que Socrate lui eut plongé les pattes dans la cire – voir aussi Hésychius (π 2005) : \grec{περσικά· εὐτελῆ ὑποδήματα}.% \eff{N'y a-t-il pas là un proverbe ?}
}%
\note{2}{Les Athéniennes pouvaient les garder au lit, en prenant garde toutefois de maintenir les jambes en l'air, cf. \italique{Lysistrata} v. 229.}%
\note{3}{\italique{Assemblée des femmes} v. 319 et Hésychius (μ 929) : \grec{μεσοπερσικαί· ὑποδήματα γυναικεῖα}.}%
\note{4}{Reproduit en annexe. Voir la note aux vers 754-757.}%
\note{5}{Pour plus d'informations, aller consulter l'indispensable Daremberg \& Saglio, \italique{Dictionnaire des antiquités grecques et romaines} à l'entrée \latin{Uter}.% \eff{+ Realencyclopädie}
}%
\\
ὦ θερμόταται γυναῖκες, ὦ ποτίσταται\testimonia{735-7}{Suidas θ 251 (ὦ et αὖ omissis, καλόν et δὲ pro ἀγαθόν et δ᾽)}\testimonia{735}{\Eustathe 1441.23 ; \EtymSym α 241 (ed. Lasserre-Livadaras) \espTest (ποτίσταται) \EtymMagnum 31.13}%\apparat{735}{\lecon{θερμόταται} R, Suid. θ 251}% : θερμοπόται Bibi}%
\commentaire{735}{ὦ θερμόταται γυναῖκες.}{θερμός évoque l'\italique{impétuosité} féminine : \anglais{hot-headed, hasty}, l'\italique{audace} des femmes\marqueur{1}.
Le mot peut aussi désigner la \italique{fièvre} ; il n'y a cependant pas d'équivalent grec à notre \italique{fièvre de Bercy} : à ma connaissance, le grec n'a jamais rapproché la fièvre de l'ivrognerie.
Quand donc le scholiaste à notre vers écrit `ὦ μέθυσοι', il glose `ὦ ποτίσταται' et non `ὦ θερμόταται'.
Suidas (θ 251) s'y est trompé qui recopie la glose en la faisant précéder du lemme θερμόταται :%
\citComm{Θερμόταται: μέθυσοι. ὦ θερμόταται γυναῖκες, μέγα καπήλοις καλόν, ἡμῖν δὲ κακόν [...]}%
Le cas n'est pas isolé : Suidas avait déjà commis un précédent du même genre : au vers 178 (\italique{confer} le commentaire \italique{ad locum}), lisant κοινῇ à place de καινῇ, il n'hésita pas à utiliser la glose afférente, μεγάλῃ καὶ θαυμαστῇ, comme nouvelles acceptions de κοινῇ (en y ajoutant le véritable sens de κοινῇ : ἀνθρωπίνῃ).
Ce sont là des méthodes qui montrent que la \italique{Souda} est un lexique à manipuler avec précaution.
%Faut-il voir avec le \LSJ , Taillardat (§313, 351) et les différents traducteurs une évocation de l'\italique{impétuosité} féminine : \anglais{hot-headed, hasty}, de leur \italique{audace} ? Comme suivent une énumération des vertus ivrognesques du beau sexe, je serais porté à l'avis du scholiaste – qui glose \grec{ὦ μέθυσοι} –, repris par Suidas à l'entrée (θ 251) :
%\citComm{Θερμόταται: μέθυσοι. ὦ θερμόταται γυναῖκες, μέγα καπήλοις καλόν, ἡμῖν δὲ κακόν [...]}
%Faut-il comprendre qu'il s'agit d'un type très particulier de bouffées de chaleur : la \italique{fièvre de Bercy}~? «~Ô les enfiévrées de Bercy~!~»\\
%Pourtant, cette acception de θέρμος n'est attestée nulle part ailleurs et n'a pas non plus été relevée par les lexicographes...\\
%Alors, un scribe aurait-il par erreur attribué la glose au mauvais lemme ; ὦ μέθυσοι se rapportant à ὦ ποτίσταται et non ὦ θερμόταται ?\\
%Je vois encore une autre hypothèse : si le texte d'Aristophane contenait à l'origine ὦ θερμοπόται\marqueur{1}.
%\eff{θερμοπότης désigne littéralement un \italique{buveur de boisson chaude}\marqueur{θερμοποτίς étant l'un des termes pour la coupe où l'on délayait \italique{à chaud} le vin.}. Avec l'usage de mélanger le vin épais à de l'eau \italique{chaude}, ce mot a pu désigner plus particulièrement les francs buveurs (voir Athénée). serait à prendre au sens de "ardent à boire" (forgé sur le même modèle que le θερμό-βουλος d'Euripide, fr. 858, qu'Aristophane parodie au vers 119 des \italique{Acharniens}). 
%"vin béni pour les débitants"}%
}%
\debutNote{735}%
\note{1}{Cf. le \LSJ , Taillardat §§313, 351.%\eff{le citer ? + Érasme}
}%
%\note{1}{À prendre au sens de "ardent à boire" ? (forgé sur le même modèle que le θερμό-βουλος d'Euripide, fr. 858, qu'Aristophane parodie au vers 119 des \italique{Acharniens}.) Le θερμοπώλιον, ou \latin{thermopolium}, désignait une \italique{taverne} chez les latin ; le θερμοπότης, \italique{qui boit chaud}, devait être particulièrement porté sur certain type de boisson, comme il appert d'un passage d'Athénée (VIII, 352b) :τοὺς δὲ Ῥοδίους ὁ αὐτὸς Στρατόνικος σπαταλῶνας καὶ θερμοπότας θεωρῶν.}%
}%
\end{groupe}%
%% v. 736-738
\begin{groupe}%
\trimIamb{κἀκ παντὸς ὑμεῖς μηχανώμεναι πιεῖν,\\
ὦ μέγα καπήλοις ἀγαθὸν ἡμῖν δ᾽ αὖ κακόν,\\
κακὸν δὲ καὶ τοῖς σκευαρίοις καὶ τῇ κρόκῃ.\\
\commentaire{738}{κακὸν δὲ καὶ τοῖς σκευαρίοις καὶ τῇ κρόκῃ.}{%\eff{Erasme 1646 et l'athénien qui tend la main jusqu'à la mort.}
Mnésiloque reproche à ces soûlaudes d'Athéniennes de contribuer à la ruine des ménages par leur inconduite biberonnesque.
Montrons-le !
\grec{τό σκευάριον}, diminutif de σκεῦος, désigne littéralement un \italique{petit récipient} (\italique{menue vaisselle}) ou \italique{petit meuble}.
Employé au pluriel, il peut désigner plus particulièrement le meuble où l'on entasse son argent. \italique{Confer} le passage du \italique{Ploutos} ci-dessous (vers 806-811), qui nous renseigne sur l'usage du mobilier grec :%
\citComm{Ἡ μὲν σιπύη μεστή ᾽στι λευκῶν ἀλφίτων,\\
οἱ δ᾽ ἀμφορῆς οἴνου μέλανος ἀνθοσμίου.\\
Ἅπαντα δ᾽ ἡμῖν ἀργυρίου καὶ χρυσίου\\
\textbf{τὰ σκευάρια} πλήρη ᾽στίν, ὥστε θαυμάσαι.\\
Τὸ φρέαρ δ᾽ ἐλαίου μεστόν· αἱ δὲ λήκυθοι\\
μύρου γέμουσι, τὸ δ᾽ ὑπερῷον ἰσχάδων.\\
\francais{Notre huche au pain regorge d'une blanche farine,\\
nos amphores d'un vin noir qui embaume.\\
Tous nos \italique{meubles}} (τὰ σκευάρια) \francais{sont pleins\\
d'argent et d'or, que c'est merveille.\\
La citerne regorge d'huile, les lécythes\\
sont emplis d'essences, le grenier de figues sèche}}
(C'est bien de \italique{pièces} d'or et d'argent dont il est question puisque les plats et autres assiettes d'airain ou d'argent sont évoqués aux vers suivants.)
Il est tout naturel que ces meubles (coffrets ?) servant de porte-monnaie aient pris par métonymie le sens de \italique{biens}, \italique{économies}.
C'est d'ailleurs si naturel que c'est effectivement le cas : aux vers 838-839, le Juste, contant ses misères passées, emploie τὰ σκευάρια en ce sens métonymique :
\citComm{Δίκαιος — κομιδῇ μὲν οὖν· / αὐχμὸς γὰρ ὢν \textbf{τῶν σκευαρίων} μ᾽ ἀπώλεσεν.\\
\francais{Le Juste — Précisément. Une sécheresse s'étant mise dans \italique{mon mobilier} ($=$ mes économies) fut ce qui me perdit.}}%
Cette acception de τὰ σκευάρια est encore plus manifeste dans un passage au début des \italique{Grenouilles}. Dionysos y cherche un moyen d'expédier en Enfer ses \italique{bagages} ; son esclave Xanthias regimbant, il hèle un mort en route pour l'Hadès (vers 170-172) :
\citComm{Καὶ γάρ τιν᾽ ἐκφέρουσι τουτονὶ νεκρόν.\\
Οὗτος, σὲ λέγω μέντοι, σὲ τὸν τεθνηκότα.\\
Ἄνθρωπε, βούλει \textbf{σκευάρι᾽} εἰς Ἅιδου φέρειν;\\
\francais{Car voilà justement un mort qu'on emporte.\\
Holà ! c'est à toi que je parle, toi le mort.\\
Homme, veux-tu porter à l'Hadès quelques \italique{petits colis} ?}}%
Il y a un évident détournement dans la plus pure veine aristophanesque de l'universel adage : on n'emporte (φέρειν) pas son \italique{argent} (σκευάρι᾽) dans l'autre monde (εἰς Ἅιδου) ; les σκευάρια de l'adage sont à prendre au sens figuré, \italique{argent} mais se retrouvent tout matériellement incarnés en \italique{bagages} dans la version illustrée d'Aristophane\marqueur{1}. Voilà pour τοῖς σκευαρίοις, passons à τῇ κρόκῃ.\\
\grec{ἡ κρόκη} est au propre \italique{le fil de trame}. 
Il possède, comme nous l'allons voir, le même sens imagé de \italique{richesses} que τὰ σκευάρια.
Le grec, comme le français, a prodigalement tiré maintes images et maints proverbes de l'activité filandière, ou tisserande. Observons l'un deux : σπαθᾶν τὴν κρόκην.
Au vers 55 des \italique{Nuées}, Strepsiade énumère les oppositions entre lui, le rustre économe, et sa femme, la raffinée dilapidatrice. Il pousse la comparaison jusque dans l'alcôve où sa femme souffre de besoins tyranniques ; il lui reproche de \italique{tirer trop sur la corde} -- le grec dit \italique{frapper sur le fil} :%
\citComm{ὦ γύναι, λίαν σπαθᾷς.}%
σπαθᾶν [τὴν κρόκην] désigne l'action de battre chaque fil de trame (κρόκη) au moyen de la σπάθη, ou \italique{battant}, pour rendre le tissu compact\marqueur{2} ; ce faisant, l'on utilise davantage de laine, d'où le fait que l'expression signifie \italique{dissiper l'argent}, \italique{gaspiller}.
Or l'on trouve les expressions σπαθᾶν τὰ πατρῷα\marqueur{3}, τὰ χρήματα\marqueur{4}, τάλαντα\marqueur{5}. 
Le sème commun à ces différents mots est celui de \italique{richesse} (matérielle ou pécuniaire) ; puisqu'ils ont été substitués à τὴν κρόκην, c'est donc la preuve -- taillardesque\marqueur{6} -- qu'ils partagent avec κρόκη ce sème.\\
Pour s'assurer que notre propos soit bien entendu, illustrons-le d'un exemple, en français cette fois. On rencontre les expressions \italique{ourdir un complot}, \italique{une intrigue}, \italique{une conspiration}, \italique{une vengeance}. Pris dans leur contexte, le sème commun aux différents compléments d'ourdir est : \italique{machination secrète}. Or l'expression tisserandière propre est «~ourdir une trame~», qui consiste à entrecroiser les fils de trame avec les fils ourdis, c'est-à-dire disposés sur l'ourdissoir. Les termes de complot, intrigue, \&c. ne sont substituables à la trame que parce qu'ils ont en commun avec elle le sème de \italique{machination secrète} (sens figuré de trame).\\
Finissons avec un amusant proverbe sur ce fléau qu'est l'alcool, surtout pour le sexe filandier\marqueur{7} :%
\citComm{ἡ γραῦς τὰς ἀγαθίδας ἐκπίνουσα\\
La vieille boit\marqueur{8} ses pelotes.}%
ἡ ἀγαθίς est une \italique{pelote de fil}. On constate que ce mot a aussi pris le sens de \italique{bien}, \italique{richesse} (on parlerait d'un \italique{bas de laine} en français) ; par permutation métaphorique avec κρόκη ? peut-être aussi par paronomase : ἀγαθῶν ἀγαθίδες est une expression désignant une grande quantité de biens.
Notre tradition populaire n'est d'ailleurs pas en reste : voici ce que l'on chantait à l'époque où l'on chantait encore\marqueur{9} :%
\citComm{\francais{Parlons un peu des filandières.\\
Celles qui filaient l'étoupe\\
Ont mangé toute la soupe.\\
Celles qui filaient le lin\\
Ont lampé tout le vin.}}%
Et encore, à propos de filles qui se sont enivrées :%
\citComm{\francais{Elles ne l'ont pas fait en filant du lin,\\
Mais elles l'ont bien fait en buvant du vin.}}%
%\eff{Problème (Ach. 451 et Guêpes) : c'est aussi les vêtements/costumes de théâtre...
%Guêpes 794 : le coq digère même l'argent :
%ἀλεκτρυόνος μ᾽ ἔφασκε κοιλίαν ἔχειν. “ταχὺ γοῦν καθέψεις τἀργύριον,” ἦ δ᾽ ὃς λέγων.
%À moins que le coq porte une krokê naturellement : Hésychius :
%κρόκος" τὸ φυτόν. καὶ ἀλεκτρυόνες οἱ τὸν τράχηλον τοιοῦτον ἔχοντες. κρόκη
%Item : οὐδ᾽ἐτρύπησε κρόκην : οὐδ᾽εἰργάσατo
%Lysistrata 897-898 :
%ΚΙ. Ὀλίγον μέλει σοι τῆς κρόκης φορουμένης
%ὑπὸ τῶν ἀλεκτρυόνων;
%EMPORTER LE FIL : gaspiller... sachant que les coqs portent un krokos(ê) (Hésychius !) proverbe : tirer le fil du vêtement ?}%
}%
\debutNote{738}%
\note{1}{Même procédé dans les \italique{Nuées} : quand Strepsiade désabusé songe au \italique{Tourbillon} (Δῖνος) de Socrate, il ne voit plus en cette nouvelle divinité qu'une vile \italique{poterie} (δῖνος). Si l'on veut savoir ce que le véritable Socrate entendait par ce \italique{Tourbillon}, il suffit d'interroger un autre grand esprit socratique : Pierre Boutang : «~La lignée [de race et d'esprit], c'est merveilleux !.. mais un peu trop linéaire pour mon goût. Donc il me faut quelque chose de plus circulaire : des lignes qui se répètent, qui reviennent sur elles-mêmes, qui se retournent ; des tourbillons ! [ ... ] Un grand tourbillon créateur !~». Quoi qu'on en ait dit, le Socrate croqué par Aristophane est loin d'être si caricatural et malgré ses défauts qu'Aristophane grossit à la loupe, il est encore à mille coudées au-dessus de notre pâle réplique moderne de Cléon, le cocaïno-mégalomane à la `pensée complexe' -- pour prendre un exemple qui nous \italique{touche}, temporairement -- heureusement !}%
\note{2}{Avec le commentaire du scholiaste repris dans la \italique{Souda} (ε 3176) : καὶ γὰρ ἐπὶ τῶν ὑφαινομένων ἱματίων λέγομεν σπαθᾶν τὸ ἄγαν κρούειν τὴν κρόκην. Taillardat (§440) développe l'exégèse de l'image. Voir aussi le lexique de Photius (ε 2014) : Ἐσπαθημένος· πεπυκνωμένος· ἀπὸ τῶν κατακρουουσῶν σπαθῶν τὴν κρόκην.}%
\note{3}{Diphil. 43, 27.}%
\note{4}{Plut., \italique{Per.} 14. Il s'agit ici des \italique{deniers publics} : τῶν δὲ περὶ τὸν Θουκυδίδην ῥητόρων καταβοώντων τοῦ Περικλέους ὡς σπαθῶντος τὰ χρήματα καὶ τὰς προσόδους ἀπολλύντος.}%
\note{5}{Lucien, \italique{Cat.} 20.}%
\note{6}{Lire et relire l'excellentissime article du regretté Jean, « Images et matrices métaphoriques » (1977), in \italique{Bulletin de l'Association Guillaume Budé} 4, pp. 344-354.}%
\note{7}{\italique{Corpus paroemiographorum Graecorum}, Appendices, III, 3.}%
\note{8}{Sur l'image de «~manger~» son bien, voir Taillardat §439.}%
\note{9}{Doit-on dire de cette époque qu'«~elle est révolue, qu'elle a fait son temps ; qu'on n'la reverra plus... et qu'c'est bien attristant~» ? Toujours est-il que j'emprunte ces chansons à une tradition qui a depuis longtemps défunté, mais qui fut néanmoins -- et heureusement ! -- bien embaumée par notre grand Patrice Coirault, aux cotes 110-05 et 110-06 de son \italique{Répertoire des chansons françaises de tradition orale}.}%
}%
\end{groupe}%
%% v. 739
\begin{groupe}%
\ledpb%
\personnageGrec{Μίκα}%
\trimIamb{\ligne{738}παράβαλλε πολλὰς κληματίδας ὦ Μανία.\\}%
\end{groupe}%
%% v. 740-741a
\begin{groupe}%
\personnageGrec{Μνησίλοχος}%
\trimIamb{παράβαλλε δῆτα· σὺ δ᾽ ἀπόκριναί μοι τοδί,\apparat{740}{ἀπόκρινέ R : \lecon{ἀπόκριναί} Μυ2}\\
τουτὶ τεκεῖν φῄς;\\\passeLigne}%
\end{groupe}%
%% v. 741b-742a
\begin{groupe}%
\personnageGrec{Μίκα}%
\trimIamb{\ligne{741}\decalage{τουτὶ τεκεῖν φῄς;}καὶ δέκα μῆνας αὔτ᾽ ἐγὼ\apparat{741}{δέκα μῆνας R : \lecon{καὶ δέκα μῆνας} Biset : δέκα μῆνας γε Scaliger}%
\commentaire{741}{δέκα μῆνας.}{Les mois lunaires grecs étant de 29 ou 30 jours, la femme portait bien environ dix mois l'enfant. Aulu-Gelle dans ses \italique{Nuits attiques} traite pendant tout un chapitre (III, 16) de la question du temps de gestation. N'en déplaise au scholiaste, la durée de 10 mois est bien attestée, notamment par Ménandre ; Aulu-Gelle cite aussi Varron expliquant que les noms des Parques \italique{Nona} et \italique{Decima} viennent des nombres de mois fixés pour l'accouchement (\latin{item Nona et Decima a partus tempestivi tempore}).}\\
ἤνεγκον.\testimonia{742}{(ἤνεγκον) \Antiatt 98.11 ; \Phryn 73.1}\\\passeLigne}%
\end{groupe}%
%% v. 742b
\begin{groupe}%
\personnageGrec{Μνησίλοχος}%
\trimIamb{\ligne{742}\decalage{ἤνεγκον.}ἤνεγκας σύ;%
\commentaire{742-66}{}{L'on possède pour ce passage un autre témoin que le \italique{Ravennas} : le premier fragment du \nieme{3839} papyrus d'Oxyrhynque\marqueur{1} (\PsiDeux \  dans l'apparat critique). Ce morceau de papyrus est bien antérieur au manuscrit de Ravenne : il est daté d'environ 200 ap. J.-C. Il présente objectivement de meilleures leçons en bien des points. Aussi, aux endroits où ces deux témoins différaient et que le bon choix à faire n'était pas évident, j'ai privilégié la version du papyrus. Nous reproduisons en annexe ce papyrus ainsi que la lecture qu'en fit Walter E. H. Coockle pour la collection \italique{The Oxyrhynchus Papyri}.%
}%
\debutNote{742-66}%
\note{1}{\italique{The Oxyrhynchus Papyri}, vol. LVI (1989), n°3839, fr. 1.}%
\\\passeLigne}%
\end{groupe}%
%% v. 742c
\begin{groupe}%
\personnageGrec{Μίκα}%
\trimIamb{\ligne{742}\decalage{ἤνεγκον.ἤνεγκας σύ;}νὴ τὴν Ἄρτεμιν.\\\passeLigne}%
\end{groupe}%
%% v. 743a
\begin{groupe}%
\personnageGrec{Μνησίλοχος}%
\trimIamb{τρικότυλον; ἢ πῶς.\apparat{743}{\_τρικο\douteux{τ}\douteux{υ}\douteux{λ}\douteux{ο}ν·η\douteux{π}\douteux{ω}\douteux{c} \  ειπ\douteux{ε}\douteux{μ}οι\douteux{τ}[ \PsiDeux : τρικότυλον ἢ πῶς; εἰπέ μοι --- τί μ᾽ εἰργάσω; R : \lecon{ἠργάσω} Hall-Geldart}
\commentaire{743}{τρικότυλον ἢ πῶς;}{C'est-à-dire du vin supérieur, à une obole les trois cotyles. Hésychius (τ 1388) : %
\citComm{\grec{τρικότυλος οἶνος· οὗ τρεῖς κοτύλαι ὀβολοῦ πωλοῦνται},}%
ce qui est cher, en effet, car cela met le métrète (144 cotyles, environ 40 litres) à $7fr68$\marqueur{1}, le double à peu près du prix ordinaire. Démosthène (1045.5) désigne de la même manière un vin plus cher encore, \grec{ὀ οἶνος δωδεκάδραχμος}, le vin à douze drachmes, c'est-à-dire à $11fr52$ le métrète\marqueur{2}. À rapprocher de τετρώβολον, désignant dans la \italique{Paix}, 254, le miel attique. [...] — Ἤ πῶς; est mis pour ἤ πῶς ὤνιος; Cf. \italique{Acharn}., 758, \italique{Caval}., 480. (Note de Willems.)\\
Mnésiloque s'était déjà adressé à Mica en lui reprochant d'avoir acheté son enfant au vers 565.\\
Comme annoncé, je suis le papyrus d'Oxyrhynque contre le \italique{Ravennas} ; on peut y déchiffrer :% 
\citComm{\_τρικο\douteux{τ}\douteux{υ}\douteux{λ}\douteux{ο}ν·η\douteux{π}\douteux{ω}\douteux{c} \ \ ειπ\douteux{ε}\douteux{μ}οι\douteux{τ}[}%
Le point médian est, dans ce papyrus, un signe de ponctuation interrogatif. Quant à l'espace blanc de deux lettres entre \douteux{π}\douteux{ω}\douteux{c} et ειπ\douteux{ε}\douteux{μ}οι, il indique que le changement d'interlocuteur a lieu à cet endroit, et non après εἰπέ μοι comme le placent le \italique{Ravennas} et les différents éditeurs.
La lecture du papyrus donne donc pour le vers 743 :%
\citComm{τρικότυλον; ἢ πῶς. --- εἰπέ μοι·τί μ᾽ ἠργάσω;}%
}%
\debutNote{743}%
\note{1}{$7fr68$ de feu notre franc germinal~; c'est-à-dire qu'on pouvait échanger à tout moment contre $7,68 \times 4,5 = 34,56$ grammes d'argent pur~; soit effectivement 48 oboles (de $0,72$ grammes dans le système éginète)... mais cela, c'était avant l'invention des monnaies de singe : chez nous, bien-sûr, et déjà en $-406$ à Athènes (\italique{confer} les \italique{Grenouilles}, vers 718-733).}%
\note{2}{72 oboles soit $51,84$ grammes d'argent.}\\\passeLigne}%
\end{groupe}%
%% v. 743b-745a
\begin{groupe}%
\personnageGrec{Μίκα}%
\trimIamb{\ligne{743}\decalage{τρικότυλον; ἢ πῶς.}εἰπέ μοι· τί μ᾽ ἠργάσω;\\
ἀπέδυσας ὦναίσχυντέ μου τὸ παιδίον\\
τυννοῦτον ὄν.\testimonia{745}{(τυννοῦτο) Photius τ 552 τυννοῦτο $\sim$ Hesychius τ 1645}\\\passeLigne}%
\end{groupe}%
%% v. 745b-746
\begin{groupe}%
\personnageGrec{Μνησίλοχος}%
\trimIamb{\ligne{745}\decalage{τυννοῦτον ὄν.}τυννοῦτο; μικρὸν νὴ Δία.\apparat{745}{\lecon{τυννοῦτον} et \lecon{τυννοῦτο} \PsiDeux, Brunck : τυνοῦτον et τυνοῦτο R \espApp Γϋ. ante μικρὸν R : \lecon{nulla nota mutatae personae} \PsiDeux \ligneApp{746} \lecon{γέγονε} \PsiDeux, Brunck : γέγονεν R \ligneApp{747} καὶ ὅσον \PsiDeux, R : ὅσον Suid. δ 1168 : \lecon{χὤσον} Bentley}\\
πόσ᾽ ἔτη δὲ γέγονε; τρεῖς χοᾶς ἢ τέτταρας;%
\commentaire{746}{πόσ᾽ ἔτη δὲ γέγονε; τρεῖς Χοᾶς ἢ τέτταρας;}{S'il faut en croire le scholiaste, l'on pouvait compter l'âge à partir du nombre de Grandes Dionysies écoulées (\grec{τοῦτο ἔθος ἦν Ἀττικοῖς, λέγειν τὰ ἔτη κατὰ τὸν ὑπερπίπτοντα ἀριθμὸν ἀπὸ τῶν Διονυσίων}) ; usage qui correspond chez nous à celui de compter les \italique{printemps}. Mnésiloque, filant la métaphore de l'outre et l'enfant, préfère se référer malicieusement aux \italique{Choés}, journée la plus arrosée des Anthestéries\marqueur{1} ; d'autant qu'un \grec{χοῦς} est aussi une unité de volume\marqueur{2}. Cela permet insidieusement à Mnésiloque de préciser que l'outre n'est pas si petite : elle est d'une capacité d'une dizaine de litres\marqueur{3}.}%
\debutNote{746}%
\note{1}{L'on trouvera d'abondants et utiles renseignements sur cette fête à l'article \latin{Dionysia} du \italique{Dictionnaire des antiquités grecques et romaines} de Daremberg \& Saglio. Cette fête des Anthestéries correspond à peu près à nos Saint-Vincent, Percée du Vin Jaune. (Ou celle du beaujolais nouveau pour les dégénérés.)}%
\note{2}{Le \grec{χοῦς} – ou \italique{conge} romain — correspond à 12 cotyles, soit environ 3 litres ¼.}%
\note{3}{D'une capacité \italique{théorique} de trois ou quatre conges : il est évident que Mica ne s'est pas encombrée du fardeau que serait cette outre pleine (Mnésiloque aurait d'ailleurs eu le plus grand mal à la lui chaparder). Une telle capacité peut étonner mais il faut bien voir que ces outres étaient faites de la peau complète d'un animal (souvent un bouc ou une chèvre) : ainsi, Strepsiade (\italique{Nuées}, v. 1238), trouvant que son créancier ferait une belle outre, l'estime de six conges. Il y avait des outres d'une capacité impressionnante pour jouer au jeu de l'\grec{ἀσκωλιασμός}, consistant à sauter -- \italique{clocher}, comme l'étymologie l'indique -- et tenir en équilibre sur une outre graissée. Voir l'article \latin{Uter} du \italique{Dictionnaire des antiquités grecques et romaines} de Daremberg \& Saglio. Il y a de toute façon une exagération comique : l'outre représentée sur le cratère de Würzbourg n'a certainement pas cette capacité et Mnésiloque est dans l'outrance verbale.}\\\passeLigne}%
\end{groupe}%
%% v. 747-748a
\begin{groupe}%
\personnageGrec{Μίκα}%
\trimIamb{σχεδὸν τοσοῦτον χὤσον ἐκ Διονυσίων.\testimonia{747}{Suidas δ 1168 (καὶ omisso)}%\apparat{747}{καὶ ὅσον \PsiDeux, R : ὅσον Suid. δ 1168 : \lecon{χὤσον} Bentley}%
\commentaire{747}{σχεδὸν τοσοῦτον χὤσον ἐκ Διονυσίων.}{L'on traduit généralement ce vers par quelque chose d'équivalent peu ou prou à : « C'est cela, plus le temps écoulé depuis les Dionysies ». Outre qu'inventer un anniversaire à l'outre est une billevesée indigente relevant davantage de l'humour des philologues que de celui d'Aristophane, il présente une impossibilité : essayez donc de conserver pendant trois ou quatre \italique{printemps} une outre de vin dans les conditions de l'Athènes de l'an $-2300$ avant le Réfrigérateur !
Et puisque Mica ne peut pas parler de l'âge de son outre, elle parle donc de sa \italique{capacité}. En l'occurrence, elle essaye de se défendre du chef d'accusation de biberonne porté par Mnésiloque~: sa chère petite outre, «~peut-être est-elle de cette taille [c'est-à-dire d'une capacité d'une dizaine de litres] mais elle est telle [c'est-à-dire aussi peu remplie] depuis les Dionysies~». Autrement dit, Mica au jour des Thesmophories prétend ne pas avoir touché au contenu de l'outre depuis les dernières Dionysies, ce qui représente une abstinence de sept mois ! (Bluff !)}\\
ἀλλ᾽ ἀπόδος αὐτό.\\\passeLigne}%
\end{groupe}%
%% v. 748b
\begin{groupe}%
\personnageGrec{Μνησίλοχος}%
\trimIamb{\ligne{748}\decalage{ἀλλ᾽ ἀπόδος αὐτό.}μὰ τὸν Ἀπόλλω τουτονί.
\commentaire{748}{μὰ τὸν Ἀπόλλω τουτονί.}{Contre les commentateurs qui se contentent d'une banale formule de ce genre pour inventer sur scène la présence d'une statue d'Apollon omniprésente,
% (alors que – faut-il le rappeler ? – Mnésiloque se trouve dans l'enceinte du Thesmophorion, consacré à Déméter et Perséphone et est déjà juché sur l'autel qu'il aurait donc bien du mal à faire passer pour un Ἀγυιεύς ; qu'à cela ne tienne, nos professeurs inventèrent une statue d'Apollon omniprésente pour le décor fixe des pièces de théâtre, statue qui se surajoute au décor particulier à chaque pièce), 
Willems avait déjà répondu par une note au vers 83 des \italique{Nuées}. Comme ce genre d'idées saugrenues se retrouve encore dans les dernières éditions les plus autorisées, nous avons cru faire œuvre utile en reproduisant ladite note en annexe. %Il est à espérer que sa lecture dissipe définitivement les nuées de Philologus. % ce genre de nuées. %(Il n'est d'ailleurs pas impossible qu'Aristophane fasse jurer Mnésiloque sur Apollon simplement pour citer un passage correspondant du \italique{Télèphe}.)
}\\\passeLigne}%
\end{groupe}%
%% v. 749a
\begin{groupe}%
\personnageGrec{Μίκα}%
\trimIamb{ἐμπρήσομεν τοίνυν σε.\\\passeLigne}%
\end{groupe}%
%% v. 749b-750
\begin{groupe}%
\personnageGrec{Μνησίλοχος}%
\trimIamb{\ligne{749}\decalage{ἐμπρήσομεν τοίνυν σε.}πάνυ γ᾽ ἐμπίμπρατε·\testimonia{749}{(ἐμπίμπρατε) Hesychius ε 2458 ἐμπίπραται $\sim$ Suidas ε 1018 $\sim$ Zonaras 709}\apparat{749}{ἐμπιμπρᾶται R : \lecon{ἐμπιμπράτε} (sic) \Mon}\\
αὕτη δ᾽ ἀποσφαγήσεται μάλ᾽ αὐτίκα.\\\passeLigne}%
\end{groupe}%
%% v. 751-752a
\begin{groupe}%
\personnageGrec{Μίκα}%
\trimIamb{μὴ δῆθ᾽, ἱκετεύω σ᾽· ἀλλ᾽ ἔμ᾽ ὅ τι χρῄζεις ποίει\\
ὑπέρ γε τούτου.\\\passeLigne}%
\end{groupe}%
%% v. 752b-753
\begin{groupe}%
\personnageGrec{Μνησίλοχος}%
\trimIamb{\ligne{752}\decalage{ὑπέρ γε τούτου.}φιλότεκνός τις εἶ φύσει.\\
ἀλλ᾽ οὐδὲν ἧττον ἥδ᾽ ἀποσφαγήσεται\\\passeLigne}%
\end{groupe}%
%% v. 754-755
\begin{groupe}%
\personnageGrec{Μίκα}%
\trimIamb{οἴμοι τέκνον. δός τὸ σφαγεῖον Μανία,\testimonia{754}{Photius σ 850 σφαγεῖον· [...] περισπᾶται ἡ δευτέρα συλλαβή}\apparat{754}{\lecon{δός τὸ} \PsiDeux, Lobeck ad Soph. \italique{Ai.} v.1066 : δός μοι τὸ R \espApp σφαγίον \Rac : \lecon{σφαγεῖον} \Rpc, \Mon}
\commentaire{754-7}{Scène d'égorgement de l'outre nourrissonne.}{Cette scène est illustrée sur le cratère de Würzbourg\marqueur{1} (reproduit en annexe). Amusante mise en abîme : le `σφαγεῖον' que tient Mica sur cette représentation n'est autre qu'un cratère de forme identique au cratère de Würbourg lui-même -- c'est-à-dire un vase propre à recevoir du vin pour l'y mélanger et non une coupe servant à recevoir le sang des victimes (σφαγεῖον).
% Les détournements de ce genre sont fort nombreux dans les \italique{Thesmophories}.
}%
\debutNote{754-7}%
\note{1}{Cratère apulien d'environ 370 av. J.-C., Martin von Wagner Museum der Universität Würzburg, Antiken-Sammlung, H5697. L'on pourra se reporter pour plus de détails à l'article d'Eric Csapo, \italique{A Note on the Würzburg Bell-Crater H5697} («~Telephus Travestitus~») in \italique{Phoenix}, Vol. 40, No. 4 (Winter, 1986), p. 379-392.}\\
\commentaire{754}{δός μοι τὸ σφαγεῖον.}{Le papyrus d'Oxyrhynque n°3839 ne comporte pas le μοι. Nous avons dit précédemment suivre le papyrus contre le \italique{Ravennas} et nous nous y sommes tenus. Mais il y a d'excellentes raisons à cela, outre bien-sûr le privilège de l'ancienneté.
D'abord, il n'est rien de plus courant que l'interpolation de petits mots tels que μοι (cf. le vers 1047, où un autre papyrus, daté du \nieme{II} siècle ap. J.-C., ne présente pas non plus le μοι du manuscrit de Ravenne).
Comment ce μοι est-il apparu ? il est vraisemblable que ce soit pour de pseudo-raisons métriques.
Faisons l'hypothèse que d'anciens manuscrits comportaient un \nieme{754} vers ainsi troussé :%
\citComm{οἴμοι τέκνον δός τὸ σφαγίον Μανία.}%
Noter la graphie σφαγίον (et non σφαγεῖον) : c'est celle de la première version du \italique{Ravennas} (avant correction en σφαγεῖον, faite d'une autre main s'il faut en croire Coockle, \latin{scholiorum manus, aequalis fere primae manui} selon von Velsen\marqueur{1}). σφαγίον était donc scandé : $ \cup \cup - $ ; il manquait par conséquent une syllabe longue à placer entre τέκνον et σφαγίον pour que le vers soit d'équerre. L'ajout d'un μοι était la solution la plus simple et à la portée de n'importe quel scribe.\\
Ce sont là des considérations métriques qui devaient passer bien au-dessus de la tête des scribes ayant retranscrit les \italique{Thesmophories}, me direz-vous peut-être. Détrompez-vous : la correction même du \italique{Ravennas} -- alors qu'à l'époque les deux graphies avaient la même prononciation -- montre bien l'intérêt pour la scansion ; sur le papyrus, Coockle releva une «~\anglais{unexplained ink like an oddly shaped perispomen accent in a second hand above} γε~» (le γε de σφαγειον) ; l'explication me paraît simple : l'annotateur confirmait que la deuxième syllabe de σφαγεῖον est longue. C'est d'ailleurs l'objet d'une glose de Photius (σ 850) :%
\citComm{σφαγεῖον· [...] περισπᾶται ἡ δευτέρα συλλαβή.}%
}%
\debutNote{754}%
\note{1}{\italique{Aristophanis Thesmophoriazusae} (1883). Adolphe von Velsen a établi définitivement dans cette édition que le second manuscrit par lequel nous sont connues les \italique{Thesmophories}, le \italique{Monacensis gr. 492} (G dans son apparat critique, \Mon dans le nôtre), n'était qu'une copie du \nieme{XV} siècle du \italique{Ravennas}, apportant quelques émendations bien vues. Son édition compare exhaustivement le \italique{Ravennas} et le \italique{Monacensis} ; ce qui est d'un grand intérêt pour déterminer si une correction du \italique{Ravennas} est ancienne ou plus récente : il suffit d'aller comparer avec la leçon du manuscrit de Münich. Son apparat critique est le plus précis et permet d'avoir la meilleure idée possible du manuscrit de Ravenne -- avec celui de l'édition Prato ; en partie une reprise du premier mais revu et amélioré en de nombreux points. Pour qui voudrait pousser le scrupule plus loin, et puisque le manuscrit de Ravenne n'a toujours pas été numérisé -- ô scandale ! le manuscrit sans qui nous ne connaîtrions que neuf des onze pièces survivantes d'Aristophane !.. --, Van Leeuwen édita en fac-similé le \italique{Ravennas} : \italique{Comoediae Undecim cum Scholiis : Codex Ravennas 137, 4, A} [$=$ \italique{Rav.} 429 (nouvelle cote)], \italique{phototypice editus} in \italique{Codices Graeci et Latini photographice depicti}, t. IX (1904).}%
ἵν᾽ οὖν τό γ᾽ αἷμα τοῦ τέκνου τοὐμοῦ λάβω.%
\commentaire{755-7}{τό γ᾽ αἷμα τοῦ τέκνου τοὐμοῦ λάβω --- ὕπεχ᾽ αὐτό, χαριοῦμαι γὰρ ἕν γε τοῦτό σοι --- κακῶς ἀπόλοι᾽, ὡς φθονερὸς εἶ καὶ δυσμενής.}{En prononçant le vers 756, les commentateurs avaient imaginé jusqu'ici que : \ \ soit Mnésiloque déversait le contenu de l'outre sur l'autel\marqueur{1} ; \ \ soit il buvait d'un trait le vin\marqueur{2}. 
Mais ces visions de la scène rentrent toutes deux en contradiction avec le texte d'Aristophane : Mica a demandé à \italique{récupérer} le sang de son enfant (τό γ᾽ αἷμα τοῦ τέκνου τοὐμοῦ λάβω) et Mnésiloque lui concède cette unique faveur (χαριοῦμαι γὰρ ἕν γε τοῦτό σοι). Il faut l'en croire.
%Or la seconde hypothèse est incompatible avec la réplique de Mnésiloque, qui lui concède cette unique faveur, χαριοῦμαι γὰρ ἕν γε τοῦτό σοι -- et l'ironie a bon dos quand il s'agit d'ignorer tout bonnement le texte que nous livre Aristophane !.. En outre, si Mnésiloque, déversait effectivement l'outre sur l'autel, comment pourrait-il dire : ὕπεχ᾽ αὐτό, "mets-le [le σφαγεῖον] sous"... sous quoi ? sous l'\italique{autel} ?! les commentateurs qui proposent cela savent-ils bien à quoi ressemble l'autel en question ? et quand bien même : il serait alors impossible à Mica d'en récupérer plus qu'une infime partie et Mnésiloque de ne pas tenir sa promesse.\\
Seulement, Mnésiloque \italique{prend au mot} Mica. Le \italique{mot}, en l'occurrence, c'est λάβω : \italique{prendre}, \italique{récupérer}, dans la bouche de Mica. Mnésiloque entend \italique{recevoir}. C'est entendu : elle n'a qu'à \italique{présenter} (ὕπεχ᾽ αὐτό) le σφαγεῖον. Il lui jette alors à la figure le contenu de l'outre. C'est précisément cet instant, où Mica tend le σφαγεῖον, qui est saisi sur le cratère de Würzbourg. On voit qu'il suffit à Mnésiloque de balancer le bras pour asséner l'outre à fragmentation liquide à la tête de sa mère.
C'est ce comportement, dépassant, outrepassant la mesure -- μηδὲν ἄγαν, Mnésiloque, μηδὲν ἄγαν. Ne te l'a-t-on pas assez répété ? --, qui explique l'emportement de Mica : κακῶς ἀπόλοι᾽, ὡς φθονερὸς εἶ καὶ δυσμενής. Ajoutez à cela l'apport bénéfique de confusion théâtrale (Mica maculée du «~sang~» de sa fille) à la confusion induite par les nombreux quiproquos\marqueur{3}, et vous comprendrez aisément les quelques précisions que demande Critylla déboussolée à son arrivée au vers 759 : `τί τῆς ἱερείας γίγνεται;' puis `ταλαντάτη Μίκα τίς ἐξεκόρησέ σε;'.
(L'exégèse de ce passage est d'ailleurs la preuve \italique{in texto} que Mica présente des signes \italique{visibles} d'agression -- agression qui n'est donc pas uniquement psychologique. Ces signes, comme nous l'allons voir à la note s'y rapportant, sont les taches de vin qui imprègnent son vêtement.)}%
\debutNote{755-7}%
\note{1}{C'est notamment l'opinion du scholiaste : τοῦτο ἅμα εἰπὼν ἐκχέει τὸν οἶνον ὁ κηδεστὴς ἐν τῷ βωμῷ. Je passe sur les autres explications alambiquées pour metteur en scène payé à l'heure de spectacle produit. (Mnésiloque en boit un peu, en verse par terre, mais pris de scrupules, finit par en tendre les dernières gouttes à Mica. Je ne caricature qu'à peine...) Soyons sérieux : une représentation d'Aristophane n'est pas le Festival de Cannes.}%
\note{2}{Austin \& Olson, conformément à leur pente, ne se prononcent pas et avancent les deux hypothèses. Il faudrait cesser de croire que la pusillanimité est une preuve d'esprit critique. «~La recherche de la vérité n'empêche pas de prendre parti~», comme disait l'autre imbécile.}%
\note{3}{Le grec dirait qu'on surajoute \italique{du feu au feu}, πῦρ ἐπὶ πῦρ.}%
\\\passeLigne}%
\end{groupe}%
%% v. 756
\begin{groupe}%
\personnageGrec{Μνησίλοχος}%
\trimIamb{ὕπεχ᾽ αὐτό, χαριοῦμαι γὰρ ἕν γε τοῦτό σοι.\\
\\\passeLigne}%
\end{groupe}%
%% v. 757
\begin{groupe}%
\personnageGrec{Μίκα}%
\trimIamb{κακῶς ἀπόλοι᾽, ὡς φθονερὸς εἶ καὶ δυσμενής.\\\passeLigne}%
\end{groupe}%
%% v. 758
\begin{groupe}%
\personnageGrec{Μνησίλοχος}%
\trimIamb{\ligne{757}τουτὶ τὸ δέρμα τῆς ἱερείας γίγνεται.\apparat{758-9}{\lecon{τῆς ἱερείας} \PsiDeux, R : τῆς ἱερέας Austin \& Olson}%
\commentaire{758}{τουτὶ τὸ δέρμα τῆς ἱερείας γίγνεται.}{τὸ δέρμα désigne l'outre vide – qui rappelons-le était faite d'une peau d'animal évidé. L'usage était que le prêtre reçoive certaines parties de la victime (\grec{τά ἱερώσυνα}).}\\\passeLigne}%
\end{groupe}%
%% v. 759a
\begin{groupe}%
\ledpb%
\personnageGrec{Κρίτυλλα}%
\trimIamb{\ligne{758}τί τῆς ἱερείας γίγνεται;\apparat{759}{nulla nota personae in R : ἄλλη γυνή \Rpc (ad v.760), \Mon : \lecon{(Κρίτ.)}. Sommerstein}%
\commentaire{759}{}{Je suis Sommerstein qui donne la parole pour ce vers à Critylla, quoiqu'aucun \italique{paragraphos} indiquant un changement d'interlocuteur ne soit présent sur le manuscrit de Ravenne (le papyrus est inexploitable sur ce point).
Faut-il prendre au premier degré la déclaration de Mnésiloque et en inférer que c'est Critylla qui jouait le rôle de prêtresse ?%
}\\\passeLigne}%
\end{groupe}%
%% v. 759b
\begin{groupe}%
\personnageGrec{Μνησίλοχος}%
\trimIamb{\ligne{759}\decalage{τί τῆς ἱερείας γίγνεται;}τουτί. λαβέ.\apparat{759}{post τουτί interpunxit Bergk}\\\passeLigne}%
\end{groupe}%
%% v. 760-761
\begin{groupe}%
\personnageGrec{Κρίτυλλα}%
\trimIamb{ταλαντάτη Μίκα τίς ἐξεκόρησέ σε;\apparat{760}{ἐξεκόρησεν R : \lecon{ἐξεκόρησε} \Mon \ligneApp{761}[\douteux{ }\douteux{ }\douteux{ }\douteux{ }\douteux{ }\douteux{ }\douteux{ }\douteux{ }]απη\douteux{τ}η\douteux{ν}\douteux{π}\douteux{ }ι\douteux{ }α\douteux{ }\douteux{ }υξ[ \PsiDeux \espApp σοὐξῃρήσατο R : \lecon{σοὐξηράσατο} Fritzsche}\\
τίς τὴν ἀγαπητὴν παῖδά σοὐξηράσατο;
\commentaire{760-1}{τίς ἐξεκόρησέ σε; / τίς τὴν ἀγαπητὴν παῖδά σου ᾿ξῃρήσατο;}{Faute d'avoir compris que Mnésiloque venait d'asperger Mica du «~sang~» de sa propre fille, ce passage n'a pas non plus été entendu.
Même le grand Willems -- une fois n'est pas coutume -- fut pris la main dans le sac à ordures : trompé par un jeu de mots en latin\marqueur{1} qu'il a cru pouvoir rapprocher de ce distique, il proposa comme interprétation une paillardise qui n'a pas lieu d'être, et qui d'ailleurs n'y est pas.\\
Figurons-nous la scène : l'outre a été égorgée et déversée sur Mica ; arrive Critylla à qui Mnésiloque jette l'outre vide. Critylla avec l'«~enfant~» éviscéré ; Mica maculée de sang face à elle.%
\citComm{Critylla --- τίς ἐξεκόρησέ σε;}%
Le verbe ἐκκορέω, \italique{balayer}, est fréquemment utilisé comme synonyme de \italique{maltraiter}, \italique{dévaster}. Cf., par exemples\marqueur{2}, Hésychius (ε 1453 ; κ 3067) :%
\citComm{ἐκκοροῦσι· φθείρουσιν, ἐκκαλλύνουσιν.\\
κορέων, κορέσκων· ἐξυβρίζων.}%
Photius (ε 412) :%
\citComm{ἐκκορηθῆναι· ἐκκαλλυνθῆναι· τὸ γὰρ κορεῖν καλλύνειν καὶ σαίρειν.}%
Ainsi dans la \italique{Paix} (vers 58-59), Trygée s'écrie\marqueur{3} :%
\citComm{Ὦ Ζεῦ, τί ποτε βουλεύει ποεῖν;\\
Κατάθου τὸ κόρημα· μὴ ᾽κκόρει τὴν Ἑλλάδα\\
\francais{Ô Zeus, que prétends-tu faire ?\\
Dépose ton balai : ne balaie pas l'Hellade de la surface de la Terre.}}%
Il n'y a à notre \nieme{760} aucun sens obscène, ni là ni ailleurs\marqueur{4}. La question de Mica revient donc à ceci : «~Qui t'a fourragée\marqueur{5} ?~», «~Qui t'a ruiné [ton habit] ?~»\\
Mais ἐξεκόρησε n'est pas mis là fortuitement ; il fait entendre le mot κόρη, la «~fille~» de Mica, avec la reprise au vers suivant : τὴν ἀγαπητὴν παῖδά σου.
Le sens premier de la question de Mica est donc de savoir qui l'a ainsi salie ; le sens allusif qui l'a \italique{privée de sa fille} (ἐξ-κόρη).\\
Pour être sûr que le spectateur saisisse bien l'amphibologie, Aristophane revient à la charge au vers 761 : τίς τὴν ἀγαπητὴν παῖδά σου ᾿ξηράσατο\marqueur{6};
Le pronom σύ a volontairement été rejeté avant le verbe pour laisser toute ambiguïté à la crase : σοὐξηράσατο $=$ σοι ἐξηράσατο\marqueur{7} d'une part ; σοὐξῃρήσατο $=$ σοῦ ἐξῃρήσατο d'autre part.\\
σοι ἐξερᾶν est le pendant de ἐκκορεῖν σε du vers précédent : Critylla demande qui a osé \italique{vider}, \italique{déverser}\marqueur{8} sur Mica ce que sa fille chérie avait dans le ventre.\\
σοῦ ἐξαιρεῖν : en raison de la proximité phonétique entre le `η' et le `α', le spectateur devait aisément entendre -- pour peu que l'acteur de Critylla s'y prêtât : τίς τὴν ἀγαπητὴν παῖδά σου ἐξῃρήσατο; «~Qui t'a ravi ta fille chérie ?~» ; écho au ἐκ-κόρη subrepticement ouï au vers précédent.\\
Puisque le sens premier de la question de Critylla -- celui de ἐκκορεῖν σε et celui en conformité avec l'action -- est bien σοι ἐξερᾶν, je reprends la conjecture de Friztsche (lire σοὐξηράσατο à la place du σοὐξῃρήσατο du \italique{Ravennas}) d'autant que ἐξῃρήσατο serait un aoriste inhabituel\marqueur{9}.%
}%
\debutNote{760-1}%
\note{1}{\latin{Quis te devirginavit} ?}%
\note{2}{Pour d'autres exemples, qui ne manquent pas, se reporter aux \italique{testimonia} des éditions des lexiques sous-cités}%
\note{3}{Voir Taillardat (§629) pour un exposé plus fourni sur cette image.}%
\note{4}{Suétone (τρὶς κεκορημένε· οἱονεὶ πολλάκις ἐκσεσαρωμένος) a donc raison contre Taillardat (Περὶ βλασφημιῶν, β 260 $=$ Eustathe p.1542.47 $=$ fr. Anacréon 366 \italique{Poetae Melici Graeci} Page) et beaucoup d'autres : l'imprécation d'Anacréon `τρὶς κεκορημένε' lancée au Thrace Smerdiès signifie chastement «~triple ordure~».}%
\note{5}{Je m'amuse, mais comme j'appréhende déjà que cette badinerie soit prétexte à la meute libidineuse des commentamenteurs pour déchaîner leur rut philololubrique, je m'empresse de préciser : à aucun moment Critylla n'évoque l'idée que le sang qui couvre Mica puisse être hymenéen. Et d'abord parce que ce serait inepte : Mica est mère. Mais si vous tenez tant à cette théorie, ô Philologus, avant d'ouvrir cette bouche qui vomit en surabondance ce que vous devriez celer de votre cœur, mentionnez-moi seulement un passage -- ou même, exhibez-moi seulement un indice, à défaut de vos polyçonneries -- dans lequel le mot παῖδα ait pu désigner l'hymen. En attendant, il serait bon prophylactiquement d'expurger -- ἐξερᾶν, mais j'anticipe... -- du Bailly, à l'entrée ἐκκορέω, l'acception \italique{déflorer} afférente à ce passage.}%
\note{6}{Je fais mienne la conjecture de Fritzsche : 'ξηράσατο (de ἐξερᾶν) à la place du 'ξῃρήσατο (de ἐξαιρεῖν) présent dans le manuscrit de Ravenne. (Je m'en explique en fin de note.)}%
\note{7}{Cf. les vers 158 et 624.}%
\note{8}{Peut-être même faut-il prendre ἐξερᾶν au sens de \italique{régurgiter}, \italique{recracher}, et comprendre que Mnésiloque, ajoutant l'insulte à l'outrage, après avoir empli sa bouche de vin, l'aura \italique{recraché} à la tête de Mica.}%
%\note{7}{\^ \^}%
\note{9}{Impossible selon le rigoriste O. Lautensach, p. 128 de sa grammaire \italique{Die Aoriste bei den Attischen Tragikern und Komikern} (1911) : \allemand{Aoriste auf} -ησα\allemand{. Der im Etymologicum Magnum p.777, 53 als "ungriechisch" verworfene und von Lobeck zu Phrynicus p. 716ff. aus der späteren Gräcität nachgewiesene Aor. }ᾕρησα \allemand{ist unmöglich bei Ar. \italique{Thesm.} 761 in} σοὐξῃρήσατο\allemand{ zu dulden, wofür Meineke} σου διεχρήσατο\allemand{, Fritzsche }σου ᾽ξηράσατο\allemand{, Dindorf} σου ᾽ξαφείλετο\allemand{ schreibt, während Lobeck, v. Velsen und Blaydes den Vers für unecht erklären.}}\\\passeLigne}%
\end{groupe}%
%% v. 762-764
\begin{groupe}%
\personnageGrec{Μίκα}%
\trimIamb{ὁ πανοῦργος οὗτος. ἀλλ᾽ ἐπειδήπερ πάρει,\\
φύλαξον αὐτόν, ἵνα λαβοῦσα Κλεισθένη\\
τοῖσιν πρυτάνεσιν ἃ πεποίηχ᾽ οὗτος φράσω.\\\passeLigne}%
\end{groupe}%
%% v. 765-775
\begin{groupe}%
\ledpb%
\personnageGrec{Μνησίλοχος}%
\trimIamb{ἄγε δὴ τίς ἔσται μηχανὴ σωτηρίας;\\
τίς πεῖρα, τίς ἐπίνοι᾽; ὁ μὲν γὰρ αἴτιος\\
κἄμ᾽ εἰσκυλίσας εἰς τοιαυτὶ πράγματα\apparat{767}{ἐσκυλίσας ἐς R : \lecon{εἰσκυλίσας εἰς} edd. pler.}%
\commentaire{767}{ἐσκυλίσας ἐς τοιαυτὶ πράγματα.}{Pour la métaphore de «~rouler dans les malheurs~», voir le vers 650 ainsi que ce qu'en dit Taillardat (§558).}\\
οὐ φαίνετ᾽, οὔπω. φέρε τίν᾽ οὖν ἂν ἄγγελον\apparat{768}{\lecon{φαίνετ᾽, οὔπω} R : φαίνεταί πω Dobree \espApp τίν᾽ οὖν R : τίν᾽ οὖν νῦν Biset : \lecon{τίν᾽ οὖν ἂν} Porson (\italique{Aristophanica} p.218) : alii alia}\\
πέμψαιμ᾽ ἐπ᾽ αὐτόν; οἶδ᾽ ἐγὼ καὶ δὴ πόρον\testimonia{769-75}{Suidas π 45 (οἶδα δ' pro οἶδ᾽, πότερον pro πόρον et γ' ἔκειντο pro γένοιντ᾽ ἄν)}\\
ἐκ τοῦ Παλαμήδους· ὡς ἐκεῖνος, τὰς πλάτας\testimonia{770-1}{Suidas π 1702}%
\commentaire{770}{ἐκ τοῦ Παλαμήδους·}{Euripide représenta son \italique{Palamède} en 415 av. J.-C. dans une tétralogie comportant aussi l'\italique{Alexandros}, les \italique{Troyennes} et le drame satyrique \italique{Sysiphe} (il fut classé second). Eschyle et Sophocle représentèrent aussi un \italique{Palamède}. Ce héros aux-mille-tours ingénieux de l'épopée troyenne (et pourtant absent de l'\italique{Iliade} ; sans doute pour ne pas faire d'ombre à Ulysse) avait procuré à l'armée grecque d'« immenses bienfaits »\marqueur{1} : il apaisa une famine à Aulis et des troubles causés à propos du partage des provisions, en enseignant aux Grecs les « lettres phéniciennes », leur apprit le jeu de dés pour se distraire, les mesures et les jetons de vote, ce qui lui valut un grand renom. Le scholiaste de l'\italique{Oreste} (v. 432) continue :%
\citComm{ἐπὶ τούτῳ δὲ φθονήσαντες οἱ περὶ Ἀγαμέμνονα καὶ Ὀδυσσέα καὶ Διομήδην τοιόνδε τι σκευωροῦσι κατ᾽ αὐτοῦ. λαβόντες γὰρ Φρύγα αἰχμάλωτον χρυσίον κομίζοντα Σαρπηδόνι ἠνάγκασαν γράψαι Φρυγίοις γράμμασι περὶ προδοσίας ὡς παρὰ Πριάμου πρὸς Παλαμήδην. καὶ τοῦτον μὲν φονεύουσι, θεράποντα δὲ Παλαμήδους πείθουσι χρήμασιν ἅμα τοῖς Τρωϊκοῖς χρήμασι καὶ τὸ γραφὲν πινάκιον ὑπὸ τὴν κλίνην θέσθαι Παλαμήδους. αὐτοὶ δὲ παρελθόντες προδοσίαν κατήγγελλον τοῦ ἥρωος καὶ φωραθῆναι τὴν σκηνὴν ἐκέλευον. εὑρεθέντος δὲ τοῦ πινακίου καὶ τῶν χρημάτων ὑπὸ τὴν κλίνην λίθοις φονεύεται Παλαμήδης\\
\francais{Devenus jaloux de ce fait, Agamemnon, Ulysse et Diomède imaginèrent ce genre de procédé : ayant fait prisonnier un Phrygien qui apportait de l'or à Sarpédon, ils le forcèrent à écrire en caractères phrygiens une lettre parlant de trahison, comme si elle avait été adressée par Priam à Palamède, puis ils le tuèrent. Ils corrompirent alors un serviteur de Palamède, pour qu'il place en même temps que l'or phrygien la tablette écrite sous la tente de son maître. Eux-mêmes, ayant comparu [dans l'assemblée des chefs], accusèrent le héros de trahison et demandèrent qu'on fouille la tente. Après la découverte de la tablette et de l'or sous la tente, Palamède fut lapidé.}}%
Anxieux de faire connaître la mort injuste de son fils à Nauplios, Œax, le frère de Palamède, imaginait le stratagème ici repris par Mnésiloque\marqueur{2} :
\citComm{ὁ γὰρ Ἐὐριπίδης ἐν τῷ Παλαμήδει ἐποίησε τὸν Οἴακα τὸν ἀδελφὸν Παλαμήδους ἐπιγράψαι εἰς τὰς ναῦς τὸν θάνατον αὐτοῦ, ἵνα φερόμεναι αὗται ἔλθωσιν εἰς τὸν Ναύπλιον τὸν πατέρα αὐτοῦ καὶ ἀπαγγείλωσι τὸν θάνατον αὐτοῦ.\\
\francais{Euripide, dans le \italique{Palamède}, a représenté Oiax, frère de Palamède, gravant sur des navires [le récit de] sa mort, afin qu'ils soient portés jusque chez Nauplios et lui apprennent sa mort.}}}%
\debutNote{770}
\note{1}{On retrouve dans une scholie à son \italique{Oreste}, au vers 432, les grandes lignes de la version euripidéenne du drame. J'en reprends la traduction de Jouan \& van Looy, à retrouver dans leur édition des fragments d'Euripide aux Belles Lettres, tome VIII, \nieme{2} partie.}
\note{2}{Scholie au vers 771 des \italique{Thesmophories}.}
\commentaire{770}{ὡς ἐκεῖνος,}{Œax, frère de Palamède.}\\
ῥίψω γράφων. ἀλλ᾽ οὐ πάρεισιν αἱ πλάται.\apparat{771}{\lecon{οὐ πάρεισιν} R : οὐ γάρ εἰσιν Fritzsche}\\
πόθεν οὖν γένοιτ᾽ ἄν μοι πλάται. πόθεν; πόθεν;\apparat{772}{\lecon{γένοιτ᾽ ἄν} R : γ᾽ ἔκειντό Suid. π 45 : γένοιντ᾽ ἄν Grynaeus (\italique{fierent} Divus) \espApp πόθεν; R : πόθεν πλάται; Biset : \lecon {πόθεν; πόθεν;} Scaliger : πλάται πόθεν; Sandbach : alii alia \ligneApp{773} εἶτα διὰ τἀγάλματ᾽ R : ταδὶ τἀγάλματ᾽ Suid. σ 742 : \lecon{εἰ ταδὶ τἀγάλματ᾽} Ellebodius}%
\commentaire{772}{πόθεν οὖν γένοιντ᾽ ἄν μοι πλάται. πόθεν;}{Le vers est incomplet. Scaliger proposa la solution la plus instinctive : redoubler le πόθεν en fin de vers : πόθεν ; <πόθεν;> Cette solution est vraisemblablement la bonne. \ \ \ 1° \ Une simple haplographie l'explique. \ \ \ 2° \ Elle est une nouvelle occurrence des railleries artiphanesques quant à la manie du redoublement chez les tragiques. Le double πόθεν se trouve d'ailleurs dans l'\italique{Ion} (vers 1453 sq.) :%
\citComm{ἰὼ γύναι, πόθεν πόθεν ἔλαβες ἐμὸν\\
βρέφος ἐς ἀγκάλας;}%
Mais ce qui faisait vraiment battre le cœur d'Euripide, c'était le redoublement du ποῦ : \italique{Électre} (vers 487 sq) :%
\citComm{ποῦ ποῦ νεᾶνις πότνι᾽ ἐμὴ δέσποινά τε,\\
Ἀγαμέμνονος παῖς, ὅν ποτ᾽ ἐξέθρεψ᾽ ἐγώ;}%
%Et encore 
%Devant ces épouvantables redoublement de ποῦ, le sang d'Aristophane n'a fait qu'un tour. La caricature ne s'est pas fait attendre :%
Pour se borner au \italique{Palamède} (fr. 2 J.-v. L. $=$ 579 Kn. $=$ 776 M.), on trouve le redoublement de l'adverbe πάλαι :%
\citComm{πάλαι πάλαι δή σ᾿ ἐξερωτῆσαι θέλων, σχολή μ᾿ ἀπεῖργε.}%
Redoublement moqué dans les \italique{Oiseaux} (vers 921), où l'ennuyeux poète qui débite le vers suivant :%
\citComm{πάλαι, πάλαι δὴ τήνδ᾽ ἐγὼ κλῄζω πόλιν,}%
ne pénétrera pas plus dans Νεφελοκοκκυγίᾳ qu'il n'aura de chance de demeurer dans la Cité de Platon.%
}\\
τί δ᾽ ἂν εἰ ταδὶ τἀγάλματ᾽ ἀντὶ τῶν πλατῶν\\
γράφων διαρρίπτοιμι; βέλτιον πολύ.
\commentaire{773-4}{τί δ᾽ ἂν εἶτα διὰ τἀγάλματ᾽ ἀντὶ τῶν πλατῶν / γράφων διαρρίπτοιμι;}{Willems (\nieme{22} note, reproduite en annexe) a expliqué en quel sens il fallait prendre \grec{τὰ ἀγάλματα} : \grec{πᾶν ἐφ᾽ ᾧ τις ἀγάλλεται} (scholie à notre vers), \italique{tout ce qu'on offre et consacre aux dieux}. En l'occurrence, ce sont les \grec{πόπανα} que Mnésiloque a apportés (au vers 285) et non pas des \italique{ex-voto} comme on en voit encore apparaître dans les modernes traductions\marqueur{1}.\\
Y a-t-il ici une allusion au proverbe grec ἄλλο γένος κώπης, «~un autre genre de rame~»\marqueur{2} ?}%
\debutNote{773-4}%
\note{1}{Après avoir encombré la scène d'une statue d'Apollon permanente d'après le vers 748, les commentateurs rajoutent ici des tablettes votives pour le décor.}%
\note{2}{Rapporté par Plutarque, \italique{Proverbes alexandrins}, 2, 52. Ou aussi dans \italique{Zu den Paroemiographen} de Cohn, p.42. Érasme l'a incorporé à ses \italique{Adages} (n°1621).}\\
ξύλον γέ τοι καὶ ταῦτα κἀκεῖν᾽ ἦν ξύλον.\\\passeLigne}%
\end{groupe}%
%% v. 776-784
\begin{groupe}%
\ledpb%
\personnageGrec{Μνησίλοχος}%
\typeVers{anap.}\anap{ὦ χεῖρες ἐμαὶ}%
\end{groupe}%
%% v. 776-784
\begin{groupe}%
\typeVers{dim. anap.}\anap{ἐγχειρεῖν χρῆν ἔργῳ πορίμῳ.\testimonia{777}{Suidas π 2064}\apparat{777}{χρή R, Suid. π 2064 : \lecon{χρῆν} Bentley}\\
ἄγε δὴ πινάκων ξεστῶν δέλτοι,\testimonia{778-80}{(πινάκων -- μόχθων) Pollux X 58 (ἄγε δὴ et δέλτοι omissis)}\testimonia{778-9}{Suidas σ 742}}%
\end{groupe}%
%% v. 776-784
\begin{groupe}%
\typeVers{(catalect.)}\anap{δέξασθε σμίλης ὁλκοὺς\\
κήρυκας ἐμῶν μόχθων· οἴμοι}%
\end{groupe}%
%% v. 776-784
\begin{groupe}%
\typeVers{(catalect.)}\anap{τουτὶ τορῶ μοχθηρόν·\apparat{781}{τὸ ῥῶ R~: \lecon{τορῶ} Bibi ap. Hesych. τ 1162}
\commentaire{781}{τουτὶ τὸ ῥῶ μοχθηρόν·}{« Voilà un `Rhô' défectueux ! »... telle est la leçon de notre unique manuscrit et telles aussi les traductions embarrassées. Car la question fuse : pourquoi Mnésiloque nous parle-t-il d'un `Ρ' ? Le scholiaste toujours prompt à lancer une ânerie propose d'y voir la troisième lettre du destinataire : Εὐρῑπίδη, ce qui ne répond pas à la question : pourquoi la troisième lettre ? pourquoi le Rhô ? il ne me semble pas plus difficile que les autres lettres à tracer. Passons sur les autres explications du même acabit (si Aristophane avait un peu plus de glosateurs, tous les mots du dictionnaire comportant un rhô y seraient effectivement passés : Εὐριπίδη, χώρει, χαίρειν, \&c.)\\
Lisez simplement : τουτὶ τορῶ μοχθηρόν. L'aoriste τορεῖν est homérique et fait partie d'un groupe dialectal, également homérique, d'aoristes radicaux thématiques où s'est dégagé un vocalisme en `o' : ἔμολον de βλώσκω, ἔθορον de θρώσκω et ἔτορον de τείρω. Ces homérismes se sont introduits chez les Tragiques, et peuvent aussi se rencontrer dans certains passages d'Aristophane. On l'avait déjà noté pour ἔμολον : on peut maintenant y ajouter ἔτορον\marqueur{1}. Hésychius (τ 1162) nous fut – comme à son habitude – d'un grand secours :%
\citComm{τορεῖν· τορῆσαι, τρῆσαι, τεμεῖν ;}%
le verbe signifie \italique{percer}, \italique{perforer}, \italique{couper}.\\
Reconstituons la scène : Mnésiloque cherche un expédient pour envoyer son message. Heureusement, il connaît son euripide par cœur : il décide de procéder à la façon d'Œax dans le \italique{Palamède}. Manque de chance : il manque les rames (ἀλ ̓ οὐ πάρεισιν αἱ πλάται) pour graver dessus son message !.. qu'à cela ne tienne : les ustensiles de théâtres pour figurer les πόπανα feront bien l'affaire (ils sont de bois, après tout !) ; voilà donc nos morceaux de bois promus en tablettes (δέλτοι). Mnésiloque peut maintenant se mettre à graver la “cire” avec son stylet – un stylet ? quel stylet ? Nouveau détournement : il appelle `stylet' le coutelas avec lequel il vient d'égorger l'outre de vin. Seulement, à force de prendre des vessies pour des lanternes, Mnésiloque se fait rattraper par la dure sylvestre réalité : graver un message au couteau sur du bois se révèle bien plus difficile qu'avec un stylet sur une tablette de cire ; ceci explique son cri de découragement : τουτὶ τορῶ μοχθηρόν, « ceci que je perce [il ne précise pas quoi, et pour cause : c'est à la fois un πόπανον, une πλάτη et un δέλτος]... est de mauvaise qualité ! »}%
\debutNote{781}%
\note{1}{Je dois cette remarque à Jean-Victor Vernhes qui, entre autres qualités, est un très fin connaisseur de la langue homérique. Que le lecteur n'hésite pas à consulter le Dictionnaire étymologique de la langue grecque de Pierre Chantraine à l'entrée τορεῖν pour plus de détails.}\\
`χώρει', `χώρει'. ποίαν αὔλακα;\apparat{782}{\lecon{χώρει} bis R : χωρεῖ bis Brunck}\\%
\commentaire{782}{χώρει, χώρει.}{Le doublement du verbe χωρεῖν se trouve aussi au vers 1076 de \italique{Médée}, grand moment d'euripidéenne psychologie où l'«~héroïne~», prise d'un sursaut d'instinct maternel, s'interroge : est-il bien sain de massacrer ses propres enfants par vengeance envers un autre ? Ne vaudrait-il pas mieux les laisser s'échapper : χωρεῖτε χωρεῖτ᾽, «~fuyez, fuyez !~». Ah ! ces fameux dilemmes euripidéens !.. \anglais{to kill ? or not to kill ?..}
Mais prenons garde à ne pas fonctionner par système ; ici, l'explication du double χωρεῖν est tout autre : Mnésiloque est en train de graver son message : χώρει, «~accours !~», enjoint-il à Euripide ; χώ...ρει χώ...ρει, «~Ac...cours ... Ac...cours~», se dicte-t-il à lui-même. Il l'écrit sur deux rames différentes, comme le pluriel αἱ πλάται le laissait présager et comme c'était aussi le cas dans la tragédie originale ; le scholiaste de notre \italique{Ravennas} précise bien : %
\citComm{ὁ γὰρ Οἴαξ ἐγχαράττει \textbf{πολλαῖς πλάταις} τὰ περὶ τὸν Παλαμήδην καὶ ἀφίησιν εἰς θάλασσαν, ὥστε μιᾷ γέ τινι τὸν Ναύπλιον περιπεσεῖν.\\
Œax grave effectivement de \italique{nombreuses rames} avec les circonstances touchant Palamède et les lance vers la mer de sorte que Nauplios tombe sur l'une d'elles.}%
À propos des ἀγάλματα, un scholiaste avait aussi écrit :%
\citComm{πίνακες γὰρ ἦσαν ἐν τῷ ἱερῷ, καὶ ἀπὸ τούτων λαβὼν γράφει, καὶ ῥίπτει αὐτὰς λέγων· `ἀπέλθετε, σημάνατε Ἐὐριπίδῄ\\
C'était des tablettes dans le temple. Les prenant, il écrivit puis les lança, en disant : «~Allez ! Informez Euripide !~»}%
}%
\commentaire{782-4}{ποίαν αὔλακα / βάσκετ᾽ ἐπείγετε πάσας καθόδους / κείνᾳ ταύτᾳ· ταχέως χρή.}{Les différents traducteurs voient dans la formule `ποίαν αὔλακα' une exclamation d'auto-contentement de Mnésiloque : «~Quelle belle lettre ! Quel sillon !~» Αὖλαξ, le \italique{sillon}, est en effet attesté dans l'\italique{Anthologie palatine} (6, 68) au sens de sillon creusé par le stylet, c'est-à-dire, par analogie, le \italique{trait d'écriture}. Soit. Mais cette interprétation confère à ποίαν un sens qu'il n'a pas. Cet adjectif interrogatif sert à demander des précisions : \italique{de quel nature ?} \italique{de quel genre ?} (Aristophane l'emploie donc très fréquemment dans ses dialogues pour débrouiller -- ou souligner ? -- un quiproquo.) Il est utilisé en reprise d'un terme, reprise directe : Ἀγάθων ὁ κλεινὸς --- Ποῖος οὗτος Ἁγάθων; (vers 30) ; Τὸν δεῖνα γιγνώσκεις --- Τὸν δεῖνα; ποῖον; (vers 621) ;  Πρωτέως τάδ᾽ ἐστὶ μέλαθρα --- Ποίου Πρωτέως; (vers 874) ou implicite : Πρὸς ποῖον ἔργον; (vers 586) ; Χρείᾳ δὲ ποίᾳ τόνδ᾽ ἐπεισφέρεις λόγον; (vers 1163). Dans le monologue de Mnésiloque, c'est la première fois qu'il est question de cette αὖλαξ ; mais reprise il y a bien\marqueur{1} : l'αὖλαξ devait sans aucun doute apparaître dans la tragédie d'Euripide. 
Réfléchissons ! Où en est Mnésiloque de l'action ? Il vient de graver en deux exemplaires l'appel au secours, χώρει, qu'il compte faire parvenir à Euripide... comment ? telle est la question ! Assurément, cette αὖλαξ désigne l'expédient qu'Œnax avait trouvé pour emporter ses messages.
On peut raisonnablement penser qu'il les déposait dans un \italique{cours d'eau} (qui envoyait ensuite les πλάται à la mer avant que l'une arrive à bon port, comme l'indique le scholiaste ? ça ne paraît pas très vraisemblable, mais la vraisemblance et une pièce d'Euripide font deux.)
C'est bien en ce sens qu'il faut prendre αὖλαξ, le \italique{sillon}, et, par extension, la \italique{rigole} ou le \italique{ruisseau}\marqueur{2}.\\
On peut supposer que Mnésiloque ne s'embarrasse pas vraiment du lièvre qu'il vient de soulever et que la question n'est là que pour souligner les invraisemblances de la mise en scène, comme Aristophane aime bien les relever. Il décide alors d'aller au plus simple : au vers suivant, il jette, rapidement et de toutes ses forces, les deux πλάται. (Une de chaque côté, vers chaque εἴσοδος, comme le laisse soupçonner πάσας καθ᾽ ὁδοὺς ?) Cette \italique{nouveauté}, \italique{originalité} (κείνᾳ ταύτᾳ) par rapport au modèle original d'Euripide a l'avantage d'accélérer l'action ; c'est ainsi qu'il faut comprendre le vers 784 : κείνᾳ ταύτᾳ· ταχέως χρή.%
}%
\debutNote{782-4}%
\note{1}{Et si l'on en doute, on peut toujours corriger le texte -- ce dont on se gardera bien ! Lire ποῖ ποτ' ἂν αὔλακα revient au même quant à l'action et dispense de supposer que l'αὖλαξ était le terme employé par Euripide. (ποτε \italique{gratia metri} mais dont la disparition relèverait d'une haplographie bien explicable. Au reste, la collocation ποῖ ποτε est très fréquente.)}%
\note{2}{Voir le \LSJ \ et le \espagnol{Diccionario Griego-Español} IV à l'entrée αὖλαξ.}\\
βάσκετ᾽ ἐπείγετε πάσας καθ᾽ ὁδοὺς\apparat{783}{καθόδους R : \lecon{καθ᾽ ὁδοὺς} Biset}}\\%
\end{groupe}%
%% v. 776-784
\begin{groupe}%
\lednopb%
\typeVers{(catalect.)}\anap{κείνᾳ ταύτᾳ· ταχέως χρή.\apparat{784}{ταῦτα R : ταύτᾳ Grynaeus (\italique{hac} Divus)}%
\\\passeLigne}%
\end{groupe}%
%% v. 785-843
\begin{groupe}%
\ledpb%
\personnageGrec{Χορός}%
\typeVers{\decalageN{a}$4$ an\cat}\tetramAnap{\ligne{785}ἡμεῖς τοίνυν ἡμᾶς αὐτὰς εὖ λέξωμεν παραβᾶσαι,
\commentaire{785-845}{Parabase.}{La parabase des \italique{Thesmophories}, à l'instar de \italique{Lysistrata} et des \italique{Grenouilles}, ne présente ni ᾠδή/ἀντῳδή ni ἀντεπίρρημα\marqueur{1}.}%
\debutNote{785-845}%
\note{1}{Mon but n'est pas de proposer une description exhaustive de la structure des comédies d'Aristophane, ni même des \italique{Thesmophories} en particulier. Ce travail a déjà été très bien fait par White, \italique{The Verse of Greek Comedy} (1912). On lira aussi avec un grand intérêt les travaux de Jean Irigoin regroupés dans son anthologie posthume \italique{Le poète grec au travail} (2009) et notamment son article «~Prologue et \italique{parodos} dans la comédie d'Aristophane (\italique{Acharniens}, \italique{Paix}, \italique{Thesmophories})~» :  «~[Aristophane], ce poète comique à l'imagination débridée, s'astreignait, dans certaines parties de ses pièces, à respecter des règles concernant le choix des vers, leur disposition, et même leur nombre~». Le lecteur intéressé se reportera à ces travaux -- et à nombre d'autres -- auxquels je ne saurais me substituer.}\\
καίτοι πᾶς τις τὸ γυναικεῖον φῦλον κακὰ πόλλ᾽ ἀγορεύει,\\
ὡς πᾶν ἐσμὲν κακὸν ἀνθρώποις κἀξ ἡμῶν ἐστιν ἅπαντα,\\
ἔριδες νείκη στάσις ἀργαλέα λύπη πόλεμος. φέρε δή νυν,\apparat{788}{στάσεισ· ἀργαλέα· R : \lecon{στάσις ἀργαλέα} Scaliger \espApp εἰ καὶ κακόν R : \lecon{εἰ κακόν} Zanetti}%
\commentaire{788}{ἔριδες νείκη στάσις ἀργαλέα}{Le vocabulaire est Homérique : \italique{Iliade}, XVII, v. 384-5 :%
\citComm{τοῖς δὲ πανημερίοις ἔριδος μέγα νεῖκος ὀρώρει / ἀργαλέης·\\
\francais{Durant tout le jour se prolongea cette lutte sanglante et terrible.}}%
}\\
εἰ κακόν ἐσμεν, τί γαμεῖθ᾽ ἡμᾶς, εἴπερ ἀληθῶς κακόν ἐσμεν,%\apparat{789}{εἰ καὶ κακόν R : \lecon{εἰ κακόν} Zanetti}
\commentaire{789}{εἰ καὶ κακόν ἐσμεν, τί γαμεῖθ᾽ ἡμᾶς, εἴπερ ἀληθῶς κακόν ἐσμεν}{La perfidie d'Aristophane est à son comble quand il donne la parole aux femmes dans la parabase où elles assurent leur propre défense. « Si nous sommes un fléau, disent-elles aux hommes, pourquoi nous épousez-vous ? » Mais le spectateur, qui connaît, depuis l'école, son Hésiode par cœur et qui a déjà vu bien des comédies, répond intérieurement, comme mu par un ressort : « Parce qu'on n'échappe pas si facilement au piège tendu par Zeus, parce que la femme est un fléau nécessaire pour la perpétuation de l'espèce\marqueur{1} ». Le poète comique ne fait ici, sous des apparences trompeuses, qu'entretenir le vieux réflexe de misogynie. Autre argument employé par les femmes pour montrer leur valeur : celui de la jalousie qu'elles excitent\marqueur{2}. Tout commentaire est ici inutile... (Note d'Alain Blanchard\marqueur{3}.)}
\debutNote{789}
\note{1}{Hésiode, les \italique{Travaux et les jours}, v. 83-84 ; Susarion, fr. 1, 3-4 Kassel-Austin. Voir aussi ce que dit Renzo Tosi (§284 et §285 de son monumental \italique{Dictionnaire des sentences latines et grecques}, 2010) : Κακὸν ἀναγκαῖον γυνή ou \latin{Femina dulce malum} : la femme est un καλὸν κακόν (\italique{Théogonie}, v. 587). Quant au vers précédent, il y une allusion à peine voilée aux pandoriennes calamités déclenchées par la mythique -- et belle, surtout, pour nos malheurs ! -- Hélène. (Voir à ce propos Tosi, §350 : \latin{Nam fuit ante Helenam cunnus taeterrima belli causa}.) Il est donc tout naturel que Mnésiloqie réponde à cette parabase par une reprise de l'\italique{Hélène}.}
\note{2}{Aux vers 792-794.}
\note{3}{Tirée de son remarquable article «~L'art de la mauvaise foi dans les \italique{Thesmophories} d'Aristophane~», in \italique{Dieux, héros et médecins grecs. Hommage à Fernand Robert.} (2001), p. 97-106. Tout amateur d'Aristophane ne peut pas ne pas aller prestement lire ces pages si riches (j'émets seulement quelques réserves sur le tout dernier paragraphe).
Blanchard y montre que « la mauvaise foi d'Aristophane s'est exercée, en ce qui concerne Euripide, sur un thème moins traditionnel que celui de la misogynie et qu'elle permet de mieux comprendre le sujet même de sa comédie des \italique{Thesmophories}. » De même que la défaite athénienne à Délion avait profité à la popularité de Socrate, qui s'était conduit héroïquement dans la débâcle comme on le sait par Platon, et que c'est ce bonheur d'une réputation rehaussée dans le malheur athénien qui avait provoqué Aristophane à l'écriture vengeresse des \italique{Nuées} ; de même, les \italique{Thesmophories} furent écrites peu après la désastreuse expédition de Sicile : alors qu'Athènes vivait ses heures les plus sombres, Euripide en profita pour redorer son blason : on fit appel à lui pour écrire un chant funèbre en l'honneur des Athéniens morts à Syracuse (Plutarque nous en a conservé deux vers). 
Anecdote plus significative encore : certains Athéniens, réduits en esclavage, ne durent leur salut qu'à Euripide, en récitant de ses prologues, si prisés des Syracusains. Que la poésie d'Euripide pût sauver qui que ce soit ?.. alors qu'elle est justement l'une des causes de la dégénérescence athénienne !? voilà encore une fumisterie qu'Aristophane n'aura de cesse de mettre à mal tout au long des \italique{Thesmophories} : les différents détournements des euripidéennes tragédies se heurteront toutes successivement au mur de l'incompréhension des geôliers de Mnésiloque. Une fois bien cerné le contexte historique, l'on comprend mieux Aristophane : ce « jeu de massacre » n'est motivé que par la passion -- brûlante chez Aristophane -- de la justice : il cogne, cogne, et cogne encore contre tous ces charognards des misères athéniennes. \anglais{Now is a time to storm !} Et si Euripide est misogyne, c'est là son moindre défaut.}\\
κἀπαγορεύετε μήτ᾽ ἐξελθεῖν μήτ᾽ ἐκκύψασαν ἁλῶναι,\apparat{790}{ἐγκύψασαν R : \lecon{ἐκκύψασαν} Reiske}\\
ἀλλ᾽ οὑτωσὶ πολλῇ σπουδῇ τὸ κακὸν βούλεσθε φυλάττειν;\\
κἂν ἐξέλθῃ τὸ γύναιόν ποι, κᾆθ᾽ εὕρητ᾽ αὐτὸ θύρασιν,\apparat{792}{κᾆθ᾽ εὕροιτ' R : \lecon{κᾆθ᾽ εὕρητ᾽} Porson}\\
μανίας μαίνεσθ᾽, οὓς χρῆν σπένδειν καὶ χαίρειν, εἴπερ ἀληθῶς\apparat{793}{μαίνεθ᾽ R : \lecon{μαίνεσθ᾽} \Mon \espApp χρή R : \lecon{χρῆν} Brunck \ligneApp{794} εὕρετε R : \lecon{ηὕρετε} Meineke \espApp καταλαμβάνετ᾽ R : \lecon{κατελαμβάνετ᾽} Brunck}\\
ἔνδοθεν ηὕρετε φροῦδον τὸ κακὸν καὶ μὴ κατελαμβάνετ᾽ ἔνδον.\\
κἂν καταδάρθωμεν ἐν ἀλλοτρίων παίζουσαι καὶ κοπιῶσαι,\\%\apparat{795}{καταδαρθῶμεν R : \lecon{καταδάρθωμεν} Dindorf}\\
πᾶς τις τὸ κακὸν τοῦτο ζητεῖ περὶ τὰς κλίνας περινοστῶν.\\\passeLigne}
\end{groupe}%
%% v. 785-843
\begin{groupe}%
\ledpb%
\tetramAnap{κἂν ἐκ θυρίδος παρακύπτωμεν, τὸ κακὸν ζητεῖτε θεᾶσθαι·\apparat{797}{ζητεῖ τὸ κακὸν τεθεᾶσθαι R : \lecon{τὸ κακὸν ζητεῖτε θεᾶσθαι} J. Kaye ap. Dobree \ligneApp{798} αἰσχυνθεὶς (ex αἰσχυνθῆς) \Rpc : \lecon{αἰσχυνθεῖσ᾽} Zanetti \ligneApp{799} παρακύψασαν ἰδεῖν τὸ κακὸν R : \lecon{τὸ κακὸν παρακύψαν ἰδεῖν} Dindorf}%
\commentaire{797}{τὸ κακὸν ζητεῖτε θεᾶσθαι}{Suivant la belle correction de Kaye : τὸ κακὸν ζητεῖτε θεᾶσθαι, v. 797, au lieu de : ζητεῖ τὸ κακὸν τε θεᾶσθαι. (Note de Willems.)}\\
κἂν αἰσχυνθεῖσ᾽ ἀναχωρήσῃ, πολὺ μᾶλλον πᾶς ἐπιθυμεῖ\\
αὖθις τὸ κακὸν παρακύψαν ἰδεῖν. οὕτως ἡμεῖς ἐπιδήλως\\
ὑμῶν ἐσμεν πολὺ βελτίους, βάσανός τε πάρεστιν ἰδέσθαι.\testimonia{800, 810}{(βελτίους $=$ βελτίονες) Moeris β 8}\\
βάσανον δῶμεν πότεροι χείρους. ἡμεῖς μὲν γάρ φαμεν ὑμᾶς,%
\commentaire{801-2}{ἡμεῖς μὲν γάρ φαμεν ὑμᾶς, ὑμεῖς δ᾽ ἡμᾶς}{Molière, \italique{Impromptu de Versailles}.%
\citComm{\francais{Il gage que c'est moi, et moi je gage que c'est lui.}}%
}\\
ὑμεῖς δ᾽ ἡμᾶς. σκεψώμεθα δὴ κἀντιτιθῶμεν πρὸς ἕκαστον,\apparat{802}{ὑμεῖς δ᾽ ὑμᾶς R : \lecon{ὑμεῖς δ᾽ ἡμᾶς} \Mon \espApp σκεψώμεσθα R : \lecon{σκεψώμεθα} \Mon \ligneApp{803} ἕκαστον R : \lecon{ἑκάστου} Grynaeus \ligneApp{804} μέν R : \lecon{μέν γ᾽} Dobree \espApp ἐστὶ R : \lecon{ἐστὶν} Küster \espApp δηλαδῆ R : \lecon{δῆλα δὲ} Grynaeus \ligneApp{805} χεῖρον R : \lecon{χείρων} Zanetti}\\
παραβάλλουσαι τῆς τε γυναικὸς καὶ τἀνδρὸς τοὔνομ᾽ ἑκάστου.\\
Ναυσιμάχης μέν γ᾽ ἥττων ἐστὶν Χαρμῖνος· δῆλα δὲ τἄργα.%
\commentaire{804}{Ναυσιμάχης μέν ἥττων ἐστὶν Χαρμῖνος·}{Charminos, stratège athénien, qui venait de se faire battre sur mer près de Symé\marqueur{1}, est rapproché de Nausimaché (\italique{bataille navale}). En quoi Charminos est-il inférieur (ἥττων) à Nausimaché ? Certes par le nom -- comme Aristophane n'avait pas manqué de prévenir : après les déboires qu'il fit subir à la marine athénienne, il serait bien indigne d'en porter un de ce genre !.. mais en rester à cet unique rapprochement serait bien léger\marqueur{2}. La précision de la scholie -- Ναυσιμάχη πόρνη καὶ Σαλαβακχώ -- n'est, pour une fois, pas une fantaisie de Byzantin ; Nausimaché ne surpasse pas Charminos dans le \italique{combat naval} que par le nom : de par sa profession, elle devait aussi y exceller \italique{en pratique} (τἄργα) -- métaphoriquement, s'entend\marqueur{3} --, ce qui lui valut à la fois son renom et son surnom.}%
\debutNote{804}%
\note{1}{Cf. Thucydide, VIII, 41.}%
\note{2}{À l'instar des précédents commentateurs qui, du haut de leur tour d'érudition, ignorent l'essentiel : Aristophane est drôle !}%
\note{3}{La langue grecque, et celle d'Aristophane en particulier, fourmille de métaphores nautiques pour rendre de façon imagée les corps-à-corps charnels ; ces images remplissent les sections 179 à 184 des \italique{Images d'Aristophane} de Taillardat et les pages 161-166 de la \italique{Maculate Muse} d'Henderson. Pour faire écho à Nausimaché, ne citons que le verbe ναυμαχεῖν usité par Aristophane aux vers 430 des \italique{Grenouilles} et 675 de \italique{Lysistrata}. Quant à τἄργα, on a déjà souligné l'acception scabreuse de ce mot (cf. le vers 204 : ἔργα νυκτερείσια).}\\
καὶ μὲν δὴ καὶ Κλεοφῶν χείρων πάντως δήπου Σαλαβακχοῦς.%
\commentaire{805}{Κλεοφῶν χείρων πάντως δήπου Σαλαβακχοῦς.}{Salabachô est un nom générique de prostituée\marqueur{1}. Aristophane le mentionne aussi au vers 765 des \italique{Cavaliers} : (c'est le Paphlagonien qui parle) :
\citComm{Κλέων --- τῇ μὲν δεσποίνῃ Ἀθηναίᾳ τῇ τῆς πόλεως μεδεούσῃ\\
εὔχομαι, εἰ μὲν περὶ τὸν δῆμον τὸν Ἀθηναίων γεγένημαι\\
βέλτιστος ἀνὴρ μετὰ Λυσικλέα καὶ Κύνναν καὶ Σαλαβακχώ,\\
ὥσπερ νυνὶ μηδὲν δράσας δειπνεῖν ἐν πρυτανείῳ·\\
Cléon --- \francais{À notre maîtresse Athéna, qui régit cette cité,\\
je fais cette prière ; si pour le peuple des Athéniens je me suis montré\\
le plus méritant des hommes après Lysiclès, Cynna et Salabacchô,\\
puissé-je comme maintenant, sans avoir rien fait, dîner dans le Prytanée.}}%
L'exégèse de ce passage nous permettra de comprendre du même coup celui qui nous intéresse.
Willems avait justement remarqué que les trois noms étaient mis ici par jeu de consonance avec les trois représentants de l'Athènes impérialiste : Périclès, Cimon et Thémistocle\marqueur{2}. 
Pour le sens, voici ce qu'il en est : la modestie -- toujours fausse -- de Cléon, qui se place après ces trois-là, n'est que pure affectation : il fait comprendre qu'il cumule en fait les mérites des trois.
Λυσι-κλῆς, il l'est depuis sa victoire de Sphactérie : il acquit la \italique{célébrité} en \italique{déliant} la situation empêtrée du siège de Sphactérie.
Κύννα, il prétend l'être, lui qui ne cesse de proclamer être le \italique{chien} de garde de la cité\marqueur{3}.
Σαλα-βακχώ se laisse aisément comprendre ; on rapprochera la première partie du nom d'une glose de Photius (σ 43) :%
\citComm{σάλα· φροντίς. οὕτως Αἰσχύλος.}%
On reprochait au braillard Cléon d'être un excité ; il a beau jeu de prétendre l'être en vertu de sa \italique{préoccupation}, du \italique{soin} qu'il porte à la cité. La seconde partie du nom, βακχώ, étant une allusion à l'agitation de tous les fidèles de Bacchus, et notamment les bacchantes. Suétone, dans son Περὶ βλασφημιῶν, \italique{Des termes injurieux}, cite βακχός au chapitre εἰς ἀλαζόνας, «~envers les bavards, les prétentieux~»\marqueur{4}.
%C'est pour la même raison qu'Aristophane compare Salabacchô et Cléophon ; l'avantage est bien-sûr donné à Salabacchô.
\textbf{Κλεοφῶν} était un démagogue dont la trop longue `carrière politique' fut brutalement interrompue en 404 av. J.-C. par une condamnation à mort. Dans les \italique{Grenouilles}, Aristophane charge ce va-t-en-guerre\marqueur{5} dont la langue par trop pendue et la lippe par trop pendante n'émettaient qu'un babil barbarisant criant ses barbaresques origines\marqueur{6}. (Pour un dossier complet sur Cléophon, le lecteur pourra consulter le commentaire de Fritzsche \italique{ad loc.}) La glossolalie de Cléon trouve un écho en Σαλαβακχώ -- mélioratif pour Σαλαβακχώ puisque Cléon n'hésite pas à se comparer à elle -- c'est ce en quoi Σαλαβακχώ emporte l'avantage sur Cléophon.}%
\debutNote{805}%
\note{1}{Hésychius (σ 88) :%
\citComm{Σαλαβακχώ· πόρνης ὄνομα. ἀπὸ δὲ ταύτης καὶ τὰς καταφερεῖς εἰς τὰ ἀφροδίσια οὕτως ἔλεγον Ἀττικοί.}%
}%
\note{2}{Platon prit les trois mêmes comme exemples d'hommes d'état honorés alors qu'en réalité, ils sont les auteurs méconnus de la crise morale athénienne :%
\citComm{ἄνευ γὰρ σωφροσύνης καὶ δικαιοσύνης λιμένων καὶ νεωρίων καὶ τειχῶν καὶ φόρων καὶ τοιούτων φλυαριῶν ἐμπεπλήκασι τὴν πόλιν.\\
sans se préoccuper de la sagesse ni de la justice, ils ont gorgé la ville de ports, d'arsenaux, de murs, de tributs \& autres niaiseries.}%
(\italique{Gorgias}, 519a.) 
Le détournement d'Aristophane n'est pas plus la marque d'une grande déférence envers ces trois-là.}%
\note{3}{Même jeu de mots avec le nom de Κύννα au vers 1032 des \italique{Guêpes} : ἀπ᾽ ὀφθαλμῶν Κύννης rappelle l'expression homérique des \italique{yeux d'un chien} (\italique{Iliade}, I, 225).}%
\note{4}{Περὶ βλασφημιῶν, Περὶ παιδιῶν, édité en 1967 (aux Belles Lettres) par Jean Taillardat -- taillardesque édition, c'est tout dire ! βακχός est la \nieme{127} βλασφημία.}%
\note{5}{Vers 1504 et 1532-3.}%
\note{6}{Vers 678-682. L'on se reportera pour l'exégèse au toujours brillant Alphonse Willems, t. III, p. 115-117.}\\
πρὸς Ἀριστομάχην δὲ χρόνου πολλοῦ, πρὸς ἐκείνην τὴν Μαραθῶνι,%
\commentaire{806-8}{Ἀριστομάχη, Στρατονίκη et Εὐβούλη}{ sont des noms attestés au moment de la représentation des \italique{Thesmophories}\marqueur{1}. Ces noms sont parlants, bien-sûr, mais l'allusion ; la raison pour laquelle Aristophane les cite, m'échappe. Les deux premières étaient-elles de puissantes princesses contre lesquelles les Athéniens n'auraient pas osé se frotter\marqueur{2} ?}
\debutNote{806-8}
\note{1}{Voir \italique{A Lexicon of Greek Personal Names} de P. M. Fraser et E. Matthews.}%
\note{2}{Στρατονίκη est porté par plusieurs membres de la famille royale macédonienne, dont l'une fut donnée en mariage à l'un des rois de Thrace Σεύθης, contemporain des \italique{Thesmophories} (voir Thucydide, II, 101 et IV, 101). Cicéron (\italique{Tusculanes}, V, 59) nous apprend que Denys lorsqu'il devint Tyran de Syracuse en 406 comptait une certaine \latin{Aristomache} parmi ses épouses.}\\\passeLigne}
\end{groupe}%
%% v. 785-843
\begin{groupe}%
\ledpb%
\tetramAnap{καὶ Στρατονίκην ὑμῶν οὐδεὶς οὐδ᾽ ἐγχειρεῖ πολεμίζειν.\apparat{807}{\lecon{Στρατονίκην} \PsiUn : Στρατωνίκην R}\\
ἀλλ᾽ Εὐβούλης τῶν πέρυσίν τις βουλευτής ἐστιν ἀμείνων\testimonia{808-9}{(τις βουλευτής -- βουλείαν) Suidas β 431 (ἐστι pro ἐστιν et ἀμείνων omisso)}\apparat{808}{τίς βουλευτής R : \lecon{τις βουλευτής} Küster}
\commentaire{808-9}{ἀλλ᾽ Εὐβούλης τῶν πέρυσίν τις βουλευτής ἐστιν ἀμείνων / παραδοὺς ἑτέρῳ τὴν βουλείαν; οὐδ᾽ αὐτὸς τοῦτό γε φήσεις.}{Ces deux vers ont été diversement interprétés. Interprétation d'autant plus importante que ce passage est peut-être le seul de la comédie qui puisse nous aider véritablement à dater la première représentation des \italique{Thesmophories}. (Je renvoie à l'Introduction sur la question de la date de notre pièce.)\\
Deux écoles s'affrontent :\\
--- celle de Wilamowitz-Willems\marqueur{1} qui situe cette représentation aux Dionysies de 411 : le Conseil des Cinq-Cents bouleutes fut rendu tout-à-fait caduc par l'instauration d'un conseil de dix vieillards, les πρόβουλοι\marqueur{2}, magistrats à compétence presque illimitée qui prirent en main la direction des affaires. Willems continue : «~C'était délivrer au Conseil des Cinq-Cents un certificat d'incapacité, le droit de formuler un avis préalable étant précisément l'essentiel de sa prérogative. Aussi, quand il sortit de charge en skirophorion (juin) 412, n'existait-il plus que de nom. Et si, ce que nous ignorons, l'usage prévalait déjà, comme au temps de Démosthène, de décerner à la Boulè démissionnaire une couronne d'or pour prix de ses services, il est bien sûr que cette fois nul ne se leva à l'Ecclésie pour proposer pareil décret [...] C'est à cette transmission annuelle de pouvoirs (παράδοσις, c'est le terme consacré), par laquelle chaque bouleute sortant cédait sa place à son successeur, que le chœur fait allusion~» ;\\
--- celle de Rogers -- et à laquelle je me rattache qui considère que la pièce est d'un an plus tardive : lorsque les Quatre Cents conjurés favorables à l'oligarchie se rendirent au βουλευτήριον au mois de mai 411, ils versèrent leur salaire aux bouleutes alors en séance. Ils leur intimèrent alors l'ordre de se retirer, ce que ces derniers firent sans même que les conjurés eussent à recourir à la violence\marqueur{3}. Ainsi périt une première fois la démocratie athénienne. C'est ce manque de courage (ἀμείνων) des bouleutes qu'Aristophane brocarde : pas un n'aura défendu la charge qu'ils avaient reçue. Que ces bouleutes couards fussent ceux de l'\italique{an dernier} (τῶν πέρυσίν) place donc notre pièce aux Dionysies de l'an 410.}%
%{Des différentes explications apportées à ce passage, nous ne retiendrons que celle de Willems\marqueur{1} -- la seule plausible\marqueur{2} : le Conseil des Cinq-Cents bouleutes fut rendu tout-à-fait caduc par l'instauration d'un conseil de dix vieillards, les πρόβουλοι\marqueur{3}, magistrats à compétence presque illimitée qui prirent en main la direction des affaires. Willems continue : «~C'était délivrer au Conseil des Cinq-Cents un certificat d'incapacité, le droit de formuler un avis préalable étant précisément l'essentiel de sa prérogative. Aussi, quand il sortit de charge en skirophorion (juin) 412, n'existait-il plus que de nom. Et si, ce que nous ignorons, l'usage prévalait déjà, comme au temps de Démosthène, de décerner à la Boulè démissionnaire une couronne d'or pour prix de ses services, il est bien sûr que cette fois nul ne se leva à l'Ecclésie pour proposer pareil décret [...] C'est à cette transmission annuelle de pouvoirs (παράδοσις, c'est le terme consacré), par laquelle chaque bouleute sortant cédait sa place à son successeur, que le chœur fait allusion~»}%
\debutNote{808-9}%
\note{1}{Voir la \nieme{23} note de Willems, reproduite entièrement à la fin.}%
%\note{2}{D'aucuns, trompés par le scholiaste, ont prétendu qu'il y avait dans ces vers une pointe d'Aristophane contre un bouleute précis, non nommé... Mais ne voit-on pas qu'avec la construction τις ... ἀμείνων, τις ne peut avoir qu'une valeur indéfinie ? « Citez m'en un seulement de meilleur », et que si Aristophane pensait à un bouleute en particulier, il aurait dû user d'une construction analogue aux vers précédents avec ἥττων, χείρων ou tout autre synonyme.}%
\note{2}{Voir Thucydide VIII, 1 et la \italique{Constitution d'Athènes}, XXIX ; Aristophane met en scène un Probule dans \italique{Lysistrata} (v. 387-613), jouée quelques mois auparavant -- là est la question : combien ? deux mois comme le pense Willems ?.. ou quatorze comme le pense Rogers ?}%
\note{3}{Événement conté par Thucydide aux chapitre 69-70 de son livre VIII.}\\
παραδοὺς ἑτέρῳ τὴν βουλείαν; οὐδ᾽ αὐτὸς τοῦτό γε φήσει.\testimonia{809}{(παραδοὺς ἑτέρῳ τὴν βουλείαν) \Antiatt 84.30 \espTest (βουλείαν) \Harp B 19 ; β Photius β 233 βουλία ; Hesychius β 921 βουλεῖα}\apparat{809}{\lecon{αὐτὸς} R : Ἄνυτος Maas \espApp φήσεις R : \lecon{φήσει} Küster \ligneApp{810} εὐχόμεσθ᾽ R : \lecon{εὐχόμεθ᾽} \Mon}\\
οὕτως ἡμεῖς πολὺ βελτίους τῶν ἀνδρῶν εὐχόμεθ᾽ εἶναι.\\
οὐδ᾽ ἂν κλέψασα γυνὴ ζεύγει κατὰ πεντήκοντα τάλαντα%\apparat{811}{ζεύγει post γυνὴ R : post οὐδ᾽ ἂν van Herwerden}
\commentaire{811-2}{οὐδ᾽ ἂν κλέψασα γυνὴ ζεύγει κατὰ πεντήκοντα τάλαντα / εἰς πόλιν ἔλθοι τῶν δημοσίων·}{Aristophane use de la même formule κατὰ πεντήκοντα τάλαντα, «~cinquante talents d'un coup~», dans les \italique{Guêpes} : Bdélycléon tance son paternel, berné par les démagogues qui, non contents de dissiper l'argent de la cité, détournent le tribut des autres cités (vers 666-671) : %
\citComm{Φιλοκλέων --- καὶ ποῖ τρέπεται δὴ ᾽πειτα τὰ χρήματα τἄλλα;\\
Βδελυκλέων --- εἰς τούτους τοὺς “οὐχὶ προδώσω τὸν Ἀθηναίων κολοσυρτόν,\\
ἀλλὰ μαχοῦμαι περὶ τοῦ πλήθους ἀεί”. σὺ γάρ, ὦ πάτερ, αὐτοὺς\\
ἄρχειν αἱρεῖ σαυτοῦ τούτοις τοῖς ῥηματίοις περιπεφθείς.\\
κᾆθ᾽ οὗτοι μὲν δωροδοκοῦσιν \textbf{κατὰ πεντήκοντα τάλαντα}\\
ἀπὸ τῶν πόλεων ἐπαπειλοῦντες τοιαυτὶ κἀναφοβοῦντες·\\
“δώσετε τὸν φόρον, ἢ βροντήσας τὴν πόλιν ὑμῶν ἀνατρέψω.”\\
Philocléon --- \francais{Et où passe alors le reste de l'argent ?}\\
Bdélycléon --- \francais{À ceux qui nous disent: « Non, je ne trahirai point la masse athénienne, toujours je combattrai pour le populaire ! » Car c'est toi, mon père, qui te les donnes pour maîtres, enjôlé par leurs boniments. Eux cependant extorquent aux cités des cinquante talents à la fois, en les menaçant, en les terrifiant de cette sorte : « Vous paierez le tribut, ou je tonne contre votre ville et la renverse. »}}%
Hérodote (VI, 132-136) rapporte une histoire semblable : celle de Miltiade. Le vainqueur de Marathon, promit aux Athéniens quantité d'or s'ils lui remettaient soixante-dix vaisseaux avec troupe et argent. Les Athéniens lui fournirent le tout et Miltiade de cingler vers Paros pour en faire le siège. Il envoya un héraut réclamer 100 talents avec le même genre de menaces que sus-citées. L'expédition fut un échec et l'on intenta un procès à Miltiade pour avoir trompé la cité (τῆς Ἀθηναίων ἀπάτης εἵνεκεν). Or Miltiade se mourait d'une mauvaise blessure. 
Compte tenu des services rendus à la nation, et que c'était certainement peine perdue de transporter ce semi-cadavre jusqu'au \italique{barathre}\marqueur{1}, l'agonisant ne fut condamné qu'à restituer \italique{cinquante talents}.\\
Avez-vous, vous aussi, été saisi par le rapprochement à faire entre l'histoire de Miltiade et les deux passages d'Aristophane, qui coïncident jusque dans le chiffre du détournement ?
Miltiade, le héros de Marathon -- τὸν Μαραθῶνι, c'est son épithète, quasi-homérique, qui lui est inlassablement accolée --, profita du crédit qu'il avait auprès de Dêmos, le cajolant des souvenirs de Marathon\marqueur{2}, pour mener une expédition dont le but véritable était d'enflouer ses propres caisses.\\
En regard des déboires de Miltiade, la construction de l'argumentation du chœur se comprend parfaitement : «~Devant Aristomaque, notre héroïne de Marathon (τὴν Μαραθῶνι), aucun d'entre vous n'entreprendrait de se battre~», et, parant d'avance à l'objection que les hommes ont aussi leur marathonien héros en la personne de Miltiade, l'on précise aussitôt : «~Une femme ayant volé cinquante talents d'un (bon) coup n'aurait pas l'outrecuidance de venir parader sur un char lors de son retour à la cité~».\\
On finira de saisir tout-à-fait le sens de ces deux vers 811-812 en allant lire la note 24 d'Alphonse Willems (reproduite en annexe) : à la fête des Panathénées, la file brillante des chars était ce qui attirait le plus les regards dans la procession (Schoemann, \italique{Antiquités grecques}, t. II, p. 549).
}%
\debutNote{811-2}%
\note{1}{Platon (\italique{Gorgias}, 516e) nous apprend que Miltiade devait être initialement précipité vers le \italique{barathre} (εἰς τὸ βάραθρον ἐμβαλεῖν), illustrant l'adage latin qu'il n'y a pas loin du Capitole à la roche tarpéienne. (Pour en apprendre davantage sur ce gouffre, que le lecteur se reporte au \italique{Dictionnaire des antiquités grecques et romaines}, à l'entrée \latin{Barathron}.)
Si Miltiade s'est vu gracier -- la Parque se consola néanmoins très rapidement : la grâce ne fut qu'un sursis de quelques jours --, la raison principale est peut-être qu'il bénéficiait de nombreuses relations, critère qui distingue souvent les condamnés à mort exécutés de ceux qui voient leur peine commuée, selon le mot de Gide.}%
\note{2}{«~Flattez-les, payez-les d'agréables mensonges : /	Quelque indignation dont leur cœur soit rempli, / Ils goberont l'appât ; vous serez leur ami.~»}%
\\%
εἰς πόλιν ἔλθοι τῶν δημοσίων· ἀλλ᾽ ἢν τὰ μέγισθ᾽ ὑφέληται\\
\lednopb%
φορμὸν πυρῶν τἀνδρὸς κλέψασ᾽, αὐθημερὸν αὔτ᾽ ἀπέδωκεν\apparat{813}{\lecon{αὔτ᾽ ἀπέδωκεν} R, M\expo{pc} : ἀνταπέδωκεν M\expo{ac}, Bentley}\\\passeLigne}%
\end{groupe}%
%% v. 814-829
\begin{groupe}%
\ledpb%
\personnageGrec{\strophe{πνῖγος}}\typeVers{\decalage{πνῖγος \ }dim. anap.}\anap{ἀλλ᾽ ἡμεῖς ἂν πολλοὺς τούτων%
\commentaire{814}{τούτων.}{Les spectateurs.}\\
ἀποδείξαιμεν ταῦτα ποιοῦντας.\apparat{815}{ἀποδείξομεν R : \lecon{ἀποδείξαιμεν} Bentley \espApp ποθοῦντας R : \lecon{ποιοῦντας} Zanetti (\italique{facientes} Divus)}%\espApp ταῦτα R : ταὐτά` Bibi 
\commentaire{815}{ποθοῦντας.}{L'erreur manifeste du scribe du \italique{Ravennas}, copiant ποθοῦντας à la place de πο(ι)οῦντας, m'est l'occasion -- ὁ καιρός ! -- d'une courte digression portant sur les différentes graphies ποι- ou πο- lorsque la lettre qui suit est un `ο' ou un `ω'\marqueur{1}.
Je m'en tiens aux leçons de l'épigraphie attique ainsi qu'à l'usage scolaire : savoir, écrire systématiquement ποι-. Toutefois, le \italique{Ravennas} adopte majoritairement la graphie πο- lorsque la syllabe est brève. Cette convention, reprise par la plupart des éditions actuelles des \italique{Thesmophories}, est celle de nombreux manuscrits remontant à la période hellénistique. La confusion du scribe -- qui aurait lu rapidement une suite de `Ο' en y ajoutant un `Θ' -- tend à faire penser que cette règle était aussi de mise pour le manuscrit qu'il recopiait.}%
\debutNote{815}%
\note{1}{Débat qui a fait déjà couler beaucoup d'encre et que je ne compte ni trancher, ni même exposer en détail. Le lecteur intéressé pourra se reporter aux pages 226-230 que lui consacre Willem J. W. Koster dans son livre \italique{Autour d'un manuscrit d'Aristophane écrit par Demetrius Triclinius} (1957).}\\
καὶ πρὸς τούτοις γάστριδας ἡμῶν\testimonia{816}{(γάστριδας) Photius γ 40 γάστρις $=$ Suidas γ 72 ; cf. \AelDion γ 3}\\
ὄντας μᾶλλον καὶ λωποδύτας\\
καὶ βωμολόχους κἀνδραποδιστάς.\\
καὶ μὲν δήπου καὶ τὰ πατρῷά γε\apparat{819}{δήπου R : δή που Denniston}%
\commentaire{819-23}{}{%\eff{Euripide Medée fr. 4 ?γυνὴ γὰρ ἐξελθοῦσα πατρῴων δόμων / οὐ τῶν τεκόντων ἐστίν, ἀλλὰ τοῦ λέχους. ἀνδρὸς Ἕλληνος  Γυνὴ μὲν“ἐξελϑοῦσαπατρῴωνδόμων οὐ τῶν “τεκόντων ἐστίν, ἀλλὰ τοῦ λέχους. τὸ δ᾽ ἄρσενἕστηκ'ἐν δόμοιςἀεὶ γέγος ϑεῶνπατρῴωνκαὶ τάφωντιμᾶ ὅρον}
Le manuscrit présente le texte suivant :%
\citComm{καὶ μὲν δήπου καὶ τὰ πατρῷά γε\\
χείρους ἡμῶν εἰσιν σῴζειν·\\
ἡμῖν μὲν γὰρ σῶν ἔτι καὶ νῦν\\
τἀντίον ὁ κανὼν οἱ καλαθίσκοι\\
τὸ σκιάδειον·}%
Les trois premiers vers ne posent pas de difficulté. Et encore : à condition de garder à l'esprit les divers sens de τὰ πατρῷα, \italique{l'héritage} : désigne les \italique{biens transmis par les ancêtres} ; \italique{héritage} tout matériel que seul une femme sera à même de conserver\marqueur{1}, mais héritage aussi immatériel, \italique{les traditions}\marqueur{2}, que les femmes perpétuent tout-de-même moins mal que les hommes (χείρους ἡμῶν εἰσιν σῴζειν).
Suit une liste de ces biens que les femmes sauvent mais que leur mari, eux, perdent. C'est là que le bât blesse.
D'abord, cette liste est étonnamment hétéroclite ; si les trois premiers éléments qu'elle comporte sont en rapport avec le filage et le tissage : τὸ ἀντίον, \italique{l'ensouple} (cyclindre du métier à tisser) ; ὁ κανὼν, \italique{barre}, \italique{règle} (avec le sens spécialisé de \italique{tige de quenouille}), οἱ καλαθίσκοι, \italique{petits paniers} (où les femmes déposaient la laine filée), l'on se demande par quel ressort \italique{l'ombrelle}, τὸ σκιάδειον, vient clôturer cette liste.
Plus important : sur ces quatre éléments, seuls deux, κανών et σκιάδειον, subsistent dans la comparaison avec ce qu'en font les hommes. Pourquoi Aristophane nous parle-t-il alors de l'ἀντίον et des καλαθίσκοι, inutiles -- donc nuisibles -- si ce n'est pas pour filer la métaphore quelques vers plus loin ? (Je ne parle que de ce qui achoppe au niveau du sens. Quant à la forme, l'on avait déjà relevé l'inhabituelle succession de quatre brèves au vers 822.)\\
Comme souvent en pareille situation, le mieux est encore de se mettre à la place d'Aristophane. Il veut nous peindre la sotte dissipation des hommes ; hé bien, comment évoque-t-on une prodigalité de cigale ou une gestion de fourmi dans sa langue imagée et proverbiale ? Interrogeons les parémiographes\marqueur{3} :%
\citComm{Αὐτῷ κανῷ: αὕτη τάττεται κατὰ τῶν ἄρδην ὁτιοῦν λυμαινομένων. Φησὶ γὰρ ὁ Δίδυμος ὅτι τὰ δεῖπνα ἐπὶ κανῶν ἐκομίζετο. Καθάπερ ἐν Γρυψὶ λέγει καὶ Πλάτων·\\
Ἀνήρπακας ἁπαξάπαντ᾿ \textbf{αὐτῷ κανῷ}.\\
Et la corbeille elle-même : ce [proverbe] se range contre ceux prompts à se ruiner de quelque manière. Didyme dit que les repas étaient emportés dans des paniers. Ainsi Platon dans ses \italique{Griffons} :\\
Tu as tout emporté, tout, \italique{et le panier avec} !..}%
Le grec ajoute «~et le panier~» pour désigner tous les biens jusqu'au dernier\marqueur{4}, comme en français on ajoute «~et son père~» pour désigner «~tout le monde~». Dans les \italique{Grenouilles} (vers 560), ce n'est pas autrement qu'il faut entendre le désarroi de l'hôtesse face à Hercule, cet ogre qui a liquidé toutes ses provisions :%
\citComm{Μὰ Δί᾽ οὐδὲ τὸν τυρόν γε τὸν χλωρόν, τάλαν,\\
ὃν οὗτος \textbf{αὐτοῖς τοῖς ταλάροις} κατήσθιεν.\\
Par Zeus, et ce fromage frais, pauvre amie, qu'il dévorait \italique{avec les éclisses elles-mêmes}.}%
C'est peut-être encore ce même proverbe qu'on retrouve au vers 170 des \italique{Guêpes} dans la bouche de Philocléon :%
\citComm{μὰ τὸν Δί᾽ οὐ δῆτ᾽, ἀλλ᾽ ἀποδόσθαι βούλομαι\\
τὸν ὄνον ἄγων \textbf{αὐτοῖσι τοῖς κανθηλίοις}·\\
νουμηνία γάρ ἐστιν.\\
Par Zeus, certes non ! Je veux aller vendre l'âne \italique{avec ses paniers} ; car c'est nouvelle lune.}%
L'«~αὐτῇ λόγχῃ~» du vers 826 -- noter la présence du pronom --, mis comme permutant métaphorique d'«~αὐτῷ κανόνι~», rappelle le détournement du proverbe avec jeu de mots sur κανών / κανοῦν.\\
Notre conclusion est que le texte original des \italique{Thesmophories} -- l'\allemand{Urtext} -- comportait certainement au vers 822 le terme κανοῦν, \italique{corbeille}, et non κανών. 
Cette hypothèse permet de faire coup double et d'expliquer la singulière présence d'«~οἱ καλαθίσκοι~», \italique{les corbeilles à laine}, qui n'est que le fruit d'une glose de κανοῦν fort importunément ajoutée au texte.\\
L'image ne porte pas sur l'activité filandière ; la présence simultanée de corbeilles et d'une ombrelle ne laisse pas de doute possible : Aristophane nous décrit une procession religieuse\marqueur{5}. 
Au sein de ces cortèges traditionnels, on voyait défiler les \italique{cané-}phores, qui avaient la charge des corbeilles à offrandes, et qu'elles portaient sur la tête\marqueur{6}, quand d'autres, les \italique{skiro-}phores, les accompagnaient\marqueur{7} en tenant de deux mains fermes et assurées, la lourde \italique{ombrelle} consacrée, τὸ σκίρον ou τὸ σκιάδειον. 
Et pour Aristophane, sa foi, ces parades d'Athéniennes défilant cahin-caha n'étaient pas tout-à-fait indignes des traditions reçues des ancêtres (τὰ πατρῷα).
À la place des vers 822-823, je propose donc de lire :%
\citComm{καὶ μὲν δήπου καὶ τὰ πατρῷά γε\\
χείρους ἡμῶν εἰσιν σῴζειν·\\
ἡμῖν μὲν γὰρ σῶν ἔτι καὶ νῦν\\
ταὐτὸ κανοῦν\marqueur{8} [καὶ] τὸ σκιάδειον·\\
Et assurément aussi, ils sont moins aptes que nous\\
à garder les traditions\\
qui sont, par notre fait, sauves encore aujourd'hui\\
-- même le panier [et] l'ombrelle.}%
%ταὐτό
%Hésychius ἐπίρρηται: ἐπικάτακται
%τανταλίζεται: σαλεύεται (Αr fr. 963 Κ.-Α.
}%
\debutNote{819-23}%
\note{1}{L'homme qui avait aliéné son patrimoine, ὁ τὰ πατρῷα κατεδηδοκώς, était frappé d'ἀτιμία. Eschine, \italique{In Timarchum}, 30 ; Diogène Laërce, \italique{Solon}, I, 55. Cette loi qui n'était certainement plus observée au temps d'Eschine subsistait pour la forme, comme un vestige de l'antique règle ; il y eut toujours une δίκη κατεδηδοκέναι τὰ πατρῷα (Bekker, \italique{Anecdota}, p. 199 et 310 ; note de Fustel de Coulanges).}%
\note{2}{Certains commentateurs n'ont pas manqué de noter que les Athéniennes n'auraient pas eu légalement de droit de propriété. Cette remarque aurait pu -- aurait dû -- leur mettre la puce à l'oreille : c'est bien en tant que championnes des \italique{traditions} que se posent les femmes du chœur dans ce passage.}%
\note{3}{Leutsch \& Schneidewim, \italique{Corpus Paroemiographorum Graecorum}, vol. I, Zenobius, II, 31. Recensé par Érasme (\italique{Adages}, n° 1088). Kassel \& Austin l'ont intégré comme \nieme{15} fragment de Platon le Comique. Suidas (α 4465 ; α 4466 ; α 4540) a relevé des expressions analogues : αὐταῖς ἀμάλαις ; αὐταῖσι διαβολαῖς ; αὐτοῖς ὅπλοις ; αὐτοῖς ταλάροις.}%
\note{4}{Biens dévorés, en l'occurrence. Mais l'image du panier comme signe d'abondance se retrouve dans un autre proverbe, transmis par Hésychius (ε 7077) : εὔπλουτον κανοῦν.}%
\note{5}{La fête des Σκίρα est évoquée quelques vers plus loin. Cf. les articles \latin{Umbella} et \latin{Skirophoria} du \italique{Dictionnaire des antiquités grecques et romaines} de Daremberg \& Saglio.}%
\note{6}{Ces canéphores ont certainement servi de modèle aux fameuses Caryatides, portant, elles, un chapiteau sur la tête.}%
\note{7}{Voir, par exemple, le vers 1550 des \italique{Oiseaux}. Les ombrelles tenaient effectivement un rôle important dans la liturgie.}%
\note{8}{Au vers 164, il y a confusion de finale -ων / -ουν du même genre entre les leçons du \italique{Ravennas} (διεκίνων) et de la \italique{Souda} ε 989 (διεκίνουν).}%
\\
χείρους ἡμῶν εἰσιν σῴζειν·\\\passeLigne}
\end{groupe}%
%% v. 814-829
\begin{groupe}%
\ledpb%
\anap{ἡμῖν μὲν γὰρ σῶν ἔτι καὶ νῦν\\
ταὐτὸ κανοῦν καὶ τὸ σκιάδειον·\testimonia{822}{(τἀντίον ὁ κανὼν) Pollux VII 35 κανὼν ἱστοῦ τὸ καλούμενον ἀντίον}\\
% \commentaire{821-3}{τἀντίον ὁ κανὼν οἱ καλαθίσκοι / τὸ σκιάδειον·}{Sur le sens exact de ces mots, voir l'excellente thèse de M. Fr. Sosset, \italique{Le Tissage dans la Grèce antique}, Bruxelles (1896), pp. 11 à 25. (Extrait de la \italique{Revue de l'Université de Bruxelles}. Note de Willems.)}%
\apparat{822-3}{τἀντίον ὁ κανὼν οἱ καλαθίσκοι / τὸ σκιάδειον· R : \lecon{ταὐτὸ κανοῦν καὶ τὸ σκιάδειον·} Bibi \ligneApp{824} ἀνδράσιν R : \lecon{ἀνδράσι} \Mon \ligneApp{825} ἀπόλωλε R : \lecon{ἀπόλωλεν} Daubuz, Bentley}%
\ligne{824}τοῖς δ᾽ ἡμετέροις ἀνδράσι τούτοις\\
ἀπόλωλεν μὲν πολλοῖς ὁ κανὼν\\
ἐκ τῶν οἴκων αὐτῇ λόγχῃ,%
\commentaire{826}{ἐκ τῶν οἴκων.}{Précisons le sens «~ἐκ τῶν οἴκων~». Ce ne sera pas du luxe : les traductions actuelles étant du plus haut comique ; bien comprises, elles signifient que le reproche des femmes porte sur ce que les maris rompent leur lance \italique{à l'extérieur des maisons}. Préféreraient-elles qu'ils le fassent \italique{à l'intérieur} ? Les commentateurs vont encore plus loin : ils nous expliquent -- que dis-je : ils nous prouvent -- que le fameux héritage paternel, τὰ πατρῷα, dissipé par les hommes n'est autre que la panoplie militaire paternelle -- ô scandale justifiant de couvrir d'infamie tous ces fils prodigues indignes !! Au garde-à-vous, je félicite ces brillants universitaires pour leurs éclaircissements : sans eux, je crois que les douceurs martiales de la poésie d'Aristophane m'auraient échappé. Notre poète roulant des variations sur le thème : «~Tu seras un soldat, mon fils !~», c'est à vous tirer des larmes (de guerre, bien-sûr !)\\
Bref !.. Sortons d'Oxford cinq minutes et remettons les pieds sur Terre. 
Οἶκος, ce mot rare et qui renvoie à des réalités n'existant pas ailleurs qu'en Grèce classique, a tout-de-même donné le terme οἰκο-νομία, qui me semble être la bonne gestion d'une \italique{maison}, perçue comme un \italique{fonds} ou un \italique{capital}. «~ἐκ τῶν οἴκων~» désigne tout ce que l'on \italique{tire de} ce fonds, ou de ce capital. La bonne équation est évidemment celle-ci : τὰ πατρῷα $ = $ ἐκ τῶν οἴκων.
La même tournure avec le même sens se trouve chez Xénophon dans son ... \italique{Économique} (XX, 21) :%
\citComm{τὰ οὖν συντρίβοντα τοὺς οἴκους πολῦ μᾶλλον ταῦτά ἐστιν ἢ αἱ λίαν ἀνεπιστημοσύναι. τὸ γὰρ τὰς μὲν δαπάνας χωρεῖν ἐντελεῖς \textbf{ἐκ τῶν οἴκων}, τὰ δὲ ἔργα μὴ τελεῖσθαι λυσιτελούντως πρὸς τὴν δαπάνην\\
Voilà qui ruine une maison beaucoup plus sûrement qu'une excessive ignorance : d'un côté, dépenser complètement \italique{ce qu'on tire des maisons}, de l'autre des travaux qui ne sont pas profitables à la dépense}%
Traduisez donc :%
\citComm{τοῖς δ᾽ ἡμετέροις ἀνδράσι τούτοις\\
ἀπόλωλεν μὲν πολλοῖς ὁ κανὼν\\
ἐκ τῶν οἴκων αὐτῇ λόγχῃ\marqueur{1}.\\
Du fait de nos ci-devant maris (elles les désignent),\\
la lance a anéanti pour beaucoup\\
ce qu'on tirait de nos maisons -- lance comprise\marqueur{1}.}%
}%
\debutNote{826}%
\note{1}{Écho au proverbe : αὐτῷ κανῷ... sauf que le \italique{canon} est remplacé par la \italique{lance}. Autre sexe, autre proverbe :%
\citComm{ἦν γάρ τις αἶνος, ὡς γυναιξὶ μὲν τέχναι\\
μέλουσι, \textbf{λόγχῃ δ ̓ ἄνδρες} εὐστοχώτεροι.\\
εἰ γὰρ δόλοισιν ἦν τὸ νικητήριον,\\
ἡμεῖς ἂν ἀνδρῶν εἴχομεν τυραννίδα.\\
Il existe en effet un proverbe qui dit qu'à la femme les artifices reviennent, tandis que \italique{les hommes sont plus habiles à atteindre le but par la lance}. Si le prix de la victoire allait à la ruse, c'est nous qui régnerions sur les hommes.}%
(Euripide, fr. 6 \italique{Danaé} Jouan-van Looy $=$ fr. 321 Kannicht $=$ fr. 413 Mette.)
}\\
πολλοῖς δ᾽ ἑτέροις ἀπὸ τῶν ὤμων\\
\{ ἐν ταῖς στρατιαῖς \}%
}%
\end{groupe}%
%% v. 814-829
\begin{groupe}%
\typeVers{(catalect.)}\anap{\lednopb%
ἔρριπται τὸ σκιάδειον.%
\commentaire{827-9}{πολλοῖς δ᾽ ἑτέροις ἀπὸ τῶν ὤμων / ἐν ταῖς στρατιαῖς / ἔρριπται τὸ σκιάδειον.}{Un bouclier est une \italique{ombrelle}. Cette métaphore unique est très claire : un bouclier est rond (cf. κύκλος, Taillardat § 642) et bombé comme une ombrelle. Mais on peut se demander si l'argot des soldats athéniens ne connaissait aussi pas cette métaphore unique\marqueur{1} : le bouclier protège des coups comme l'ombrelle du soleil, cf. schol. : τὴν ἀσπίδα λέγει· σκιάζει\marqueur{2} γὰγ τὸν ἄνδρα. Cette hypothèse est malheureusement invérifiable.
(Note de Taillardat, §643. L'utilisation deux vers plus haut de κανών comme synonyme de \italique{lance}, λόγχη, pourrait relever de ce même argot militaire dont Taillardat fait la très astucieuse hypothèse.)\\
La charge soldatesque des hommes ? elle leur passe par-dessus la tête et les épaules lorsqu'ils en délestent leur bouclier (ἀπὸ τῶν ὤμων ἔρριπται τὸ σκιάδειον), bien encombrant il est vrai lorsqu'il faut prendre la poudre d'escampette en plein milieu des combats... 
Mica, la femme de Cléonyme -- Cléonyme qui semble avoir fait école avec sa célèbre ῥίψασπις (voir le commentaire au vers 605) -- devait probablement à ce moment être pointée, au moins des yeux.}%
\debutNote{827-9}%
\note{1}{À propos de θύλακος (§268), on [l'«~on~», ou l'«~om~» en question n'est autre que le grand Taillardat !] a déjà relevé une métaphore ironique qui a dû appartenir à l'argot des soldats athéniens. Cette tendance à ironiser sur ses armes a toujours existé : on sait que les soldats de César appelaient, par plaisanterie, \latin{cippi} et \latin{lilia} certaines défenses qu'ils construisaient devant Alésia (\italique{Bell. Gall.} 7, 73, 4 et 8). C'est de la même façon que dans l'argot militaire français, un fusil peut s'appeler une \italique{arbalète}, une \italique{seringue}, une \italique{canne à pêche}, etc. (voir A. Dauzat, \italique{L'argot de la guerre}\expo{2} (1919).)}%
\note{2}{Σκιαζειν, \italique{couvrir} : Hésiod. \italique{Théog.} 716 ; Hdt. 6, 117, etc.}%
\\\passeLigne}%
\end{groupe}%
%% v. 830-845
\begin{groupe}%
\ledpb%
\typeVers{$4$ troch\cat}\personnageGrec{Κορυφαῖος}\tetramTroch{πόλλ᾽ ἂν αἱ γυναῖκες ἡμεῖς ἐν δίκῃ μεμψαίμεθ᾽ ἂν\\
τοῖσιν ἀνδράσιν δικαίως, ἓν δ᾽ ὑπερφυέστατον.\\
χρῆν γάρ, ἡμῶν εἰ τέκοι τις ἄνδρα χρηστὸν τῇ πόλει,\apparat{832}{χρὴ R : \lecon{χρῆν} Brunck}\\
ταξίαρχον ἢ στρατηγόν, λαμβάνειν τιμήν τινα,\\
προεδρίαν τ᾽ αὐτῇ δίδοσθαι Στηνίοισι καὶ Σκίροις\testimonia{834}{(Στηνίοισι) Suidas σ 1089 Στήνια καὶ Σκίρα ($=$ Scholium) ; Hesychius σ 1825 στήνια}\apparat{834}{Τηνίοισι R : \lecon{Στηνίοισι} Schol., Suid. σ 1089}
\commentaire{834}{Στηνίοισι καὶ Σκίροις.}{Fêtes athéniennes réservées aux femmes. C'est notamment aux Sciries que fut élaboré le projet politique de Praxagora dans l'\italique{Assemblée des femmes} (vers 18).}\\
ἔν τε ταῖς ἄλλαις ἑορταῖς αἷσιν ἡμεῖς ἤγομεν·\\
εἰ δὲ δειλὸν καὶ πονηρὸν ἄνδρα τις τέκοι γυνή,\apparat{836}{ἤν δὲ R : \lecon{εἰ δὲ} Brunck}\\
ἢ τριήραρχον πονηρὸν ἢ κυβερνήτην κακόν,%
\commentaire{833-7}{ταξίαρχον ἢ στρατηγόν ... ἢ τριήραρχον πονηρὸν ἢ κυβερνήτην κακόν}{Les exemples d'Aristophane sont tirés de l'actualité la plus brûlante pour ce patriote pacifiste -- pacifiste parce que patriote. Les mauvais triérarques et les mauvais commandants de vaisseaux sont ceux responsables de la défaite navale d'Érétrie à l'été 411, défaite qui plongea les Athéniens dans une consternation, nous rapporte Thucydide (VIII, 95 et 96), plus grande encore que celle consécutive au désastre de Sicile. Défaite cinglante qui ne rattrape pas aux yeux d'Aristophane la toute récente et lointaine victoire de Κυνὸς Σῆμα (VIII, 104-108), dans l'Hellespont. (Cette victoire devait être bien amère à Aristophane, puisque son contemporain Eupolis y mourut. Nous y reviendrons.)}\\
ὑστέραν αὐτὴν καθῆσθαι σκάφιον ἀποκεκαρμένην\testimonia{838}{(σκάφιον ἀποκεκαρμένην) Photius σ 291 σκάφιον ἀποκεκαρμένην ; cf. Hesychius σ 879 σκάφιον ; Photius σ 292 ($\simeq$ Scholium) ; \Eustathe 907.40, 1292.61 ; \AelDion σ 23}\apparat{838}{\lecon{ὑστέραν} (ex ὕστερον) \Rpc : ὑστάτην Schol.}%
\commentaire{838}{σκάφιον ἀποκεκαρμένην.}{C'est l'humiliation que le tyran de Cumes, Aristodème, avait fait subir aux filles de son pays : περιτρόχαλα κείρεσθαι, tandis qu'il forçait les garçons à porter des cheveux longs. Plutarque, \italique{De mul. virtut}, 26, p. 361f. (Note de Willems.)\\
La scholie précise à propos de σκάφιον : %
\citComm{εἶδος κουρᾶς δουλικῆς.\\
\francais{Genre de coupe rase qui convient à l'esclave.}}%
}\\
τῆς τὸν ἀνδρεῖον τεκούσης· τῷ γὰρ εἰκὸς ωὖπολι\apparat{839}{ὦ πόλεις R : ὦ πόλις Gelenius : \lecon{ωὖπολι} Bibi}%
\commentaire{839}{ὦ πόλεις.}{Dans son livre \italique{Aristophanes und Eupolis} (2007), Natalia Kyriakidi a montré qu'Eupolis, le poète comique rival d'Aristophane, est selon toute vraisemblance mort en 411 lors de la victoire navale de Κυνὸς Σῆμα dont on a déjà parlé plus haut. Pour arriver à cette conclusion, elle récuse très justement la plupart des témoignages, qui exhalent l'anecdote recyclée, pour ne conserver que le témoignage de Suidas (ε 3657) -- ou de la \italique{Souda}, si l'on préfère :%
\citComm{Εὔπολις, Σωσίπολις, Ἀθηναῖος, κωμικός. ὃς ἑπτακαίδεκα ἔτη γεγονὼς ἤρξατο ἐπιδείκνυσθαι: καὶ ἐδίδαξε μὲν δράματα ιζ, ἐνίκησε δὲ ζ. καὶ ἀπέθανε ναυαγήσας κατὰ τὸν Ἑλλήσποντον ἐν τῷ πρὸς Λακεδαιμονίους πολέμῳ: καὶ ἐκ τούτου ἐκωλύθη στρατεύεσθαι ποιητήν. τὰ δὲ δράματα αὐτοῦ Αἶγες, Ἀστράτευτος ἢ Ἀνδρόγυνοι, καὶ ἄλλα.\\
Fils de Sosipolis, Athénien, un poète comique. Il a commencé à monter des pièces de théâtre à l'âge de dix-sept ans. Il a produit 17 pièces et en a gagné 7. Il est mort dans le naufrage de son navire dans l'Hellespont pendant la guerre contre les Lacédémoniens ; à la suite de cela, il fut interdit à un poète de servir en campagne. Ses pièces sont les \italique{Chèvres}, l'\italique{Exempté} ou les \italique{Androgynes}, et d'autres.}%
Natalia continue (p. 11) : \allemand{Die zwei großen Schlachten, die am Hellespont während des Peloponnesischen Krieges stattfanden, sind der Sieg der Athener bei} Κυνὸς Σῆμα \allemand{411 und die endgültige Niederlage bei }Αἰγὸς Ποταμοί \allemand{405. Der Tod des Eupolis während einer Seeschlacht könnte auch die Entstehung der Geschichte vom Ertränken des Eupolis durch Alkibiades erklären} [$=$ Eupolis, selon certaines légendes, aurait été noyé en mer par Alcibiade au retour de Sicile], \allemand{weil beide Geschichten als Vorfälle im Meer bei einer Kriegshandlung wenige Jahre vor dem Ende des Peloponnesischen Krieges überliefert sind. Storey vermutet, dass Ailians aitiologische Erklärung des Ortsnamens} Κυνὸς Θρῆνος [$=$ le lieu commun du chien qui meurt de chagrin sur la tombe de son maître] \allemand{auf }Κυνὸς Σῆμα \allemand{zurückzuführen ist. Da kein Werk des Eupolis nach 411 zu datieren ist und für die Zeit danach keine Indizien für eine weitere Tätigkeit des Dichters überliefert sind, erscheint die Information der Suda als das einzige zuverlässige Zeugnis über den Tod des Eupolis. So enden sein Leben und seine dichterische Karriere wahrscheinlich im Sommer 411 während der Seeschlacht bei} Κυνὸς Σῆμα.\\
Le texte du manuscrit de \italique{Ravenne} comporte au vers 839 une adresse à la Cité :%
\citComm{τῷ γὰρ εἰκὸς ὦ πόλεις (sic.)}%
% "Est-ce à raison, ô ville, etc." dont l'emphase conviendrait parfaitement à un pastiche de tragédie, mais qui ne paraît point à sa place ici. 
La vraisemblable quasi-coïncidence des dates entre la mort d'Eupolis et la représentation des \italique{Thesmophories}\marqueur{1} m'incite à penser qu'il y a là plutôt une évocation d'Eupolis, auteur intarissable sur la mère d'Hyperbolos, et à laquelle il a même consacré une pièce entière : son \italique{Maricas} (voir le commentaire du vers 840).
S'il faut attribuer l'idée d'accorder les sièges des Athéniennes avec la valeur de leur progéniture à feu Eupolis (et peut-être aussi celle de les tondre le cas décevant échéant), il n'aurait pas manqué de s'indigner lui aussi que la mère d'Hyperbolos spécule avec le fruit de la rapine de son fils. 
Cette hypothèse se peut-elle prouver ? on ne possède malheureusement que peu de choses d'Eupolis et du \italique{Maricas} en particulier. 
On sait que la comédie dut être représentée en $-421$, que le chœur était formé de deux demi-chœurs : celui des riches et celui des pauvres citoyens -- qui s'opposent en cela aux \italique{Cavaliers} d'Aristophane ? --, chœurs qui débattaient avec le démagogue Hyperbolos.
La comédie est une violente charge contre ce dernier.
Elle se terminait par une danse licencieuse de la mère du démagogue.
Nombre de commentateurs, à l'instar des scholiastes, ont rapproché les \italique{Cavaliers} du \italique{Maricas}. Les \italique{Thesmophories} étant plus pauvres en scholies anciennes, personne n'a vu jusqu'ici qu'Aristophane reprenait certaines idées d'Eupolis. Nous l'allons voir tout-à-l'heure.
Mettons en garde tout-de-suite : les hypothèses -- comme les explications -- seront fort spéculatives.\\
Un commentaire du \italique{Maricas}, daté du \nieme{II} siècle ap. J.-C., nous est parvenu de façon très morcelaire ($=$ fr. 192 Kassel-Austin). On y lit cette citation (fr. 192.96-99 Kassel-Austin) :%
\citComm{τοῦτ᾽ ἐκδανείζει καὶ κυκᾷς [τὸ]\douteux{ν} ναυτικὸν \douteux{ἄ}[ργυρον \ \ \ \ \ (\textbf{ἄργυρον} Bibi)\\
ἐπὶ τῷ πέμπτῳ μέρε[ι \ \ \ \ \ \ \ ]τοὺς τόκ[ους οἱ\\
ναυτικοί, ἡμεῖς δ᾽ ἄρ᾽ οἴκ[ \ \ \ \ \ \ \ \ \ \ \ ]εν· ὁ τῶν\\
πλουσίων λέγει χορός.}%
(J'interprète sans véritable boussole et navigue à vue ; des erreurs se glissent certainement : lecteurs, soyez attentifs !)
Le chœur des riches (ὁ τῶν πλουσίων χορός) reproche à Maricas ($=$ Hyperbolos) de prêter (ἐκδανείζει) en concoctant (κυκᾷς) des prêts (τὸν ναυτικὸν ἄργυρον) au cinquième (ἐπὶ τῷ πέμπτῳ μέρει\marqueur{2}), c'est-à-dire avec un taux d'intérêt de $20 \%$.
Il est questions ensuite des τόκους, `enfants' ($=$ intérêts) de ces ναυτικοί (cf. vers 845).
Continuant le jeu de mot nautique (les prêts sont littéralement de l'argent \italique{marin}), est-ce à ce moment qu'Eupolis, par la voix du chœur, proposait de faire tondre \italique{nautiquement} cet usurier ou sa mère et de lui faire porter une \italique{barque} (σκάφιον vers 838) ?\\
Revenons aux \italique{Thesmophories} : en guise d'oraison funèbre, Aristophane rend hommage à la pertinence des propositions d'Eupolis :%
%Ce \italique{Maricas} remonte à la période où Lamachos était encore perçu par les comiques comme un bravache. Eupolis aura pu mettre dans le même panier la mère de Lamachos avec celle d'Hyperbolos. Rétrospectivement, au vu des derniers événements, voilà qui semble bien peu pertinent.
%S'il faut attribuer l'idée d'accorder les sièges des Athéniennes (et peut-être aussi celle de les tondre) avec la valeur de leur progéniture à feu Eupolis, 
%Une preuve supplémentaire -- et donc, nous espérons, définitive -- de la quasi-coïncidence de la mort d'Eupolis et de la représentation des \italique{Thesmophories} est l'adresse d'Aristop  :% dont l'oraison funèbre est réduite à la portion congrue
\citComm{τῷ γὰρ εἰκὸς\marqueur{3} ωὖπολι\\
Comment tolérer\marqueur{3}, ô Eupolis, que \&c.}%
%Aristote multiples exemples : τῷ γὰρ τὸ πρᾶγμα εἶναι ἢ μὴ εἶναι
%(L'expression `τῷ γὰρ εἰκὸς' se trouve telle quelle aussi au vers 703 des \italique{Acharniens}.)
Dernier coup de chapeau d'Aristophane à son défunt rival, mort des suites de la folie des Athéniens et de l'incompétence des amiraux.%
}%
\debutNote{839}%
\note{1}{Voir l'introduction et la question de la date. Avec Rogers, je penche pour l'hypothèse que les \italique{Thesmophories} furent représentées aux Dionysies de 410. Si nous avons vu juste avec Natalia, Eupolis serait donc mort l'été précédent la première représentation de la pièce.}%
\note{2}{Sens bien attesté d'ἐπί. Cf. Dém. 1250, 20 : ἐπ᾽ ὀκτὼ ὀβολοῖς δανείζειν τοῦ μηνὸς τὴν μνᾶν, «~prêter de l'argent moyennant un intérêt de huit oboles la mine par mois~», c'est-à-dire à $16 \%$.}%
\note{3}{Le tour elliptique apparaît aussi au vers 703 des \italique{Acharniens} où il est glosé `πῶς δίκαιόν ἐστιν' par une scholie. Le ton est effectivement indigné : «~comment peut-il être juste/raisonnable que...~», c'est-à-dire «~quelle indignité que...~». Voir aussi Sophocle (\italique{Ajax}, vers 1366-1367), non indigné mais ironique :%
\citComm{Ἀγαμέμνων -- ἡ πάνθ᾽ ὅμοια πᾶς ἀνὴρ αὑτῷ πονεῖ.\\
Ὀδυσσεύς -- τῷ γάρ με μᾶλλον εἰκὸς ἢ ᾽μαυτῷ πονεῖν;\\
Agamemnon -- En toute chose, chacun ne travaille que pour lui-même.\\
Ulysse -- Comment peut-il être plus raisonnable que de travailler pour moi-même ?}%
}}
\end{groupe}%
%% v. 830-845
\begin{groupe}%
\ledpb%
\tetramTroch{τὴν Ὑπερβόλου καθῆσθαι μητέρ᾽ ἠμφιεσμένην
\commentaire{840-1}{Ὑπερβόλου ... Λαμάχου,}{Hyperbolos, le démagogue athénien -- marchand de lampes de son état\marqueur{1} -- devint démagogue en chef après la mort de Cléon devant Amphipolis en 422 av. J.-C.\marqueur{2} ; sa mère et lui furent lors une cible privilégiée d'Aristophane et des poètes comiques\marqueur{3} (suite sans doute à quelque lapsus à relent freudien du démagogue, lapsus qui fut repris à l'encan par les Comiques). Cet ὀξίνην Ὑπερβόλον\marqueur{4} réunissait tout pour être détesté d'Aristophane : de caractère exécrable, il était un partisan, comme son mentor Cléon, de la guerre à outrance\marqueur{5} ; c'était un intenteur de procès\marqueur{6}, un fourbe\marqueur{7}, qui s'est enrichi par ses coquineries\marqueur{8}. Il eut l'insigne honneur d'être le dernier des Athéniens ostracisés en 417\marqueur{9}. Il fut assassiné à Samos, son lieu d'exil\marqueur{10}, et son corps jeté en mer\marqueur{11}... cette \italique{inhudation} lui valut une dernière raillerie de la part d'Aristophane\marqueur{12}. 
%\eff{Qu'est-ce qui valut à la mère d'Hyperbolos d'être raillée ainsi par les poètes comiques ? Il me semble qu'on peut inférer un début de réponse du fameux passage des \italique{Nuées} (vers 551-556) où Aristophane fait preuve d'une belle assurance de posséder sa Muse, pourtant encore bien neuve : %
%\citComm{οὗτοι δ’, ὡς ἅπαξ παρέδωκεν λαβὴν Ὑπέρβολος,\\
%τοῦτον δείλαιον κολετρῶσ’ ἀεὶ καὶ τὴν μητέρα.\\
%Εὔπολις μὲν τὸν Μαρικᾶν πρώτιστον παρείλκυσεν\\
%ἐκστρέψας τοὺς ἡμετέρους Ἱππέας κακὸς κακῶς,\\
%προσθεὶς αὐτῷ γραῦν μεθύσην τοῦ κόρδακος οὕνεχ’, ἣν\\
%Φρύνιχος πάλαι πεπόηχ’, ἣν τὸ κῆτος ἤσθιεν.\\
%εἶθ᾽ Ἕρμιππος αὖθις ἐποίησεν εἰς Ὑπέρβολον,\\
%ἄλλοι τ᾽ ἤδη πάντες ἐρείδουσιν εἰς Ὑπέρβολον,\\
%τὰς εἰκοὺς τῶν ἐγχέλεων τὰς ἐμὰς μιμούμενοι.\\
%ὅστις οὖν τούτοισι γελᾷ, τοῖς ἐμοῖς μὴ χαιρέτω:\\
%ἢν δ᾽ ἐμοὶ καὶ τοῖσιν ἐμοῖς εὐφραίνησθ᾽ εὑρήμασιν,\\
%εἰς τὰς ὥρας τὰς ἑτέρας εὖ φρονεῖν δοκήσετε.\\
%Mes rivaux, depuis qu'une fois Hyperbolos leur a donné prise, ne cessent de piétiner le malheureux et sa mère. Eupolis, tout le premier, traîna sur la scène son \italique{Maricas}, où, piètre qu'il est, il retourna piètrement nos \italique{Cavaliers}, en y ajoutant en vue de la cordace une vieille drôlesse, celle-là même qu'imagina jadis Phrvnichos et que le monstre marin voulait manger. Puis, à son tour, Hermippos prit à partie Hyperbolos, et tous les autres de se buter contre Hyperbolos, en m'empruntant ma comparaison des anguilles. Que celui donc qui rit à leurs pièces ne se plaise pas aux miennes ! Mais si vous vous laissez égayer par moi et mes inventions, dans les temps à venir vous passerez pour sensés}%
%La clef est dans "παρέδωκεν λαβὴν Ὑπέρβολος"}\\
\textbf{Lamachos} était un stratège athénien\marqueur{13} et l'un des représentants athéniens de la Paix de Nicias\marqueur{14}. Avec Nicias et Alcibiabe, il est investi des pleins pouvoirs lors de l'expédition de Sicile\marqueur{15} où il trouve la mort\marqueur{16}. Lamachos présente une originalité unique parmi les personnages historiques d'Aristophane puisque le poète opère un volte-face à son endroit : du va-t-en-guerre au panache ridicule dépeint dans les \italique{Acharniens} et la \italique{Paix}\marqueur{17} (accusé même de s'enrichir grâce à la guerre\marqueur{18}), Aristophane lui rend hommage ici comme dans les \italique{Grenouilles}\marqueur{19}. La mort héroïque de Lamachos\marqueur{20} est sans doute pour beaucoup dans ce changement radical d'opinion.}
\debutNote{840-1}
\note{1}{\italique{Cavaliers}, v. 739 (non nommé) ; \italique{Nuées}, v. 1064 ; \italique{Paix}, v. 690.}
\note{2}{Voir la \italique{Paix}, v. 679-692, Paix qui fout le camp, révoltée que le peuple se soit donné un si indigne patron.}
\note{3}{Cf. les \italique{Nuées}, v. 551-562 : Aristophane daube sur ses plagiaires de contemporains, notamment Eupolis réchauffant leurs \italique{Cavaliers} pour son \italique{Maricas} et y rajoutant un autre plat déjà concocté par Phrynicos en son temps : une mère d'Hyperbolos, ivrognesse et dansant le licencieux cordax. Sur les attaques des comiques contre la mère d'Hyperbolos, présentée comme une débauchée notoire, voir aussi les fragments 7-12 des \italique{Vendeuses de pain} (αἱ Ἀρτοπώλιδες) d'Hermippos (dans l'édition Kassel-Austin) ainsi que les fragments de Platon le Comique dans sa pièce éponyme \italique{Hyperbolos}.}
\note{4}{\italique{Cavaliers}, v. 1303. Thucydide, qui n'estimait guère plus ce μοχθὴρον ἄνθρωπον, ne parle qu'une seule fois d'Hyperbolos : pour nous annoncer sa mort et en brosser un portrait expéditif (VIII, 73.3) : %
\citComm{καὶ Ὑπέρβολόν τέ τινα τῶν Ἀθηναίων, μοχθηρὸν ἄνθρωπον, ὠστρακισμένον οὐ διὰ δυνάμεως καὶ ἀξιώματος φόβον, ἀλλὰ διὰ πονηρίαν καὶ αἰσχύνην τῆς πόλεως, ἀποκτείνουσι μετὰ Χαρμίνου\\
S'en prenant à Hyperbolos, un Athénien méprisable qu'on avait ostracisé non par peur de son influence et de son prestige, mais parce que c'était un malhonnête homme qui déshonorait la cité, ils [les conjurés favorables à l'oligarchie] l'assassinèrent avec l'aide de Charminos.}%
}
\note{5}{\italique{Cavaliers}, v. 1303-1315 ; il cherchait à renforcer la flotte pour que l'expédition de Sicile fût poussée jusqu'à Carthage.}
\note{6}{\italique{Acharniens}, v. 846-847.}
\note{7}{\italique{Guêpes}, v. 1007.}
\note{8}{\italique{Nuées}, v. 1064-1066. Tel fils, telle mère ? L'usure brocardée par Aristophane dont fait preuve la mère d'Hyperbolos n'était-elle pas la continuation des activités louches du fils, depuis que ce dernier se trouvait en exil à Samos ?}
\note{9}{Plut., \italique{Nicias} 11 ; \italique{Alcib.} 13.}
\note{10}{%\italique{Après} la représentation des \italique{Thesmophories}. Willems avait déjà remarqué que l'attitude de la mère d'Hyperbolos, prenant part en costume d'apparat aux fêtes célébrées par les femmes, était incompatible avec le deuil d'une mère éplorée. Ce détail confirme, ce nous semble, la date de 411 pour la première représentation des \italique{Thesmophories} et infirme celle de 410 soutenue notamment par Rogers.
Thucydide, VIII, 73.3.}
\note{11}{Théopompe 115F69, in \italique{Die Fragmente der griechischen Historiker}, Teil 2, \italique{Zeitgeschichte - B}.}
\note{12}{\italique{Grenouilles}, v. 570.}
\note{13}{Thucydide IV, 75. 1 : il commande une dizaine de navires dans le Pont.}
\note{14}{Thucydide V 19.2 et 24.1}
\note{15}{Thucydide (VI, 8) ; il s'opposa à Alcibiade sur la stratégie (V 49) : Lamachos voulait aller droit au but et attaquer Syracuse sitôt le débarquement fait ; il finit par se ranger à l'opinion d'Alcibiade (VI, 50).}
\note{16}{Thucydide (VI, 101) : Lamachos courant à la rescousse de l'aile droite de l'armée athénienne se retrouva isolé avec quelques-uns de ses camarades ; il fut tué par les Syracusains qui récupérèrent son corps avant de le restituer (VI, 103) aux Athéniens. Nicias fut dès lors le seul commandant en Sicile.}
\note{17}{Lamachos arborait sur son casque trois aigrettes qui lui donnaient des faux airs de Gorgone (\italique{Acharniens}, v. 574 ; v. 964-965). Aristophane rapproche cet accoutrement d'épouvante à d'autre monstre-croquemitaine au nom voisin (Μορμώ, \italique{Acharniens} v. 582 ; \italique{Paix} v. 474). Enfin, cette crinière en bataille -- de circonstance pour un militaire ! -- permettait un jeu analogue en grec entre λόφοι et λόχοι, trop tentant pour qu'Aristophane n'y cèdât pas (\italique{Acharniens}, v. 567, 574, 1074).}
\note{18}{Il est traité de μισθαρχίδης au v. 597 des \italique{Acharniens} ; attaque renouvelée au vers 619.}
\note{19}{\italique{Grenouilles}, v. 1039.}
\note{20}{Voir -- outre le passage de Thucydide (VI, 101) -- Plutarque, \italique{Nicias} 18 (et aussi le \italique{Lachès} (19c) de Platon où il est question du courage de Lamachos.)}\\
λευκὰ καὶ κόμας καθεῖσαν πλησίον τῆς Λαμάχου·\\
καὶ δανείζειν χρήμαθ᾽, ἥν χρῆν, εἰ δανείσειέν τινι\apparat{842}{χρήματα τ' R : χρήματ' \Mon : \lecon{χρήμαθ᾽} Scaliger \espApp \lecon{ἥν χρῆν} R : ᾗ χρῆν Scaliger (\italique{cui} Divus) \espApp χρῆν ἦν R : \lecon{χρῆν εἰ} Brunck \espApp δανείσειε R : \lecon{δανείσειέν} Biset \ligneApp{844} ἀφερεῖσθαι R : \lecon{ἀφαιρεῖσθαι} \Mon \espApp χρήματατ᾽ R : \lecon{χρήματ᾽} Grynaeus \espApp τ\auDessus{ὰ}{ο}δὶ \Rpc : \lecon{τοδί} \Mon}%
%\commentaire{842}{ᾗ χρῆν}{\eff{Imprécation ? (Austin)}}
\\
καὶ τόκον πράττοιτο, διδόναι μηδέν᾽ ἀνθρώπων τόκον,\\
ἀλλ᾽ ἀφαιρεῖσθαι βίᾳ τὰ χρήματ᾽ εἰπόντας τοδί,\\
`ἀξία γοῦν εἶ τόκου τεκοῦσα τοιοῦτον τόκον.'
\commentaire{845}{ἀξία γοῦν εἶ τόκου τεκοῦσα τοιοῦτον τόκον.}{Jeu de mots qui roule sur le double sens de τόκος, \latin{usura} et \latin{proles}. C'est ainsi que l'orateur Eschine disait de Démosthène qui, souffleté par Midias, s'était désisté moyennant 30 mines : %
\citComm{ὁ γὰρ ἄνθρωπος οὐ κεφαλήν, ἀλλὰ πρόσοδον κέκτηται.\\
Ce n'est pas une tête (κεφαλήν), mais un capital que possède cet homme.}%
(Cf. \italique{Contre Ctésiphon}, 212 ; note de Willems.)}\\\passeLigne}%
\end{groupe}%
%% v. 846-851
\begin{groupe}%
\ledpb%
\personnageGrec{Μνησίλοχος}%
\typeVers{trim. iamb.}\trimIamb{ἰλλὸς γεγένημαι προσδοκῶν· ὁ δ᾽ οὐδέπω.\testimonia{846}{Suidas ι 323 $\sim$ Zonaras 1103 ἰλλός \espTest (ἰλλός) Pollux II 52 ; \Eustathe 768.9 ; \AelDion ι 6 ; cf. Moeris ι 21 (ἰλλόν $=$ στραβόν) ; Hesychius ι 578 ἰλλός ; Hesychius η 257 (ἠίλλοντο $=$ ἐστράβιζον)}\\
τί δῆτ᾽ ἂν εἴη τοὐμποδών; οὐκ ἔσθ᾽ ὅπως\\
οὐ τὸν Παλαμήδη ψυχρὸν ὄντ᾽ αἰσχύνεται.\apparat{848}{Παλαμήδην R : \lecon{Παλαμήδη} Elmsley (\italique{Quart. Rev.} 7 (1812) p.454)}\\
τῷ δῆτ᾽ ἂν αὐτὸν προσαγαγοίμην δράματι;\\
ἐγᾦδα· τὴν καινὴν Ἑλένην μιμήσομαι.
\commentaire{850}{τὴν καινὴν Ἑλένην.}{« La nouvelle \italique{Hélène} ». \italique{Nouvelle} parce qu'Euripide l'avait présentée il y a peu de temps, sous l'archontat de Cléocrite (printemps 412). Mais l'adjectif καινή est aussi à prendre au sens d'\italique{innovante}, \italique{imprévue}, \italique{étrange}. Voici ce que dit Victor-Henry Debidour de cette « \italique{Hélène} nouveau style »\marqueur{1} :\\
Dans la tradition légendaire, Hélène de Sparte, femme de Ménélas, était le parangon de la beauté. Mais aussi le plus scandaleux exemple d'infidélité féminine, cause, par son inconduite, des malheurs sans nombre de la guerre de Troie. Le personnage s'était stéréotypé en « femme fatale », d'autant plus haïssable qu'elle était plus prestigieuse ; et \italique{Euripide lui-même ne la traite pas autrement} -- sauf justement dans la tragédie très peu tragique qui porte son nom. Certes, elle y reste, bien-sûr, la belle Hélène, et si charmeuse que le fils du sage Protée l'aime comme un fou, et comme un sot. Mais la voici transformée en épouse sans tache, victime innocente de la longue perfidie des deux déesses, dont l'une l'a prostituée et l'autre séquestrée : elles l'ont engagée dans une rocambolesque histoire où elle mérite toute notre sympathie. Pour se faire désigner par Pâris comme « la plus belle », Cypris avait promis au prince-berger troyen de lui donner Hélène ; Héra, furieuse, contre sa rivale et son juge, fit en sorte que Pâris n'emmenât à Troie qu'un fantôme d'Hélène, une « nuée » toute semblable à la vraie et suffisamment consistante pour leurrer le ravisseur pendant les dix ans qu'il la garda, et le légitime époux après ce qu'il crut être sa reconquête. La vraie Hélène, elle, miraculeusement transportée en Égypte, y était tenue cachée durant toutes ces années, confiée à la vertueuse protection de Protée. Ainsi, c'est pour un « crime sans criminelle » que des milliers de preux se sont entre-tués et que Troie a été anéantie.\\
Telle est l'\italique{étonnante} variante que recueille Euripide sur la plus célèbre entreprise guerrière de la légende grecque. Car il ne l'a pas inventée : Hérodote en fait état bien avant lui\marqueur{2}. De ce point de départ, Euripide a tiré un canevas sur lequel il a brodé une aventure romanesque et humoristique qui souvent ferait croire qu'avant d'être parodiée elle-même par Aristophane, l'\italique{Hélène} donnait dans la parodie.}
\debutNote{850}
\note{1}{Dans l'introduction à sa traduction de l'\italique{Hélène}. Oui ! car Debidour le grand traducteur d'Aristophane, qui a su rendre un peu de la verve de notre poète dans la langue de Rabelais, Debidour ne s'est pas limité à la Comédie Ancienne ; il est aussi l'auteur d'une traduction des pièces conservées des trois grands Tragiques : \italique{Les Tragiques -- Eschyle, Sophocle, Euripide, théâtre complet}, La Pochothèque (1999). Traduction à laquelle s'ajoute l'excellent dossier et les index fort utiles de Paul Demont et Anne Lebeau. Victor-Henry Debidour vérifie ainsi le paradoxe de Socrate dont Platon n'a pas cru utile de nous donner la démonstration à la fin de son \italique{Banquet}.}
\note{2}{Hérodote II, 112-120. Par la suite, du reste, le thème de la réhabilitation d'Hélène fut maintes fois repris en exercice d'école : la paradoxe devient à son tour lieu commun.}\\
πάντως δ᾽ ὑπάρχει μοι γυναικεία στολή.\apparat{851}{πάντως R : πάντως δ᾽ Bentley, cf. v.1012}\\\passeLigne}%
\end{groupe}%
%% v. 852-854
\begin{groupe}%
\personnageGrec{Κρίτυλλα}%
\trimIamb{τί αὖ σὺ κυρκανᾷς; τί κοικύλλεις ἔχων;\testimonia{852}{Suidas κ 2534 \espTest (κοικύλλεις) Hesychius κ 3232 κοικύλλειν $\sim$ Photius κ 860}\apparat{852}{ἢ τί κοικύλλεις R, Suid. κ 2534 : \lecon{τί κοικύλλεις} Faber \ligneApp{853} Ἑλένη R : Ἑλένην Bonini : \lecon{ἑλένην} Bibi, ap. Hesych. ε 1995}%
\commentaire{852}{τί κοικύλλεις ἔχων;}{Κοικύλλειν est un verbe à redoublement expressif\marqueur{1} qui tire son origine de τὰ κύλα, \italique{les poches sous les yeux}, d'où le sens d'« écarquiller les yeux »\marqueur{2}. Voici pour le sens en lequel l'emploie Critylla.
Le verbe est particulièrement bien adapté à la situation : notre anti-héros vient de lancer `à la mer' ses `rames' dans l'espoir de rameuter Euripide ; il est là maintenant, le regard hagard, devant cette mer imaginaire, à attendre... pareil à ce personnage comique, surnommé Κοικυλίων, type de l'idiot qui comptait les vagues de la mer\marqueur{3}, le comble de l'inutilité pour un grec : κύματα μετρεῖν était même passé en proverbe (détourné par Aristophane au vers 1257 de la \italique{Paix}).}%
\debutNote{852}%
\note{1}{Sur le même modèle que μοιμύλλειν, δενδίλλειν, δανδαίνειν, παπταίνειν, \&c.}%
\note{2}{Cf. Hésychius κ 3232 : κοικύλλειν· περιβλέπειν· κύλα γὰρ τὰ ἐπάνω τῶν ὀφθαλμῶν [$=$ la partie supérieure des paupières, mais voir aussi la glose d'Hésychius sous κύλα, où le terme désigne la partie \italique{inférieure} des paupières]. οἱ δὲ τὸ περιστρέφεσθαι καὶ διατρίβειν περιβλεπόμενον. ἄλλοι φθονεῖν.}%
\note{3}{Ælien, \italique{Varia Historia}, XIII, 15 ($=$ \latin{Adespota comica}, fr. 72 K.-A.) : 
\citComm{φασὶ παχύτατον γενέσθαι τὴν διάνοιαν οἱ τῆς κωμῳδίας ποιηταὶ Πολύδωρόν τινα, καὶ ἄλλον Κοικυλίωνα ὄνομα, ὅσπερ τὰ κύματα ἠρίθμει ὑπὸ τῆς ἄγαν ἀνοίας.\\
\francais{Les poètes comiques disent qu'un certain Polydore (aussi [connu sous] un autre nom : Cœcylion) affolit tant qu'il comptait les vagues sous le coup de son trop plein de folie.}}%
Ce passage, déniché par Fritzsche (à l'érudition et la bibliothèque toujours impressionnantes !), est à rapprocher du chapitre du Περὶ βλασφημιῶν de Suétone (édité, pour ce qu'il a pu en reconstituer, en 1967 aux Belles Lettres par le grand Taillardat) consacré à quelques idiots célèbres. Περὶ βλασφημιῶν, β 189 $=$ Eustathe \italique{Comm. Od.} p. 1669.53 : ὥσπερ δὲ τουτούς ἡ ἱστορία λόγου ἠξίωσεν ἐν οἷς καὶ τὸν τὰ κύματα μετροῦντα Πολύ[δ]ωρον.}\\
πικρὰν ἑλένην ὄψει τάχ᾽, εἰ μὴ κοσμίως
\commentaire{853}{πικρὰν ἑλένη.}{Le détournement du \italique{Palamède} ayant échoué (après celui du \italique{Télèphe}), Mnésiloque songe maintenant à recourir à la récente \italique{Hélène} d'Euripide. Seulement, de toutes les \italique{Thesmophories}, aucun des geôliers de Mnésiloque ne rentre jamais dans le jeu de ses parodies d'Euripide. Il est étonnant que pas un des nombreux commentateurs d'Aristophane ne s'en soit avisé~: quand Kritylla prononce `πικρὰν ἑλένην ὄψει τάχ᾽', elle ne songe pas un seul instant à la mythique Hélène, mise nouvellement en scène par Euripide avec réécriture du scénario. C'est Hésychius encore une fois – décidément le plus utile des lexicographes pour les \italique{Thesmophories}~! – qui nous le confirme (ε 1995)~:%
\citComm{ἑλένη· λαμπάς, δετή\\
\francais{hélène~: flambeau, torche.}}%
Kritylla menace Mnésiloque d'un bon coup de flambeau, flambeau qu'elle avait lors de la procession et l'arrivée au temple. (Coup de flambeau qu'elle assénera au vers 858 en évoquant Hécate \textbf{φωσ}φόρος.)}\\
ἕξεις, ἕως ἂν τῶν πρυτάνεών τις φανῇ.\\\passeLigne}%
\end{groupe}%
%% v. 855-857
\begin{groupe}%
\ledpb%
\personnageGrec{Μνησίλοχος}%
\trimIamb{Νείλου μὲν αἵδε καλλιπάρθενοι ῥοαί,
\commentaire{855-7}{Début du prologue d'\italique{Hélène}.}{Ainsi commence la véritable pièce d'Euripide :%
\citComm{Ἑλένη --- Νείλου μὲν αἵδε καλλιπάρθενοι ῥοαί,\\
ὃς ἀντὶ δίας ψακάδος Αἰγύπτου πέδον\\
λευκῆς \difference{τακείσης χιόνος ὑγραίνει γύας.}\\
\francais{Du Nil voici les rives aux belles Nymphes,\\
Nil qui à la place de la pluie dispensée par Zeus\\
arrose la plaine d'Égypte une fois la blanche neige fondue et irrigue les terres.}}%
Comme tous les semi-habiles -- euphémisme pascalien : comprenez `semi-débiles' --, Euripide aime à bavasser ; il ne peut s'empêcher dès le troisième vers de la pièce d'exposer son opinion sur la raison des crues du Nil -- Ah ça ! se croit-il obligé d'en avoir une quel que soit le sujet ? et croit-il par-dessus le marché qu'elle nous intéresse ? Selon donc le très savant Euripide, ces crues sont dues à la fonte des neiges. 
Théorie fantaisiste et traitée avec dédain par Hérodote (II, 19-25), condescendant vis-à-vis de cette thèse soutenue par Anaxagore\marqueur{1}, le maître d'Euripide -- raison pour laquelle ce dernier se croit tenu de la défendre ? En tout cas il récidive quelques années après l'\italique{Hélène} et les \italique{Thesmophories}, l'an 408/407, dans \italique{Archélaos}\marqueur{2}:%
\citComm{Δαναὸς ὁ πεντήκοντα θυγατέρων πατὴρ\\
Νείλου λιπὼν καλλήϊον Γαιᾶ ὕδωρ, \ \ \ \ κάλλιστον ἐκ γαίας (vel γαίης) codd. : \lecon{καλλήϊον Γαιᾶ} Bibi\\
ὃς ἐκ μελαμβρότοιο πληροῦται ῥοὰς\\
Αἰθιοπίδος γῆς, ἡνίκ' ἂν τακῇ χιὼν\\
τέθριπφ' ἱέντος ἡλίου κατ' αἰθέρα \ \ \ \ \ \ τέθριππεύοντος : \lecon{τέθριπφ' ἱέντος} Papathomopoulos\\
Danaos, le père de cinquante filles $[ =$ les Danaïdes ; un oracle avait prédit que Danaos serait tué par un de ses gendres, d'où son exil d'Égypte pour ne pas que ses filles épousent leurs cousins et leur arrivée en \italique{Suppliantes} à Argos$]$,
du Nil abandonnant l'eau fertile du Géon,
Nil qui emplit ses berges depuis la terre
au visage sombre d'Éthiopie chaque fois que vient à fondre la neige
lorsque le Soleil élance ses quadriges au fil de l'éther.}%
(Il y a quelques problèmes textuels : au second vers, les manuscrits donnent †κάλλιστον ἐκ γαίας† (ou ἐκ γαίης) ; Euripide use encore d'un néologisme formé sur κάλλι- comme avec καλλιπάρθενος ou καλλιδόναξ. Il faut rétablir καλλήϊον, \italique{à la belle récolte}, \italique{fructueuse}, d'après la glose d'Hésychius (κ 468) elle-même corrompue :%
\citComm{†κάλλης· καρπίμου}%
(Dans un cas comme dans l'autre, l'origine de la corruption vient d'un iota I mélu en sigma lunaire C.)\\
La fin du vers se corrige sans difficulté à l'aide d'une autre glose d'Hésychius (γ 66) :%
\citComm{Γαιών· ποταμός, παρ᾽ Ἕλλησι Νεῖλος}%
Au cinquième vers, le †τέθριππεύοντος† fautif doit se lire τέθριπφ' ἱέντος conformément à l'image grecque de \italique{lancer} une cavalcade\marqueur{3} et l'usage d'Euripide d'employer τέθριππον au pluriel.\\
On constate au passage qu'encore une fois il n'y avait nul besoin de tournebouler le texte transmis pour le remettre d'aplomb.)\\
Remarquons qu'en ces deux passages, les ῥοαί ne doivent pas se traduire par \italique{flots} mais bien par \italique{rives} ou \italique{berges}\marqueur{4}. (L'expression «~καλλιπάρθενοι ῥοαί~» du prologue d'\italique{Hélène} suggère peut-être un jeu de mots avec les Ῥοαί, les \italique{nymphes des grenadiers}\marqueur{5}.)\\
Attaquons-nous enfin au détournement du prologue par Aristophane !
Pour traîner dans la boue -- ou pis -- la muse d'Euripide, notre poète ne recule devant rien. 
La syrméa dont il est question au vers 857 était une espèce de rave très abondante en Égypte. 
On en tirait un suc portant le même nom utilisé comme purgatif\marqueur{6}. Suivant Hérodote\marqueur{7}, les Égyptiens se purgeaient menstruellement, trois jours de suite. 
Ce curieux hygiénisme est assez plaisant pour qu'Aristophane n'ait pas manqué de le railler.\\
Au vers 570, nous avons vu que les gâteaux de miel grecs avaient, au contraire, des vertus astringentes, ce qui expliquait qu'Aristophane ait certainement mis en scène une confrontation entre le vendeur de syrméa et le vendeur de pâtisseries dans ses \italique{Danaïdes}, comédie aujourd'hui perdue.\\
De même, dans la \italique{Paix}, Trygée suggère au marchand d'armes menacé par la ruine de couper ses lances pour les proposer à la vente comme échalas. Autre idée : recycler ses casques d'une bien étrange manière (vers 1253-1254) :%
\citComm{Πώλει βαδίζων αὐτὰ τοῖς Αἰγυπτίοις·\\
ἔστιν γὰρ ἐπιτήδεια συρμαίαν μετρεῖν.\\
Va les vendre en Égypte.\\
Ils sont bons à mesurer la syrméa.}%
(C'est-à-dire à mesurer \italique{les effets} de la syrméa ; Trygée ne voit donc en ces casques que de vulgaires pots de chambre. Je dirais même plus : ils ne sont bons à rien, car Trygée fait entendre le proverbe κύματα μετρεῖν, \italique{compter les vagues}, comble de l'inutilité.)\\
Même genre de détournement dans le présent passage des \italique{Thesmophories} : les Égyptiens se servent du Nil pour se \italique{nettoyer} (ici νοτίζειν $=$ βαλανεύειν, σπογγίζειν) des désagréments de leur régime. 
On peut conclure de ce passage qu'en grec le mot πέδον, le \italique{sol} ; littéralement, ce sur quoi repose le pied, a fourni la même image qu'en français le \italique{fondement} et sert à désigner ἡ πυγή\marqueur{8}.\\
%\eff{
%Ce sens permet de comprendre le passage suivant des \italique{Oiseaux} (vers 504-507) où il est encore question d'Égypte : CITER COIRAULT%
%\citComm{Πισθέταιρος --- Αἰγύπτου δ’ αὖ καὶ Φοινίκης πάσης κόκκυξ βασιλεὺς ἦν·\\
%χὠπόθ’ ὁ κόκκυξ εἴποι «κόκκυ,» τότ’ <ἂν> οἱ Φοίνικες ἅπαντες\\
%τοὺς πυροὺς ἂν καὶ τὰς κριθὰς ἐν τοῖς πεδίοις ἐθέριζον.\\
%Εὐελπίδης --- Τοῦτ’ ἄρ’ ἐκεῖν’ ἦν τοὔπος ἀληθῶς· «Κόκκυ, ψωλοί, πεδίονδε.»\\
%Pisthétaire --- En Égypte, d'autre part, et dans toute la Phénicie, le coucou était roi ; et quand le coucou criait : coucou ! tous les Phéniciens allaient moissonner le froment et l'orge dans les champs\\
%Évelpide --- Le voilà donc le dicton dans toute sa vérité : coucou, déprépucés, aux champs}%
%Voici, selon Willems, comment il faut entendre ce passage. Le coucou apparaissait en Grèce au printemps. Pline nous apprend que les laboureurs qui n'avaient pas fini de tailler leurs vignes à l’équinoxe de cette saison étaient en butte à des propos salés, \latin{petulantiæ sales}, qu’on leur jetait en imitant le chant du coucou\marqueur{\italique{Hist. nat.}, XVII, 66.}. 
%Le dicton consacré par l'usage, nous l’avons ici en original : κόκκυ, ψωλοὶ, πεδίονδε. Il avait cours, ce qui n'a rien d'étonnant, en Italie aussi bien qu'en Grèce, et Horace y fait également allusion\marqueur{\italique{Satir}., I, 7, 29}. I'o)ó; était le terme cru pour désigner un débauché ; les Latins disaient verpus, et nous renverrons pour le sens du mot, tant au propre qu'au figuré, à la définition de Forcellini citée ci-dessus. En français on traduirait assez exactement : « Coucou, déprépucés, aux champs! (*) » On conçoit que pareil lardon, jeté en passant auvigneron négligent, soit appelé par Pline une exprobratio fœda. C'est, n'en doutons pas, pour amener cette équivoque cynique, familière à la généralité de ses auditeurs, qu'Aristophane a fait du coucou le roi des Égyptiens et des Phéniciens, lesquels, comme on le sait par Hérodote, étaient des circoncis (!o)oi, au sens de repitogo.). Il ne sera pas inutile d'observer, ce que certains paraissent avoir oublié, qu'Évelpide est un pur campagnard, habitant à Halimus (vv. 496 et 585), et que les exemples qu'il cite sont empruntés à la vie rustique. Nous savons, depuis Mannhardt, combien les mêmes croyances et pratiques se retrouvent chez les populations rurales de tout legroupe indo-européen. C'est ainsi que dans certaines parties de l'Allemagne les paysans ont encore coutume de se rouler par terre à la première apparition, non du milan, comme dans Aristophane (v. 502), mais du coucou (º). La présente note complète ce que dit Mannhardt à propos de l'usage mentionné par Pline (º). Entre autres rapprochements, le savant mytho graphe cite l'épithète de Horbuck lancée aux passants par les ruraux du Schleswig. Cemot depatois (en allemand Hurenbock) est bien l'équivalent du po)ó; des Grecs.
%Mais avec le sophiste des sophistes Euripide, les mots perdent tout sens ; la cause du discrédit total de la parole publique et de tous les malheurs qui s'ensuivent pour Athènes, c'est dans ses tragédies qu'il fut aller les chercher. (Aristophane rajouterait aussi dans le même panier de crabes tous les socratiques ergoteurs.) Quelle sera l'épithète assez vigoureuse pour qualifier le scribouilleur qui, pour décrire les flots (ῥοαί) du Nil, cet afflux fangeux charriant des tonnes de limon noir, recourt à l'adjectif ... καλλιπάρθενοι, \italique{à la beauté virginale} ? Il a osé ! c'est trop fort ! je vais te le réécrire, moi, ton prologue, mon salaud. Ça va être du propre ! parce que j'en ai une autre d'explication sous la main. Ta masse noirâtre, des \italique{ondes virginales} ? il nous prend vraiment pour des perdrix ! je vais te dire ce qu'est c'est : des  ῥοαί ? que non pas ; des διάρροιαι, oui !\\
%Le lecteur sent venir l'explication alternative d'Aristophane (qui n'a jamais eu peur d'appeler un chat un chat et une pièce d'Euripide de la merde). (L'image du flux d'un fleuve comparable à une colique se trouve dans un fragment du \italique{Gérytadès}\marqueur{Athénée, XII, 551a $=$ fr. 156(149.150) Kassel-Austin.} : τούτους γὰρ, ἢν πολλῷ ξυνέλθῃ, ξυλλαβὼν ὁ τῆς διαρροίας ποταμὸς οἰχήσεται, "le fleuve de fange, s'il est beaucoup gonflé, va les emporter au passage [les poètes souffreteux]".)\\
%Le verbe νοτίζειν, \italique{mouiller}, est attesté avec le sens de οὐρειν\marqueur{Voir le \LSJ }. Ici, νοτίζειν $=$ χέζειν. 
%Aristophane reprend-il une hypothèse plus populaire que celle d'Hérodote pour expliquer la crue du Nil ? ou use-t-il de mots crus pour une invention de son cru tout-à-fait incongrue ? En tout cas, le texte est clair : comme dans la chanson, les Égyptiens, ce \italique{peuple à la noire syrméa}, μελανοσυρμαῖον λεών, auront tant consommé de syrméa « que trois moulins en ont tourné / Petits ruisseaux, grandes rivières / Pendant trois jours ont débordé. »}%
}%
\debutNote{855-7}%
\note{1}{59A91 D.-K. $=$ fr. 300 Eschyle (ed. Radt). Eschyle avait en son temps aussi évoqué cette fonte des neiges dans les \italique{Suppliantes}, v. 559 sq.% \eff{ainsi qu'au fragment}
}%
\note{2}{fr. 1 Ἀρχέλαος Jouan-van Looy $=$ fr. 228 Kannicht. Commentaire de Tibérius le rhéteur : οἱ γὰρ τρεῖς περὶ τοῦ ποταμοῦ στίχοι περιττοί.}%
\note{3}{Dans le jeu enfantin de la χελιχελώνη, la «~tortue~» bondissait sur la troupe aux cris de «~λευκᾶν ἀφ' ἵππων εἰς θάλασσαν ἅλατο~», «~du haut de ses blanches cavales il a bondi dans la mer~». Voir l'édition des Περὶ βλασφημιῶν et Περὶ παιδιῶν par Taillardat, χελιχελώνη $=$ π 19.}%
\note{4}{Hésychius (ρ 415) :%
\citComm{ῥοία· κυλίστρα τῶν ἵππων παρὰ τῷ ποταμῷ καὶ ψάμμῳ.}%
Rien que pour l'\italique{Hélène}, les ῥοαί sont à prendre au sens de \italique{rives} au moins aux vers 53, 124 (avec jeu de mots sur Εὐρώτας), 250, 492.}%
\note{5}{Eustathe, \italique{Commentaire à l'}Odyssée, p. 1572, 35 : ῥοιαὶ ὁμωνύμως τῷ δένδρῳ καὶ αἱ περὶ αὐτὰς μυθικαὶ νύμφαι.}%
\note{6}{\italique{Souda} (σ 1667) : 
\citComm{Συρμαία: τὸ καθαρτικὸν δι' ἐμέτων καὶ κοιλίας. καὶ Συρμαί̈ζειν τὸ κενοῦσθαι, ἐκ ζειῶν καὶ ὕδατος. οὕτως Ἡρόδοτος. λάχανον σελίνῳ ἐοικός: ἢ βρῶμα διὰ στέατος καὶ μέλιτος.}%
Et voici ce qu'en dit Hésychius (σ 2782) :
\citComm{σύρμαια· [...] ἔστι δὲ βρωμάτιον διὰ στέατος καὶ μέλιτος. λέγεται δὲ καὶ συρμαϊσμός, καὶ ἔστι πρὸς κάθαρσιν. καὶ πόμα δι᾽ ὕδατος καὶ ἁλῶν.}}%
\note{7}{Hérodote II, 77 ; et aussi II, 88 ; II, 125.}%
\note{8}{Noter que si πέδον, la \italique{plaine}, renvoie au verso ; πεδίον, la \italique{petite plaine}, elle, fait allusion au recto. Cf. \italique{Lysistrata}, v. 88 ; \italique{Oiseaux}, v. 507 avec le commentaire de Willems. Voir aussi l'emploi de λειμών au vers 171 du \italique{Cyclope} d'Euripide.}\\
ὃς ἀντὶ δίας ψακάδος Αἰγύπτου πέδον\\
λευκῆς νοτίζει μελανοσυρμαῖον λεών.\\\passeLigne}%
\end{groupe}%
%% v. 858
\begin{groupe}%
\personnageGrec{Κρίτυλλα}%
\trimIamb{πανοῦργος εἶ νὴ τὴν Ἑκάτην τὴν φωσφόρον.%
\commentaire{858}{τὴν Ἑκάτην τὴν φωσφόρον.}{Ce n'est pas par hasard que Critylla adjoint à Hécate l'épithète de φωσφόρος, \italique{porte-torche} : elle est elle-même équipée d'une telle torche et il est à parier qu'en jurant de la sorte, Critylla met sa menace du vers 853 à exécution en assénant à Mnésiloque le coup de flambeau promis\marqueur{1}.}%
\debutNote{858}%
\note{1}{Cf. \italique{Lysistrata} vers 453 ; fr. 592.5 Kassel-Austin [$=$ Athénée XV, 701b] : ἐκφέρετε πεύκας κατ᾽ Ἀγάθωνα φωσφόρους.}\\\passeLigne}%
\end{groupe}%
%% v. 859-860a
\begin{groupe}%
\personnageGrec{Μνησίλοχος}%
\trimIamb{ἐμοὶ δὲ γῆ μὲν πατρὶς οὐκ ἀνώνυμος%
\commentaire{859-60}{}{\italique{Hélène}, v. 16-17a :%
\citComm{\difference{ἡμῖν} δὲ γῆ μὲν πατρὶς οὐκ ἀνώνυμος\\
Σπάρτη, πατὴρ δὲ Τυνδάρεως.}%
}\\
Σπάρτη, πατὴρ δὲ Τυνδάρεως.\\\passeLigne}%
\end{groupe}%
%% v. 860b-861
\begin{groupe}%
\personnageGrec{Κρίτυλλα}%
\trimIamb{\ligne{860}\decalage{Σπάρτη, πατὴρ δὲ Τυνδάρεως.}σοί γ᾽ ὦλεθρε\apparat{860}{σύ (ex σί) \Rdeux : \lecon{σοί} Zanetti (\italique{tibi} Divus) \espApp τ' R : \lecon{γ᾽} Gelenius \ligneApp{861} ἐστιν R : \lecon{ἐστι} \Mon}\\
πατὴρ ἐκεῖνός ἐστι; Φρυνώνδας μὲν οὖν.\testimonia{861}{(Φρυνώνδας) Suidas φ 700 ; \Syn φ 211 Φρυνώνδας $\sim$ Hesychius φ 938}%
\commentaire{861}{Φρυνώνδας.}{Phrynondas et Eurybatos, deux coquins dont le nom était devenu appellatif, comme en français Cartouche et Mandrin. (Note de Willems ; qui a en tête le \italique{Protagoras} (327d) de Platon lorsque ces deux truands de la tradition populaire sont rapprochés.)\\
\italique{Souda}, φ 770 :%
\citComm{Φρυνώνδας· τῶν ἐπὶ πονηρίᾳ διαβεβοημένων. ὃς ξένος ὢν κατὰ τὰ Πελοποννησιακὰ διέτριβεν Ἀθήνησιν. Ἀριστοφάνης Ἀμφιάρεῳ· ὦ μιαρὲ καὶ Φρυνώνδα καὶ πονηρὲ σύ. ἐκ τούτου τοὺς πονηροὺς Φρυνώνδας καλοῦσι. κακοήθης ὁ Φρυνώνδας οὗτος καὶ πανοῦργος, μισθοῦ τινας μηχανὰς πραγμάτων ἐπὶ κακοῖς ῥᾳδίως συντιθείς, ὥς φησιν Εὔδημος. καὶ αὖθις· ἀλλ' οἶμαι οὔτε Φρυνώνδας οὔτ' Εὐρύβατος οὐδ' ἄλλος πώποτε τῶν πάλαι πονηρῶν, τοιοῦτος μάγος καὶ γόης.\\
\francais{Phrynondas : de ceux connus pour leur méchanceté. C'était un étranger qui vivait à Athènes au moment de la guerre du Péloponnèse. Aristophane dans son \italique{Amphiaraos} : « ô misérable et Phrynondas et vil que tu es ». On appelle les méchants de son nom : Phrynondas.
Ce Phrynondas [était] mauvais et méchant ; quelqu'un qui, pour de l'argent, élaborait facilement quelque intrigue pour de mauvaises fins, comme le dit Eudemos.
Et ailleurs : « mais je pense que ni Phrynondas ni Eurybatos ni aucun autre des vilains d'autrefois ne se sont jamais comportés comme un filou et un tricheur ».}%
}%
}\\\passeLigne}%
\end{groupe}%
%% v. 862a
\begin{groupe}%
\ledpb%
\personnageGrec{Μνησίλοχος}%
\trimIamb{Ἑλένη δ᾽ ἐκλήθην.%
\commentaire{862}{}{\italique{Hélène}, v. 22 :%
\citComm{Ἑλένη δ᾽ ἐκλήθην. ἃ δὲ πεπόνθαμεν κακὰ}%
}\\\passeLigne}%
\end{groupe}%
%% v. 862b-863
\begin{groupe}%
\personnageGrec{Κρίτυλλα}%
\trimIamb{\ligne{862}\decalage{Ἑλένη δ᾽ ἐκλήθην.}αὖθις αὖ γίγνει γυνή,\apparat{862}{γίγνη R : \lecon{γίγνει} Brunck}\\
πρὶν τῆς ἑτέρας δοῦναι γυναικίσεως δίκην;\testimonia{863}{Suidas γ 496 γυναίκισις et Zonaras 459 πρὶν τῆς ἑτέρας γυναικίσεως δίκην δοῦναι}\apparat{863}{\lecon{δοῦναι post ἑτέρας} R : post δίκην Suid. γ 496}\\\passeLigne}%
\end{groupe}%
%% v. 864-865a
\begin{groupe}%
\personnageGrec{Μνησίλοχος}%
\trimIamb{ψυχαὶ δὲ πολλαὶ δι᾽ ἔμ᾽ ἐπὶ Σκαμανδρίοις\apparat{864}{Σκαμανδρίαις R : \lecon{Σκαμανδρίοις} Fritzsche ex Eur.\italique{Hel} 52 et 609}%
\commentaire{864-5}{}{\italique{Hélène}, v. 52-53a :%
\citComm{ψυχαὶ δὲ πολλαὶ δι᾽ ἔμ᾽ ἐπὶ Σκαμανδρίοις\\
ῥοαῖσιν ἔθανον· \difference{ἡ δὲ πάντα τλᾶσ᾽ ἐγὼ}}%
}\\
ῥοαῖσιν ἔθανον.\\\passeLigne}%
\end{groupe}%
%% v. 865b
\begin{groupe}%
\personnageGrec{Κρίτυλλα}%
\trimIamb{\ligne{865}\decalage{ῥοαῖσιν ἔθανον.}ὤφελες δὲ καὶ σύ γε.\\\passeLigne}%
\end{groupe}%
%% v. 866-868a
\begin{groupe}%
\personnageGrec{Μνησίλοχος}%
\trimIamb{κἀγὼ μὲν ἐνθάδ᾽ εἴμ᾽· ὁ δ᾽ ἄθλιος πόσις\\
οὑμὸς Μενέλεως οὐδέπω προσέρχεται.\apparat{867}{Μενέλαος R : \lecon{Μενέλεως} Brunck, cf. v.901}
\commentaire{866-8}{}{\italique{Hélène}, v. 49 :%
\citComm{κἀγὼ μὲν ἐνθάδ᾽ εἴμ᾽, ὁ δ᾽ ἄθλιος πόσις}%
\italique{Hélène}, v. 56a : %
\citComm{τί δῆτ᾽ ἔτι ζῶ; }%
}\\
τί οὖν ἔτι ζῶ;\\\passeLigne}%
\end{groupe}%
%% v. 868b
\begin{groupe}%
\personnageGrec{Κρίτυλλα}%
\trimIamb{\ligne{868}\decalage{τί οὖν ἔτι ζῶ;}τῶν κοράκων πονηρίᾳ.\testimonia{868}{Suidas τ 648}\apparat{868}{τῶν κορ. πον. continuat (Μνη.) R : \lecon{(Γυνή)} Schol., Ellebodius}\\\passeLigne}%
\end{groupe}%
%% v. 869-870
\begin{groupe}%
\personnageGrec{Μνησίλοχος}%
\trimIamb{ἀλλ᾽ ὥσπερ αἰκάλλει τι καρδίαν ἐμήν.\testimonia{869}{(αἰκάλλει τι καρδίαν ἐμήν) Suidas αι 161 ; Zonaras 91 (τι omisso)}\\
μὴ ψεῦσον ὦ Ζεῦ τῆς ἐπιούσης ἐλπίδος.\testimonia{870}{\GregCor pp. 14-15 ($=$ scholium\expo{Mosch.} Euripides, \italique{Hecub.} 225) (sed ἡμέρας pro ἐλπίδος) ; Suidas ψ 62 \espTest (μὴ ψεῦσον) \Antiatt 107.31}\apparat{post v. 870}{< οἶσθ᾽ οὖν ὃ δρᾶσον; $ - \cup - | \times - \cup - $ > \lecon{om.} R : restituit ap. Schol. Bibi (cf. Eur. \italique{Hel.}325)}%\apparat{870}{\lecon{ἐλπίδος} R : ἡμέρας ap. Greg. Cor. p.14-15, schol. ad Eur. \italique{Hec}.225 \espApp post v. 870 < οἶσθ᾽ οὖν ὃ δρᾶσον; $ - \cup - | \times - \cup - $ > \lecon{om.} R : restituit ap. Schol. Bibi (cf. Eur. \italique{Hel.}325)}%
%οἶσθ᾽ οὖν ὃ δρᾶσον; $ - \cup - | \times - \cup - $ %
\commentaire{870}{μὴ ψεῦσον ὦ Ζεῦ τῆς ἐπιούσης ἐλπίδος.}{Voici la scholie au vers 870 :% -- malheureusement lacunaire :%-- laisse à penser que le texte lui-même du manuscrit de Ravenne admet aussi quelque lacune. Voici la scholie :%
\citComm{σύνηθες τὸ σχῆμα. Μένανδρος·\\
\alinea οἶσθ᾽ ὃ ποίησον.\\
Σοφοκλῆς Πηλεῖ·\\
\alinea μὴ ψεῦσον ὦ Ζεῦ· μή μ᾽ ἕλης ἄνευ δορός.\\
\alinea οἶσθ᾽ οὖν ὃ δρᾶσον.\\
Εὐριπίδης Πολυΐδῳ.\\
\francais{La construction est usuelle. Ménandre :\\
\alinea Tu sais quoi faire. \\
Sophocle dans son \italique{Pelée} :\\
\alinea Ne me trompe pas, ô Zeus ; tu ne me prendras pas sans lutte.\\
\alinea Tu sais donc ce qu'il te reste à faire.\\
(Euripide dans son \italique{Polyidos}.)}}%
Cette scholie ne se peut comprendre que si l'on la met en regard d'un passage du Περὶ διαλέκτων de Grégoire de Corinthe (\nieme{XII} siècle, p. 14 ed. Schaefer) ou de la scholie de l'\italique{Hécube} d'Euripide (vers 225 -- nos trois commentaires remontent assurément à la même source ; cf. aussi Suidas ψ 62). À propos de la tournure `οἶσθ᾽ οὖν ὃ δρᾶσον;', voici ce que le scholiaste de l'\italique{Hécube} écrit :%
\citComm{δρᾶσον: τὰ καλὰ τῶν ἀντιγράφων δρᾶσον ἔχει. οἱ Ἀττικοὶ χρῶνται τοῖς προστακτικοῖς ἀντὶ ὑποτακτικῶν ἀπαγορευτικῶν, καὶ αὐτοῖς τοῖς προστακτικοῖς ἀντὶ ὁριστικῶν.\\
τοῦ μὲν οὖν πρώτου τὸ Ἀριστοφάνους ἐν Θεσμοφοριαζούσαις (v. 870) παράδειγμα `μὴ ψεῦσον, ὦ Ζεῦ, τῆς ἐπιούσης ἡμέρας (sic).'\\
τοῦ δευτέρου δὲ Μένανδρος `οἶσθ' ὅ τι ποίησον' ἀντὶ 
%<note type="footnote">21.
%ποιήσεις Matth. pro ποιήσῃς. 27. οἱ Ἀττικοὶ —] Scholion ex 26.
%δρᾶσον, superscripto εις, in Greg. Cor. p. 14. compositum. textu
%habet M. In Gr. est τὰ ἀκριβέστερα 30. Θεσμοφ.] 870. ubi ἐπιούσης
%τῶν ἀντιγράφων οἶσθ᾿ οὖν ὃ ἐλπίδος. δρᾶσον ἔχουσιν.</note>
τοῦ `ὅ τι ποιήσεις'.\\
τὸν αὐτὸν τρόπον καὶ Εὐριπίδης `οἶσθ᾿ οὖν ὃ δρᾶσον' ἀντὶ τοῦ `δράσεις' ἐξήνεγκεν.\\
δρᾶσον: les meilleures copies ont `δρᾶσον'.
Les Attiques recourent aux impératifs à place de subjonctifs de prohibition. Et de nouveau aux (je lis : καὶ αὖ τοῖς) impératifs à la place d'indicatifs.\\
L'exemple d'Aristophane sur le premier point : `μὴ ψεῦσον κτλ.'\\
Ménandre sur le second point : `οἶσθ' ὅ τι ποίησον' à la place de `ὅ τι ποιήσεις'.\\
Euripide a le même tour : `οἶσθ᾿ οὖν ὃ δρᾶσον' à la place de `δράσεις'.}%
Noter que ce tour d'Euripide se trouve justement dans l'\italique{Hélène} (vers 325) :%
\citComm{Χορός --- οἶσθ᾽ οὖν ὃ δρᾶσον; μνήματος λιποῦσ᾽ ἕδραν.\\
\francais{Hé bien, sais-tu ce que tu dois faire ? Quitte ton poste près du tombeau.}}%
Faut-il en induire qu'Aristophane le détournait à ce moment des \italique{Thesmophories} ? 
Cela ne me paraît pas impossible : ces mots seraient même tout-à-fait de saison dans la bouche de Mnésiloque : Euripide et lui ne se sont pas revus depuis la scène de travestissement chez Agathon ; comment Euripide saurait-il le plan échafaudé par Mnésiloque (détourner de ses pièces pour s'évader) ? il me paraît vraisemblable que Mnésiloque voyant apparaître Euripide lui demande s'il connaît bien son rôle -- rôle de Ménélas qu'il connaît forcément, l'ayant lui-même écrit. Façon habile de suggérer à Euripide quoi faire sans éveiller les soupçons de Critylla. Un vers de l'\allemand{Urtext} des \italique{Thesmophories} se sera-t-il perdu dans les scholies ? la conjecture est malheureusement invérifiable et restera circonscrite au commentaire :%
\citComm{Μνησίλοχος -- ἀλλ᾽ ὥσπερ αἰκάλλει τι καρδίαν ἐμήν.\\
μὴ ψεῦσον ὦ Ζεῦ τῆς ἐπιούσης ἐλπίδος.\\
οἶσθ᾽ οὖν ὃ δρᾶσον; $ - \cup - | \times - \cup - $}%
%Le point commun de ces citations est qu'elles comportent toutes deux « οἶσθ᾽ὃ ποίησον / δρᾶσον» ; si le scholiaste les a rapprochées du présent passage, il paraît probable que l'expression se fût aussi trouvée dans le texte original. 
%D'ailleurs, ces mots seraient tout-à-fait de saison dans la bouche de Mnésiloque : Euripide et lui ne se sont pas revus depuis la scène de travestissement chez Agathon ; comment Euripide saurait-il le plan échafaudé par Mnésiloque (détourner de ses pièces pour s'évader) ? il nous paraît vraisemblable que Mnésiloque  voyant apparaître Euripide lui demande s'il connaît bien son rôle -- rôle de Ménélas qu'il connaît forcément, l'ayant lui-même écrit. Façon habile de suggérer à Euripide quoi faire sans éveiller les soupçons de Critylla.\\
%Peut-être Aristophane détournait-il légèrement le vers 325 de l'\italique{Hélène} :
%\citComm{Χορός --- οἶσθ᾽ οὖν ὃ δρᾶσον; μνήματος λιποῦσ᾽ ἕδραν.\\
%\francais{Hé bien, sais-tu ce que tu dois faire ? Quitte ton poste près du tombeau.}}%
}%\\
\\\passeLigne}%
\end{groupe}%
%% v. 871-873
\begin{groupe}%
\ledpb%
\personnageGrec{Εὐριπίδης}%
\trimIamb{\ligne{871}τίς τῶνδ᾽ ἐρυμνῶν δωμάτων ἔχει κράτος,%\apparat{871}{ὡς Περσεύς in mg. R : \lecon{ὡς Μενέλαος} Zanetti \ligneApp{872} ξένος R : \lecon{ξένους} Zanetti \ligneApp{873} \lecon{κάμνοντας} R : καμόντας Lenting}%
\commentaire{871}{}{\italique{Hélène}, v. 68 :%
\citComm{Τεῦκρος --- τίς τῶνδ᾽ ἐρυμνῶν δωμάτων ἔχει κράτος;}%
\italique{Hélène}, v. 449 :%
\citComm{Μενελέως --- ναυαγὸς \difference{ἥκω} ξένος, \difference{ἀσύλητον γένος}.}}\\
ὅστις ξένους δέξαιτο ποντίῳ σάλῳ\apparat{872}{ξένος R : \lecon{ξένους} Zanetti}\\
κάμνοντας ἐν χειμῶνι καὶ ναυαγίαις;\apparat{873}{\lecon{κάμνοντας} R : καμόντας Lenting}\\\passeLigne}%
\end{groupe}%
%% v. 874a
\begin{groupe}%
\personnageGrec{Μνησίλοχος}%
\trimIamb{Πρωτέως τάδ᾽ ἐστὶ μέλαθρα.%
\commentaire{874-8}{}{\italique{Hélène}, v. 460 : %
\citComm{Πρωτεὺς τάδ᾽ οἰκεῖ δώματ᾽, Αἴγυπτος δὲ γῆ.}}\\\passeLigne}%
\end{groupe}%
%% v. 874b-876
\begin{groupe}%
\personnageGrec{Κρίτυλλα}%
\trimIamb{\ligne{874}\decalage{Πρωτέως τάδ᾽ ἐστὶ μέλαθρα.}ποίου Πρωτέως;%
\commentaire{874-6}{ποίου Πρωτέως; ... ἐπεὶ τέθνηκε Πρωτέας ἔτη δέκα.}{Elle confond Protée avec Protéas\marqueur{1}, stratège athénien dont il est question dans Thucydide\marqueur{2}. (Note de Willems.)
%Hésychius : Αἰγύπτιος: ὄνομα ἐθνικόν: “ “ἀθάνατος Πρωτεὺς Αἰγύπτιος citant Homère d 385
}%
\debutNote{874-6}%
\note{1}{Identiques au génitif.}%
\note{2}{Thucydide I, 45.2 ; II, 23.2.}\\
ὤ τρισκακοδαίμων, ψεύδεται νὴ τὼ θεώ,\apparat{875}{ὦ τρισκακόδαιμον R : \lecon{ὤ τρισκακοδαίμων} Wilamowitz ad \italique{Lys.}449}\\
ἐπεὶ τέθνηκε Πρωτέας ἔτη δέκα.\\\passeLigne}%
\end{groupe}%
%% v. 877
\begin{groupe}%
\personnageGrec{Εὐριπίδης}%
\trimIamb{ποίαν δὲ χώραν εἰσεκέλσαμεν σκάφει;%\commentaire{877}{εἰσεκέλσαμεν}{\eff{lire un epekelsamen ; cf. Hésychius ? erreur inverse de l'autre ? attention à la scansion !}}
\\\passeLigne}%
\end{groupe}%
%% v. 878a
\begin{groupe}%
\personnageGrec{Μνησίλοχος}%
\trimIamb{Αἴγυπτον.\\\passeLigne}%
\end{groupe}%
%% v. 878b
\begin{groupe}%
\personnageGrec{Εὐριπίδης}%
\trimIamb{\ligne{878}\decalage{Αἴγυπτον.}ὦ δύστηνος οἷ πεπλώκαμεν.\testimonia{878}{Suidas οι 1 ὦ δύστηνος, οἷ πεπλώκαμεν;}%
\commentaire{878}{}{\italique{Hélène}, v. 461 :%
\citComm{Αἴγυπτος; ὦ δύστηνος, οἷ \difference{πέπλευκ᾽ ἄρα.}}}\\\passeLigne}%
\end{groupe}%
%% v. 879-880
\begin{groupe}%
\personnageGrec{Μνησίλοχος}%
\trimIamb{πείθει τι τῷ κακῷ κακῶς ἀπολουμένῳ\apparat{879}{τῷ R : τούτῳ τῷ Grynaeus (1544) : \lecon{τῷ κακῷ} Scaliger \ligneApp{880} \lecon{Θεσμοφόριον} R, cf. v.278 : Θεσμοφορεῖον Meineke}\\
ληροῦντι λῆρον; Θεσμοφόριον τουτογί.\\\passeLigne}%
\end{groupe}%
%% v. 881
\begin{groupe}%
\personnageGrec{Εὐριπίδης}%
\trimIamb{αὐτὸς δὲ Πρωτεὺς ἔνδον ἔστ᾽ ἢ 'ξώπιος;\apparat{881}{ἠξώπιος R, id. v.884 : \lecon{ἢ 'ξώπιος;} Biset}%
\commentaire{881}{}{\italique{Hélène}, v. 465 :%
\citComm{Μενελέως --- \difference{ποῦ δῆτ᾽ ἂν εἴη; πότερον ἐκτὸς ἢ 'ν δόμοις;}}}\\\passeLigne}%
\end{groupe}%
%% v. 882-884
\begin{groupe}%
\personnageGrec{Κρίτυλλα}%
\trimIamb{οὐκ ἔσθ᾽ ὅπως οὐ ναυτιᾷς ἔτ᾽ ὦ ξένε,\testimonia{882}{(ναυττιᾶν $=$ ναυσιᾶν) Moeris ν 9}\\
ὅστις γ᾽ ἀκούσας ὅτι τέθνηκε Πρωτέας\apparat{883}{ὅστις R : \lecon{ὅστις γ᾽} Scaliger}\\
ἔπειτ᾽ ἐρωτᾷς `ἔνδον ἔστ᾽ ἢ 'ξώπιος;'\\\passeLigne}%
\end{groupe}%
%% v. 885
\begin{groupe}%
\personnageGrec{Εὐριπίδης}%
\trimIamb{αἰαῖ τέθνηκε. ποῦ δ᾽ ἐτυμβεύθη τάφῳ;\apparat{885}{τέθνηκεν R : \lecon{τέθνηκε} Bentley}\\\passeLigne}%
\end{groupe}%
%% v. 886
\begin{groupe}%
\personnageGrec{Μνησίλοχος}%
\trimIamb{τόδ᾽ ἐστὶν αὐτοῦ σῆμ᾽, ἐφ᾽ ᾧ καθήμεθα.%
\commentaire{886}{}{\italique{Hélène}, v. 466 :%
\citComm{Γραῦς --- τόδ᾽ ἐστὶν αὐτοῦ \difference{μνῆμα, παῖς δ᾽ ἄρχει χθονός}.}}\\\passeLigne}%
\end{groupe}%
%% v. 887-888
\begin{groupe}%
\personnageGrec{Κρίτυλλα}%
\trimIamb{κακῶς τ᾽ ἄρ᾽ ἐξόλοιο κἀξολεῖ γέ τοι,\apparat{887}{\lecon{κακῶς τ᾽ ἄρ᾽} R : κακῶς γ᾽ ἄρ᾽ Denniston (\italique{CR} XLIV(1930), p.215)}\\
ὅστις γε τολμᾷς σῆμα τὸν βωμὸν καλεῖν.\\\passeLigne}%
\end{groupe}%
%% v. 889-890a
\begin{groupe}%
\personnageGrec{Εὐριπίδης}%
\trimIamb{τί δὴ σὺ θάσσεις τάσδε τυμβήρεις ἕδρας\apparat{889}{τί δὲ σὺ R : τί δαὶ σὺ Faber : \lecon{τί δὴ σὺ} Scaliger}\\
φάρει καλυπτὸς ὦ ξένη;\\\passeLigne}%
\end{groupe}%
%% v. 890b-891
\begin{groupe}%
\personnageGrec{Μνησίλοχος}%
\trimIamb{\ligne{890}\decalage{φάρει καλυπτὸς ὦ ξένη;}βιάζομαι\\
γάμοισι Πρωτέως παιδὶ συμμεῖξαι λέχος.\apparat{891}{συμμῖξαι R : \lecon{συμμεῖξαι} van Leeuwen}\\\passeLigne}%
\end{groupe}%
%% v. 892-894
\begin{groupe}%
\personnageGrec{Κρίτυλλα}%
\trimIamb{τί ὦ κακόδαιμον ἐξαπατᾷς αὖ τὸν ξένον;\apparat{892}{αὐτὸν (ex αὖτὸν) \Rpc : \lecon{αὖ τὸν} Casaubon}\\
οὗτος πανουργῶν δεῦρ᾽ ἀνῆλθεν ὦ ξένε\\
ὡς τὰς γυναῖκας ἐπὶ κλοπῇ τοῦ χρυσίου.%
\commentaire{894}{ὡς τὰς γυναῖκας ἐπὶ κλοπῇ τοῦ χρυσίου,}{«~Subtiliser les bijoux des femmes~» ; voilà un chef d'accusation nouveau contre Mnésiloque, d'autant plus étonnant que rien jusqu'ici ne l'a laissé présager. 
Le fait que Critylla porte un flambeau indique pourtant qu'elle est présente depuis l'arrivée de la procession. Ment-elle pour éviter d'avoir à raconter la trop longue affaire ? il pourrait y avoir une autre explication. D'abord, l'on dispose de lois du \nieme{V} siècle av. J.-C. interdisant tout habillement luxueux aux couleurs vives lors des cultes à Déméter\marqueur{1} :
\citComm{[Εἰ γυ]νὰ Fέσετοι ζτεραῖον λ\chap{ο}πος, | [ἰερὸ]ν \chap{ἐ}ναι ταῖ Δάματρι ταῖ Θεσμοφόροι· | [εἰ δὲ] μὲ ὑνιερόσει δυμενὲς ἔασα ἐπὲ Fέργο | [κακ\chap{ο}ς] ζ᾽ ἐξόλοιτυ, κὰ ὄζις τότε δαμιοFοργ\chap{ε} | [ἀφάε]σται δαρχμὰς τριάκοντα· εἰ δὲ μὲ ἀφάετοι, | [μεδ᾽ ἰν] τὰν ἀσέβειαν ἔχε ὄδε κῦρος, δέκο Fέτεα \chap{ἐ}να[ι | ἰναγὲς] τ\chap{ο}δε.\\
Lorsqu'une femme met un vêtement de couleur\marqueur{2}, celui-ci doit être dédié à Déméter Thesmophore. Mais si elle ne le consacre pas, en étant hostile à la cérémonie, elle périra, et celui qui sera alors le fonctionnaire civil aura trente drachmes. Mais si ce dernier n'a pas l'autorité pour empêcher l'indignité, il sera souillé par elle pendant dix ans.}
Seules les femmes tenant un rôle actif lors de ces cérémonies portaient des bijoux (prêtresses, canéphores, \&c.)\marqueur{3}.\\
En fait, le terme χρυσίον, \italique{bijou}, \italique{ornement} semble avoir parfois un sens imagé chez Aristophane et désigner les \italique{appâts} féminins. Ainsi, lorsque Pisthétaire s'exclame à propos de la belle \& court vêtue joueuse de flûte Procné\marqueur{4} :%
\citComm{Πισθέταιρος --- ὅσον δ᾽ ἔχει τὸν χρυσόν, ὥσπερ παρθένος,}%
il est bien évident qu'il fait l'éloge de sa beauté, son galbe, ou que sais-je encore\marqueur{5} et non pas de sa verroterie. Ce blason sera d'ailleurs continué quelques vers plus loin par le chœur\marqueur{6}. 
Autre exemple : lorsque Dicéopolis fait venir la jeune canéphore en lui prodiguant ce conseil (vers 257-258) :%
\citComm{πρόβαινε, κἀν τὤχλῳ φυλάττεσθαι σφόδρα\\
μή τις λαθών σου περιτράγῃ τὰ χρυσία.\\
\francais{Avance, et surtout prends garde que dans la foule\\
on ne te grignote tes bijoux.}}%
Certes, comme on l'a vu, les canéphores portaient des bijoux ; mais on devine aisément avec Aristophane qu'en voyant cette jeune beauté dont Dicéopolis vient de dire : « Heureux celui qui t'épousera et te procréera des belettes », la foule va ressentir quelque pulsion ne relevant pas uniquement de la cleptomanie.
Le verbe περιτρώγειν paraît d'ailleurs bien impropre pour une image de vol à la tire\marqueur{7}.
Quant à κλοπῇ, Mnésiloque nous a déjà fait remarquer au vers 205 que le verbe κλέπτειν pouvait par euphémisme désigner une activité qui n'a de commun avec le vol que de se passer au plus noir de la nuit.\\
Terminons en indiquant que cette accusation de tentative de viol lors d'une réunion exclusivement féminine se retrouve ailleurs ; ces assemblées féminines à huis-clos et que la rumeur publique voyait comme des séances de débauche sexuelle avaient tout pour attiser la curiosité de quelque luxurieux. C'est à une telle tentative d'un certain Critias qu'un fragment d'Hipponax\marqueur{8} fait allusion :%
\citComm{οὔ μοι δικαίως μοιχὸς ἁλῶναι δοκέει\\
Κριτίης ὁ Χῖος ἐν $ \cup - \times - $ δούμῳ\\
Ce n'est pas selon moi par justice\marqueur{9} qu'en tant que débauché fût saisi\\
Critias de Khios dans $[$ … $]$ une assemblée féminine.}}%
\debutNote{894}%
\note{1}{R. Koerner, \italique{Inschriftliche Gesetztexte der Früher Grieschischen Polis} (1993), 35 $=$ \italique{Lois sacrées des cités grecques. Supplément} 32 $=$ \italique{Supplementum Epigraphicum Graecum} XI, 1112. Je reprends le texte établi par Koerner.}%
\note{2}{Et pourtant Mnésiloque rejoint les Thesmophories avec sa crocote couleur safran !.. au milieu d'une féminine assemblée sobrement vêtue ? (de crocotes et de mitres tout-de-même, cf. le vers 941, (ἐ)ν κροκωτοῖς καὶ μίτραις ; mais la crocote était-elle nécessairement de la couleur du κρόκος dont elle tire son nom ?) 
Quoi qu'il en soit, une arrivée cocasse de Mnésiloque avec une robe chatoyante et qui dénote avec celles du reste de l'assemblée ne ferait que renforcer par contraste le comique de situation.}%
\note{3}{Ce que semble confirmer un passage de \italique{Lysistrata} (v. 1189-1193) : c'est le chœur des femmes qui parle :%
\citComm{στρωμάτων δὲ ποικίλων καὶ\\
χλανιδίων καὶ ξυστίδων καὶ\\
\textbf{χρυσίων}, ὅσ᾽ ἐστί μοι,\\
οὐ φθόνος ἔνεστί μοι πᾶσι παρέχειν φέρειν\\
τοῖς παισίν, \textbf{ὁπόταν} τε θυγάτηρ τινὶ \textbf{κανηφορῇ}.\\
Et de surtouts brodés, petites chlamides, xystides, \italique{bijoux}, tout ce que je possède, il ne me coûte pas d'offrir à chacun de l'emporter, pour ses enfants et pour sa fille, \italique{si elle devenait canéphore}.}}%
\note{4}{\italique{Oiseaux}, v. 669.}%
\note{5}{Hésychius (χ 792) : χρυσίον· τὸ τῶν παιδίων αἰδοῖον. Henderson (\italique{Maculate Muse}, p. 131) y flaire un jeu de mot entre τὸν χρυσόν et τὸν κυσόν.}%
\note{6}{Vers 676-684.}%
\note{7}{Si impropre que l'on en serait tenté de comprendre λαβών à la place de λαθών.}%
\note{8}{Fragment 30 de l'édition au fort utile commentaire d'Olivier Masson ($=$ fragment 42 de l'édition Degani). Je reprends la correction proposée par Masson : δούμῳ à la place de δούλῳ, qui s'appuie sur une correction identique chez Hésychius (δ 2256) :%
\citComm{δοῦμος· ἡ οἰκία, ἢ τὴν ἐπὶ τὸ αὐτὸ συνέλευσιν τῶν γυναικῶν,\\
doumos : la maison. \italique{Item} : une réunion en commun des femmes.}}%
\note{9}{Je lis : οὔ μοι δίκᾳ ὡς μοιχὸς ἁλῶναι δοκέει κτλ.}%
%\note{1}{M. Dillon, \italique{Girls and Women in Classical Greek Religion}, 2004, p. 261-262}
\\\passeLigne}%
\end{groupe}%
%% v. 895
\begin{groupe}%
\personnageGrec{Μνησίλοχος}%
\trimIamb{βάϋζε τοὐμὸν σῶμα βάλλουσα ψόγῳ.\apparat{895}{βαϋζα R : \lecon{βάϋζε} Zanetti}%
\commentaire{895}{βαϋζα τοὐμὸν σῶμα βάλλουσα ψόγῳ.}{Taillardat (§569) avait observé que le tour βάλλειν τινὰ ψόγῳ, « accabler quelqu'un de reproches » est d'un ton tragique fait sur le modèle de βάλλειν τινὰ λίθοις\marqueur{1}. (Et réciproquement, les mots sont bien souvent considérés comme des projectiles, cf. Taillardat §502.) 
De même, la périphrase τοὐμὸν σῶμα pour ἐγώ est d'un style élevé et particulièrement prisée d'Euripide\marqueur{2}.
La correction de la leçon du \italique{Ravennas} universellement admise, βαϋζα en βάϋζε, me semble des plus douteuses. 
Las ! on ne détruit que ce que l'on remplace...%\\
%\eff{ Hésychius 352 βᾶϋ· εἶδος ἄνθους
%ζᾶ: μέγα.ἰσχυρόν.πολύ. ὅθενκαὶ ζάπλουτοςὁ μεγαλόπλουτος Cf. Photius zêta 12}
%Anthologie Palatine : τέχνας εἵνεκα σεῖο καὶ ἁ λίθος οἶδε βρυάζειν,
%Πραξίτελες: λῦσον, καὶ πάλι κωμάσομαι.
%νῦν δ᾽ ἡμῖν οὐ γῆρας ἔτ᾽ ἀδρανές, ἀλλ᾽ ὁ πεδητὰς
%Σειληνοῖς κώμων βάσκανός ἐστι λίθος.
%βρυάζει:πάλιν τρυφερῶςδιάκειται,θάλλει,εὐφραίνεται Synagoge cf Photius
%	βρύα: ἃ γίνεται μὲν ἐπὶ πετρῶν, λέγεται δὲ καὶ σφακχὸςἢ
%	Τσκαφίς.ἐπετίθετοδὲ ἁγνισμοῦχάριν
%	αἰσχρὰ καθʼ ἡμετέρης ἔβλυσε παρθενίης => BLUZA
%	βαυκαλίζω jouer à la prude
%	ὄβρυζος
%	βρῦ(ν) ζαλάω  βρῦν εἰπεῖν Ar.Nu.1382
%	βαυβίζω
%	Archil. 46
%	χορυζᾷ Τξαίνει.μωραίνει vere en zato avec haplographie ? ZATOTOEMON
%Le βαϋζα du \italique{Ravennas} a été universellement corrigé en βάϋζε, correction qui soulève plusieurs objections. \ \ \  1° \ \ \ Au niveau du sens : après les graves accusations que Critylla vient de porter à l'encontre de Mnésiloque-Hélène, on s'attend à une défense vigoureuse : iel vient d'être diffamé.e ! le simple reproche de forme : βάϋζε τοὐμὸν σῶμα, "tu jappes un peu fort après moi" nous paraît une défense un peu légère. \ \ \ 2° \ \ \ Si la confusion αι/ε, est courante (en raison d'une prononciation identique), ce vers serait l'unique exemple de la pièce d'une substitution finale de ε par α.\\
%Faisons l'hypothèse que la fin du mot transmis soit correcte. Le suffixe -υζᾰ a été productif de quelques noms de végétaux en grec\marqueur{ἄλυζα $=$ ἄλυπον ; κνύζα	 $=$ κόνυζα $=$ σκόνυζα, \italique{conyze} ou \italique{herbe aux puces} ou \italique{encensière} (\latin{erigeron viscosum L.}) ; ἡ μώλυζα, \italique{la tête d'ail} ; ἡ ὄρυζα, \italique{le riz}.} ; ἄλυζα concorderait avec la leçon du manuscrit. C'est chez Hésychius qu'on trouve le mot (α 3287) :%
%\citComm{ἄλυζα· ἄλυπον\\
%Alyza : Turbith blanc (\latin{Globularia alypum})}%
%Le sens s'éclaire ainsi : Mnésiloque se compare à une plante -- une fleur dirai-t-on en français -- en proie aux jets de reproches de reproches de Critylla. C'est là le point : ψόγῳ à la fin du vers est mis πάρα προσδοκίαν pour λίθῳ. Or il nous est justement parvenu un proverbe grec :%
%\citComm{πόᾳ λίθον ἐκβάλλει· ἐπὶ τῶν τἀναντία τοῖς ἐπιβούλοις ποιούντων}%
%\eff{Hésychius : ἄλυζα: ἄλυπον   et  κνύζα: beaucoup de noms de fleurs en -uza
%οὔτ᾽ἐκ χερὸς μεθέντα καρτερὸνλίθον on ne peut pas plus rattraper le mot lâché que la pierre (Eur.)
%Jeter une pierre à une fleur Πόᾳ λίϑον ἐκβάλλει: ἐπὶ τῶνταναντίατοῖςἐστιβούλοιςποιούντων. Ὃ γὰρ δρυοκολάπτης τὸ ζῷον,
%ἔχειμὲν ῥάμφοςἐπίκυρτον,κολάσιτεεδὲ ἄρα τούτῳτὰς
%δρῦς,καὶ ἐνταυϑοῖκοιλάναςτὸν τόπον»,ὡς δἰς καλλιὰ}
%
%Élien : Aelian, De Natura Animalium I 45 : γυπῶν πτερὰ εἰ θυμιάσειέ τις, ὡς ἀκούω, καὶ ἐκ φωλεῶν καὶ ἐξ εἰλυῶν τοὺς ὄφεις προάξει ῥᾷστα. τὸ ζῷον ὁ δρυοκολάπτης ἐξ οὗ δρᾷ καὶ κέκληται. ἔχει μὲν γὰρ ῥάμφος ἐπίκυρτον, κολάπτει δὲ ἄρα τούτῳ τὰς δρῦς, καὶ ἐνταυθοῖ ὡς ἐς καλιὰν τοὺς νεοττοὺς ἐντίθησιν, οὐ δεηθεὶς καρφῶν καὶ τῆς ἐξ αὐτῶν πλοκῆς καὶ οἰκοδομίας οὐδὲ ἕν. οὐκοῦν εἴ τις λίθον ἐνθεὶς ἐπιφράξειε τῷ ὀρνέῳ τῷ προειρημένῳ τὴν ἔσδυσιν, ὃ δὲ συμβαλὼν τὴν ἐπιβουλὴν [p. 25] κομίζει πόαν ἐχθρὰν τῷ λίθῳ: ὃ δὲ οἷα βαρούμενος καὶ μὴ φέρων ἐξάλλεται, καὶ ἀνέῳγεν αὖθις τῷ προειρημένῳ ἡ φίλη ὑποδρομή.
%J’ai appris qu’en faisant des fumigations de plumes de vautour on parvient sans difficulté à faire sortir les serpents de leurs repaires ou de leurs tanières. L’animal qu’on nomme pivert doit son nom à ses façons 44 : il a un bec recourbé à son extrémité, dont il se sert pour becqueter les chênes, et c’est à l’intérieur de l’arbre qu’il loge ses petits comme si c’était un nid ; et il n’a nul besoin de se procurer des brindilles, de les tresser ensemble et de construire un nid. Or, si l’on obstrue l’ouverture que s’est ménagée notre oiseau en y mettant une pierre, il prend acte de la manœuvre et va chercher une herbe qui a un pouvoir antipathique sur la pierre, et il la place dessus. La pierre, accablée par ce poids et incapable de le supporter, s’éjecte, et voilà de nouveau ouvert à notre oiseau l’accès à son précieux abri.
%\\
}%
\debutNote{895}%
\note{1}{Βάλλειν κακοῖς (Sophocle, \italique{Ajax} 1244), αἰτίᾳ (\italique{id}. \italique{Trach.} 940), φθόνῳ (Eur. \italique{Élect.} 902). Cf. ἀτιμίῃσιν ἰάλλειν (Homère, \italique{Odyssée} XIII, 142), λόγοις ἰάπτειν (Sophocle, \italique{Ajax} 501).}%
%\note{2}{ de ses \italique{Images d'Aristophane}.}%
\note{2}{Voir le \LSJ \ à l'entrée σῶμα II.}\\\passeLigne}%
\end{groupe}%
%% v. 896
\begin{groupe}%
\personnageGrec{Εὐριπίδης}%
\trimIamb{ξένη τίς ἡ γραῦς ἡ κακορροθοῦσά σε;
\commentaire{896}{κακορροθοῦσά.}{Le vocabulaire est euripidéen -- du véritable Euripide s'entend\marqueur{1}. Le verbe κακορροθεῖν, substitut expressif de κακολογεῖν, est un terme qu'on retrouve exclusivement dans ses tragédies -- sauf reprise aristophanesque (\italique{Acharniens}, v. 576).}%
\debutNote{896}%
\note{1}{Aristophane possède un talent sans pareil pour nous croquer ses contemporains -- au moins ceux qui en valaient la peine à ses yeux. Pour s'en assurer, on peut relire le \italique{Banquet}, l'\italique{Apologie} et les autres pièces de Platon où Socrate apparaît pour constater avec amusement que les traits grossiers qu'Aristophane dessine, Platon les a esquissés discrètement. Mais lus avec attention, ses dialogues nous révèlent le même Socrate que celui des \italique{Nuées} : un va-nu-pieds à l'hygiène si douteuse qu'il lui faut des circonstances exceptionnelles pour prendre un bain ; un caractère irascible et violent ; menant une vie essentiellement recluse ; faisant montre d'un mépris affiché envers les «~faiseurs de comédies~» ; etc., etc. Oui : Aristophane était dans le vrai dans ses \italique{Nuées} ! et Socrate n'a été grand que dans son dernier combat pour affronter la rage plébéienne, dernier sursaut d'humeur d'Athènes la décadente qui ignorait alors qu'elle avait déjà sombré définitivement, coulée par sa démocratie aux relents impérialistes -- démocratie en rut qui avait au moins le mérite de donner quelque signe de vie contrairement à celles qu'on connaît. Si Socrate n'avait pas été héroïquement retranché du jeune Platon, il ne fait à mon sens aucun doute que le plus grand des Grecs aurait fini par s'éloigner du maître de sa jeunesse. Au lieu de devenir Platon, serait-il devenu Shakespeare ? y aurions-nous gagné ? Il ne se serait en tout cas certainement pas abîmé avec l'âge arrivant dans une métaphysique aporétique, et si détachée des réalités qu'elle révolta Aristote -- comment comprendre la boulimie de ce dernier pour les sciences naturelles, si ce n'est qu'il allait à rebours de son maître ? Ce très curieux phénomène psychologique, ce mysticisme de la soixantaine qui s'élève au-dessus de nos réalités bassement terrestres -- chez des hommes d'action (!) et d'une trempe exceptionnelle (!) -- a touché de la même manière notre Platon français, Pierre Boutang. À ma connaissance, la plus belle description qui en a été faite -- outre l'expérience personnelle qui ne se partage pas -- se trouve dans le chef-d'œuvre de Jean Boudou, \italique{Lo libre de Catòia}, avec le personnage du grand-père.}%
\\\passeLigne}%
\end{groupe}%
%% v. 897a
\begin{groupe}%
\personnageGrec{Μνησίλοχος}%
\trimIamb{αὕτη Θεονόη Πρωτέως.\\\passeLigne}%
\end{groupe}%
%% v. 897b-899a
\begin{groupe}%
\personnageGrec{Κρίτυλλα}%
\trimIamb{\ligne{897}\decalage{αὕτη Θεονόη Πρωτέως.}μὰ τὼ θεὼ\\
εἰ μὴ Κρίτυλλά γ᾽ Ἀντιθέου Γαργηττόθεν·
\commentaire{898}{εἰ μὴ Κρίτυλλά γ᾽ Ἀντιθέου Γαργηττόθεν·}{Je signale à l'attention d'un éventuel exégète parcourant ces pages que la réplique de Critylla trouve un écho dans un fragment de la comédie perdue d'Aristophane des Δαιταλεῖς, \italique{Convives}, transmise par Pollux (X, 158 $=$ fr. 226(217) Kassel-Austin) :%
\citComm{εἰ μὴ δικῶν τε γυργάθους ψηφισμάτων τε θωμοὺς\\
< φέροντες >}%
Le rapport est flagrant : même formule commençant par «~εἰ μὴ~» ; Κρίτ-υλλά, nom formé du suffixe -υλλα\marqueur{1} et du radical κρίτ- qui fait entendre le \italique{jugement} a le même champs sémantique que les δικῶν ; Γαργηττόθεν et γυργάθους sont très proches phonétiquement ainsi que -τιθέου (de Ἀντιθέου) et τε θωμοὺς.\\
Athénée (I, 4d) évoque ce même fragment d'Aristophane dans son premier livre ; ce qui nous permet d'en préciser le contexte. Dans un chapitre assez décousu\marqueur{2} consacré aux propos de table, on comprend qu'il est certainement question de sycophantes -- pique-assiettes et autres croque-lardons --, dont la vie est douce, bonne et grasse :%
\citComm{τὸν βίον, φησίν, εὐσταθεῖς· οὐκ ἐγχειρογάστορες, γυργάθους ψηφισμάτων φέροντες, Ἀριστοφάνης φησίν}%
%Avec ces faibles indices, on peut raisonnablement penser que -τιθέου et τε θωμοὺς font entendre un "τις θεός" ou "τις θεῶν" dans la locution originale. C'est bien cela : Critylla jure ses grands dieux et invoque la justice... On en arrive à la phase la plus délicate : à quoi renvoie Γαργηττόθεν ou γυργάθους ? (noter le suffixe -θεν de Γαργηττόθεν, sûrement mis à dessein.) 
%Or l'oreille grecque devait percevoir comme particulièrement proches les sons `γ' et `κ'\marqueur{\eff{Retrouver ; c'est dans Érasme avec des histoires de carcinos et de gargalos karkanos ?}}. On rétablit donc, d'après le sens et l'homophonie, la présence d'un κατώρθωσεν dans le proverbe détourné. \\
Comme l'oreille grecque percevait comme particulièrement proches les sons `γ' et `κ', faut-il entendre la locution grecque :%
\citComm{οὐκ ... εἰ μὴ τὴν δίκην τις θεῶν κατώρθωσεν,\\
À moins qu'un dieu ne rétablisse la justice,}%
conservée par Isocrate 
\citComm{οὐκ ἄν ποτέ σοι γενέσθαι τηλικαῦτα τὸ μέγεθος, \textbf{εἰ μή τις θεῶν} αὐτὰ \textbf{συγκατώρθωσεν}\\
\italique{Orationes}, V (Φίλιππος), 152.}%
\citComm{οὐ \textbf{δίκαιός} ἐστιν μετέχειν \textbf{εἴ τι Τιμόθεος} πράττων \textbf{μὴ κατώρθωσεν}\\
	\italique{Orationes}, XV (Περὶ ἀντιδόσεως), 105.}%
%Chance : Isocrate semble avoir fait allusion à la même locution :%
%\citComm{οὐκ εἴ τις μιᾷ τύχῃ τηλικοῦτόν τι κατώρθωσεν\\
%\italique{Orationes}, XV (Περὶ ἀντιδόσεως), 128.}%
Qu'en peut-on conclure ? pas grand-chose pour l'heure... %, si ce n'est qu'il devait exister un tour grec, s'en remettant à la justice divine ressemblant peu ou prou à : %
Il est vraisemblable qu'Aristophane y fasse aussi allusion au vers 1360 des \italique{Cavaliers}, quand un synégore passé maître dans le chantage menace :%
\citComm{Οὐκ ἔστιν ὑμῖν τοῖς δικασταῖς ἄλφιτα,\\
εἰ μὴ καταγνώσεσθε ταύτην τὴν δίκην}%
%\citComm{οὐκ ... εἰ μὴ τὴν δίκην τις θεῶν κατώρθωσεν\\
%À moins qu'un dieu ne rétablisse la justice.}%
%Si nous avons vu juste, alors il est vraisemblable qu'Aristophane y fasse aussi allusion au vers 1360 des \italique{Cavaliers}, quand un synégore passé maître dans le chantage menace :%
%\citComm{Οὐκ ἔστιν ὑμῖν τοῖς δικασταῖς ἄλφιτα,\\
%εἰ μὴ καταγνώσεσθε ταύτην τὴν δίκην}%
%Je vous passe le relais, ô lecteur(s), et m'en remets à vous pour pousser plus avant la recherche. 
%Cavaliers 1360 εἰ μὴ καταγνώσεσθε ταύτην τὴν δίκην / dikê = trois(oboles) ?
%NONLa métaphore semble halieutique : ἐγχειρογάστορες, \italique{ceux qui vivent du travail de leurs mains}, est certainement à prendre au sens spécialisé de \italique{pêcheurs}\marqueur{Hésychius}. Le 
%εἰ μή τις θεῶν αὐτὰ συγκατώρθωσεν isocrate
%εἰ μὴ δικῶν τε γυργαθοὺς ψηφισμάτων τε θωμοὺς φέροντες
%et des Cavaliers 979
% Καίτοι πρεσβυτέρων τινῶν
%οἵων ἀργαλεωτάτων
%ἐν τῷ δείγματι τῶν δικῶν
%ἤκουσ’ ἀντιλεγόντων, 
%Athénée  :
%On devine que les ἐγχειρογάστορες au sens de pêcheurs, le gurgathos étant un diktuon
%τὸν βίον, φησίν, εὐσταθεῖς, οὐκ ἐγχειρογάστορες. γυργάθους ψηφισμάτων φέροντες, Ἀριστοφάνης φησίν
%	Ach 684 : οὐχ ὁρῶντες οὐδὲν εἰ μὴ τῆς δίκης τὴν ἠλύγην.
%	Cav 1360  «Οὐκ ἔστιν ὑμῖν τοῖς δικασταῖς ἄλφιτα,
%	εἰ μὴ καταγνώσεσθε ταύτην τὴν δίκην»
%Κρίτυλλά
%Aristophane emploie ailleurs le prénom de Critylla : au vers 324 de \italique{Lysistrata} :
%\citComm{Χορὸς Γυναικῶν --- πέτου πέτου Νικοδίκη,\\
%πρὶν ἐμπεπρῆσθαι Καλύκην\\
%τε καὶ Κρίτυλλαν περιφυσήτω\\
%\francais{$[$Demi-$]$chœur des femmes --- Vole, vole, Nicodicé,\\
%avant que, léchées de toutes parts par les flammes, Calycé\\
%et Critylla périssent.}}
%Le demi-ch\oe ur des hommes est en train de bouter le feu aux barricades féminines ; le demi-ch\oe ur des femmes sonne alors le tocsin et prend ses dispositions « avant que ne soient brûlées Calycé et Critylla ».\\
%Calycé est mise ici car son nom évoque le bois\marqueur{1}, bois qui ne va pas tarder à se consumer si les femmes ne réagissent promptement. Il en va de même pour Critylla : son nom est formé du suffixe -υλλα\marqueur{2} et du radical κρίτ- qui fait entendre le pistachier térébinthe (\latin{Pistacia Terebinthus}) : Hésychius 4154 :
%\citComm{κρίτανος\marqueur{3}· τέρμινθος\\
%\francais{Critanos : térébinthe.}}
%Le pistachier térébinthe -- bois dur et dense -- en font entre autres choses un bois de chauffage. C'est pour cette raison que le choeur des femmes l'évoque dans \italique{Lysistrata} et pour la même raison qu'Aristophane la baptise ainsi ici. Car Critylla est la manieuse de torche\marqueur{4} ; et elle a déjà montré à Mnésiloque de quel bois elle se chauffait avec sa torche qu'elle a reconverti en gourdin.
}%
\debutNote{898}%
%\note{1}{\eff{peuplier noir ? Daremberg ?}}
\note{1}{Le suffixe en -υλλα/ιλλα est courant pour les noms de femme : Κύδιλλα, Ἥρυλλα, ou Ξένυλλα qui est apparue plus haut \&c. (voir \italique{La formation des noms en grec ancien} de Chantraine, p. 252.)}%
\note{2}{Le premier livre, duquel le passage est tiré, ne nous est connu que par l'\italique{épitomé} des \italique{Deipnosophistes}, ce qui complexifie encore la délimitation et l'attribution des citations... Voir l'article de Benoît Louyest, «~L'épitomé du \italique{Banquet des sophistes} d'Athénée~», in \italique{Rursus} 8 (2012).}%
%\note{3}{Noter l'emplacement de l'accent : que Κρίτυλλα soit proparoxyton permet de rapprocher son nom de κρίτανος \eff{la ville de Critalla ?} tandis que d'autres noms comme Κιρτίας ou Κριτόλα sont à rapprocher du verbe κρίνειν.\eff{Et Cri'ton ?}}
%\note{4}{Voir les vers 853 et 858.}
\commentaire{898}{Γαργηττόθεν·}{La scholie :%
\citComm{ἀπὸ δήμου τῆς Αἰγηίδος φυλῆς, ἔνθα τὸ Εὐρυσθέως ἀπόκειται σῶμα ἐν τῷ Ἀρατηρίῳ. ἐπίρρημα δὲ τὸ Γαργηττόθεν.\\
D'un dème de la tribu Égéide, c'est là que repose le corps d'Eurysthée dans l'Arêtérion. D'où le nom de `Gargêttothen'.}%
}\\
σὺ δ᾽ εἶ πανοῦργος.\\\passeLigne}%
\end{groupe}%
%% v. 899b-901
\begin{groupe}%
\personnageGrec{Μνησίλοχος}%
\trimIamb{\ligne{899}\decalage{σὺ δ᾽ εἶ πανοῦργος.}ὁπόσα τοι βούλει λέγε.\\
οὐ γὰρ γαμοῦμαι σῷ κασιγνήτῳ ποτὲ\\
προδοῦσα Μενέλεων τὸν ἐμὸν ἐν Τροίᾳ πόσιν.\apparat{901}{Μενέλαον R : \lecon{Μενέλεων} Scaliger}%
%\commentaire{901}{προδοῦσα Μενέλαον τὸν ἐμὸν ἐν Τροίᾳ πόσιν.}{\italien{Staccato e allegro} : comptez donc le nombre de résolutions de ce vers. Μενέλεων à la place de Μενέλαον -- avec synalèphe, synérèse, synizèze, ou tout ce que vous voudrez -- est la solution la plus simple pour rétablir le mètre.}
\\\passeLigne}%
\end{groupe}%
%% v. 902
\begin{groupe}%
\personnageGrec{Εὐριπίδης}%
\trimIamb{γύναι τί εἶπας; στρέψον ἀνταυγεῖς κόρας.\testimonia{902}{(στρέψον ἀνταυγεῖς κόρας) Suidas α 2617 \espTest (ἀνταυγεῖς κόρας) Photius α 2059 ἀνταυγεῖς κόρας}\\\passeLigne}%
\end{groupe}%
%% v. 903
\begin{groupe}%
\personnageGrec{Μνησίλοχος}%
\trimIamb{αἰσχύνομαί σε τὰς γνάθους ὑβρισμένη.
\commentaire{903}{αἰσχύνομαί σε τὰς γνάθους ὑβρισμένη.}{Rappelons que lors de la séance de travestissement, le pauvre et anciennement barbu Mnésiloque s'était fait flamber et raser la barbe. Sur un fameux cratère\marqueur{1} représentant la scène d'égorgement de la `petite' de Mika, on le voit en porter de très visibles stigmates. Et le scholiaste de commenter ce vers :%
\citComm{ἐπειδὴ ξυρηθεὶς ἦν,\\
puisqu'il a été rasé.}%
}%
\debutNote{903}%
\note{1}{Le « cratère de Würzburg », Martin-von-Wagner-Museum der Universitàt Wurzburg, H 5697. Reproduit en annexe ; cf. la note aux vers 754-757.}\\\passeLigne}%
\end{groupe}%
%% v. 904-905
\begin{groupe}%
\ledpb%
\personnageGrec{Εὐριπίδης}%
\trimIamb{τουτὶ τί ἔστιν; ἀφασία τίς τοί μ᾽ ἔχει.\testimonia{904}{\Antiatt 83.9 ; cf. Hesychius α 8572 ἀφασία}%
\commentaire{904-5}{}{\italique{Hélène}, v. 549 : %
« \difference{ἔκπληξιν ἡμῖν} ἀφασίαν \difference{τε προστίθης} » ; v. 557 : « τίς εἶ; τίν᾽ ὄψιν \difference{σήν}, γύναι, \difference{προσδέρκομαι}; » ; v. 72 : « ὦ θεοί, τίν᾽ \difference{εἶδον} ὄψιν; \difference{ἐχθίστην ὁρῶ} ».}\\
ὦ θεοὶ τίν᾽ ὄψιν εἰσορῶ; τίς εἶ γύναι;\\\passeLigne}%
\end{groupe}%
%% v. 906
\begin{groupe}%
\personnageGrec{Μνησίλοχος}%
\trimIamb{σὺ δ᾽ εἶ τίς; αὑτὸς γὰρ σὲ κἄμ᾽ ἔχει λόγος.\apparat{906}{αὐτὸς R : \lecon{αὑτὸς} Boissonade}%
\commentaire{906-12}{}{\italique{Hélène}, v. 558 sq. :%
\citComm{Ἑλένη --- σὺ δ᾽ εἶ τίς; αὑτὸς γὰρ σὲ κἄμ᾽ ἔχει λόγος.\\
\difference{Μενελέως --- οὐπώποτ᾽ εἶδον προσφερέστερον δέμας.\\
Ἑλένη --- ὦ θεοί· θεὸς γὰρ καὶ τὸ γιγνώσκειν φίλους.}\\
< Μενελέως --- Ἑλληνὶς εἶ τις ἢ ἐπιχωρία γυνή; > (Ce vers a été restitué à la tragédie d'après Aristophane.)\\
Ἑλένη --- Ἑλληνίς· ἀλλὰ καὶ τὸ σὸν θέλω μαθεῖν.\\
Μενελέως --- Ἑλένῃ σ᾽ ὁμοίαν δὴ μάλιστ᾽ εἶδον, γύναι.\\
Ἑλένη --- ἐγὼ δὲ Μενέλεῴ γε σέ· \difference{οὐδ᾽ ἔχω τί φῶ}.\\
Μενελέως --- ἔγνως γὰρ ὀρθῶς ἄνδρα δυστυχέστατον.\\
Ἑλένη --- ὦ χρόνιος ἐλθὼν σῆς δάμαρτος \difference{ἐς χέρας}.}}\\\passeLigne}%
\end{groupe}%
%% v. 907
\begin{groupe}%
\personnageGrec{Εὐριπίδης}%
\trimIamb{Ἑλληνὶς εἶ τις ἢ 'πιχωρία γυνή;\\\passeLigne}%
\end{groupe}%
%% v. 908
\begin{groupe}%
\personnageGrec{Μνησίλοχος}%
\trimIamb{Ἑλληνίς. ἀλλὰ καὶ τὸ σὸν θέλω μαθεῖν.\\\passeLigne}%
\end{groupe}%
%% v. 909
\begin{groupe}%
\personnageGrec{Εὐριπίδης}%
\trimIamb{Ἑλένῃ σ᾽ ὁμοίαν δὴ μάλιστ᾽ εἶδον γύναι.\apparat{909}{Ἐλένης R : \lecon{Ἑλένῃ σ᾽} Eur. \italique{Hel.} 563 \espApp μάλισθ᾽ ἴδον R : \lecon{μάλιστ᾽ εἶδον} \italique{Hel}.563}\\\passeLigne}%
\end{groupe}%
%% v. 910
\begin{groupe}%
\personnageGrec{Μνησίλοχος}%
\trimIamb{ἐγὼ δὲ Μενελέῳ σ᾽ ὅσα γ᾽ ἐκ τῶν λαιφίων.\testimonia{910}{Suidas ι 778 (sed Μενέλεων pro Μενέλαον et ἰφύων pro ἀφύων)}\apparat{910}{Μενέλαον R : Μενέλεων Suid. ι 778 : \lecon{Μενελέῳ σ᾽} Thiersch (cf. Eur. \italique{Hel}.564) \espApp \lecon{ὅσα γ᾽} R : ὄσον Suid. \espApp ἀφύων R : ἰφύων Schol., Suid. : ἀμφίων Grégoire : φυκίων Austin \& Olson : \lecon{λαιφίων} Bibi}%
\commentaire{910}{ἐγὼ δὲ Μενέλαον ὅσα γ᾽ ἐκ τῶν ἰφύων.}{Scène de reconnaissance. Lorsqu'elle avait reconnu son vieux mari face à elle, la véritable Hélène d'Euripide s'était exclamé (v. 564) :%
\citComm{Ἐγὼ δὲ Μενέλεῴ γε σέ· οὐδ' ἔχω τί φῶ.\\
Nul ne fut plus pareil à Ménélas. Que dire ?\\
(traduction d'Henri Grégoire.)}
Chez Aristophane, à quoi Hélène reconnaît-elle son naufragé de mari ? « Οὐδ' ἔχω τί φῶ » devient dans le manuscrit de Ravenne : ὅσα γ᾽ ἐκ τῶν ἀφύων, « d'après ton menu fretin ». 
On lit dans le lexique de la \italique{Souda} (ι 778) : ὄσον ἐκ τῶν ἰφύων, « à tes lavandes ». 
(C'est cette dernière version que commente le scholiaste, recopié en marge du manuscrit de Ravenne.)
Qu'il s'agisse de \italique{menu fretin} ou de \italique{lavande}, de telles leçons sont ineptes\marqueur{1} !
Colin Austin a bien résumé la situation\marqueur{2} :
« L'alpha d'ἀφύη étant bref, cette leçon est métriquement impossible. Quant à la ``lavande'' de la \italique{Souda}, le scholiaste y voit une allusion aux herbes de la mère d'Euripide, ce qui semble bien tiré par les cheveux ! Comme le note H. Grégoire\marqueur{3} « les herbes sont ici plus ridicules encore que les anchois », moins ridicules cependant que les \italique{oignons marins} de Gag***\marqueur{4} (τιφύων, « \italique{your squills}... \anglais{here a metaphor for male genitals} ») : cette saloperie-là\marqueur{5} se réfute d'elle-même. En fait la solution du problème se trouve au v. 935. Là Euripide est accusé d'être un ἱστιορράφος, \italique{un ravaudeur de voiles}. Fritzsche (p. 363)\marqueur{6} avait bien vu qu'Aristophane se moque ici de l'accoutrement de Ménélas dans l'\italique{Hélène} d'Euripide: tout au long de cette pièce il est fait mention des haillons, des lambeaux de voile dont le roi de Sparte est pitoyablement affublé, comme un gueux, après son naufrage en Égypte (v. 416 τὰς ἐμὰς δυσχλαινίας, 422 ναὸς ἐκβόλοις ἃ ἀμπίσχομαι, 544 sq. ἄγριος δέ τις / μορφὴν ὅδ᾽ ἐστίν, 554 στολήν γ᾽ ἄμορφον, 790 ὥσπερ πτωχὸς, 1204 ὡς ἐσθῆτι δυσμόρφῳ πρέπει, 1282 τῆς ἀχλαινίας, 1297 ἐσθῆτά τ᾽ ἐξάλλαξον). En particulier au ν. 423 il s'écrie: « Mes vêtements d'antan, mes habits somptueux, mes anciennes splendeurs, la vague m'a tout pris », πέπλους δὲ τοὺς πρὶν λαμπρά τ᾽ ἀμφιβλήματα / χλιδάς τε πόντος ἥρπασ᾽, et au ν. 1079 « les haillons qui m'affublent d'eux-mêmes parleront, en témoins du naufrage », καὶ μὴν τάδ᾽ ἀμφίβληστρα σώματος ῥάκη / ξυμμαρτυρήσει ναυτικῶν ἐρειπίων. Dans les \italique{Grenouilles} (ν. 842) Eschyle a bien raison d'appeler Euripide \italique{un faiseur de mendiants}, \italique{un ravaudeur de guenille}, καὶ πτωχοποιὲ καὶ ῥακιοσυρραπτάδη. Pour en revenir au v. 910 des « Thesmo », on est bien tenté de croire que c'est à ses \italique{hardes} que le pauvre naufragé est reconnu par sa femme. »\\ 
Grégoire\marqueur{7} proposa de lire : `ἐκ τῶν ἀμφίων' ; τὸ ἄμφιον étant une abréviation de ἀμφίεσμα, le \italique{vêtement}. Cette proposition expliquerait la présence des ἀφύων, par simple chute d'une lettre suivie d'iotacisme. En revanche, elle passe à l'as la leçon `ἰφύων' conservée par le scholiaste du manuscrit de Ravenne et par Suidas (ι 778), témoins tout aussi dignes de confiance.
(Re-)prends ton habit Grégoire ! Mais creusons son idée, qui est la bonne ; il nous suffit de trouver un synonyme de vêtement (avec, de préférence, une connotation péjorative) qui puisse rendre compte des deux leçons\marqueur{8}. 
Ce bon vieil Hésychius ne nous ayant jamais trahi, c'est encore à lui que nous avons posé la question. Cet affable vieillard n'a pas tardé à répondre à notre requête (λ 171) :%
\citComm{λαίφια· ῥάκη.\\
Laiphia : hardes.}%
Hésychius nous confirme l'existence d'un diminutif de λαῖφος, \italique{vêtement en lambeaux}, \italique{tissu}, \italique{voile}, ainsi que son sens.\\
Tentons à présent une explication paléographique aux ἀφύων et autres ἰφύων : si l'on fait la simple hypothèse de la disparition de la première lettre du mot λαιφίων\marqueur{9}, le lambda, le scribe se retrouvait alors devant l'inconnu αιφιων, qu'il lisait indifféremment – iotacisme oblige ! – αιφυων. En butte à ce mot qui n'existe pas, il pouvait alors raisonnablement écrire, selon l'inspiration du moment, ἀφύων (changeant d'un iota le texte) ou ἰφύων (en supprimant l'alpha).}%
\debutNote{910}%
\note{1}{Il se trouve néanmoins de nombreux défenseurs de la leçon « ἐκ τῶν ἰφύων », leçon qui fait beaucoup rire de savants hellénistes (?!) – décidément très bon public ! Nous avons expliqué, aux vers 455-456, les λάχανα ἄγρια, le σκάνδιξ et le verbe διασκανδικίζειν ; la présence de ces ἴφυα, elle, ne s'explique que par une mauvaise compréhension des quolibets des Comiques visant la mère d'Euripide.}%
\note{2}{«~Observations critiques sur les \italique{Thesmophories} d'Aristophane~», in \italique{Dodone} 19.2 (1990), p. 27.}%
\note{3}{Tome V de l'édition d'Euripide aux Belles Lettres (1950), p. 65 n. 1.}%
\note{4}{« Une telle saloperie ?.. ou plutôt un Gag, non ? » C'est ce mot d'ordre que je reprends à Austin, modifiant au passage le nom – \latin{nomen omen} ! – du farceur qui va recevoir sa fessée, conformément à ce que réclame la justice !}%
\note{5}{Pour réussir à faire perdre son sang-froid au pourtant toujours modéré Colin Austin, il fallait effectivement que ce fut une belle salauderie… Nous aurions volontiers laissé la misérable hypo-foutaise « qui se réfute d'elle-même » de Gag*** dans un oubli bien mérité, malheureusement, il nous a été infligé de la voir exhumée et prise au sérieux dans plusieurs articles récents… je crois donc de santé publique de l'autopsier une bonne fois pour toute et d'enterrer – définitivement, espérons le ! – cette charogne ! Ne nous attardons pas dans cette opération et allons droit au but : la saloperie de Gag*** présente à la fois le quadruple exploit : \ \ \ 1°) \ de ne pas rendre compte de la leçon `ἀφύων' du témoin principal de la pièce : le manuscrit de Ravenne ; \ \ \ 2°) \ d'élucubrer sur un soi-disant sens obscène de τίφυον qui n'est attesté nulle part ; \ \ \ 3°) \ d'être intrinsèquement absurde : si Ménélas est réputé avoir bien peu d'ardeur (raison du départ de l'insatisfaite madame ?) ; comment se fait-il qu'il soit en pleine \italique{érection} et que ce soit à cela qu'Hélène le reconnaisse ? \ \ \ 4°) \ \italique{last but not least}, d'être anti-comique à souhait : un cocu impuissant n'est drôle \italique{que} parce qu'il est impuissant : représentez-le en pleine érection et vous gâterez à coup sûr la farce ! \ \ \ \
Reconnaissons au moins à Gag*** le mérite (pédagogique pour l'édification du lecteur) d'avoir réussi à concentrer en un seul passage les quatre vices cardinaux du \italique{misthophore} de la recherche, savant homme qui se paye – littéralement ! – de mots. Qu'y a-t-il d'étonnant à ce qu'un individu salarié pour palabrer palabre ? Rien, à priori, mais on leur saurait tout-de-même gré d'occuper leur jactance à autre chose que de salir Aristophane.}%
\note{6}{Voir aussi G. Lange, \italique{Quaest. in Ar. Thesm., Diss. Gott}. (1891), p. 8.}%
\note{7}{Au tome V de son édition d'Euripide aux Belles Lettres, déjà mentionnée. Mais son hypothèse est de douze ans plus ancienne : il l'avait présentée en anglais dans l'article « Modern Greek in Aristophanes » paru dans le premier numéro de \italique{The Link, A Review of Mediaeval and Modern Greek} (1938), pp. 16-20. Une anecdote qui montre comme le monde (grec) est petit : Henri Grégoire n'est autre que celui qui reprit la suite du grand Alphonse Willems à la chaire de philologie grecque de l'Université libre de Bruxelles en 1909, lorsque ce dernier dut prendre sa retraite.}%
\note{8}{Il y avait une autre possibilité de recherche : trouver un mot dont les multiples sens auraient pu être glosés à la fois par ἴφυον et ἀφύη ; les gloses explicatives ayant fini par prendre la place de la leçon originale (le cas s'est déjà vu). Creusant dans cette voie, Austin \& Olson avaient suggéré une émendation très ingénieuse : φυκίων. Émendation ingénieuse sur le plan paléographique, beaucoup moins sur le plan comique…}%
\note{9}{Erreur courante dans notre manuscrit de Ravenne. Citons pour exemple le vers 834 où le manuscrit comporte « δίδοσθαι \textbf{Τ}ηνίοισι » à la place de « δίδοσθαι \textbf{Στ}ηνίοισι ». Cette erreur est d'autant plus fréquente lorsque il y a deux lettres identiques consécutives, ainsi au vers 555 du manuscrit de Ravenne : « ὡς \textbf{τ}λεγγίδας λαβοῦσαι » en place de « ὡς \textbf{στ}λεγγίδας λαβοῦσαι ». 
Une telle faute s'explique aisément si l'on imagine le manuscrit copié en onciales et où les mots n'étaient pas même espacés. Dans le mot « λαιφίων », les deux premières lettres, le `λ' et le `α' sont de formes semblables en onciales. La `chute' du `λ' me paraît donc une hypothèse loin d'être invraisemblable.}\\\passeLigne}%
\end{groupe}%
%% v. 911
\begin{groupe}%
\ledpb%
\personnageGrec{Εὐριπίδης}%
\trimIamb{ἔγνως ἄρ᾽ ὀρθῶς ἄνδρα δυστυχέστατον.\\\passeLigne}%
\end{groupe}%
%% v. 912-916a
\begin{groupe}%
\personnageGrec{Μνησίλοχος}%
\trimIamb{ὦ χρόνιος ἐλθὼν σῆς δάμαρτος ἐσχάρας\testimonia{912}{(ἐσχάρας) Photius ε 2044 ἐσχάρας $\subset$ Hesychius ε 6447 ἐσχάραι}%
\commentaire{912}{ἐλθὼν σῆς δάμαρτος ἐσχάρας / λαβέ.}{Le vers correspondant de l'\italique{Hélène} (le vers 566) finit par : ἐς χέρας ; Hélène reproche à Ménélas d'avoir trop tardé à rejoindre ses \italique{mains} (ses \italique{bras}, dira-ton en français). Occasion par trop belle pour Aristophane de renouer avec la tradition populaire qui ne voit en Hélène qu'une pandorienne dévergondée, à l'origine de tant de maux : il fait donc prononcer à Mnésiloque `ἐσχάρας'. Le mot ἐσχάρα désigne proprement le \italique{foyer} ou le \italique{réchaud}\marqueur{1}. Ici comme ailleurs\marqueur{2}, le terme peut désigner le \italique{mont de Vénus}. Jean Taillardat\marqueur{3} rapproche cette métaphore d'un passage de la \italique{Paix}\marqueur{4} où le mot ὀπτάνιον, \italique{cuisine} est pris comme synonyme de κύσθος.}
\debutNote{912}
\note{1}{\italique{Acharniens}, v. 888 ; \italique{Guêpes}, v. 938 ; \&c.}
\note{2}{\italique{Cavaliers} v. 1286 : κυκῶν τὰς ἐσχάρας avec le commentaire du scholiaste : τὰ χείλη τῶν γυναικείων αἰδοίων.}
\note{3}{\italique{Les images d'Aristophane}, §117.}
\note{4}{Vers 891.}}
\end{groupe}%
%% v. 913-915
\begin{groupe}%
\typeVers{2. dochm.}\diDoch{λαβέ με λαβέ με πόσι, \ligne{914}περίβαλε δὲ χέρας.\apparat{914}{περίβαλλε R : \lecon{περίβαλε} Biset}
\commentaire{913-5}{λαβέ με λαβέ με πόσι, / περίβαλε δὲ χέρας. / φέρε σὲ κύσω. ἄπαγέ μ᾽ ἄπαγ᾽ ἄπαγέ με.}{Les quatre dochmiaques sont tous intégralement résolus (et Lætitia\marqueur{1} note la «~\anglais{correption }κύσ\lettreBreve{ω}. ἄπαγε \anglais{harsh and unusual in dochmiacs}~») ; signe d'une accélération -- voire une précipitation en l'occurrence -- du rythme.\\
Aristophane nous présente Hélène pour ce qu'elle est : une nymphomane chronique. (Loin donc de la pseudo-réhabilitation d'Euripide, plus risible que crédible.)
Comment Hélène nous fait-elle entendre les besoins tyranniques qui la tiraillent ? par le même procédé que les deux esclaves au début des \italique{Cavaliers} (vers 26) : répéter un bout de phrase, en l'accélérant... et faire entendre autre chose que le motif initial. 
La luronne détachait-elle le `πάγε με' de la fin de son exhortation afin de suggérer le verbe πήγνυμι, \italique{ficher} ?\\
La source d'inspiration : \italique{Hélène}, v. 625 sq. :%u
\citComm{Ἑλένη --- ὦ φίλτατ᾽ ἀνδρῶν Μενέλεως, ὁ μὲν χρόνος\\
παλαιός, ἡ δὲ τέρψις ἀρτίως πάρα.\\
ἔλαβον ἀσμένα πόσιν ἐμόν, φίλαι,\\
περί τ᾽ ἐπέτασα χέρα\\
φίλιον ἐν μακρᾷ φλογὶ φαεσφόρῳ.\\
Μενελέως --- κἀγὼ σέ· πολλοὺς δ᾽ ἐν μέσῳ λόγους ἔχων\\
οὐκ οἶδ᾽ ὁποίου πρῶτον ἄρξωμαι τὰ νῦν.\\
Ἑλένη --- γέγηθα, κρατὶ δ᾽ ὀρθίους ἐθείρας\\
ἀνεπτέρωκα καὶ δάκρυ σταλάσσω,\\
περὶ δὲ γυῖα χέρας ἔβαλον, ἡδονάν,\\
ὦ πόσις, ὡς λάβω.}%
}%
\debutNote{913-5}%
\note{1}{\italique{The Songs of Aristophanes}, p. 429.}\\
φέρε σὲ κύσω. ἄπαγέ μ᾽ ἄπαγ᾽ ἄπαγέ με%
%\commentaire{914}{ἄπαγέ μ᾽ ἄπαγ᾽ ἄπαγέ με.}{Il n'est pas impossible qu'Aristophane recoure au même procédé }%
}
\end{groupe}%
%% v. 912-916a
\begin{groupe}%
\typeVers{trim. iamb.}\trimIamb{λαβὼν ταχὺ πάνυ.\\\passeLigne}%
\end{groupe}%
%% v. 916b-917
\begin{groupe}%
\personnageGrec{Κρίτυλλα}%
\trimIamb{\ligne{916}\decalage{λαβὼν ταχὺ πάνυ.}κλαύσετ᾽ ἄρα νὴ τὼ θεὼ\\
ὅστις σ᾽ ἀπάξει τυπτόμενος τῇ λαμπάδι.%\commentaire{917}{τῇ λαμπάδι.}{ \eff{Scholie} : didascalie d'Aristophane ?}
\\\passeLigne}%
\end{groupe}%
%% v. 918-919
\begin{groupe}%
\personnageGrec{Εὐριπίδης}%
\trimIamb{σὺ τὴν ἐμὴν γυναῖκα κωλύεις ἐμέ,\\
τὴν Τυνδάρειον παῖδ᾽, ἐπὶ Σπάρτην ἄγειν;\\\passeLigne}%
\end{groupe}%
%% v. 920-923
\begin{groupe}%
\ledpb%
\personnageGrec{Κρίτυλλα}%
\trimIamb{οἴμ᾽ ὡς πανοῦργος καὐτὸς εἶναί μοι δοκεῖς\\
καὶ τοῦδέ τις ξύμβουλος. οὐκ ἐτὸς πάλαι\\
ᾐγυπτιάζετ᾽. ἀλλ᾽ ὅδε μὲν δώσει δίκην.\testimonia{922}{(ᾐγυπτιάζετ᾽) \SynB α 494 αἰγυπτιάζειν $=$ Photius α 518 $\subset$ Suidas αι 75}%
\commentaire{922}{ᾐγυπτιάζετ᾽.}{Αἰγυπτιάζειν, faire l'Égyptien, comme Σιφνιάζειν, Χιάζειν, \italique{faire le Chiote}, Ἑλληνίζειν, \italique{faire l'Hellène}, Λεσβιάζειν, Βακίζειν, \italique{faire le Bakis} (\italique{Paix}, 1072). Les Égyptiens avaient la réputation d'être très rusés. On a conservé un vers d'Eschyle\marqueur{1} : %
\citComm{δεινοί πλέκειν τοι μηχανὰς Αἰγύπτιοι\\
Les Égyptiens habiles à ourdir des trames.}%
Cf. Properce (III, 9, 33) : \latin{Noxia Alexandria, dolis aptissima lellus.} (Note de Willems ; cf. aussi Taillardat §409.)}%
\debutNote{922}%
\note{1}{Tiré d'une scholie à Théocrite (15, 48b). C'est le fragment 373 de l'édition de Radt.}\\
προσέρχεται γὰρ ὁ πρύτανις χὠ τοξότης.\\\passeLigne}%
\end{groupe}%
%% v. 924
\begin{groupe}%
\personnageGrec{Εὐριπίδης}%
\trimIamb{τουτὶ πονηρόν· ἀλλ᾽ ὑπαποκινητέον.\\\passeLigne}%
\end{groupe}%
%% v. 925a
\begin{groupe}%
\personnageGrec{Μνησίλοχος}%
\trimIamb{ἐγὼ δ᾽ ὁ κακοδαίμων τί δρῶ;\\\passeLigne}%
\end{groupe}%
%% v. 925b-927
\begin{groupe}%
\personnageGrec{Εὐριπίδης}%
\trimIamb{\ligne{925}\decalage{ἐγὼ δ᾽ ὁ κακοδαίμων τί δρῶ;}μέν᾽ ἥσυχος.\\
οὐ γὰρ προδώσω σ᾽ οὐδέποτ᾽, ἤνπερ ἐμπνέω,\testimonia{926}{(ἤνπερ ἐμπνέω) Suidas η 379 (ἤν pro ἤνπερ)}\\
ἢν μὴ προλίπωσ᾽ αἱ μυρίαι με μηχαναί.\apparat{927}{προλείπωσ᾽ R : \lecon{προλίπωσ᾽} Scaliger}\\\passeLigne}%
\end{groupe}%
%% v. 928
\begin{groupe}%
\personnageGrec{Μνησίλοχος}%
\trimIamb{\lednopb%
αὕτη μὲν ἡ μήρινθος οὐδὲν ἔσπασεν.\testimonia{928}{Suidas α 4475, μ 978 ; \Eustathe 760.44 ; scholium\expo{RV} \italique{Vespae} 175b}\apparat{928}{(Γυ.) R : \lecon{(Μνη.)} Scaliger}
\commentaire{928}{αὕτη μὲν ἡ μήρινθος οὐδὲν ἔσπασεν.}{Le tour est proverbial et l'expression nous est bien connue grâce à une scholie au vers 175 des \italique{Guêpes} où Aristophane joue avec l'halieutique proverbe :%
\citComm{παρὰ τὴν παροιμίαν « αὕτη μὲν ἡ μήρινθος οὐδὲν ἔσπασεν ». ἐλέγετο δὲ ἐπὶ τῶν πείρᾳ καθιέντων ἄγκιστρον μετὰ σχοίνου, καὶ μὴ ἐπιτυγχανόντων ἰχθύος τινός.\\
\francais{En regard du proverbe : « voici une ligne qui n'a rien tiré hors de l'eau ». On le dit à propos de ceux qui jettent un appât à la pointe acérée (un hameçon avec un roseau), mais qui ne rencontrent pas un poisson.}}%
Ce proverbe est dit d'une ruse qui échoue\marqueur{1}. En français : un échec sur toute la ligne.}
\debutNote{928}
\note{1}{Pour un exposé plus développé sur ces deux passages avec à la clef l'explication du jargon halieutique, le lecteur pourra aller consulter le §395 des \italique{Images d'Aristophane} de J. Taillardat. Érasme a recensé l'adage (n° 845).}\\\passeLigne}%
\end{groupe}%
%% v. 929-934a
\begin{groupe}%
\ledpb%
\personnageGrec{Πρύτανις}%
\trimIamb{ὅδ᾽ ἔσθ᾽ ὁ πανοῦργος ὃν ἔλεγ᾽ ἡμῖν Κλεισθένης;\apparat{929}{ἔστ᾽ R : \lecon{ἔσθ᾽} \Mon \espApp ἔλεγεν R : \lecon{ἔλεγ᾽} \Mon}\\
οὗτος τί κύπτεις; δῆσον αὐτὸν εἰσάγων%\commentaire{}{}{\eff{ἔκυψεν ὥσπερ Τήλεφος} cf Taillardar 349}
\\
ὦ τοξότ᾽ ἐν τῇ σανίδι, κἄπειτ᾽ ἐνθαδὶ\testimonia{931}{(σανίδι) Photius σ 68 σανίς, σ 65 σανίδα ; Hesychius σ 169 σανίς ; \AelDion σ 7}\\
στήσας φύλαττε καὶ προσιέναι μηδένα\\
ἔα πρὸς αὐτόν, ἀλλὰ τὴν μάστιγ᾽ ἔχων\\
παῖ᾽ ἢν προσίῃ τις.\\\passeLigne}%
\end{groupe}%
%% v. 934b-935
\begin{groupe}%
\personnageGrec{Κρίτυλλα}%
\trimIamb{\ligne{934}\decalage{παῖ᾽ ἢν προσίῃ τις.}νὴ Δί᾽ ὡς νῦν δή γ᾽ ἀνὴρ\apparat{934}{νῦν δῆτ' R : \lecon{νῦν δή γ᾽} Dobree : νυνδή γ' Cobet}\\
ὀλίγου μ᾽ ἀφείλετ᾽ αὐτὸν ἱστιορράφος.\testimonia{935}{Suidas ι 691 ὀλίγου ταύτην ἱστιορράφος \espTest (ἱστιορράφος) Pollux VII 161}%
\commentaire{935}{ἱστιορράφος.}{Allusion aux hardes dont était couvert le prétendu Ménélas. (Note de Willems ; ἱστιορράφος est mis παρὰ προσδοκίαν pour μηχανορράφος\marqueur{1}.)}%
\debutNote{935}%
\note{1}{Cf. Euripide, \italique{Andromaque}, vv. 447, 1116 ; Sophocle, \italique{Œdipe Roi}, v. 387.}\\\passeLigne}%
\end{groupe}%
%% v. 936-938
\begin{groupe}%
\personnageGrec{Μνησίλοχος}%
\trimIamb{ὦ πρύτανι πρὸς τῆς δεξιᾶς, ἥνπερ φιλεῖς\\
κοίλην προτείνειν ἀργύριον ἤν τις διδῷ,%
\commentaire{936-7}{ὦ πρύτανι πρὸς τῆς δεξιᾶς, ἥνπερ φιλεῖς / κοίλην προτείνειν ἀργύριον ἤν τις διδῷ.}{%\eff{Schol. + Hésychius : c'était l'usage d'attacher nu au sanis ?}
Précis à l'usage du justiciable confronté à un prytane.
Les deux répliques que Mnésiloque adresse au prytane nous donnent un bon aperçu de l'estime en laquelle Aristophane tenait ces magistrats -- car nous allons voir que Mnésiloque tente de le soudoyer, comme cela semble l'usage à Athènes.
On trouve le même reproche de concussion dans la \italique{Paix} (vers 905-908) :%
\citComm{Ἀλλ᾽, ὦ πρυτάνεις, δέχεσθε τὴν Θεωρίαν.\\
Θέασ᾽ ὡς προθύμως ὁ πρύτανις παρεδέξατο.\\
Ἀλλ᾽ οὐκ ἄν, εἴ τι προῖκα προσαγαγεῖν σ᾽ ἔδει,\\
ἀλλ᾽ ηὗρον ἄν σ᾽ ὑπέχοντα τὴν ἐκεχειρίαν.\\
Allons, prytanes, recevez Théoria. Voyez l'accueil empressé que lui fait ce prytane. Quelle différence s'il t'avait fallu introduire une affaire gratis. Je t'aurais vu te retrancher sur les vacances.}%
(ὑπέχοντα τὴν ἐκεχειρίαν, «~alléguant la trêve~», est la version bégayée de ce qu'il fallait vraiment entendre : ὑπέχοντα τὴν χεῖρα, «~tendant la main~».)\\
En effet, pour soumettre une requête à la Boulé l'on devait préalablement obtenir l'aval des prytanes ; le cas échéant, les bouleutes discutaient et rédigeaient une proposition de loi (προβούλευμα) soumise ensuite à l'Ecclésia. L'auteur de la \italique{République des Athéniens} (III, 1) constate qu'il fallait parfois patienter jusqu'à une année avant que les prytanes ne donnent leur accord.
De même, pour intenter une action en justice, un citoyen devait régler des \italique{frais de justice}, τά πρυτανεῖα\marqueur{1}, dont le nom indique clairement à qui cette somme devait être remise.\\ 
Par ses fonctions, le prytane \italique{s'engageait} à servir d'intermédiaire : ce genre de serment passait très certainement par une \italique{poignée de main}\marqueur{2}.
C'est le sens de «~ἥνπερ φιλεῖς / κοίλην προτείνειν~», «~[cette main] que tu as l'habitude de tendre ouverte~».\\
Mais toutes ces transactions étaient autant d'occasions pour les prytanes de détourner à leur profit une partie du bien public.
C'est entendu : le prytane est concussionnaire mais il serait suicidaire pour Mnésiloque de le lui dire !%
\citComm{Flattez-les, payez-les d'agréables mensonges,\\
Quelque indignation dont leur cœur soit rempli,\\
Ils goberont l'appât, vous serez leur ami.}%
Au contraire : Mnésiloque tente de mettre le prytane dans sa poche -- il y parvient d'ailleurs au vu de la bienveillante réponse qu'il reçoit : τί σοι χαρίσωμαι;\\
Pour cela, jeu de scène\marqueur{3} : tout en déclamant «~πρὸς τῆς δεξιᾶς κτλ.~», Mnésiloque, qui n'est pas encore lié, tend sa main vers le prytane et s'offre par ce geste en suppliant\marqueur{4}. 
(Et d'ajouter -- en aparté, évidemment ! mais ce que le prytane a bien compris et qui explique ses bonnes dispositions à l'égard de Mnésiloque -- : «~ἀργύριον ἤν τις διδῷ~», «~[cette main que tu aimes à tendre] lorsque l'on t'offre de l'argent~».)\\
Le prytane accordera bien cette petite faveur à Mnésiloque (χάρισαι βραχύ τί μοι), à lui qui est sur le point de mourir ; et Mnésiloque d'insister sur ce dernier point (καίπερ ἀποθανουμένῳ).
Aristophane joue avec une locution proverbiale, selon sa pente maintenant bien établie par notre commentaire : les Athéniens, même à l'article de la mort, étaient réputés continuer de quémander la main tendue\marqueur{5} :%
\citComm{Ἀττικὸς ὑπέχει τὴν χεῖρα ἀποθνήσκων· ἐπὶ τῶν φιλαργύρων. φιλοκερδεῖς γὰρ οἱ Ἀθηναῖοι.}%
Le \italique{verre du condamné} que demande Mnésiloque (καίπερ ἀποθανουμένῳ) n'est autre que cette faveur mendiée la main tendue vers le prytane.}%
\debutNote{936-7}%
\note{1}{Voir à ce sujet le Daremberg \& Saglio à l'entrée \latin{Prytaneia}. Ces informations aident à interpréter la glose d'Hésychius (π 1511) : %
\citComm{πεπρυτανευκώς· ἐπιμεμηνιευκώς. καὶ προανηλωκώς,}%
rangée parmi les fragments comiques adespotes (408) dans l'édition Kassel-Austin.}%
\note{2}{On pourra lire l'article de Francesco Mari : «~Les sens de la poignée de main en Grèce ancienne du VIIIe au Ve siècle avant J.-C.~», in \italique{Ktèma}, 43 (2018), pp. 105-131.}%
\note{3}{Il faut rendre aux jeux de scènes toute l'importance qu'ils avaient à l'époque ; empêtrés dans leurs costumes et d'imposants masques de théâtre, il est à parier que les acteurs compensaient cette gêne par des jeux de mains des plus expressifs.}%
\note{4}{(La \italique{droite} en question est bien celle de Mnésiloque et non du prytane.) Pour cette signification de la main tendue, cf. \LSJ \ à l'entrée προτείνω II.}%
\note{5}{Recensé par Diogénianus (III, 12), \italique{Corpus paroemiographorum Graecorum}, I, p. 216. C'est l'adage n° 1646 d'Érasme. À noter qu'Aristophane y fait aussi une possible allusion aux vers 170 et suivants des \italique{Grenouilles}, lorsque Dionysos hèle un mort pour qu'il porte ses bagages, ce dernier lui répond : Δύο δραχμὰς μισθὸν τελεῖς; Voir plus haut le commentaire du vers 738.}\\
χάρισαι βραχύ τί μοι καίπερ ἀποθανουμένῳ.\\\passeLigne}%
\end{groupe}%
%% v. 939a
\begin{groupe}%
\personnageGrec{Πρύτανις}%
\trimIamb{τί σοι χαρίσωμαι;\apparat{939}{χαρίσομαι R : \lecon{χαρίσωμαι} Casaubon}\\\passeLigne}%
\end{groupe}%
%% v. 939b-942
\begin{groupe}%
\personnageGrec{Μνησίλοχος}%
\trimIamb{\ligne{939}\decalage{τί σοι χαρίσωμαι;}γυμνὸν ἀποδύσαντά με,\\
κέλευε πρὸς τῇ σανίδι δεῖν τὸν τοξότην,\testimonia{940}{(πρὸς τῇ σανίδι δεῖν) Hesychius δ 926 δῆσαι πρὸς σανίδι· Ἀττικῶς· βασανίσαι ἐν τῷ δεσμωτηρίῳ}\\
ἵνα μὴ 'ν κροκωτοῖς καὶ μίτραις γέρων ἀνὴρ\apparat{941}{μὴν R : \lecon{μὴ 'ν} Küster}\\
γέλωτα παρέχω τοῖς κόραξιν ἑστιῶν.\testimonia{942}{Suidas ε 3215 \espTest (ἑστιῶν) Hesychius ε 6411 ἑστιῶν· ἀριστοποιῶν $=$ Photius ε 2031 $=$ \Syn ε 875 $=$ Suidas ε 3227 ; cf. \Phryn 66.7}\apparat{942}{\inRas{γε}{} τοῖς \Rdeux}% \espApp ἑστιῶν R : \lecon{ἡστίων} Bibi}%
\commentaire{939-42}{γυμνὸν ἀποδύσαντά με / κέλευε πρὸς τῇ σανίδι δεῖν τὸν τοξότην, / ἵνα μὴ 'ν κροκωτοῖς καὶ μίτραις γέρων ἀνὴρ / γέλωτα παρέχω τοῖς κόραξιν ἑστιῶν.}{%
%1004 γυμνῷ φυλακὴν ἐπιτάττειν· παροιμία ἐπὶ τῶν <μὴ> δεομένων προστάξεως, διὰ τὸ ἀναγκαίως καὶ χωρὶς ἐπικελεύσεως τοῦτο πράττειν. Φερεκράτης Τυραννίδι (fr.154 Κ.Ἁ.)
%Suidas : Γυμνῷ φυλακὴν ἐπιτάττεις: ἐπὶ τῶν ἀδυνάτως ἐχόντων τὰ προστεταγμένα πληροῦν, διὰ τὰ ἀναγκαῖα. ὅτι τὸ Γυμνῶ μετὰ γενικῆς: μετὰ αἰτιατικῆς δὲ, ὡς τὸ, γυμνώσαντες αὐτὸν ἐξήγαγον ἔξω. καὶ, γυμνῶν τὰ ἐσφραγισμένα.
%
%\marqueur{Pour se rendre bien compte de ce que pouvait être cette omniprésence envahissante -- et salissante ! -- des oiseaux lors des sacrifices ou, comme c'est notre cas, d'exécution à ciel ouvert, le lecteur pourra se reporter au Daremberg \& Saglio, à l'entrée \latin{Meniskos} (Μηνίσκος).}
%\eff{Hésychius δῆσαι πρὸς σανίδι:᾿Αττικῶς:βασανίσαιἐν τῷ δεσμωτηρίῳ (Αr. Τhesm. 940 ?) en fait Cratinus cf. scholiaste de notre passage  ἡλωτός: cloué }
%	pros tê sanidi : face à la planche, pour cacher sa nudité... mais voir aussi gerron Ael. Dion.
%\eff{fr. 639 K.-A. : ἐπὶ τῶι ταρίχει τὸν γέλωτα κατέδομαι.\\
%fr. 803 γέρρον ...  ̓Αριοτοφάνης ὧς ἐπὶ φυλακῆς τινος καὶ μοχλοῦ τροπικῶς παρέλαβε τὴν λέξιν}
Ces quatre vers sont peut-être parmi les plus mécompris de toute la comédie. 
Les commentateurs avaient effectivement bien peu de chance de les comprendre puisqu'ils commencent par un contre-sens : ils construisent κέλευε τὸν τοξότην γυμνὸν ἀποδύσαντά με, «~ordonne à l'archer de me déshabiller tout nu~»...
Comment Mnésiloque pourrait-il faire cette requête alors que Clisthène l'a déjà mis à nu\marqueur{1} ? En réalité, le groupe `γυμνὸν ἀποδύσαντά με' est un accusatif absolu et se traduit par : «~puisque j'ai été mis à nu, ordonne à l'archer \&c.~»\\
Mnésiloque pratique avec un art consommé le double langage : toute parole qu'il adresse au prytane est à double sens : un clair, mais où il est dit que le sens véritable est caché. 
Que dit en apparence Mnésiloque ? il fait profession honorable de pudeur ; ce pudibond de la onzième heure tient à cacher ce qui doit être caché. 
Voilà pourquoi il demande à être retourné \italique{contre} (πρός) \italique{la planche} du supplice (τῇ σανίδι)\marqueur{2}. %\marqueur{Le scholiaste ne disait pas autre chose NON : `ἐπίτρεψον οὖν τῷ τοξότῃ', φησὶ, `γυμνόν με δῆσαι πρὸς τῇ σανίδι'.}.
À qui tient-il à cacher son anatomie ? aux femmes évidemment -- et non aux corbeaux\marqueur{3} !.. Au vers 942, les `κροκωτοῖς καὶ μίτραις' sont les \italique{crocotes} et \italique{mitres} des femmes réunies, rappelons-le, en assemblée dans le but de perdre Euripide.\marqueur{4}. Dans ce contexte, ἐν + datif a le sens qu'il a toujours\marqueur{6} : \italique{au milieu de}, \italique{parmi}, \italique{devant} : «~ἐν κροκωτοῖς καὶ μίτραις~» $=$ «~au milieu des (en)crocotées et des (en)mitrées~».
Voici pour le sens clair.\\
Mais que demande réellement Mnésiloque au prytane ?
L'expression `δῆσαι πρὸς σανίδι' a heureusement été relevée par Hésychius (δ 926) -- qui se montre encore une fois beaucoup plus pertinent que notre scholiaste :%
\citComm{δῆσαι πρὸς σανίδι· Ἀττικῶς· βασανίσαι ἐν τῷ δεσμωτηρίῳ.\\
Lier à la planche : Atticisme. Interroger dans la prison.}%
Mnésiloque veut un procès équitable (il sent que la justice d'exception risque de lui être par trop expéditive) : il est préférable pour lui de s'en remettre aux masculines institutions.
La σανίς de l'expression désigne sans doute la \italique{planche} sur laquelle était écrite inscrite l'ordre du jour des tribunaux.
Cf. en \italique{Guêpes} (v. 848) et Photius (σ 68) : %
\citComm{σανίς· καὶ ἐν ᾗ τοὺς κακούργους ἔδουν καὶ ἐν ᾗ τὰς δίκας ἐνεχάραττον· τὸ µὲν ἐν Θεσµοφοριαζούσαις, τὸ δὲ ἐν Σφηξὶ (848) δεδήλωται.}%
%L'on reste consterné devant l'interprétation des professionnels du commentaire pour ce passage : Mnésiloque se proposerait de graisser la patte du prytane – qui semble avoir quelque inclination naturelle pour la chose – en avançant la raison officielle qu'il ne veut pas prêter à rire aux corbeaux (sic !?!)... \\
%Soufflé, bouche bée, cloué au bec par une telle exégèse !
%Il y a effectivement de quoi rire, mais ce n'est pas les corbeaux qui le feront.
%Et d'abord, leur fantaisie ne tient même pas du point de vue du grec. Nous ne sommes dans un pastiche de style tragique que je sache ; le pluriel ἐν κροκωτοῖς καὶ μίτραις n'aurait donc rien à faire ici s'il désignait \italique{la}\marqueur{1} crocote ou \italique{la} mitre de Mnésiloque.
%Et si ces crocotes et mitres ne sont pas celles de Mnésiloque ; ce sont donc celles des femmes réunies, rappelons-le, en assemblée dans le but de perdre Euripide. Dans ce contexte, ἐν + datif a le sens qu'il a toujours : \italique{au milieu de}, \italique{parmi}, \italique{devant} : "ἐν κροκωτοῖς καὶ μίτραις" $=$ "au milieu des (en)crocotées et des (en)mitrées".\\
%Ce n'est donc évidemment pas devant les corbeaux mais bien devant les \italique{femmes} que Mnésiloque dit craindre d'être ridicule.\\
%Il n'y a plus qu'à rétablir un augment oublié à la fin du vers 942\marqueur{2} : ἡστίων (imparfait) en place de l'incompréhensible ἑστιῶν et voilà notre passage parfaitement intelligible : %
%\citComm{γυμνὸν ἀποδύσαντά με\\
%κέλευε πρὸς τῇ σανίδι δεῖν τὸν τοξότην,\\
%ἵνα μὴ 'ν κροκωτοῖς καὶ μίτραις γέρων ἀνὴρ\\
%γέλωτα παρέχω τοῖς κόραξιν ἡστίων.\\
%Ordonne à l'archer que je sois déshabillé avant que d'être attaché à la planche afin qu'au milieu des porteuses de crocotes et de mitres je ne prête pas à rire avec les corbeaux qui dîneraient.}%
%Figure porte absence et présence, dit Pascal ; l'on soupçonne un double sens : un clair, que nous venons d'établir, mais où il est dit que le sens est caché. 
%Puisqu'il ne semble pas qu'il y ait d'allusion proverbiale, littéraire ou à la chronique du temps, reste le jeu de mots (cas le plus favorable puisque nous possédons tout le matériel pour le saisir).\\
%Ne perdons pas de vue l'objectif de Mnésiloque : s'enfuir ; or par analogie κόραξ désigne tout instrument de forme recourbée comme un bec de corbeau et peut être synonyme de κατακλείς : \italique{cliquet}, \italique{crochet} ; cette fois, nous tenons le bon bout ! Mnésiloque demande bien-sûr, en loucédé, à être libéré de ses entraves.\\
%Poursuivons. Le verbe παρέχω, \italique{fournir}, ne diffère que bien peu du verbe παρα-χόω, \italique{verser sur le côté}. 
%En faisant quelques efforts de désarticulation, le protagoniste pouvait certainement faire entendre à la fois παρέχω, \italique{que je prête} [à rire] et παραχοῖ\marqueur{3}, \italique{qu'il verse de côté} [mes crochets].\\
%Le double sens se précise ; il demeure cependant le dernier mot, ἑστιῶν. 
%Notre correction en ἡστίων nous facilite la tâche : il suffit de faire la liaison avec τοῖς κόραξιν, la prononciation devait être très proche de celle de τοῖς κόραξι νήστιων, \italique{les crochets [au niveau] de mes viscères}. 
%(Pour cette acception de νῆστις, Photius la donne à l'entrée νῆστιν en citant Aristophane lui-même : νῆστιν· τὸ ἔντερον· `ἢ σπλῆνα ἢ γλῶτταν'. Ἀριστοφάνης\marqueur{4}.)\\
%Sans forcer le texte, ce nous semble, voici donc une lecture alternative :%
%\citComm{παρέχω τοῖς κόραξιν ἡστίων\\
%παραχοῖ τοῖς κόραξι νήστιων\\
%Qu'il [l'archer] me débarrasse  des crochets [au niveau] de mes reins.}%
%Admirez, s'il vous plaît, ce petit bijou d'ambiguïté qui ne peut être comparé, chez nous, qu'aux perles du génial Raymond Devos.
}%
\debutNote{939-42}%
\note{1}{Cf. vers 637 sqq. et les notes afférentes ; Mnésiloque n'a plus sur lui que sa crocote tout-à-fait transparente.}%
\note{2}{L'expression grecque correcte pour dire \italique{lier à la planche} est ἐν σανίδι et non πρός σανίδι, cf. vers 934 et 1165.}%
\note{3}{Je prie le lecteur d'excuser la présence de remarques aussi évidentes mais il faut bien rectifier la clique des commentateurs capable d'écrire à l'unisson et sans sourciller ce genre d'âneries.}%
\note{4}{Et non les habits de Mnésiloque comme le prétend l'internationale du commentaire !.. d'abord, on les lui a retirés ; ensuite cette interprétation fantaisiste ne tient même pas du point de vue du grec : nous ne sommes dans un pastiche de style tragique que je sache ; le pluriel ἐν κροκωτοῖς καὶ μίτραις n'aurait donc rien à faire ici s'il désignait \italique{la}\marqueur{5} crocote ou \italique{la} mitre de Mnésiloque.}%
\note{5}{Ou \italique{le} si l'on voulait rester conforme au genre du mot grec ; mais l'oreille française interdit toute féminisation de «~crocote~».}%
\note{6}{Cf. les vers 90-91 et la note qui s'y rapporte.}%
%\note{2}{Même oubli que, par exemples, aux vers 488, 521 ou 794.}%
%\note{3}{D'autant que les règles classiques de contraction auraient dû fournir la forme παραχῷ.}%
%\note{4}{Photius, ν 201. Cf. Athénée III, 96c $=$ fr. 520 Kassel-Austin.}
\\\passeLigne}%
\end{groupe}%
%% v. 943-944
\begin{groupe}%
\personnageGrec{Πρύτανις}%
\trimIamb{ἔχοντα ταῦτ᾽ ἔδοξε τῇ βουλῇ σε δεῖν,\apparat{943}{ἔδοξεν R : \lecon{ἔδοξε} Biset}%
\commentaire{943-56.}{Le papyrus d'Oxyrhynque n° 3839\marqueur{1},}{dont il a été déjà question aux vers 742-766 contient contient également un fragment de ce passage. Il est si gâté qu'il ne nous apporte aucune information (si ce n'est une colométrie que Lætitia suivit \anglais{without strong conviction}\marqueur{2}). Le lecteur le trouvera néanmoins reproduit en annexe ; gage à lui d'en tirer quelque chose.}%
\debutNote{943-56}%
\note{1}{\italique{The Oxyrhynchus Papyri}, vol. LVI (1989), n°3839, fr. 2. Analysé par W. E. H. Coockle avec le concours de Colin Austin.}%
\note{2}{L. P. E. Parker, \italique{The Songs of Aristophanes} (1997), p. 430.}\\
ἵνα τοῖς παροῦσι δῆλος ᾖς πανοῦργος ὤν.\apparat{944}{\lecon{παροῦσι} R : παριοῦσι Brunck}\\\passeLigne}%
\end{groupe}%
%% v. 945-946
\begin{groupe}%
\personnageGrec{Μνησίλοχος}%
\trimIamb{ἰαππαπαιάξ· ὦ κροκώθ᾽ οἷ᾽ εἴργασαι·\testimonia{945}{(ἰαππαπαιάξ) Suidas ε 2807}\apparat{945}{\lecon{ἰαππαπαιάξ} R, Suid. ε 2807 : ἰατταταιάξ Biset \ligneApp{946} ἔστιν R : \lecon{ἔστ᾽} Bentley}%
\commentaire{945}{ὦ κροκώθ᾽ οἷ᾽ εἴργασαι·}{Y a-t-il un jeu de mots entre κρόκαλον et κροκωτός ? 
Cf. Hésychius (κ 4170 et 4270) :%
\citComm{κρόκαλον· τὸ πανοῦργον παιδίον.}%
}\\
κοὐκ ἔστ᾽ ἔτ᾽ ἐλπὶς οὐδεμία σωτηρίας.%\apparat{946}{ἔστιν R : \lecon{ἔστ᾽} Bentley}
\\\passeLigne}%
\end{groupe}%
%% v. 947-952
\begin{groupe}%
\ledpb%
\personnageGrec{Κορυφαῖος}%
\typeVers{\decalageN{4}$4$ an\cat}\tetramAnap{ἄγε νυν ἡμεῖς παίσωμεν ἅπερ νόμος ἐνθάδε ταῖσι γυναιξίν,\apparat{947}{πέσωμεν R : \lecon{παίσωμεν} Zanetti, cf. Schol. \espApp ταῖσ R : \lecon{ταῖσι} Daubuz \ligneApp{948} \lecon{θεαῖν} R : θεοῖν Cobet}%
\commentaire{947-1000}{}{Pour une analyse plus fine du chant mêlé de danses de tout ce passage, je renvoie à l'article d'Ester Cerbo\marqueur{1}.
Quant à sa structure, celle d'Enger, assez communément admise -- et notamment par Lætitia\marqueur{2} --, me semble juste  :%
\citComm{--- anapestes (v. 947-952) ;\\
--- proode (v. 953-958) ;\\
--- première paire strophique (v. 959-962 $=$ v. 963-965) ;\\
--- première mésode (v. 966-968) ;\\
--- seconde paire strophique (v. 969-976 $=$ v. 977-984) ;\\
--- seconde mésode (v. 985-989) ;\\
--- tierce paire strophique (v. 990-994 $=$ v. 995-1000).}%
}%
\debutNote{947-1000}%
\note{1}{«~Uno spettacolo di danza alla festa delle donne: a proposito di Ar. \italique{Thesm.} 947-1000~» (2021).}%
\note{2}{Dans son indispensable \italique{The Songs of Aristophanes}, pp. 428-436.}\\
παραχώρει ὅταν ὄργια σεμνὰ θεαῖν ἱεραῖς ὥραις ἀνέχωμεν, ἅπερ καὶ\testimonia{947}{(ὅταν -- ἀνέχωμεν) Suidas α 2364 (θεαῖν codd. GFS\expo{a.c.} vel θεῶν codd. ATS\expo{p.c} vel θεεῖν cod. I vel θεαῖνων cod. M pro θεοῖν)}}%
\end{groupe}%
\begin{groupe}%
\typeVers{dim. anap.}\dimAnap{Παύσων σέβεται καὶ νηστεύει,%
\commentaire{949-52}{Παύσων σέβεται καὶ νηστεύει κτλ.}{Pauson était connu pour son indigence. Les parémiographes nous ont rapporté à son sujet deux proverbes\marqueur{1} :%
\citComm{Παύσωνος πτωχότερος· οὗτος ζωγράφος ἦν καὶ ἐπὶ πενίᾳ διατεθρύλλητο.\\
Παύσων καὶ Ἶρος: ὀνόματα κύρια πενήτων. ὁ Ἀσκληπιὸς Παύσωνα καὶ Ἶρον κἂν ἄλλον τινὰ τῶν ἀπόρων ἰάσαιτο. κτλ.\\
(Iros, le mendiant odieux d'Homère, \italique{Odyssée} XVIII.)}%
L'on trouvera tous les témoignages sur Pauson d'Athènes dans le \italique{Recueil Milliet} d'Adolphe Reinach\marqueur{2}. Pour ce qui est d'Aristophane, la scholie à notre passage indique :%
\citComm{Παύσων ὁ ζωγράφος εἰς πενίαν διαβάλλεται, καὶ διὰ τοῦτο εἶπε `νηστεύει'.}%
Pauson l'importun est aussi mentionné au vers 853 des \italique{Acharniens} :%
\citComm{οὐδ᾽ αὖθις αὖ σε σκώψεται\\
Παύσων ὁ παμπόνηρος\\
Ni les brocards de Pauson, le gueux achevé, tu ne les essuieras [sur le marché].\\
Schol. : οὗτος ὁ Παύσων ζωγράφος πένης σχωπτολόγος.}%
Enfin au vers 602 du \italique{Ploutos} :%
\citComm{Χρέμυλος --- Παρὰ τῆς Ἑκάτης ἔξεστιν τοῦτο πυθέσθαι,\\
εἴτε τὸ πλουτεῖν εἴτε τὸ πεινῆν βέλτιον. Φησὶ γὰρ αὕτη\\
τοὺς μὲν ἔχοντας καὶ πλουτοῦντας δεῖπνον κατὰ μῆν᾽ ἀποπέμπειν,\\
τοὺς δὲ πένητας τῶν ἀνθρώπων ἁρπάζειν πρὶν καταθεῖναι.\\
Ἀλλὰ φθείρου καὶ μὴ γρύξῃς\\
ἔτι μηδ᾽ ὁτιοῦν.\\
Οὐ γὰρ πείσεις, οὐδ᾽ ἢν πείσῃς.\\
Πενία --- Ὦ πόλις Ἄργους, κλύεθ᾽ οἷα λέγει.\\
Χρέμυλος --- \textbf{Παύσωνα} κάλει τὸν ξύσσιτον.\\
Chrémyle --- Il n'est que de s'informer chez Hécate, lequel vaut mieux d'être riche ou d'être pauvre. Car elle prétend que les riches et les opulents lui font hommage chaque mois d'un dîner, et que les pauvres s'en emparent avant qu'on l'ait mis là. Mais la peste soit de toi ; ne souffle plus le mot. Car tu ne me persuaderas pas, quand tu me persuaderais.\\
La Pauvreté --- « Cité des Argiens, écoutez ce qu'il dit ! »\\
Chrémyle --- Invoque \textbf{Pauson}, ton commensal.\\
Schol. : ὁ Παύσων δὲ ἐπὶ πενίᾳ κωμῳδεῖται ζωγράφος ὦν. --- μετακαλοῦ σύντροφον τὸν Παύσωνα.}%
On constate que seuls les commentateurs indiquent que Pauson fût peintre (ζωγράφος) : il n'en n'est jamais question chez Aristophane.
Il paraît douteux que ce Pauson-là, le mendiant fâcheux à la dévotion fort intéressée\marqueur{4}, fût le peintre connu qu'évoque Aristote\marqueur{3} ; la confusion venant peut-être de ce que la façon de peindre les figures du ζωγράφος ainsi que la seule anecdote nous ayant été rapportée à son sujet témoignent d'un caractère porté à la misanthropie\marqueur{5}.}%
\debutNote{949-52}%
\note{1}{Respectivement Michel Apostolius (XIII, 2) et Suidas (π 823).}%
\note{2}{\italique{Recueil Milliet : Textes grecs et latins relatifs à l'histoire de la peinture ancienne}, tome I, publiés, traduits et commentés par Adolphe Reinach (1921) ; §§ 131-132 ; 173-179.}%
\note{3}{Aristote, \italique{Politique}, 1340a36 ; \italique{Poétique}, 1448a6.}%
\note{4}{Si sa dévotion prend les formes habituelles du souhait de renouvellement des prières (cf. le vers 287 ; Euripide, \italique{Électre}, vers 805), on se doute bien de ce qui attirait réellement Pauson-crève-la-faim. Voir ce que nous avons déjà dit aux vers 558-9 de la κρεανομία -- ou orgie institutionnalisée -- qui seule explique la réputation de piété des Athéniens.}%
\note{5}{Pauson le peintre peignait des figures pires que ce qu'elles étaient en réalité et Aristote le déconseille pour l'éducation des jeunes gens. On lui avait commandé le tableau d'un cheval se roulant dans la poussière ; l'artiste dessine un cheval courant et soulevant un nuage de poussière autour de lui. Il y travaillait quand survient l'amateur. Celui-ci récrimine : ce n'est pas le sujet commandé. Pauson fait retourner par son garçon le tableau sens dessus dessous ; on voit alors un cheval renversé, qui se roule dans la poussière.}\\
πολλάκις αὐταῖν ἐκ τῶν ὡρῶν\apparat{950}{\lecon{αὐταῖν} R : αὐτοῖν Cobet}\\
ἐς τὰς ὥρας ξυνεπευχόμενος}
\end{groupe}%
\begin{groupe}%
\typeVers{paroem.}\dimAnap{τοιαῦτα μέλειν θάμ᾽ ἑαυτῷ.\apparat{952}{μέλλειν R : \lecon{μέλειν} Zanetti (\italique{curae esse} Divus) : μολεῖν Casaubon}\\\passeLigne}%
\end{groupe}%
%% v. 953-958
\begin{groupe}%
\ledpb%
\personnageGrec{Χορός}%
\typeVers{anap.}\chantChoeur{ὅρμα χώρει·}
\end{groupe}%
%% v. 953-958
\begin{groupe}%
\typeVers{lecyth.}\chantChoeur{κοῦφα ποσὶν ἄγ᾽ ἐς κύκλον,}
\end{groupe}%
%% v. 953-958
\begin{groupe}%
\typeVers{aristoph.}\chantChoeur{χειρὶ σύναπτε χεῖρα,\apparat{955}{\lecon{χειρὶ} R : χερὶ Dobree}%
\commentaire{955}{χειρὶ σύναπτε χεῖρα.}{La danse que commande le chœur est probablement ce que Maurice Emmanuel\marqueur{1} appelle une \italique{danse des mains jointes}. Il s'agit donc bien d'une danse individuelle et non en couple : chaque danseur joignant la main à son autre main.}%
\debutNote{955}%
\note{1}{\italique{La danse grecque antique} (1896), §319.}}%
\end{groupe}%
%% v. 953-958
\begin{groupe}%
\chantChoeur{†ῥυθμὸν χορείας ὕπαγε πᾶσα†·\apparat{956}{\lecon{ὕπαγε πᾶσα} R : πᾶσ' ὕπαγε Austin}}
\end{groupe}%
%% v. 953-958
\begin{groupe}%
\typeVers{glyc.}\chantChoeur{\ligne{956}βαῖνε καρπαλίμοιν ποδοῖν.}
\end{groupe}%
%% v. 953-958
\begin{groupe}%
\typeVers{dim. iamb.}\chantChoeur{ἐπισκοπεῖν δὲ πανταχῇ}
\end{groupe}%
%% v. 953-958
\begin{groupe}%
\typeVers{trim. iamb.}\chantChoeur{κυκλοῦσαν ὄμμα χρὴ χοροῦ κατάστασιν.\\\passeLigne}%
\end{groupe}%
%% v. 959-961 strophe a
\begin{groupe}%
\Strophe{(στροφή α)}\typeVers{cret. $+$ lec.}\chantChoeur{ἅμα δὲ καὶ \ligne{960}γένος Ὀλυμπίων θεῶν}
\end{groupe}%
%% v. 959-961 strophe a
\begin{groupe}%
\typeVers{$4$ troch\cat}\chantChoeur{μέλπε καὶ γέραιρε φωνῇ πᾶσα χορομανεῖ τρόπῳ.\\\passeLigne}%
\end{groupe}%
%% v. 962-965 antistrophe a
\begin{groupe}%
\Strophe{(ἀντιστρ. α)}\typeVers{cret. $+$ lec.}\chantChoeur{εἰ δέ τις \ligne{963}προσδοκᾷ κακῶς ἐρεῖν}
\end{groupe}%
%% v. 962-965 antistrophe a
\begin{groupe}%
\typeVers{$4$ troch\cat}\chantChoeur{ἐν ἱερῷ γυναῖκά μ᾽ οὖσαν ἄνδρας,\ligne{965} οὐκ ὀρθῶς φρονεῖ.\\\passeLigne}%
\end{groupe}%
%% v. 966-968 mesodos a
\begin{groupe}%
\Strophe{(μεσῳδός α)}\typeVers{cret.}\chantChoeur{ἀλλὰ χρή\apparat{966}{\lecon{χρὴ} R : χρῆν Bothe : χρή μ᾽ Meineke}}
\end{groupe}%
%% v. 966-968 mesodos a
\begin{groupe}%
\typeVers{dim. troch.}\chantChoeur{ὥσπερ ἔργον αὖ τι καινὸν\apparat{967}{\lecon{ὥσπερ} R : ὡς πρὸς Austin}}
\end{groupe}%
%% v. 966-968 mesodos a
\begin{groupe}%
\typeVers{$4$ troch\cat}\chantChoeur{πρῶτον εὐκύκλου χορείας εὐφυᾶ στῆσαι βάσιν.\testimonia{968}{(εὐκύκλου) Hesychius ε 6929 εὔκυκλος· εὔχωρος $=$ Photius ε 2236}\apparat{968}{εὐφυῆ R : \lecon{εὐφυᾶ} Brunck}\\\passeLigne}%
\end{groupe}%
%% v. 969-976 strophe b
\begin{groupe}%
\ledpb%
\Strophe{(στροφή β)}\typeVers{dim. iamb.}\chantChoeur{\ligne{969}πρόβαινε ποσὶ τὸν εὐλύραν\apparat{969}{ποσὶν R : \lecon{ποσὶ} Bentley \espApp ἐλύραν R : \lecon{εὐλύραν} Zanetti}}
\end{groupe}%
%% v. 969-976 strophe b
\begin{groupe}%
\typeVers{dim. iamb.}\chantChoeur{μέλπουσα καὶ τὴν τοξοφόρον}
\end{groupe}%
%% v. 969-976 strophe b
\begin{groupe}%
\typeVers{$2$ iamb\cat}\chantChoeur{Ἄρτεμιν ἄνασσαν ἁγνήν.}
\end{groupe}%
%% v. 969-976 strophe b
\begin{groupe}%
\typeVers{reiz.}\chantChoeur{χαῖρ᾽ ὦ ἑκάεργε,}
\end{groupe}%
%% v. 969-976 strophe b
\begin{groupe}%
\typeVers{reiz.}\chantChoeur{ὄπαζε δὲ νίκην·}
\end{groupe}%
%% v. 969-976 strophe b
\begin{groupe}%
\typeVers{$2$ iamb\cat}\chantChoeur{Ἥραν τε τὴν τελείαν\testimonia{973}{(τελείαν) Hesychius τ 400 Τελεία}}
\end{groupe}%
%% v. 969-976 strophe b
\begin{groupe}%
\typeVers{$2$ iamb\cat}\chantChoeur{μέλψωμεν ὥσπερ εἰκός,}
\end{groupe}%
%% v. 969-976 strophe b
\begin{groupe}%
\typeVers{trim. iamb.}\chantChoeur{ἣ πᾶσι τοῖς χοροῖσιν ἐμπαίζει τε καὶ}
\end{groupe}%
%% v. 969-976 strophe b
\begin{groupe}%
\typeVers{$2$ iamb\cat}\chantChoeur{κλῇδας γάμου φυλάττει.\\\passeLigne}%
\end{groupe}%
%% v. 977-984 antistrophe b
\begin{groupe}%
\Strophe{(ἀντιστρ. β)}\typeVers{dim. iamb.}\chantChoeur{Ἑρμῆν τε νόμιον ἄντομαι}
\end{groupe}%
%% v. 977-984 antistrophe b
\begin{groupe}%
\typeVers{dim. iamb.}\chantChoeur{καὶ Πᾶνα καὶ Νύμφας φίλας}
\end{groupe}%
%% v. 977-984 antistrophe b
\begin{groupe}%
\typeVers{$2$ iamb\cat}\chantChoeur{\ligne{979}ἐπιγελάσαι προθύμως}
\end{groupe}%
%% v. 977-984 antistrophe b
\begin{groupe}%
\typeVers{reiz.}\chantChoeur{ταῖς ἡμετέραισι\apparat{980}{ἡμετέραις R : \lecon{ἡμετέραισι} Hermann}}
\end{groupe}%
%% v. 977-984 antistrophe b
\begin{groupe}%
\typeVers{reiz.}\chantChoeur{χαρέντα χορείαις.}
\end{groupe}%
%% v. 977-984 antistrophe b
\begin{groupe}%
\typeVers{$2$ iamb\cat}\chantChoeur{ἔξαιρε δὴ προθύμως}
\end{groupe}%
%% v. 977-984 antistrophe b
\begin{groupe}%
\typeVers{$2$ iamb\cat}\chantChoeur{\ligne{982}διπλῆν χάριν χορείας.\apparat{982}{χαίρειν R : \lecon{χάριν} Ellebodius}%
\commentaire{982}{διπλῆν.}{Pollux (IV, 105) compte la διπλῆ parmi les figures de danse pour la tragédie (τραγικῆς ὀρχήσεως σχήματα) :%
\citComm{Καὶ μὲν τραγικῆς ὀρχήσεως σχήματα σιμή χείρ, καλαθίσκος, χεὶρ καταπρανής, ξύλου παράληψις, \textbf{διπλῆ}, θερμαυστρίς, κυβίστησις, παραβῆναι τὰ τέτταρα}%
Hésychius (δ 1948) ne nous apprend rien davantage :%
\citComm{διπλῆ· ὀρχήσεως εἶδος, ἢ κρούματος.}%
Il nous est difficile de reconstituer, même grossièrement, la chorégraphie de ce choral puisqu'Aristophane recourt au vocabulaire technique de la danse de son époque, dont on ne sait plus grand chose.}}
\end{groupe}%
%% v. 977-984 antistrophe b
\begin{groupe}%
\typeVers{trim. iamb.}\chantChoeur{παίσωμεν ὦ γυναῖκες οἷάπερ νόμος,}
\end{groupe}%
%% v. 977-984 antistrophe b
\begin{groupe}%
\typeVers{$2$ iamb\cat}\chantChoeur{νηστεύομεν δὲ πάντως.\apparat{984}{νηστεύωμεν R : \lecon{νηστεύομεν} Bentley}\\\passeLigne}%
\end{groupe}%
%% v. 985-989 mesodos b
\begin{groupe}%
\Strophe{(μεσῳδός β)}\typeVers{trim. iamb.}\chantChoeur{\ligne{985}ἀλλ᾽ εἶα πάλλ᾽ ἀνάστρεφ᾽ εὐρύθμῳ ποδί,\apparat{985}{ειαπαλλ᾽ R : \lecon{εἶα πάλλ᾽} Blaydes}}
\end{groupe}%
%% v. 985-989 mesodos b
\begin{groupe}%
\typeVers{$2$ iamb\cat}\chantChoeur{τόρευε πᾶσαν ᾠδήν·\testimonia{986}{Suidas τ 789}}
\end{groupe}%
%% v. 985-989 mesodos b
\begin{groupe}%
\typeVers{ia. $+$ mol.}\chantChoeur{ἡγοῦ δέ γ᾽ ὧδ᾽ αὐτὸς σὺ%\apparat{987}{γ᾽ ὧδ᾽ suspectum}
}
\end{groupe}%
%% v. 985-989 mesodos b
\begin{groupe}%
\typeVers{cret. $+$ mol.}\chantChoeur{κισσοφόρε Βακχεῖε}
\end{groupe}%
%% v. 985-989 mesodos b
\begin{groupe}%
\typeVers{aristoph.}\chantChoeur{δέσποτ᾽· ἐγὼ δὲ κώμοις}
\end{groupe}%
%% v. 985-989 mesodos b
\begin{groupe}%
\typeVers{$2$ iamb\cat}\chantChoeur{\ligne{989}σὲ φιλοχόροισι μέλψω.\apparat{989}{φιλοχόροισιν R : \lecon{φιλοχόροισι} Scaliger}\\\passeLigne}%
\end{groupe}%
%% v. 990-994 strophe g
\begin{groupe}%
\Strophe{(στροφή γ)}\typeVers{aristoph.}\chantChoeur{\ligne{990}†εὔιον ὦ Διόνυσε\apparat{990}{\lecon{εὔιον} R : Εὔιε Enger : Εὐοῖ Hermann (1845) \espApp \lecon{Διόνυσε} R : Διὸς σὺ Enger : Διός τε Fritzsche}%
%\commentaire{990}{}{\eff{Hésychius ε 6909 εὐΐων· Διονυσιακῶν} (Ευρ. Βαϲϲ 791...) εὐϊάζω = εὐάζω Eur. Cycl. 495.}
}
\end{groupe}%
%% v. 990-994 strophe g
\begin{groupe}%
\typeVers{pherecr.}\chantChoeur{Βρόμιε καὶ Σεμέλας παῖ,}
\end{groupe}%
%% v. 990-994 strophe g
\begin{groupe}%
\typeVers{hexasyll. aeol.}\chantChoeur{χοροῖς τερπόμενος}
\end{groupe}%
%% v. 990-994 strophe g
\begin{groupe}%
\typeVers{ia. $+$ arist.}\chantChoeur{κατ᾽ ὄρεα νυμφᾶν ἐρατοῖς ἐν ὕμνοις,\apparat{993}{\lecon{ἐρατοῖς ἐν ὕμνοις} R : ἐρατοῖσιν ὕμνοις Wilamowitz (\italique{Gr. Versk. p.476}) \ligneApp{994} \lecon{εὔιον εὔιον} R : ὦ Εὔι Εὔι' Hermann \espApp lacunam statuit Hermann : ἥκεις γὰρ von Velsen : ἡδόμενος Austin cf. \italique{Plut}.288}}
\end{groupe}%
%% v. 990-994 strophe g
\begin{groupe}%
\typeVers{$2$ iamb\cat}\chantChoeur{†εὔιον εὔιον εὐοῖ%\apparat{994}{\lecon{εὔιον εὔιον} R : ὦ Εὔι Εὔι' Hermann \espApp lacunam statuit Hermann : ἥκεις γὰρ von Velsen : ἡδόμενος Austin cf. \italique{Plut}.288}
}
\end{groupe}%
%% v. 990-994 strophe g
\begin{groupe}%
%\times \  \overarc{\smile \smile} \widearc{\smile \smile} \overset{\huge{\frown}}{\cup \cup} \overset{\huge{\frown}}{\smile \smile \:} \overset{\huge{\frown}}{\smile \smile} \ \cup
% \widearc{\smile \smile}
\typeVers{$2$ iamb\cat}\chantChoeur{\ligne{994}< $ \times \ \widearc{\cup \cup} \ \cup $ > ἀναχορεύων.\\\passeLigne}%
\end{groupe}%
%% v. 995-1000 antistrophe g
\begin{groupe}%
\Strophe{(ἀντιστρ. γ)}\typeVers{aristoph.}\chantChoeur{ἀμφὶ δὲ σοὶ κτυπεῖται\apparat{995}{συῒ κτυπεῖται R : \lecon{σοὶ κτυπεῖται} Zanetti : συγκτυπεῖται Wilamowitz \ligneApp{996} Κιθαρώνιος R : \lecon{Κιθαιρώνιος} Zanetti)}}
\end{groupe}%
%% v. 995-1000 antistrophe g
\begin{groupe}%
\typeVers{pherecr.}\chantChoeur{Κιθαιρώνιος ἠχώ,}
\end{groupe}%
%% v. 995-1000 antistrophe g
\begin{groupe}%
\typeVers{hexasyll. aeol.}\chantChoeur{μελάμφυλλά τ᾽ ὄρη}
\end{groupe}%
%% v. 995-1000 antistrophe g
\begin{groupe}%
\typeVers{(?)}\chantChoeur{δάσκια †καὶ νάπαι πετρώδεις† βρέμονται·\apparat{998}{\lecon{καὶ νάπαι πετρώδεις} R : πετρώδεις τε νάπαι Enger}}
\end{groupe}%
%% v. 995-1000 antistrophe g
\begin{groupe}%
\typeVers{$2$ iamb\cat}\chantChoeur{κύκλῳ δὲ περὶ σὲ κισσὸς}
\end{groupe}%
%% v. 995-1000 antistrophe g
\begin{groupe}%
\typeVers{$2$ iamb\cat}\chantChoeur{εὐπέταλος ἕλικι θάλλει.\lednopb%
\\\passeLigne}%
\end{groupe}%
%% v. 1001
\begin{groupe}%
\ledpb%
\personnageGrec{Τοξότης}%
\typeVers{trim. iamb.}\trimIamb{ἐνταῦτα νῦν οἰμῶξι πρὸς τὴν αἰτρίαν.\apparat{1001}{οἰμώξει R : \lecon{οἰμῶξι} Brunck}\\\passeLigne}%
\end{groupe}%
%% v. 1002a
\begin{groupe}%
\personnageGrec{Μνησίλοχος}%
\trimIamb{ὦ τοξόθ᾽ ἱκετεύω σε.%
\commentaire{1001}{τοξότης.}{Aristophane a donné à l'archer un jargon mi-scythe mi-grec, une sorte de parler nègre. On sent qu'il est impossible de rendre cela dans une traduction. (Note de Willems.)\\
À Athènes, comme partout ailleurs, les personnes chargées des missions de police étaient recrutées parmi la lie de l'humanité. En l'occurrence, il y avait en guise de police un contingent de 300 mercenaires Scythes aux ordres des prytanes\marqueur{1}. C'est cette inquiétant contraste, entre l'importance de la mission (les forces de l'«~ordre~», comme on dit...) et la bassesse de ceux qui y sont préposés, qui inspirera à Platon de si longs développements sur le bon choix et l'éducation des gardiens de sa Cité. \latin{Quis custodiet custodes?}}%
\debutNote{1001}%
\note{1}{Que les archers Scythes fussent directement aux ordres des prytanes, c'est assez bien établi par notre pièce ainsi que le vers 54 des \italique{Acharniens}, cf. aussi \italique{Protagoras} (319c).
On pourra lire l'article de Jean-Christophe Couvenhes «~L'introduction des archers scythes, esclaves publics, à Athènes~», in \italique{Transferts culturels et droits dans le monde grec et hellénistique} (2012), pp. 99-118, ou l'article plus ancien d'Agnès Plassart «~Les archers d'Athènes~», in \italique{Revue des Études Grecques}, t. 26 (1913), pp. 151-213 dont quelques détails sont discutés dans le premier article.}%
%\note{2}{}%
\commentaire{1001}{ἐνταῦτα νῦν οἰμῶξι πρὸς τὴν αἰτρίαν}{$=$ ἐνταῦθα νῦν οἴμωζε πρὸς τὴν αἰθρίαν.
(Noter le sadisme proprement argousin de l'archer.)% : il dit à Mnésiloque qu'il va pouvoir gémir au grand air (οἰμῶξι πρὸς τὴν αἰτρίαν) au moment même où il lui coupe le sifflet et serre la vis qui l'étrangle.
}\\\passeLigne}%
\end{groupe}%
%% v. 1002b
\begin{groupe}%
\personnageGrec{Τοξότης}%
\trimIamb{\ligne{1002}\decalage{ὦ τοξόθ᾽ ἱκετεύω σε.}μή μ᾽ ἰκετεῦσι σύ.\apparat{1002}{ἱκετεῦση R : \lecon{ἱκετεῦσι} Brunck}\\\passeLigne}%
\end{groupe}%
%% v. 1003a
\begin{groupe}%
\personnageGrec{Μνησίλοχος}%
\trimIamb{χάλασον τὸν ἧλον.%\commentaire{}{}{\eff{Souda} Σχάσον}
\\\passeLigne}%
\end{groupe}%
%% v. 1003b
\begin{groupe}%
\personnageGrec{Τοξότης}%
\trimIamb{\ligne{1003}\decalage{χάλασον τὸν ἧλον.}ἀλλὰ ταῦτα δρᾶσ᾽ ἐγώ.\apparat{1003}{δρᾷς R : \lecon{δρᾶσ᾽} Voss}\\\passeLigne}%
\end{groupe}%
%% v. 1004
\begin{groupe}%
\personnageGrec{Μνησίλοχος}%
\trimIamb{οἴμοι κακοδαίμων, μᾶλλον ἐπικρούεις σύ γε.\apparat{1004}{ἐπικρούσεις R : \lecon{ἐπικρούεις} \Mon}\\\passeLigne}%
\end{groupe}%
%% v. 1005a
\begin{groupe}%
\personnageGrec{Τοξότης}%
\trimIamb{ἔτι μᾶλλο βοῦλις;\apparat{1005}{μᾶλλον R : \lecon{μᾶλλο} Bentley}\\\passeLigne}%
\end{groupe}%
%% v. 1005b-1006a
\begin{groupe}%
\personnageGrec{Μνησίλοχος}%
\trimIamb{\ligne{1005}\decalage{ἔτι μᾶλλο βοῦλις;}ἀτταταῖ ἰατταταῖ·\apparat{1005}{τατταταί R : \lecon{ἰατταταῖ} Faber}\\
κακῶς ἀπόλοιο.\\\passeLigne}%
\end{groupe}%
%% v. 1006b-1007
\begin{groupe}%
\personnageGrec{Τοξότης}%
\trimIamb{\ligne{1006}\decalage{κακῶς ἀπόλοιο.}σῖγα κακοδαίμων γέρον.\apparat{1006}{κακοδαίμω\expo{ν} γέρ\inRas{ω}{ο}ν \Rpc, M\expo{pc}}\\
πέρ᾽ ἐγὼ 'ξινίγκι πορμός, ἴνα πυλάξι σοι.\apparat{1007}{ξεινίγκι R : \lecon{'ξινίγκι} Enger \espApp ἵνα R : \lecon{ἴνα} Brunck \espApp \lecon{πυλάξι} R : πυλάξῃ Schol.}%
\commentaire{1007}{πέρ᾽ ἐγὼ 'ξινίγκι πορμός, ἴνα πυλάξι σοι}{$=$ φέρ᾽ ἐγὼ 'ξενέγκω φορμόν, ἵνα φυλάξω σε.}\\\passeLigne}%
\end{groupe}%
%% v. 1008-1014
\begin{groupe}%
\personnageGrec{Μνησίλοχος}%
\trimIamb{ταυτὶ τὰ βέλτιστ᾽ ἀπολέλαυκ᾽ Εὐριπίδου.\\
ἔα· θεοί, Ζεῦ σῶτερ, εἰσὶν ἐλπίδες.%\apparat{1009}{εἰσὶν R : εἴς ἔτ᾽ Naber (Faber ?) \cf 946}
\\
ἁνὴρ ἔοικεν οὐ προδώσειν, ἀλλά μοι\apparat{1010}{ἀνὴρ R : \lecon{ἁνὴρ} Bothe \espApp \lecon{ἔοικεν} R : μ᾽ ἔοικε von Velsen \ligneApp{1011} ὑπεδήλωσεν R : \lecon{ὑπεδήλωσε} Biset}\\
σημεῖον ὑπεδήλωσε Περσεὺς ἐκδραμών,\\
ὅτι δεῖ με γίγνεσθ᾽ Ἀνδρομέδαν· πάντως δέ μοι\\
τὰ δέσμ᾽ ὑπάρχει. δῆλον οὖν τοῦτ᾽ ἔσθ᾽ ὅτι\apparat{1013}{οὖν R : οὖν νῦν Bizet : \lecon{οὖν τοῦτ᾽} Dobree}\\
ἥξει με σώσων· οὐ γὰρ ἂν παρέπτετο.\apparat{1014}{σῶσον R : \lecon{σώσων} \Mon}\\\passeLigne}%
\end{groupe}%
%% v. 1015-1055
\begin{groupe}%
\ledpb%
\personnageGrec{Μνησίλοχος}%
\typeVers{bacch. $+$ ia.}\paratragiques{\ligne{1015}φίλαι παρθένοι φίλαι,\apparat{1015-21}{(Εὐρ.) R : \lecon{(Μνη.)} Tyrwhitt}%
\commentaire{1015-135}{%Hésychius + Synagoge : αἰδέσασθαι: τὸ ἔνοχον ὄντα [καὶ] φόνῳ ἀκουσίῳ καὶ πεφυκαδευμένον ἐφ᾽ ὡρισμένῳ χρόνῳ... τοῦ τετελευτηκότος ἐξιλᾶσθαι, ὡς εἰληφότος ἤδη τιμωρίαν
%αἰδῶσσα· les murs d'une cour Hsch
%ἴτας:ὅρκος.καὶ ὁ “Αιδης
Parodie de l'\italique{Andromède}.}{La pièce perdue d'Euripide \italique{Andromède} fut peut-être représentée la même année que l'\italique{Hélène}\marqueur{1}, en 412, un avant les \italique{Thesmophories}. L'édition des fragments d'Euripide de Jouan \& van Looy\marqueur{2} propose un dossier complet auquel le lecteur pourra se reporter. Le récit le plus circonstancié du mythe d'Andromède se trouve chez Apollodore (II, 4, 3) :%
\citComm{Παραγενόμενος δὲ εἰς Αἰθιοπίαν, ἧς ἐβασίλευε Κηφεύς, εὗρε τὴν τούτου θυγατέρα Ἀνδρομέδαν παρακειμένην βορὰν θαλασσίῳ κήτει. Κασσιέπεια γὰρ ἡ Κηφέως γυνὴ Νηρηίσιν ἤρισε περὶ κάλλους, καὶ πασῶν εἶναι κρείσσων ηὔχησεν· ὅθεν αἱ Νηρηίδες ἐμήνισαν, καὶ Ποσειδῶν αὐταῖς συνοργισθεὶς πλήμμυράν τε ἐπὶ τὴν χώραν ἔπεμψε καὶ κῆτος. Ἄμμωνος δὲ χρήσαντος τὴν ἀπαλλαγὴν τῆς συμφορᾶς, ἐὰν ἡ Κασσιεπείας θυγάτηρ Ἀνδρομέδα προτεθῇ τῷ κήτει βορά, τοῦτο ἀναγκασθεὶς ὁ Κηφεὺς ὑπὸ τῶν Αἰθιόπων ἔπραξε, καὶ προσέδησε τὴν θυγατέρα πέτρα. Ταύτην θεασάμενος ὁ Περσεὺς καὶ ἐρασθεὶς ἀναιρήσειν ὑπέσχετο Κηφεῖ τὸ κῆτος, εἰ μέλλει σωθεῖσαν αὐτὴν αὐτῷ δώσειν γυναῖκα. Ἐπὶ τούτοις γενομένων ὅρκων, ὑποστὰς τὸ κῆτος ἔκτεινε καὶ τὴν Ἀνδρομέδαν ἔλυσεν. Ἐπιβουλεύοντος δὲ αὐτῷ Φινέως, ὃς ἦν ἀδελφὸς τοῦ Κηφέως ἐγγεγυημένος πρῶτος τὴν Ἀνδρομέδαν, μαθὼν τὴν ἐπιβουλήν, τὴν Γοργόνα δείξας μετὰ τῶν συνεπιβουλευόντων αὐτὸν ἐλίθωσε παραχρῆμα.\\
Parvenu en Éthiopie, dont Céphée était roi, [Persée] trouva sa fille Andromède exposée pour être dévorée par un monstre marin. Cassiopée, épouse de Céphée, avait osé se comparer aux Néréides pour la beauté, et s'était même vantée de l'emporter sur elles. Les Néréides en furent irritées. Neptune partagea leur indignation, submergea le pays, et y envoya un monstre marin. L'oracle d'Ammon ayant annoncé que ces désastres cesseraient si on exposait Andromède, fille de Cassiopée, pour être dévorée par le monstre, les Éthiopiens forcèrent Céphée à faire ce que l'oracle ordonnait et à attacher sa fille à un rocher. Persée l'ayant vu, il devint amoureux, et promit à Céphée de tuer le monstre s'il voulait la lui donner en mariage. Céphée s'y étant engagé par serment, il attendit le monstre, le tua et délivra Andromède. Phinée, frère de Céphée, à qui Andromède avait été promise avant cet événement, conspira contre lui pour le faire périr ; mais Persée l'ayant découvert, lui montra la tête de la Gorgone, et le changea en pierre, ainsi que tous ceux qui avaient pris part à son complot.}%
}%
\debutNote{1015-135}%
\note{1}{C'est ce qu'indiquent les scholies au vers 1012, πιθανῶς· συνδεδίδακται γὰρ τῇ Ἑλένη, ainsi que celle du vers 53 des \italique{Grenouilles}. Le début du prologue d'\italique{Andromède} est repris au vers 1065, «~ὦ νὺξ ἱερὰ~», -- oui : Aristophane a chamboulé l'ordre de la pièce ; scholie à ce vers 1065 : τοῦ προλόγου Ἀνδρομέδας εἰσβολή --, laisse à penser qu'\italique{Andromède} fut jouée au lever du jour, avant donc \italique{Hélène}. Matthew Whright (\italique{Euripides' Escape-Tragedies}, 2005) fait l'hypothèse contraire et avance qu'\italique{Iphigénie en Tauride} concluait la trilogie d'Euripide de 412, agrémenté du drame satyrique du \italique{Cyclope}. Voir la recension que feu Alain Blanchard en fit dans la \italique{Revue des Études Grecques}, tome 119 (2006), pp. 798-800.}%
\note{2}{Tome VIII (fragments), première partie, de l'édition d'Euripide aux Belles Lettres (1998).}%
}
\end{groupe}%
%% v. 1015-1055
\begin{groupe}%
\typeVers{chor. $+$ cret.}\paratragiques{πῶς ἂν ἀπέλθοιμι καὶ\apparat{1016}{\lecon{ἀπέλθοιμι} R : ἐπέλθοιμι Ellebodius}% \espApp ἂν <οὖν> Hermann ut numeri dochmiaci fierent}
\commentaire{1015-55}{Monodie d'Andromède-Mnésiloque.}{La structure de ce chant a fait couler beaucoup d'encre... et nous échappe encore\marqueur{1}. Mais bien loin de l'anarchie décrite par les commentateurs, je suis tenté d'y voir une illustration sonore des échos : chaque phrase musicale semble être répétée : l'on peut ainsi regrouper par paires successives les différentes structures métriques (le thème puis son écho, de mètre identique ou très proche). On regroupe ainsi les vers : 1015-1071 \& 1018-1021 ; 1022-1023 \& 1024-1025 ; 1026-1029 \& 1030-1031 ; 1032-1033 \& 1034-1035 ; 1036-1037 \& 1038-1039a ; 1043 \& 1045 ; 1050-1051 \& 1052-1053.}%
\debutNote{1015-55}%
\note{1}{On pourra aller voir ce qu'en dit Lætitia, pp. 436-445.}%
}
\end{groupe}%
%% v. 1015-1055
\begin{groupe}%
\typeVers{$+$ ityph.}\paratragiques{τὸν Σκύθην λάθοιμι;\apparat{1017}{λάβοιμι R : \lecon{λάθοιμι} Ellebodius}}
\end{groupe}%
%% v. 1015-1055
\begin{groupe}%
\typeVers{$4$ bacch.}\paratragiques{\ligne{1019}κλύεις; ὦ προσᾴδουσ' ἀϋτὰς ἐν ἄντροις,\apparat{1019}{προσᾳδοῦσσαι τὰς R, Schol. : \lecon{προσᾴδουσ'} Elmsley, \lecon{ἀϋτὰς} Burges (ἀϋταῖς Sommerstein) : alii alia}%
\commentaire{1018}{ὦ προσᾳδοῦσσαι τὰς ἐν ἄντροις.}{Comme dans la pièce d'Euripide, la nymphe Écho répond aux plaintes d'Andromède. (Dans la reprise comique, c'est Euripide qui se métamorphose en Écho.)
Willems a consacré sa \nieme{25} note (reproduite en annexe) à ce vers et aux diverses retouches qu'il a subies.
Je reprends la lecture de Burges\marqueur{1}, προσᾴδουσ' ἀϋτάς, satisfaisante quant au sens et qui offre l'avantage de rester très proche du texte du manuscrit.%\eff{αἱ Λιταί}
}%
\debutNote{1018}%
\note{1}{G. Burges, \italique{Classical Journal} 14 (1816) 217. Ou bien voir W. Mitsdörffer, «~Das Mnesichoslied in Aristophanes' \italique{Thesmophoriazusen}~», \italique{Philologus} 98 (1954-55), p. 69.}}
\end{groupe}%
%% v. 1015-1055
\begin{groupe}%
\typeVers{telesill.}\paratragiques{κατάνευσον, ἔασον ὡς}
\end{groupe}%
%% v. 1015-1055
\begin{groupe}%
\typeVers{$+$ ithyph.}\paratragiques{τὴν γυναῖκά μ᾽ ἐλθεῖν.}
\end{groupe}%
%% v. 1015-1055
\begin{groupe}%
\typeVers{dim. iamb.}\paratragiques{ἄνοικτος ὅς μ᾽ ἔδησε τὸν\apparat{1022}{(Μνη.) praefixit \Rpc, \Mon}}
\end{groupe}%
%% v. 1015-1055
\begin{groupe}%
\typeVers{$+$ lecyth.}\paratragiques{πολυπονώτατον βροτῶν·\apparat{1023}{\lecon{πολυπονώτατον} R, Schol. : πολυστονώτατον Burges}}
\end{groupe}%
%% v. 1015-1055
\begin{groupe}%
\typeVers{dim. iamb.}\paratragiques{\ligne{1024}μόλις δὲ γραῖαν ἀποφυγὼν}
\end{groupe}%
%% v. 1015-1055
\begin{groupe}%
\typeVers{dim. iamb.}\paratragiques{σαπρὰν ἀπωλόμην ὅμως.}
\end{groupe}%
%% v. 1015-1055
\begin{groupe}%
\typeVers{do. $+$ ia. $+$ do.}\paratragiques{\ligne{1026}ὅδε γὰρ ὁ Σκύθης φύλαξ πάλαι ἐφέστηκ᾽· ὀλοὸν\apparat{1026}{ὁ Σκύθης R : del. Jackson}%
\commentaire{1026-8}{ὅδε γὰρ ὁ Σκύθης φύλαξ / πάλαι ἐφέστηκ᾽ ὀλοὸν / ἄφιλον ἐκρέμασε κόραξι δεῖπνον.}{Ce passage est sans doute légèrement corrompu. Les tentatives pour le corriger ont toutes été infructueuses\marqueur{1}. }%
\debutNote{1026-8}%
\note{1}{Cf. l'inventaire sans concession qu'en fait Lætitia \italique{The Songs of Aristophanes}, p. 443-444 : son regard infaillible de métricienne aguerrie détecte les invraisemblances métriques des «~corrections~». Sa proposition n'en est pas moins aporétique : elle voit dans ces vers une suite de trois dochmiaques se terminant par un dimètre iambique catalectique tout en concédant que la correption πάλα\lettreBreve{ι} (\latin{vocalis ante vocalem corripitur}) est suspecte : les seuls autres exemples chez Aristophane étant au vers 915, «~κύσ\lettreBreve{ω}. ἄπαγε~», déjà évoqué plus haut, et le peu convaincant «~ἰὼ ἰώ~» au vers 1170 des \italique{Nuées}.}%
}
\end{groupe}%
%% v. 1015-1055
\begin{groupe}%
\typeVers{$2$ cr. $+$ ba.}\paratragiques{\ligne{1029}ἄφιλον ἐκρέμασε κόραξι δεῖπνον.\apparat{1027}{\lecon{ἐφέστηκ᾽ ὁλοὸν} R : ἐφεστὼς ὀλοὸν Meineke : ἐφέστηκε κὠλοὸν Austin}% \ligneApp{1029} \lecon{κρέμασε} R : κρέμασεν Bothe : κρέμασε με Mehler}
}
\end{groupe}%
%% v. 1015-1055
\begin{groupe}%
\typeVers{do. $+$ ia. $+$ do.}\paratragiques{ὁρᾷς; οὐ χοροῖσιν οὐδ᾽ ὑφ᾽ ἡλίκων νεανίδων}
\end{groupe}%
%% v. 1015-1055
\begin{groupe}%
\typeVers{$2$ cr. $+$ ba.}\paratragiques{\ligne{1031}ψῆφον ἕστηκ᾽ ἔχουσα κημόν,\apparat{1031}{κημὸν ἕστηκ᾽ ἔχουσ᾽ R : \lecon{ἕστηκ᾽ ἔχουσα κημόν} Bibi, cf. schol. ad \italique{Vesp.} 674c}%
\commentaire{1031}{ψῆφον κημὸν ἕστηκ᾽ ἔχουσ᾽.}{On mettra ce vers en regard d'un fragment d'Hermippe :%
\citComm{καὶ πρὸς κύβους ἕστηκ' ἔχων <τὸ> κήθιον.}%
%ξιφισμὸς (σχῆμα ἕστηκ᾽ ἔχουσ᾽)
τὸ κήθιον et ὁ κημός sont synonymes : ils désignent dans l'univers ludique le \italique{cornet} dans lequel on secouait %(λαγαρίζεσθαι)\eff{Non !!} 
les dés\marqueur{1}. ἡ ψῆφος, comme λίθος, est dans ce contexte un \italique{pion}.
On constate avec ces deux exemples que `ἕστηκεν ἔχων' devait être l'expression consacrée pour le joueur qui s'apprêtait à lancer les dés.
On a soupçonné une altération du vers du fait de l'étrange succession initiale de trois longues. 
L'exemple d'Hermippe suscite naturellement l'émendation suivante :%
\citComm{ὁρᾷς; οὐ χοροῖσιν οὐδ᾽ ὑφ᾽ ἡλίκων ν\synizese{εα}νίδων\\
ψῆφον ἕστηκ᾽ ἔχουσα κημόν.}%
(κημός placé après ἔχω, les mots sont ainsi agencés dans le même ordre que dans le fragment\marqueur{2}.)\\
L'avantage est qu'on obtient (avec la synizèze de ν\synizese{εα}νίδων) le schéma métrique suivant :%
\citComm{$ \cup \ - \ - \ \cup \ - \ | \ \cup \ - \ \cup \ - \ | \ \cup \ - \ - \ \cup \ -$\\
$ - \ \cup \ - \ | \ - \ \cup \ - \ | \ \cup \ - \ \smile $\\
dochmiaque $+$ iambe $+$ dochmiaque \\
crétique $+$ crétique $+$ bacchée,}%
de scansion identique à celle des deux vers précédents, annulant du même coup le surprenant abrègement de πάλα\lettreBreve{ι} en seconde position d'une résolution (qui le faisait scander $ \widearc{\cup \cup}$, cf. commentaire v. 1026-8 n.1). 
Je ne sais ce que vaut cette émendation ; je laisse le soin d'en juger à des métriciens compétents.\\
Il est en tout cas bien évident qu'Andromède-Mnésiloque, se lamentant de son sort présent et regrettant les douces années de son enfance, ne fait aucunement allusion à une quelconque \italique{urne de vote} (κημός ?). Les chœurs de danse et les jeux enfantins, voilà qui suffit à sa mélancolie. Le marigot électoral ? alors que la jeune fille ne pouvait même pas y participer activement ? Mais où donc l'ensemble des commentateurs -- à commencer par le scholiaste -- est-il allé cherché cette abracadabrantesque explication ?
Qu'Aristophane ait pu utiliser le mot κημός, nom d'une sorte d'entonnoir posé sur l'urne pour en celer le vote, comme synecdoque pour l'\italique{urne} elle-même, cela est vrai mais avec une réserve de taille ayant échappé aux commentateurs : c'est pour permettre un jeu de mots, comme nous l'allons voir. (Ce qui explique qu'il soit le seul auteur recourant à ce vocable avec le sens d'\italique{urne}.)\\
% l'inadmissible traduction du terme κημός par "urne" [pour le vote]
Cette inadmissible traduction de κημός par «~urne~» [pour le vote] nous oblige à rebondir sur un passage des \italique{Guêpes}. %, étudié lors de l'un de vos cours. 
%Le proposez-vous encore à vos élèves ? auquel cas, j'ai quelques éclaircissement à vous partager -- et si vous ne l'étudiez plus, tout ce qui est gratuit est de toute façon plus beau !
Il s'agit du début des \italique{Guêpes} avec la traditionnelle scène des deux esclaves qui présentent la situation : en l'occurrence, le mal dont souffre leur maître Philocléon (et une bonne partie des Athéniens), pris qu'il est de \italique{philéliastite}. Et lorsqu'il voit une dédicace au beau Dêmos, le fils de Pyrilampe, notre malade ne peut s'empêcher d'écrire en marge (παρέγραψε) : κημὸς καλός. Vers 97-99 :
\citComm{καὶ νὴ Δί᾽ ἢν ἴδῃ γέ που γεγραμμένον\\
υἱὸν Πυριλάμπους ἐν θύρᾳ Δῆμον καλόν,\\
ἰὼν παρέγραψε πλησίον `κημὸς καλός.'}%
Qui était ce Dêmos Πυριλάμπους ? On possède quelques autres témoignages à son sujet : Calliclès, l'horrible démagogue du \italique{Gorgias} (481d), en avait fait son éromène. 
Platon souligne malicieusement que, à l'instar de l'autre dêmos que lorgne le démagogue, Calliclès finissait toujours par céder devant son cadet aimé (αἰσθάνομαι οὖν σου ἑκάστοτε, καίπερ ὄντος δεινοῦ, ὅτι ἂν φῇ σου τὰ παιδικὰ καὶ ὅπως ἂν φῇ ἔχειν, οὐ δυναμένου ἀντιλέγειν). 
Eupolis (fr. 227(213) Kassel-Austin) nous le décrit comme \italique{bouché} ; littéralement : il est plein de cérumen (κυψέλη) : 
\citComm{καὶ τῷ Πυριλάμπους ἆρα Δήμῳ κυψέλη ἔνεστιν; }%
Enfin, Athénée (IX, 397c-d), après avoir décrit le paon comme un animal bruyant qui réveille les dormeurs (τοιόνδε ταὧν, ὃς τοὺς εὕδοντας ἐγείρει) nous apprend que le fils de Pyrilampe en faisait élevage, dont il était d'ailleurs très fier...\\
Que peut-on en conclure ? assurément, que le beau Dêmos était un invétéré babillard, obtus à tout conseil, qui en remontrait à Calliclès au rebours d'une relation maître-élève classique et qui, pas plus que les paons qu'il élevait, ne s'arrêtait jamais de parler.
Comment lire maintenant le graffiti : %
\citComm{κημὸς καλός Δῆμον καλόν}% ? 
(noter l'accusatif de Δῆμον καλόν) ?
Tout simplement en conformité avec le sens premier de κημός, savoir : une \italique{muselière}.%
\citComm{Une bonne muselière [fera taire] le beau Dêmos.}%
Que la muselière (κημός), ait pu désigner aussi une \italique{urne de vote}, voilà qui devait laisser Aristophane bien rêveur !..
La pique contre la démocratie est évidente :
\citComm{Une bonne muselière [musellera] le bon peuple [pour cinq ans, pourrait-on ajouter à propos de la nôtre, encore plus décadente que celle d'Athènes].}%
Mais Philocléon dans sa naïveté -- qui serait touchante si elle n'était pas criminelle -- ne voit tout ce qu'il y a d'avilissant dans cette mascarade démocratique ; il est satisfait que le bon peuple puisse urner à son aise :
κημὸς καλός δῆμον καλόν.\\
Autre exemple, un peu plus loin dans la même pièce : Philocléon préside au procès du général Lachès\marqueur{3}, animalisé pour l'occasion en chien \italique{Chapardeur} (Λάβης). L'accusateur public n'est autre que le chien de Dêmos, Cléon. Bdélycléon fait entrer le prévenu et son accusateur ; c'est-à-dire les deux chiens, la canaille et le roquet. Réaction de Philocléon devant ces deux patibulaires (\italique{Guêpes}, vers 852-853): %
\citComm{Φιλοκλέων -- ἐς κόρακας, ὡς ἄχθομαι\\
ὁτιὴ 'πελαθόμην τοὺς καδίσκους ἐκφέρειν.\\
Aux corbeaux ! Quel ennui ! j'ai oublié d'apporter les καδίσκοι.}%
On traduit καδίσκοι par \italique{urnes} ; et il n'est pas douteux que dans la parodie de procès, les καδίσκοι soient mis en place d'\italique{urnes}... mais on voit que Philocléon frémit tout simplement devant les deux molosses qu'on lui présente et regrette l'absence de \italique{muselières}, c'est sans doute aussi le sens propre de καδίσκος.\\
%\eff{Qu'il me soit permis pour cela un petit détour par les autres comédies d'Aristophane. Commençons par le passage des \italique{Guêpes} déjà évoqué plus haut\marqueur{Voir le commentaire des vers 811-812, à propos de l'histoire de Miltiade qui détourna à son profit cinquante talents d'un coup, κατὰ πεντήκοντα τάλαντα.} dont c'est justement une scholie qui contient le fragment d'Hermippe cité au début. Ce passage est fort mécompris (vers 669-676) : Bdélycléon parle des démagogues qui se payent sur la bête athénienne -- le gros animal aurait dit Platon :%
%\citComm{κᾆθ᾽ οὗτοι μὲν δωροδοκοῦσιν κατὰ πεντήκοντα τάλαντα\\
%ἀπὸ τῶν πόλεων ἐπαπειλοῦντες τοιαυτὶ κἀναφοβοῦντες,\\
%`δώσετε τὸν φόρον, ἢ βροντήσας τὴν πόλιν ὑμῶν ἀνατρέψω.'\\
%σὺ δὲ τῆς ἀρχῆς ἀγαπᾷς τῆς σῆς τοὺς ἀργελόφους περιτρώγων.\\
%οἱ δὲ ξύμμαχοι, ὡς ᾔσθηνται τὸν μὲν σύρφακα τὸν ἄλλον\\
%ἐκ κηθαρίου λαγαριζόμενον (voir Taillardat § 130) καὶ τραγαλίζοντα τὸ μηδέν,\\ 
%σὲ μὲν ἡγοῦνται Κόννου ψῆφον, τούτοισι δὲ δωροφοροῦσιν\\
%ὕρχας, οἶνον, δάπιδας, τυρόν, μέλι, σήσαμα, προσκεφάλαια, κτλ.\\
%Eux [les démagogues] cependant extorquent aux cités des cinquante talents à la fois,\\
%en les menaçant, en les terrifiant de cette sorte :\\
%« Vous paierez le tribut, ou je tonne contre votre ville et la renverse. »\\
%Toi, quoique tu sois le chef, tu te contentes de grignoter les rognures.\\
%Et quant aux alliés, te voyant buvotant le reste de lie\\
%de la saucière et te nourrissant de rien,\\
%ils t'estiment autant que la voix de Konnos. Tandis qu'à eux ils offrent\\
%salaisons, vin, tapis, fromage, sésame, coussins, etc.}%
%Cette traduction appelle quelques explications. \\
%%Κόννου ψῆφον : par l'intermédiaire du scholiaste\marqueur{L'érudite scholie -- cela nous change de celles des \italique{Thesmophories} ! -- est de rédactions multiples et d'avis contradictoires, altérée par endroits et commentant des textes différents (que le lecteur aille consulter une édition avec apparat critique) ... aussi est-il délicat d'en piocher le bon pour en laisser le rebut. Au le lecteur qui voudra juger de mes choix, la voici telle que le manuscrit nous l'a transmise : \citComm{\eff{scholie}}}, Callistrate nous apprend l'existence d'un proverbe : Κόννου θρῖον (ἔστι δὲ παροιμία, ἐπὶ τῶν μηδὲν ἀξίων). τὸ θρῖον, \italique{la feuille de figuier}, et par extension le vain bruit qu'elle fait en brûlant, pouvait être permuté métaphoriquement avec ψόφος\marqueur{Κοννᾶς est décrit comme un vieillard divaguant au vers 536 des \italique{Cavaliers}. Cf. θρίων ψόφους et Taillardat §509, à propos de  vers 436 : ὡς ἐγὼ πολλῶν ἀκούσας οἶδα θρίων τὸν ψόφον. Mais il y a un jeu de mots non relevé : θρίων est aussi le génitif des θριαί, \italique{galets}, écho à la grève que cultive Philocléon pour n'être jamais à court de cailloux pour le vote.}. Si notre hypothèse est juste, le proverbe Κόννου ψόφος devait aussi exister, et c'est bien ce dernier que fait entendre Bdélycléon.\\
%τό κηθάριον, un hapax vu comme le diminutif de κηθίς, est au propre une \italique{saucière} pour le vinaigre\marqueur{Cf. la scholie afférente, de rédactions multiples et contradictoires, mais dont l'avis le plus pertinent semble être : ... ἣ κηθ<άρ>ια τὰ ὀξύβαφα, ἐξ ὧν τοὺς κύβους ἠ\{μ\}φίεσαν. πλεκτὰ δὲ ἦν ταῦτα. Ἄλλως. ἐκ κηθαρίου· τὰ ἐκ πέταλα τρύβλια... C'est cette interprétation qu'a retenue Hésychius : κήθεια· κηθάρια. τὰ ὀξύβαφα, ἐν οἷς τοὺς κύβους ἔβαλλον.} ; κηθίς, κήθιον ou κηθίδιον ayant pris par la suite le sens spécialisé de \italique{cornet à dés}.\\
%λαγαρίζεσθαι, c'est \italique{jouer des coudes}, peut-être \italique{gratter} ou \italique{renverser}\marqueur{Cf. Hésychius (λ 37) : λαγαριζόμενοι· σκαλεύοντες. δηλοῖ δὲ τὸ πρὸς τὰς λαγόνας [τὰ λαγαρὰ ?] τὸν ἀγκῶνα προσάγειν, πυκνὰ διατείνοντα τὴν χεῖρα. Photius (λ 15), λαγαριζόμενοι: σκαλεύοντες, τὰς λαγόνας τύπτοντες, est déjà plus vague : \italique{battre les flancs}... d'où le sens imagé de \italique{secouer} un récipient, c'est-à-dire lui remuer les flancs ? Voir la scholie au vers 674 \italique{Guêpes} ($=$ fr. 126(121) Phérécrate Kassel-Austin) qui précise : τὸ δὲ λαγαρίζεσθαι [λαγαρύζεσθαι cod.] ἕτερόν τι [autre que l'explication précédente] ἔοικεν ἐμφαίνειν. Φερεκράτης Μυρμηκανθρώποις. --- ἀντὶ τοῦ σιδαλεύην [σκαλεύειν Meineke] : ἄρα ποθ᾽ὑμεῖς καὶ τῆς ὀροφῆς τὸν χοῦν κατὰ τῆς κεφαλῆς καταμήσονται λαγαριζόμενοι. [fragment et donc d'interprétation peu sûre ; χοῦς, la \italique{conge} (unité de mesure) semble ici désigner un \italique{amas} de terre, comme la congère en français est un amas de neige.]}. Le scholiaste nous semble être dans le vrai lorsqu'il sous-entend que λαγαριζόμενον est ici un synonyme de ῥοφοῦντα\marqueur{ἐκ κηθαρίου: Τὰ ἐκ πέταλα τρύβλια, <ἃ> Ἐὐφρόνιος κήθια. μικρὸν οὖν ῥοφοῦντα ἐκ τρυβλίου καὶ τρωγαλίζοντα. Et encore : καὶ τὸ λαγαρίζειν [λαρυγίζειν cod.], τὸ ἀπο τρώγοντα θηλάζειν.}, \italique{sirotant}, \italique{buvotant}, \italique{sifflant} ; \italique{levant le coude}, dira-t-on.\\
%Bdélycléon reproche à son père un comportement de crève-la-faim : il grignote les rognures (τοὺς ἀργελόφους περιτρώγων) et infuse le reste des ordure (τὸν μὲν σύρφακα τὸν ἄλλον) dans du vinaigre qu'il sirote ensuite (ἐκ κηθαρίου λαγαριζόμενον).\\
%Pour comprendre la réaction qu'inspire une telle attitude, il faut avoir en tête une anecdote rapportée par Théophraste (Περὶ φυτικῶν ἱστορίων, IX, 17, 2-3) :%
%\citComm{Εὔδημος γοῦν ὁ φαρμακοπώλης εὐδοκιμῶν σφόδρα κατὰ τὴν τέχνην συνθέμενος μηδὲν πείσεσθαι πρὸ ἡλίου δύναντος κατέφαγε [sc. ἐλλέβορον] μέτριον πάνυ καὶ οὐ κατέσχεν οὐδ' ἐκράτησεν. ὁ δὲ Χῖος Εὔδημος\marqueur{} πίνων ἐλλέβορον οὐκ ἐκαθαίρετο. καί ποτε ἔφη πιεῖν ἐν μιᾷ ἡμέρᾳ δύο καὶ εἴκοσι πόσεις ἐν τῇ ἀγορᾷ καθήμενος ἐπὶ τῶν σκευῶν καὶ οὐκ ἐξαναστῆναι πρὸ τοῦ δείλην γενέσθαι· τότε δ' ἐλθὼν καὶ λούσασθαι καὶ δειπνεῖν ὥσπερ εἰώθει καὶ οὐκ ἐξεμέσαι· πλὴν οὗτός γε βοήθειάν τινα παρασκευασάμενος κατέσχε· \textbf{κίσσηριν} γὰρ ἐπιπάττων ἐπ' ὄξος δριμὺ πιεῖν ἔφη μετὰ τὴν ἑβδόμην πόσιν, καὶ πάλιν ὕστερον ἐν οἴνῳ τὸν αὐτὸν τρόπον.\\
%Eudême, le vendeur de drogue, qui avait une haute réputation dans son commerce, après avoir fait le pari qu'il n'éprouverait aucun effet avant le coucher du soleil, but une dose assez modérée [d'hellébore, plante purgative, cf. La Fontaine et le proverbe "se purger avec quatre grains d'hellébore"], et elle s'avéra trop forte pour son pouvoir de résistance. Mais une fois le Chiote Eudême\marqueur{La critique a vu en ce `ὁ δὲ Χῖος Εὔδημος' la mention d'un second Eudême, homonyme au premier, le φαρμακοπώλης, et qui se serait lui-aussi lancé dans des expériences pour contrer les remontées helléboriennes... fameuses coïncidences !.. il est beaucoup plus simple de considérer que c'est bien le même Eudême, mais qu'après son premier échec, Aristote se soit plu à préciser son origine, Χῖος, qui désigne l'\italique{as} du dé, c'est-à-dire le coup perdant. Apollonius le paradoxographe ne s'y était pas trompé, lui qui a compilé cette anecdote dans ses ἱστορίαι θαυμάσιαι, et qui qualifie Eudême (le manuscrit porte fautivement Εὔνυμος) de Chios de φαρμακοπώλης. Cf. l'édition de O. Keller, \italique{Rerum Naturalium Scriptores, Apollonii} Histor. mir. \italique{L}, p. 53.} a pris un bouillon d'hellébore et n'a pas été purgé. À une occasion, il a dit avoir bu en une seule journée vingt-deux potions [d'hellébore] sur la place du marché alors qu'il était assis à son étal, et n'a pas quitté l'endroit jusqu'au soir, puis il est rentré chez lui, a pris un bain et dîné comme à son habitude et n'a pas rendu. Cependant cet homme a pu tenir parce qu'il s'était muni d'un antidote ; car il dit qu'après la septième potion, il prit une gorgée de vinaigre acidulé saupoudrée de \italique{pierre ponce}, et plus tard en prit une gorgée dans du vin de la même manière.}%
%Cet Eudême de Chios, φαρμακοπώλης de son état, est mentionné ailleurs par Aristophane (\italique{Ploutos}, vers 884). L'anecdote contée par Théophraste devait être fameuse pour qu'à son époque encore il en ait connaissance. Les vertus de la pierre ponce, κίσηρις, étaient déjà connues des Anciens\marqueur{À la suite du passage cité, Théophraste parle du fait que quelque grains de pierre ponce permettent de réguler l'ébullition d'un liquide. Cette propriété est encore enseignée et utilisée dans nos lycées au laboratoires de chimie.}. Il n'est à mon avis pas fortuit qu'Aristophane juxtapose dans un de ses fragments\marqueur{fr. 332(320) Kassel-Austin $=$ Pollux, VII, 95.} κίσηρις et στρόφος, dont l'un des multiples sens est la \italique{colique}.\\
%Avec l'image que développe Bdélycléon, voyant son père tremper étrangement les rognures les plus immangeables, semblables à de petits cailloux, dans une solution vinaigrée, puis de siroter cette "potion", Philocléon devait évoquer celle bien connue du fantasque Eudême et de ses décoctions à la pierre ponce. Il nous paraît vraisemblable qu'au lieu de la leçon des manuscrits, Κόννου ψῆφον, Aristophane avait écrit Χίων ψῆφον, cf. Paix 171 et Plutarque 232f
%
%Le scholiaste a pressenti qu'Aristophane filait la métaphore avec la cubomanie : que dit-il en substance ? que Philocléon est un \italique{petit joueur}, lui qui se contente du triobole.
%(Ajoutez à cela un qu'un autre hapax -- décidément ! --, "καὶ τραγαλίζοντα τὸ μηδέν", est peut-être bien mis là pour faire entendre "κἀστραγαλίζοντα τὸ μηδέν", "jouer aux astragales pour rien"\marqueur{κυβεύειν ne désigne pas le seul fait de jouer aux dés puisqu'on le trouve employé pour des jeux ne nécessitant pas de dés. Stephen Kidd propose de voir derrière κυβεία un terme générique pour tous les jeux d'argent quand ἀστραγαλίζειν signifierait \italique{jouer pour du beurre}, cf. son article : "How to Gamble in Greek : the Meaning of Kubeia", \italique{Journal of Hellenic Studies} 137 (2017), p. 119-134.}.)\\
%C'est en ce sens qu'il faut comprendre la fin de la tirade de Bdélycléon et la réplique de Philocléon (vers 679-680) :%
%\citComm{οὐδεὶς οὐδὲ σκορόδου κεφαλὴν τοῖς ἑψητοῖσι δίδωσιν.\\
%Φιλοκλέων --- μὰ Δί᾽ ἀλλὰ παρ᾽ Εὐχαρίδου καὐτὸς τρεῖς γ᾽ ἄγλιθας μετέπεμψα.\\
%Pas un ne te donne ne fût-ce qu'une tête d'ail pour assaisonner ta poissonnaille.\\
%Philocléon --- Mais non par Zeus : }}%
%
%
%
%σφυρά: τὰ περὶ τοὺς ἀστραγάλους, αἱ περιφέρειαι τῶν ποδῶν ἢ τῶν ὀρῶν τὰ ἐπίπεδα καὶ κατώτατα ἢ ἡ παρά τισι λεγομένη κέστρα, ἰχθὺς θαλάσσιος
%σάγλη : schlechter Wurf Charis : le récipient
%se nourrissant de rien en  creuse, ils font état de toi comme du suffrage de Connos (2); tandis qu'à eux ils envoient en présent jarrons de salaison, vin, lapis, fromage, miel, sésame, coussins, coupes, manteaux, couronnes, colliers, vases à boire, richesse et santé. Et de ceux à qui tu commandes, « après tant de labeurs sur la terre et sur l'onde », pas un no te donnera pour la poissonnaille une simple tète d'ail.%
%
%
%fr. Ar. 556(543)
%\marqueur{On se trouve bien en peine de déterminer son mètre ; il y a cet énigmatique `παρὰ τὴν ξιφοθήκην' en marge du manuscrit de Ravenne. Est-ce une simple déformation de `παρὰ τὴν ψηφοθήκην' glosant κημόν ?  Ou faut-il en conclure qu'un ξιφο- se trouvait dans le texte originel des \italique{Thesmophories} ?}, 
%Κυβευτοῦ σκεύη · ἄβαξ ἀβάκιον, κόσκινον, κύβοι, διάσειστοι κύβοι, ἀστράγαλοι, φιμοί, κημοί, τηλία, κήθια κηθίδια, ψῆφοι, πεττοί.
%.
%
%Objets du joueur de dés : plateau, petit plateau, crible, dés, dés à secouer, osselets, cornets à dés, entonnoirs couvrant l’urne, table à rebord, cornets à secouer les dés, petits cornets à secouer les dés, pions, jetons.
%(Les cornets désignent aussi des muselières.)
%Il nous
% Voir dans ce vers une allusion aux activités philéliastes avec les diversLes commentateurs, à commencer par le scholiaste
%
%
%
%κήτειον:μέγαν. [καὶ (εἰς) ὃ τὰς ψήφους διωθοῦσιν ἐν τοῖς κληρωτηρίοις].καὶ ἐμετήριονἀπὸ λαχάνου ἀγρίου,ὅπερ ἀντὶ πτερῶνκαθίεσανεἰςτὸ στόμα,ὡςΚρατῖνοςἐν “Ὥραις(΄.292 Κ.-Α.) σημαίνει,καὶ ἥδυσμα. καὶ τὸ νεαρὸν χρόμμυον. [ἔνιοι μέλοςτι] οἱ δὲτὰὈψαρτυτικὰ συνθέντεςδιὰτοῦγ φασὶγήτιο (Hésychius)
%		Voir Lamer Realencyclopedia : 
%		κηθίδιον, Poll. 10.150.—Wrongly written κητίον, Alciphr. 1.39, Ath. 11.477d
%φράζε τοίνυν: ὡς ἐγώ σοι πᾶς ἀνέρριμμαι κύβος Aristophane fr. 929(673) où pas a le sens de décisif comme dans panta lithon kinein)
%
%Grenouilles 968-970 :
%Διόνυσος --- Θηραμένης; σοφός γ᾽ ἀνὴρ καὶ δεινὸς ἐς τὰ πάντα,
%ὃς ἢν κακοῖς που περιπέσῃ καὶ πλησίον παραστῇ,
%970πέπτωκεν ἔξω τῶν κακῶν, οὐ Χῖος ἀλλὰ Κεῖος.
%
%
%proskeimai : être face ontre terre? cf. Oiseaux Χίοισιν ἥσθην πανταχοῦ προσκειμένοις. + Souda Κῷος πρὸς Χῖον
%ἀστράγαλος = σφόνδυλος (vertèbres et non kondulois)
%Hésychius :  
%ἀγάς· ἡ πτῶσις τοῦ ἀστραγάλου
%ἀστραγάλας· τὰς πονηρὰς κύνας. κακοήθεις κύνες
%βουβάλιες: οἱ βούβαλοι(Ηdt.,4,192,1).ἢ ἀστράγαλοι
%γλάματα: ἀστράγαλοι
%γραῦς:τὸ συναγόμενον καὶ ἐπιπηγνύμενον ἐπάνω, ὅταν γάλα ἕψηται, ἤ τι τοιοῦτον ὑγρόν (Αρ.ΡΙυτ. 1206).καὶ ἡ ἐν τοῖς χείλεσι τῶν ποτηρίων γραμμή. καὶ βόλος τις ἀστραγάλων
%δικέντρων: βόλος τις ἀστραγάλων οὕτως ἐκαλεῖτο
%δορχελοί:ἀστράγαλοι
%Crux : ἴδα· τὸ εἰκὴν ἐν τῷ ἀστρωγαλίζειν (keinon : le 1 = Xhios ?)
%κίλλαι:ἀστράγαλοι.ἢ ὄνοι
%κόνδοι:κεραῖαι.ἀστράγαλο
%κυνοῦραι: ἀστράγαλοι
%κῳάζειν:ἀστραγαλίζειν.ἐνεχυράζειν
%κῷοι ἀστράγαλοι
%Κῶος Χῖον: ὁ Κῶος ἀστράγαλος, ὁ ἕξ. ὁ μὲν γὰρ Χῖος ἐδύνατο ἕν ὁ δὲ Κῶος ἕξ
%Crux λίγδοι:οἱ προστετριμμένοι τῶν ἀστραγάλων
%λίσποι· οἱ ἐκτετριμμένοι τῶν ἀστραγάλων. καὶ ᾿Αθηναῖοι: ἀπὸ Θησέως, ἀπὸ τοῦ ἐν Ἅιδου προσσχεθῆναι τῇ πέτρᾳ ἀπὸ τῶν γλουτῶν
%λωγάλιοι: ἀστράγαλοι.ἢ πόρνοι
%Crux μελαιαί:ἀστράγαλοι.ἢ νωθροί
%-> μήλιαι: ἀστράγαλοι (kata- ?)
%πέλλυτρα: οἱ δεσμοὶ οἱ περὶ τὰ[ς] σφυρὰ (spheiras ?) καὶ τοὺς ἀστραγάλους τῶν δρομέων περιελισσόμενοι εἰς τὸ μὴ ἐκ<σ>τραφῆναι
%πλειστοβολεῖν: παίζειν ἀστραγάλοις
%στερεαί:τινὲς τῶν δικαστικῶν ψήφων [ἄλλοι δὲ ἀστέρας astragales ?]
%σφυρά: τὰ περὶ τοὺς ἀστραγάλους, αἱ περιφέρειαι τῶν ποδῶν ἢ τῶν ὀρῶν τὰ ἐπίπεδα καὶ κατώτατα ἢ ἡ παρά τισι λεγομένη κέστρα, ἰχθὺς θαλάσσιος
%σφυρόν:[σφοδρόν.ἢ]τὸ ἐπάνω τῶν ἀστραγάλωn
%ταλαρίς: ἐξ οὗ βάλλουσι τοὺς ἀστραγάλους
%τρόπα: εἶδος παιδιᾶς, καθ᾽ἣν στρέφουσι τοὺς ἀστραγάλους εἰς τὸ ἕτερον μέρος (Krati fr 180 Κ.Ἁ.)
%ὑπόκυκλα: τοὺς ἀστραγάλους τοὺς ὑποτιθεμένους τῷ πυθμένι τῶν τριπόδων (Σ 375 ν. 1.)
%La femme qui découpe son mari ? ἀγερσικύβηλις: Κρατῖνος ἐν Δραπέτισιν (. 66 Κ.Ἁ.) ἐπὶ  Λάμπωνος. τὸν αὐτὸν ἀγύρτην καὶ κυβηλιστὴν εἶπεν, οἱονεὶ θύτην καὶ μάντιν. Κύβηλιν γὰρ ἔλεγοντὸν πέλεκυν ... ὅθεν καὶ Λύσιππος ἐν Βάκχαις (.6 Κ.Ἁ.) τὸν αὐτὸν (ὡς) ἀγύρτην κωμωδεῖ [σικυβη ?] cf. suivant
%Sur surphetwdes mendiant : ἀγύρτης: ὀχλαγωγός ,προσαίτης, ἐπαίτης, συρφετώδης, ἰδιώτης(Εὐρ. Ρες. 503). ἔστι δὲ καὶ ὁ ἀγείρων ὄχλον. καὶ ὄνομα βόλου κυβευτικοῦ (Ευβ  .57,5 Κ.Ἁ.)
%-> À PROPOS DE AÏX : αἴγλη: λαμπηδών , αὐγή ,φῶς, λαμπρότης  ... ἔστι δὲ καὶ βόλος φαῦλος κυβευτικός (Κρατ. 405 Κ Α.)
%--> κριοῦ διακονία: ὅταν προστάσσωμεν παιδίῳ διακονῆσαι, λέγομεν· δώσω σοι εἰ τύχοι ἀστραγάλους, ἢ ἰσχάδας. καὶ ὁ κριὸς οὖν εἶπε· καταλείψω τοὺς ἀστραγάλους σοι [δώσω]
%--> Konnos : Βερενίκης [Βερονίκη: πόλις] πλόκαμος: τοῦτον κατηστερίσθαι φησὶ Κόνων. καὶ βόλος δέ τις ἀστραγάλων οὕτω καλεῖται
%᾿Αργεῖος: κυβευτικοῦ βόλου ὄνομα(Ευbul.57,3Κ.-Α.)
%--> λαγαριζόμενοι: σκαλεύοντες. δηλοῖ δὲ τὸ πρὸς τὰς λαγόνας (creux de chaque côté) τὸν ἀγκῶνα (courbure => coude) προσάγειν,πυκνὰ διατείνοντα τὴν χεῖρα(Ρher.126,4Κ.-Α.)
%Photius : λαγαριζόμενοι: σκαλεύοντες, τὰς λαγόνας (les flancs) τύπτοντες.
%(Pollux II 142 : καὶ τὸ παίειν τῷ ἀγκῶνι κυβιτίζειν λέγει)
%γυλλός:κύβος,ἢ τετράγωνος λίθος
%ἐκκεκομμένος: ὁ διὰ τοῦ κυβεύειν τὰ αὑτοῦ ἀπολέσας, ὅπερ νῦν ἐκκεκοπισμένος φαμέν. ἢ ὁ ἀποκεχωρισμένος
%“ἢ τρὶς ἕξ, ἢ τρεῖς κύβοι᾽ παροιμία ἐπὶ τῶν ἡ ἐπιτυγχανόντων ἢ ἀποτυγχανόντων. πάλαι γὰρ τρισὶν ἐχρῶντο πρὸς τὰς παιδιὰς κύβοις, καὶ οὐχ, ὡς οἱ νῦν, δύο. ἔστι δὲ ὁμωνυμία
%κήθεια: κηθάρια.τὰ ὀξύβαφα, ἐν οἷς τοὺς κύβους ἔβαλλον
%κίγκασος:κυβευτικός τις βόλος
%κυβεῦσαι: παῖξαι.ταβλίσαι, χλευάσαι
%κυβευτήν: σοφὸν κυβευτήν, ἀλλὰ μὴ στένειντύχην' (Σοφ. 947) τὸν τοῖς κύβοις χρώμενον. ἐπιστήμην γάρ τινα ἐδόκουν τὸ [δὲ] δεξιῶς χρῆσθαι. ἐπιστημονικωτέραν τὴν πεττείαν τῆς κυβείας ἐνόμιζον
%Crux : (κυ?)βῆναις: γλαῦξαις
%κύων: ὁ ἐλαυνομένου τοῦ σιδήρου τοῦ ἀργοῦ ἐξ αλλόμενος σπινθήρ. οἱ δὲ τὴν Ἐρινῦν. οἱ δὲ τὴν νόσον τὴν οὕτως λεγομένην κύνα. δηλοῖ δὲ καὶ τὸ ἀνδρεῖον μόριον. καὶ τὸ ὑλακτοῦν ζῶον. καὶ τὸν ἀναιδῆ. καὶ τὸ ἄστρον.καὶ τὸ θαλάσσιον ζῶον. καὶ ὑπὸ τῶν κυβευτῶν βόλον τινά, ὃς Χῖος καλεῖται
%-> Pour Mnésiloque/andromède : κύων ἐπὶ δεσμά: παροιμία ἐπὶ τῶν ἑαυτοὺς εἰς κολάσεις ἐπιδιδόντων. λέγεται δὲ καὶ βοῦς ἐπὶ δεσμά
%κύων: παρ᾽ ἐντέροισι παροιμία ἐπὶ τῶν μὴ δυναμένων ἀπολαύειν τῶν παρακειμένων
%-> μάγνης: κυβευτικοῦ βόλου προσηγορία καὶ ἄλλων πολλῶν λεγόμενα ὀνόματα βόλων
%Μάνης: κυβευτικοῦ βόλου ὄνομα [καὶ] βαρβαρικόν,καθάπερ Εὔβουλος ἐν Κυβευταῖς(fr 59 Κ.-Α.): παρίστησι καὶ ἄλλων κατα λεγομένων βόλων <ὀνόματα>
%μίδας: κυβευτικοῦ βόλου ὄνομα. καὶ θηρίδιόν τι διεσθίον τοὺς κυάμους. καὶ ὁ πλούσιος: ἀπὸ Μίδα, τοῦ βασιλέως
%οἴνη:ἡ ἄμπελος .καὶ κυβευτικὸς λέγεται βόλος,ὁ κενὸς καὶ ἀντικείμενος τῷ ἕξ, ὃ καὶ Χῖος. οἰνίζειν γὰρ τὸ μονάζειν κατὰ γλῶσσαν
%-> Πέρσης: κυβευτικοῦ βόλου ὄνομα
%-> qui donne raison à Taillardat sur la diff petteia/kubeia (cf. aussi Aristote) : πεσσὰ πεντέ γραμμα[τα]καὶ κύβων βολαί: Σοφοκλῆς Ναυπλίῳ Πυρκαεῖ (φρ. 429 Ἀ.). παρόσον πέντε γραμμαῖς ἔπαιζον. διαφέρει δὲ πεττεία κυβείας. ἐν ἣ μὲν γὰρ τοὺς κύβους ἀναρρίπτουσιν: ἐν δὲ τῇ πεττείᾳ αὐτὸ μόνον τὰς ψήφους μετακινοῦσι
%Mais : πεττεία: διὰ κύβων παιδιά(Ρlat. Leg.820d)
%σκειράφιον: τὸ κυβευτήριον ,ἴσως διὰ τὸ ἐν Σκ[ε]ίρῳ τὴν διατριβὴν ἔχειν
%σχαδών: κυβευτικὸς βόλος
%τρὶς ἕξ:νικητήριος βόλος(Αesch. Αg. 33). καὶ παροιμία: τρὶς ἕξ ἢ τρεῖς κύβοι. ἐλέγετο δὲ ἐπὶ τοῦ εὐστόχου, ἀπὸ τῶν κυβευτῶν (Pherecr.129Κ.- Α.)
%-> Sur mesos : τροχαῖα: μέσα ἐν κύβοις. ἢ ὁδός, ὡς Ῥίνθων (fr 23 Κ.-Α.)
%χαλκίνδα: τὸ εἰς χαλκὸν κυβεύειν
%ἀναθέσθαι: ὁμολογῆσαι ἐπὶ τῶν πεσσῶν ἐλέγετο: οἱ γὰρ παίζοντες καὶ τὰς ψήφους διώρθουν
%ἀρτοπόπος: ὁ πέσσων ἐν ἐργαστηρίῳ
%Crux ἄρτος:βόλος τις
%-> sur surfax :  παλινδορία: τὸ σκῦτος. Πλάτων Σύρφαχι (φρ. 180Κ.Ἀ.):  'σὲ μὲν, ὦ μοχθηρέ, παλινδορίᾳ παίσας αὐτοῦ καταθήσω᾽
%ET : ἐνκατασυρφάζειν: ἐνυβρίζειν
%ἐς βόθυνον: εἶδος παιδιᾶς παρὰ Κρατίνῳ (fr. 219 Κ.-Α.)
%-> έκδότω δέ τις 
%καί ψηφολόγιον ώδε καί δίφρω δύο (Ar. fr. ...)
%Diphilus 77 K.-A. : ἕλκ᾽ἐς μέσον τὸν φιμὸν ὡς ἂν ἐμβάλῆ
%Photius : Τρὶς ἓξ ἢ τρεῖς κύβοι: οἱ μὲν τρὶς ἓξ νίκην· οἱ δὲ τρεῖς κύβοι κενοί· καὶ τὸ παρ’ Εὐριπίδηι τοιοῦτον·   (10) Βέβληκ’ Ἀχιλλεὺς δύο κύβω καὶ τέτταρα· τριῶν γὰρ ὄντων τῶν ἀναρριπτουμένων βόλων, δύο μὲν κενοὺς αὐτὸν φησὶ βάλλειν· ἕνα δὲ τὸν τέταρτον· διττὸς οὖν ὁ κύβος ὅ, τε ἀναρριπτούμενος, εἴτε κενὸς, εἴτε πλήρης εἴη· καὶ ἰδίως ὁ κενός· ὅτι δὲ ὁ    (15) τρὶς ἓξ βαλὼν κατώρθου, καὶ Αἰσχύλος ἐν Ἀγα- μέμνονι παρίστησι· Τὰ δεσποτῶν γὰρ εὖ πεσόντα θήσομαι, Τρὶς ἓξ βαλούσης τῆς ἐμῆς φρυκ- τωρίας.
%Λάκωνες: κυβευτικοῦ βόλου ὄνομα. Phot
%Ploutos 816 artiasmos (pair ou impair)
%σκιραφεῖα ̓ τὰ κυβεῖα (les gains ?) Photius
%σκιρζάφιτζα ̓ τὰ κυβευτήριά φησιν ̓ἐπειδὴ διέτριβον ἐν Σκίρῳοἱ κυβεύοντες, ὡς Θεόπομπος(FrgHist 115F228)
%Photius : ἐκκόψαι· ἀντὶ τοῦ νικῆσαι. οὕτως Ἄλεξις
%
%Schol. Vesp. 674 (il y a un unique psêfoforwn !) avec cette citation : καὶ πρὸς κύβους ἕστηκ' ἔχων τὸ κήθιον. (Hermippe, fr. 27 K.-A.) :
%ἐκ κηθαρίου: Σύρφακα. τὸ ἱκανὸν πλῆθος τῶν δικαστῶν τὸ συρφετῶδες, ἐξ εὐτελῶν τρεφόμενον. 
%κηθάριον γὰρ πλέγμα ἐστὶ κανισκῶδες [ἐστὶ κοσκινῶδες Dobrée], ἐπιτιθέμενον τῇ κληρωτ<ρ>ίδι τῶν ψήφων. λαγαριζόμενον δὲ τὰ λάγαρα ἐσθίοντα, ὅ ἐστιν εὔθρα<υ>στα καὶ εὐτελῆ τινα. ἣ κηθ<άρ>ια τὰ ὀξύβαφα, ἐξ ὧν τοὺς κύβους ἠ\{μ\}φίεσαν. πλεκτὰ δὲ ἦν ταῦτα. Ἄλλως. ἐκ κηθαρίου: Τὰ ἐκ πέταλα τρύβλια, <ἃ> Ἐὐφρόνιος κήθια. μικρὸν οὖν ῥοφοῦντα ἐκ τρυβλίου καὶ τρωγαλίζοντα, οἷον κυάμους καὶ ὅσα τοιαῦτα. οὐκ ἔστι δὲ, ἀλλ᾽ἀγγεῖον πλεκτὸν, εἰς ὃ τὰς ψήφους καθιᾶσι, τὸ κήθιον, οἷον ὡς ἀπὸ τοῦ ψηφηφορεῖν μόνον. καὶ Χαῖρις δὲ τὸ κήθιον τρυβλίου τι εἶδος. καὶ τὸ λαγαρίζειν [λαρυγίζειν V.], τὸ ἀπο τρώγοντα θηλάζειν. μήποτε μέντοι καὶ πρὸς τὴν τῶν κύβων παιδιὰν ὁμοίως ἐχρῶντό τινι βάλλοντες ἐξ αὐτοῦ. καὶ γὰρ Ἕρμιππος ἐν Θεοῖς 
%καὶ πρὸς κύβους ἕστηκ' ἔχων <τὸ> κήθιον.
%τὸ δὲ λαγαρίζεσθαι [λαγαρύζεσθαι V. Ald.] ἕτερόν τι ἔοικεν ἐμφαίνειν.
%Φερεκράτης Μυρμηκανθρώποις. --- ἀντὶ τοῦ σιδαλεύειν [ σκαλεύειν ;] ἄρα ποθ᾽ὑμεῖς καὶ τῆς ὀροφῆς τὸν χοῦν κατὰ τῆς κεφαλῆς καταμήσονται λαγαριζόμενοι. Υ. --- ἴσως ἂν οὖν λέγοι ἦτοι ἐκ τῶν ἀγγείων τῆς ψήφου ἀποξύοντα τὸν μισθὸν καὶ ἐξ ἄγγους τινὸς πλεχτοῦ, οἷον ὁ κάρταλος, τρωγαλιζόμενον, ὡς ἂν πένητα. λαγαριζόμενον δὲ ἀντὶ τοῦ ἀποξύοντα, ὡς ἀπὸ τῶν δικαστικῶν ἀγγείων. ταῦτα δὲ ὡς περὶ πενήτων.
%Surfax : 
%
%
%
%
%(1) Le sens de cette locution a été contesté. Il me semble pleinement confirmé par le vers 562 du Ploutos. 
%(2) 11 s'agit de ce pauvre diable de Connos, l'ancien maître de musique de Socrate, dont il a été question dans les Cavaliers (v. 534) et qui certes ne s'inquiétait guère de voter. (Peut-etre allusion à Koos aux dés ?)
%
%
%SUR KÊMOS : 
%
%
%Acharniens :
%Χορός
%ἀπολεῖς ἄρ᾽ ὁμήλικα τόνδε φιλανθρακέα;
%Δικαιόπολις
%οὐδ᾽ ἐμοῦ λέγοντος ὑμεῖς ἀρτίως ἠκούσατε.
%Χορός
%ἀλλὰ γὰρ νῦν λέγ᾽, εἴ σοι δοκεῖ, τόν τε Λακεδαιμόνιον
%† αὐτὸν ὅτι τῷ † τρόπῳ σοὐστὶ φίλος:
%340ὡς τόδε τὸ λαρκίδιον οὐ προδώσω ποτέ.
%Δικαιόπολις
%τοὺς λίθους νύν μοι χαμᾶζε πρῶτον ἐξεράσατε.
%Χορός
%οὑτοιί σοι χαμαί, καὶ σὺ κατάθου πάλιν τὸ ξίφος.
%Δικαιόπολις
%ἀλλ᾽ ὅπως μὴ ν᾽ τοῖς τρίβωσιν ἐγκάθηνταί που λίθοι.
%Χορός
%ἐκσέσεισται χαμᾶζ᾽: οὐχ ὁρᾷς σειόμενον;
%345ἀλλὰ μή μοι πρόφασιν, ἀλλὰ κατάθου τὸ βέλος.
%ὡς ὅδε γε σειστὸς ἅμα τῇ στροφῇ γίγνεται.
%
%Equites : 
%
%ἐν τῇ πυκνί, κᾆθ’ ὅταν
%μή σοι τύχῃ ὄψον ὄν, @1
%τούτων ὃς ἂν ᾖ παχύς,
%θύσας ἐπιδειπνεῖς.   (1140)
%ΔΗ.                       Σκέψασθε δέ μ’, εἰ σοφῶς
%αὐτοὺς περιέρχομαι
%τοὺς οἰομένους φρονεῖν
%κἄμ’ ἐξαπατύλλειν.
%Τηρῶ γὰρ ἑκάστοτ’ αὐ-   (1145)
%τοὺς οὐδὲ δοκῶν ὁρᾶν
%κλέπτοντας· ἔπειτ’ ἀναγ-
%κάζω πάλιν ἐξεμεῖν
%ἅττ’ ἂν κεκλόφωσί μου, (keklotosi : secouer [le cornet à dés] ?)
%κημὸν καταμηλῶν.    (1150)  (mêlos : la luette (pour se faire vomir) ? kêmmmon ici la muselière ? pour Willems, elle désigne l'enquête judiciaire + jeu de mots : cf. schol. guêpes 674 : τὸν χοῦν κατὰ τῆς κεφαλῆς καταμήσονται λαγαριζόμενοι :  en renversant l'urne ?)
%Ar. fr. 625(614) K.-A. : τὴν φάρυγα μηλῶν δύο δραχμὰς ἕξει μόνας
%Phot. = Suid. μ 940 μηλῶσαι· τὸ καθεῖναί τι εἰς βάθος. καὶ τὴν φάρυγγα μηλῶσαι,τὸ διαχρῖςαι (-ῆσαι Phot.) τῷ δακτύλῳ.᾿Αριsτοφάνης· τὴν ----- μόνας
%
%
%ΠΑ.                 Ἄπαγ’ ἐς μακαρίαν ἐκποδών.    (1150)
%ΑΛ.                                            Σύ γ’, ὦ φθόρε.    (1151)
%ΠΑ.                 Ὦ Δῆμ’, ἐγὼ μέντοι παρεσκευασμένος
%τρίπαλαι κάθημαι βουλόμενός σ’ εὐεργετεῖν.
%ΑΛ.                 Ἐγὼ δὲ δεκάπαλαι γε καὶ δωδεκάπαλαι
%
%
%Vespae : 
%
%C'est aussi avec le double sens de \italique{muselière} qu'il faut prendre κημός au vers 99 des \italique{Guêpes} : Dêmos le fils de Pyrilampe était si beau que Calliclès, le démagogue du \italique{Gorgias} (481d), en avait fait son éromène. Mais Dêmos était aussi un obtus babillard : Platon le souligne malicieusement : Calliclès, malgré qu'il en ait, finissait toujours par céder devant son cadet aimé (αἰσθάνομαι οὖν σου ἑκάστοτε, καίπερ ὄντος δεινοῦ, ὅτι ἂν φῇ σου τὰ παιδικὰ καὶ ὅπως ἂν φῇ ἔχειν, οὐ δυναμένου ἀντιλέγειν). 
%Voilà pourquoi Eupolis (fr. 227(213) Kassel-Austin) nous le décrit comme \italique{bouché} ; littéralement : il est plein de cérumen (κυψέλη) : καὶ τῷ Πυριλάμπους ἆρα Δήμῳ κυψέλη ἔνεστιν;
%Il n'est donc pas indifférent qu'Athénée (IX, 397c-d), après avoir décrit le paon comme un animal bruyant qui réveille les dormeurs (τοιόνδε ταὧν, ὃς τοὺς εὕδοντας ἐγείρει) nous apprenne que le fils de Pyrilampe en faisait élevage, dont il était très fier.
%On comprend mieux les raisons -- outre sa philéliastite -- qui poussèrent Philocléon à rajouter (παρέγραψε πλησίον) : "κημὸς καλός" à côté de "Δῆμον καλόν" (noter l'accusatif) : une bonne muselière fera utilement taire le beau Démos.
%φιληλιαστής ἐστιν ὡς οὐδεὶς ἀνήρ·
%ἐρᾷ τε τούτου τοῦ δικάζειν, καὶ στένει
%ἢν μὴ ’πὶ τοῦ πρώτου καθίζηται ξύλου.   (90)
%ὕπνου δ’ ὁρᾷ τῆς νυκτὸς οὐδὲ πασπάλην·
%ἢν δ’ οὖν καταμύσῃ κἂν ἄχνην, ὅμως ἐκεῖ
%ὁ νοῦς πέτεται τὴν νύκτα περὶ τὴν κλεψύδραν.
%ὑπὸ τοῦ δὲ τὴν ψῆφόν γ’ ἔχειν εἰωθέναι
%τοὺς τρεῖς ξυνέχων τῶν δακτύλων ἀνίσταται,   (95)
%ὥσπερ λιβανωτὸν ἐπιτιθεὶς νουμηνίᾳ.
%καὶ νὴ Δί’ ἢν ἴδῃ γέ που γεγραμμένον
%υἱὸν Πυριλάμπους ἐν θύρᾳ Δῆμον καλόν, 
%ἰὼν παρέγραψε πλησίον “κημὸς καλός”.
%τὸν ἀλεκτρυόνα δ’, ὃς ᾖδ’ ἀφ’ ἑσπέρας, ἔφη @1    (100)
%ὄψ’ ἐξεγείρειν αὐτὸν ἀναπεπεισμένον,
%παρὰ τῶν ὑπευθύνων ἔχοντα χρήματα·
%εὐθὺς δ’ ἀπὸ δορπηστοῦ κέκραγεν ἐμβάδας,
%κἄπειτ’ ἐκεῖσ’ ἐλθὼν προκαθεύδει πρῲ πάνυ,
%
%Non, par le chien, ô Nicostratos, car Pliiloxène est un ruffian. — Vaine dépense de paroles, vous ne trouverez pas. Si vous tenez à savoir, faites silence, je vais vous expliquer la maladie du maître. Il a le goût de l'Héliée, comme personne. Son envie, ce n'est que cela : être dicaste ; et il se lamente s'il ne siège pas sur le premier banc. De sommeil, il n'en goûte pas un brin de la nuit, ou s'il ferme l'œil, fût-ce un instant, son esprit voltige là-bas, nuitamment, autour de la clepsydre. Oui, habitué qu'il est à tenir le suffrage, il se lève en serrant les trois doigts, comme s'il mettait l'encens sur l'autel à la nouvelle lune. Et, par Zeus, voit-il quelque part inscrit sur une porte le nom du fils de Pyrilampe: « Bravo, Démos! » (2), il s'en va écrire à côté: « Bravo, Kêmos! ». Le coq, qui chantait à l'entrée de la nuit (4), il prétendit qu'il s'était laissé séduire
%
%
%(2) L'épithète kalos, ajoutée à un nom, se rencontre fréquemment sur les vases peints et ne signifie nullement beau, pas plus que dans l'expression kalos xàgagôc, littér. bon et brave Elle correspond exactement au bravo des Italiens, et se dit pour marquer qu'on fait grand cas d'un homme ou d'une chose. On pourrait aussi traduire : « Vive Démos! » 
%
%
%κριοῦ διακονία: ὅταν προστάσσωμεν παιδίῳ διακονῆσαι, λέγομεν· δώσω σοι εἰ τύχοι ἀστραγάλους, ἢ ἰσχάδας. καὶ ὁ κριὸς οὖν εἶπε· καταλείψω τοὺς ἀστραγάλους σοι
%Prov. : Κόννου θρῖον (mis pour ἰσχάς ), bagatelle,  θρίου ψόφος (Vespae 436 : menaces vides)
%Postérieur : Βερενίκης [Βερονίκη: πόλις] πλόκαμος: τοῦτον κατηστερίσθαι φησὶ Κόνων. καὶ βόλος δέ τις ἀστραγάλων οὕτω καλεῖται
%Φιλοκλέων
%καὶ ποῖ τρέπεται δὴ 'πειτα τὰ χρήματα τἄλλα;
%Βδελυκλέων
%ἐς τούτους τοὺς ‘οὐχὶ προδώσω τὸν Ἀθηναίων κολοσυρτόν,
%ἀλλὰ μαχοῦμαι περὶ τοῦ πλήθους ἀεί.’ σὺ γὰρ ὦ πάτερ αὐτοὺς
%ἄρχειν αἱρεῖ σαυτοῦ τούτοις τοῖς ῥηματίοις περιπεφθείς.
%κᾆθ᾽ οὗτοι μὲν δωροδοκοῦσιν κατὰ πεντήκοντα τάλαντα
%670ἀπὸ τῶν πόλεων ἐπαπειλοῦντες τοιαυτὶ κἀναφοβοῦντες,
%‘δώσετε τὸν φόρον, ἢ βροντήσας τὴν πόλιν ὑμῶν ἀνατρέψω.’
%σὺ δὲ τῆς ἀρχῆς ἀγαπᾷς τῆς σῆς τοὺς ἀργελόφους περιτρώγων.
%οἱ δὲ ξύμμαχοι ὡς ᾔσθηνται τὸν μὲν σύρφακα τὸν ἄλλον
%ἐκ κηθαρίον λαγαριζόμενον καὶ τραγαλιζοντα τὸ μηδέν,
%675σὲ μὲν ἡγοῦνται Κόννου ψῆφον, τούτοισι δὲ δωροφοροῦσιν (sur le verbe hegoumai, voir Lamer 1939)
%ὕρχας οἶνον δάπιδας τυρὸν μέλι σήσαμα προσκεφάλαια
%φιάλας χλανίδας στεφάνους ὅρμους ἐκπώματα πλουθυγιείαν:
%σοὶ δ᾽ ὧν ἄρχεις, πολλὰ μὲν ἐν γῇ πολλὰ δ᾽ ἐφ᾽ ὑγρᾷ πιτυλεύσας,
%οὐδεὶς οὐδὲ σκορόδου κεφαλὴν τοῖς ἑψητοῖσι δίδωσιν.
%
%
%
%
%Φι.        ἰώ μοί μοι.   (750)
%Βδ.        οὗτος, τί βοᾷς;
%Φι.        μή μοι τούτων μηδὲν ὑπισχνοῦ.
%κείνων ἔραμαι, κεῖθι γενοίμαν,
%ἵν’ ὁ κῆρύξ φησι· “τίς ἀψήφιστος; ἀνιστάσθω.”  (Hésych. : ἄψηφον: πολύ. μέγα.ἰσχυρόν)
%κἀπισταίην ἐπὶ τοῖς κημοῖς    (755)
%ψηφιζομένων ὁ τελευταῖος. (de la dernière chance ? cf. Taillardat p. 153 n.1)
%“σπεῦδ’, ὦ ψυχή—”. ποῦ μοι ψυχή; “πάρες, ὦ σκιερά—”. (..fon : c'est le lieu de jeu ? voir Lamer 1936... ou un élément de Jeu ? (taillardat))
%μὰ τὸν Ἡρακλέα μή νυν ἔτ’ ἐγὼ ’ν τοῖσι δικασταῖς
%κλέπτοντα Κλέωνα λάβοιμι.
%Βδ.        ἴθ’, ὦ πάτερ, πρὸς τῶν θεῶν ἐμοὶ πιθοῦ.   (759)
%
%
%
%       ἐγὼ γὰρ ἀπεδάκρυσα νῦν γνώμην ἐμὴν
%οὐδέν ποτέ γ’ ἀλλ’ ἢ τῆς φακῆς ἐμπλήμενος.
%Βδ.        οὔκουν ἀποφεύγει δῆτα;    (985)
%Φι.                              χαλεπὸν εἰδέναι.   (985)
%Βδ.        ἴθ’, ὦ πατρίδιον, ἐπὶ τὰ βελτίω τρέπου.
%τηνδὶ λαβὼν τὴν ψῆφον ἐπὶ τὸν ὕστερον
%μύσας παρᾷξον κἀπόλυσον, ὦ πάτερ.
%Φι.        οὐ δῆτα· κιθαρίζειν γὰρ οὐκ ἐπίσταμαι.
%Βδ.        φέρε νύν σε τῃδὶ τὴν ταχίστην περιάγω.    (990)
%Φι.        ὅδ’ ἔσθ’ ὁ πρότερος;
%Βδ.                            οὗτος.    (991)
%
%
%Βδ.        ἀπολεῖς με. ποίας χάρακας; ἀλλ’ ὡς ἢ κάπρον
%ἐδιώκαθές ποτ’ ἢ λαγών, ἢ λαμπάδα
%ἔδραμες, ἀνευρὼν ὅ τι νεανικώτατον.
%Φι.        ἐγᾦδα τοίνυν τό γε νεανικώτατον· @1   (1205)
%ὅτε τὸν δρομέα Φάυλλον ὢν βούπαις ἔτι
%εἷλον διώκων λοιδορίας ψήφοιν δυοῖν.
%Βδ.        παῦ’· ἀλλὰ δευρὶ κατακλινεὶς προσμάνθανε
%ξυμποτικὸς εἶναι καὶ ξυνουσιαστικός.
%Φι.        πῶς οὖν κατακλινῶ; φράζ’ ἁνύσας.    (1210)
%Βδ.                                        εὐσχημόνως.   (1210)
%Φι.        ὡδὶ κελεύεις κατακλινῆναι;
%
%
%
%ἁθρόοι γὰρ ἥξομέν σε προσκαλούμενοι.
%Φι.        ἰὴ ἰεῦ, “καλούμενοι”.   (1335)
%ἀρχαῖά γ’ ὑμῶν. ἆρά γ’ ἴσθ’
%ὡς οὐδ’ ἀκούων ἀνέχομαι
%δικῶν; ἰαιβοῖ αἰβοῖ.
%τάδε μ’ ἀρέσκει· βάλλε κημούς.
%οὐκ ἄπει; ποῦ ’στ’ ἠλιαστής; ἐκποδών.   (1340)
%ἀνάβαινε δεῦρο, χρυσομηλολόνθιον,  (il y a un jeu qui un nom proche)
%τῇ χειρὶ τουδὶ λαβομένη τοῦ σχοινίου.
%ἔχου· φυλάττου δ’, ὡς σαπρὸν τὸ σχοινίον·
%ὅμως γε μέντοι τριβόμενον οὐκ ἄχθεται.
%
%
%
%Ranea : compter avec des pierres ?
%
%       ΔΙ.                 Καὶ μὴν λογιοῦμαι ταῦτα τῶν ψήφων λαβών.
%ΕΥ.               Φθιῶτ’ Ἀχιλλεῦ, τί ποτ’ ἀνδροδάικτον ἀκούων @1
%ἰὴ κόπον οὐ πελάθεις ἐπ’ ἀρωγάν;
%Ἑρμᾶν μὲν πρόγονον τίομεν γένος οἱ περὶ λίμναν.   (1265)
%Ἰὴ κόπον οὐ πελάθεις ἐπ’ ἀρωγάν;
%ΔΙ.                   Δύο σοι κόπω, Αἰσχύλε, τούτω.
%
%LSJ lagara : τὴν πόλιν ἀντὶ λαγαρᾶς καὶ ὑποσόμφου μεστὴν ἐποίησεν ἀγλαΐας Them. Or. 18.222d
%
%Platon Lysis 206e (sur l'artiasmos) : Ταῦτα, ἦν δ' ἐγώ, χρὴ ποιεῖν. Καὶ ἅμα λαβὼν τὸν [206e] Κτήσιππον προσῇα εἰς τὴν παλαίστραν· οἱ δ' ἄλλοι ὕστεροι ἡμῶν ᾖσαν. Εἰσελθόντες δὲ κατελάβομεν αὐτόθι τεθυκότας τε τοὺς παῖδας καὶ τὰ περὶ τὰ ἱερεῖα σχεδόν τι ἤδη πεποιημένα, ἀστραγαλίζοντάς τε δὴ καὶ κεκοσμημένους ἅπαντας. Οἱ μὲν οὖν πολλοὶ ἐν τῇ αὐλῇ ἔπαιζον ἔξω, οἱ δέ τινες τοῦ ἀποδυτηρίου ἐν γωνίᾳ ἠρτίαζον ἀστραγάλοις παμπόλλοις, ἐκ φορμίσκων τινῶν προαιρούμενοι· τούτους δὲ περιέστασαν ἄλλοι θεωροῦντες. Ὧν δὴ καὶ ὁ Λύσις ἦν, καὶ εἱστήκει ἐν [207a] τοῖς παισί τε καὶ νεανίσκοις ἐστεφανωμένος
%
%Vesp. 74
%
%
%Ploutos :
%
%Χρεμύλος
%555ὡς μακαρίτην ὦ Δάματερ τὸν βίον αὐτοῦ κατέλεξας,
%εἰ φεισάμενος καὶ μοχθήσας καταλείψει μηδὲ ταφῆναι.
%Πενία
%σκώπτειν πειρᾷ καὶ κωμῳδεῖν τοῦ σπουδάζειν ἀμελήσας,
%οὐ γιγνώσκων ὅτι τοῦ Πλούτου παρέχω βελτίονας ἄνδρας
%καὶ τὴν γνώμην καὶ τὴν ἰδέαν. παρὰ τῷ μὲν γὰρ ποδαγρῶντες
%560καὶ γαστρώδεις καὶ παχύκνημοι καὶ πίονές εἰσιν ἀσελγῶς,
%παρ᾽ ἐμοὶ δ᾽ ἰσχνοὶ καὶ σφηκώδεις καὶ τοῖς ἐχθροῖς ἀνιαροί.
%Χρεμύλος
%ἀπὸ τοῦ λιμοῦ γὰρ ἴσως αὐτοῖς τὸ σφηκῶδες σὺ πορίζεις.
%Πενία
%περὶ σωφροσύνης ἤδη τοίνυν περανῶ σφῷν κἀναδιδάξω
%ὅτι κοσμιότης οἰκεῖ μετ᾽ ἐμοῦ, τοῦ Πλούτου δ᾽ ἐστὶι ὑβρίζειν.
%Χρεμύλος
%565πάνυ γοῦν κλέπτειν κόσμιόν ἐστιν καὶ τοὺς τοίχους διορύττειν.
%Βλεψίδημος
%†νὴ τὸν Δί᾽, εἰ δεῖ λαθεῖν αὐτόν, πῶς οὐχὶ κόσμιόν ἐστι;†
%Πενία
%σκέψαι τοίνυν ἐν ταῖς πόλεσιν τοὺς ῥήτορας, ὡς ὁπόταν μὲν
%ὦσι πένητες, περὶ τὸν δῆμον καὶ τὴν πόλιν εἰσὶ δίκαιοι,
%πλουτήσαντες δ᾽ ἀπὸ τῶν κοινῶν παραχρῆμ᾽ ἄδικοι γεγένηνται,
%570ἐπιβουλεύουσί τε τῷ πλήθει καὶ τῷ δήμῳ πολεμοῦσιν.
%
%
%Est-ce la TROPA qui est décrite ?? 
%
%Pollux IX 103 :
%τὸ μὲν οὖν ἀστραγάλοις παίζειν καὶ ἀστραγαλίζειν καὶ ἀστερίζειν ἔνιοι τῶν ποιητῶν εἰρήκασιν, ὅτι τοὺς ἀστραγάλους καὶ ἀστρίας εἰσὶν οἳ ὠνόμαζον, Ἀντιφάνης δὲ καὶ ἀστρίχους. 100 Τὸ δὲ σχῆμα τοῦ κατὰ τὸν ἀστράγαλον πτώματος ἀριθμοῦ δόξαν εἶχεν, καὶ τὸ μὲν μονάδα δηλοῦν καλεῖται κύων, τὸ δὲ ἀντικείμενον χιάς, καὶ Χῖος οὗτος ὁ βόλος. Δυὰς δὲ καὶ πεντὰς ἐν ἀστραγάλοις, ὥσπερ ἐν κύβοις, οὐκ ἔνεστιν. Οἱ δὲ πλείους τὸν μὲν ἑξίτην Κῷον, τὸν δὲ κύνα Χῖον καλεῖσθαι λέγουσιν. Καὶ μὴν καὶ Στησίχορος ἐκαλεῖτό τις παρὰ τοῖς ἀστραγαλίζουσιν ἀριθμός, ὃς ἐδήλου τὰ ὀκτώ· τὸν γὰρ ἐν Ἱμέρᾳ τοῦ ποιητοῦ τάφον ἐξ ὀκτὼ πάντων συντεθέντα πεποιηκέναι τὴν 'πάντ'ὀκτώ΄ φασι παροιμίαν9. 101 Ἐπεὶ δὲ τοῖς τετταράκοντα τοῖς μετὰ τοὺς τριάκοντα προστᾶσιν Ἀθήνησι συνῆρξεν Εὑριππίδης, εἰ τετταράκοντα συνήθροιζεν ἀστραγάλων βολή, τὸν ἀριθμὸν τοῦτον Εὑριππίδην ὠνόμαζον. Καὶ μὴν καὶ ἀρτιάζειν ἀστραγάλους ἐκ φορμίσκων κατερωμένους ἐν τῷ ἀποδυτηρίῳ τοὺς παῖδας ὁ Πλάτων (Lys. 206Ε) ἔφη· τὸ δ΄ἀρτιάζειν ἐν ἀστραγάλων πλήθει, κεκρυμμένων ὑπὸ ταῖν χεροῖν, μαντεῖαν εἶχε τῶν ἀρτίων ἢ καὶ περιττῶν. Τἀυτὸν δὲ τοῦτο καὶ κυάμοις ἢ καρύοις ἢ ἀμυγδάλαις, οἱ δὲ καὶ ἀργυρίῳ πράττειν ἠξίουν, εἰ πιστὸς Ἀριστοφάνης ἐν Πλούτῳ (816) λέγων· στατῆρσι δ'οἱ θεράποντες ἀρτιάζομεν. 102 Εἰ μὲν οὖν κύκλου περιγραφέντος ἀφιέντες ἀστράγαλον ἐστοχάζοντο τοῦ μεῖναι τὸν βληθέντα ἐν τῷ κύκλῳ, ταύτην εἰς ὤμιλλαν τὴν παιδιὰν ὠνόμαζον. […] 103 ἡ δὲ τρόπα καλουμένη παιδιὰ γίνεται μὲν ὡς τὸ πολὺ δι΄ἀστραγάλων, οὓς ἀφιέντες στοχάζονται βόθρου τινὸς εἰς ὑποδοχὴν τῆς τοαιύτης ῥίψεως ἐξεπίτηδες πεποιημένου· πολλάκις δὲ καὶ ἀκύλοις καὶ βαλάνοις ἀντὶ τῶν ἀστραγάλων οἱ ῥίπτοντες ἐχρῶντο.
%Alcuni poeti hanno chiamato anche astrizein il fatto di giocare agli astragali, perché chiamano gli astragali astragaloi e astriai10, Antifane dice anche astrichoi. 100 Lo schema secondo la maniera di cadere dell’astragalo ha un valore numerico e il lato dell’unità si chiama il cane (kyôn), quello opposto chias e questo colpo Chio. Non ci sono valori di due o di cinque negli astragali e nemmeno nei dadi. I più dicono che si dà alla faccia che vale 6 il nome di Koos e al cane quello di Chios. E un numero era anche detto Stesicoro dai giocatori agli astragali, era l´otto. Si dice, infatti, che la tomba del poeta a Imera era un assemblaggio di otto parti, donde il proverbio: «tutto fa otto.» 101 Poiché Eurippide dirigeva i Quaranta che governavano ad Atene dopo i Trenta, se il tiro di astragali dava in totale 40, era detto Eurippide. E Platone ha detto che i bambini giocavano a pari e dispari nel vestiario dei bagni con gli astragali che prendevano nei bussolotti. Giocare a pari e dispari sulla somma degli astragali nascosti nelle mani era la divinazione dei numeri pari e dispari. E alcuni facevano lo stesso con fave, noci, mandorle o monete, se si dà credito ad Aristofane, che dice nel Pluto: «Noialtri, gli schiavi giochiamo a pari e dispari con gli stateri». 102 Se si descrive un cerchio, lanciando un astragalo con tanta abilità per farlo restare nel cerchio, questo gioco prende nome di ômilla. [...] 103 Il gioco detto tropa si fa prevalentemente con gli astragali, che vengono lanciati nello sforzo di raggiungere una fossetta scavata in modo da poter contenere tutti gli astragali che vi sono gettati dentro. Sovente quelli che lanciano si servono di ghiande e castagne invece che di astragali. (traduzione mia)
%
%
%
%Sur les faces de l'osselet :
%Aristotele nella nota discussione anatomica dell’Historia Animalium II 1 499b 26-31: 
%Πάντα δὲ τὰ ἔχοντα ἀστραγάλους ἐν τοῖς ὄπισθεν ἔχει σκέλεσιν. Ἔχει δ’ ὀρθὸν τὸν ἀστράγαλον ἐν τῇ καμπῇ, τὸ μὲν πράνης ἔξω, τὸ δ'ὕπτιον εἴσω, καὶ τὰ μὲν κῷα ἐντὸς ἐστραμμένα πρὸς ἄλληλα, τὰ δὲ χῖα καλούμενα ἔξω, καὶ τὰς κεραίας ἄνω. Ἡ μὲν οὖν θέσις τῶν ἀστραγάλων τοῖς ἔχουσι πᾶσι τοῦτον ἔχει τὸν τρόπον. 
%Tutti gli animali che hanno gli astragali li presentano nelle zampe posteriori. Hanno l’astragalo posto verticalmente nell’articolazione, la parte superiore posizionata verso l’esterno, quella inferiore verso l’interno; i lati di Coo sono rivolti all’interno uno verso l’altro; i lati che sono detti di Chio verso l’esterno, e le parti con le corna verso l’alto. La posizione degli astragali è dunque di tal sorta in tutti gli animali che ne siano provvisti. (traduzione mia)
%
%Ancora nel De caelo II, 12, 292a 28-292b 1 una similitudine sulla legge delle probabilità e dei grandi numeri è desunta dal gioco degli astragali: 
%Ἔστι δὲ τὸ κατορθοῦν χαλεπὸν ἢ τὸ πολλὰ ἢ τὸ πολλάκις, οἷον μυρίους ἀστραγάλους Χίους βαλεῖν ἀμήχανον, ἀλλ’ ἕνα ἢ δύο ῥᾷον. Καὶ πάλιν ὅταν τοδὶ μὲν δέῃ τοῦδ’ ἕνεκα ποιῆσαι, τοῦτο δ’ ἄλλου καὶ τοῦτο ἑτέρου, ἐν μὲν ἑνὶ ἢ δυσὶ ῥᾴδιον ἐπιτυχεῖν, ὅσῳ δ’ ἂν διὰ πλειόνων, χαλεπώτερον. 
%È difficile, peraltro, avere un buon esito in molte cose o di frequente, così per esempio agli astragali non riesce di fare innumerevoli volte il colpo di Chio, mentre è facile farlo una volta o due. E ancora di nuovo, quando si deve fare una tal cosa a un tal fine, la talaltra ad un secondo fine ed un’altra ancora ad un fine ancora differente, è facile riuscirvi in un caso o due, ma ciò è tanto più difficile, quanto più si moltiplicano gli impegni da assolvere. (traduzione mia)
%Dunque, per riassumere i quattro lati dell’astragalo di forma pressoché cubica ricevono differenti denominazioni e punteggi che possiamo schematizzare (cfr. Nollé 2007: 8-9): 
%1) Χῖος, Chios, o κύων, cane, lato stretto a forma di S): 1 punto; 
%2) ὕπτιος dorso (lato largo, concavo): 3 punti; 
%3) πρανής, ventre (lato largo, convesso): 4 punti; 
%4) Κῷος, Cos, 6 punti (lato stretto, piatto): 6 punti.
%Moirai, alle Tychai (Tyche egemone, Tyche Salvatrice o ancora Tyche Felice o fonte di felicità) e alle Nikai (Vincitrice e Vittoria gioiosa). Altre divinità del destino si riconoscono in Daimon, Agathos Daimon e Nemesis,
%
%
%Τὰς ἐν τῇ φαρέτρᾳ ψηφῖδας: Φύλαρχος φησὶ τοὺς Σκύθας μέλλοντας καθεύδειν ἄγειν τὴν φαρέτραν, καὶ εἰ μὲν ἀλύπως τύχοιεν τὴν ἡμέραν ἐκείνην διαγαγόντες, καθιέναι εἰς τὴν φαρέτραν ψηφῖδα λευκήν· εἰ δὲ ὀχληρῶς, μέλαιναν. Ἐπὶ τοίνυν τῶν ἀποθνησκόντων ἐκφέρειν τὰς φαρέτρας καὶ ἀριθμεῖν τὰς ψήφους· καὶ εἰ εὑρεθείησαν πλείους αἱ λευκαί, εὐδαιμονίζειν τὸν ἀπογενόμενον. Ὅθεν παροιμιασθῆναι, τὴν ἀγαθὴν ἥμῶν ἡμέραν λεγόντων τῶν ἐκ τῆς φαρέτρας εἶναι. Καὶ Μένανδρος δὲ φησὶν ἐν Λευκαδίᾳ, τὴν ἀγαθὴν ἡμέραν λευκὴν καλεῖσθαι.
%Xifothếkè dans l'original euripidéen ?
}%
\debutNote{1031}%
\note{1}{Cf. Pollux (X, 150) sur les ustensiles de jeu : %
\citComm{Κυβευτοῦ σκεύη · ἄβαξ ἀβάκιον, κόσκινον, κύβοι, διάσειστοι κύβοι, ἀστράγαλοι, φιμοί, κημοί, τηλία, κήθια κηθίδια, ψῆφοι, πεττοί.}%
φιμός et κημός sont littéralement des \italique{muselières}, avec le sens spécialisé de \italique{cornets} (à dés ou astragales) et sont synonymes de κήθιον. On pourra par exemple se reporter à l'article de Michel Casevitz, «~Sur les termes de jeu : compléments d'histoire des mots~», \italique{Kentron} 36 (2021), p. 209-212. Pour des ouvrages beaucoup plus conséquents, le lecteur pourra lire \italique{Les jeux des Anciens}\expo{2} (1873) de Louis Becq de Fouquières, auteur à la vaste érudition et au style des plus agréables comme le \nieme{XIX} siècle en a été prodigue ; l'édition par Jean Taillardat déjà mentionnée du Περὶ βλασφημιῶν, \italique{Des termes injurieux} et du Περὶ παιδιῶν, \italique{Des jeux grecs}, de Suétone ; la somme de référence qu'est l'article de Hans Lamer «~\latin{Lusoria tabula}~» de la \italique{Paulys Realencyclopädie der classischen Altertumswissenschaft}, \allemand{Band} XIII 2 (1927), col. 1900-2029 ; plus récemment, les \latin{excerpta de ludis} chez Pollux ont été publiés par Salvatore Costanza : \italique{Giulio polluce, onomasticon: excerpta de ludis} accompagnés d'une traduction italienne. Outre l'archéologie, nos témoins principaux quant aux jeux grecs sont Pollux, Hésychius et Eustathe (qui s'appuie sur les traités perdus de Suétone) -- non exempts de contradictions entre eux d'ailleurs.}%
\note{2}{Van Leeuwen a recensé d'autres erreurs similaires où l'ordre des mots est inversé dans le \italique{Ravennas} (la comparaison avec les neuf pièces du \italique{Ravennas} qui nous sont connues par d'autres témoins est préférable : la correction est davantage sûre). Cf. l'édition en fac-similé qu'il fit de notre manuscrit : \italique{Comoediae Undecim cum Scholiis : Codex Ravennas 137, 4, A} [$=$ \italique{Rav.} 429 nouvelle cote], \italique{phototypice editus} in \italique{Codices Graeci et Latini photographice depicti}, t. IX (1904), pp. XVIII-XIX.}%
\note{3}{La scène suivante est une amusante parodie du procès intenté, ou censé intenté, au stratège athénien Lachès, qui était en effet du dême d'Æxone. Le chien du peuple, c'est Cléon, du dême de Cydathène. Lachès avait commandé la flotte envoyée en Sicile en 427. Rappelé au commencement de 425, il fut accusé d'avoir malversé. Le procès (s'il eut lieu) fut présidé suivant l'usage par un thesmothète (v. 935). Cléon, comme synégore, soutint l'accusation. Lachès fut acquitté, car il remplit encore diverses fonctions publiques. Il est intéressant de constater que si le poète n'épargne pas Cléon, il n'est guère plus indulgent pour Lachès, encore que celui-ci appartînt au parti aristocratique. (Note de Willems.)}}
\end{groupe}%
%% v. 1015-1055
\begin{groupe}%
\typeVers{$+$ trim. iamb.}\paratragiques{ἀλλ᾽ ἐν πυκνοῖς δεσμοῖσιν ἐμπεπλεγμένη}
\end{groupe}%
%% v. 1015-1055
\begin{groupe}%
\typeVers{$3$ ia. sync\cat}\paratragiques{\ligne{1033}κήτει βορὰ Γλαυκέτῃ πρόκειμαι,%
\commentaire{1033}{κήτει βορὰ Γλαυκέτῃ πρόκειμαι.}{On sait par le scholiaste des \italique{Oiseaux} (vers 348) que la pièce originale d'Euripide comportait à cet endroit :%
\citComm{ἐκθεῖναι κήτεϊ φορβάν,}%
Γλαυκέτης, l'Athénien glouton (τένθης) est cité au vers 1008 de la \italique{Paix}.
%\citComm{καὶ τὴν ἀγορὰν ἡμῖν ἀγαθῶν\\
%ἐμπλησθῆναι, μεγάλων σκορόδων,\\
%σικύων πρῴων, μήλων, ῥοιῶν.\\
%δούλοισι χλανισκιδίων μικρῶν·\\
%κἀκ Βοιωτῶν γε φέροντας ἰδεῖν\\
%χῆνας νήττας φάττας τροχίλους·\\
%καὶ Κωπᾴδων ἐλθεῖν σπυρίδας,\\
%καὶ περὶ ταύτας ἡμᾶς ἁθρόους\\
%ὀψωνοῦντας τυρβάζεσθαι\\
%Μορύχῳ Τελέᾳ \textbf{Γλαυκέτῃ}, ἄλλοις\\
%τένθαις πολλοῖς.\\
%Fais que notre marché regorge de bonnes choses : denrées de Mégare, aulx, concombres précoces, coings, grenades, petits manteaux pour esclaves ; de chez les Béotiens, qu'on y voie affluer des gens chargés d'oies, de canards, de ramiers, de pluviers ; que les anguilles du Copaïs viennent par paniers et que, massés autour d'elles pour faire nos provisions, nous nous coudoyions avec Morychos, Téléas, \textbf{Glaucetès} et bien d'autres gourmands.}%
Trygée y fait une prière pour que la paix revienne à Athènes ; la ville est sens dessus dessous du fait de la guerre : les embargos et autres razzias ont littéralement renversé le marché : au lieu de toutes les bonnes denrées, l'Athénien n'y rencontre que des rappels de ses privations. 
Autre inversion : si la paix revenait, ce serait alors aux citoyens de bousculer les trois parasites du marché aux anguilles. Aristophane joue sur une image qu'il a développée ailleurs\marqueur{1} : comme le chasseur d'anguille trouble l'eau pour capturer sa proie, le politicien véreux trouble les affaires publiques pour mieux s'enrichir. (Remarquer l'utilisation du verbe τυρβάζεσθαι, qu'Aristophane emploie aussi au vers 257 des \italique{Guêpes} : τὸν πηλὸν ὥσπερ ἀτταγᾶς τυρβάσεις βαδίζων.) 
On peut mettre notre main au feu que le groupe des trois affreux étaient de francs coquins qui se gobergeaient sur le dos d'Athènes pour mener une vie douce, bonne et grasse\marqueur{2}.
Γλαυκέτης est mis par paronomase ; \italique{Nomen, omen} : lorsqu'on s'appelle Γλαυκέτης, on est prédestiné à devenir ce poisson vorace γλαῦκος-κῆτος\marqueur{3} !
(L'image de la femme livrée en pâture (βορά) à la poissonnaille est récurrente chez les Comiques, \italique{confer} Athénée (VII, 315b) citant Platon le Comique ($=$ fr. 57(56) Kassel-Austin)) :%
\citComm{σὲ γάρ, γραῦ, συγκατῴκισεν σαπρὰν\\
ὀρφῷσι σελαχίοις τε καὶ φάγροις βοράν}%
(βορά désigne ici la menue fretaille qui sert de pâture aux poissons carnivores comme l'orphe (ὄρφος / ὀρφώς), qu'Athénée dans le même chapitre rapproche du \italique{glauque}, γλαῦκος.)\\
}%
\debutNote{1033}%
\note{1}{Cavaliers, vers 864-867. Voir aussi le vers 361. Érasme a relevé l'image dans ses \italique{Adages} : n° 2576 (\latin{Anguillas captare}).}%
\note{2}{Cf. pour Morykhos les \italique{Acharniens}, v. 887 ainsi que les \italique{Guêpes}, v. 506 \& 1142. On notera au passage qu'un Téléas fils de Télénikos apparaît à la même époque dans les \italique{Inscriptiones Graecae} (I\expo{3}, 307 ; 308 ; 331 ; 332 ; 370) à la secrétairie (ἐγραμμάτευε) du Trésor (ℎοι ταμίαι το͂ν ℎιερο͂ν χρεμάτον τε͂ς Ἀθεναίας). Voir ce qu'en disent les scholies aux vers 167-8 des \italique{Oiseaux}.}%
\note{3}{Autre jeu du même acabit sur le nom de Glaucétès : (Platon le Comique, fr. 114(106) K.-A.) :%
\citComm{ὦ θεῖε Μόρυχε· πῶς γὰρ οὐ δαίμων ἔφυς;\\
καὶ Γλαυκέτης ἡ ψῆττα, καὶ Λεωγόρας,\\
οἵ ζῆτε τερπνῶς οὐδὲν ἐνθυμούμενοι.}}}
\end{groupe}%
%% v. 1015-1055
\begin{groupe}%
\typeVers{dim. ia\cat}\paratragiques{γαμηλίῳ μὲν οὐ ξὺν\apparat{1034-5}{ξυμπαιῶνι R : \lecon{ξὺν παιῶνι} Zanetti (\italique{cum carmuine} Divus)}}
\end{groupe}%
%% v. 1015-1055
\begin{groupe}%
\typeVers{$+$ 2 ia\cat}\paratragiques{παιῶνι δεσμίῳ δὲ}
\end{groupe}%
%% v. 1015-1055
\begin{groupe}%
\typeVers{dim. iamb.}\paratragiques{γοᾶσθέ μ᾽ ὦ γυναῖκες ὡς}
\end{groupe}%
%% v. 1015-1055
\begin{groupe}%
\typeVers{$2$ troch.}\paratragiques{μέλεα μὲν πέπονθα μέλεος,}
\end{groupe}%
%% v. 1015-1055
\begin{groupe}%
\typeVers{lecyth.}\paratragiques{ὢ τάλας ἐγὼ τάλας,}
\end{groupe}%
%% v. 1015-1055
\begin{groupe}%
\typeVers{dochm. $+$ ...}\paratragiques{\ligne{1039}ἀπὸ δὲ συγγόνων ἀλλ᾽ ἄνομα πάθεα,\apparat{1039}{ἀλλὰν R : \lecon{ἀλλ᾽} Scaliger}}
\end{groupe}%
%% v. 1015-1055
\begin{groupe}%
\typeVers{}\paratragiques{†φῶτά λιτομένα†,\apparat{1040}{λιτομέναν R : ἀντομένα (vel -αν) Schol. : \lecon{λιτομένα} Enger \ligneApp{1041} φεύγουσαν R : χέουσα Rau (\italique{Paratragödia}, p.75) : φλέγουσαν Mugrave (ad Eur. \italique{Or.}1402$=$1396) : \lecon{ψεφαυγοῦς} Bibi, cf Schol. ad v.1054, Hesych. ψ 135}% : ἀλιτόμην ἄν Bibi}%
%\commentaire{1040}{φῶτα λιτομέναν.}{
	%ἠλιτόμηνιen celui qui appelle en vain Hésychius ἠλιτομήνιεν
%264 φέγγουσαν· λαμπρύνουσαν
%1104 φῶντα:λάμποντα
%2839 πολυδάκρυτον· πολλῶν δακρύων αἴτιον
%Alceste :
%  εἴθ’ ἐπ’ ἐμοὶ μὲν εἴη,       [στρ.    (454)
%δυναίμαν δέ σε πέμψαι    (455)
%φάος ἐξ Ἀίδα τεράμνων
%καὶ Κωκυτοῖο ῥεέθρων
%Hippolyte :
%κῶμος λέλακεν, Ἄρτεμιν τιμῶν θεὰν    (55)
%ὕμνοισιν· οὐ γὰρ οἶδ’ ἀνεῳγμένας πύλας
%Ἅιδου, φάος δὲ λοίσθιον βλέπων τόδε.
%Hécube :
%  ὡς πρὶν σφαγῆναί γ’ ἐκτέτηκα καρδίαν
%θρήνοισι μητρὸς τήνδε τ’ ἐκτήκω γόοις.
%ὦ φῶς· προσειπεῖν γὰρ σὸν ὄνομ’ ἔξεστί μοι,    (435)
%μέτεστι δ’ οὐδὲν πλὴν ὅσον χρόνον ξίφους
%βαίνω μεταξὺ καὶ πυρᾶς Ἀχιλλέως.
%
%Une scholie nous apprend une variante parmi les anciens manuscrits :%
%\citComm{γρ[άφεται] καὶ `φῶτα ἀντομένα'}%
%%Et une autre :%
%%\citComm{καὶ τοῦτο ἔχει νοῦν πρὸς τὰ ἑξῆς, `φῶτ' ἀντομέναν, ὃς ἔμ᾽ ἀπεξύρησε πρῶτον }%
%Précieuse scholie, car si le nu Ν de ἀντομένα provient très certainement d'une mélecture du digramme ΛΙ ; elle nous apprend que ledit nu était précédé d'un double alpha ΑΑ dont seul un a subsisté dans la leçon du \italique{Ravennas}. (On a déjà évoqué ce genre de coquilles, très fréquent dans le manuscrit de Ravenne. Cf. par exemples les vers 834 ou 555.)\\
%Je propose donc pour l'\allemand{urtext} des \italique{Thesmophories} :%
%\citComm{ΦΩΤΑΑΛΙΤΟΜΗΝΑΝ\\
%φῶτα ἀλιτόμην ἄν}%
%Cette leçon -- qu'on espère être l'originale -- ayant dérivé en :%
%\citComm{ΦΩΤΑΛΙΤΟΜΕΝΑΝ (\italique{Ravennas})\\
%φῶτα λιτομέναν.\\
%ΦΩΤΑΑΝΤΟΜΕΝΑ(Ν) (Scholiaste)\\
%φῶτα ἀντομένα(ν)}%
%ἀλιτόμην est l'aoriste du verbe ἀλιταίνω, \italique{commettre une faute}, pris ici au sens d'\italique{expier} (\eff{Cf. Chantraine à  aleiths et Aristophane et Hésych alitwmai : atimw}). fwta (accusatif de durée) est pris dans son acception courante de \italique{jours}.
%Traduisez donc :%
%\citComm{ἀπὸ δὲ συγγόνων ἀλλ᾽ ἄνομα πάθεα\\
%φῶτα ἀλιτόμην ἄν,\\
%πολυδάκρυτον Ἀίδα γόον φρύγουσαν\marqueur{Pour la correction : φρύγουσαν à la place de la leçon du manuscrit φεύγουσαν, voir le commentaire au vers suivant.}.\\
%Par la faute de mes parents, j'endurerais des tortures impies pendant des jours, me desséchant pour l'Hadès d'un torrent de larmes.}%
%vἀλιτόμην ἄν Par la faute de sa mère/Euripide...
%	j'expierais des jours durant
%Eschyle suppliantes 416 : ὃς οὐδ’ ἐν Ἅιδου τὸν ἀλιτόντ’ ἐλευθεροῖ. (ed. Murray du TLG... délire ? ton thanonta alt. edd.)
%Hésychius : “ἐκ δ᾽ἄγαγε πρὸ φόωσδε᾽: προ εξήγαγε δὲ εἰς τὸ φῶς τὸν ἀλιτόμηνον Homéri
% ἠλιτόμηνα: ἐλαττούμενα τοῖς μησίν
% ἠλῐτό-μηνις· ὁ μάτην ἐγκαλῶν
%Crux : ἀφήγουσα:μεταστρέφουσα
%φεγγουσα
%φρύγ
%φθεγγουσα <- n'existe pas
%L'expression πολυδάκρυτος γόος est homérique (\italique{Odyssée}, XIX, 213).
%}%
}
\end{groupe}%
%% v. 1015-1055
\begin{groupe}%
\typeVers{do. $+$ cr. $+$ ba.}\paratragiques{πολυδάκρυτον Ἀίδα γόον ψεφαυγοῦς.%\apparat{1041}{φεύγουσαν R : χέουσα Rau (\italique{Paratragödia}, p.75) : φλέγουσαν Mugrave (ad Eur. \italique{Or.}1402$=$1396) : \lecon{ψεφαυγοῦς} Bibi, cf Schol. ad v.1054 e Hesych. ψ 135}%
\commentaire{1041}{Ἀίδα,}{Aristophane use du même génitif dorique au vers 1333 des \italique{Grenouilles} (parodie d'Euripide).}%
\commentaire{1041}{πολυδάκρυτον Ἀίδα γόον φεύγουσαν.}{Le manuscrit porte l'incompréhensible †φεύγουσαν†. Y suppléer par χέουσα(ν), à l'instar de certains éditeurs, c'est faire fi de sa leçon. Je dirais même plus : χέουσα est le mot banal qu'on attendait, il paraît alors très difficile qu'un scribe s'y soit trompé. (σφριγῶσα, \italique{gonflée}, aurait au moins le mérite de se justifier paléographiquement mais le terme paraît impropre en grec à l'image d'yeux gonflés de larmes.)
Une autre proposition, `φλέγουσα(ν)', a aussi emporté la faveur de nombreuses éditions. Elle est reprise par Jean Taillardat à propos de l'image de la Muse-flamme (§736 n.2)... mais le grand Taillardat me semble s'abuser : si un poème, un chant, un \italique{thrêne} (qui est avant tout sonore ; la flamme étant celle de la voix selon moi) peuvent être enflammés, cela ne prouve pas pour autant que des pleurs peuvent l'être.\\
Le scholiaste nous fournit une information qui pourrait peut-être bien débrouiller l'écheveau: une glose, σκοτεινήν, \italique{obscur}, se trouve au vers 1054 pour qualifier soi-disant αἰόλαν, \italique{rapide}, \italique{mobile}, \italique{changeant}. 
Sauf qu'αἰόλος n'a jamais eu le sens de σκοτεινός et que Willems a clairement montré qu'il fallait prendre αἰόλος en son sens courant\marqueur{1}.
Suidas (αι 245) -- ou l'un des lexicographes de la \italique{Souda}, si l'on préfère -- n'a pas manqué néanmoins de relever la glose à l'entrée Αἰόλαν. Mais Suidas est coutumier de ce genre de fausses attributions (voir la note au vers 735).
Or en cherchant des mots proches de †φεύγουσαν†, je suis retombé sur ce cher Hésychius (ψ 135) :%
\citComm{ψεφαυγοῦς· σκοτεινῆς,}%
Génitif dorique (comme Ἀίδα) d'un certain *ψεφαυγής, \italique{au sombre regard}.
Euripide a une prédilection pour forger des composés en -αυγής. (\italique{Confer} le Chantraine, \italique{sub voce} αὐγή.) 
Si l'on fait l'hypothèse de la chute de `ψε', le φαυγους restant aurait bien pu être complété d'un `αν' en regard de l'autre participe †λιτομέναν† juste au-dessus, et donner ainsi notre leçon †φεύγουσαν†.
%Point tout-à-fait exacte, car ce n'est pas l'image du \italique{feu} qu'évoque Andromède, mais celle du \italique{dessèchement}.
%Lisez donc `φρύγουσα' ; un epsilon Ε se sera substitué au rhô Ρ\marqueur{Se reporter au tableau comparatif de l'évolution des crusives grecques du livre déjà cité\eff{The Greek Manuscripts ?} pour constater à quel point ces deux lettres ont pu se ressembler sur certains manuscrits}.\\
%Le verbe φρύγειν, littéralement : \italique{faire griller}, est d'un emploi assez large. Ainsi, les dérivés φρύγιον et φρύγανον désignent du \italique{bois sec}. Une épigramme funéraire de l'\italique{Anthologie Palatine} (VII, 293) fait de vents le sujet du verbe φρύγειν, qui ont \italique{desséché} Nicophềme au point de le faire périr de soif :%
%\citComm{οὐ χεῖμα Νικόφημον, οὐκ ἄστρων δύσις ἁλὸς Λιβύσσης κύμασιν κατέκλυσεν ἀλλ ἐν γαλήνῃ, φεῦ τάλας,\\
%ἀνηνέμῳ πλόῳ πεδηθείς, \textbf{ἐφρύγη δίψευς ὕπο}}%
%frugein = mourir ?
%Asséchant son torrent lacrimal (πολυδάκρυτος γόος\marqueur{L'expression est homérique (\italique{Odyssée}, XIX, 213) ; on trouve chez Eschyle (\italique{Choéphores}, vers 449) :%
%\citComm{χέουσα πολύδακρυν γόον κεκρυμμένα.}%
%Et chez Euripide (\italique{Suppliantes}, vers 773):%
%\citComm{Ἅιδου τε μολπὰς ἐκχέω δακρυρρόους.}}) qu'Andromède déverse l'\italique{assèchent l}
%σφριγάω σφριγῶσα
%Hésychius : σφριγῶσα· πεπυκνωμένη, συνεστηκυῖα, ἰσχύουσα
%cf. Taillardat §62.
}%
\debutNote{1041}%
\note{1}{Note n° 27, reproduite ci-après.}}
\end{groupe}%
%% v. 1015-1055
\begin{groupe}%
\typeVers{extr. metr.}\paratragiques{αἰαῖ αἰαῖ ἒ ἒ,}
\end{groupe}%
%% v. 1015-1055
\begin{groupe}%
\typeVers{dim .troch.}\paratragiques{ὃς ἔμ᾽ ἀπεξύρησε πρῶτον,%
\commentaire{1043-51}{Papyrus d'Oxyrhynque n° 4935.}{Un maigre fragment (n° 1) du Papyrus d'Oxyrhynque 4935\marqueur{1} (\nieme{II} siècle ap. J.-C.), reproduit en annexe, nous apporte une leçon nouvelle au vers 1047 : le `μοι' du \italique{Ravennas} n'y figure pas. De plus, il confirme -- comme l'avait déjà fait Willems en avance sur son temps -- le `βάρβαρον' du vers 1051 qui laissait de nombreux éditeurs sceptiques. (Voir sa \nieme{26} note, très instructive sur ce passage et le rôle de l'étoile Sirius.)}%
\debutNote{1043-51}%
\note{1}{\italique{The Oxyrhynchus Papyri}, LXXIII (2009), ed. A. Benaissa.}}
\end{groupe}%
%% v. 1015-1055
\begin{groupe}%
\typeVers{}\paratragiques{\ligne{1044}ὃς ἐμὲ κροκόεντ᾽ ἐνέδυσεν·\apparat{1044}{\lecon{κροκόεντ᾽} R : κροκόεν τόδ' Bergk \espApp ἐνέδυσεν R : ἀμφέδυσεν Hermann ap Schol.}%
%\commentaire{1044}{}{Phéniciennes d'Euripide 1491 : 
%  Ανtigone.                 οὐ προκαλυπτομένα βοτρυχώδεος   (1485)
%ἁβρὰ παρηίδος οὐδ’ ὑπὸ παρθενί-    (1486-1487)
%ας τὸν ὑπὸ βλεφάροις φοίνικ’, ἐρύθημα προσώπου,   (1488)
%αἰδομένα φέρομαι βάκχα νεκύ-
%ων, κράδεμνα δικοῦσα κόμας ἀπ’ ἐ-    (1490)
%μᾶς, στολίδος κροκόεσσαν ἀνεῖσα τρυφάν,
%ἁγεμόνευμα νεκροῖσι πολύστονον. αἰαῖ, ἰώ μοι.
%ὦ Πολύνεικες, ἔφυς ἄρ’ ἐπώνυμος· ὤμοι μοι, Θῆβαι·   (1493-1494)
%σὰ δ’ ἔρις—οὐκ ἔρις, ἀλλὰ φόνῳ φόνος—    (1495)
%Οἰδιπόδα δόμον ὤλεσε κρανθεῖσ’
%αἵματι δεινῷ, αἵματι λυγρῷ.
%τίνα προσῳδὸν
%ἢ τίνα μουσοπόλον στοναχὰν ἐπὶ
%δάκρυσι δάκρυσιν, ὦ δόμος, ὦ δόμος,   (1500)
%ἀγκαλέσωμαι,
%τρισσὰ φέρουσα τάδ’ αἵματα σύγγονα,}
}
\end{groupe}%
%% v. 1015-1055
\begin{groupe}%
\typeVers{dim. troch.}\paratragiques{ἐπὶ δὲ τοῖσδε τόδ᾽ ἀνέπεμψεν}
\end{groupe}%
%% v. 1015-1055
\begin{groupe}%
\typeVers{pherecr.}\paratragiques{ἱερόν, ἔνθα γυναῖκες.}
\end{groupe}%
%% v. 1015-1055
\begin{groupe}%
\typeVers{dochm.}\paratragiques{ἰώ μοι μοίρας\apparat{1047}{[ι]\douteux{ω}μοι\douteux{ρ}[ας \PsiQuatre : ἰώ μοι μοίρας R : ἰώ μοίρας Hermann \espApp δαίμων R : δαῖμον Burges \espApp ἀνέτικτε R : \lecon{ἄτεγκτε} Ellebodius : ἅν ἔτικτε Casaubon}%
%\commentaire{1047}{ἰώ μοι μοίρας.}{\italique{Andromaque}, vers 824 : \citComm{ἰώ μοί μοι.}}
}
\end{groupe}%
%% v. 1015-1055
\begin{groupe}%
\typeVers{ithyph.}\paratragiques{ἄτεγκτε δαίμων·%\apparat{1047}{δαίμων R : δαῖμον Burges \espApp ἀνέτικτε R : \lecon{ἄτεγκτε} Ellebodius : ἅν ἔτικτε Casaubon}
}
\end{groupe}%
%% v. 1015-1055
\begin{groupe}%
\typeVers{hem.}\paratragiques{ὦ κατάρατος ἐγώ·}
\end{groupe}%
%% v. 1015-1055
\begin{groupe}%
\typeVers{lecyth.}\paratragiques{\ligne{1048}τίς ἐμὸν οὐκ ἐπόψεται}
\end{groupe}%
%% v. 1015-1055
\begin{groupe}%
\typeVers{trim. iamb.}\paratragiques{\ligne{1049}πάθος ἀμέγαρτον ἐπὶ κακῶν παρουσίᾳ;\testimonia{1049}{(παρουσίᾳ) Moeris π 14}}
\end{groupe}%
%% v. 1015-1055
\begin{groupe}%
\typeVers{$4$ dactyl.}\paratragiques{εἴθε με πυρφόρος αἰθέρος ἀστὴρ—}
\end{groupe}%
%% v. 1015-1055
\begin{groupe}%
\typeVers{paroem.}\paratragiques{τὸν βάρβαρον ἐξολέσειεν.}
\end{groupe}%
%% v. 1015-1055
\begin{groupe}%
\typeVers{$4$ dactyl.}\paratragiques{οὐ γὰρ ἔτ᾽ ἀθανάταν φλόγα λεύσσειν\apparat{1052}{λεύσειν R : \lecon{λεύσσειν} Biset}}
\end{groupe}%
%% v. 1015-1055
\begin{groupe}%
\typeVers{$+$ $4$ dact.}\paratragiques{ἐστὶν ἐμοὶ φίλον, ὡς ἐκρεμάσθην,}
\end{groupe}%
%% v. 1015-1055
\begin{groupe}%
%\typeVers{hypod. $+$ doch.}
\paratragiques{λαιμότμητ᾽ ἄχη δία μένων αἰόλαν\apparat{1054}{δαιμόνων R, Schol. : δαιμονῶν Brunck : δαιμόνι᾽ amicus Fritzschii : \lecon{διὰ μένων} Willems}%
\commentaire{1054-5}{λαιμότμητ᾽ ἄχη δαιμόνων αἰόλαν / νέκυσιν ἐπὶ πορείαν.}{%ὡς ἐκρεμάσθην = depuis que j'ai été attaché ? 
Willems a consacré une note à ce passage\marqueur{1}. Si les vues qu'il y développe sont ingénieuses, elle n'emporte pas pour autant ma conviction. %D'abord, la traduction λαιμότμητ᾽ ἄχη reporter à la  note de Willems, reproduite en annexe --, le texte en est corrompu, certainement en cet étonnant †δαιμόνων†. 
Willems propose de construire ἄχη δία μένων là où le \italique{Ravennas} porte ἄχη †δαιμόνων†. À défaut d'autres propositions plus convaincantes, c'est néanmoins la correction que je retiendrai.%\marqueur{2} ; je préfère y voir le verbe ἀδημονέω (ἀδαιμονέω DGE)
%\eff{Retoucher ! MAsculin adhmonwn ? ὡς marque l'exclamation ?} 
}%
\debutNote{1054-5}%
\note{1}{\nieme{27} note aux \italique{Thesmophories} d'Alphonse Willems, reproduite en annexe.}%
%\note{2}{Si le sens est tout-à-fait satisfait par la conjecture de Willems ; qu'il se justifie habilement de la construction peu habituelle; ; ma réserve vient de ce que je m'explique difficilement le passage du epsilon de μένων au omicron de δαιμόνων. Tandis que la transformation du êta en le digramme ΑΙ est beaucoup plus naturelle en supposant l'étape intermédiaire d'un epsilon : Η $\xrightarrow{}$ Ε  $\xrightarrow{prononciat. ident.}$ ΑΙ.}%
}%
\end{groupe}%
%% v. 1015-1055
\begin{groupe}%
\typeVers{ithyphall.}\paratragiques{νέκυσιν ἐπὶ πορείαν.\\\passeLigne}%
\end{groupe}%
%% v. 1056-1057
\begin{groupe}%
\personnageGrec{Εὐριπίδης}%
\trimIamb{\ligne{1056}χαῖρ᾽ ὦ φίλη παῖ· τὸν δὲ πατέρα Κηφέα\apparat{1056}{ευρ\auDessus{ι}{π} ηχω in mg. \Rpc, \Mon}\\
ὅς σ᾽ ἐξέθηκεν ἀπολέσειαν οἱ θεοί.\\\passeLigne}%
\end{groupe}%
%% v. 1058
\begin{groupe}%
\personnageGrec{Μνησίλοχος}%
\trimIamb{σὺ δ᾽ εἶ τίς ἥτις τοὐμὸν ᾤκτιρας πάθος;\apparat{1058}{ὤκτειρας R : \lecon{ᾤκτιρας} van Herwerden \espApp πάθος R : \lecon{πάθος;} Biset}\\\passeLigne}%
\end{groupe}%
%% v. 1059-1063a
\begin{groupe}%
\personnageGrec{Εὐριπίδης}%
\trimIamb{Ἠχὼ, λόγων ἀντῳδὸς ἐπικοκκάστρια,\testimonia{1059}{Suidas η 686 (ἐπικοκκάστρια codd. GM\expo{a.c.} vel ἐπικοκάστρια cod. I vel ἐπικαγχάστρια cod. A vel ἐπικαχάστρια codd. FV\expo{mg.} vel ἐπικοστράτ textu V) \espTest (ἀντῳδὸς ἐπικοκκάστρια) \Eustathe 1761.26 ἐκτίθεται καινοφώνους λέξεις οἷον [ ... ] καὶ τὸ ἐπικοκκάζειν ὕβρει ἤ κατά τινας ἐπιμηκάζειν, ὃθεν παρὰ Ἀριστοφάνει τὸ `ἀοιδὸς ἐπικοκκάστρια'}%
\commentaire{1059}{ἀντῳδὸς ἐπικοκκάστρια.}{Ceci rappelle la jolie épigramme d'Archias (\italique{Anthol. Palat.}, VII, 191) sur une pie (κίσσα\marqueur{1})%
\citComm{qui souvent, pareille à l'Écho, faisait entendre, d'un bec prompt à répliquer, une équisonnance bruyante et moqueuse.\\
πολλάκι δὲ κρέξασα πολύθροον, οἷά τις ἀχώ,\\
κέρτομον ἀντῳδοῖς χείλεσιν ἁρμονίαν.}%
(Note de Willems.)}%
\debutNote{1059}%
\note{1}{Si la pie est un oiseau capable de parler, selon Descartes, l'oiseau dont il question dans cette épigramme serait à en croire le \LSJ \ un corvidé cousin, de la famille des geais : un \italique{geai des chênes}, \latin{Garrulus glandarius}. (NDÉ.)}%
\commentaire{1059}{Ἠχὼ λόγων ἀντῳδὸς ἐπικοκκάστρια.}{Seule compagnie d'Andromède : la nymphe Écho, cachée dans un antre non loin, et qui répète les paroles. L'occasion d'un détournement comique de ces échos était trop belle pour qu'Aristophane n'en profitât pas !
Euripide – qui joue dans ce passage la nymphe Écho – répond effectivement à l'appel de l'Andromède postiche qu'essaye d'interpréter le peu crédible vieillard Mnésiloque. 
Dans la tragédie éponyme d'Euripide d'antan, l'acteur d'Écho ne se montrait pas sur scène mais renvoyait les répliques depuis un antre à proximité de la falaise où était enchaînée Andromède. %\eff{sûr ? cf. représentation}. 
Ajoutez à cela la μηχανή permettant une entrée aérienne de Persée et vous comprendrez quelque peu l'impression forte et durable que la pièce suscita auprès des spectateurs.\\
Revenons au vers 1059 de notre comédie : Ἠχὼ, λόγων ἀντῳδὸς ἐπικοκκάστρια.\\
On a passé un peu vite sur l'adjectif ἐπικοκκάστρια ; si le sens de λόγων ἀντῳδὸς est clair : « la rediseuse en chantant des paroles », qualificatifs prévisibles de la nymphe Écho (notez tout-de-même que les échos étaient probablement renvoyés en \italique{chantant}, ou au moins en \italique{persiflant}), Aristophane ajoute l'hapax ἐπικοκκάστρια. 
Scrutons ce mot : ἐπι-κοκκάσ-τρια ; le préfixe ἐπι- apporte selon nous sa nuance d'hostilité, ici de causticité %(en fait, les autres mots composés sur ἐπι- jouent avec les différentes nuances du préfixe qui peut marquer la succession temporelle ainsi que l'hostilité) 
; le suffixe -τρια est courant, principalement attesté chez les comiques et en prose\marqueur{1} ; reste κοκκάς, qu'une oreille française ne peut s'empêcher de rapprocher de notre `cocasse', risible mais avec une certaine violence comme le coq dont il tire son origine… 
Si le volatile d'Aristophane n'est pas un coq ; κοκκάς renvoie au κόκκυξ ou \italique{coucou}, les Grecs ne prêtaient pas moins à cet oiseau de malheur une méchanceté bien supérieure. \italique{Confer} le magistrale note de Willems sur le \nieme{507} vers des \italique{Oiseaux} et les propos salés (\latin{petulantiæ sales}) qu'on jetait aux laboureurs en imitant le cri du coucou. (Soit dit en passant, on retrouve les mêmes traits dans feu nos traditions orales : le coucou chantait les événements fâcheux (dont en premier lieu, bien-sûr, le cocufiage)\marqueur{2}.\\
Le scholiaste nous est ici précieux\marqueur{3} : il définit ainsi ἐπικοκκάστρια : εἰωθυῖα γελᾶν, γελάστρια, « ayant l'habitude de rire, moqueuse ». Aristophane nous prévient donc que les échos auront quelque chose de piquant et que si on les trouve d'une « \anglais{annoying loquacity} »\marqueur{4} , c'est que, peut-être, le sel attique nous en a échappé.\\
Une autre scholie nous permet de préciser un peu la teneur des échos :%
\citComm{ἐπεὶ εἰσήγαγε κακοστένακτον τὴν Ἡχὼ ὁ Εὐριπίδης ἐν τῇ Ανδρομέδᾳ, εἰς τοῦτο παίζει [...]}%
La bonne interprétation de cette scholie réside dans le sens à accorder à l'adjectif κακοστένακτος : \italique{à la méchante lamentation} ?\\
Il est étonnant que les modernes éditions n'aient pas cherché à percer ce que les répliques d'Écho-Euripide pouvaient avoir d'amusantes ou de méchantes… Pourtant, elles disposaient déjà d'un exemple de détournement par la « rediseuse goguenarde » : à la fin de la scène, au vers 1094, alors que l'archer Scythe tente d'attraper Écho-Euripide et que ce dernier s'enfuit sans demander son reste, une scholie nous indiquait clairement ce que disait en son jargon l'archer Scythe :%
\citComm{ἀντὶ τοῦ `οὐ χαιρήσεις'\\
À la place de : « tu ne vas pas te réjouir ».}%
On peut déduire de cette scholie que le Scythe – qui ne tient aucun compte dans son parler des aspirations – avait menacé Écho-Euripide d'un : `οὐ καιρήσεις'. Or, le manuscrit de Ravenne comportait effectivement la réponse ἐπικοκκάστρια d'Euripide-Écho\marqueur{5} : `οὐκ᾽ αἰρήσεις', c'est-à-dire en bon attique : `οὐχ αἱρήσεις', « tu ne m'attraperas pas ! ». Le cas était simple mais il aura fallu tout de même des siècles de philologiques hésitations entre les répartitions : Τοξότης (l'Archer) : « οὐ καιρήσεις » — Εὐριπίδης : « οὐ καιρήσεις » et Τοξότης : « οὐκ αἰρήσεις » — Εὐριπίδης : « οὐκ αἰρήσεις » avant que le grand Wilamowitz\marqueur{6} ne s'avise de la solution conciliatrice, la seule sensée et la seule comique : Τοξότης : « οὐ καιρήσεις » — Εὐριπίδης : « οὐχ αἱρήσεις ».\\
Bref ! les répétitions de la nymphe servent à Aristophane pour son jeu de massacre de la tragédie tant appréciée d'Euripide\marqueur{7} : par de légers détournements phonétiques (facilités par les barbarismes de l'archer Scythe), Écho transforme ce qui devait être une scène d'émotion grave en une foire d'empoigne générale.\\
C'est -- reconnaissons-le -- un exercice des plus périlleux que de s'aventurer dans des tentatives de reconstruction ; laisser ainsi libre cours à l'imagination peut exposer à de rudes retours de bâton ; pour Aristophane, nous devons cependant impérieusement nous y livrer puisqu'il vient de nous prévenir que les échos seraient \italique{moqueurs} ! Avec une réserve cependant : lorsque l'ἐπικοκκάστρια d'une réplique nous est demeurée celée, et que malgré nos efforts les seules propositions que nous avions à présenter n'étaient pas même dignes d'intégrer l'apparat critique -- le cas n'est pas isolé --, nous lui avons apposé la \italique{croix du désespoir} †.. Charge à d'autres de compléter notre ébauche ! je pense, pour ma part, ne pas avoir démérité d'Aristophane.}%
\debutNote{1059}%
\note{1}{Cf. Chantraine pour une longue liste des mots formés sur le même modèle, \italique{La formation des noms en grec ancien}, p. 106-7.}%
\note{2}{Pour s'en convaincre, il n'y a qu'à aller feuilleter le \italique{Répertoire} Coirault \italique{des chansons françaises de tradition orale}.}%
\note{3}{%Une fois n'est pas coutume ! ne boudons pas notre plaisir : pour une fois que dans pièces, si pauvre en scholies de qualité par rapport aux autres pièces d'Aristophane, les commentaires relèvent d'une érudition tout alexandrine ; connaissance lexicographique, donc, d'une part ; scholies qui font état parfois état d'un texte mieux conservé que celui qui nous est parvenu d'autre part ; enfin, possibilité pour le commentateur d'aller comparer avec l'original d'Euripide.
Les glosateurs avaient sans doute un texte des \italique{Thesmophories} en meilleur état que le nôtre (ce passage est le moins bien conservé de toute la pièce ; et pour cause : le comique a très tôt échappé aux scribes qui le recopiaient). De plus, ils avaient un avantage indéniable sur nous : ils pouvaient aller comparer la parodie avec l'original d'Euripide.}%
\note{4}{Austin-Olson… et pis encore (p. LXX) : « \anglais{much of the effect of the parodies of Helen and Andromeda depends on the fact that the tragic heroines are played by an individual wearing a prominent stagephallus} » (!!!) Mais pour qui prennent-ils Aristophane ?! le plus gras des vaudevillistes serait pris de scrupules devant les procédés que ces brillants professeurs prêtent à Aristophane, le plus fin des Attiques ! Il faudrait leur expliquer qu'il y a un temps pour tout, et qu'un pastiche d'euripidéenne tragédie n'est pas une scène burlesque de déshabillage ou encore une pièce sur la grève du sexe ; c'est calomnier Aristophane que de croire – et plus grave : faire croire à ses lecteurs, circonstance aggravante lorsque l'on prétend à être l'édition de référence des \italique{Thesmophories} ! – qu'il ne soit capable que de jouer d'un seul registre. Au reste, cette opinion est un expédient bien pratique qui dispense de se pencher sur ce que signifie véritablement le texte… ou de reconnaître honnêtement qu'on ne le comprend pas. Enfin, le véritable Aristophane nous dit lui-même ce qu'il pense du recours abusif à ces « phallus de cuir, épais, rouges par le bout, pour donner à rire aux gamins », c'est au vers 538 des \italique{Nuées} (voir aussi la \italique{Paix}, v. 739-753) ; que les futurs exégètes d'Aristophane ruminent bien ce passage avant de prétendre que notre poète concoctait à toutes les sauces des scènes de sexes sémillants. Et qu'ils lisent aussi l'article que Willems y avait consacré, « Le nu dans la Comédie ancienne », t. III, p. 381-400 (reproduit dans les annexes) : Aristophane recourait à l'\italique{olisbos} avec la plus grande parcimonie, et toujours pour de bonnes raisons.}%
\note{5}{En attribuant, il est vrai, la réplique à l'archer. Mais les indications de personnages sont de peu de valeur : les anciens parchemins ne comportaient pour la plupart que de légers signes (\italique{paragraphos} ou parfois simple indentation) pour indiquer le changement d'interlocuteur, et ce sans préciser quels étaient ces interlocuteurs ; les indications de personnages furent ajoutées par des scribes très postérieurement… et sont parfois assez folkloriques. (Pour celles du \italique{Ravennas}, quelques indications sont de la main d'Euphrosynus Boninus, celui qui compulsa le manuscrit pour la première édition imprimée des \italique{Thesmophories} -- la \italique{Juntina} (Venise, 1516) des apparats critiques. Von Velsen sut identifier les annotations de la main d'Euphrosynus Boninus sur le \italique{Ravennas} et ne les fit pas apparaître dans son apparat critique. Nous l'avons suivi en cela. Cf. Albert Martin, \italique{Les scolies du manuscrit d'Aristophane à Ravenne} (1882), p. V ; A. von Velsen, \italique{Über den Codex Urbinas der Lysistrata und der Thesmophoriazusen des Aristophanes} (1871).)}%
\note{6}{En simple note de ses \italique{Kleine Schriften}, IV (1962), p. 486. Mais ces écrits n'ont été publiés qu'après la mort du grand philologue. Si une partie des travaux modernes rétablissent correctement la menace du Scythe et son écho, il semble que le mérite en revienne plutôt à Sommerstein, qui eut la bonne intuition indépendamment de Wilamowitz.}%
\note{7}{Rappelons, par exemple, que c'est en relisant \italique{Andromède} que le Dionysos des \italique{Grenouilles} est pris d'un soudain désir d'Euripide.}\\
ἥπερ πέρυσιν ἐν τῷδε ταὐτῷ χωρίῳ\apparat{1060}{ταυτῷ R : \lecon{ταὐτῷ} Dindorf (1830)}%
\commentaire{1060}{πέρυσιν.}{L'\italique{Andromède} avait été produite l'année précédente, cf. la scholie :%
\citComm{ἐπεὶ πέρυσιν ἐδιδάχθη <ἡ> Ἀνδρομέδα.}
(Voir l'introduction sur la question de la date des \italique{Thesmophories}.)}\\
Εὐριπίδῃ καὐτὴ ξυνηγωνιζόμην.\apparat{1061}{καὶ αὐτὴ R : \lecon{καὐτὴ} Biset}\\
ἀλλ᾽ ὦ τέκνον σὲ μὲν τὸ σαυτῆς χρὴ ποιεῖν,\apparat{1062}{τοσαύτης R : \lecon{τὸ σαυτῆς} Dindorf (1839)}\\
κλαίειν ἐλεινῶς.\apparat{1063}{\lecon{κλαίειν} et \lecon{ἐπικλαίειν} R : κλάειν et ἐπικλάειν Brunck}\\\passeLigne}%
\end{groupe}%
%% v. 1063b
\begin{groupe}%
\personnageGrec{Μνησίλοχος}%
\typeVers{trim. iamb.}\trimIamb{\ligne{1063}\decalage{κλαίειν ἐλεινῶς.}σὲ δ᾽ ἐπικλαίειν ὕστερον.%
\commentaire{1063}{σὲ δ᾽ ἐπικλάειν ὕστερον.}{Puisqu'Écho répète tout, il est logique qu'elle répète \italique{en retour} (ἐπι-) les pleurs versés par Andromède... il y a là cependant une menace de vengeance clairement exprimée par Mnésiloque pour après sa libération (avec l'emploi à double sens du préfixe ἐπι- : hostilité et succession temporelle). 
%C'est en tout cas en ce sens que le comprit le scholiaste : il nous précise que Mnésiloque répond \italique{à part} (ἠρέμα) -- ce qui n'est d'ailleurs pas le cas puisqu'Euripide lui répond -- :%
%\citComm{ἠρέμα ὁ κηδεστὴς τοῦτο πρὸς τὸν Εὐριπίδην.}%
%Mais le verbe ἐπικλαίειν n'est pas autrement attesté à l'époque classique, au contraire d'ἐπικλάειν\marqueur{1}
}%
%\debutNote{1063}\note{1}{Remarque : la syllabe qui distingue les deux verbes (par sa longueur) tombe dans les deux cas qui se présentent au vers 1063 en position qui peut être indifféremment brève ou longue.}
%\commentaire{1063}{σὲ δ᾽ ἐπικλαίειν ὕστερον.}{\eff{Émouvoir ?}}
\\\passeLigne}%
\end{groupe}%
%% v. 1064
\begin{groupe}%
\personnageGrec{Εὐριπίδης}%
\trimIamb{ἐμοὶ μελήσει ταῦτά γ᾽. ἀλλ᾽ ἄρχου λόγων.\\\passeLigne}%
\end{groupe}%
%% v. 1065-1069a
\begin{groupe}%
\personnageGrec{Μνησίλοχος}%
\typeVers{metr. anap.}\anap{ὦ νὺξ ἱερὰ\apparat{1065}{ὁ μνησίλοχος ὡς ἀνδρομέδα in mg. \Rpc, \Mon}%
\commentaire{1065-72}{ὦ νὺξ ἱερὰ κτλ.}{Nous savons par une scholie que ces vers sont les premiers d'\italique{Andromède} (τοῦ προλόγου Ἀνδρομέδας εἰσβολή). Le mètre est anapestique\marqueur{1} ; les poètes tragiques y recourent effectivement pour certaines scènes de lamentation\marqueur{2}.}%
\debutNote{1065-72}%
\note{1}{De même que dans le prologue d'\italique{Iphigénie en Aulis}.}%
\note{2}{Voir \italique{The Songs of Aristophanes}, p. 57.}\\
ὡς μακρὸν ἵππευμα διώκεις\testimonia{1066-8}{(ὡς μακρὸν -- αἰθέρος) Suidas ι 531 $\sim$ Zonaras 1117}\\
ἀστεροειδέα νῶτα διφρεύουσ᾽\testimonia{1067}{(νῶτα) Moeris ν 6}\testimonia{1067-8}{Suidas αι 122 $\sim$ Zonaras 67, 78 (sed ἀστρο- pro ἀστερο-)}\commentaire{1067}{ἀστεροειδέα νῶτα.}{La voûte solide du firmament incorruptible. (Note de Willems.)}\\
αἰθέρος ἱερᾶς\\
τοῦ σεμνοτάτου δι᾽ Ὀλύμπου;%\commentaire{1068}{δι᾽ Ὀλύμπου.}{Hésychius (ο 654) :\citComm{Ὀλύμπιος· οὐράνιος.}}
\\\passeLigne}%
\end{groupe}%
%% v. 1069b
\begin{groupe}%
\personnageGrec{Εὐριπίδης}%
\anap{\ligne{1069}\decalage{τοῦ σεμνοτάτου δι᾽ Ὀλύμπου;}†δι᾽ Ὀλύμπου;†\apparat{1069}{\lecon{δι᾽ Ὀλύμπου} R : †δι᾽† ὀλυμπίας Bibi ap. Hesych. (ο 649)}%
\commentaire{1069}{τοῦ σεμνοτάτου δι᾽ Ὀλύμπου; --- δι᾽ Ὀλύμπου.}{%Prov : Oλυμπικὴ στοά: ἐπὶ τοῦ λάλου, ἢ ὅτι εἰκὸς ἔχειν τὴν στοὰν ἀντῳδὸn ἠχὼ, ἢ διὰ τὰςτῶνἐκεῖσυνηγμένωνβοάς
En bonne rigueur philologique, nous devons nous demander sur quels éléments probants s'appuyer dans notre aventureuse entreprise de rétablissement de l'\allemand{Urtext} des \italique{Thesmophories}. Car il ne suffit de partir à l'aventure, toutes brides de l'imaginations abattues ; encore faut-il apporter des indices de ce que nous avançons. 
Nous avons déjà évoqué des bribes de scholies ; nous en croiserons d'autres en route. 
Une variante orthographique entre un mot et son écho nous fournira aussi une amorce philologique. 
Cela est bien maigre... est-ce qu'Hésychius, ou les commentateurs qu'il recopie\marqueur{1}, dont nous avons assez montré toute l'importance, avait entre les mains une édition d'Aristophane assez bien conservée pour qu'y apparaisse certains détournements de la nymphe Écho ? Réponse : ce n'est pas impossible ! %\eff{Hum..} 
voyons cela avec première réplique ἐπικοκκάστρια de la nymphe Écho. \\
Voici la glose, parcellaire, d'Hésychius (ο 649) qui éclaire et justifie peut-être l'interruption d'Écho :%
\citComm{ὀλυμπία· δοκός\\
olympia : poutre.}%
La langue d'Aristophane possède la matrice métaphorique : \italique{poutre} $=$ \italique{poutre à laquelle on se pend} $=$ \italique{pendaison}. (Le français dispose d'une matrice semblable avec le mot \italique{corde}.) J'en veux pour preuves : \ \ \ 1° \ ce commentaire\marqueur{2} : [δοκὸς] παρὰ Ἀριστοφάνει ἐν Ταγηνισταῖς σημαίνει τὴν ἀγχόνην, commentaire repris par Hésychius (δ 2129) : δοκόν· τὴν ἀγχόνην.
\ \ \ 2° \ Dans les \italique{Oiseaux}, après avoir poutré l'importun chresmologue Hiéroclès, on l'envoie se faire pendre εἰς Ἐλύμνιον (vers 1126), ville d'Eubée dont le nom fait entendre à l'oreille grecque sa \italique{poutre} : ἐλύμνιαι· δοκοὶ ὀρόφιναι (Hésychius δ 2227).\\
Rappeler à Andromède sa condition de crucifiée \italique{pendue} à la falaise, voilà une incise qui pourrait être qualifiée de κακοστένακτος.\\
Je m'arrête ici, au bord de la reconstruction tentatrice. Car grande est la tentation de retoucher au texte même d'Andromède en y intégrant le iota d'ὀλυμπία et voir dans `δι᾽ Ὀλύμπου' une mention de Ζεύς Ὀλύμπιος\marqueur{3} ; qu'on se souvienne du vers 272 : Euripide jurait sur l'\italique{Éther, demeure de Zeus} (αἰθέρ᾽, οἴκησιν Διός\marqueur{4}). Le vrai Euripide parle ailleurs du \italique{ciel étoilé de Zeus} (\italique{Ion}, v. 1077 sq.) :%
\citComm{καὶ Διὸς ἀστερωπὸς\\
ἀνεχόρευσεν αἰθήρ.}%
La mention de Zeus lorsqu'un poète évoque le ciel constellé est d'ailleurs très naturelle puisque le quatrième astre le plus lumineux, après le Soleil, la lune et Vénus, n'est autre que \italique{Jupiter}, Διὸς ἀστήρ.
%279 crux ἱέρας· σημεῖον
%La première réplique ἐπικοκκάστρια de la nymphe Écho est-elle une consigne de διδάσκαλος, \italique{metteur en scène} ?\\
%En tout cas, le verbe qui répond le mieux au δι᾽ Ὀλύμπου est le verbe διαλάμπειν, \italique{briller}. (On verra que les assonances reposent essentiellement sur les consonnes.) Le verbe λάμπειν peut s'employer en parlant de la voix : παιὰν λάμπει (Soph.), \italique{le péan retentit}. Aristote (\italique{Problemata}, 922a31-39) utilise le composé dans l'un des ses \italique{problèmes} : Pourquoi ceux qui chantent en grand nombre gardent-ils mieux le rythme que ceux qui sont peu nombreux ? Διὰ τί οἱ πολλοὶ ᾄδοντες σώζουσι μᾶλλον τὸν ῥυθμὸν ἢ οἱ ὀλίγοι; 
%διολύγιον: ἐξηχοῦν ἐπὶ πολύ 
%μογ-είω, = μογέω, Dorian acc. to Hsch. (leg. μογίοντι).
% Exemples d'erreurs dans un lemme chez Hésychius où le 'ou' final est transformé en autre chose ('on')
%606 ἀγλαοπίστου (on):λαμπροπίστου
%1450 ἀζείρου:ἀζώστου. ἢ πολυζώστου (ἀζείριος: ἄζωτος Κ)
%3614 ἀποδοῦ (uson):ἀπόλυσον 
%8974 ἀψορρόου (on) :ἀνάπαλινῥέοντος,ἰταχυρρόου
%7255 εὐτυκάζου (on):εὔτυκον ἔχε, ἕτοιμον ()
%7355 "ἐφ᾽ἑαυτοῦ:καθ᾽ἑαυτοῦ (on)
%3693 προσβαλοῦ (on) :προσκατακλίθητι
%110  χαλκοπαρήου (reih):χαλκᾶςπαρειὰςἐχούσης
%En nombre, aucun des chanteurs ne \italique{brillerait}, c'est-à-dire ne \italique{se ferait remarquer} en s'élevant au-dessus de la masse : οὐδεὶς ἂν αὐτῶν διαλάμψειεν ὑπεράρας τὸ πλῆθος. Au contraire, moins nombreux, il leur est plus facile de \italique{briller} (et les chanteurs de rivaliser entre eux plutôt que de se régler sur le chef), ἐν δὲ τοῖς ὀλίγοις μᾶλλον διαλάμπουσιν· διὸ καθ’ αὑτοὺς ἐν αὐτοῖς μᾶλλον ἀγωνίζονται ἢ πρὸς τὸν ἡγεμόνα. Le verbe διαλάμπειν s'emploie donc pour des chanteurs qui \italique{brillent}, notamment par des modulations de rythme. \citComm{τοῦ σεμνοτάτου δι᾽ Ὀλύμπου; --- διαλάμπου\\Plus haut ! }%
}%
\debutNote{1069}%
\note{1}{Cf. le passionnant livre posthume de Pierre Boudreaux, \italique{Le texte d'Aristophane et ses commentateurs}.}%
\note{2}{Fragment d'Aristophane 528 K.-A.(515 Kock). Cf. l'édition Kassel-Austin des \italique{Poetae Comici Graeci} pour la juste distinction entre δοκός oxyton et δόκος paroxyton.% \eff{(Le \LSJ est à corriger sur ce point.)}
}%
\note{3}{Διός Ὀλυμπίου \italique{contra metrum} cependant.}%
\note{4}{Aristophane y revient au vers 100 des \italique{Grenouilles} : αἰθέρα Διὸς δωμάτιον.}%
%\commentaire{1069-72}{δι᾽ Ὀλύμπου — μέρος ἐξέλαχον — θανάτου τλήμων.}{Justifions notre modification du texte ainsi que la traduction. En philologie, tout finit en consultant son Hésychius (ε 2507) :\eff{Jeu du diable boîteux voir Louis Beck de Fouqières}\citComm{Ἔμπουσα· φάσμα δαιμονιῶδες ὑπὸ Ἑκάτης ἐπιπεμπόμενον καί, ὥς τινες, ἑνὶ ποδὶ χρώμενον. Ἀριστοφάνης δὲ τὴν Ἑκάτην ἔφη Ἔμπουσαν.\\Empuse : phantasme démonique envoyé par Hécate et, selon certains, n'ayant qu'une seule jambe. Aristophane appelle Hécate Empuse.}\italique{Nomen omen}, l'Empuse n'a qu'une seule jambe ; ἔμπους (singulier) suggère une entrave à \italique{une seule} jambe quand ἐμποδίζειν consiste à entraver les deux jambes d'\italique{impedementa}. Par dérivation de sens, que l'\italique{empuse} ait pu désigner un \italique{carcan}, cela nous paraît hautement probable au vu de la glose suivante d'Hésychius (ε 1264) :\citComm{ἑκάτη· ξύλον ἐν τοῖς φυλακίοις ᾧ τοὺς κακούργους προσδεσμεύοντες ἐμαστίγουν.\\hécate : morceau de bois dans les corps de garde auquel on attachait les malfaiteurs pour la flagellation.}L'explication qui me paraît la plus simple est bien qu'ἑκάτη ait servi de permutant métaphorique à ἔμπουσα -- Ἔμπουσα et Ἑκάτη étant parfois synonymes. (De toute façon, un glissement métaphorique dans l'autre sens, de ἑκάτη vers ἔμπουσα, suffirait tout autant à notre explication.)\\Tout cela pour justifier notre traduction : Euripide enjoint Mnésiloque de se libérer de ses entraves (qu'il s'aide et du ciel où il est attaché en costume d'Écho, il pourra l'aider ?) :\citComm{δι᾽ Ὀλύμπου — διάλυ(ε) (ἔ)μπουσ(αν).\\Sépare toi de cette empuse.}Voilà un exercice périlleux que de se livrer à des tentatives de reconstruction de texte sur ces trois répliques d'Écho pour faire apparaître d'hypothétiques jeux de mots ; laisser ainsi libre cours à l'imagination peut exposer à de rudes retours de bâton ; pour Aristophane, nous devons cependant impérieusement nous y livrer puisqu'il vient de nous prévenir que les échos seraient \italique{moqueurs}. Il nous semble bien ici qu'Aristophane joue sur le gore pour décrédibiliser la solennité de la pièce d'Euripide ; l'intuition provient de μέρος ἐξέλαχον qu'il est tentant de découper ainsi : μέρος ἐκ σέλαχος, \italique{la part du requin} (τό σέλαχος désignant tout type de poisson cartilagineux, raies, requins.) ; en effet, n'oublions pas qu'Andromède est attachée à une falaise pour servir de pâture à un monstre marin : le κῆτος dont il est question quelques vers plus haut. Aristophane, ce nous semble, s'adonne à un jeu de massacre en décrivant ce qu'est le sort attendu d'Andromède. Le reste se reconstruit selon cette perspective : δι᾽ Ὀλύμπου μέρος ἐξέλαχον θανάτου τλήμων $\longrightarrow$ διό λυποῦ μέρος ἐκ σέλαχος θαλάττης τέμνον « Et pour cette raison (quelle raison ?), tu seras chagrinée ; la part du requin cisailleur sorti de mer ». Pas très satisfaisant... à retoucher ! Pour ce qui est de l'apposition entre μέρος et σέλαχος, fréquente chez Euripide, voir Austin en français}%\debutNote{1069}%\note{1}{Voir, par exemple, le vers 293 des \italique{Grenouilles} et le scholiaste \italique{ad. loc.}.
%Tossi 2202 ou encore les vers 700-705 de \italique{Lysistrata} : Aristophane joue entre l'invocation à Hécate et le proverbe (Érasme n°394) ἀπ' οὐρᾶς τὴν ἔγχελυν ἔχεις, \italique{tu tiens l'anguille par la queue}, expliquant le singulier τοῦ σκέλους (ἐκτραχηλίσῃ était-il le terme propre d'halieutique pour la mise à mort du poisson pêché ?).}%
\\\passeLigne}%
\end{groupe}%
%% v. 1070-1071a
\begin{groupe}%
\personnageGrec{Μνησίλοχος}%
\anap{τί ποτ᾽ Ἀνδρομέδα περίαλλα κακῶν\apparat{1070}{περὶ ἄλλα R : \lecon{περίαλλα} Ellebodius}\\
μέρος ἐξέλαχον;\\\passeLigne}%
\end{groupe}%
%% v. 1071b
\begin{groupe}%
\personnageGrec{Εὐριπίδης}%
\anap{\ligne{1071}\decalage{μέρος ἐξέλαχον;}μόρος ἔξελ᾽ ἄχος.\apparat{1071}{μέρος ἐξέλαχον; R : \lecon{μόρος ἔξελ᾽ ἄχος.} Bibi}%
\commentaire{1071}{μέρος ἐξέλαχον;}{Sans \italique{tordre} le texte, il me semble que la réplique ἐπικοκκάστρια d'Écho s'entend assez bien -- soit dit avec toute la prudence qu'une telle chirurgie du texte impose :%
\citComm{μέρος ἐξέλαχον --- μόρος ἔξελ(ε)\marqueur{1} ἄχος\\
La Mort\marqueur{2} extirpe la douleur.}
Andromède se plaint des malheurs qui l'accablent ; Écho la rassure à sa manière : la mort l'en délivrera !.. Était-ce ce genre de répliques que le scholiaste qualifiait de `κακοστένακτος' ?\\
On trouve un tour analogue chez Sophocle\marqueur{3} :%
\citComm{Ἀλλ᾽ ἔσθ᾽ ὁ θάνατος λοῖσθος ἰατρὸς νόσων.}%
%Ταῖς νόσοις ὁ ϑάνατος ἀλλήλαις ἐπιλαβανούσαις οὐδ᾽ αὐτὸς παρεῖναι ἀναβάλλεται:.Μοσχίωνο
Et dans \italique{la langue des dieux} :%
\citComm{\textbf{Le trépas vient tout guérir} ;\\
Mais ne bougeons d'où nous sommes.\\
Plutôt souffrir que mourir,\\
C'est la devise des hommes.}
%\eff{maux sans guère de répit
%La mort guérit des maux
%la malheur t'a guère de répit
%la mâle-mort à son heure/ oeuvre}
}%
\debutNote{1071}%
\note{1}{Aoriste \italique{gnomique}.}%
\note{2}{Cette acception de μόρος a été relevée par à peu près tous les lexicographes anciens : Hésychius μ 1683 ; Photius μ 538 ; Συναγωγὴ λέξεων χρησίμων μ 267 ; Suidas μ 1257 ; \latin{Etymologicum genuinum} \italique{sub voce} μόρος· ὁ θάνατος ; \&c.}%
\note{3}{$=$ fr. 698 Radt $=$ fr. 636 Dindorf. Stobée (IV, 52, 26) a inséré la citation à son \italique{Florilège} et Michel Apostolius (II, 52c) la reprit pour ses \italique{Parœmiæ}.}\\\passeLigne}%
\end{groupe}%
%% v. 1072a
\begin{groupe}%
\personnageGrec{Μνησίλοχος}%
\anap{θανάτου τλήμων\\\passeLigne}%
\end{groupe}%
%% v. 1072b
\begin{groupe}%
\personnageGrec{Εὐριπίδης}%
\anap{\ligne{1072}\decalage{θανάτου τλήμων}†θανάτου τλήμων†.%\apparat{1072}{θανάτου τλήμων R : θαλάττης τέμνον Bibi}%
\commentaire{1072}{θανάτου τλήμων...}{L'on connaît par une scholie la fin du vers original d'\italique{Andromède} :%
\citComm{λείπει `μέλλουσα τυχεῖν'.}}\\\passeLigne}%
\end{groupe}%
%% v. 1073
\begin{groupe}%
\personnageGrec{Μνησίλοχος}%
\anap{ἀπολεῖς μ᾽ ὦ γραῦ στωμυλλομένη.%
\commentaire{1073-5}{ὦ γραῦ στωμυλλωμένη — στωμυλωμένη — νὴ Δί᾽.}{Le manuscrit de Ravenne (qui, rappelons-le, est le seul témoin complet de notre pièce) comporte deux orthographes différentes pour στωμυλλωμένη : avec un double lambda -λλ- au vers 1073 et un unique -λ- au vers 1074... Étourderie ou fantaisie de notre scribe jonglant avec une orthographe dont il est peu sûr ? Possible... ou bien preuve q'il recopiait scrupuleusement un texte sans parfois – comme ici – le comprendre ? c'est ce que nous allons tenter de montrer : il se trouve que le verbe στωμύλλω, \italique{babiller}, \italique{déblatérer}, s'écrit effectivement avec deux lambdas en grec classique. Cependant, l'adjectif στωμύλος (\italique{qui parle beaucoup}, et souvent avec nuance méliorative par rapport au verbe) qui en dérive n'en comporte qu'un seul ! Si donc, l'on fait le pari que le scribe ne s'y est pas trompé en conservant l'orthographe, savoir : on est en présence d'un participe moyen-passif au vers 1073 et de l'adjectif au vers 1074, il est alors aisé de retrouver la véritable réponse \italique{ἐπικοκκάστρια} d'Écho : στωμυλιο-μενής, « [je suis] de tempérament prolixe ».
Le composé est de la même facture que στωμυλιο-συλλεκτάδης, hapax qu'Aristophane mit dans la bouche d'Eschyle au vers 839 des \italique{Grenouilles} pour clore le bec à Euripide, cet intarissable \italique{collectionneur de commérages}. La seconde partie du composé est forgée avec le suffixe -μενής, \italique{au tempérament}, \italique{à la disposition...} (sur le modèle d'εὐ-μενής, δυσ-μενής, πρευ-μενής \&c).\\
Notre conjecture est renforcée par la concessive de Mnésiloque au vers suivant : il abonde, νὴ Δί᾽, ... néanmoins, γ᾽ : « ça oui par Zeus ! du moins... ».%
% Même scansion qu'aithéros iéras
}%
\apparat{1073}{στωμυλλωμένη R : \lecon{στωμυλλομένη} Grynaeus}\\\passeLigne}%
\end{groupe}%
%% v. 1074
\begin{groupe}%
\personnageGrec{Εὐριπίδης}%
\anap{στωμυλιομενής.\apparat{1074}{στωμυλωμένη R : \lecon{στωμυλιομενής} Bibi}\\\passeLigne}%
\end{groupe}%
%% v. 1075-1076a
\begin{groupe}%
\personnageGrec{Μνησίλοχος}%
\anap{νὴ Δί᾽, ὀχληρά γ᾽ εἰσήρρηκας\\
λίαν.\\\passeLigne}%
\end{groupe}%
%% v. 1076b
\begin{groupe}%
\personnageGrec{Εὐριπίδης}%
\anap{\ligne{1076}\decalage{λίαν.}†λίαν†.\\\passeLigne}%
\end{groupe}%
%% v. 1077-1078a
\begin{groupe}%
\personnageGrec{Μνησίλοχος}%
\anap{ὦγάθ᾽ ἔασόν με μονῳδῆσαι,\\
καὶ χαριεῖ μοι παῦσαι.%\apparat{1078}{Μνη. χαριεῖ μοι παῦσαι --- Εὐρ. παῦσαι R : \lecon{Μνη. χαριεῖ μοι ᾽παΰσαι --- simul Μνη. et Εὐρ. παῦσαι} Bibi}%
\commentaire{1078}{καὶ χαριεῖ μοι παῦσαι --- παῦσαι.}{Mnésiloque-Andromède, qui ne peut chanter sa monodie correctement, demande à Euripide-Écho de cesser de l'interrompre :%
\citComm{καὶ χαριεῖ μοι παῦσαι...}%
«~Tu m'obligeras d'arrêter...~» -- «~De crier~» interrompt facilement Écho. (Et, bien-sûr, en le lui criant à l'oreille :)%
\citComm{ (ἐ)παΰσαι.}%
(Le verbe ἐπαΰω, attesté par Eschyle dans ses \italique{Choéphores} (v. 828), est forgé avec le préfixe ἐπι- à l'instar d'ἐπικοκκάστρια et le verbe αὔω, \italique{crier} (cf. le vers 1019).)\\
Une telle provocation explique la réaction d'emportement de Mnésiloque :%
\citComm{βάλλ᾽ ἐς κόρακας.}%
%Andromède, qui ne peut chanter sa monodie correctement, demande à Écho de cesser de l'interrompre.
%	 par \italique{retour} (ἐπι-) de \italique{cri} (αὔω, cf. le vers 1019). 
%Mnésiloque s'apprêtait donc à dire :%
%\citComm{καὶ χαριεῖ μοι (ἐ)παΰσαι παῦσαι.}%
%Évidemment, cette répétition de deux `παυσαι' est entrecoupée par Écho.\\
%Pour le sens qu'il faut donner à son bien peu compatissant παῦσαι ; l'on le rapprochera du vers 1382 des \italique{Oiseaux} :%
%\citComm{παῦσαι μελῳδῶν, ἀλλ' ὅ τι λέγεις εἰπέ μοι.}%
%Ou encore de ce fragment attribué à Aristophane (962 (adesp. 601) Kassel-Austin = Photius π 519) :%
%\citComm{παῦσαι μελῳδοῦσ', ἀλλὰ πεζῇ μοι φράσον}
%%
%La réplique déclenche pas moins la fureur de Mnésiloque : βάλλ᾽ ἐς κόρακας.%
%παῦσαι peut être à la fois un infinitif et un impératif aoriste ; et c'est la-dessus qu'Aristophane nous semble jouer. La plupart des éditeurs ponctuent avec une pause entre χαριεῖ μοι et παῦσαι – et inclinent donc vers un double impératif avec le sens de « finis ! », c'est-à-dire « tais-toi ! ». Pourtant, juste après, Mnésiloque-Andromède demande, à être libéré de ses entraves, il est donc bien plus indiqué de comprendre un appel au secours – interrompu par Écho – de Mnésiloque, une supplique pour une délivrance morale avant la délivrance physique du vers suivant (ainsi, chez Sophocle, \italique{Électre}, 987, παῦσον ἐκ κακῶν ἐμέ ; ou dans l'\italique{Odyssée}, XXIV, 384, παύσαντο πόνου). On peut donc traduire ainsi : « Et tu m'obligeras en abrégeant...[mon tourment ? mes peines ?] — Abrège ! »
}\\\passeLigne}%
\end{groupe}%
%% v. 1078b
\begin{groupe}%
\personnageGrec{Εὐριπίδης}%
\anap{\ligne{1078}\decalage{καὶ χαριεῖ μοι παῦσαι.}᾽παΰσαι.\apparat{1078}{(Εὐρ.) παῦσαι R : \lecon{(Εὐρ.) ᾽παΰσαι} Bibi}%
\\\passeLigne}%
\end{groupe}%
%% v. 1079a
\begin{groupe}%
\personnageGrec{Μνησίλοχος}%
\anap{βάλλ᾽ ἐς κόρακας.%
\commentaire{1079}{βάλλ᾽ ἐς κόρακας — βάλλ᾽ ἐς κόρακας.}{Dans notre entreprise fort périlleuse de dénicher des jeux de mots tout au long de ce passage, en voici un qui est moins spéculatif et qui va nous permettre de reprendre haleine ! %
%\eff{Nous avons déjà analysé au vers 942 comment Aristophane jouait avec le mot κόραξ et ses deux sens possibles : le corbeau ou bien le crochet qui entrave Mnésiloque. }
Les Grecs usaient d'une formule d'imprécation : ἐς κόρακας, équivalente de notre « [va] au diable ! ». Aristophane l'emploie fréquemment avec le verbe βάλλειν : \italique{Guêpes} v. 835 ; \italique{Ploutos} v. 782 ; \italique{Nuées} v. 133 – digressons d'ailleurs sur ce passage des \italique{Nuées} puisqu'Aristophane recourt au même jeu de mots qu'ici : Strepsiade arrive devant le pensoir de Socrate et toque à la porte, il se fait recevoir par un peu accueillant :%
\citComm{Βάλλ᾽ ἐς κόρακας· / Τίς ἐσθ᾽ ὁ κόψας τήν θύραν;}%
signalons à toutes fins utiles que les κόρακες dont parle le disciple de Socrate sont des marteaux de porte\marqueur{1}. 
Le disciple tint donc à peu près ce langage : %
\citComm{Βάλλε κόρακας·/ Τίς ἐσθ᾽ ὁ κόψας τὴν θύραν;\\
Bas les marteaux ! Qui est ce tambourineur de porte ?}%
(compatible avec le mètre iambique et surtout avec le bon sens : qui – fût-il disciple de Socrate – s'aviserait de demander l'identité de quelqu'un qu'on vient d'envoyer se faire voir chez les corbeaux grecs ?)\\
Il en va de même dans le passage qui nous intéresse présentement : Écho conseille frivolement à Mnésiloque de retirer les crochets qui le retiennent : « Mais va te faire voir ! — Et toi mets bas le fermoir ! »}%
\debutNote{1079}%
\note{1}{Sens attesté par Pollux, X, 22.20.}%
\\\passeLigne}%
\end{groupe}%
%% v. 1079b
\begin{groupe}%
\personnageGrec{Εὐριπίδης}%
\anap{\ligne{1079}\decalage{βάλλ᾽ ἐς κόρακας.}βάλλε σύ κόρακας.\apparat{1079}{βάλλ᾽ ἐς κόρακας R : \lecon{βάλλε σύ κόρακας} Bibi}\\\passeLigne}%
\end{groupe}%
%% v. 1080a
\begin{groupe}%
\personnageGrec{Μνησίλοχος}%
\anap{†τί τὸ κακόν†;\apparat{1080}{\lecon{τί τὸ κακόν} bis R : τί κακόν bis Bentley}%
\commentaire{1080}{τί τὸ κακόν;}{La leçon du \italique{Ravennnas}, «~τί τὸ κακόν; --- τί τὸ κακόν;~», ne respecte pas le mètre anapestique. Gardons-nous pour autant de supprimer avec Bentley l'article `τὸ'. 
Si la formule «~τί τὸ κακόν;~» (ou d'autres très approchantes) se rencontre couramment chez Aristophane\marqueur{1}, sa version lapidaire, «~τί κακόν;~», elle, ne se trouve jamais\marqueur{2}. La bonne solution n'est donc pas celle de Bentley, et pour notre part nous nous résignons à \latin{crux}-désespérer †%
%κάγκελον, ου (τὸ) c. le suiv. Sch.-Ar. Eq. 638, 672.
%καγχάζω
%Hésychius 346 †κακόρας: κατακόψας. παρὰ Εὔκλῳ 
}%
\debutNote{1080}%
\note{1}{\italique{Thesmophories}, v. 610 ; \italique{Paix}, v. 322 ; \italique{Oiseaux}, v. 1213 ; \italique{Guêpes} v. 1136, τουτὶ τὸ κακὸν τί ἐστι ; \italique{Paix}, v. 181.}%
\note{2}{Certains éditeurs prétendent que si : au vers 1085. En vérité, ces faussaires commettent là une double forfaiture en «~corrigeant~» le présent vers \italique{metri gratia} ainsi que le vers 1085 ; le second forfait servant d'alibi au premier : au vers 1085, ils font apparaître «~une seconde fois~» le tour `τί κακόν' alors que le manuscrit porte `σί κακόν', confirmé par le scholiaste (ἀντὶ τοῦ `σοὶ κακόν'). Et quand bien même : c'est dans la bouche d'un Scythe qu'Aristophane la met au vers 1085, l'argument pour la faire dire aussi à Mnésiloque est irrecevable. En dernière analyse, ces éditeurs manquent cruellement de bon sens : le Scythe en son jargon ne tient déjà aucun compte des aspirations, de la grammaire, il ignore sa conjugaison et ses déclinaisons, il iotacise à l'envi... et l'on voudrait encore qu'il sigmatise ses `τ' ?! un tel massacre du grec aurait achevé d'être tout-à-fait inintelligible aux spectateurs.}\\\passeLigne}%
\end{groupe}%
%% v. 1080b
\begin{groupe}%
\personnageGrec{Εὐριπίδης}%
\anap{\ligne{1080}\decalage{†τί τὸ κακόν†;}†τί τὸ κακόν;†%
\commentaire{1080-2}{τί τὸ κακόν; [...] ὀτότυζ᾽.}{Quel est le jeu de scène de tout ce passage, dont les répliques d'Écho demeurent obscures ? Il me semble qu'on peut peut y répondre avec vraisemblance : Euripide-Écho s'est approchée de Mnésiloque (cf. v. 1090, γυνὴ πλησίον αὕτη) ; la raison en est évidente : Euripide tente de libérer Mnésiloque. 
Écho parvient cependant à se cacher dans un premier temps du Scythe (οὖτος τί λαλῖς; --- πῶτε τὸ πωνή;) ; c'est l'iconographie qui nous fournit la réponse : l'acteur jouant Écho se cachait derrière un grand voile\marqueur{1}.
Et lorsque l'archer Scythe veut mettre la main sur Euripide, ce dernier \italique{s'envole} (ποῖ ποῖ πεύγει -- certainement à l'aide de la μηχανή), lançant un narquois `οὐχ αἱρήσεις'.}%
\debutNote{1080-2}%
\note{1}{Sur un vase attique à figure rouge du \nieme{IV} siècle av. J.-C. (n° 7864 au catalogue de Beazley) représentant Dionysos et sa suite, on voit la nymphe Écho, ailée, se voiler le visage derrière son manteau. Ce détail est reproduit en annexe.}\\\passeLigne}%
\end{groupe}%
%% v. 1080c
\begin{groupe}%
\personnageGrec{Μνησίλοχος}%
\anap{\ligne{1080}\decalage{†τί τὸ κακόν†;†τί τὸ κακόν;†}ληρεῖς.\\\passeLigne}%
\end{groupe}%
%% v. 1080d
\begin{groupe}%
\personnageGrec{Εὐριπίδης}%
\anap{\ligne{1080}\decalage{†τί τὸ κακόν†;†τί τὸ κακόν;†ληρεῖς.}†ληρεῖς†.\\\passeLigne}%
\end{groupe}%
%% v. 1081a
\begin{groupe}%
\personnageGrec{Μνησίλοχος}%
\anap{οἴμωζ᾽.\\\passeLigne}%
\end{groupe}%
%% v. 1081b
\begin{groupe}%
\personnageGrec{Εὐριπίδης}%
\anap{\ligne{1081}\decalage{οἴμωζ᾽.}†οἴμωζ᾽†.\\\passeLigne}%
\end{groupe}%
%% v. 1081c
\begin{groupe}%
\personnageGrec{Μνησίλοχος}%
\anap{\ligne{1081}\decalage{οἴμωζ᾽.†οἴμωζ᾽†.}ὀτότυζ᾽.\\\passeLigne}%
\end{groupe}%
%% v. 1081d
\begin{groupe}%
\personnageGrec{Εὐριπίδης}%
\anap{\ligne{1081}\decalage{οἴμωζ᾽.†οἴμωζ᾽†.ὀτότυζ᾽.}†ὀτότυζ᾽†.%\apparat{1082}{\lecon{ὀτότυζ᾽} R : ὀλόφυγ᾽ Bibi}%
\commentaire{1082}{ὀτότυζ᾽ --- ὀτότυζ᾽.}{Signalons à l'attention du lecteur -- que nous espérons exégète à ses heures -- qu'Aristophane invente le gentilé Ὀτοτύξιοι au vers 1043 des \italique{Oiseaux} par jeu de consonance avec Ὀλοφύξιοι (v. 1042), \italique{habitants d'Olophyxos}. Y a-t-il ici un jeu analogue ?}\\\passeLigne}%
\end{groupe}%
%% v. 1083a
\begin{groupe}%
\personnageGrec{Τοξότης}%
\anap{\ligne{1083}οὖτος τί λαλῖς;\apparat{1083}{οῦτος οὗτος R : \lecon{οὖτος} bis Brunck \espApp λαλεῖς bis \Rac : \lecon{λαλῖς} bis \Rdeux, Schol.}%
\commentaire{1083}{λαλῖς.}{Scholie à propos de la graphie :%
\citComm{χωρὶς τοῦ ε γράφεται· ὁ γὰρ Σκύθης βαρβαρίζει.}}\\\passeLigne}%
\end{groupe}%
%% v. 1083b
\begin{groupe}%
\personnageGrec{Εὐριπίδης}%
\anap{\ligne{1083}\decalage{οὖτος τί λαλῖς;}†οὖτος τί λαλῖς;†\\\passeLigne}%
\end{groupe}%
%% v. 1084a
\begin{groupe}%
\personnageGrec{Τοξότης}%
\anap{πρυτάνεις καλέσω;\\\passeLigne}%
\end{groupe}%
%% v. 1084b
\begin{groupe}%
\personnageGrec{Εὐριπίδης}%
\anap{\ligne{1084}\decalage{πρυτάνεις καλέσω;}†πρυτάνεις καλέσω;†\\\passeLigne}%
\end{groupe}%
%% v. 1085a
\begin{groupe}%
\personnageGrec{Τοξότης}%
\anap{σί κακόν;%\apparat{1085}{σί bis R, sch. : τί Coulon}
\commentaire{1085}{σί κακόν;}{Commentaire du scholiaste :%
\citComm{ἀντὶ τοῦ `σοὶ κακόν'.}}%
\\\passeLigne}%
\end{groupe}%
%% v. 1085b
\begin{groupe}%
\personnageGrec{Εὐριπίδης}%
\anap{\ligne{1085}\decalage{σί κακόν;}†σί κακόν;†\\\passeLigne}%
\end{groupe}%
%% v. 1086a
\begin{groupe}%
\personnageGrec{Τοξότης}%
\anap{†πῶτε† τὸ πωνή;\apparat{1086}{π\inRas{ρ}{}ωτε bis \Rpc : πωτε bis \Mon}%
\commentaire{1086}{πρωτε το πωνη;}{Le scholiaste interprète ainsi le jargon scythe :%
\citComm{πόθεν ἡ φωνή; θαυμάζει τὴν ἠχώ.}%
On lui objectera que le \italique{Ravennas} comporte `πρωτε'. Et quand bien même on adopterait la correction `πωτε' (qu'une main postérieure a corrigé sur le manuscrit -- peut-être d'après la scholie elle-même), qu'il n'en resterait pas moins qu'à aucun autre endroit le Scythe ne confond un omicron ο avec un oméga ω\marqueur{1}.%
%962 γρωθώνη:σαπρὰγραῦς.οἱ δὲτὴνπαλαιὰνὄνονἢ σῦν
}%
\debutNote{1086}%
\note{1}{Voir la seconde note du commentaire au vers 1080 à propos du nombre déjà bien suffisant de fautes commises par le Scythe. 
%Faut-il aller chercher une explication du côté d'Hésychius ? (γ 962) :%
%\citComm{γρωθώνη· σαπρὰ γραῦς· οἱ δὲ τὴν παλαιὰν ὄνον ἢ σῦν.}
}%
\\\passeLigne}%
\end{groupe}%
%% v. 1086b
\begin{groupe}%
\personnageGrec{Εὐριπίδης}%
\anap{\ligne{1086}\decalage{†πῶτε† τὸ πωνή;}†πῶτε τὸ πωνή;†%\commentaire{1086}{}{Hésychius : πρόπωνα: vεὐχρατῆ. εὔφημα. πρόχειρα, ἕτοιμα, ἀνεμπόδιστα}
\\\passeLigne}%
\end{groupe}%
%% v. 1087a
\begin{groupe}%
\personnageGrec{Τοξότης}%
\anap{σὺ λαλῖς;\apparat{1087}{\lecon{λαλῖς} bis R : λαλεῖς bis \Rpc}\\\passeLigne}%
\end{groupe}%
%% v. 1087b
\begin{groupe}%
\personnageGrec{Εὐριπίδης}%
\anap{\ligne{1087}\decalage{σὺ λαλῖς;}†σὺ λαλῖς;†\\\passeLigne}%
\end{groupe}%
%% v. 1088a
\begin{groupe}%
\personnageGrec{Τοξότης}%
\anap{κλαύσαι.\apparat{1088}{κλαύσαιμι bis R : κλαύσει bis Brunck : \lecon{κλαύσαι} bis Rogers : κλαῦσ' ἔτι bis Austin}\\\passeLigne}%
\end{groupe}%
%% v. 1088b
\begin{groupe}%
\personnageGrec{Εὐριπίδης}%
\anap{\ligne{1088}\decalage{κλαύσαι.}†κλαύσαι.†%\commentaire{1088}{}{Hésychius 2856 κλᾶις:μοχλός 2881 χλαῦσιν:ἀγῶνα,ἄνθοςτι}%
\\\passeLigne}%
\end{groupe}%
%% v. 1089a
\begin{groupe}%
\personnageGrec{Τοξότης}%
\anap{κἀκκάσκι μοι;\apparat{1089}{κακκασκι bis R : κἀγκάσκις bis van Leeuwen (<ς> add. Fritzsche) : κἀγκάσκι bis Coulon : alii alia}%
\commentaire{1089}{κακκασκι μοι;}{$=$ καὶ ἐγχάσκει μοι. Cf. \italique{Guêpes}, vers 721 et le scholiaste : καταγελᾷς μοι.}\\\passeLigne}%
\end{groupe}%
%% v. 1089b
\begin{groupe}%
\personnageGrec{Εὐριπίδης}%
\anap{\ligne{1089}\decalage{κἀκκάσκι μοι;}†κἀκκάσκι μοι;†\\\passeLigne}%
\end{groupe}%
%% v. 1090
\begin{groupe}%
\personnageGrec{Χορός}%
\anap{μὰ Δί᾽ ἀλλὰ γυνὴ πλησίον αὕτη.\apparat{1090}{(Μνη.) edd. plerique : \lecon{(Χορ.)} Bibi, cf. Schol.}%
\commentaire{1090}{μὰ Δί᾽ ἀλλὰ γυνὴ πλησίον αὕτη.}{C'est le chœur qui donne cette indication à l'archer désemparé, comme le suggérait timidement le scholiaste :%
\citComm{τοῦτο ὁ χορὸς ἢ ὁ γέρων.}%
Aucun éditeur cependant n'a suivi la première hypothèse du scholiaste, quoique la seule vraisemblable -- aussi agacé soit-il, Mnésiloque ne peut en aucun cas dénoncer son complice, son dernier espoir de survie.}\\\passeLigne}%
\end{groupe}%
%% v. 1091
\begin{groupe}%
\personnageGrec{Εὐριπίδης}%
\anap{πλήξασ᾿ αὕτη.\apparat{1091}{πλησίον αὕτη R : \lecon{πλήξασ᾿ αὕτη} Bibi}%
\commentaire{1091}{πλησίον αὕτη.}{Après un long calvaire parsemé de \latin{cruces desperationis}, il nous semble enfin saisir le sel ἐπικοκκάστρια de la réplique d'Écho, fournissant par la même occasion un indice sur l'action.\\
Sur l'indication du chœur, l'archer Scythe est en train d'inspecter \italique{à proximité} (πλησίον) de Mnésiloque. Il se rapproche effectivement d'Euripide (ce dernier est toujours camouflé derrière son manteau). 
Quand l'archer arrive à portée de bras, Euripide lui assène alors un coup : πλήξασ᾿ αὕτη.\\
Ce coup explique la fureur du Scythe, qui change de langage et passe à l'offensive : ποῦ ᾿στ᾽ ἠ μιαρά;}\\\passeLigne}%
\end{groupe}%
%% v. 1092-1094a
\begin{groupe}%
\personnageGrec{Τοξότης}%
\anap{ποῦ ᾿στ᾽ ἠ μιαρά;\apparat{1092}{'σθ' ἡ R : ᾿στ᾽ ἠ Brunck}\\\passeLigne}%
\end{groupe}%
%% v. 1092-1094a
\begin{groupe}%
\personnageGrec{Χορός}%
\anap{\decalage{ποῦ ᾿στ᾽ ἠ μιαρά;}καὶ δὴ φεύγει.\apparat{1092}{(Μνη.) edd. plerique : \lecon{(Χορ.)} Bibi}\\\passeLigne}%
\end{groupe}%
%% v. 1092-1094a
\begin{groupe}%
\personnageGrec{Τοξότης}%
\anap{ποῖ ποῖ πεύγεις;\apparat{1093}{φεύγεις R : \lecon{πεύγεις} Enger}\\\passeLigne}%
\end{groupe}%
%% v. 1092-1094a
\begin{groupe}%
\personnageGrec{Εὐριπίδης}%
\anap{\ligne{1092}\decalage{ποῖ ποῖ πεύγεις;}ποῖ ποῖ πέτει.\apparat{1093}{(Εὐρ.) ποῖ ποῖ πέτει om. R : \lecon{suppl.} Bibi}%
\commentaire{1093}{ποῖ ποῖ φεύγεις;}{Il y a longtemps que les scribes chargés de recopier cette scène n'en saisissaient plus le comique ; ne nous étonnons donc pas qu'elle soit la plus altérée de toute la comédie : il manque plusieurs répliques (cf. le vers 1094 ; où la lacune est certaine et peut-être comblée grâce à une ancienne scholie). C'est le cas aussi ici.\\
Euripide s'est approché de Mnésiloque pour le délivrer. Sur les indications du chœur (μὰ Δί᾽ ἀλλὰ γυνὴ πλησίον αὕτη), le Scythe a fini par repérer Euripide, et ce malgré son camouflage. C'est ce coup qu'il est bon de partir ; et Euripide de déloger sans trompette (καὶ δὴ φεύγει) mais non sans abandonner sa verve moqueuse. \\
Remettons-nous bien la scène : Euripide incarne la nymphe ailée Écho -- sa fuite est donc aérienne ; et sur scène, l'acteur est tiré dans les airs à l'aide de la μηχανή. La réplique ἐπικοκκάστρια d'Écho fait vraisemblablement allusion à cette \italique{envolée}. Or pour la circonstance, nous ne sommes pas tout-à-fait démunis puisque nous possédons une comédie entière, les \italique{Oiseaux}, avec des jeux théâtraux similaires.\\
En \italique{Oiseaux} 226, lorsque la Huppe se dispose à chanter, elle entonne : ἐποποποῖ ποποῖ, ποποποποῖ ποποῖ.\\
De même, lorsque le chœur au vers 310 répond à l'appel de la Huppe, c'est au cri de : ποποποποποπο ποῦ μ' ὃς ἐκάλεσε;\\
C'est ainsi qu'Aristophane transcrivait les onomatopées d'oiseaux\marqueur{1} ; on comprend mieux pourquoi lorqu'Iris descend du ciel au moyen de la μηχανή, Pisthétaire lui lance ces paroles ailées (vers 1199) : αὕτη σύ, ποῖ ποῖ ποῖ πέτει;\\
Il en va de même dans notre passage, alors que le Scythe poursuit Écho en train de s'envoler, il lui jette à la figure des cris d'oiseaux : ποῖ ποῖ πεύγεις; -- qui fait entendre ποῖ ποῖ πέτει -- qu'on rétablit donc comme réplique ἐπικοκκάστρια manquante d'Écho.}%
\debutNote{1093}%
\note{1}{Pour d'autres exemples de telles onomatopées chez Aristophane, voir l'article de Loredana Di Virgilio «~Some thoughts on humor and music in Aristophanes' plays~», in \italique{Mediterranean World}, 24 (2019), pp. 36 sqq.}\\\passeLigne}%
\end{groupe}%
%% v. 1092-1094a
\begin{groupe}%
\personnageGrec{Τοξότης}%
\anap{οὐ καιρήσεις.\apparat{1094}{(Τοξ.) οὐ καιρήσεις om. R : \lecon{suppl.} Brunck ap. Schol.}%
\commentaire{1094}{<οὐ καιρήσεις.> --- οὐκ᾽ αἰρήσεις.}{Le texte manquant a été rétabli à partir de la scholie :%
\citComm{ἀντὶ τοῦ οὐ χαιρήσεις.}%
\italique{Confer} la note \italique{supra} au vers 1056.}\\\passeLigne}%
\end{groupe}%
%% v. 1094b
\begin{groupe}%
\personnageGrec{Εὐριπίδης}%
\anap{\ligne{1094}\decalage{οὐ καιρήσεις.}οὐχ αἱρήσεις.\apparat{1094}{οὐκ᾽ αἰρήσεις; R : \lecon{οὐχ αἱρήσεις.} Wilamowitz (\italique{Kleine Schriften} IV, p.486n.1)}\\\passeLigne}%
\end{groupe}%
%% v. 1095a
\begin{groupe}%
\personnageGrec{Τοξότης}%
\anap{ἔτι γὰρ γρύζεις;\\\passeLigne}%
\end{groupe}%
%% v. 1095b
\begin{groupe}%
\personnageGrec{Εὐριπίδης}%
\anap{\ligne{1095}\decalage{ἔτι γὰρ γρύζεις;}†ἔτι γὰρ γρύζεις;†\\\passeLigne}%
\end{groupe}%
%% v. 1096a
\begin{groupe}%
\personnageGrec{Τοξότης}%
\anap{λαβὲ τὴ μιαρά.\apparat{1096}{τῇ et τῆ R : \lecon{τὴ} bis Bothe}\\\passeLigne}%
\end{groupe}%
%% v. 1096b
\begin{groupe}%
\personnageGrec{Εὐριπίδης}%
\anap{\ligne{1096}\decalage{λαβὲ τὴ μιαρά.}†λαβὲ τὴ μιαρά.†\\\passeLigne}%
\end{groupe}%
%% v. 1097
\begin{groupe}%
\personnageGrec{Τοξότης}%
\anap{λάλο καὶ κατάρατο γύναικο.\\\passeLigne}%
\end{groupe}%
%% v. 1098-1102a
\begin{groupe}%
\personnageGrec{Εὐριπίδης}%
\typeVers{trim. iamb.}\trimIamb{ὦ θεοὶ τίν᾽ εἰς γῆν βαρβάρων ἀφίγμεθα\apparat{1098}{ευρ\auDessus{ι}{π} ως περσευς in mg. add. \Rpc, \Mon}%
\commentaire{1098-102}{ὦ θεοὶ ... κάρα κομίζων.}{La scholie à propos de la tirade de Persée\marqueur{1} :%
\citComm{εἰς Περσέα ἐξ Ἀνδρομέδας τρία τὰ πρῶτα. καὶ λοιπὸν (i.e. vv. 1101-2) ἐπέζευξεν τὰ ἑξῆς.}%
}%
\debutNote{1098-102}%
\note{1}{Fr. 9 de l'\italique{Andromède} dans l'édition Jouan-van Looy $=$ fr. 124 Kannicht.}\\
ταχεῖ πεδίλῳ; διὰ μέσου γὰρ αἰθέρος\\
τέμνων κέλευθον πόδα τίθημ᾽ ὑπόπτερον\\
Περσεὺς πρὸς Ἄργος ναυστολῶν τὸ Γοργόνος\apparat{1101}{ναυτολῶν R : \lecon{ναυστολῶν} Canini (\italique{navigans} iam Divus)}\\
κάρα κομίζων.\\\passeLigne}%
\end{groupe}%
%% v. 1102b-1103a
\begin{groupe}%
\personnageGrec{Τοξότης}%
\trimIamb{\ligne{1102}\decalage{κάρα κομίζων.}τί λέγι; τὴ Γόργος πέρι\apparat{vv. 1102-3}{om R, \Mon : in mg. inferiore add. \Rdeux \ligneApp{1102} Γοργόνος R : \lecon{Γόργος} Schol., Bothe \ligneApp{1103} κεφαλή R : \lecon{κεπαλή} Bothe}%
\commentaire{1102-3}{τί λέγι; τὴ Γόργος πέρι / τὸ γραμματέο σὺ τὴ κεπαλή;}{$=$ τί λέγεις; τοῦ Γόργου φέρεις / τοῦ γραμματέως σὺ τὴν κεφαλήν;\\
Et la scholie (plus que douteuse) :%
\citComm{ὁ δὲ Γόργος γραμματεὺς, ἀλλὰ καὶ βάρβαρος.}%
%\eff{mais la scansion ?..}
Il n'y a en grec qu'une différence d'accent entre Γοργών (la Gorgone) et Γόργων (Gorgon).}\\
τὸ γραμματέο σὺ τὴ κεπαλή;\\\passeLigne}%
\end{groupe}%
%% v. 1103b-1104a
\begin{groupe}%
\ledpb%
\personnageGrec{Εὐριπίδης}%
\trimIamb{\ligne{1103}\decalage{τὸ γραμματέο σὺ τὴ κεπαλή;}τὴν Γοργόνος\\
ἔγωγε φημί.\\\passeLigne}%
\end{groupe}%
%% v. 1104b
\begin{groupe}%
\personnageGrec{Τοξότης}%
\trimIamb{\ligne{1104}\decalage{ἔγωγε φημί.}Γόργο τοι κἀγὼ λέγι.\\\passeLigne}%
\end{groupe}%
%% v. 1105-1106
\begin{groupe}%
\personnageGrec{Εὐριπίδης}%
\trimIamb{ἔα· τίν᾽ ὄχθον τόνδ᾽ ὁρῶ καὶ παρθένον%
\commentaire{1105}{ἔα· τίν᾽ ὄχθον τόνδ᾽ ὁρῶ καὶ παρθένον.}{On connaît le texte original d'Euripide grâce à une scholie\marqueur{1} :%
\citComm{πάλιν ἐξ Ἀνδρομέδας·\\
ἔα, τίν᾽ ὄχθον τόνδ᾽ ὁρῶ περίρρυτον\\
ἀφρῷ θαλάσσης παρθένου τ᾽ εἰκώ τίνα\\
$[$ ἐξ αὐτομόρφων λαΐνων τυκισμάτων,
σοφῆς ἄγαλμα χειρός; $]$ }%
}%
\debutNote{1105}%
\note{1}{Dixième fragment de l'\italique{Andromède} dans l'édition Jouan-van Looy $=$ fr. 125, 1-2 Kannicht. La scholie rapporte les seuls deux premiers vers ; le fragment a pu être complété en le recoupant avec d'autres témoignages.}\\
θεαῖς ὁμοίαν ναῦν ὅπως ὡρμισμένην;%
\commentaire{1106}{ναῦν ὅπως ὡρμισμένην;}{Quoique le scholiaste ne l'indique pas, le vers 1106 est probablement lui aussi repris de l'\italique{Andromède} d'Euripide. On trouve en tout cas la même image chez l'auteur d'\italique{Hercule Furieux} (vers 1094-1097), Hercule dont la situation était analogue à celle d'Andromède :%
\citComm{τί δεσμοῖς ναῦς ὅπως ὡρμισμένος\\
νεανίαν θώρακα καὶ βραχίονα,\\
πρὸς ἡμιθραύστῳ λαΐνῳ τυκίσματι\\
ἧμαι, νεκροῖσι γείτονας θάκους ἔχων;}%
}\\\passeLigne}%
\end{groupe}%
%% v. 1107-1108a
\begin{groupe}%
\personnageGrec{Μνησίλοχος}%
\trimIamb{ὦ ξένε κατοίκτιρόν με τὴν παναθλίαν,\apparat{1107}{κ\auDessus{η}{δ} ως ανδρομ\auDessus{ε}{δ} in mg. \Rpc, \Mon \espApp κατοικτειρόν R : \lecon{κατοίκτιρόν} van Herwerden}\\
λῦσόν με δεσμῶν.\\\passeLigne}%
\end{groupe}%
%% v. 1108b-1109
\begin{groupe}%
\personnageGrec{Τοξότης}%
\trimIamb{\ligne{1108}\decalage{λῦσόν με δεσμῶν.}οὐκὶ μὶ λαλῆσι σύ;\apparat{1108}{\lecon{μὶ λαλῆσι} R : μὴ λαλῆσι Brunck}\\
κατάρατο τολμᾷς ἀποτανουμένη λαλῖς;\apparat{1109}{λαλᾶς R : \lecon{λαλῖς;} Bergk (cf. vv. 1082 et 1087) : λαλᾷς; Hall-Geldart}\\\passeLigne}%
\end{groupe}%
%% v. 1110
\begin{groupe}%
\personnageGrec{Εὐριπίδης}%
\trimIamb{ὦ παρθέν᾽ οἰκτίρω σὲ κρεμαμένην ὁρῶν.\apparat{1110}{οἰκτείρω R : \lecon{οἰκτίρω} van Herwerden}\\\passeLigne}%
\end{groupe}%
%% v. 1111-1112a
\begin{groupe}%
\personnageGrec{Τοξότης}%
\trimIamb{οὐ παρτέν᾽ ἐστίν, ἀλλ᾽ ἀμαρτωλὴ γέρων\\
καὶ κλέπτο καὶ πανοῦργο.\\\passeLigne}%
\end{groupe}%
%% v. 1112b-1113
\begin{groupe}%
\personnageGrec{Εὐριπίδης}%
\trimIamb{\ligne{1112}\decalage{καὶ κλέπτο καὶ πανοῦργο.}ληρεῖς ὦ Σκύθα.\\
αὕτη γάρ ἐστιν Ἀνδρομέδα παῖς Κηφέως.\\\passeLigne}%
\end{groupe}%
%% v. 1114
\begin{groupe}%
\personnageGrec{Τοξότης}%
\trimIamb{σκέψαι τὸ κύστο· μή τι μικκὸν παίνεται;\apparat{1114}{σκύτο R : \lecon{κύστο} Scaliger grat. metr. : σῦκο Sommerstein \espApp μῆτι R : \lecon{μή τι} \Mon (\italique{numquid} Divus) \espApp μικτὸν R : \lecon{μικκὸν} Bentley : μικρὸ Enger \ligneApp{1115} δεῦρο δεῦρο R : \lecon{δεῦρο} Biset}%
\commentaire{1114}{σκέψαι τὸ σκύτο·}{Fait remarquable : l'erreur du scribe, copiant τὸ σκύτο en place de τὸ κύστο ($=$τὸν κύσθον), est la même qu'au vers 557 : métathèse du `s' (cf. le commentaire à ce vers). 
Le scholiaste devait avoir devant les yeux la bonne leçon qui précise : %
\citComm{δείκνυσιν αὐτῷ τὸ αἰδοῖον.}%
Second fait remarquable : en première rédaction, le \italique{Ravennas} faisait apparaître un `σ' supplémentaire au début d'un mot du vers 1187 : τὸ \textbf{σ}ποστίον à la de τὸ πόστιον ($=$ τὸ πόσθιον)...  on constate qu'il s'agit là-encore d'un organe sexuel -- mâle cette fois.}\\\passeLigne}%
\end{groupe}%
%% v. 1115-1118a
\begin{groupe}%
\personnageGrec{Εὐριπίδης}%
\trimIamb{φέρε δεῦρό μοι τὴν χεῖρ᾽, ἵν᾽ ἅψωμαι κόρης·\\
φέρε Σκύθ᾽· ἀνθρώποισι γὰρ νοσήματα\\
ἅπασίν ἐστιν· ἐμὲ δὲ καὐτὸν τῆς κόρης\\
ταύτης ἔρως εἴληφεν.\\\passeLigne}%
\end{groupe}%
%% v. 1118b-1120
\begin{groupe}%
\ledpb%
\personnageGrec{Τοξότης}%
\trimIamb{\ligne{1118}\decalage{ταύτης ἔρως εἴληφεν.}οὐ ζηλῶσί σε·\\
ἀτὰρ εἰ τὸ πρωκτὸ δεῦρο περιεστραμμένον,\apparat{1119}{τῶ πρωκτῶ R : \lecon{τὸ πρωκτὸ} Biset}%
%\commentaire{1019}{τῶ πρωκτῶ}{Le prokton peut donc être (mal-)employé au pluriel. \cf la scène où Mnésiloque a le feu aux fesses.}
\\
οὐκ ἐπτόνησά σ᾽ αὐτὸ πυγίζεις ἄγων.\apparat{1120}{ἐπτόνησας R : \lecon{ἐπτόνησά σ᾽} Bentley}\\\passeLigne}%
\end{groupe}%
%% v. 1121-1122
\begin{groupe}%
\personnageGrec{Εὐριπίδης}%
\trimIamb{τί δ᾽ οὐκ ἐᾷς λύσαντά μ᾽ αὐτὴν ὦ Σκύθα\\
πεσεῖν ἐς εὐνὴν καὶ γαμήλιον λέχος;\apparat{1122}{ἐστ᾽ R : ἐς τ᾽ \Mon : \lecon{ἐς} Grynaeus}\\\passeLigne}%
\end{groupe}%
%% v. 1123-1124
\begin{groupe}%
\personnageGrec{Τοξότης}%
\trimIamb{εἰ σπόδρ᾽ ἐπιτυμεῖς τὴ γέροντο πύγισο,\\
τὴ σανίδο τρήσας ἐξόπιστο πρώκτισον.\apparat{1124}{ἐξόπισθο R : \lecon{ἐξόπιστο} Daubuz}\\\passeLigne}%
\end{groupe}%
%% v. 1125a
\begin{groupe}%
\personnageGrec{Εὐριπίδης}%
\trimIamb{μὰ Δί᾽ ἀλλὰ λύσω δεσμά.\testimonia{1125}{(δεσμά) Moeris δ 28}\apparat{1125}{δέμας R : \lecon{δεσμά} Grynaeus}\\\passeLigne}%
\end{groupe}%
%% v. 1125b
\begin{groupe}%
\personnageGrec{Τοξότης}%
\trimIamb{\ligne{1125}\decalage{μὰ Δί᾽ ἀλλὰ λύσω δεσμά.}μαστιγῶ σ᾽ ἄρα.\\\passeLigne}%
\end{groupe}%
%% v. 1126a
\begin{groupe}%
\personnageGrec{Εὐριπίδης}%
\trimIamb{καὶ μὴν ποιήσω τοῦτο.\\\passeLigne}%
\end{groupe}%
%% v. 1126b-1127
\begin{groupe}%
\personnageGrec{Τοξότης}%
\trimIamb{\ligne{1126}\decalage{καὶ μὴν ποιήσω τοῦτο.}τὸ κεπαλή σ᾽ ἄρα\apparat{1126}{τὸ κεπαλῆς (ex καὶ παλῆς) \Rpc, \Mon : \lecon{τὸ κεπαλή σ᾽} Brunck \ligneApp{1127} ἀποκέκοιψ᾽ \Rac : ἀποκεκόψοι \Rpc : ἀποκεκόψο \Mon : \lecon{ἀποκεκόψι} Brunck}\\
τὸ ξιπομάκαιραν ἀποκεκόψι τουτοΐ.%\apparat{1127}{ἀποκέκοιψ᾽ \Rac : ἀποκεκόψοι \Rpc : ἀποκεκόψο \Mon : \lecon{ἀποκεκόψι} Brunck}%
\commentaire{1127}{τὸ ξιπομάκαιραν.}{La harpé de Persée. (Note de Willems.)}\\\passeLigne}%
\end{groupe}%
%% v. 1128-1132
\begin{groupe}%
\personnageGrec{Εὐριπίδης}%
\trimIamb{αἰαῖ· τί δράσω; πρὸς τίνας στρεφθῶ λόγους;\\
ἀλλ᾽ οὐκ ἂν ἐνδέξαιτο βάρβαρος φύσις.\apparat{1129}{ἂν δέξαιτο R : ἂν ἀνδέξαιτο Ellebodius : \lecon{ἂν ἐνδέξαιτο} Küster}\\
σκαιοῖσι γάρ τοι καινὰ προσφέρων σοφὰ%
\commentaire{1130}{σκαιοῖσι γάρ τοι καινὰ προσφέρων σοφὰ.}{Cf. le reproche qu'Euripide mit dans la bouche de Médée dans la tragédie de ce nom (vers 298) :%
\citComm{σκαιοῖσι μὲν γὰρ καινὰ προσφέρων σοφὰ.}}\\
μάτην ἀναλίσκοις ἄν, ἀλλ᾽ ἄλλην τινὰ\\
τούτῳ πρέπουσαν μηχανὴν προσοιστέον.\apparat{1132}{τοῦτο R : \lecon{τούτῳ} Zanetti (\italique{huic} Divus)}\\\passeLigne}%
\end{groupe}%
%% v. 1133
\begin{groupe}%
\personnageGrec{Τοξότης}%
\trimIamb{μιαρὸς ἀλώπηξ, οἶον ἐπιτήκιζί μοι.\apparat{1133}{οἷον R : \lecon{οἶον} Brunck \espApp ἐπιτηκίζει R : \lecon{ἐπιτήκιζί} anonymus Parisinus}%
\commentaire{1133}{οἷον ἐπιτηκίζει μοι}{$=$ οἷον ἐπιθήκιζέ μοι. À propos des images animales du trompeur : du renard, ἡ ἀλώπηξ, et du singe, ὁ πίθηκος, vus comme les représentant de la ruse et de la fourbe, le grand Taillardat a consacré deux paragraphes auxquels je renvoie le lecteur : §§405-406. Notons que sur le même modèle que πιθηκίζειν (qui se rencontre aussi au vers 1290 des \italique{Guêpes}), il existe chez Aristophane le verbe ἀλωπεκίζειν (\italique{Guêpes}, vers 1241).}\\\passeLigne}%
\end{groupe}%
%% v. 1134
\begin{groupe}%
\personnageGrec{Μνησίλοχος}%
\trimIamb{μέμνησο Περσεῦ μ᾽ ὡς καταλείπεις ἀθλίαν.\\\passeLigne}%
\end{groupe}%
%% v. 1135
\begin{groupe}%
\personnageGrec{Τοξότης}%
\trimIamb{ἔτι γὰρ σὺ τὴ μάστιγαν ἐπιτυμεῖς λαβεῖν;\apparat{1135}{τῆ μαστιγᾶν (ex μαστηγαν) \Rpc \espApp λαβεῖν· R : \lecon{λαβεῖν;} Scaliger (\italique{consequi?} iam Ellebodius in vers. Lat.)}\\\passeLigne}%
\end{groupe}%
%% v. 1136-1159
\begin{groupe}%
\ledpb%
\personnageGrec{Χορός}%
\typeVers{ibyc.}\chantChoeur{Παλλάδα τὴν φιλόχορον ἐμοὶ%
\commentaire{1136-59}{Chant ibycéen du chœur.}{Il a été beaucoup discuté du caractère strophique ou a-strophique de ce chant. Si la balance des commentateurs penche plutôt en faveur d'un chant astrophique\marqueur{1}, il est tentant, à l'instar d'Hermann, von Velsen ou Prato, de décomposer ce chant en deux parties avec pour chacune la division classique du chœur en deux demi-chœurs (strophe -- antistrophe). Voici la structure que je propose :%
\begin{alignat*}
\text{1\expo{ère} partie } & \text{ : invocation d'Athéna, } && \text{ vv. 1140-1142 } &&&\text{ } = \text{ vv. 1145-1147 ; }\\
\text{2\expo{nde} partie } & \text{ : invocation des Thesmophores, } && \text{ vv. 1148-1153 } &&&\simeq\text{ vv. 1155-1159.}
\end{alignat*}
(Nous n'avons pas cependant forcé la concordance métrique au point de retoucher le texte lorsqu'il ne s'y prêtait que de mauvaise grâce ; si les vers 1140 ou 1148 peuvent se voir comme des glyconéens à base trissyllabique\marqueur{2}, ainsi, à la rigueur, que le vers 1155\marqueur{3}, il subsiste un problème de correspondance entre le vers 1153, ὄργια σεμνὰ θεαῖν ἵνα λαμπάσι, et le vers 1159, νῦν ἀφίκεσθον ἱκετεύομεν ἐνθάδ᾽ ἡμῖν.)\\
Mis de cette façon en ordre, le chant est entrecoupé entre chaque paire strophique d'un intermède, μεσῳδός, commençant par un impératif de φαίνω : aux vers 1143-4, φάνηθ᾽ ὦ τυράννους στυγοῦσ᾽ ὥσπερ εἰκός, et au vers 1154, φαίνετον ἄμβροτον ὄψιν.}%
\debutNote{1136-59}%
\note{1}{L'analyse la plus poussée de ce chant est celle de Lætitia, \italique{The Songs of Aristophanes}, pp. 446-453 ; on pourra aussi consulter les non moins magistraux ouvrages : \italique{The Lyric Metres of Greek Drama}\expo{2} (1968) pp. 166-167 d'Amy M. Dale ainsi que \italique{The Verse of Greek Comedy} (1912), §387, de White. White ne perçoit de répons antistrophique que pour les $6 = 3 + 3$ premiers vers.}%
\note{2}{Cf. Euripide, \italique{Iphigénie en Aulis}, vers 759.}%
\note{3}{Voir ce qu'en dit Amy, \italique{The Lyric Metres of Greek Drama}\expo{2}, p. 166 n.2.}}
\end{groupe}%
%% v. 1136-1159
\begin{groupe}%
\typeVers{ibyc.}\chantChoeur{δεῦρο καλεῖν νόμος εἰς χορόν,\apparat{1138}{ἐς R : \lecon{εἰς} von Velsen}}
\end{groupe}%
%% v. 1136-1159
\begin{groupe}%
\typeVers{ibyc\cat}\chantChoeur{\ligne{1139}παρθένον ἄζυγα κούρην,\testimonia{1139}{(ἄζυγα) \Phryn 100.3}\apparat{1139}{κόρην R : \lecon{κούρην} Bothe}\\\passeLigne}
\end{groupe}%
%% v. 1136-1159
\begin{groupe}%
\Strophe{(στροφή α)}\typeVers{glyc.}\chantChoeur{ἣ πόλιν ἡμετέραν ἔχει}
\end{groupe}%
%% v. 1136-1159
\begin{groupe}%
\typeVers{glyc.}\chantChoeur{καὶ κράτος φανερὸν μόνη}
\end{groupe}%
%% v. 1136-1159
\begin{groupe}%
\typeVers{pherecr.}\chantChoeur{κλῃδοῦχός τε καλεῖται.\apparat{post v. 1143}{στυγνὰς ὡσσε \Rac (cf. στυγοῦσ᾽ ὥσπερ v. sq. ?) : \lecon{del.} \Rpc : στυγνὰς ὡς ἐ \Mon}\\\passeLigne}
\end{groupe}%
%% v. 1136-1159
\begin{groupe}%
\Strophe{(μεσῳδός α)}\typeVers{$4$ bacch.}\chantChoeur{\ligne{1144}φάνηθ᾽ ὦ τυράννους στυγοῦσ᾽ ὥσπερ εἰκός.\\\passeLigne}
\end{groupe}%
%% v. 1136-1159
\begin{groupe}%
\Strophe{(ἀντιστρ. α)}\typeVers{glyc.}\chantChoeur{δῆμός τοί σε καλεῖ γυναι-}
\end{groupe}%
%% v. 1136-1159
\begin{groupe}%
\typeVers{$+$ glyc.}\chantChoeur{\alinea κῶν· ἔχουσα δέ μοι μόλοις}
\end{groupe}%
%% v. 1136-1159
\begin{groupe}%
\typeVers{pherecr.}\chantChoeur{εἰρήνην φιλέορτον.\\\passeLigne}
\end{groupe}%
%% v. 1136-1159
\begin{groupe}%
\ledpb%
\Strophe{(στροφή β)}\typeVers{glyc.}\chantChoeur{ἥκετέ τ᾽ εὔφρονες ἵλαοι,\apparat{1148}{ἥκετ᾽ R : \lecon{ἥκετέ τ᾽} Enger : ἥκετε δ᾽ Hermann}}
\end{groupe}%
%% v. 1136-1159
\begin{groupe}%
\typeVers{ibyc.}\chantChoeur{\ligne{1149}πότνιαι, ἄλσος ἐς ὑμέτερον,}
\end{groupe}%
%% v. 1136-1159
\begin{groupe}%
\typeVers{dact. aeol.}\chantChoeur{οὗ δὴ ἀνδράσιν οὐ θεμιτεύετον εἰσορᾶν\apparat{1150-1}{\lecon{οὗ δὴ ἀνδράσιν} R : οὗ δῆτ' ἀνδράσιν Burges : ἄνδρας ἵν(α) Hermann \espApp θεμιτὸν R : θέμις Burges : θεμιτ' Bothe : \lecon{θεμιτεύετον} Bibi \ligneApp{1153} λαμπάσι R : \lecon{λαμπάσι.} Bibi}%
\commentaire{1150-3}{οὗ δὴ ἀνδράσιν οὐ / θεμιτὸν εἰσορᾶν / ὄργια σεμνὰ θεαῖν / ἵνα λαμπάσι φαί-/}{Les mots †οὐ θεμιτὸν† ont été diversement retouchés, jamais de façon satisfaisante. Pour ma part, j'y verrais le verbe θεμιτεύω au duel.
La première partie du chant est une prière adressée à Athéna ; un scribe -- point bien réveillé -- n'a pas tout-de-suite vu que la seconde partie s'adressait aux deux déesses thesmophores.
Au premier vers de cette seconde partie (1148), il aura omis de mettre le verbe au pluriel : ἥκετ<έ τ᾽> εὔφρονες (simple \italique{haplographie}, comme on dit).
\italique{Bis repetita} : non encore averti par la longue énumération des duels aux vers 1154 et suivants, il persévère dans son erreur et le premier duel lui échappe tout aussi haplographiquement :%
\citComm{θεμι<τεύε>τον.}
Le verbe θεμιτεύω, \italique{permettre selon la loi ou le rite}, est utilisé par Euripide avec le même complément, ὄργια, au vers 79 des \italique{Bacchantes} :%
\citComm{τά τε ματρὸς μεγάλας ὄρ-\\
\alinea για Κυβέλας θεμιτεύων\\
Vénérant les orgies de la mère Cybèle la grande}\\
Cerise sur le gâteau : cette émendation permet de rétablir la correspondance métrique avec le vers 1158 : %
\citComm{οὗ δ\synizese{ὴ ἀ}νδράσιν\marqueur{1} οὐ θεμιτεύετον εἰσορᾶν\\
εἰ καὶ πρότερόν ποτ᾽ ἐπηκόω ἤλθετον\\
$ - \ - \ \cup \ \cup \ - \ \cup \ \cup \ - \ \cup \ \cup \ - \ \cup \ -$\marqueur{2}}%
}%
\debutNote{1150-53}%
\note{1}{Signalons la synalèphe au début du vers : δ\synizese{ὴ ἀ}νδράσιν -- je reprends le vocable de Koster dans son \italique{Traité de métrique grecque}. Il la distingue de la synizèse, terme qu'il applique lorsque la contraction entre deux voyelles a lieu au sein d'un même mot. Subtilité de grammairien ? pas que : la synizèse semble obéir à des lois différentes quant à la longueur des voyelles. Prato, lui, évoque à ce vers la présence d'une \italique{crase}.% \eff{italien ?}.
}%
\note{2}{Séquence interprétée par Lætitia comme un \italique{dactyle éolien}, c'est-à-dire dans son jargon une base (ici unissyllabique $ \times $, permise par l'éolisme du vers) sur laquelle viennent se fixer un nombre arbitraire de dactyles purs ($ - \cup \cup $) et finissant par un \italique{dodrans} ($ - \cup \cup - \cup - $), incontournable en contexte choriambique.}}
\end{groupe}%
%% v. 1136-1159
\begin{groupe}%
\typeVers{tetram. dact.}\chantChoeur{ὄργια σεμνὰ θεαῖν ἵνα λαμπάσι.%\apparat{1152-3}{\lecon{θεαῖν} R : θεοῖν Cobet \espApp λαμπάσι R : \lecon{λαμπάσι.} Bibi}%
\commentaire{1151}{ἵνα,}{Une glose à l'intérieur du texte du manuscrit\marqueur{} précise quel sens donner à (ἵνα : ) ὅπου\marqueur{1}.}%
\debutNote{1151}%
\note{1}{On retrouve une glose identique dans de nombreux lexiques : Συναγωγή ι 82 ; Hésychius ι 648 ; Photius ι 133 ; Suidas ι 360 ; \&c.}\\\passeLigne}
\end{groupe}%
%% v. 1136-1159
\begin{groupe}%
\Strophe{(μεσῳδός β)}\typeVers{trim. dact.}\chantChoeur{\ligne{1154}φαίνετον ἄμβροτον ὄψιν.\\\passeLigne}
\end{groupe}%
%% v. 1136-1159
\begin{groupe}%
\Strophe{(ἀντιστρ. β)}\typeVers{glyc.}\chantChoeur{μόλετον, ἔλθετον, ἀντόμεθ᾽ ὦ}
\end{groupe}%
%% v. 1136-1159
\begin{groupe}%
\typeVers{ibyc.}\chantChoeur{Θεσμοφόρω πολυποτνία,\apparat{1156}{πολὺ πότνια R : \lecon{πολυποτνία} Hermann}}
\end{groupe}%
%% v. 1136-1159
\begin{groupe}%
\typeVers{dact. aeol.}\chantChoeur{εἰ καὶ πρότερόν ποτ᾽ ἐπηκόω ἤλθετον}
\end{groupe}%
%% v. 1136-1159
\begin{groupe}%
\chantChoeur{\ligne{1159}νῦν ἀφίκεσθον †ἱκετεύομεν ἐνθάδ᾽ ἡμῖν†.\apparat{1159}{\lecon{ἐνθάδ᾽ ἡμῖν} R : ἐνθάδε χἠμῖν Hermann}\\\passeLigne}%
\end{groupe}%
%% v. 1160-1163
\begin{groupe}%
\ledpb%
\personnageGrec{Εὐριπίδης}%
\typeVers{trim. iamb.}\trimIamb{γυναῖκες εἰ βούλεσθε τὸν λοιπόν χρόνον\\
σπονδὰς ποιήσασθαι πρὸς ἐμέ, νυνὶ πάρα,\\
ἐφ᾽ ᾧτ᾽ ἀκοῦσαι μηδὲν ὑπ᾽ ἐμοῦ μηδαμὰ\\
κακὸν τὸ λοιπόν. ταῦτ᾽ ἐπικηρυκεύομαι.\\
\\\passeLigne}%
\end{groupe}%
%% v. 1164
\begin{groupe}%
\personnageGrec{Κορυφαῖος}%
\trimIamb{\ligne{1164}χρείᾳ δὲ ποίᾳ τόνδ᾽ ἐπεσφέρεις λόγον;\\\passeLigne}%
\end{groupe}%
%% v. 1165-1169
\begin{groupe}%
\personnageGrec{Εὐριπίδης}%
\trimIamb{ὅδ᾽ ἐστὶν οὑν τῇ σανίδι κηδεστὴς ἐμός.\apparat{1165}{οὖν R : \lecon{οὑν} Biset}\\
ἢν οὖν κομίσωμαι τοῦτον, οὐδὲν μή ποτε\apparat{1166}{\lecon{μή} R : μου Bentley}\\
κακῶς ἀκούσητ᾽· ἢν δὲ μὴ πίθησθέ μοι,\apparat{1167}{ἀκούσαιτ' R : \lecon{ἀκούσητ᾽} Elmsley (ad \italique{Ach.} v.295) \espApp πείθησθέ R : \lecon{πίθησθέ} Hirschig (ad \italique{Vesp} v.155)}\\
ἃ νῦν ὑποικουρεῖτε τοῖσιν ἀνδράσιν\\
ἀπὸ τῆς στρατιᾶς παροῦσιν ὑμῶν διαβαλῶ.\testimonia{1169}{(στρατιᾶς) Photius σ 608 στρατιάν}\commentaire{1168-9}{ἃ νῦν ὑποικουρεῖτε τοῖσιν ἀνδράσιν / ἀπὸ τῆς στρατιᾶς παροῦσιν.}{À propos de la légèreté des Athéniennes pendant le service de leur mari, cf. le vers 495 et le commentaire \italique{ad locum}.}\\\passeLigne}%
\end{groupe}%
%% v. 1170-1171
\begin{groupe}%
\personnageGrec{Κορυφαῖος}%
\trimIamb{τὰ μὲν παρ᾽ ἡμῶν ἴσθι σοι πεπεισμένα·\commentaire{1170}{τὰ μὲν παρ᾽ ἡμῶν ἴσθι σοι πεπεισμένα·}{Une scholie intérieure au texte nous précise :%
\citComm{Ἀττικῶς ἀντὶ τοῦ `πεπείσμεθα'.}}\\
τὸν βάρβαρον δὲ τοῦτον αὐτὸς πεῖθε σύ.\apparat{1171}{πεῖσαι R : \lecon{πεῖθε} Biset}\\\passeLigne}%
\end{groupe}%
%% v. 1172-1175
\begin{groupe}%
\ledpb%
\personnageGrec{Εὐριπίδης}%
\trimIamb{ἐμὸν ἔργον ἐστίν· καὶ σὸν ὦλάφιον ἅ σοι\apparat{1172}{Εὐριπίδης ἐν σχήματι προαγωγοῦ γραός Schol. \espApp \lecon{ἐμὸν} \inRas{γ᾽}{} \Rdeux \ligneApp{1173} \lecon{ἔφραζον} (ex ἔφρον) \Rpc \ligneApp{1174} κἀνακόλπασον R : \lecon{κἀνακάλπασον} Hermann, cf. Schol., Suid. κ 623, Zonar. p.1171 %: κᾆτα δ' ὄκλασον Moreau
\ligneApp{1175} Τερηδὼν R, Suid. τ 340 AGFM\expo{ac} : \lecon{Τερηδὸν} Suid. VM\expo{pc}}\\
καθ᾽ ὁδὸν ἔφραζον ταῦτα μεμνῆσθαι ποιεῖν.\\
πρῶτον μὲν οὖν δίελθε κἀνακάλπασον.\testimonia{1174}{(κἀνακάλπασον) Photius α 1514 ἀνακαλπάζει ; cf. Zonaras 1171 καλπάζειν· τὸ ἄκρως βαδίζειν $\simeq$ Scholium $=$ Suidas κ 623 καλπάζειν}%
\commentaire{1174}{πρῶτον μὲν οὖν δίελθε κἀνακόλπασον.}{(Que le lecteur curieux aille faire son miel de la note qu'a consacré Jacques Moreau à ce vers\marqueur{1}.)
Le \italique{Ravennas} porte au vers 1174 la leçon : πρῶτον μὲν οὖν δίελθε κἀνακ\textbf{ό}λπασον.
Le verbe ἀνακολπάζειν est un hapax, certes, mais dont le sens se laisserait entendre -- s'il n'y avait lieu comme nous l'allons voir de le soupçonner : ὁ κόλπος désigne au propre un \italique{pli}, \italique{creux} ; aussi bien le pli d'un vêtement -- notamment au niveau de la ceinture -- que, chez la mère ou la nourrice, le \italique{genou} ou le \italique{sein}. On trouve chez Homère l'expression κόλπον ἀνιεμένη\marqueur{2} lorsqu'Hécube fait tomber le haut de son vêtement et découvre son sein.
Il n'y aurait rien d'étonnant à ce qu'ἀνακολπάζειν ait pu signifier \italique{retrousser un vêtement}\marqueur{3}.
Pour définir les ζειραί, ces vêtements qu'on \italique{resserre au niveau de la ceinture}, le rédacteur anonyme de l'\latin{Etymologicum Magnum} use du verbe ἀνακολπόω :%
\citComm{ζειραί: χιτῶνες ἀνακεκολπωμένοι\marqueur{4}, ἢ ἀνάκωλοι.}%
Si le \italique{Ravennas} ne comportait en sus du texte au vers 1174 une scholie marginale venant renverser tout notre fragile édifice, ἀνακολπάζειν ferait partie de ces nombreux hapax dont le sens n'a été inféré qu'à grand renfort d'imagination sur de maigres indices étymologiques.\\
Cependant un ancien commentateur eut l'heureuse idée de gloser la leçon que son manuscrit donnait :%
\citComm{σ\douteux{η}[ * * * ] τοῦτο. ση[μαίνει] δὲ τὸ `ἁβρῶς βαδίζειν'.}%
Exit donc la leçon ἀνακολπάζειν, sans rapport avec le \italique{déplacement langoureux}\marqueur{5} dont parle le scholiaste, nullement étayée et d'ailleurs inepte : la danseuse court-vêtue n'a aucun vêtement pesant à retrousser ; au moment d'enflammer le Scythe, Euripide lui fera sauter les préliminaires en lui ordonnant à brûle-pourpoint de retirer tout bonnement son léger ἱμάτιον.\\
Hermann\marqueur{6} proposa de lire le verbe ἀνακαλπάζειν\marqueur{7}. 
Son hypothèse a été reprise depuis lors par toutes les éditions ; elle s'appuyait effectivement sur un indice de taille : à l'adresse κ 623 de la \italique{Souda}, l'on retrouve notre glose précédée du lemme καλπάζειν :%
\citComm{καλπάζειν· τὸ ἁβρῶς βαδίζειν.}%
%Ne chicanons pas l'émendation : transposition d'un `α' et d'un `ο', lettres parmi les plus dissemblables -- graphiquement comme phonétiquement -- et dont la confusion reste rare. Regardons plutôt ce que le sens y gagnerait :
Le verbe est formé à partir de ἡ κάλπη, \italique{le trot} (ou le \italique{galop} ?), qui désigne une sorte de course aux jeux olympiques, où le cavalier, en approchant du but, devait prendre le trot, sauter à terre et suivre son cheval.
%On retrouve le simple καλπάζειν chez Hésychius (κ 519) : καλπάζει· ὀξυποδεῖ ; ὀξυ- est à prendre au sens de ταχύ-, \italique{rapide} (cf. le passage de l'\italique{Oreste} d'Euripide, vers 1550, où Ménélas furieux accourt ὀξύπους $=$ ταχύπους). 
Quant au composé, ἀνακαλπάζειν, Photius (α 1514) nous en a conservé une trace en recopiant la glose d'un grammairien versé dans le \italique{Bon usage} attique :%
\citComm{ἀνακαλπάζει· τινὲς μὲν ὡς οὐ δόκιμον ἐφυλάξαντο τὴν φωνήν, Αἰσχύλος δὲ ἐχρήσατο Μυσοῖς \{ὡς δόκιμον\}· λέγει γάρ `εἶδον καλπάζοντας ἐν αἰχμαῖς'. ὁμοίως Σοφοκλῆς καὶ Ἀριστοφάνης καὶ Πλάτων καὶ ἕτεροι.\\
ἀνακαλπάζει· certains défendirent [l'usage de] de ce mot, comme non attesté. Eschyle pourtant l'utilisa dans ses \italique{Mysiens}. Il dit en effet : «~Je les vis trotter au milieu des lances~». De même Sophocle, Aristophane, Platon et d'autres.}%
C'est bien le sens de \italique{trotter}, \italique{galoper}, c'est-à-dire \italique{courir} qu'il faut donner au verbe <ἀνα>καλπάζειν, témoin le fragment d'Eschyle cité par Photius\marqueur{8} : il y fait allusion à une locution proverbiale\marqueur{9} :%
\citComm{`δι᾽ ὀξείας δραμεῖν'· ἐπὶ τῶν διακινδυνευόντων. ὀξεῖαν γὰρ λέγουσι τὴν λόγχην\\
«~Courir en traversant les pointes~» : à propos de ceux qui prennent des risques ; on appelle pointe la lance.}%
Dans le passage d'Eschyle, καλπάζοντας est un permutant métaphorique -- donc synonyme -- de δραμόντας.\\
Willems \& nombre d'exégètes voulurent rattacher <ἀνα>καλπάζειν à un terme technique de danse : \italique{se déplacer} ou \italique{danser sur les pointes}. Le sens leur a-t-il été suggéré par l'image du cheval au trot ? Dans le célèbre roman de Longus (II, 37), lorsque Daphnis imite le dieu Pan avec ses sabots, il se met à courir %
\citComm{ἐπ' ἄκρων τῶν ὀνύχων.}%
Là est le point : car l'on sait le type d'entraînement auquel se soumet la danseuse : au vers 1180, le Scythe qui la voit si \italique{légère} se mettre à \italique{sauter} («~comme une puce~», dixit le Scythe...), il s'exclame : ὡς ἐλαπρός, ὥσπερ ψύλλο κατὰ τὸ κῴδιο.
Moreau avait donc raison de souligner que la glose présentant notre danseuse aguichant l'Archer \italique{mollement} (ἁβρῶς βαδίζειν) ne tenait pas ; mais l'erreur venait de la glose et non de la correction d'Hermann : on eut le tort de vouloir conserver à tout prix cette glose telle que rapportée par le \italique{Ravennas} et Suidas ; elle n'était pourtant en conformité ni avec l'action ni avec le type d'entraînement -- hyperactif -- de la danseuse.
Le lexique de Jean Zonaras\marqueur{10} nous en livre la bonne version (simple substitution de ἄ\textbf{κ}ρως à ἁ\textbf{β}ρῶς) :%
\citComm{καλπάζειν· τὸ ἄκρως βαδίζειν.}%
Willems avait donc vu juste : l'exercice de la danseuse devait évoquer la course de la κάλπη : commencer par un \italique{trot}, c'est-à-dire par un déplacement en \italique{pointes} (τὸ ἄκρως βαδίζειν) qui s'accélérait pour terminer par un saut ; ce qu'Euripide désigne par le verbe ἀνακαλπάζειν.\\
Cette interprétation nous est confirmée par Eustathe (\italique{Commentaire. à l'}Iliade, XX, v. 228 p. 1206, à propos du mot `σκιρτῷεν') :%
\citComm{%τὸ δὲ σκιρτᾶν ἐοικὸς μὲν πρὸς τὸ σκαίρειν, καίριον δὲ ὂν εἰς ἀλλομένους ἵππους δὶς ἔφη, ἐπιμείνας ὡς ἐχρῆν. 
ἔστι δὲ σκιρτᾶν οὐ τὸ ἁπλῶς θέειν, ἀλλὰ τὸ ἀνέτως καὶ δίχα ἐπιταγῆς καὶ πηδητικῶς καὶ οἷον ὀρχηστικῶς \textbf{κατὰ τοὺς καλπάζοντας}.\\
%«~Bondir~» (σκιρτᾶν) est semblable à «~sauter~» (σκαίρειν), utilisé à bon escient par deux fois pour des chevaux bondissants, ἐπιμείνας ὡς ἐχρῆν (?) 
«~Bondir~» ((σκιρτᾶν)) n'est pas le simple fait de courir, mais de le faire librement, sans ordre, fougueusement, ainsi que le font les \textbf{καλπάζοντες} en dansant.}%
Terminons par Hésychius (κ 519)\marqueur{11}, qui ne le comprenait pas autrement :%
\citComm{καλπάζει· ὀξυποδεῖ.}%
ὀξυποδέω signifie \italique{accélérer le pas}\marqueur{12}. Notons aussi l'autre glose d'Hésychius (σ 822) :%
\citComm{σκαλ\{α\}πάζειν· ῥεμβωδῶς βαδίζειν.\\
σκαλπάζειν : se déplacer en tournoyant.}%
%Ce terme technique n'aurait de toute façon pas de rapport avec la définition, ἁβρῶς βαδίζειν\marqueur{6}, qu'en donne le scholiaste.\\
%Moreau propose de voir derrière `κἀνα' une déformation de κᾆτα ; dont la présence serait tout-à-fait naturelle après πρῶτον μέν dans le premier membre de la phrase\marqueur{7}.\\
%Si l'on admet cette émendation, la suite de la restitution paraît couler de source : remplacez le verbe κολπάζειν par μολπάζειν\marqueur{8}. ἡ μολπή est chez Homère un \italique{chant mêlé de danse}, le verbe se rencontre dans les \italique{Grenouilles} (vers 379 ; partie lyrique) ainsi que chez Athénée (XIII, 598f ; un poète qui chante son amour). Le verbe s'inscrirait parfaitement dans l'action et expliquerait l'ordre d'Euripide à la flûtiste de commencer l'accompagnement musical. Ainsi que le réveil en sursaut au vers suivant de l'archer qui, Scythe qu'il est, confond la douce mélopée avec un κῶμος dionysiaque. Toutefois, nous avons assez évoqué la scholie à ce vers pour ne pas reconnaître que μολπάζειν serait étrangement glosé de la sorte. L'explication réside peut-être dans la lacune de cette scholie (le début en est corrompu). Pour confirmer mon hypothèse, il faudrait qu'elle fût peu ou prou rédigée ainsi :%
%\citComm{ση[μαίνει] `ᾄδειν' τοῦτο. ση[μαίνει] δὲ τὸ `ἁβρῶς βαδίζειν'.}%
% Faute de grives on mange des merles : hésychius η δεϊ χελωνης χρέα φαγεϊν, ή μη φαγεϊν.
}%
\debutNote{1174}%
\note{1}{«~Note sur le vers 1174 des \italique{Thesmophories} d'Aristophane~», \italique{L'antiquité classique}, tome 7, fasc. 2 (1938) pp. 227-231. À peine âgé de vingt ans, le jeune étudiant d'alors fait montre d'une érudition et d'un esprit critique peu communs ; il faut dire qu'il avait pour professeur à l'Université de Bruxelles Henri Grégoire -- le successeur d'Alphonse Willems -- et dont l'assistant n'était autre que Roger Goossens ! bouillon intellectuel dont on imagine sans peine l'extraordinaire émulation qui devait en résulter parmi les élèves.}%
\note{2}{\italique{Iliade}, XX, 80. Et l'explication d'Hésychius (κ 3393) : κόλπον ἀνιεμένη· τὸ κατὰ τὸ στῆθος μέρος ἀνέλκουσα, ἵνα τὸν μαστὸν δείξῃ.}%
\note{3}{Ou, à la rigueur, «~tenir haut les seins~» comme le traduit Germaine Prudhommeau.}%
\note{4}{†ἀνακεκολλαμένοι† dans divers lexiques, notamment le lexique platonicien de Timée le Sophiste ou celui de Photius. La correction est-elle du fait du rédacteur de l'\latin{Etymologicum Magnum} ?}%
\note{5}{Ou \italique{lascif}, dira-t-on, au sens que le \nieme{XXI} siècle prête à ce mot ; c'est-à-dire avec \italique{langueur}, un \italique{manque de vigueur} qui trahit un fort penchant pour la sensualité. Est-ce l'influence de \italique{las} qui aura changé radicalement le sens originel de \italique{lascif} (\italique{lubrique}, \italique{folâtre} ; sens que rapportent encore exclusivement nos dictionnaires) ?}%
\note{6}{Dans sa recension des \italique{Thesmophories} éditées par Enger en 1844 (\italique{Aristophanis Comoediae cum scholiis}), in \italique{Zeitschrift für altertumswissenschaft} (1845), coll. 905-920. Article réédité par Fritzsche dans les \italique{Godofredi Hermanni Opuscula}, vol. 8 ; Hermann y discute de notre vers à la page 307.}%
%\note{3}{\eff{Retrouver dans le bouquin de danse attique la réf ; avec aussi la danse des épées.}}%
\note{7}{Au vers 1201, un `α' est lu fautivement de la même manière en `ο'.}%
\note{8}{$=$ fr. 144a Radt.}%
\note{9}{Transmis par les parémiographes Zénobius (III, 13), Suidas (δ 1186 ; ο 414) et Érasme (2014). Il n'est pas impossible que le proverbe soit lui-même tiré de la danse, et d'une éventuelle \italique{danse des épées}. Sur l'utilisation d'épées pour certaines acrobaties périlleuses, voir \italique{La danse grecque antique} de Maurice Emmanuel (1896), fig. 556 et 557 ou l'article «~Histoire anthropologique des danses acrobatiques~» de	Myriam Peignist, \italique{Corps} (2009/2), n° 7, pp. 29-38. La danseuse se livre-t-elle à ce genre d'acrobaties guerrières pour attirer l'attention du Scythe ? cela paraît peu vraisemblable ; il est préférable d'interpréter le fragment d'Eschyle de façon imagée et ne pas rattacher directement le verbe καλπάζειν à une quelconque danse des épées.
Un proverbe analogue et que l'on connaît en français, \italique{être sur le fil du rasoir}, apparaît déjà chez Homère (\italique{Iliade} X, 173).}%
%\note{9}{Et à sa suite \eff{citer les bouquins sur la danse grecque attique}}%
\note{10}{Édité par Tittmann sous le titre de \italique{Iohannis Zonarae Lexicon} (1808), p. 1171. Ce lexique de la fin du \nieme{XII} ou début du \nieme{XIII} siècle, s'il faut en croire Klaus Alpers, serait plutôt à attribuer à un certain Nicéphore. Le Pseudo-Zonaras a puisé abondamment dans la \italique{Souda} (ou des sources très voisines), notamment parmi les adresses de la lettre Α à Κ et les gloses touchant à Aristophane. Notre cas est intéressant, puisqu'il est peu vraisemblable que la transposition de β et κ soit une erreur ou correction du copiste ; il faut donc en conclure que si le Pseudo-Zonaras est allé piocher sa glose dans la \italique{Souda}, il disposait d'un meilleur manuscrit que ceux qui nous sont parvenus. Klaus Alpers a consacré un article très riche au lexique du Pseudo-Zonaras dans la monumentale \italique{Realencyclopädie der classischen Altertumswissenschaft}, 10/1 (1972), à l'entrée \allemand{'Zonarae' Lexikon}, coll. 732-763. Il le résume dans \italique{Der Neue Pauly}, 12/2 (2003), col. 831. Pour une étude plus générale de l'ensemble du corpus des anciens lexiques en langue grecque, l'on consultera la synthèse du même auteur dans l'\italique{Historisches Wörterbuch der Rhetorik}, 5 (2001) à l'entrée \allemand{Lexikographie} B, coll. 194-210.}%
\note{11}{Contrairement à l'usage de notre commentaire, nous n'avons réservé Hésychius que pour la fin de cette note. La raison en est que la glose de l'unicodex -- un \italique{Marcianus} du \nieme{XV} siècle -- est légèrement corrompue et donc sujette à caution :%
\citComm{καλπάγει ὀξυπόδης}%
(Le lexique de Cyrille ajoute : σακκιζει (sic, immo βαδίζει).)}%
\note{12}{Cf. le passage de l'\italique{Oreste} d'Euripide, vers 1550, où Ménélas furieux accourt ὀξύπους ($=$ ταχύπους). Ce sens est celui rapporté avec raison par le \LSJ. Le mot latin \latin{addensator}, \italique{celui qui rend plus serré} (en parlant du foulon et de son pas), est glosé dans certains manuscrits gréco-latin ὀξυποδητής.}%
%\note{6}{Et quoi qu'en pense Suidas (κ 623), qui aura eu la même idée qu'Hermann -- on a suffisamment vu en chemin à quoi s'en tenir avec certaines entrées de la \italique{Souda} :%
%\citComm{καλπάζειν· τὸ ἁβρῶς βαδίζειν.}%
%Mais digressons un tout petit peu, si le lecteur me le permet, et constatons amusés à quel point les modernes n'ont rien inventé : Jean Zonaras en son temps (\nieme{XII} siècle) était déjà gêné par la définition que Suidas donnait à καλπάζειν, incompatible avec le sens de "faire des pointes en danse" qu'il voulait lui plaquer. Il trouva une bonne astuce pour son propre lexique : substituer ἄκρως à ἁβρῶς -- solution si ingénieuse qu'on s'étonne que les modernes n'y aient pas pensé ! On ne sera donc pas surpris de lire à la page  : %
%\citComm{καλπάζειν· τὸ ἄκρως βαδίζειν.}}%
%\note{7}{Cf. par exemple l'arrivée tendue des ambassadeurs laconiens dans \italique{Lysistrata} (vers 1073-1074 ; avec la présence concomitante de πρῶτα μέν / εἶτα et de deux impératifs) :%
%\citComm{Ἄνδρες Λάκωνες, πρῶτα μέν μοι χαίρετε,\\
%εἶτ’ εἴπαθ’ ἡμῖν πῶς ἔχοντες ἥκετε.}}%
%\note{8}{L'on a déjà au vers 570 la forte ressemblance dans certains manuscrits du mu et du kappa.}
\\
σὺ δ᾽ ὦ Τερηδὸν ἐπαναφύσα Περσικόν.\testimonia{1175}{Suidas τ \espTest (Περσικόν) Pollux IV 100 ; cf. Hesychius π 2005 (?)}\\
\commentaire{1175}{Τερηδὼν,}{ἡ αὐλητρίς. (Note du scholiaste ; pourquoi pense-t-il que Térédon soit \italique{une} joueuse de flûte ?)}%
\commentaire{1175}{Περσικόν.}{Euripide demande au flûtiste d'accompagner la danseuse avec un air \italique{persan} -- les oreilles ? Or il existe une danse du même nom\marqueur{1} que Xénophon nous décrit dans son \italique{Anabase} (VI, 1, 10) : %
\citComm{μετὰ τοῦτο Μυσὸς\marqueur{2} εἰσῆλθεν ἐν ἑκατέρᾳ τῇ χειρὶ ἔχων πέλτην, καὶ τοτὲ μὲν ὡς δύο ἀντιταττομένων μιμούμενος ὠρχεῖτο, τοτὲ δὲ ὡς πρὸς ἕνα ἐχρῆτο ταῖς πέλταις, τοτὲ δ᾽ ἐδινεῖτο καὶ ἐξεκυβίστα ἔχων τὰς πέλτας, ὥστε ὄψιν καλὴν φαίνεσθαι. τέλος δὲ \textbf{τὸ περσικὸν ὠρχεῖτο} κρούων τὰς πέλτας καὶ ὤκλαζε καὶ ἐξανίστατο· καὶ ταῦτα πάντα ἐν ῥυθμῷ ἐποίει πρὸς τὸν αὐλόν.\\
Un Mysien\marqueur{2} entra ensuite ; il tenait, dans chacune de ses mains, un bouclier léger ; quelquefois il s'en servait en dansant, comme s'il eût eu à se défendre contre deux adversaires ; quelquefois, comme s'il n'eût eu affaire qu'à un seul. Souvent il tournait et faisait le saut périlleux sans lâcher ses boucliers : c'était un spectacle agréable qu'il donnait ; il finit par \textbf{danser la Persique} ; frappant d'un bouclier sur l'autre ; il s'accroupissait, se relevait, et exécutait tous ces mouvements en mesure et au son de la flûte.}%
L'accroupissement (ὀκλάζω) est le mouvement principal de cette danse et lui a donné son nom ; c'est le scholiaste aux \italique{Thesmophories} qui nous l'apprend : %
\citComm{Περσικὸν τὸ ὄκλασμα καλεῖται.}%
Pollux (IV, 100), compulsant un peu vite les \italique{Thesmophories} et ses commentateurs, croit que le terme ὄκλασμα apparaît dans le texte-même d'Aristophane : il range l'ὄκλασμα auprès du κνισμός parmi les danses guerrières et frénétique :%
\citComm{καὶ ὄκλασμα· οὕτω γὰρ ἐν Θεσμοφοριαζούσαις ὀνομάζεται τὸ ὄρχημα τὸ Περσικὸν καὶ σύντονον. τὴνδ' αὐτὴν καὶ ὑγρὰν ὠνόμαζον.}%
L'on trouve d'ailleurs l'une des scholies des \italique{Thesmophories} les plus savantes à notre vers 1175\marqueur{3} :%
\citComm{Περσικὸν ὄρχημα, περὶ οὗ Ἰόβας μακρὸν πεποίηται λόγον ἐν τοῖς περὶ τῆς θεατρικῆς ἱστορίας. ὥστε λελύσθαι τὴν Σελεύκοὺ πρότασιν. προτείνει γὰρ ἐν τῷ πρὸς Ζήνωνα προτατικῷ τὸ Περσικὸν ὄρχημα.}%
Pour ma part, je doute que l'exercice auquel se livre notre danseuse Ἐλάφιον soit cette danse guerrière illustrée par le Mysien de Xénophon ; l'adjectif Περσικός était assez commun pour qu'il pût désigner à la fois un air de flûte et une forme de danse sans qu'il y eut nécessairement de rapport entre les deux.}%
\debutNote{1175}%
\note{1}{Hésychius (π 2005) : περσικά· [ ... ] καὶ ὄρχησις Περσική.}%
\note{2}{Je cite longuement le passage de Xénophon pour faire apparaître la race du danseur. Dans le cas où la danse persique aurait effectivement un rapport avec le verbe ἀνακαλπάζειν du vers 1174 -- ce que je ne crois pas --, on notera qu'il faisait partie du peuple d'Asie qui servit de titre à la pièce d'Eschyle, citée par Photius à l'adresse ἀνακαλπάζει et dont on a parlé au vers précédent.}%
\note{3}{Elle est discutée par Jacques Moreau dans l'article déjà mentionné : «~Note sur le vers 1174 des \italique{Thesmophories} d'Aristophane~», \italique{L'antiquité classique}, tome 7, fasc. 2 (1938) pp. 227-231.}}%
\end{groupe}%
%% v. 1176
\begin{groupe}%
\ledpb%
\personnageGrec{Τοξότης}%
\trimIamb{\ligne{1176}τί τὸ βόμβο τοῦτο; κῶμό τις ἀνεγεῖρί μοι;\apparat{1176}{κωμοτίς ἀν- R : κωμοδρίαν Schol. : \lecon{κῶμό τις ἀν-} van Velsen}%
\commentaire{1176}{κωμοτίς ἀνεγεῖρί μοι;}{Variante de lecture du scholiaste : διχῶς κωμοδρίαν.}
%\commentaire{1176}{}{Le bourdonnement Hésychius : κόμμωσις:ὁ τοῦ σώματος ὥὡραϊσμός. καὶ ὑπὸ τῶν μελισσουργῶν ἡ τοῦ σμήνους διάχρισις}%
\\\passeLigne}%
\end{groupe}%
%% v. 1177-1178
\begin{groupe}%
\personnageGrec{Εὐριπίδης}%
\trimIamb{\ligne{1177}ἡ παῖς ἔμελλε προμελετᾶν ὦ τοξότα.\\
ὀρχησομένη γὰρ ἔρχεθ᾽ ὡς ἄνδρας τινάς.\apparat{1178}{ἔρχετ᾽ R : \lecon{ἔρχεθ᾽} Brunck}\\\passeLigne}%
\end{groupe}%
%% v. 1179-1180
\begin{groupe}%
\personnageGrec{Τοξότης}%
\trimIamb{ὀρκῆσι καὶ μελετῆσι, οὐ κωλῦσ᾽ ἐγώ.\apparat{1179}{κωλύσ᾽ R : \lecon{κωλῦσ᾽} Blaydes (κωλῦσι $=$ κωλύσω), cf. ad 1003, 1196, 1201, 1225 \ligneApp{1180} ὡς et ὥσπερ R : \lecon{ὠς} et \lecon{ὤσπερ} Brunck}\\
ὠς ἐλαπρός, ὤσπερ ψύλλο κατὰ τὸ κῴδιο.%
\commentaire{1172 et 1180}{ὦλάφιον / ὡς ἐλαπρός, ὥσπερ ψύλλο κατὰ τὸ κῴδιο.}{Les verbes signifiant \italique{sauter}, \italique{danser} (πηδᾶν, ἄλλεσθαι) attirent presque inévitablement une comparaison, et le plus souvent avec un animal.
Les jeunes filles de Sparte dansent en bondissant comme des pouliches, ἇτε πῶλοι ταὶ κόραι / πὰρ τὸν Εὐρώταν / ἀμπαδίοντι (ἀμπαδίοντι Coulon : ἀμπάλλοντι codd.) (\italique{Lysistr.} 1308)\marqueur{1}. 
La Muse laconienne danse et saute comme une biche, ποδοῖν τε πάδη / ἇ τις ἔλαφος (\italique{ib}. 1317)\marqueur{2}.
La danseuse des \italique{Thesmophories} porte d'ailleurs le sobriquet de \italique{Bichette} (Ἐλάφιον, v. 1172).
Bichette danse avec tant de grâce devant l'archer Scythe qu'il s'écrie : « qu'elle est légère ! une vraie puce sur ma peau de mouton » (\italique{Thesm}. 1180). Cette comparaison dénuée de poésie sent le corps de garde et ses peaux de mouton où grouille la vermine (on dormait communément sur des \italique{peaux de mouton}, κῴδια). (Note de Taillardat, §791.)}%
\debutNote{1172 et 1180}%
\note{1}{Comparaison banale, cf. Eur. \italique{Or}. 45 (à propos d'Oreste furieux): πηδᾷ δρομαῖος, πῶλος ὣς ὑπὸ ζυγοῦ ; Eubul. 75, 6 [$=$ Athénée, III, 108ab]. Le πωλίον symbolise la vivacité, Ar. \italique{Guêp}. 188. En outre, πῶλος désigne couramment les jeunes femmes : Anacréon 88 \italique{Diehl} [$=$ 417 \italique{Poetae Melici Graeci} Page]; Eur. \italique{Héc.} 142, \italique{Hipp}. 516, etc. ; Cratin. 87 [$=$ 94 Kassel-Austin] ; Épicrat. 9, 4; Eubul. 84 [$=$ 82 Kassel-Austin].}%
\note{2}{Le cerf (ἔλαφος) est le symbole de l'agilité, \italique{Nuées} 351. (Et aussi, comme chez Homère, de la lâcheté.)}\\\passeLigne}%
\end{groupe}%
%% v. 1181-1183a
\begin{groupe}%
\personnageGrec{Εὐριπίδης}%
\trimIamb{φέρε θοἰμάτιον ἄνωθεν ὦ τέκνον τοδί·%
\commentaire{1181}{φέρε θοἰμάτιον ἄνωθεν.}{Sur le texte et le sens de ce vers 1181, voir la note 28 d'Alphonse Willems, reproduite ci-après.}\\
καθιζομένη δ᾽ ἐπὶ τοῖσι γόνασι τοῦ Σκύθου\apparat{1182}{τοῖς R : \lecon{τοῖσι} Bentley}\\
τὼ πόδε πρότεινον, ἵν᾽ ὑπολύσω.\testimonia{1183}{(τὼ πόδε) Hesychius τ 1758 τὼ πόδε· τοὺς πόδας δυϊκῶς}\\\passeLigne}%
\end{groupe}%
%% v. 1183b-1185
\begin{groupe}%
\personnageGrec{Τοξότης}%
\trimIamb{\ligne{1183}\decalage{τὼ πόδε πρότεινον, ἵν᾽ ὑπολύσω.}ναῖκι ναὶ\\
\lednopb%
κάτησο κάτησο, ναίκι ναί, τυγάτριον.\apparat{1184}{\lecon{ναίκι ναὶ} R : ναίκι ναίκι Bothe}\\
\lednopb%
οἴμ᾽ ὠς στέριπο τὸ τιττί᾽, ὤσπερ γογγυλί.%
\apparat{1185}{ὡ στεριπὸ (ex σπεριπο) \Rpc : \lecon{ὡς στεριπὸ} (sic) \Mon \espApp ὥσπερ R : \lecon{ὤσπερ} Brunck \espApp \lecon{γογγυλί} \PsiTrois, Enger : γογγύλῃ R}%
\commentaire{1185}{ὡς στεριπὸ τὸ τιττί᾽}{$=$ ὡς στέριφον τὸ τιτθίον.}%
\commentaire{1185-93}{Papyrus d'Oxyrhynque n° 3840.}{Un mince débris contenant quelques lettres de notre passage fut déniché parmi les papyrus d'Oxyrhynque\marqueur{1}. Il y a néanmoins quelques informations à en tirer : il confirme la conjecture d'Enger au vers 1185 (γογγυλί à la place de γογγύλῃ), ainsi que μέλι au vers 1192. Parsons note la présence d'une ligne supplémentaire après le vers 1187, là où le \italique{Ravennas} possède l'une de ses rares didascalies\marqueur{2} : ἀνακύπτη· καὶ παρακύπτι· ἀπεψωλημένος· (même jeu de scène avec l'olisbos qu'aux vers 643-648).}%
\debutNote{1185-93}%
\note{1}{\italique{The Oxyrhynchus Papyri}, vol. LVI (1989), n°3840. Étudié par Peter J. Parsons, aidé de Colin Austin -- déjà en train de préparer son édition des \italique{Thesmophories}. La transcription qu'en fit Parsons et son analyse sont recopiées en annexe. Voir aussi le papyrus d'Oxyrhynque n° 5132 (\italique{The Oxyrhynchus Papyri} LXXVIII) pour un semblant d'\italique{addendum} : la présence de doubles points `$ : $', comme au vers 1190, pour indiquer un changement de locuteur est confirmée par un autre fragment du même papyrus au vers 1209.}%
\note{2}{Pour notre pièce où les scholies très anciennes et érudites se font rares, on ne relève d'autres didascalies qu'aux vers 100, 129 et 276.}\\\passeLigne}%
\end{groupe}%
%% v. 1186
\begin{groupe}%
\ledpb%
\personnageGrec{Εὐριπίδης}%
\trimIamb{αὔλει σὺ θᾶττον· ἔτι δέδοικας τὸν Σκύθην;\\\passeLigne}%
\end{groupe}%
%% v. 1187-1189
\begin{groupe}%
\personnageGrec{Τοξότης}%
\trimIamb{καλό γε τὸ πυγή. κλαῦσί γ᾽ ἢν μὴ 'νδον μένῃς.\testimonia{1187}{(πυγή) Hesychius π 4279 πυγή· τὸ κάθισμα}\apparat{1187}{κλαυσ\inRas{᾽ ε}{}ϊ γ᾽ \Rdeux : κλαυσεῖ γ᾽ \Mon \espApp ἂν R : \lecon{ἢν} Dindorf (1839) \ligneApp{post 1187} μ]\douteux{ε}νο\douteux{ }[ \PsiTrois : ἀνακύπτη· καὶ παρακύπτι· ἀπεψωλημένος· R : \lecon{del.} Ellebodius \ligneApp{1188} τὸ \inRas{σ}{}ποστίον \Rpc : τὸ ποστίον \Mon \ \ligneApp{1189}  'στιν R : \lecon{'στὶ} \Mon}%
\commentaire{1187}{κλαῦσί γ᾽ ἂν μὴ 'νδον μένῃς.}{πρὸς τὸ πέος λέγει ἐπιπλήττων ἐντεινομένῳ. (Note du scholiaste.)\\
L'archer suit d'un œil émérillonné tous les mouvements de la danseuse. Son agitation se trahit d'une manière par trop visible. Le Scythe, qui s'en aperçoit, reprend « gare à toi, etc .. »}\\
\ligne{1188}εἶεν· καλὴ τὸ σκῆμα περὶ τὸ πόστιον.%
%\commentaire{1188}{}{\eff{Cf. fr. ... περι το σπηλαιον / σπειον ... et vers ... περι το τραμιν}}
\\\passeLigne}%
\end{groupe}%
%% v. 1189-1190a
\begin{groupe}%
\personnageGrec{Εὐριπίδης}%
\trimIamb{καλῶς ἔχει. λαβὲ θοἰμάτιον· ὥρα 'στὶ νῷν\\%\apparat{1189}{'στιν R : \lecon{'στὶ} \Mon}\\%
ἤδη βαδίζειν.\\\passeLigne}%
\end{groupe}%
%% v. 1190b
\begin{groupe}%
\personnageGrec{Τοξότης}%
\trimIamb{\ligne{1190}\decalage{ἤδη βαδίζειν.}οὐκὶ πιλῆσι πρῶτά με;\apparat{1190}{τί οὐκ ἐπιλήσει R : \lecon{οὐκὶ} (Biset) \lecon{πιλῆσι} Porson (ad \italique{Aristophanica} p.222), cf. Schol. `οὐχὶ φιλήσειν'}\\\passeLigne}%
\end{groupe}%
%% v. 1191a
\begin{groupe}%
\personnageGrec{Εὐριπίδης}%
\trimIamb{πάνυ γε· φίλησον αὐτόν.\\\passeLigne}%
\end{groupe}%
%% v. 1191b-1193a
\begin{groupe}%
\personnageGrec{Τοξότης}%
\trimIamb{\ligne{1191}\decalage{πάνυ γε· φίλησον αὐτόν.}ὂ ὂ ὂ παπαπαπαῖ,\apparat{1191}{ὂ R : \lecon{ὂ ὂ ὂ} \Mon}\\
ὠς γλυκερὸ τὸ γλῶσσ᾽, ὤσπερ Ἀττικὸς μέλι.\apparat{1192}{ὡς et ὥσπερ R : \lecon{ὠς} et \lecon{ὤσπερ} Brunck \espApp \lecon{μέλι} \PsiTrois, \Mon : μέλις R}%
\commentaire{1192}{ὠς γλυκερὸ τὸ γλῶσσ᾽, ὤσπερ Ἀττικὸς μέλις.}{Si l'image des baisers plus doux que le miel\marqueur{1} ou le nectar nous paraît banale, c'est à une longue tradition héritée de ces culs-terreux de latins que nous le devons\marqueur{2} ; pour ce qui est des auteurs de langue grecque, on ne peut guère citer que Alciphron (I, 38, 7), postérieur d'environ six cents ans à Aristophane. 
Les spectateurs d'Aristophane devaient ressentir avec cette comparaison une confusion burlesque : l'épithète poétique μελίγλωσσος existe en grec, mais l'image d'une \italique{langue pleine de miel} signifie tout autre chose. 
Ainsi parlait Taillardat (§§739-743) :\\
«~Le miel étant le symbole de la douceur\marqueur{3}, les poètes se plaisent à dire que leurs chants sont de miel\marqueur{4}. Cette image est chez Aristophane : μελιγλώσσων ἐπέων ... ἀοιδάν (\italique{Ois}. 908) donne un exemple de plus de l'épithète banale μελίγλωσσος\marqueur{5} [ ... ]
On dira donc tout naturellement que la lèvre d'un poète est enduite de miel : c'est le cas de la bouche de Sophocle, Σοφοκλέους τοῦ μέλιτι κεχριµένου... τὸ στόμα (\italique{fr}. Aristophane 581 [$=$ 598 Kassel-Austin]). Composant des poèmes doux comme le miel, les poètes aiment à se comparer aux abeilles, ainsi qu'en témoigne Platon (\italique{Ion}, 534b).~»\\
Aristophane fait rire son public aux dépens de ces images trop usées : le Scythe rustique vante la \italique{langue douce comme le miel} de la danseuse, mais c'est une langue toute matérielle -- et qui appartient d'ailleurs à une actrice muette. Il y a des impropriétés de termes qui sont des malpropretés du cœur.}%
\debutNote{1192}%
\note{1}{Et l'Attique est souvent citée pour son miel. Cf. la \italique{Paix}, vers 252-254 ; Antiphane fr. 177(179) Kassel-Austin $=$ Athénée II, 43b \& III, 74d ; le corpus hippocratique mentionne maintes fois du miel Attique.}%
\note{2}{Cf. le \italique{Dictionnaire des sentences latines et grecques} de Renzo Tosi, §1385 \latin{Nectar dulcius} : Catulle 99, 1-2 ; Horace, \italique{Carmina} 1, 13, 14-16 ; \&c. Honoré de Balzac a fameusement repris cette longue tradition : « Les baisers d'une femme sincère ont un miel divin qui semble mettre dans cette caresse une âme, un feu subtil par lequel le cœur est pénétré. Les baisers dénués de cette onction savoureuse sont âpres et secs. »}%
\note{3}{\italique{Guêpes} 878 (voir §361), \italique{Thesm.} 1192, fr. 587, 622[$=$ 681, 633 K.-A.]; Alexis 146, 6 [$=$ 150.6 K.-A.]; Ménand. 708 (?), etc.}%
\note{4}{Pind. \italique{Ol.} 11, 4 : μελιγάρυες ὕμνοι ; id. \italique{Ném.} 11,18 : μελιγδούποισι ... μελιζέμεν ἀοιδαῖς ; id. \italique{Isthm.} 2, 32 : μελικόμπων ἀοιδᾶν ; Alcman 41 D [$=$ Athénée XIV, 633a $=$ fr. adesp. 954b \italique{Poetae Melici Graeci}, ed. Page] : μέλεα μελιπτέρωτα ; Simonid. 40 D [$=$ fr 595.3 \italique{Poetae Melici Graeci} Page] : μελιαδέα γᾶρυν, etc. On notera que, déjà chez Homère (\italique{Il.} I, 249), la voix de Nestor est dite [ἀπὸ γλώσσης] μέλιτος γλυκίων. D'autre part Pindare, \italique{Ol.} 7, 8, appelle encore la poésie un nectar, νέκταρ.}%
\note{5}{Eschyl. \italique{Prom.} 172 ; Bacchyl. 3, 97 ; fr. 4, 25 Snell (μελιγλώσσων ἀοιδᾶν ἄνθεα).}\\
τί οὐ κατεύδει παρ᾽ ἐμέ;\\\passeLigne}%
\end{groupe}%
%% v. 1193b-1194a
\begin{groupe}%
\personnageGrec{Εὐριπίδης}%
\trimIamb{\ligne{1193}\decalage{τί οὐ κατεύδει παρ᾽ ἐμέ;}χαῖρε τοξότα,\\
\lednopb%
οὐ γὰρ γένοιτ᾽ ἂν τοῦτο.\\\passeLigne}%
\end{groupe}%
%% v. 1194b-1195a
\begin{groupe}%
\personnageGrec{Τοξότης}%
\trimIamb{\ligne{1194}\decalage{οὐ γὰρ γένοιτ᾽ ἂν τοῦτο.}ναὶ ναὶ γρᾴδιο.\apparat{1194}{ναὶ R : ναίκι Bentley : \lecon{ναὶ ναὶ} Brunck \espApp γράδιον R : \lecon{γρᾴδιο} Bothe \ligneApp{1195} \inRas{σὺ}{} κάρισοσ οὐ \Rpc : κάρισο σοῦ \Mon : \lecon{κάρισο σὺ} Brunck}\\
ἐμοὶ κάρισο σὺ τοῦτο.%
\commentaire{1195}{ἐμοὶ κάρισο σὺ τοῦτο}{$=$ ἐμοὶ χάρισαι σὺ τοῦτο.}\\\passeLigne}%
\end{groupe}%
%% v. 1195b
\begin{groupe}%
\personnageGrec{Εὐριπίδης}%
\trimIamb{\ligne{1195}\decalage{ἐμοὶ κάρισο σὺ τοῦτο.}δώσεις οὖν δραχμήν;%
\\\passeLigne}%
\end{groupe}%
%% v. 1196a
\begin{groupe}%
\personnageGrec{Τοξότης}%
\trimIamb{ναὶ ναῖκι δῶσι.\\\passeLigne}%
\end{groupe}%
%% v. 1196b
\begin{groupe}%
\personnageGrec{Εὐριπίδης}%
\trimIamb{\ligne{1196}\decalage{ναὶ ναῖκι δῶσι.}τἀργύριον τοίνυν φέρε.\\\passeLigne}%
\end{groupe}%
%% v. 1197
\begin{groupe}%
\personnageGrec{Τοξότης}%
\trimIamb{ἀλλ᾽ οὐκ ἔκὠδέν· ἀλλὰ τὸ συμβήνην λαβέ.\testimonia{1197}{Photius σ 660 συβήνη $=$ Suidas σ 1273 (sed ἐχω οὐδέν pro έκ' ὠδέν, τὸ omisso et συβήνην pro συμβήνην) \espTest (συμβήνην) Hesychius σ 2137 συμβίνης· καπροβόλον. ἐμβόλιον}\apparat{1197}{έκ' ὠδέν R : ἐχω οὐδέν Suid. σ 1273, Photius σ 660 : \lecon{ἔκὠδέν} Bergk \espApp \lecon{συμβήνην} R : συβήνην Schol., Suid., Photius \ligneApp{1198} κ\inRas{ω}{ο}μίζεις αὐτοῖς \Rpc : κομιεῖς αὐτό (αὖθις Brunck) Biset : (Τοξότ.) κομίσι σ' αὐτός (αὖτις anonym. Parisinus) Enger : (Τοξότ.) \lecon{κομίζι σ᾽ αὐτίκ᾽} Bibi \espApp ἀκουλούτι R : \lecon{ἀκολούτι} \Mon : ἀκουλούτει Brunck}%
\commentaire{1197-216}{ἀλλὰ τὸ συμβήνην λαβέ [...] ὀρτῶς δὲ συβήνη 'στί· καταβηνῆσι γάρ.}{Ce passage a été si maltraité par les commentateurs -- anciens comme modernes -- que nous nous sommes cru obligé à commettre cette note rectificatrice. \italique{Recte}, descendons-donc dans l'arène (ὀρτῶς δὲ συμβήνη 'στί· καταβηνῆσι γάρ.)\\
L'origine de tous les malentendus est la glose d'un commentateur ancien qui sera allé un peu vite : le \italique{Ravennas} comporte deux scholies pour définir la συμβήνην (la première scholie se trouve à l'intérieur même du texte) :%
\citComm{τὴν τοξοθήκην\\
συβήνη· αὐλοθήκη. λέγουσι δὲ καὶ τὸν φαρετρεῶνα συβήνην.}%
On retrouve cette scholie chez Hésychius (σ 2136) :%
\citComm{συβήνη· αὐλοθήκη, ἢ τοξοθήκη. ἢ ὁ ναυτικὸς χιτών.}%
Suidas cite notre passage dans une glose qu'on croise à l'identique, ou peu s'en faut, dans d'autres lexiques\marqueur{1}:%
\citComm{Συβήνη: ἡ δερματίνη αὐλοθήκη\marqueur{2}. ἢ ἡ φαρέτρα. Θεσμοφοριαζούσαις Ἀριστοφάνης· ἀλλ' οὐκ ἔχω οὐδέν. ἀλλὰ συβήνην λάβε.}%
Cette unanimité ne doit pas nous tromper : elle prouve seulement l'ancienneté de la première rédaction qui aura diversement infusé dans de nombreux lexiques. Je dirais même plus : ces gloses concordent si bien qu'il est certain qu'elles remontent toutes à une même source -- certes ancienne -- mais unique, et donc faillible.\\
Que le terme συβήνη -- qui n'est d'ailleurs pas celui transmis par notre manuscrit comportant συ\textbf{μ}βήνην --, que συβήνη ait pu désigner un \italique{étui à flûte}, je n'en disconviens pas. À la rigueur, que par dérivation le mot ait aussi pu être appliqué à un \italique{carquois}, pourquoi pas\marqueur{3} ? Mais il serait absurde de lui donner ce sens dans notre passage : croit-on que l'Archer, tout Scythe qu'il fût, pensait un seul instant qu'une vieille maquerelle pût accepter son carquois dont elle n'aurait que faire au lieu d'une drachme sonnante et trébuchante ? ce serait tout bonnement ridicule\marqueur{4} ! Le Scythe, Euripide, Aristophane, le savent bien ; en fait n'importe quel homme, pourvu qu'il ne fût pas versé dans l'exégèse d'Aristophane : %
\citComm{Chez la belle Suzon, pas d'argent, pas de cuisse...\marqueur{5}}%
En réalité, le Scythe demande à Euripide de garder son prisonnier quelques instants, le temps qu'il aille chercher l'argent nécessaire.
Encore une fois, c'est Hésychius (σ 2137) qui détenait la solution :%
\citComm{συμβίνης· καπροβόλον. ἐμβόλιον.}%
Le mot est apparenté à σιβύνη, une \italique{lance} (Hésychius nous précise pour la chasse au \italique{sanglier}, κάπρος), dont il est une légère déformation qui prépare le jeu de mots scabreux du vers 1215.
Le texte est aussi clair qu'un ciel d'azur au milieu de l'été : le Scythe confie sa lance à Euripide pour la garde de Mnésiloque.
Et nous nous garderons bien de biffer le μ de la leçon du manuscrit, συ\textbf{μ}βήνην, qui aura survécu à plus de mille ans de tradition manuscrite mais pas à la griffe implacable des éditeurs d'Aristophane.\\
Le début du vers suivant, altéré, ἔπειτα †κ\inRas{ω}{ο}μίζεις αὐτοῖς†, se laisse lire sans difficulté :%
\citComm{ἔπειτα κομίζι ($=$ κομίζω) σ(οὶ) αὐτίκ(α)\\
Ensuite je te l'apporte [la drachme] sur-le-champ.}%
La dernière version du manuscrit attribue, il est vrai, cette réplique à Euripide. Ces indications marginales de personnages ne doivent pas nous arrêter : elles sont récentes et même pour certaines du \nieme{XVI} siècle\marqueur{6} (c'est le cas de cette annotation, «~γϋ.~», que porte le manuscrit (la même qu'au vers 1189) et mis pour γυνή -- c'est-à-dire la \italique{vieille}, γραῦς ; c'est-à-dire Euripide. Von Velsen ne l'a pas intégrée à son apparat critique et avait donc reconnu l'écriture d'Euphrosynus Boninus).\\
Notre émendation appelle quelque justification : l'hypothétique confusion αυτικ / αυτοις peut aisément s'expliquer en constatant la ressemblance entre le digramme `ις' et le `κ'\marqueur{7}. Continuons.\\
Le Scythe part chercher l'argent, accompagné de la jeune beauté, qui se montre si caressante selon les instructions d'Euripide (vers 1172 ; l'attitude de la danseuse est la version muette de celle de la Myrrhine dans \italique{Lysistrata}) qu'il en revient tout émoustillé et en proie à une promesse de félicité qui point\marqueur{8}. Il déchante très vite : Euripide et Mnésiloque ont pris la poudre d'escampette. Élaphion décampe à ce moment sous les imprécations du Scythe (ἀπότρεκ᾽ ὠς τάκιστα σύ).
Le Scythe s'arme donc de sa lance, laissée sur place, et s'apprête à descendre (καταβαίνειν) à leur poursuite :%
\citComm{ὀρτῶς δὲ συμβήνη ᾽στί· καταβηνῆσι γάρ.}%
Je n'explique pas le trait d'esprit -- d'esprit \italique{Scythe} ; si vous ne le comprenez pas, laissez-là Aristophane et allez faire des mathématiques !\\
Nul besoin de retailler le vers : καταβαίνω est le terme consacré lorsqu'un acteur descend de la σκηνή vers l'ὀρχήστρα, d'autant plus adapté quand la scène se situe sur un lieu en hauteur comme la colline de la Pnyx\marqueur{9} ; cf. les \italique{Guêpes} (vers 1514) ou l'\italique{Assemblée des femmes} (vers 1152)\marqueur{10}.
%Ces savants hommes ne manquent jamais de mettre leur érudition en branle dès qu'ils subodorent derrière la plus anodine des répliques le relent du pornographe Aristophane... et passent sans s'arrêter quand le plus grossier de ses personnages paye une maquerelle d'un écot -- à défaut d'un écu -- au nom sonore : écho de *συμβινεῖν. (\italique{Confer} le vers 1215.)
}%
\debutNote{1197-216}%
\note{1}{Photius (σ 660), \latin{Etymologicum magnum}, \latin{genuinum}, Zonaras, \&c. Pollux (VII, 153 ; X, 153) : καὶ τὸ μὲν ἀγγεῖον τὸ τὰς γλώττας ἔχον γλωττοκομεῖον, ἣ δὲ τῶν αὐλῶν θήκη συβήνη. Pour un dossier plus complet, je renvoie le lecteur à la \allemand{Realencyclopädie} : \italique{Paulys Realencyclopädie der classischen Altertumswissenschaft}, IV 1 (1931), coll. 1011-1012.}%
\note{2}{Dans certains manuscrits du \italique{Zonarae Lexicon} (p. 1686), on trouve βελοθήκη en place d'αὐλοθήκη.}%
\note{3}{Encore qu'il me semble plus vraisemblable que ce sens ait été inféré par notre scholiaste, puis scrupuleusement intégré à tous les lexiques.}%
\note{4}{À tous les vulgarisateurs d'Aristophane, tournebouleurs et avilisseurs de ses onze chefs-d'œuvre, rappelons la maxime de bon goût qu'Aristophane s'imposait :%
\citComm{ὡς ἐγὼ φοβοῦμαι περὶ τῶν μελλόντων ῥηθήσεσθαι, οὔ τι μὴ γελοῖα εἴπω — τοῦτο μὲν γὰρ ἂν κέρδος εἴη καὶ τῆς ἡμετέρας μούσης ἐπιχώριον — ἀλλὰ μὴ καταγέλαστα.\\
Ce n'est pas de provoquer le rire que je redoute -- cela serait tout gain, et conforme à ma muse -- mais de dire des choses ridicules.\\
(Trad. P. Boutang.)}%
(Tradition reçue en droite ligne d'Aristodème, témoin des événements, qui en fit le récit à Apollodore, qui lui-même le transmit à Glaucon, frère de Platon, qui -- on peut l'imaginer ainsi -- confia bientôt, après peut-être encore quelques méandres, à son frère Platon le précieux compte-rendu, lequel Platon finit par coucher par écrit cette fameuse nuit de \italique{Banquet}, 189b, quelques quarante années après les faits.)}%
\note{5}{Le proverbe d'origine détourné par Brassens (\italique{Grand-père}) étant «~point d'argent, point de Suisse~». Quel sera l'homme d'assez bonne volonté pour nous offrir le colossal équivalent brassinien du maître-œuvre de Taillardat, quel est le preux qui prendra à bras-le-corps les quelques deux cents chansons de Brassens que ce dernier a bien voulu nous conserver ? Espérons qu'il ne soit pas besoin d'attendre encore vingt-cinq siècles comme pour Aristophane. (C'est quand le vers est troussé qu'il faut l'analyser.) Je tiens à la disposition de cette bonne âme éventuelle un matériau qui poserait de solides fondations à ces \italique{Images de Brassens} en fomentation.}%
%\note{6}{\eff{Retrouver la réf sur l'armement des Scythes ; Hall ?}}%
\note{6}{Cf. \italique{supra}, v. 1059 n. 5.}%
\note{7}{Il faudrait alors supposer l'existence d'une forme intermédiaire αυτιις ou αυτις, bien vite corrigée en αυτοις. Une forme intermédiaire, ou un \italique{chaînon manquant}, pour reprendre le jargon des darwiniens. Ne poussons donc pas plus loin la «~preuve~» afin d'éviter l'écueil de ces demi-savants, capables de tout prouver et même son contraire !}%
\note{8}{Savez-vous ce qu'on peut lire dans le commentaire le plus universitairement reçu, le plus académiquement académisant -- celui d'Austin \& Olson, puisqu'il faut le nommer ? Je vous le donne en mille : le Scythe a consommé la belle Élaphion... il revient donc, précisent leurs auteurs, avec l'\italique{olisbos} abaissé (!!!) Quand on sait -- ou qu'on ne sait pas, comme Austin \& Olson -- combien Aristophane adorait ses chères actrices, qu'il ne manquait jamais de rendre grâce à leur beauté, toujours plein d'attention et de prévenance vis-à-vis de la voracité de la lie athénienne... un archer Scythe... la lie de la lie... suggérer une telle monstruosité ?! «~Paraît qu'elle s'vend même à des flics. / Quelle décadence ! / Y'a plus d'moralité publique / Dans notre France...~» En France peut-être, mais au théâtre d'Aristophane, certainement pas ! (Lisez la note de Willems, t. III pp. 394-395, in « Le nu dans la Comédie ancienne des Grecs », article reproduit en annexe ; purgatif si nécessaire quand on est souillé par de telles salauderies !) A\&O ne nous indiquent pas si le public pouvait se délecter en sus de bruitages émanant des coulisses, quand ça y fornique... Mais pour qui prennent-ils Aristophane ?.. voici donc la littérature qu'on inflige à nos pauvres étudiants... je tiens le cas pour pendable !%
\citComm{et si, par un malheur, j'en avois fait autant,\\
je m'irois, de regret, pendre tout à l'instant.}}%
\note{9}{Le Scythe s'interrompt quelques instants de parler, pause soulignée par l'interjection \italique{extra metrum} οἴμοι qu'il lance une fois descendu : il se demande vers laquelle des deux entrées de l'orchestre se tourner (τί δρᾶσι; ποῖ τὸ γρᾴδι᾽;)}%
\note{10}{\italique{Confer} surtout la note d'Alphonse Willems au vers 1341 des \italique{Guêpes}. Lisez Willems ! aucune excuse : la note est reproduite en annexe.}\\
ἔπειτα κομίζι σ᾽ αὐτίκ᾽. ἀκολούτι τέκνον.%\apparat{1198}{%γϋ. ante v. 1198 et σκ'. ante ἀκολ. in R \espApp 
%κ\inRas{ω}{ο}μίζεις αὐτοῖς \Rpc : κομιεῖς αὐτό (αὖθις Brunck) Biset : (Τοξότ.) κομίσι σ' αὐτός (αὖτις anonym. Parisinus) Enger : (Τοξότ.) \lecon{κομίζι σ᾽ αὐτίκ᾽} Bibi \espApp ἀκουλούτι R : \lecon{ἀκολούτι} \Mon : ἀκουλούτει Brunck}
\\%
%\commentaire{1198}{ἔπειτα κ\inRas{ω}{ο}μίζεις αὐτοῖς.}{Le †αὐτοῖς† du \italique{Ravennas}, en plus d'être \italique{contra metrum}, n'offre aucun sens acceptable. 
%Les diverses tentatives de correction se résument à prolonger les pourparlers. Tantôt Euripide demande son reste ; l'on traduit alors péniblement par "tu le récupéreras une autre fois \auDessus{=}{?} après" :%
%\citComm{(Εὐρ.) ἔπειτα κομιεῖς αὐτό (Biset ; αὐτό \auDessus{=}{?} τὴν συμβήνην)\\
%(Εὐρ.) ἔπειτα κομιεῖς αὖθις (Faber, Brunck, von Velsen, Austin)}%
%Tantôt attribue-t-on au Scythe la réplique qui promet d'apporter à Euripide le reste de la drachme plus tard, ou bien s'enquiert de pouvoir récupérer son carquois en gage... mais le tour paraît bien elliptique :%
%\citComm{(Τοξότ.) ἔπειτα κομίσι σ' αὐτός (... τὸ λοιπόν ?.. τὸ ἀργύριον ?.. ; Enger)\\
%(Τοξότ.) ἔπειτα κομίζεις αὐτόσ᾽ (Fritzsche)\\
%(Τοξότ.) ἔπειτα κομίσι σ' αὖτις (Coulon, Prato)}%
%Ces interprétations me semblent forcer le grec ; elles perdent aussi de vue l'action : Euripide ne cherche assurément pas à être payé rubis sur l'ongle,
%comme dans la chanson, %
%\citComm{Chez la belle Suzon, pas d'argent, pas de cuisse...\marqueur{1}}%
%C'est à cette réplique que se concluent les négociations puisqu'aussitôt le Scythe emmène Élaphion (ἀκολούτι τέκνον).\\
%Euripide, qui ne cherche qu'à éloigner le Cerbère de Mnésiloque, accepte donc l'offre ridicule du Scythe et lui ordonne de quitter la scène -- l'impératif est ici de rigueur : %
%\citComm{ἔπειτα κόμιζ(ε) εἰσαυτίκ(α).\\
%Eh bien alors emmène-la sur-le-champ.}%
%L'émendation appelle quelques justifications : d'abord, εἰσαυτίκα est autrement attesté au vers 367 de la \italique{Paix} ; quant à l'hypothétique confusion αυτικ / αυτοις, elle n'est pas invraisemblable en constatant la ressemblance entre le digramme `ις' et le `κ' ; il faudrait alors supposer l'existence d'une forme intermédiaire\marqueur{2} αυτιις ou αυτις, bien vite corrigée en αυτοις.}%
%\debutNote{1198}%
%\note{1}{}%
%\note{2}{}
σὺ δὲ τοῦτο τήρει τὴ γέροντο, γρᾴδιο.\\
ὄνομα δέ σοι τί ἔστιν;\\\passeLigne}%
\end{groupe}%
%% v. 1200b
\begin{groupe}%
\personnageGrec{Εὐριπίδης}%
\trimIamb{\ligne{1200}\decalage{ὄνομα δέ σοι τί ἔστιν;}Ἀρτεμισία.\\\passeLigne}%
\end{groupe}%
%% v. 1201
\begin{groupe}%
\personnageGrec{Τοξότης}%
\trimIamb{μεμνῆσι τοίνυν τοὔνομ᾽· Ἀρταμουξία.\apparat{1201}{ἀρτομουξία R : \lecon{Ἀρταμουξία} ed. Lipsiensis (1814)}%
\commentaire{1201}{Ἀρταμουξία.}{L'archer scythifie le nom d'Ἀρτεμισία : il l'entend comme un composé de ἀρτα-, multiplie les `α' et ajoute un `ξ'\marqueur{1}. Cf. le parler de Ψευδ\textbf{αρτά}βας dans les \italique{Acharniens} (v. 100) :%
\citComm{ἰαρταμὰν ἐξάρξαν ἀπισσόνα σάτρα.}%
}%
\debutNote{1201}%
\note{1}{Voir le commentaire vv. 91-2 d'Olson aux \italique{Acharniens} ainsi que l'article «~The Archer Scene in Aristophanes' \italique{Thesmophoriazusae}~» d'Édith M. Hall, in \italique{Philologus} 133 (1989), p. 40. On peut citer Ἀρταβάτας, Ξέρξης, Ἀρταξέρξης, Ἀρταφρένης, Ἀρτάβαζος \&c.}\\\passeLigne}%
\end{groupe}%
%% v. 1202-1206
\begin{groupe}%
\personnageGrec{Εὐριπίδης}%
\trimIamb{Ἑρμῆ δόλιε ταυτὶ μὲν ἔτι καλῶς ποιεῖς.\\
σὺ μὲν οὖν ἀπότρεχε παιδάριον τουτὶ λαβών·\apparat{1203}{\lecon{τουτὶ} R : ταυτὶ Dobree, cf. Schol.}%
\commentaire{1203-11}{Papyrus d'Oxyrhynque n° 4935 fr. 2 \& n° 5132.}{Deux fragments d'un même papyrus d'Oxyrhynque\marqueur{1} contiennent quelques lettres de ce passage. (Voir en annexe.) Ils n'apportent aucune leçon originale. Au vers 1208, l'espace manque pour être compatible avec la leçon \italique{ante correctionem} λέλυσον du \italique{Ravennas} ; l'état du fragment ne permet pas pour autant de trancher entre les corrections λέλυσο (\Rpc) et λέλυσαι (Bentley).}%
\debutNote{1203-11}%
\note{1}{N° 4935 fr. 2, in \italique{The Oxyrhynchus Papyri}, LXXIII (2009), ed. A. Benaissa. N° 5132, \italique{The Oxyrhynchus Papyri}, LXXVIII (2012), ed. W. B. Henry.}\\
ἐγὼ δὲ λύσω τόνδε. σὺ δ᾽ ὅπως ἀνδρικῶς\\
ὅταν λυθῇς τάχιστα φεύξει καὶ τενεῖς\\
ὡς τὴν γυναῖκα καὶ τὰ παιδί᾽ οἴκαδε.\\\passeLigne}%
\end{groupe}%
%% v. 1207
\begin{groupe}%
\personnageGrec{Μνησίλοχος}%
\trimIamb{ἐμοὶ μελήσει ταῦτά γ᾽ ἢν ἅπαξ λυθῶ.\\\passeLigne}%
\end{groupe}%
%% v. 1208-1209a
\begin{groupe}%
\personnageGrec{Εὐριπίδης}%
\trimIamb{λέλυσο. σὸν ἔργον, φεῦγε πρὶν τὸν τοξότην\apparat{1208}{ca. 6 litt. $]$· σ\douteux{ο}$[$νεργον \PsiQuatre : \lecon{λέλυσο}\inRas{ν}{} \Rpc : λέλυσαι Bentley}\\
ἥκοντα καταλαβεῖν.\\\passeLigne}%
\end{groupe}%
%% v. 1209b
\begin{groupe}%
\personnageGrec{Μνησίλοχος}%
\trimIamb{\ligne{1209}\decalage{ἥκοντα καταλαβεῖν.}ἐγὼ δὴ τοῦτο δρῶ.\\\passeLigne}%
\end{groupe}%
%% v.1210-1216
\begin{groupe}%
%\ledpb%
\personnageGrec{Τοξότης}%
\trimIamb{ὦ γρᾴδι᾽ ὠς καρίεντό σοι τὸ τυγάτριο,\apparat{1210}{τυγάτριον ex τυγάθριον ut videtur \Rpc}\\
κοὐ δύσκολ᾽ ἀλλὰ πρᾶο. ποῦ τὸ γρᾴδιο;\apparat{1211}{δύσκολλ᾽ R : \lecon{δύσκολ᾽} Grynaeus}\\
οἴμ᾽ ὠς ἀπόλωλο· ποῦ τὸ γέροντ᾽ ἐντευτενί;\\
ὦ γρᾴδι᾽, ὦ γρᾶ᾽. οὐκ ἐπαινῶ γρᾴδιο.\apparat{1213}{γρᾶίο R : \lecon{γρᾶ᾽} Bentley}\\
Ἀρταμουξία.\\
\ligne{1214}διέβαλλέ μ᾽ ἡ γραῦς. ἀπότρεκ᾽ ὠς τάκιστα σύ·\testimonia{1214}{Suidas δ 892 διέβαλέ μ' ἡ γραῦς ; \Eustathe 1406.19}\apparat{1214}{\lecon{διέβαλλε} R : διέβαλέ Schol., Suid. δ 892, Eust. p. 1406.19 \espApp ὁ γραῦς R : \lecon{ἡ γραῦς} Suid., Eust. \espApp ὡς τάχιστα R : ὠς τάκιστα Brunck \ligneApp{1215} συβήνη R : \lecon{συμβήνη} Bibi, cf. v.1197, Hesych. σ 2137 \ligneApp{1216} δράσει R : \lecon{δρᾶσι} Blaydes (1842) \espApp γράδιο R : \lecon{γρᾴδιο} \Mon}\\
ὀρτῶς δὲ συμβήνη 'στί· καταβηνῆσι γάρ.\\
οἴμοι, τί δρᾶσι; ποῖ τὸ γρᾴδιο;\commentaire{1216}{οἴμοι, τί δράσει; ποῖ τὸ γράδιο·}{Il manque un pied à ce vers.}\\
Ἀρταμουξία.\\\passeLigne}%
\end{groupe}%
%% v. 1217
\begin{groupe}%
\personnageGrec{Κορυφαῖος}%
\trimIamb{\ligne{1217}τὴν γραῦν ἐρωτᾷς, ἣ 'φερεν τὰς πηκτίδας;\testimonia{1217}{Suidas π 1502}\apparat{1217}{γραῦ R : \lecon{γραῦν} Suid. π 1502}\\\passeLigne}%
\end{groupe}%
%% v. 1218a
\begin{groupe}%
\personnageGrec{Τοξότης}%
\trimIamb{ναὶ ναῖκι. εἶδες αὐτό;\\\passeLigne}%
\end{groupe}%
%% v. 1218b-1219
\begin{groupe}%
\personnageGrec{Κορυφαῖος}%
\trimIamb{\ligne{1218}\decalage{ναὶ ναῖκι. εἶδες αὐτό;}ταύτῃ γ᾽ οἴχεται\\
αὐτή τ᾽ ἐκείνη καὶ γέρων τις εἵπετο.\apparat{1219}{%Σκ. ante versum in mg. et Χορ. ante καὶ γέρον in R \espApp 
γέρον R : \lecon{γέρων} \Mon}\\\passeLigne}%
\end{groupe}%
%% v. 1220a
\begin{groupe}%
\personnageGrec{Τοξότης}%
\trimIamb{κροκῶτ᾽ ἔκοντο τὴ γέροντο;\\\passeLigne}%
\end{groupe}%
%% v. 1220b-1221
\begin{groupe}%
\personnageGrec{Κορυφαῖος}%
\trimIamb{\ligne{1220}\decalage{κροκῶτ᾽ ἔκοντο τὴ γέροντο;}φήμ᾽ ἐγώ.\\
ἔτ᾽ ἂν καταλάβοις, εἰ διώκοις ταυτῃί.\\\passeLigne}%
\end{groupe}%
%% v. 1222a-1222b
\begin{groupe}%
%\ledpb%
\personnageGrec{Τοξότης}%
\trimIamb{ὦ μιαρὸ γρᾷο· πότερα τρέξι τὴν ὀδό;\apparat{1222}{γραῦ et ὁδό R : \lecon{γρᾷο} et \lecon{ὀδό} Brunck}\\
\ligne{1222}Ἀρταμουξία.\\\passeLigne}%
\end{groupe}%
%% v. 1223-1224
\begin{groupe}%
\personnageGrec{Κορυφαῖος}%
\trimIamb{ὀρθὴν ἄνω δίωκε. ποῖ θεῖς; οὐ πάλιν\\
τῃδὶ διώξει; ᾽ς τοὔμπαλιν τρέχεις σύ γε.\apparat{1224}{δ᾽ ὤξεις R : \lecon{διώξει} Elmsley (ad \italique{Ach.} 278) \lecon{᾽ς} Cobet (\italique{Mnemos.} 2 (1853), p.213)}\\\passeLigne}%
\end{groupe}%
%% v. 1225
\begin{groupe}%
\personnageGrec{Τοξότης}%
\trimIamb{κακόδαιμον, ἀλλὰ τρέξι. Ἀρταμουξία.\\\passeLigne}%
\end{groupe}%
%% v. 1226-1231
\begin{groupe}%
\personnageGrec{Κορυφαῖος}%
\typeVers{anap.}\anap{τρέχε νυν κατ᾽ αὐτοὺς ἐς κόρακας ἐπουρίσας.\apparat{1226}{κατὰ τοὺς R : \lecon{κατ᾽ αὐτοὺς} Jackson (p.82) \lecon{ἐς} O. Bachmann \ligneApp{1127} πέπυσται R : \lecon{πέπαισται} Schol.}\\
ἀλλὰ πέπαισται μετρίως ἡμῖν·\\%\apparat{1227}{πέπυσται R : \lecon{πέπαισται} Schol.}\\
ὥσθ᾽ ὥρα δή 'στι βαδίζειν\\
οἴκαδ᾽ ἑκάστῃ. τὼ Θεσμοφόρω δ᾽\commentaire{1129-31}{τὼ Θεσμοφόρω δ᾽ / ἡμῖν ἀγαθὴν / τούτων χάριν ἀνταποδοίτην.}{Selon son habitude -- qui était peut-être d'usage dans l'Ancienne comédie --, Artistophane conclue sa pièce par des souhaits de succès (il vise bien-sûr la victoire au concours comique).}\\
ἡμῖν ἀγαθὴν\\
τούτων χάριν ἀνταποδοίτην.\testimonia{1231}{Subscriptio \PsiUn \ fr. 6 ΑΡΙΣΤΟΦΑΝ$]$ΟΥΣ    ΘΕΣΜΟΦΟΡΙ$]$ΑΖΥΣΑΙ   $]$ΧΗΗΙ ($=$ 1201, versuum numerus)}\apparat{1231}{ἀνταδοῖτον R : \lecon{ἀνταποδοῖτον} Grynaeus}\\\passeLigne}%
\end{groupe}%
\endnumbering
\end{Leftside}%
%%
%%
%%
%%
%%
%%
%%
%%
\begin{Rightside}%
\beginnumbering
%% v. 1-4
\begin{groupe}%
\personnageFr{Mnésiloque}%
\trimIambFr{Ô Zeus, est-ce qu'à la fin l'étoile du matin paraîtra ?\\
Cet homme me tuera en me faisant tribuler à l'aube.\\
Y a-t-il moyen, avant que perdre tout de bon la rate,\\
de savoir de toi où tu me mènes, Euripide ?\\\passeLigne}%
\end{groupe}%
%% v. 5-6a
\begin{groupe}%
\personnageFr{Euripide}%
\trimIambFr{Eh ! il n'est pas besoin que tu entendes tout ce que tout à l'heure\\
tu verras de tes yeux.\\\passeLigne}%
\end{groupe}%
%% v. 6b-7a
\begin{groupe}%
\personnageFr{Mnésiloque}%
\trimIambFr{\ligne{6}\espace{2.5cm}Comment dis-tu ? Répète.\\
Il n'est pas besoin que j'entende.\\\passeLigne}%
\end{groupe}%
%% v. 7b
\begin{groupe}%
\personnageFr{Euripide}%
\trimIambFr{\espace{2.7cm}\ligne{7}Non, ce que tu vas voir.\\\passeLigne}%
\end{groupe}%
%% v. 8a
\begin{groupe}%
\personnageFr{Mnésiloque}%
\trimIambFr{Ni donc que je voie...\\\passeLigne}%
\end{groupe}%
%% v. 8b
\begin{groupe}%
\personnageFr{Euripide}%
\trimIambFr{\espace{3.1cm}\ligne{8}Non, ce qu'il faut que tu entendes.\\\passeLigne}%
\end{groupe}%
%% v. 9-10
\begin{groupe}%
\personnageFr{Mnésiloque}%
\trimIambFr{Singulière recommandation ! Ma foi, c'est finement dit.\\
Tu prétends que je ne dois ni entendre ni voir.\\\passeLigne}%
\end{groupe}%
%% v. 11
\begin{groupe}%
\personnageFr{Euripide}%
\trimIambFr{C'est que distincte est l'essence de chacune de ces choses.\\\passeLigne}%
\end{groupe}%
%% v. 12a
\begin{groupe}%
\personnageFr{Mnésiloque}%
\trimIambFr{Ne pas entendre et ne pas voir ?\\\passeLigne}%
\end{groupe}%
%% v. 12b
\begin{groupe}%
\personnageFr{Mnésiloque}%
\trimIambFr{\ligne{12}\espace{4.4cm}Tiens-le pour assuré.\\\passeLigne}%
\end{groupe}%
%% v. 13a
\begin{groupe}%
\personnageFr{Mnésiloque}%
\trimIambFr{Comment distincte ?\\\passeLigne}%
\end{groupe}%
%% v. 13b-18
\begin{groupe}%
\personnageFr{Euripide}%
\trimIambFr{\ligne{13}\espace{1.8cm}Voici comment elles furent départies à l'origine.\\
L'Éther, quand d'abord il se dégagea\\
et procréa en soi-même des animaux doués de mouvement,\\
commença par fabriquer pour qui doit voir\\
un œil analogique de l'orbe solaire\\
et fora une oreille à l'image des godets d'un moule.\\\passeLigne}%
\end{groupe}%
%% v. 19-21
\begin{groupe}%
\personnageFr{Mnésiloque}%
\trimIambFr{Et à cause de ce moule je ne dois ni entendre ni voir ?\\
Par Zeus, je suis ravi d'apprendre encore cela.\\
Quelle fameuse chose tout de même, que les doctes accointances !\\\passeLigne}%
\end{groupe}%
%% v. 22a
\begin{groupe}%
\personnageFr{Euripide}%
\trimIambFr{Auprès de moi, tu en apprendras bien d'autres du même genre.\\\passeLigne}%
\end{groupe}%
%% v. 22b-24
\begin{groupe}%
\personnageFr{Mnésiloque}%
\trimIambFr{\ligne{22}\espace{5.3cm}Ah ! que ne puis-je\\
en outre apprendre aussi à être perclus des jambes ?\\\passeLigne}%
\end{groupe}%
%% v. 25a
\begin{groupe}%
\personnageFr{Euripide}%
\trimIambFr{Avance ici et prête-moi ton attention.\\\passeLigne}%
\end{groupe}%
%% v. 25b
\begin{groupe}%
\personnageFr{Mnésiloque}%
\trimIambFr{\ligne{25}\espace{5.2cm}J'y suis.\\\passeLigne}%
\end{groupe}%
%% v. 26a
\begin{groupe}%
\personnageFr{Euripide}%
\trimIambFr{Tu vois cette petite porte ?\\\passeLigne}%
\end{groupe}%
%% v. 26b-27a
\begin{groupe}%
\personnageFr{Mnésiloque}%
\trimIambFr{\ligne{26}\espace{3.4cm}Oui, par Héraclès,\\
à ce que je crois, du moins.\\\passeLigne}%
\end{groupe}%
%% v. 27b
\begin{groupe}%
\personnageFr{Euripide}%
\trimIambFr{\ligne{27}\espace{1.4cm}Maintenant, motus.\\\passeLigne}%
\end{groupe}%
%% v. 27c
\begin{groupe}%
\personnageFr{Mnésiloque}%
\trimIambFr{\ligne{27}\espace{3.4cm}Que je ne dise mot de la petite porte ?\\\passeLigne}%
\end{groupe}%
%% v. 28a
\begin{groupe}%
\personnageFr{Euripide}%
\trimIambFr{Écoute, là !\\\passeLigne}%
\end{groupe}%
%% v. 28b
\begin{groupe}%
\personnageFr{Mnésiloque}%
\trimIambFr{\ligne{28}\espace{1.4cm}Que j'écoute et ne dise mot de la petite porte \\\passeLigne}%
\end{groupe}%
%% v. 29-30a
\begin{groupe}%
\personnageFr{Euripide}%
\trimIambFr{C'est là que se trouve demeurer l'illustre Agathon,\\
le poète tragique.\\\passeLigne}%
\end{groupe}%
%% v. 30b
\begin{groupe}%
\personnageFr{Mnésiloque}%
\trimIambFr{\ligne{30}\espace{2.4cm}Qui est-ce, cet Agathon ?\\
Est-ce l'un de ces honnêtes hommes... celui au teint hâlé, plein d'allant ?\\\passeLigne}%
\end{groupe}%
%% v. 32
\begin{groupe}%
\personnageFr{Euripide}%
\trimIambFr{Non, c'est un autre. Tu ne l'as jamais vu ?\\\passeLigne}%
\end{groupe}%
%% v. 33a
\begin{groupe}%
\personnageFr{Mnésiloque}%
\trimIambFr{L'homme à la barbe touffue ?\\\passeLigne}%
\end{groupe}%
%% v. 33b
\begin{groupe}%
\personnageFr{Euripide}%
\trimIambFr{\ligne{33}\espace{3cm}Tu ne l'as jamais vu !..\\\passeLigne}%
\end{groupe}%
%% v. 34
\begin{groupe}%
\personnageFr{Mnésiloque}%
\trimIambFr{Non, par Zeus, pas du moins que je sache.\\\passeLigne}%
\end{groupe}%
%% v. 35-38
\begin{groupe}%
\personnageFr{Euripide}%
\trimIambFr{À coup sûr tu l'as enfourché, mais peut-être ne t'en souvient-il pas.\\
\lednopb%
Mais allons nous tapir à l'écart. Car sort de chez lui\\
\lednopb%
un serviteur avec du feu et des myrtes.\\
\lednopb%
Il va sacrifier, ce semble, en faveur du succès de sa poésie.\\\passeLigne}%
\end{groupe}%
%% v. 39-45a
\begin{groupe}%
\personnageFr{Serviteur d'Agathon}%
\dimAnapFr{Que le peuple entier se recueille\\
et ferme la bouche, car est venu séjourner\\
un thiase de muses dans la demeure}%
\end{groupe}%
%% v. 42-45a
\begin{groupe}%
\dimAnapFr{de mon maître pour composer des chants.\\
Que retienne ses haleines l'éther paisible,\\
que le flot de la mer point ne retentisse\\
glauque...\\\passeLigne}%
\end{groupe}%
%% v. 45b
\begin{groupe}%
\personnageFr{Mnésiloque}%
\dimAnapFr{\ligne{45}\espace{1.5cm}Bombax !\\\passeLigne}%
\end{groupe}%
%% v. 45c
\begin{groupe}%
\personnageFr{Euripide}%
\dimAnapFr{\ligne{45}\espace{2.8cm}Chut !\\\passeLigne}%
\end{groupe}%
%% v. 45d
\begin{groupe}%
\personnageFr{Mnésiloque}%
\dimAnapFr{\ligne{45}\espace{3.6cm}Dit-il quelque chose ?\\\passeLigne}%
\end{groupe}%
%% v. 46-48a
\begin{groupe}%
\personnageFr{Serviteur}%
\dimAnapFr{Que la gent ailée s'assoupisse,\\
que les pieds des bêtes sauvages vaguant par les bois\\
point ne se meuvent...\\\passeLigne}%
\end{groupe}%
%% v. 48b
\begin{groupe}%
\personnageFr{Mnésiloque}%
\dimAnapFr{\ligne{48}\espace{1.9cm}Bombalo bombax !\\\passeLigne}%
\end{groupe}%
%% v. 49-50a
\begin{groupe}%
\personnageFr{Serviteur}%
\dimAnapFr{... Car il s'apprête l'homme au beau langage, Agathon,\\
il se penche notre...\\\passeLigne}%
\end{groupe}%
%% v. 50b
\begin{groupe}%
\personnageFr{Mnésiloque}%
\dimAnapFr{\ligne{50}\espace{2.6cm}Et se fait enfiler ?\\\passeLigne}%
\end{groupe}%
%% v. 51a
\begin{groupe}%
\personnageFr{Serviteur}%
\dimAnapFr{Qui a parlé?\\\passeLigne}%
\end{groupe}%
%% v. 51b
\begin{groupe}%
\personnageFr{Mnésiloque}%
\dimAnapFr{\ligne{51}\espace{2.3cm}L'éther paisible.\\\passeLigne}%
\end{groupe}%
%% v. 52-57a
\begin{groupe}%
\personnageFr{Serviteur}%
\dimAnapFr{... Il va poser des colombiers, fondement d'un drame.\\
Il cambre, en jantes nouvelles, des vers.\\
Il met les uns sur le tour, il assemble les autres.\\
Il frappe des sentences, fait s'opposer des mots,\\
malaxe pour modeler, compactifier,\\
puis introduit le moule...\\\passeLigne}%
\end{groupe}%
%% v. 57b
\begin{groupe}%
\personnageFr{Mnésiloque}%
\dimAnapFr{\ligne{57}\espace{2.2cm}Et avale.\\\passeLigne}%
\end{groupe}%
%% v. 58
\begin{groupe}%
\personnageFr{Serviteur}%
\dimAnapFr{Quel est le manant qui approche cette enceinte ?\\\passeLigne}%
\end{groupe}%
%% v. 59-61
\begin{groupe}%
\personnageFr{Mnésiloque}%
\dimAnapFr{Quelqu'un de prêt, pour toi et pour ton poète\\
au beau langage, par votre enceinte,\\
à introduire cet outil-ci}%
\end{groupe}%
%% v. 62
\begin{groupe}%
\dimAnapFr{bien compact et ramassé..\\\passeLigne}%
\end{groupe}%
%% v. 63
\begin{groupe}%
\personnageFr{Serviteur}%
\trimIambFr{Dans ta jeunesse tu devais être un ribaud, vieillard..\\\passeLigne}%
\end{groupe}%
%% v. 64-65
\begin{groupe}%
\personnageFr{Euripide}%
\trimIambFr{(À Mnésiloque :) Diantre d'homme, laisse-le tranquille. (Au serviteur :) Toi,\\
va m'appeler Agathon bon gré, mal gré.\\\passeLigne}%
\end{groupe}%
%% v. 66-69
\begin{groupe}%
\personnageFr{Serviteur}%
\trimIambFr{Inutile d'insister. Lui-même ne tardera pas à sortir,\\
car il commence un chant lyrique. Et de fait, en hiver,\\
il n'est pas facile de contourner ses strophes,\\
sans aller se mettre dehors au soleil.\\\passeLigne}%
\end{groupe}%
%% v. 70a
\begin{groupe}%
\personnageFr{Mnésiloque}%
\trimIambFr{Et moi, que vais-je faire ?\\\passeLigne}%
\end{groupe}%
%% v. 70b
\begin{groupe}%
\personnageFr{Euripide}%
\trimIambFr{\ligne{70}\espace{2.6cm}Demeure, puisqu'il arrive.\\\passeLigne}%
\end{groupe}%
%% v. 71
\begin{groupe}%
\personnageFr{Mnésiloque}%
\trimIambFr{Ô Zeus, que prétends-tu faire de moi aujourd'hui ?}%
\end{groupe}%
%% v. 72-74
\begin{groupe}%
\trimIambFr{Par les dieux, je tiens à savoir de quoi il retourne.\\
Qu'as-tu à geindre ? à te contrarier ?\\
Il n'y a pas de mystère à faire avec moi, puisque tu es mon gendre.\\\passeLigne}%
\end{groupe}%
%% v. 75
\begin{groupe}%
\personnageFr{Euripide}%
\trimIambFr{Un grand malheur, manigancé d'avance, me menace.\\\passeLigne}%
\end{groupe}%
%% v. 76a
\begin{groupe}%
\personnageFr{Mnésiloque}%
\trimIambFr{De quelle sorte ?\\\passeLigne}%
\end{groupe}%
%% v. 76b-77
\begin{groupe}%
\personnageFr{Euripide}%
\trimIambFr{\ligne{76}\espace{1.7cm}Ce jour décidera\\
si Euripide vit encore ou s'il est mort.\\\passeLigne}%
\end{groupe}%
%% v. 78-80
\begin{groupe}%
\personnageFr{Mnésiloque}%
\trimIambFr{Comment cela ? Car ni les tribunaux\\
ne sont dans le cas de juger, ni la Boulé ne siège,\\
puisque nous sommes au jour intermédiaire des Thesmophories.%\commentaire{80}{}{Les difficultés que soulève ce passage me semblent avoir été résolues avec autant de clarté que de savoir par Guignaut, dans une note sur CREUZER, Religions de l'Antiquité, t. III, 3, 1150 et suiv.}%
\\\passeLigne}%
\end{groupe}%
%% v. 81-84a
\begin{groupe}%
\personnageFr{Euripide}%
\trimIambFr{C'est précisément ce qui fait que je m'attends à périr.\\
Car les femmes ont monté un coup contre moi.\\
Elles doivent, dans le temple des Thesmophories, à mon sujet aujourd'hui,\\
tenir assemblée afin de me perdre.\\\passeLigne}%
\end{groupe}%
%% v. 84b
\begin{groupe}%
\personnageFr{Mnésiloque}%
\trimIambFr{\ligne{84}\espace{4cm}Pourquoi cela ?\\\passeLigne}%
\end{groupe}%
%% v. 85
\begin{groupe}%
\personnageFr{Euripide}%
\trimIambFr{Parce que je les théâtralise et parle mal d'elles.\\\passeLigne}%
\end{groupe}%
%% v. 86-87
\begin{groupe}%
\personnageFr{Mnésiloque}%
\trimIambFr{Par Poseidôn, tu serais châtié en justice.\\
\lednopb%
Mais, pour y échapper, de quelle machine disposes-tu ?\\\passeLigne}%
\end{groupe}%
%% v. 88-89a
\begin{groupe}%
\personnageFr{Euripide}%
\trimIambFr{De persuader à Agathon, le poète tragique,\\
de se rendre aux Thesmophories.\\\passeLigne}%
\end{groupe}%
%% v. 89b
\begin{groupe}%
\personnageFr{Mnésiloque}%
\trimIambFr{\ligne{89}\espace{3.5cm}Pour quoi faire, dis-moi ?\\\passeLigne}%
\end{groupe}%
%% v. 90-91a
\begin{groupe}%
\personnageFr{Euripide}%
\trimIambFr{Pour intervenir à l'assemblée devant les femmes et, autant qu'il est nécessaire,\\
me défendre.\\\passeLigne}%
\end{groupe}%
%% v. 91b
\begin{groupe}%
\personnageFr{Mnésiloque}%
\trimIambFr{\ligne{91}\espace{2.8cm}Ouvertement ou en tapinois ?\\\passeLigne}%
\end{groupe}%
%% v. 92
\begin{groupe}%
\personnageFr{Euripide}%
\trimIambFr{En tapinois, revêtu d'un costume féminin.\\\passeLigne}%
\end{groupe}%
%% v. 93-94
\begin{groupe}%
\personnageFr{Mnésiloque}%
\trimIambFr{C'est ingénieux et bien dans tes façons.\\
En fait d'artifices, à nous le gâteau.\\\passeLigne}%
\end{groupe}%
%% v. 95a
\begin{groupe}%
\personnageFr{Euripide}%
\trimIambFr{Silence.\\\passeLigne}%
\end{groupe}%
%% v. 95b
\begin{groupe}%
\personnageFr{Mnésiloque}%
\trimIambFr{\ligne{95}\espace{1cm}Qu'y a-t-il ?\\\passeLigne}%
\end{groupe}%
%% v. 95c
\begin{groupe}%
\personnageFr{Euripide}%
\trimIambFr{\ligne{95}\espace{2.7cm}Agathon sort.\\\passeLigne}%
\end{groupe}%
%% v. 96a
\begin{groupe}%
\personnageFr{Mnésiloque}%
\trimIambFr{Quel homme est-ce ?\\\passeLigne}%
\end{groupe}%
%% v. 96b
\begin{groupe}%
\personnageFr{Euripide}%
\trimIambFr{\ligne{96}\espace{2cm}Celui-là qu'on roule dehors.\\\passeLigne}%
\end{groupe}%
%% v. 97-98
\begin{groupe}%
\personnageFr{Mnésiloque}%
\trimIambFr{Est-ce donc que je suis aveugle ? Je ne vois pas qu'il y ait\\
un homme ici. C'est Cyrène que je vois.\\\passeLigne}%
\end{groupe}%
%% v. 99
\begin{groupe}%
\personnageFr{Euripide}%
\trimIambFr{Chut ! Il se dispose à chanter.\\\passeLigne}%
\end{groupe}%
%% v. 100
\begin{groupe}%
\personnageFr{Mnésiloque}%
\trimIambFr{Des sentes de Fourmi, ou que sont ces fredons?\\\passeLigne}%
\end{groupe}%
%% v. 101-129
\begin{groupe}%
\personnageFr{Agathon}%
\chantAgathFr{1° — Des deux Chtoniennes, recevez}%
\end{groupe}%
%% v. 101-129
\begin{groupe}%
\chantAgathFr{la torche consacrée, vierges,}%
\end{groupe}%
%% v. 101-129
\begin{groupe}%
\chantAgathFr{\ligne{101}et avec votre libre patrie}%
\end{groupe}%
%% v. 101-129
\begin{groupe}%
\chantAgathFr{entonnez un chœur de danse.\\}%
\end{groupe}%
%% v. 101-129
\begin{groupe}%
\chantAgathFr{\ligne{104}2° — À quelle divinité ce cômos ?}%
\end{groupe}%
%% v. 101-129
\begin{groupe}%
\chantAgathFr{Parle. Je me laisse aisément induire}%
\end{groupe}%
%% v. 101-129
\begin{groupe}%
\chantAgathFr{à honorer les dieux.\\}%
\end{groupe}%
%% v. 101-129
\begin{groupe}%
\chantAgathFr{3° — Allons, Muse, glorifie}%
\end{groupe}%
%% v. 101-129
\begin{groupe}%
\chantAgathFr{l'archer à l'arc d'or,}%
\end{groupe}%
%% v. 101-129
\begin{groupe}%
\chantAgathFr{\ligne{109}Phébus, qui fonda les murs}%
\end{groupe}%
%% v. 101-129
\begin{groupe}%
\chantAgathFr{de cette cité dans la contrée du Simois.\\}%
\end{groupe}%
%% v. 101-129
\begin{groupe}%
\chantAgathFr{4° — Sois salué de nos plus beaux chants,}%
\end{groupe}%
%% v. 101-129
\begin{groupe}%
\chantAgathFr{Phébus, qui dans les fêtes des Muses}%
\end{groupe}%
%% v. 101-129
\begin{groupe}%
\chantAgathFr{\ligne{112}emportes le prix sacré.\\}%
\end{groupe}%
%% v. 101-129
\begin{groupe}%
\chantAgathFr{\ligne{114}5° — Et chantez la Vierge qui se plaît dans}%
\end{groupe}%
%% v. 101-129
\begin{groupe}%
\chantAgathFr{les montagnes mères des chênes, Artémis la chasseresse.\\}%
\end{groupe}%
%% v. 101-129
\begin{groupe}%
\chantAgathFr{6° — À ta suite je célèbre et glorifie }%
\end{groupe}%
%% v. 101-129
\begin{groupe}%
\chantAgathFr{l'auguste rejeton de Léto,}%
\end{groupe}%
%% v. 101-129
\begin{groupe}%
\chantAgathFr{\ligne{118}Artémis ignorant l'hymen.\\}%
\end{groupe}%
%% v. 101-129
\begin{groupe}%
\chantAgathFr{7° — Ainsi que Léto et les percussions de la cithare d'Asie,}%
\end{groupe}%
%% v. 101-129
\begin{groupe}%
\chantAgathFr{aux beaux rythmes contrariés marqués par le pied}%
\end{groupe}%
%% v. 101-129
\begin{groupe}%
\chantAgathFr{sous la direction des Charites phrygiennes.\\}%
\end{groupe}%
%% v. 101-129
\begin{groupe}%
\chantAgathFr{\ligne{123}8° — J'honore Léto la souveraine}%
\end{groupe}%
%% v. 101-129
\begin{groupe}%
\chantAgathFr{et la cithare mère des hymnes,}%
\end{groupe}%
%% v. 101-129
\begin{groupe}%
\chantAgathFr{à la voix mâle et éprouvée.\\}%
\end{groupe}%
%% v. 101-129
\begin{groupe}%
\chantAgathFr{9° — À ces accents, une lumière jaillit}%
\end{groupe}%
%% v. 101-129
\begin{groupe}%
\chantAgathFr{dans les yeux des dieux\\
\ligne{126}et traverse notre soudain regard.}%
\end{groupe}%
%% v. 101-129
\begin{groupe}%
\chantAgathFr{Aussi bien, glorifie le seigneur Phébus.\\}%
\end{groupe}%
%% v. 101-129
\begin{groupe}%
\chantAgathFr{10° — Salut, heureux fils de Léto.\\\passeLigne}%
\end{groupe}%
%% v. 130-145
\begin{groupe}%
\personnageFr{Mnésiloque}%
\trimIambFr{Quel chant suave, ô Génétyllides vénérées\\
et qui sent la féminie et les coups de langue\\
et les baisers lascifs. En l'écoutant\\
j'ai ressenti un prurit jusque sous le fondement.\\
Jouvenceau, si tu en es un, à la façon d'Eschyle\\
dans sa \italique{Lycurgie}, je veux t'interroger.\\
D'où es-tu, ô femmelette ? Quelle est ta patrie ? Qu'est-ce que cet accoutrement,}
\end{groupe}%
%% v. 130-145
\begin{groupe}%
\trimIambFr{cette incohérence dans tes façons ? Un barbitos,\\
qu'a-t-il à dire à une crocote, une lyre à une résille ?\\
Quoi ! une pomme d'Adam avec un strophion ? Quelle disconvenance !\\
Le rapport d'un miroir à une épée ?\\
Et toi-même, ô jeunesse, qui es-tu ? Un homme ?\\
Mais où est ton cas ? où ton manteau ? où tes laconiennes ?\\
Alors tu es femme ? Mais où sont tes seins ?\\
Que dis-tu ? Pourquoi te taire ? Alors d'après ton chant,\\
je te déchiffre puisque tu ne veux pas t'expliquer toi-même.\\\passeLigne}%
\end{groupe}%
%% v. 146-152
\begin{groupe}%
\personnageFr{Agathon}%
\trimIambFr{Vieillard, vieillard, ces reproches dictés par l'envie,\\
je les ai entendus, mais n'en ai point senti l'atteinte.\\
Je porte la mise congruente à mes pensées.\\
Car c'est une loi que le poète, suivant les drames\\
qu'il lui faut faire, modèle sur eux ses mœurs.\\
Et par exemple, si l'on compose des drames féminins,\\
il faut que le corps participe aux mœurs de la femme..\\\passeLigne}%
\end{groupe}%
%% v. 153
\begin{groupe}%
\personnageFr{Mnésiloque}%
\trimIambFr{Ainsi tu chevauches quand tu composes une Phèdre ?\\\passeLigne}%
\end{groupe}%
%% v. 154-156
\begin{groupe}%
\personnageFr{Agathon}%
\trimIambFr{Compose-t-on des drames mâles, dans le corps\\
réside ce qui est requis. Mais ce que nous ne possédons pas,\\
l'imitation est là qui l'attrape.\\\passeLigne}%
\end{groupe}%
%% v. 157-158
\begin{groupe}%
\personnageFr{Mnésiloque}%
\trimIambFr{Eh bien, quand tu feras des drames satyriques, appelle-moi,\\
et je coopérerai avec toi, de derrière, en toute raideur.\\\passeLigne}%
\end{groupe}%
%% v. 159-167
\begin{groupe}%
\personnageFr{Agathon}%
\trimIambFr{D'ailleurs, c'est chose choquante à voir qu'un poète\\
grossier et hirsute. Avise-toi que\\
le fameux Ibycos, et Anacréon de Téos\\
et Alcée, qui ont donné du montant à l'harmonie,\\
portaient la mitre ainsi que le dikanon à l'ionienne.\\
Et Phrynichos%\commentaire{154}{Φρύνιχος}{Le poète Phrynichos, qui vivait avant et pendant les guerres médiques, portait naturellement le riche costume de son temps. Cf. Caval., 1331.}
, — celui-là, en effet, tu l'as entendu, —\\
beau lui-même, il s'habillait de beaux vêtements\\
et c'est aussi ce qui fait que ses drames étaient beaux.\\
Car c'est une nécessité que l'on compose suivant sa nature.\\\passeLigne}%
\end{groupe}%
%% v. 168-170
\begin{groupe}%
\personnageFr{Mnésiloque}%
\trimIambFr{Voilà aussi pourquoi Philoclès, qui est laid, compose de laides choses ;\\
\lednopb%
pourquoi Xénoclès, qui est mauvais, en compose de mauvaises ;\\
\lednopb%
pourquoi Théognis, qui est froid, compose froidement.\\\passeLigne}%
\end{groupe}%
%% v. 171-172a
\begin{groupe}%
\personnageFr{Agathon}%
\trimIambFr{De toute rigueur. C'est, en effet, pour avoir su cela\\
que j'ai soigné ma personne.\\\passeLigne}%
\end{groupe}%
%% v. 172b
\begin{groupe}%
\personnageFr{Mnésiloque}%
\trimIambFr{\ligne{172}\espace{3.2cm}Comment, au nom des dieux ?\\\passeLigne}%
\end{groupe}%
%% v. 173-174
\begin{groupe}%
\personnageFr{Euripide}%
\trimIambFr{Cesse de japper. Car tel j'étais aussi moi-même\\
à son âge, quand je commençais à produire.\\\passeLigne}%
\end{groupe}%
%% v. 175
\begin{groupe}%
\personnageFr{Mnésiloque}%
\trimIambFr{Par Zeus ! je ne t'envie pas ton éducation.\\\passeLigne}%
\end{groupe}%
%% v. 176a
\begin{groupe}%
\personnageFr{Euripide}%
\trimIambFr{Mais laisse-moi dire ce qui m'amène.\\\passeLigne}%
\end{groupe}%
%% v. 176b
\begin{groupe}%
\personnageFr{Agathon}%
\trimIambFr{\ligne{176}\espace{5.5cm}Parle.\\\passeLigne}%
\end{groupe}%
%% v. 177-180a
\begin{groupe}%
\personnageFr{Euripide}%
\trimIambFr{Agathon, « c'est le propre d'un sage de savoir résumer en bref\\
de longs discours ».\\
Frappé d'une calamité inouïe,\\
je viens à toi en suppliant.\\\passeLigne}%
\end{groupe}%
%% v. 180b
\begin{groupe}%
\personnageFr{Agathon}%
\trimIambFr{\ligne{180}\espace{3.5cm}Que te faut-il ?\\\passeLigne}%
\end{groupe}%
%% v. 181-182
\begin{groupe}%
\personnageFr{Euripide}%
\trimIambFr{Les femmes pensent consommer ma perte, aujourd'hui,\\
\lednopb%
aux Thesmophories, parce que je dis du mal d'elles.\\\passeLigne}%
\end{groupe}%
%% v. 183
\begin{groupe}%
\personnageFr{Agathon}%
\trimIambFr{Et que pouvons-nous pour toi ?\\\passeLigne}%
\end{groupe}%
%% v. 184-187
\begin{groupe}%
\personnageFr{Euripide}%
\trimIambFr{Tout. Si siégeant en tapinois\\
parmi les femmes, en feignant d'être l'une d'elles,\\
tu prends ma défense, tu ne manqueras pas à me sauver.\\
Car seul tu parleras d'une manière digne de moi.\\\passeLigne}%
\end{groupe}%
%% v. 188
\begin{groupe}%
\personnageFr{Agathon}%
\trimIambFr{Alors, que ne vas-tu te défendre en personne ?\\\passeLigne}%
\end{groupe}%
%% v. 189-192
\begin{groupe}%
\personnageFr{Euripide}%
\trimIambFr{Je te le dirai. D'abord on me connaît.\\
Ensuite je suis tout chenu, et j'ai de la barbe.\\
Toi, tu es avenant de figure, blanc, rasé,\\
doué d'une voix de femme, mignon, plaisant à voir.\\\passeLigne}%
\end{groupe}%
%% v. 193a
\begin{groupe}%
\personnageFr{Agathon}%
\trimIambFr{Euripide...\\\passeLigne}%
\end{groupe}%
%% v. 193b
\begin{groupe}%
\personnageFr{Euripide}%
\trimIambFr{\ligne{193}\espace{1.5cm}Eh bien ?\\\passeLigne}%
\end{groupe}%
%% v. 193c-194
\begin{groupe}%
\personnageFr{Agathon}%
\trimIambFr{\ligne{193}\espace{3cm}Tu écrivis un jour : « Tu prends plaisir\\
à voir cette lumière, penses-tu que ton père y prenne moins plaisir ? »\\\passeLigne}%
\end{groupe}%
%% v. 195a
\begin{groupe}%
\personnageFr{Euripide}%
\trimIambFr{Sans doute.\\\passeLigne}%
\end{groupe}%
%% v. 195b-199
\begin{groupe}%
\personnageFr{Agathon}%
\trimIambFr{\ligne{195}\espace{1.2cm}N'espère donc pas que je prenne sur mon compte ton infortune.\\
Il faudrait que je fusse fou.\\
Charge-toi personnellement de tes propres affaires.\\
Car les revers, il convient de les subir à force, non d'artifices,\\
mais d'endurance.\\\passeLigne}%
\end{groupe}%
%% v. 200-201
\begin{groupe}%
\personnageFr{Mnésiloque}%
\trimIambFr{Au fait toi-même, bougre de cul, tu as l'anus béant à force,\\
non de discours, mais d'endurance.\\\passeLigne}%
\end{groupe}%
%% v. 202
\begin{groupe}%
\personnageFr{Euripide}%
\trimIambFr{Mais qu'est-ce qui te fait craindre d'aller là ?\\\passeLigne}%
\end{groupe}%
%% v. 203a
\begin{groupe}%
\personnageFr{Agathon}%
\trimIambFr{Je finirais plus misérablement que toi.\\\passeLigne}%
\end{groupe}%
%% v. 203b
\begin{groupe}%
\personnageFr{Euripide}%
\trimIambFr{\ligne{203}\espace{4cm}Pourquoi ?\\\passeLigne}%
\end{groupe}%
%% v. 203c-205
\begin{groupe}%
\personnageFr{Agathon}%
\trimIambFr{\ligne{203}\espace{5cm}Pourquoi ?\\
Je passe pour ravir les travaux nocturnes des femmes,\\
et dérober la Cypris féminine.\\\passeLigne}%
\end{groupe}%
%% v. 206-207
\begin{groupe}%
\personnageFr{Mnésiloque}%
\trimIambFr{Mon cul, ravir ! Eh par Zeus! dis de te faire enfiler.\\
Mais, par Zeus, la défaite est bonne.%
\commentaire{207}{la défaite}{: prétexte fallacieux pour se défaire de ce qui est importun, embarrassant. (J'ai laissé le mot choisi par Willems.)}\\\passeLigne}%
\end{groupe}%
%% v. 208a
\begin{groupe}%
\personnageFr{Euripide}%
\trimIambFr{Eh bien ? le feras-tu ?\\\passeLigne}%
\end{groupe}%
%% v. 208b
\begin{groupe}%
\personnageFr{Agathon}%
\trimIambFr{\ligne{208}\espace{3.5cm}N'y compte pas.\\\passeLigne}%
\end{groupe}%
%% v. 209
\begin{groupe}%
\personnageFr{Euripide}%
\trimIambFr{Ô triple malheureux, comment périt Euripide.\\\passeLigne}%
\end{groupe}%
%% v. 210
\begin{groupe}%
\personnageFr{Mnésiloque}%
\trimIambFr{Mon ami, mon gendre, n'abandonne point.\\\passeLigne}%
\end{groupe}%
%% v. 211a
\begin{groupe}%
\personnageFr{Euripide}%
\trimIambFr{Que puis-je faire ?\\\passeLigne}%
\end{groupe}%
%% v. 211b-212
\begin{groupe}%
\personnageFr{Mnésiloque}%
\trimIambFr{\ligne{211}\espace{3.4cm}Envoie-le se faire\\
pendre, lui, et dispose de moi comme tu voudras.\\\passeLigne}%
\end{groupe}%
%% v. 213-214a
\begin{groupe}%
\personnageFr{Euripide}%
\trimIambFr{Voyons, puisque tu te remets entre mes mains,\\
ôte ce manteau.\\\passeLigne}%
\end{groupe}%
%% v. 214b-215a
\begin{groupe}%
\personnageFr{Mnésiloque}%
\trimIambFr{\ligne{214}\espace{3.7cm}Le voilà par terre.\\
Mais que penses-tu faire de moi ?\\\passeLigne}%
\end{groupe}%
%% v. 215b-216a
\begin{groupe}%
\personnageFr{Euripide}%
\trimIambFr{\ligne{215}\espace{3.7cm}Te raser cela\\
et te flamber par dessous.\\\passeLigne}%
\end{groupe}%
%% v. 216b-217
\begin{groupe}%
\personnageFr{Mnésiloque}%
\trimIambFr{\ligne{216}\espace{3cm}Eh bien, fais à ta guise ;\\
autrement je n'aurais jamais dû me mettre à ta merci.%\commentaire{217}{}{Cf. sur le texte de ce passage, v. 217, Bull. cité, note 4, ci-après.}.
\\\passeLigne}%
\end{groupe}%
%% v. 218-219a
\begin{groupe}%
\personnageFr{Euripide}%
\trimIambFr{Agathon, pour sûr, tu portes toujours sur toi des rasoirs.\\
Prête-moi donc un peu un rasoir.\\\passeLigne}%
\end{groupe}%
%% v. 219b-220a
\begin{groupe}%
\personnageFr{Agathon}%
\trimIambFr{\ligne{219}\espace{3.8cm}Prends-le toi-même,\\
là, dans l'étui.\\\passeLigne}%
\end{groupe}%
%% v. 220b-221
\begin{groupe}%
\personnageFr{Euripide}%
\trimIambFr{\ligne{220}\espace{4cm}Tu es galant homme. \\
(À Mnésiloque.) Assieds-toi. Enfle la joue droite. (Il le rase.)\\\passeLigne}%
\end{groupe}%
%% v. 222a
\begin{groupe}%
\personnageFr{Mnésiloque}%
\trimIambFr{Aïe !\\\passeLigne}%
\end{groupe}%
%% v. 222b-223a
\begin{groupe}%
\personnageFr{Euripide}%
\trimIambFr{\ligne{222}\espace{1cm}Qu'as-tu à crier ? Je t'enfonce un bâillon\\
si tu ne te tais.\\\passeLigne}%
\end{groupe}%
%% v. 223b
\begin{groupe}%
\personnageFr{Mnésiloque}%
\trimIambFr{\ligne{223}\espace{2.2cm}Aïe, aïe, aïe ! aie, aïe, aïe !\\\passeLigne}%
\end{groupe}%
%% v. 224a
\begin{groupe}%
\personnageFr{Euripide}%
\trimIambFr{Holà ! où cours-tu ?\\\passeLigne}%
\end{groupe}%
%% v. 224b-226a
\begin{groupe}%
\personnageFr{Mnésiloque}%
\trimIambFr{\ligne{224}\espace{2.6cm}Au temple des augustes\\
déesses. Car, par Déméter, je ne resterai pas là à me faire\\
taillader.\\\passeLigne}%
\end{groupe}%
%% v. 226b-227
\begin{groupe}%
\personnageFr{Euripide}%
\trimIambFr{\ligne{226}\espace{1.7cm}Mais tu seras grotesque\\
avec ta moitié de figure dégarnie de poil.\\\passeLigne}%
\end{groupe}%
%% v. 228a
\begin{groupe}%
\personnageFr{Mnésiloque}%
\trimIambFr{Peu m'en chaut.\\\passeLigne}%
\end{groupe}%
%% v. 228b-229a
\begin{groupe}%
\personnageFr{Euripide}%
\trimIambFr{\ligne{228}\espace{2.7cm}Au nom des dieux,\\
ne m'abandonne pas. Arrive ici.\\\passeLigne}%
\end{groupe}%
%% v. 229b
\begin{groupe}%
\personnageFr{Mnésiloque}%
\trimIambFr{\ligne{229}\espace{3.6cm}Misère de moi !\\\passeLigne}%
\end{groupe}%
%% v. 230
\begin{groupe}%
\personnageFr{Euripide}%
\trimIambFr{Tiens-toi coi et lève la tête. Où te tournes-tu?\\\passeLigne}%
\end{groupe}%
%% v. 231a
\begin{groupe}%
\personnageFr{Mnésiloque}%
\trimIambFr{Mumû.\\\passeLigne}%
\end{groupe}%
%% v. 231b
\begin{groupe}%
\personnageFr{Euripide}%
\trimIambFr{\ligne{231}\espace{1.1cm}Pourquoi mumû ? Tout est fini, et selon les règles.\\\passeLigne}%
\end{groupe}%
%% v. 232
\begin{groupe}%
\personnageFr{Mnésiloque}%
\trimIambFr{Ah malheur ! Je vais resservir dans l'infanterie légère.\\\passeLigne}%
\end{groupe}%
%% v. 233-234a
\begin{groupe}%
\personnageFr{Euripide}%
\trimIambFr{Ne t'inquiète pas. Tu auras un air tout avenant.\\
Veux-tu te mirer ? (Il lui tend un miroir.)\\\passeLigne}%
\end{groupe}%
%% v. 234b
\begin{groupe}%
\personnageFr{Mnésiloque}%
\trimIambFr{\ligne{234}\espace{3.5cm}Si c'est ton plaisir. Passe.\\\passeLigne}%
\end{groupe}%
%% v. 235a
\begin{groupe}%
\personnageFr{Euripide}%
\trimIambFr{Te vois-tu ?\\\passeLigne}%
\end{groupe}%
%% v. 235b
\begin{groupe}%
\personnageFr{Mnésiloque}%
\trimIambFr{\ligne{235}\espace{2.4cm}Non, pas Zeus. Je vois Clisthène.\\\passeLigne}%
\end{groupe}%
%% v. 236
\begin{groupe}%
\personnageFr{Euripide}%
\trimIambFr{Lève-toi, que je te flambe. Tiens-toi penché.\\\passeLigne}%
\end{groupe}%
%% v. 237
\begin{groupe}%
\personnageFr{Mnésiloque}%
\trimIambFr{Ah malheur ! Je vais être cochon de lait.\\\passeLigne}%
\end{groupe}%
%% v. 238-239
\begin{groupe}%
\personnageFr{Euripide}%
\trimIambFr{Qu'on m'apporte de dedans une torche ou une lampe.\\
Courbe-toi, et gare au bout de la queue.\\\passeLigne}%
\end{groupe}%
%% v. 240-242
\begin{groupe}%
\personnageFr{Mnésiloque}%
\trimIambFr{J'y veillerai, par Zeus... N'empêche que je brûle.\\
Ah misère ! de l'eau, de l'eau, voisins,\\
avant que la fesse d'à côté ne s'enflamme.\\\passeLigne}%
\end{groupe}%
%% v. 243a
\begin{groupe}%
\personnageFr{Euripide}%
\trimIambFr{Du sang froid.\\\passeLigne}%
\end{groupe}%
%% v. 243b
\begin{groupe}%
\personnageFr{Mnésiloque}%
\trimIambFr{\ligne{243}\espace{1.4cm}Du sang froid ? quand je suis dévoré par les flammes !..\\\passeLigne}%
\end{groupe}%
%% v. 244-245a
\begin{groupe}%
\personnageFr{Euripide}%
\trimIambFr{Mais tu n'as plus à t'inquiéter. Car te voilà dépétré\\
du principal.\\\passeLigne}%
\end{groupe}%
%% v. 245b-246
\begin{groupe}%
\personnageFr{Mnésiloque}%
\trimIambFr{\ligne{245}\espace{2.3cm}Phu ! Ah! cette suie !\\
Je suis tout feu dans l'entrefesson.\\\passeLigne}%
\end{groupe}%
%% v. 247
\begin{groupe}%
\personnageFr{Euripide}%
\trimIambFr{Ne t'inquiète pas. Le voisin d'à côté t'épongera cela.\\\passeLigne}%
\end{groupe}%
%% v. 248
\begin{groupe}%
\personnageFr{Mnésiloque}%
\trimIambFr{Il en coûtera à celui qui s'attaquera à mon derrière.\\\passeLigne}%
\end{groupe}%
%% v. 249-251
\begin{groupe}%
\personnageFr{Euripide}%
\trimIambFr{Agathon, puisque tu refuses de te dévouer à moi,\\
prête-nous du moins un vêtement\\
et un strophion. Tu ne nieras pas que tu en possèdes..\\\passeLigne}%
\end{groupe}%
%% v. 252a
\begin{groupe}%
\personnageFr{Agathon}%
\trimIambFr{Prenez et usez-en ; je ne m'y oppose pas.\\\passeLigne}%
\end{groupe}%
%% v. 252b
\begin{groupe}%
\personnageFr{Mnésiloque}%
\trimIambFr{\ligne{252}\espace{5cm}Que faut-il prendre ?\\\passeLigne}%
\end{groupe}%
%% v. 253
\begin{groupe}%
\personnageFr{Euripide}%
\trimIambFr{Quoi ? prends d'abord cette crocote et revêts-la.\\\passeLigne}%
\end{groupe}%
%%% v. 254
\begin{groupe}%
\personnageFr{Mnésiloque}%
\trimIambFr{Par Aphrodite, elle fleure bon le priape.\\\passeLigne}%
\end{groupe}%
%% v. 255a
\begin{groupe}%
\personnageFr{Agathon}%
\trimIambFr{Serre-t'en vite la taille.\\\passeLigne}%
\end{groupe}%
%% v. 255b
\begin{groupe}%
\personnageFr{Euripide}%
\trimIambFr{\ligne{255}\espace{2.5cm}Maintenant passe-nous un strophion.\\\passeLigne}%
\end{groupe}%
%% v. 255c
\begin{groupe}%
\personnageFr{Agathon}%
\trimIambFr{\ligne{255}\espace{5.5cm}Là.\\\passeLigne}%
\end{groupe}%
%% v. 256
\begin{groupe}%
\personnageFr{Mnésiloque}%
\trimIambFr{Maintenant arrange-moi les jambes.\\\passeLigne}%
\end{groupe}%
%% v. 257a
\begin{groupe}%
\personnageFr{Euripide}%
\trimIambFr{Une résille et un bandeau ?\\\passeLigne}%
\end{groupe}%
%% v. 257b-258
\begin{groupe}%
\personnageFr{Agathon}%
\trimIambFr{\ligne{257}\espace{4.2cm}Plutôt cette\\
perruque, que je porte la nuit.\\\passeLigne}%
\end{groupe}%
%% v. 259
\begin{groupe}%
\personnageFr{Euripide}%
\trimIambFr{Eh ! par Zeus, c'est même tout juste l'affaire.\\\passeLigne}%
\end{groupe}%
%%% v. 260a
\begin{groupe}%
\personnageFr{Mnésiloque}%
\trimIambFr{M'ira-t-elle ?\\\passeLigne}%
\end{groupe}%
%%% v. 260b-261a
\begin{groupe}%
\personnageFr{Euripide}%
\trimIambFr{\ligne{261}\espace{2.7cm}Mais c'est parfait, par Zeus.\\
Un encycle ?\\\passeLigne}%
\end{groupe}%
%%% v. 261b
\begin{groupe}%
\personnageFr{Agathon}%
\trimIambFr{\ligne{261}\espace{2.4cm}Prends celui-là, sur la couchette.\\\passeLigne}%
\end{groupe}%
%%% v. 262a
\begin{groupe}%
\personnageFr{Euripide}%
\trimIambFr{Des chaussures ?\\\passeLigne}%
\end{groupe}%
%%% v. 262b
\begin{groupe}%
\personnageFr{Agathon}%
\trimIambFr{\ligne{262}\espace{2.6cm}Voici les miennes, prends-les.\\\passeLigne}%
\end{groupe}%
%% v. 263
\begin{groupe}%
\personnageFr{Mnésiloque}%
\trimIambFr{Est-ce qu'elles m'iront ? Tu n'aimes pas à te chausser large\\\passeLigne}%
\end{groupe}%
%% v. 264-265
\begin{groupe}%
\personnageFr{Agathon}%
\trimIambFr{À toi de t'en assurer. Or çà, tu as ce qu'il te faut.\\
Qu'on me roule au plus vite à l'intérieur.\\\passeLigne}%
\end{groupe}%
%% v. 266-268a
\begin{groupe}%
\personnageFr{Euripide}%
\trimIambFr{Voilà notre homme devenu femme,\\
du moins en apparence. Si tu parles, tâche par ton timbre de voix\\
de simuler la femme, avec aisance et naturel.\\\passeLigne}%
\end{groupe}%
%% v. 268b
\begin{groupe}%
\personnageFr{Mnésiloque}%
\trimIambFr{\ligne{268}\espace{4.1cm}(Voix de fausset.) J'y tâcherai.\\\passeLigne}%
\end{groupe}%
%% v. 269a
\begin{groupe}%
\personnageFr{Euripide}%
\trimIambFr{Va maintenant.\\\passeLigne}%
\end{groupe}%
%% v. 269b-270a
\begin{groupe}%
\personnageFr{Mnésiloque}%
\trimIambFr{\ligne{269}\espace{2.3cm}Non, par Apollon, à moins que\\
tu ne me jures...\\\passeLigne}%
\end{groupe}%
%% v. 270b
\begin{groupe}%
\personnageFr{Euripide}%
\trimIambFr{\ligne{270}\espace{2.3cm}Quoi ?\\\passeLigne}%
\end{groupe}%
%% v. 270c-271
\begin{groupe}%
\personnageFr{Mnésiloque}%
\trimIambFr{\ligne{270}\espace{4cm}D'aider à me sauver\\
par toutes tes rubriques, s'il me tombe dessus quelque malheur.\\\passeLigne}%
\end{groupe}%
%%% v. 272
\begin{groupe}%
\personnageFr{Euripide}%
\trimIambFr{Eh bien, « je jure, par l'Éther, la demeure de Zeus ».\\\passeLigne}%
\end{groupe}%
%% v. 273
\begin{groupe}%
\personnageFr{Mnésiloque}%
\trimIambFr{Pourquoi pas aussi bien par le logis de famille d'Hippocratès ?\\\passeLigne}%
\end{groupe}%
%% v. 274
\begin{groupe}%
\personnageFr{Euripide}%
\trimIambFr{Eh bien! je jure par tous les dieux en bloc.\\\passeLigne}%
\end{groupe}%
%% v. 275-276
\begin{groupe}%
\personnageFr{Mnésiloque}%
\trimIambFr{Au surplus, souviens-toi de ceci : c'est ton cœur qui a juré,\\
et non ta langue ; je ne la tiens pas pour liée.\\\passeLigne}%
\end{groupe}%
%% v. 277-279a
\begin{groupe}%
\personnageFr{Euripide}%
\trimIambFr{Hâte-toi de partir. Car le signal de l'assemblée\\
paraît sur le Thesmophorion. Moi, je m'en vais.\\
\ligne{278}(Tous sortent. Le décor du fond change et représente le Thesmophorion.\\
Entrent Mnésiloque suivi d'une servante portant une ciste.)\\\passeLigne}%
\end{groupe}%
%% v. 279b-294
\begin{groupe}%
\personnageFr{Mnésiloque}%
\trimIambFr{\ligne{279}\espace{2.3cm}Par ici, Thratta, suis-moi.\\
Vois, Thratta, sous la lueur fumeuse des torches ardentes,\\
vois la foule qui monte.\\
— Thesmophores, merveilles de beauté, accueillez-moi\\
sous d'heureux auspices, ici et à mon retour au logis.\\\passeLigne}
\end{groupe}%
%% v. 279b-294
\begin{groupe}%
\trimIambFr{— Thratta, dépose la ciste, et tires-en\\
les gâteaux, pour que je l'offre aux deux déesses.\\
— Maîtresse très honorée, chère Déméter,\\
et toi Perséphone, puissé-je mainte et mainte fois\\
être à même de te sacrifier. (À part :) À tout le moins, n'être pas découvert aujourd'hui.
(Fort :) Que ma porcelette trouve un mari riche\\
fortuné, et au surplus benêt et stupide !\\
Qu'elle m'ait esprit et intelligence quant à la cassette.\\
— Où donc, où m'asseoir en bonne place pour entendre les orateurs ?\\
— Toi, Thratta, retire-toi à l'écart,\\
car il n'est pas permis aux esclaves d'ouïr ce qui se dit ici. \\\passeLigne}%
\end{groupe}%
%% v. 295-311
\begin{groupe}%
\personnageFr{Coryphée}%
\proseFr{(La prêtresse :) Recueillez-vous, recueillez-vous.\\
Priez les Thesmophores,\\
et Ploutos et Calligénie,\\
et la Nourricière,\\
et Hermès et les Charites,\\
de rendre cette ecclésie et présente réunion\\
la meilleure et la plus avantageuse qu'il se pourra,\\
pour le plus grand profit de la cité des Athéniens\\
et pour notre propre bien.\\
Que celle qui aura agi et parlé\\
le mieux en faveur du peuple des Athéniens\\
et du peuple des femmes voie prévaloir ses avis.\\
Demandez aux dieux cette faveur, et pour vous-mêmes toutes les prospérités.\\
Io Péan ! Io Péan ! Soyons en Joie.\\\passeLigne}%
\end{groupe}%
%% v. 312-330
\begin{groupe}%
\personnageFr{Chœur}%
\chantChoeurFr{\ligne{312}Nous acceptons ! Et nous supplions}%
\end{groupe}%
%% v. 312-330
\begin{groupe}%
\chantChoeurFr{la gent des dieux d'apparaître}%
\end{groupe}%
%% v. 312-330
\begin{groupe}%
\chantChoeurFr{et de prendre en gré ces prières.}%
\end{groupe}%
%% v. 312-330
\begin{groupe}%
\chantChoeurFr{Zeus au grand renom, et toi, dieu à la lyre d'or,}%
\end{groupe}%
%% v. 312-330
\begin{groupe}%
\chantChoeurFr{qui tiens Délos la sainte,}%
\end{groupe}%
%% v. 312-330
\begin{groupe}%
\chantChoeurFr{et toi, vierge toute-puissante,}%
\end{groupe}%
%% v. 312-330
\begin{groupe}%
\chantChoeurFr{aux yeux glauques, à la lance d'or,}%
\end{groupe}%
%% v. 312-330
\begin{groupe}%
\chantChoeurFr{qui habites une cité enviée, viens ici.}%
\end{groupe}%
%% v. 312-330
\begin{groupe}%
\chantChoeurFr{Et toi aussi, vierge aux noms multiples, meurtrière de bêtes, rejeton}%
\end{groupe}%
%% v. 312-330
\begin{groupe}%
\chantChoeurFr{de Léto aux yeux d'or.}%
\end{groupe}%
%% v. 312-330
\begin{groupe}%
\chantChoeurFr{Et toi, auguste Poseidon, thalassien,}%
\end{groupe}%
%% v. 312-330
\begin{groupe}%
\chantChoeurFr{modérateur de la mer, }%
\end{groupe}%
%% v. 312-330
\begin{groupe}%
\chantChoeurFr{quitte le gouffre poissonneux }%
\end{groupe}%
%% v. 312-330
\begin{groupe}%
\chantChoeurFr{fouetté des vents.}%
\end{groupe}%
%% v. 312-330
\begin{groupe}%
\chantChoeurFr{Et vous, filles de Nérée le marin,}%
\end{groupe}%
%% v. 312-330
\begin{groupe}%
\chantChoeurFr{et vous, Nymphes des montagnes.}%
\end{groupe}%
%% v. 312-330
\begin{groupe}%
\chantChoeurFr{Que la phorminx d'or}%
\end{groupe}%
%% v. 312-330
\begin{groupe}%
\chantChoeurFr{accompagne nos prières,}%
\end{groupe}%
%% v. 312-330
\begin{groupe}%
\chantChoeurFr{\ligne{329}et puissions-nous mener à bonne fin notre ecclésie,}%
\end{groupe}%
%% v. 312-330
\begin{groupe}%
\chantChoeurFr{femmes de pure lignée athénienne.\\\passeLigne}%
\end{groupe}%
%% v. 331-351
\begin{groupe}%
\personnageFr{Coryphée}%
\trimIambFr{\ligne{331}Invoquez les dieux et les déesses de l'Olympe,\\
les dieux et les déesses de Pytho,\\
les dieux et les déesses de Délos\\
et les autres dieux.\\
Si quelqu'un forme quelque dessein contre le peuple\\
des femmes, fait des ouvertures\\
à Euripide et aux Mèdes au détriment\\
des femmes, vise à être tyran\\
ou à ramener le tyran,\\
dénonce une femme qui a supposé un enfant; si une esclave\\
ayant servi d'entremetteuse l'est allé chuchoter au mari,\\
ou si chargée d'un message, elle fait un rapport mensonger ;}
\end{groupe}%
%% v. 331-351
\begin{groupe}%
\trimIambFr{si un galant suborne par de belles paroles,\\
et ne tient pas les promesses qu'il a faites ;\\
si une vieille fait des présents à un galant,\\
ou en reçoit en livrant à une courtisane son galant ;\\
si un débitant ou une débitante\\
falsifie la mesure légale des chous ou des cotyles,\\
souhaitez qu'il périsse misérablement avec sa maison ;\\
mais à vous autres toutes priez les dieux\\
d'accorder des biens en abondance.\\\passeLigne}%
\end{groupe}%
%% v. 352-354
\begin{groupe}%
\personnageFr{Chœur}%
\chantChoeurFr{Nous unissons nos prières aux tiennes}%
\end{groupe}%
%% v. 352-354
\begin{groupe}%
\chantChoeurFr{pour le bien de la cité, pour le bien du peuple}%
\end{groupe}%
%% v. 352-354
\begin{groupe}%
\chantChoeurFr{afin que ces vœux s'accomplissent,}%
\end{groupe}%
%% v. 355-357
\begin{groupe}%
\chantChoeurFr{et que la victoire demeure, comme de raison,}%
\end{groupe}%
%% v. 352-354
\begin{groupe}%
\chantChoeurFr{à celles qui auront le mieux parlé. Mais toutes celles}%
\end{groupe}%
%% v. 352-354
\begin{groupe}%
\chantChoeurFr{qui en imposent ou transgressent}%
\end{groupe}%
%% v. 358-361
\begin{groupe}%
\chantChoeurFr{les serments établis,}%
\end{groupe}%
%% v. 352-354
\begin{groupe}%
\chantChoeurFr{dans leur propre intérêt, mais au détriment commun,}%
\end{groupe}%
%% v. 352-354
\begin{groupe}%
\chantChoeurFr{celles qui cherchent à bouleverser}%
\end{groupe}%
%% v. 362-364
\begin{groupe}%
\chantChoeurFr{les décrets et la loi,}%
\end{groupe}%
%% v. 352-354
\begin{groupe}%
\chantChoeurFr{révèlent}%
\end{groupe}%
%% v. 352-354
\begin{groupe}%
\chantChoeurFr{nos secrets à nos ennemis,}%
\end{groupe}%
%% v. 365-367
\begin{groupe}%
\chantChoeurFr{ou se positionnent vis-à-vis des Mèdes pour asseoir notre empire...}%
\end{groupe}%
%% v. 352-354
\begin{groupe}%
\chantChoeurFr{et empirent notre position}%
\end{groupe}%
%% v. 352-354
\begin{groupe}%
\chantChoeurFr{celles-là sont coupables de lèse-majesté divine et de lèse-nation.}%
\end{groupe}%
%% v. 368-371
\begin{groupe}%
\chantChoeurFr{Ô tout-puissant}%
\end{groupe}%
%% v. 352-354
\begin{groupe}%
\lednopb%
\chantChoeurFr{Zeus, ratifie ces vœux et fais\\\passeLigne}%
\end{groupe}%
%% v. 352-354
\begin{groupe}%
\chantChoeurFr{que les dieux nous assistent,}%
\end{groupe}%
%% v. 352-354
\begin{groupe}%
\chantChoeurFr{quoique nous soyons femmes.\\\passeLigne}%
\end{groupe}%
%% v. 372-379
\begin{groupe}%
\lednopb%
\personnageFr{Coryphée}%
\trimIambFr{\ligne{373}Écoutez toutes. La Boulé\\
des femmes, Timocléa épistate,\\
Lysilla greffière, Sostraté auteur de la motion, a décrété ce qui suit :\\
Une assemblée aura lieu le matin du jour intermédiaire\\
des Thesmophories, celui où nous sommes le plus de loisir ;\\
et on délibérera d'abord sur Euripide\\
et la peine qu'il doit subir. Car il est tenu coupable\\
par nous toutes. Qui veut prendre la parole ?\\\passeLigne}%
\end{groupe}%
%% v. 380a
\begin{groupe}%
\personnageFr{Mika}%
\trimIambFr{Moi.\\\passeLigne}%
\end{groupe}%
%% v. 380b-382
\begin{groupe}%
\personnageFr{Coryphée}%
\trimIambFr{\ligne{380}\espace{1cm}Alors commence par ceindre cette couronne, avant que parler.}%
\end{groupe}%
\begin{groupe}%
\tetramIambFr{Paix ! Silence! Fais attention ! Car voici qu'elle tousse
\lednopb%
\\
\lednopb%
comme font les orateurs. Elle paraît en avoir long à dire.\\\passeLigne}%
\end{groupe}%
%% v. 383-432
\begin{groupe}%
\personnageFr{Mika}%
\trimIambFr{Aucun motif d'ambition, non par les deux déesses,\\
ne m'a fait lever pour prendre la parole, ô femmes. Mais\\
je suis excédée, depuis longtemps, hélas !\\
de nous voir traînées dans la boue par \\
Euripide, le fils de la vendeuse d'herbes,\\
et constamment en butte à des injures de tout genre.\\
Car quelle est l'injure dont il ne nous salisse ?\\
Où s'est-il fait faute de nous calomnier, pour peu simplement\\
qu'il y eût des spectateurs, des tragédiens et des choeurs ?\\
nous traitant de galantes, passionnées d'hommes,\\
sujettes au vin, traitresses, bavardes,\\
rien qui vaille, grand fléau des hommes.\\
Si bien qu'à peine rentrés du théâtre,\\
ils nous soupçonnent, et guettent aussitôt\\
s'il n'est pas quelque amant caché.\\
Et plus rien ne nous est permis de ce que nous faisions autrefois,\\
tant il a enseigné de vilaines choses\\\passeLigne}
\end{groupe}%
%% v. 383-432
\begin{groupe}%
\trimIambFr{à nos époux. Une femme tresse-t-elle\\
une couronne, on la croit amoureuse ; laisse-t-elle tomber\\
un ustensile en allant et venant par la maison,\\
le mari demande : « Pour qui brise-t-on la marmite ?\\
Nul doute que ce ne soit pour l'étranger de Corinthe ».\\
Une fille est-elle souffrante, aussitôt son frère de dire :\\
« Ce teint ne me plaît guère chez la fillette ».\\
Mais passons. Une femme privée d'enfants veut\\
en supposer un ; même cela ne se peut celer,\\
car les maris se tiennent tout proche.\\
Et auprès des vieillards, qui auparavant épousaient les jeunesses,\\
il nous a calomniées, tellement que pas un homme d'âge\\
ne consent à prendre femme, en raison de ce vers :\\
« Pour un époux âgé, la femme est un tyran ».\\
En outre, à cause de lui, sur les appartements de femmes,\\
ils appliquent des cachets et des barres\\\passeLigne}
\end{groupe}%
%% v. 383-432
\begin{groupe}%
\trimIambFr{pour nous garder, et par là-dessus ils entretiennent des chiens molosses,\\
épouvantails pour les galants.\\
Encore lui passerait-on cela : Mais ce qu'on nous permettait autrefois,\\
de diriger nous-mêmes le ménage et de prendre dans le cellier\\
de la farine, de l'huile, du vin, cela même ne nous est plus\\
possible : Car les maris portent sur eux de petites clefs\\
secrètes, les plus malicieuses du monde,\\
je ne sais quelles laconiennes à trois dents.\\
Autrefois du moins il tenait à nous d'ouvrir furtivement la porte,\\
en nous faisant graver un anneau de trois oboles.\\
Aujourd'hui ce tâte-poule d'Euripide\\
leur a enseigné l'emploi de cachets vermoulus\\
qu'ils portent suspendus sur eux. Maintenant donc j'opine\\
que nous combinions de manière ou d'autre un procédé pour le détruire,\\
soit par des drogues, soit par un moyen quelconque,\\
\lednopb%
afin qu'il meure. Ceci, je le dis à pleine bouche.\\
\lednopb%
Le reste, je le coucherai par écrit, avec l'aide de la greffière.\\\passeLigne}%
\end{groupe}%
%% v. 433-442
\begin{groupe}%
\personnageFr{Chœur}%
\chantChoeurFr{Jamais je n'entendis}%
\end{groupe}%
%% v. 433-442
\begin{groupe}%
\chantChoeurFr{femme plus féconde en retours}%
\end{groupe}%
%% v. 433-442
\begin{groupe}%
\chantChoeurFr{ni qui sache mieux parler.}%
\end{groupe}%
%% v. 433-442
\begin{groupe}%
\chantChoeurFr{Tout ce qu'elle dit est juste.}%
\end{groupe}%
%% v. 433-442
\begin{groupe}%
\chantChoeurFr{Elle a tout envisagé sous tous égards,}%
\end{groupe}%
%% v. 433-442
\begin{groupe}%
\chantChoeurFr{tout soupesé dans son esprit : et avec discrétion,}%
\end{groupe}%
%% v. 433-442
\begin{groupe}%
\chantChoeurFr{elle a trouvé un heureux}%
\end{groupe}%
%% v. 433-442
\begin{groupe}%
\chantChoeurFr{choix d'adroites raisons.}%
\end{groupe}%
%% v. 433-442
\begin{groupe}%
\chantChoeurFr{À tel point que si}%
\end{groupe}%
%% v. 433-442
\begin{groupe}%
\chantChoeurFr{Xénoclès, fils de Carkinos, parlait tout de suite après elle,}%
\end{groupe}%
%% v. 433-442
\begin{groupe}%
\chantChoeurFr{il vous paraîtrait je pense,}%
\end{groupe}%
%% v. 433-442
\begin{groupe}%
\ligne{443}\chantChoeurFr{à tous}%
\end{groupe}%
%% v. 433-442
\begin{groupe}%
\chantChoeurFr{parler sans rien dire.\\\passeLigne}%
\end{groupe}%
%% v. 443-458
\begin{groupe}%
\personnageFr{Marchande de couronnes}%
\trimIambFr{\ligne{443}Ce n'est que pour ajouter quelques mots que je m'avance à mon tour.\\
Les autres chefs d'accusation ont été traités à souhait.\\
Mais ce que j'ai souffert personnellement, je tiens à le dire.\\
Mon mari mourut alors que j'étais enceinte,\\
me laissant cinq petits enfants, que je nourrissais à grand'peine\\
\lednopb%
en tressant des couronnes sur le marché aux myrtes.\\
En attendant, je gagnais ma vie vaille que vaille.\\
Mais aujourd'hui cet homme, avec ses tragédies,\\
\lednopb%
a persuadé aux gens qu'il n'y a pas de dieux ;%
\\\passeLigne}
\end{groupe}%
%% v. 443-458
\begin{groupe}%
\trimIambFr{de sorte que notre commerce est réduit environ de moitié.\\
C'est pourquoi je vous recommande à toutes et vous requiers\\
de le punir, cet homme, pour quantité de raisons.\\
Car il se porte contre nous, ô femmes, à de sauvages violences,\\
élevé qu'il fut parmi les herbes sauvages.\\
Mais je m'en vas à l'agora.\\
Car j'ai vingt couronnes à tresser à prix convenu.\\\passeLigne}%
\end{groupe}%
%% v. 459-465
\begin{groupe}%
\personnageFr{Chœur}%
\chantChoeurFr{Cette autre a montré un tour d'esprit}%
\end{groupe}%
%% v. 459-465
\begin{groupe}%
\chantChoeurFr{plus piquant encore}%
\end{groupe}%
%% v. 459-465
\begin{groupe}%
\chantChoeurFr{que la première.}%
\end{groupe}%
%% v. 459-465
\begin{groupe}%
\chantChoeurFr{Comme elle a débité,}%
\end{groupe}%
%% v. 459-465
\begin{groupe}%
\chantChoeurFr{non du hors-d'oœvre — elle a du sens}%
\end{groupe}%
%% v. 459-465
\begin{groupe}%
\chantChoeurFr{et un esprit délié — ni}%
\end{groupe}%
%% v. 459-465
\begin{groupe}%
\chantChoeurFr{de l'abstrus, mais toutes choses emportant conviction !}%
\end{groupe}%
%% v. 459-465
\begin{groupe}%
\chantChoeurFr{Oui, il faut que de cet outrage}%
\end{groupe}%
%% v. 459-465
\begin{groupe}%
\chantChoeurFr{l'auteur}%
\end{groupe}%
%% v. 459-465
\begin{groupe}%
\chantChoeurFr{subisse un châtiment exemplaire.\\\passeLigne}%
\end{groupe}%
%% v. 466-519
\begin{groupe}%
\personnageFr{Mnésiloque}%
\trimIambFr{\ligne{466}Que vous soyez, ô femmes, vivement irritées\\
contre Euripide, qui si mal parla de vous,\\
n'est pas pour étonner, ni que votre bile déborde.\\
Car moi-même, — ainsi puissé-je être heureuse dans mes enfants, —\\
je hais cet homme, ou je serais folle.\\
Néanmoins il convient que nous nous expliquions entre nous,\\
car nous sommes seules, et rien ne sera rapporté de ce qui se dit ici.\\
Pourquoi allons-nous l'accusant\\
et sommes-nous indignées de ce que, sachant deux ou trois\\
de nos méfaits, il les a divulgués, quand nous en commettons d'innombrables ?\\
Moi-même toute la première, pour ne rien dire des autres,\\
j'en ai sur la conscience un grand nombre, mais voici\\
le plus fort. J'étais mariée de trois jours,\\\passeLigne}
\end{groupe}%
%% v. 466-519
\begin{groupe}%
\trimIambFr{et mon mari dormait à mes côtés. Or j'avais un amant\\
qui m'avait déflorée à sept ans.\\
Pressé de désirs, il s'en vient gratter à la porte.\\
Je comprends sur-le-champ et descends à petit bruit.\\
Mon mari m'interroge : « Où descends-tu? — Où ?\\
Je suis prise de tranchées, mon homme, et j'ai mal.\\
Aussi vais-je aux latrines. — Va donc ».\\
Et là-dessus il se met à broyer des grains de genévrier, de l'aneth, de la sauge.\\
Pour moi, ayant versé de l'eau sur les gonds,\\
j'allai trouver mon galant et subir son assaut\\
à côté de l'Aguieus, courbant l'échine et me tenant au laurier.\\
Cela au grand jamais, voyez-vous, Euripide n'en a parlé.\\
Non plus qu'il ne dit qu'à défaut d'autres,\\
\lednopb%
nous nous faisons hocher par nos esclaves et par des muletiers ;\\\passeLigne}
\end{groupe}%
%% v. 466-519
\begin{groupe}%
\trimIambFr{ni qu'après avoir paillardé à cœur joie\\
la nuit, nous mâchons de l'ail le matin,\\
afin que saisi par l'odeur, notre mari, à son retour du rempart,\\
n'ait aucun soupçon de notre conduite. Cela, voyez-vous,\\
il n'en a jamais parlé. Que s'il invective Phèdre,\\
que nous importe ? Il n'a pas raconté non plus le cas\\
de cette femme qui, faisant voir son encycle à son mari,\\
\ligne{501}fit partir son galant ; non, il ne l'a pas raconté.\\
Et j'en sais une autre qui prétendit être en travail\\
dix jours durant, jusqu'à ce qu'elle eût acheté un enfant.\\\passeLigne}
\end{groupe}%
%% v. 520-530
\begin{groupe}%
\trimIambFr{Le mari courait à la ronde faire emplette de drogues pour hâter la délivrance.\\
L'enfant fut apporté par une vieille dans une marmite,\\
la bouche bourrée d'un rayon de miel pour qu'il ne criât point.\\
Puis, sur un signe de la porteuse, la femme aussitôt s'écria :\\
« Va-t'en, va-t'en, mon homme, je crois que je vais\\
accoucher ». Car il avait donné des pieds contre la panse de la marmite.\\
Tandis qu'il s'éloigne tout joyeux, elle débarrasse\\
la bouche de l'enfant qui se met à crier.\\
La vieille canaille qui l'a apporté\\
court en souriant à l'époux, et lui dit :\\
« un lion, un lion t'est né, ton portrait moulé ;\\
outre tout le reste absolument, son membre\\
est pareil au tien et tortu comme une pigne ».\\
Ne sont-ce pas là de nos tours? Si fait, par Artémis,\\
c'en sont. Et nous nous fâchons contre Euripide,\\
« sans avoir enduré par-delà nos mérites ».;\\\passeLigne}%
\end{groupe}%
%% v. 520-530
\begin{groupe}%
\personnageFr{Chœur}%
\chantChoeurFr{Une chose me confond :}%
\end{groupe}%
%% v. 520-530
\begin{groupe}%
\chantChoeurFr{où a-t-on trouvé cette créature ?}%
\end{groupe}%
%% v. 520-530
\begin{groupe}%
\chantChoeurFr{Quel pays a formé}%
\end{groupe}%
%% v. 520-530
\begin{groupe}%
\chantChoeurFr{une telle effrontée ?}%
\end{groupe}%
%% v. 520-530
\begin{groupe}%
\chantChoeurFr{La coquine ! tenir ce langage}%
\end{groupe}%
%% v. 520-530
\begin{groupe}%
\chantChoeurFr{au grand jour, avec cette impudence ?}%
\end{groupe}%
%% v. 520-530
\begin{groupe}%
\chantChoeurFr{je n'aurais pas cru que parmi nous}%
\end{groupe}%
%% v. 520-530
\begin{groupe}%
\chantChoeurFr{elle l'eût jamais osé.}%
\end{groupe}%
%% v. 520-530
\begin{groupe}%
\chantChoeurFr{Mais maintenant tout arrive,}%
\end{groupe}%
%% v. 520-530
\begin{groupe}%
\chantChoeurFr{et j'approuve l'antique proverbe,}%
\end{groupe}%
%% v. 520-530
\begin{groupe}%
\chantChoeurFr{quasiment sous chaque pierre,%\commentaire{}{}{Euripide fait allusion à ce proverbe, Héraclès, 1002 : « Il n'y a pas de pierre sous laquelle ne se glisse un scorpion ».}
}%
\end{groupe}%
%% v. 520-530
\begin{groupe}%
\chantChoeurFr{il faut regarder}%
\end{groupe}%
%% v. 520-530
\begin{groupe}%
\chantChoeurFr{de crainte d'être mordu... par un orateur..\\\passeLigne}%
\end{groupe}%
%% v. 531-532
\begin{groupe}%
\personnageFr{Coryphée}%
\tetramIambCatFr{\ligne{531}Mais les femmes au naturel impudent,\\
\lednopb%
il n'est rien de pire à tous égards, si ce n'est des femmes.%\commentaire{531-2}{}{Cf. sur ce passage : Bull. cité, note 15, ci-après.}.
\\\passeLigne}%
\end{groupe}%
%% v. 533-539
\begin{groupe}%
\personnageFr{Mika}%
\tetramIambCatFr{Non, par Aglaure, ô femmes, vous n'êtes pas dans votre bon sens,\\
mais vous êtes maléficiées, ou vous avez pris quelque autre mal,\\
pour que vous laissiez cette peste nous braver toutes à outrance.\\
(Elle se tourne vers le bureau de l'assemblée.) Ainsi qu'on se décide,...\\
sinon nous-mêmes et nos esclaves\\
irons chercher quelque part de la cendre,\\
et lui épilerons le bas-ventre, afin que, femme, elle apprenne\\
à ne plus médire des femmes à l'avenir.\\\passeLigne}%
\end{groupe}%
%% v. 540-543
\begin{groupe}%
\personnageFr{Mnésiloque}%
\tetramIambCatFr{Le bas-ventre ? Non pas, ô femmes. Car la franchise\\
étant de droit et nous toutes citoyennes ci-présentes ayant congé de parler,\\
si j'ai déduit les raisons que je savais en faveur d'Euripide,\\
est-ce un motif pour m'épiler, afin de m'en faire porter la peine ?\\\passeLigne}%
\end{groupe}%
%% v. 544-548
\begin{groupe}%
\personnageFr{Mika}%
\tetramIambCatFr{Quoi! tu n'en porterais pas la peine ? toi qui seule osas\\
répliquer en faveur d'un homme qui nous a tant malmenées,\\
trouvant à plaisir des sujets où une femme se montra mauvaise,\\
produisant des Mélanippes, des Phèdres ; mais jamais il ne produisit une Pénélope,\\
vu qu'elle passait pour sage.\\\passeLigne}%
\end{groupe}%
%% v. 549-550
\begin{groupe}%
\personnageFr{Mnésiloque}%
\tetramIambCatFr{\ligne{549}J'en sais bien la raison. Tu n'en citerais pas une seule,\\
parmi les femmes d'aujourd'hui, de Pénélope. Toutes Phèdres, sans exception.\\\passeLigne}%
\end{groupe}%
%% v. 551-552a
\begin{groupe}%
\personnageFr{Mika}%
\tetramIambCatFr{Vous entendez, femmes, comment la coquine\\
nous arrange toutes, une fois de plus.\\\passeLigne}%
\end{groupe}%
%% v. 552b-553
\begin{groupe}%
\personnageFr{Mnésiloque}%
\tetramIambCatFr{\ligne{552}\espace{2.8cm}Et, par Zeus, je n'ai pas encore\\
dit tout ce que je sais. En voulez-vous davantage ?\\\passeLigne}%
\end{groupe}%
%% v. 554
\begin{groupe}%
\personnageFr{Mika}%
\tetramIambCatFr{Tu ne pourrais. Ce que tu savais, tu l'as épanché.\\\passeLigne}%
\end{groupe}%
%% v. 555-557a
\begin{groupe}%
\personnageFr{Mnésiloque}%
\tetramIambCatFr{Non, par Zeus, pas le dix-millième de ce que nous faisons.\\
Car enfin, je n'ai pas dit – tu vois ? – comment nous prenons des strigiles\\
pour ensuite nous siphonner le con...\\\passeLigne}%
\end{groupe}%
%% v. 557b
\begin{groupe}%
\personnageFr{Mika}%
\tetramIambCatFr{\ligne{557}\espace{4.5cm}La peste soit de toi.\\\passeLigne}%
\end{groupe}%
%% v. 558-559a
\begin{groupe}%
\personnageFr{Mnésiloque}%
\tetramIambCatFr{... que les viandes des Apaturies% \commentaire{558}{τὰ κρέ᾽ ἐξ Ἀπατουρίων}{Cf. note sur la Paix, v. 893, p. 112, ci-devant}
, nous les distribuons à nos procureuses,\\
puis que nous accusons la belette...\\\passeLigne}%
\end{groupe}%
%% v. 559b
\begin{groupe}%
\personnageFr{Mika}%
\tetramIambCatFr{\ligne{559}\espace{3.9cm}Merci de moi ! Tu radotes.\\\passeLigne}%
\end{groupe}%
%% v. 560-562a
\begin{groupe}%
\personnageFr{Mnésiloque}%
\tetramIambCatFr{... ni que telle autre occit son époux d'un coup de hache,\\
je ne l'ai pas dit ; ni qu'une autre rendit fou son mari par des drogues,\\
ni que naguère sous la baignoire...\\\passeLigne}%
\end{groupe}%
%% v. 562b
\begin{groupe}%
\personnageFr{Mika}%
\tetramIambCatFr{\ligne{562}\espace{6cm}Crève...\\\passeLigne}%
\end{groupe}%
%% v. 563a
\begin{groupe}%
\personnageFr{Mnésiloque}%
\tetramIambCatFr{... cette Acharnienne enfouit son père.\\\passeLigne}%
\end{groupe}%
%% v. 563b
\begin{groupe}%
\personnageFr{Mika}%
\tetramIambCatFr{\ligne{563}\espace{3.5cm}Supporterons-nous ce langage ?\\\passeLigne}%
\end{groupe}%
%% v. 564-565
\begin{groupe}%
\personnageFr{Mnésiloque}%
\tetramIambCatFr{Ni que toi, quand ta servante mit au jour un garçon,\\
tu substituas l'enfant, et livras ta petite fille.\\\passeLigne}%
\end{groupe}%
%% v. 566-567a
\begin{groupe}%
\personnageFr{Mika}%
\tetramIambCatFr{Non, par les deux déesses, tu ne l'auras pas dit gratuitement.\\
Je t'arracherai la toison.\\\passeLigne}%
\end{groupe}%
%% v. 567b
\begin{groupe}%
\personnageFr{Mnésiloque}%
\tetramIambCatFr{\ligne{567}\espace{4.6cm}Non, par Zeus, tu ne me toucheras pas.\\\passeLigne}%
\end{groupe}%
%% v. 568a
\begin{groupe}%
\personnageFr{Mika}%
\tetramIambCatFr{Eh bien, tiens !\\\passeLigne}%
\end{groupe}%
%% v. 568b
\begin{groupe}%
\personnageFr{Mnésiloque}%
\tetramIambCatFr{\ligne{568}\espace{2.1cm}Eh bien, tiens !\\\passeLigne}%
\end{groupe}%
%% v. 568c
\begin{groupe}%
\personnageFr{Mika}%
\tetramIambCatFr{\ligne{568}\espace{4.3cm}Prends mon manteau, Philista.\\\passeLigne}%
\end{groupe}%
%% v. 569a
\begin{groupe}%
\personnageFr{Mnésiloque}%
\tetramIambCatFr{Atteins-moi seulement, et par Artémis...\\\passeLigne}%
\end{groupe}%
%% v. 569b
\begin{groupe}%
\personnageFr{Mika}%
\tetramIambCatFr{\ligne{569}\espace{6.2cm}Que feras-tu ?\\\passeLigne}%
\end{groupe}%
%% v. 570
\begin{groupe}%
\personnageFr{Mnésiloque}%
\tetramIambCatFr{Le gâteau de sésame que tu as dévoré, je te le ferai évacuer.\\\passeLigne}%
\end{groupe}%
%% v. 571-573
\begin{groupe}%
\personnageFr{Coryphée}%
\tetramIambCatFr{Trêve d'invectives. Car une femme\\
accourt vers nous en toute hâte. Avant qu'elle soit proche,\\
faites silence. Écoutons en bonne contenance ce qu'elle va nous dire.\\\passeLigne}%
\end{groupe}%
%% v. 574-581
\begin{groupe}%
\personnageFr{Clisthène}%
\trimIambFr{Femmes aimées, qui m'êtes unies par la parenté des mœurs,\\
que je vous suis ami, mes joues le prouvent assez.\\
Oui, j'aime les femmes à la rage, et toujours je m'emploie pour vous.\\
Aussi, ayant appris sur votre sujet une chose capitale,\\
dont on jasait tout à l'heure sur l'agora,\\
je suis venu la raconter et vous donner avis\\
de veiller et d'être en garde, de crainte\\
qu'une grave et fâcheuse affaire ne vous tombe sur les bras au dépourvu.\\\passeLigne}%
\end{groupe}%
%% v. 582-583
\begin{groupe}%
\personnageFr{Coryphée}%
\trimIambFr{Qu'y a-t-il mon enfant ? Car il est naturel que je t'appelle enfant,\\
tant que tu auras ainsi les joues lisses.\\\passeLigne}%
\end{groupe}%
%% v. 584-585
\begin{groupe}%
\personnageFr{Clisthène}%
\trimIambFr{On dit qu'Euripide a envoyé ici même, aujourd'hui,\\
un de ses parents, un vieillard.\\\passeLigne}%
\end{groupe}%
%% v. 586
\begin{groupe}%
\personnageFr{Coryphée}%
\trimIambFr{Pourquoi faire, dans quel dessein ?\\\passeLigne}%
\end{groupe}%
%% v. 587-588
\begin{groupe}%
\personnageFr{Clisthène}%
\trimIambFr{À cette fin que, quoi que vous agitiez ou décidiez de faire,\\
\lednopb%
il puisse épier vos discours.\\\passeLigne}%
\end{groupe}%
%% v. 589
\begin{groupe}%
\personnageFr{Coryphée}%
\trimIambFr{Et comment s'est-il dissimulé au milieu de femmes ? Un homme !\\\passeLigne}%
\end{groupe}%
%% v. 590-591
\begin{groupe}%
\personnageFr{Clisthène}%
\trimIambFr{Euripide l'a flambé et épilé ;\\
et pour tout le reste l'a équipé en femme.\\\passeLigne}%
\end{groupe}%
%% v. 592-594
\begin{groupe}%
\personnageFr{Mnésiloque}%
\trimIambFr{Vous en fiez-vous à lui? Quel homme est assez\\
niais pour s'être laissé épiler.\\
Je n'en crois rien, pour ma part, ô déesses très honorées.\\\passeLigne}%
\end{groupe}%
%% v. 595-596
\begin{groupe}%
\personnageFr{Clisthène}%
\trimIambFr{Tu veux rire. Je ne serais pas venu l'annoncer\\
si je ne le tenais de ceux qui le savent bien.\\\passeLigne}%
\end{groupe}%
%% v. 597-602
\begin{groupe}%
\personnageFr{Coryphée}%
\trimIambFr{C'est une grave affaire qu'on nous annonce là.\\
Or ça, femmes, il s'agit, non de se morfondre,\\
mais de guetter l'homme, de chercher où\\
il se dérobe à nous et se tient caché.\\
Toi, aide-nous à le trouver, que nous t'ayons\\
encore cette autre obligation, ô notre patron.\\\passeLigne}%
\end{groupe}%
%% v. 603a
\begin{groupe}%
\personnageFr{Clisthène}%
\trimIambFr{Voyons, toi d'abord, qui es-tu ?\\\passeLigne}%
\end{groupe}%
%% v. 603b
\begin{groupe}%
\personnageFr{Mnésiloque}%
\trimIambFr{\ligne{603}\espace{4cm}(À part :) Où se tourner ?\\\passeLigne}%
\end{groupe}%
%% v. 604a
\begin{groupe}%
\personnageFr{Clisthène}%
\trimIambFr{... Car on va vous passer en revue.\\\passeLigne}%
\end{groupe}%
%% v. 604b
\begin{groupe}%
\personnageFr{Mnésiloque}%
\trimIambFr{\ligne{604}\espace{2.5cm}(À part :) Malheureux moi !\\\passeLigne}%
\end{groupe}%
%% v. 605
\begin{groupe}%
\personnageFr{Mika}%
\trimIambFr{Tu me demandais qui je suis ? la femme de Cléonyme.\\\passeLigne}%
\end{groupe}%
%% v. 606
\begin{groupe}%
\personnageFr{Clisthène}%
\trimIambFr{Vous savez, vous autres, qui est cette femme ?\\\passeLigne}%
\end{groupe}%
%% v. 607
\begin{groupe}%
\personnageFr{Coryphée}%
\trimIambFr{Nous le savons. Examine les autres.\\\passeLigne}%
\end{groupe}%
%% v. 608-609a
\begin{groupe}%
\personnageFr{Clisthène}%
\trimIambFr{Et celle-là, qui est-elle ? Celle qui porte\\
l'enfant ?\\\passeLigne}%
\end{groupe}%
%% v. 609b
\begin{groupe}%
\personnageFr{Mika}%
\trimIambFr{\ligne{609}\espace{1.3cm}C'est ma nourrice, par Zeus.\\\passeLigne}%
\end{groupe}%
%% v. 609c
\begin{groupe}%
\personnageFr{Mnésiloque}%
\trimIambFr{\ligne{609}\espace{3.8cm}Je suis perdu. (Il va pour sortir.)\\\passeLigne}%
\end{groupe}%
%% v. 610
\begin{groupe}%
\personnageFr{Clisthène}%
\trimIambFr{Holà; toi, où vas-tu ? Reste ici. (Mnésiloque s'accroupit.) Qu'est-ce qui t'arrive ?\\\passeLigne}%
\end{groupe}%
%% v. 611a
\begin{groupe}%
\personnageFr{Mnésiloque}%
\trimIambFr{Laisse-moi lâcher l'eau. Tu es un impudent.\\\passeLigne}%
\end{groupe}%
%% v. 611b-612
\begin{groupe}%
\personnageFr{Clisthène}%
\trimIambFr{%\ligne{611}\espace{2.9cm}Tu es une impudente.\\
Eh bien, fais. J'attendrai ici.\\\passeLigne}%
\end{groupe}%
%% v. 613-614
\begin{groupe}%
\personnageFr{Coryphée}%
\trimIambFr{Oui, attends, et surveille-la bien.\\
Elle est la seule que nous ne connaissions pas.\\\passeLigne}%
\end{groupe}%
%% v. 615a
\begin{groupe}%
\personnageFr{Clisthène}%
\trimIambFr{Tu es bien longue à faire de l'eau.\\\passeLigne}%
\end{groupe}%
%% v. 615b-616
\begin{groupe}%
\personnageFr{Mnésiloque}%
\trimIambFr{\ligne{615}\espace{4.2cm}Par Zeus, l'ami,\\
c'est que j'ai une rétention : hier j'ai mangé du cardame.\\\passeLigne}%
\end{groupe}%
%% v. 617
\begin{groupe}%
\personnageFr{Clisthène}%
\trimIambFr{Que parles-tu de cardame ? Veux-tu venir ici, près de moi ? (Il la saisit)\\\passeLigne}%
\end{groupe}%
%% v. 618a
\begin{groupe}%
\personnageFr{Mnésiloque}%
\trimIambFr{Pourquoi me tirer ainsi ? une valétudinaire !\\\passeLigne}%
\end{groupe}%
%% v. 618b-619a
\begin{groupe}%
\personnageFr{Clisthène}%
\trimIambFr{\espace{4.5cm}Réponds-moi.\\
Quel est ton mari ?\\\passeLigne}%
\end{groupe}%
%% v. 619b-620
\begin{groupe}%
\personnageFr{Mnésiloque}%
\trimIambFr{\espace{2.7cm}Mon mari, dis-tu ?\\
Connais-tu... chose, du dème de Cothocides.\\\passeLigne}%
\end{groupe}%
%% v. 621a
\begin{groupe}%
\personnageFr{Clisthène}%
\trimIambFr{Chose ? Qui ça ?\\\passeLigne}%
\end{groupe}%
%% v. 621b-622a
\begin{groupe}%
\personnageFr{Mnésiloque}%
\trimIambFr{\espace{2cm}C'est chose, qui naguère...\\
avec chose, fils de chose...\\\passeLigne}%
\end{groupe}%
%% v. 622b-623a
\begin{groupe}%
\personnageFr{Clisthène}%
\trimIambFr{\espace{3.7cm}Tu m'as l'air de baliverner.\\
Es-tu déjà venue ici ?\\\passeLigne}%
\end{groupe}%
%% v. 623b-624a
\begin{groupe}%
\personnageFr{Mnésiloque}%
\trimIambFr{\espace{4.5cm}Oui, par Zeus,\\
tous les ans.\\\passeLigne}%
\end{groupe}%
%% v. 624b
\begin{groupe}%
\personnageFr{Clisthène}%
\trimIambFr{\espace{1.5cm}Et quelle est ta compagne de tente ?\\\passeLigne}%
\end{groupe}%
%% v. 625a
\begin{groupe}%
\personnageFr{Mnésiloque}%
\trimIambFr{La mienne, c'est chose...\\\passeLigne}%
\end{groupe}%
%% v. 625b
\begin{groupe}%
\personnageFr{Clisthène}%
\trimIambFr{\espace{2.1cm}Misère de moi. Tu ne dis rien.\\\passeLigne}%
\end{groupe}%
%% v. 626-629
\begin{groupe}%
\personnageFr{Mika}%
\trimIambFr{Retire-toi. C'est moi qui vais l'éplucher comme il faut\\
sur les cérémonies de l'an dernier. — Tiens-toi à distance !\\
tu es un homme, tu ne dois rien entendre. — Toi, dis-moi\\
lequel des rites sacrés nous fut révélé le premier.\\\passeLigne}%
\end{groupe}%
%% v. 630
\begin{groupe}%
\personnageFr{Mnésiloque}%
\trimIambFr{Voyons. Que fût-ce encore que le premier? ... Nous bûmes.\\\passeLigne}%
\end{groupe}%
%% v. 631a
\begin{groupe}%
\personnageFr{Mika}%
\trimIambFr{Et après ? le second ?\\\passeLigne}%
\end{groupe}%
%% v. 631b
\begin{groupe}%
\personnageFr{Mnésiloque}%
\trimIambFr{\espace{4cm}Nous bûmes à nos santés.\\\passeLigne}%
\end{groupe}%
%% v. 632
\begin{groupe}%
\personnageFr{Mika}%
\trimIambFr{Tu l'as appris de quelqu'une. Et le troisième ?\\\passeLigne}%
\end{groupe}%
%% v. 633
\begin{groupe}%
\personnageFr{Mnésiloque}%
\trimIambFr{Xenylla demanda une nef, car il n'y avait pas de pot de chambre.\\\passeLigne}%
\end{groupe}%
%% v. 634-635a
\begin{groupe}%
\personnageFr{Mika}%
\trimIambFr{Cela ne veut rien dire. Ici, viens ici, ô Clisthène.\\
Voilà l'homme dont tu parles.\\\passeLigne}%
\end{groupe}%
%% v. 635b
\begin{groupe}%
\personnageFr{Clisthène}%
\trimIambFr{\espace{3.7cm}Que dois-je faire ?\\\passeLigne}%
\end{groupe}%
%% v. 636
\begin{groupe}%
\personnageFr{Mika}%
\trimIambFr{Déshabille-le. Ce qu'il dit n'a pas de sens.\\\passeLigne}%
\end{groupe}%
%% v. 637-8
\begin{groupe}%
\personnageFr{Mnésiloque}%
\trimIambFr{Eh quoi! déshabiller une mère de neuf enfants !\\
Il défait mon strophion avec promptitude. Impudent !\\\passeLigne}%
\end{groupe}%
%% v. 639-640
\begin{groupe}%
\personnageFr{Mika}%
\trimIambFr{Qu'elle a l'air robuste et vigoureuse !\\
Et, par Zeus, elle n'a pas de tétons comme nous.\\\passeLigne}%
\end{groupe}%
%% v. 641
\begin{groupe}%
\personnageFr{Mnésiloque}%
\trimIambFr{C'est que je suis stérile. Je n'ai jamais conçu.\\\passeLigne}%
\end{groupe}%
%% v. 642
\begin{groupe}%
\personnageFr{Clisthène}%
\trimIambFr{À présent. Mais tout à l'heure tu étais mère de neuf enfants.\\
Tiens-toi debout. Où pousses-tu ton cas là-dessous ?\\\passeLigne}%
\end{groupe}%
%% v. 644
\begin{groupe}%
\personnageFr{Mika}%
\trimIambFr{(Derrière Mnésiloque). Le voilà qui émerge, et beau de couleur, le coquin.\\\passeLigne}%
\end{groupe}%
%% v. 645a
\begin{groupe}%
\personnageFr{Clisthène}%
\trimIambFr{(Échangeant de place avec Mika). Où est-il ?\\\passeLigne}%
\end{groupe}%
%% v. 645b
\begin{groupe}%
\personnageFr{Mika}%
\trimIambFr{\ligne{645}\espace{2.2cm}Il est reparti en avant.\\\passeLigne}%
\end{groupe}%
%% v. 646a
\begin{groupe}%
\personnageFr{Clisthène}%
\trimIambFr{(Revenant devant.) Non, pas de ce côté.\\\passeLigne}%
\end{groupe}%
%% v. 646b
\begin{groupe}%
\personnageFr{Mika}%
\trimIambFr{\ligne{646}\espace{2.2cm}Eh! le voilà revenu par ici.\\\passeLigne}%
\end{groupe}%
%% v. 647-648
\begin{groupe}%
\personnageFr{Clisthène}%
\trimIambFr{Une façon d'isthme que tu as là, homme : tu fais passer ton cas de-ci de-là,\\
plus dru que les Corinthiens.\\\passeLigne}%
\end{groupe}%
%% v. 649-650a
\begin{groupe}%
\personnageFr{Mika}%
\trimIambFr{Ah! la canaille! Voilà donc pourquoi au profit d'Euripide\\
il nous invectivait.\\\passeLigne}%
\end{groupe}%
%% v. 650b-651
\begin{groupe}%
\personnageFr{Mnésiloque}%
\trimIambFr{\ligne{650}\espace{2.5cm}Malheureux que je suis !\\
dans quelle méchante affaire je me suis empêtré !\\\passeLigne}%
\end{groupe}%
%% v. 652a
\begin{groupe}%
\personnageFr{Mika}%
\trimIambFr{Voyons, que faisons-nous ?\\\passeLigne}%
\end{groupe}%
%% v. 652b-654
\begin{groupe}%
\personnageFr{Clisthène}%
\trimIambFr{\ligne{652}\espace{2.6cm}Veillez bien sur lui,\\
qu'il ne se dérobe point par la fuite.\\
Moi, je vais rapporter la chose aux prytanes.\\\passeLigne}%
\end{groupe}%
%% v. 655-662
\begin{groupe}%
\personnageFr{Chœur}%
\tetramAnapFr{A nous maintenant, après avoir allumé nos torches,\\
serré bien et résolument nos tuniques et mis bas nos manteaux,\\
de chercher si d'aventure il n'est pas entré quelqu'autre homme. Tournons et virons\\
tout à l'entour du Pnyx, inspectons avec soin les tentes et les passages.\\\passeLigne}%
\end{groupe}%
\begin{groupe}%
\tetramTrochFr{Or sus! tout d'abord élançons-nous au dehors d'un pied léger\\
et épions partout en silence. Surtout ne tardons pas : le cas ne souffre plus de délai.\\
Il faut que notre coryphée coure au plus tôt faire la ronde.\\
(Ici commence un chœur dansé.)\\\passeLigne}%
\end{groupe}%
%% v. 663-666
\begin{groupe}%
\chantChoeurFr{Or sus ! suis chaque piste, quête vivement,}%
\end{groupe}%
%% v. 663-666
\begin{groupe}%
\chantChoeurFr{au cas où quelqu'autre en ces lieux}%
\end{groupe}%
%% v. 663-666
\begin{groupe}%
\chantChoeurFr{se serait tapi à notre insu.}%
\end{groupe}%
%% v. 663-666
\begin{groupe}%
\chantChoeurFr{Jette les yeux partout.}%
\end{groupe}%
%% v. 663-666
\begin{groupe}%
\chantChoeurFr{Que rien ici, rien là-bas,}%
\end{groupe}%
%% v. 663-666
\begin{groupe}%
\chantChoeurFr{n'échappe à ton \oe il vigilant..\\\passeLigne}%
\end{groupe}%
%% v. 667-686
\begin{groupe}%
\chantChoeurFr{Car s'il est pris après ce sacrilège,}%
\end{groupe}%
%% v. 663-666
\begin{groupe}%
\chantChoeurFr{il paiera son crime, et de surcroît}%
\end{groupe}%
%% v. 663-666
\begin{groupe}%
\chantChoeurFr{prémunira tous les autres,}%
\end{groupe}%
%% v. 663-666
\begin{groupe}%
\chantChoeurFr{par son exemple, contre la violence,}%
\end{groupe}%
%% v. 663-666
\begin{groupe}%
\chantChoeurFr{les actes injurieux}%
\end{groupe}%
%% v. 663-666
\begin{groupe}%
\chantChoeurFr{et les façons impies.}%
\end{groupe}%
%% v. 663-666
\begin{groupe}%
\chantChoeurFr{Il confessera hautement qu'il est des dieux,}%
\end{groupe}%
%% v. 663-666
\begin{groupe}%
\chantChoeurFr{il apprendra}%
\end{groupe}%
%% v. 663-666
\begin{groupe}%
\chantChoeurFr{à tous les hommes à les honorer,}%
\end{groupe}%
%% v. 663-666
\begin{groupe}%
\chantChoeurFr{à s'efforcer, par une conduite pieuse et un constant souci des lois,}%
\end{groupe}%
%% v. 663-666
\begin{groupe}%
\chantChoeurFr{de ne rien faire que de bien.}%
\end{groupe}%
%% v. 663-666
\begin{groupe}%
\chantChoeurFr{S'ils se comportent autrement, voici ce qui se passera :}%
\end{groupe}%
%% v. 663-666
\begin{groupe}%
\chantChoeurFr{celui d'eux qui sera pris à faire un sacrilège,}%
\end{groupe}%
%% v. 663-666
\begin{groupe}%
\chantChoeurFr{enflammé de fureur, délirant de rage,}%
\end{groupe}%
%% v. 663-666
\begin{groupe}%
\chantChoeurFr{témoignera aux femmes et à tous les mortels}%
\end{groupe}%
%% v. 663-666
\begin{groupe}%
\chantChoeurFr{que la divinité châtie}%
\end{groupe}%
%% v. 663-666
\begin{groupe}%
\chantChoeurFr{les atteintes aux lois et à la piété et ne se fait pas attendre.\\\passeLigne}%
\end{groupe}%
%% v. 687-688
\begin{groupe}%
\personnageFr{Coryphée}%
\tetramTrochFr{\ligne{687}Mais il nous semble avoir bien cherché quasi-partout.\\
Et nous ne voyons pas d'autre homme caché parmi nous.\\
(Mnésiloque enlève l'enfant des bras de Mica et se réfugie sur l'autel.)\\\passeLigne}%
\end{groupe}%
%% v. 689-691
\begin{groupe}%
\personnageFr{Mika}%
\trimIambFr{Ah ! Ah !}%
\end{groupe}%
%% v. 689-691
\begin{groupe}%
%\personnageFr{Mika}%
\trimIambFr{Où ? Où te sauves-tu ? holà ! holà ! Arrête. \\
Malheur ! Ah ! malheur ! le voilà parti avec mon enfant\\
arraché de ma mamelle.\\\passeLigne}%
\end{groupe}%
%% v. 692-695a
\begin{groupe}%
\personnageFr{Mnésiloque}%
\trimIambFr{Hurle. Jamais plus tu ne l'appateras,\\
si vous ne me laissez pas libre. Mais ici même « sur ces oblations,\\
frappé de ce couteau, du pur sang de ses veines\\
il rougira l'autel ».\\\passeLigne}%
\end{groupe}%
%% v. 695b-698
\begin{groupe}%
\personnageFr{Mika}%
\trimIambFr{\ligne{695}\espace{3cm}Malheur à moi !\\
Femmes, au secours ! Poussez une puissante clameur,\\
une clameur à faire lâcher pied. Souffrirez-vous\\
qu'on me ravisse mon unique enfant?\\\passeLigne}%
\end{groupe}%
%% v. 699-701
\begin{groupe}%
\personnageFr{Chœur}\chantChoeurFr{Ah! Ah !}%
\end{groupe}%
%% v. 699-701
\begin{groupe}%
\chantChoeurFr{Augustes M\oe res, quel prodige}%
\end{groupe}%
%% v. 699-701
\begin{groupe}%
\chantChoeurFr{nouveau vient frapper mes yeux !\\\passeLigne}%
\end{groupe}%
%% v. 702-703
\begin{groupe}%
\tetramTrochFr{Tout n'est donc qu'audace doublée d'effronterie !\\
Quel forfait encore, amies, quel forfait il a commis !\\\passeLigne}%
\end{groupe}%
%% v. 704
\begin{groupe}%
\personnageFr{Mnésiloque}%
\tetramTrochFr{Quel forfait ? je vous extirperai votre outrecuidance !\\\passeLigne}%
\end{groupe}%
%% v. 705
\begin{groupe}%
\personnageFr{Chœur}%
\tetramTrochFr{N'est-ce pas une abomination, et pis encore ?\\\passeLigne}%
\end{groupe}%
%% v. 706
\begin{groupe}%
\personnageFr{Mika}%
\trimIambFr{Une abomination, oui, qu'on m'ait arraché mon enfant.\\\passeLigne}%
\end{groupe}%
%% v. 707-708
\begin{groupe}%
\personnageFr{Chœur}%
\dimAnapFr{Que dire après cela, quand\\
l'auteur d'une telle action a perdu toute honte ?\\\passeLigne}%
\end{groupe}%
%% v. 709
\begin{groupe}%
\personnageFr{Mnésiloque}%
\dimAnapFr{Et je n'ai pas fini.\\\passeLigne}%
\end{groupe}%
%% v. 710
\begin{groupe}%
\personnageFr{Chœur}%
\chantChoeurFr{Au moins, reviens d'où tu es parti ;}%
\end{groupe}%
%% v. 711-713
\begin{groupe}%
\chantChoeurFr{Ce n'est pas avec une fuite facile que tu ne t'en sortiras.}%
\end{groupe}%
%% v. 710
\begin{groupe}%
\chantChoeurFr{Avec tout ce que tu as fait en t'immisçant ici,}%
\end{groupe}%
%% v. 710
\begin{groupe}%
\chantChoeurFr{malheur t'en prendra !\\\passeLigne}%
\end{groupe}%
%% v. 714
\begin{groupe}%
\personnageFr{Mnésiloque}%
\tetramTrochFr{Puisse-t-il n'en être rien, j'en conjure le ciel.\\\passeLigne}%
\end{groupe}%
%% v. 715-716
\begin{groupe}%
\personnageFr{Chœur}%
\diDochFr{Et qui donc, qui des dieux immortels s'associera\\
à des actes injurieux ?\\\passeLigne}%
\end{groupe}%
%% v. 717
\begin{groupe}%
\personnageFr{Mnésiloque}%
\trimIambCatFr{Inutiles bavardages ! Je ne la lâcherai pas.\\\passeLigne}%
\end{groupe}%
%% v. 718-725
\begin{groupe}%
\personnageFr{Chœur}%
\chantChoeurFr{Ah! par les deux déesses,\\
peut-être n'auras-tu pas lieu de rire\\
de tes outrages et des discours insensés que tu profères.}
\end{groupe}%
%% v. 718-725
\begin{groupe}%
\chantChoeurFr{Car toi, avec tes impiétés,}%
\end{groupe}%
%% v. 718-725
\begin{groupe}%
\chantChoeurFr{nous te châtierons de ta conduite, comme de juste.}%
\end{groupe}%
%% v. 718-725
\begin{groupe}%
\chantChoeurFr{Et puis la Fortune, changeant brusquement d'humeur aura vite fait de }%
\end{groupe}%
%% v. 718-725
\begin{groupe}%
\chantChoeurFr{poursuivre un méchant homme qui a pris le chemin du crime.\\\passeLigne}%
\end{groupe}%
%% v. 726-727
\begin{groupe}%
\personnageFr{Coryphée}%
\tetramTrochFr{\ligne{726}Mais prends ces femmes avec toi, et va chercher du bois.\\
Il faut le brûler vif, ce fourbe, l'anéantir au plus tôt par le feu.\\\passeLigne}%
\end{groupe}%
%% v. 728-729
\begin{groupe}%
\personnageFr{Mika}%
\trimIambFr{Allons chercher les fagots, Mania.\\
Je ferai de toi un tison aujourd'hui..\\\passeLigne}%
\end{groupe}%
%% v. 730-732
\begin{groupe}%
\personnageFr{Mnésiloque}%
\trimIambFr{Allumez, brûlez. Toi, ma petite, ôte vite ta crétique\\
et n'accuse de ta mort que ta mère.\\
Qu'est ceci ? La fillette n'est plus qu'une outre}%
\end{groupe}%
%% v. 733-735
\begin{groupe}%
\trimIambFr{pleine de vin, ... et avec cela chaussée de persiques.\\
Ô femmes archi-enfiévrées et archi-biberonnes,\\
ingénieuses à trouver partout matière à boire !}%
\end{groupe}%
%% v. 736-738
\begin{groupe}%
\trimIambFr{Ô providence des débitants, et notre fléau à nous !\\
Fléau aussi des ustensiles de ménage, et fléau de la trame !\\
(Les femmes reviennent avec des fagots.)\\\passeLigne}%
\end{groupe}%
%% v. 739
\begin{groupe}%
\personnageFr{Mika}%
\trimIambFr{Entasse force fagots autour de lui, Mania.\\\passeLigne}%
\end{groupe}%
%% v. 740-741a
\begin{groupe}%
\personnageFr{Mnésiloque}%
\trimIambFr{Entasse, c'est bon. Mais toi, réponds à ceci.\\
Tu te dis mère de cet enfant ?\\\passeLigne}%
\end{groupe}%
%% v. 741b-742a
\begin{groupe}%
\personnageFr{Mika}%
\trimIambFr{\ligne{741}\espace{2.8cm}Dix mois durant\\
je l'ai porté.\\\passeLigne}%
\end{groupe}%
%% v. 742b
\begin{groupe}%
\personnageFr{Mnésiloque}%
\trimIambFr{\ligne{742}\espace{1.5cm}Tu l'as porté, toi ?\\\passeLigne}%
\end{groupe}%
%% v. 742c
\begin{groupe}%
\personnageFr{Mika}%
\trimIambFr{\ligne{742}\espace{3.5cm}Oui, par Artémis.\\\passeLigne}%
\end{groupe}%
%% v. 743a
\begin{groupe}%
\personnageFr{Mnésiloque}%
\trimIambFr{(Il porte l'outre à ses lèvres.) Du trois cotyles ? Ou de quel prix ?\\\passeLigne}%
\end{groupe}%
%% v. 743b-745a
\begin{groupe}%
\personnageFr{Mika}%
\trimIambFr{\ligne{743}\espace{4.2cm}Dis-moi : quel tour vas-tu me faire ?\\
Impudent, tu as dévêtu mon enfant,\\
qui est si menu.\\\passeLigne}%
\end{groupe}%
%% v. 745b-746
\begin{groupe}%
\personnageFr{Mnésiloque}%
\trimIambFr{\ligne{745}\espace{2.2cm}Si menu ? (Il le regarde.) Il est petit, par Zeus.\\
Et quel âge a-t-il ? trois ou quatre choés ?\\\passeLigne}%
\end{groupe}%
%% v. 747-748a
\begin{groupe}%
\personnageFr{Mika}%
\trimIambFr{À peu près, mais elle est ainsi depuis les dionysiaques.\\
Allons ! rends-le-moi.\\\passeLigne}%
\end{groupe}%
%% v. 748b
\begin{groupe}%
\personnageFr{Mnésiloque}%
\trimIambFr{\ligne{748}\espace{2.8cm}Non, par le glorieux Apollon.%
%\commentaire{748}{μὰ τὸν Ἀπόλλω τουτονί.}{Cf. la note sur Nuées, v. 83, t. I, p.380.\eff{À reproduire en annexe !}}.
\\\passeLigne}%
\end{groupe}%
%% v. 749a
\begin{groupe}%
\personnageFr{Mika}%
\trimIambFr{Alors nous allons te brûler.\\\passeLigne}%
\end{groupe}%
%% v. 749b-750
\begin{groupe}%
\personnageFr{Mnésiloque}%
\trimIambFr{\ligne{749}\espace{3cm}Parfaitement, brûlez.\\
Mais j'égorge celle-ci à l'instant.\\\passeLigne}%
\end{groupe}%
%% v. 751-752a
\begin{groupe}%
\personnageFr{Mika}%
\trimIambFr{Non, non, je t'en conjure. Fais à moi tout ce que tu voudras, \\
plutôt qu'à elle.\\\passeLigne}%
\end{groupe}%
%% v. 752b-753
\begin{groupe}%
\personnageFr{Mnésiloque}%
\trimIambFr{\ligne{752}\espace{2.5cm}Tu as pour tes enfants une tendresse de naturel.\\
Mais elle n'en sera pas moins égorgée.\\\passeLigne}%
\end{groupe}%
%% v. 754-755
\begin{groupe}%
\personnageFr{Mika}%
\trimIambFr{Ah! mon enfant ! Donne-moi une patère, Mania,\\
que je recueille au moins le sang de ma fille..\\\passeLigne}%
\end{groupe}%
%% v. 756
\begin{groupe}%
\personnageFr{Mnésiloque}%
\trimIambFr{Tiens-la là-dessous. Je veux bien t'accorder cette seule grâce.\\
(Il asperge Mika du vin de l'outre.)\\\passeLigne}%
\end{groupe}%
%% v. 757
\begin{groupe}%
\personnageFr{Mika}%
\trimIambFr{Puisses-tu périr de male mort! Que tu es jaloux et mauvais !\\\passeLigne}%
\end{groupe}%
%% v. 758
\begin{groupe}%
\personnageFr{Mnésiloque}%
\trimIambFr{Cette peau revient à la prêtresse.\\\passeLigne}%
\end{groupe}%
%% v. 759a
\begin{groupe}%
\personnageFr{Kritylla}%
\trimIambFr{\ligne{758}Qu'est-ce qui revient à la prêtresse ?\\\passeLigne}%
\end{groupe}%
%% v. 759b
\begin{groupe}%
\personnageFr{Mnésiloque}%
\trimIambFr{\ligne{759}\espace{4cm}Tiens, prends. (Il lui jette l'outre vide.)\\\passeLigne}%
\end{groupe}%
%% v. 760-761
\begin{groupe}%
\personnageFr{Kritylla}%
\trimIambFr{Infortunée Mica, qui t'a souillée ?\\
Qui t'a soustrait ta fille chérie ?\\\passeLigne}%
\end{groupe}%
%% v. 762-764
\begin{groupe}%
\personnageFr{Mika}%
\trimIambFr{C'est ce fourbe. Mais puisque te voilà,\\
garde-le ; que j'aille, en compagnie de Clisthène,\\
raconter aux prytanes ce qu'il a fait.\\\passeLigne}%
\end{groupe}%
%% v. 765-775
\begin{groupe}%
\personnageFr{Mnésiloque}%
\trimIambFr{\ligne{765}Voyons, quel sera l'instrument de mon salut ?\\
Quelle tentative? Quelle combinaison ? Car l'auteur vrai,\\
celui qui m'impliqua dans cette méchante affaire,\\
n'apparaît pas encore. Allons, quel messager\\
lui dépêcherai-je ? (Il réfléchit)... Il me vient à l'esprit un expédient\\
tiré du \italique{Palamède}. Comme cet autre, je lancerai mes avirons\\
chargés d'écritures... Mais je ne les ai pas sous la main, ces avirons...\\
Comment donc faire pour me procurer des avirons ? Comment ?...\\
Mais quoi, si à défaut d'avirons, j'écrivais sur ces offrandes\\
et les jetais de tous côtés ? Cela vaut bien mieux.\\
Elles sont de bois après tout; ceux-là aussi étaient de bois.\\\passeLigne}%
\end{groupe}%
%% v. 776-784
\begin{groupe}%
\anapFr{À l'œuvre, ô mes mains,
}%
\end{groupe}%
%% v. 776-784
\begin{groupe}%
\anapFr{il s'agit d'un travail expédient.\\
Allons, feuillets de mes tablettes polies,}%
\end{groupe}%
%% v. 776-784
\begin{groupe}%
\anapFr{\ligne{779}recevez les entailles du ciseau,\\
messagères de mes misères. Aïe,}%
\end{groupe}%
%% v. 776-784
\begin{groupe}%
\anapFr{c'est pénible à percer !..\\
«~Ac...cours~», «~Ac...cours !~» et quel sera notre cours d'eau ?\\
(Il se tourne des deux côtés et expédie les tablettes par voie aérienne.)\\
\ligne{782}Partez, volez par toutes les voies}%
\end{groupe}%
%% v. 776-784
\begin{groupe}%
\lednopb%
\anapFr{\ligne{784}avec cette nouvelle façon de faire : il faut aller vite.\\\passeLigne}%
\end{groupe}%
%% v. 785-843
\begin{groupe}%
\personnageFr{Chœur}%
\tetramAnapFr{Maintenant nous allons faire nous-mêmes notre éloge dans la parabase.\\
Chacun pourtant médit à l'envi de l'espèce féminine :\\
nous sommes un pur fléau pour les hommes ; de nous viennent tous les maux,\\
disputes, querelles, affreuses séditions, chagrins, guerres. Çà, voyons,\\
si nous sommes un fléau, pourquoi nous épousez-vous, ce fléau que nous sommes ?\\
Pourquoi nous interdire de sortir, de nous laisser prendre à pousser la tête dehors ?\\
Pourquoi vous piquer de le garder avec un soin si jaloux, ce fléau ?\\
Arrive-t-il à votre petite femme d'être en course, et que vous la trouviez absente,\\
vous entrez dans des fureurs, quand il vous faudrait faire libation et vous réjouir\\
de ce qu'il a disparu de chez vous, le fléau, et que vous ne le trouviez pas au logis.\\
Nous endormons-nous chez les autres, fatiguées de danser,\\
chacun se met après le fléau et rode autour des lits.\\\passeLigne}
\end{groupe}%
%% v. 785-843
\begin{groupe}%
\tetramAnapFr{Nous penchons-nous par une fenêtre, vous cherchez à le dévisager, le fléau,\\
et si par pudeur elle se retire, chacun n'en est que plus impatient\\
de la voir se pencher de nouveau, ce fléau. Tant il est vrai que nous valons bien mieux\\
que vous. Si l'on veut une preuve, il n'y a qu'à regarder.\\
Venons-en à la preuve et sachons qui sont les pires. Car nous disons que c'est vous,\\
et vous dites que c'est nous. Vérifions donc et mettons-les en présence isolément\\
en rapprochant chaque fois le nom de la femme de celui de l'homme.\\
Charminos est inférieur à Nausimaché : les faits parlent d'eux-mêmes,\\
et à coup sûr aussi Cléophon est pire en somme que Salabaccho.\\
Avec Aristomaché, l'héroïne de Marathon,\\\passeLigne}
\end{groupe}%
%% v. 785-843
\begin{groupe}%
\tetramAnapFr{et avec Stratonice, il y a beau jour que nul d'entre vous ne tenterait de se mesurer.\\
Et Euboulé, parmi les bouleutes de l'an dernier, en est-il un qui,\\
après avoir résigné à un autre son office, l'emporte sur elle ? Lui-même s'en défendra.\\
Ainsi nous nous prévalons d'une grande supériorité sur les hommes.\\
Joignez qu'une femme, après avoir diverti par cinquante talents les deniers publics,\\
ne se rendra pas à l'Acropole sur un char ; mais eût-elle commis le plus grave des larcins,\\
\lednopb%
un boisseau de froment volé à son mari, elle le restituera le jour même.\\\passeLigne}%
\end{groupe}%
%% v. 814-829
\begin{groupe}%
\anapFr{\ligne{814}Mais nous pourrions en désigner plusieurs parmi ceux-là\\
qui en font autant,\\
et qui, par là-dessus, sont gloutons\\
plus que nous, tire-laine,\\
charlatans et marchands d'hommes.\\
Et assurément aussi ils savent moins que nous\\
conserver les traditions\\\passeLigne}
\end{groupe}%
%% v. 814-829
\begin{groupe}%
\anapFr{qui sont par notre fait sauves encore aujourd'hui\\
– panier et ombrelle compris.\\
Tandis que chez nombre de nos époux que voilà,\\
la lance a fait disparaître\\
ce que l'on tirait de nos maisons – et même la lance ;\\
et que maint autre,\\
\{dans ses campagnes\},}%
\end{groupe}%
%% v. 814-829
\begin{groupe}%
\anapFr{\lednopb%
\ligne{829}a jeté à bas de ses épaules... l'ombrelle.\\\passeLigne}%
\end{groupe}%
%% v. 830-845
\begin{groupe}%
\personnageFr{Coryphée}%
\tetramTrochFr{Nous autres femmes aurions droit en toute justice\\
d'imputer aux hommes bien des torts, un surtout passant toute mesure.\\
Il serait de votre devoir, quand l'une de nous a donné à la patrie un homme utile,\\
taxiarque ou stratège, d'avoir pour elle quelque considération,\\
de lui accorder la proédrie dans les Sténies, les Scirophories\\
et les autres fêtes que nous célébrons.\\
Celle, au contraire, qui aurait donné le jour à un pleutre, à un propre à rien,\\
méchant triérarque ou lâche commandant de navire,\\
s'assiérait, les cheveux taillés au bol, derrière}
\end{groupe}%
%% v. 830-845
\begin{groupe}%
\tetramTrochFr{la femme qui enfanta le brave. Comment tolérer, ô Eupolis,\\
que la mère d'Hyperbolos siège, vêtue\\
de blanc, les cheveux flottants, à côté de la mère de Lamachos,\\
et qu'elle prête de l'argent à usure... le faut-il supporter ? Si elle prête de l'argent\\
et exige un profit, loin de lui payer un intérêt,\\
il faudrait lui ôter de force ses richesses, en lui disant :\\
« Tu vaux bien que ton argent produise, toi qui produisis un tel fruit. »\\\passeLigne}%
\end{groupe}%
%% v. 846-851
\begin{groupe}%
\personnageFr{Mnésiloque}%
\trimIambFr{Je suis devenu louche à lorgner de ce côté. Et lui ?... pas encore.\\
Qu'est-ce donc qui l'empêche ? Nul doute,\\
le \italique{Palamède} est froid et il en a honte.\\
Par quel drame alors l'attirer ...\\
J'y suis : je vais contrefaire sa récente \italique{Hélène}\\
Aussi bien, je suis accoutré d'un costume de femme.\\\passeLigne}%
\end{groupe}%
%% v. 852-854
\begin{groupe}%
\personnageFr{Kritylla}%
\trimIambFr{Que manigances-tu encore ? Qu'as-tu à écarquiller les yeux ?\\
On t'en fera voir une Hélène, amère celle-là, si tu ne te tiens décemment\\
jusqu'à la venue de l'un des prytanes.\\\passeLigne}%
\end{groupe}%
%% v. 855-857
\begin{groupe}%
\personnageFr{Mnésiloque}%
\trimIambFr{(Rôle d'Helène.) « Du Nil ce sont ici les rives aux belles Nymphes,\\
Nil qui à la place de l'eau du ciel, humecte le fondement d'Égypte\\
la blanche » et du peuple à la noire syrméa.\\\passeLigne}%
\end{groupe}%
%% v. 858
\begin{groupe}%
\personnageFr{Kritylla}%
\trimIambFr{Tu es un effronté, par Hécate la porte-torche. (Elle lui porte un coup de sa torche.)\\\passeLigne}%
\end{groupe}%
%% v. 859-860a
\begin{groupe}%
\personnageFr{Mnésiloque}%
\trimIambFr{« J'ai pour patrie un lieu non sans renom,\\
c'est Sparte, et mon père est Tyndare. »\\\passeLigne}%
\end{groupe}%
%% v. 860b-861
\begin{groupe}%
\personnageFr{Kritylla}%
\trimIambFr{\ligne{860}\espace{4.5cm}Lui ton père, être de malheur ?\\
Phrynondas, veux-tu dire.\\\passeLigne}%
\end{groupe}%
%% v. 862a
\begin{groupe}%
\personnageFr{Mnésiloque}%
\trimIambFr{« Je fus nommée Hélène. »\\\passeLigne}%
\end{groupe}%
%% v. 862b-863
\begin{groupe}%
\personnageFr{Kritylla}%
\trimIambFr{\ligne{862}\espace{2.8cm}Voilà qu'encore tu redeviens femme,\\
avant d'avoir expié ton premier déguisement.\\\passeLigne}%
\end{groupe}%
%% v. 864-865a
\begin{groupe}%
\personnageFr{Mnésiloque}%
\trimIambFr{« Et à cause de moi, bien des guerriers sont morts\\
aux rives du Scamandre. »\\\passeLigne}%
\end{groupe}%
%% v. 865b
\begin{groupe}%
\personnageFr{Kritylla}%
\trimIambFr{\ligne{865}\espace{2.4cm}Que n'as-tu fait comme eux ?\\\passeLigne}%
\end{groupe}%
%% v. 866-868a
\begin{groupe}%
\personnageFr{Mnésiloque}%
\trimIambFr{« Moi, me voici. Mais lui l'infortuné », mon époux\\
Ménélas n'apparaît toujours point.\\
« Ah ! pourquoi vis-je encore? »\\\passeLigne}%
\end{groupe}%
%% v. 868b
\begin{groupe}%
\personnageFr{Kritylla}%
\trimIambFr{\ligne{868}\espace{2.3cm}C'est la faute aux corbeaux.\\\passeLigne}%
\end{groupe}%
%% v. 869-870
\begin{groupe}%
\personnageFr{Mnésiloque}%
\trimIambFr{(Mnésiloque aperçoit Euripide.)
Mais quelque chose vient flatter mon cœur.\\
Ô Zeus, ne démens pas ma naissante espérance !%\\
%(À Euripide :) Sais-tu donc ce qu'il te reste à faire ? [ ... ] 
\\\passeLigne}%
\end{groupe}%
%% v. 871-873
\begin{groupe}%
\personnageFr{Euripide}%
\trimIambFr{(Rôle de Ménélas.) « Quel est le chef puissant de ces fortes demeures ? » \\
Voudra-t-il accueillir des hôtes dérangés par la houle marine\\
au cours d'une tempête et d'un naufrage ?\\\passeLigne}%
\end{groupe}%
%% v. 874a
\begin{groupe}%
\personnageFr{Mnésiloque}%
\trimIambFr{C'est à Protée qu'est ce palais.\\\passeLigne}%
\end{groupe}%
%% v. 874b-876
\begin{groupe}%
\personnageFr{Kritylla}%
\trimIambFr{\ligne{874}\espace{4.2cm}Qui ça, Protée ?\\
Oh le triple misérable! Il ment par les deux déesses.\\
Il est mort, Protéas, depuis dix ans.\\\passeLigne}%
\end{groupe}%
%% v. 877
\begin{groupe}%
\personnageFr{Euripide}%
\trimIambFr{Et quel est ce pays où ma nef a pris terre ?\\\passeLigne}%
\end{groupe}%
%% v. 878a
\begin{groupe}%
\personnageFr{Mnésiloque}%
\trimIambFr{L'Égypte.\\\passeLigne}%
\end{groupe}%
%% v. 878b
\begin{groupe}%
\personnageFr{Euripide}%
\trimIambFr{\ligne{878}\espace{1.7cm}« Vers quel lieu, malheureux, nous avons navigué ! »\\\passeLigne}%
\end{groupe}%
%% v. 879-880
\begin{groupe}%
\personnageFr{Kritylla}%
\trimIambFr{T'en fies-tu à ce coquin — la peste soit de lui —\\
qui te raconte des bourdes ? Ceci, c'est le Thesmophorion.\\\passeLigne}%
\end{groupe}%
%% v. 881
\begin{groupe}%
\personnageFr{Euripide}%
\trimIambFr{Et lui, Protée, est-il céans, ou s'il est hors ?\\\passeLigne}%
\end{groupe}%
%% v. 882-884
\begin{groupe}%
\personnageFr{Kritylla}%
\trimIambFr{Il faut que tu aies encore mal au c\oe ur, étranger.\\
On te dit qu'il est mort, Protéas,\\
et tu demandes s'il est céans ou dehors.\\\passeLigne}%
\end{groupe}%
%% v. 885
\begin{groupe}%
\personnageFr{Euripide}%
\trimIambFr{Hélas! il a vécu. Où donc est le sépulcre où il fut inhumé ?\\\passeLigne}%
\end{groupe}%
%% v. 886
\begin{groupe}%
\personnageFr{Mnésiloque}%
\trimIambFr{C'est ici son tombeau, tu nous y vois assise.\\\passeLigne}%
\end{groupe}%
%% v. 887-888
\begin{groupe}%
\personnageFr{Kritylla}%
\trimIambFr{Puisses-tu périr misérablement ! Et tu périras,\\
toi qui oses appeler l'autel une tombe.\\\passeLigne}%
\end{groupe}%
%% v. 889-890a
\begin{groupe}%
\personnageFr{Euripide}%
\trimIambFr{D'où vient qu'assise sur ce siège sépulcral,\\
tu dérobes tes traits sous un voile, étrangère ?\\\passeLigne}%
\end{groupe}%
%% v. 890b-891
\begin{groupe}%
\personnageFr{Mnésiloque}%
\trimIambFr{\espace{3.8cm}On prétend sans mon v\oe u\\
que le fils de Protée me reçoive en son lit.\\\passeLigne}%
\end{groupe}%
%% v. 892-894
\begin{groupe}%
\personnageFr{Kritylla}%
\trimIambFr{Pourquoi, misérable, tromper encore cet étranger ?\\
Il est entré en fraude, ô étranger,\\
auprès des femmes afin de toucher subrepticement à leur bijou.\\\passeLigne}%
\end{groupe}%
%% v. 895
\begin{groupe}%
\personnageFr{Mnésiloque}%
\trimIambFr{Jappe et verse à ton gré le blâme sur ma tête.\\\passeLigne}%
\end{groupe}%
%% v. 896
\begin{groupe}%
\personnageFr{Euripide}%
\trimIambFr{Étrangère, dis-moi, qui donc est cette vieille, qui t'incrimine ainsi ?\\\passeLigne}%
\end{groupe}%
%% v. 897a
\begin{groupe}%
\personnageFr{Mnésiloque}%
\trimIambFr{Elle ? C'est Théonoé, la fille de Protée.\\\passeLigne}%
\end{groupe}%
%% v. 897b-899a
\begin{groupe}%
\personnageFr{Kritylla}%
\trimIambFr{\ligne{897}\espace{3.6cm}Non, par les déesses,\\
je ne suis autre que Critylla, fille d'Antithéos, du dème de Gargettos.\\
Toi, tu es un fripon.\\\passeLigne}%
\end{groupe}%
%% v. 899b-901
\begin{groupe}%
\personnageFr{Mnésiloque}%
\trimIambFr{\ligne{899}\espace{3cm}Dis ce que tu voudras.\\
En épousant ton frère,\\
je ne trahirai point Ménélas, mon époux, qui reste à Troie.\\\passeLigne}%
\end{groupe}%
%% v. 902
\begin{groupe}%
\personnageFr{Euripide}%
\trimIambFr{Qu'as-tu dit ? Femme, tourne vers moi ta brillante prunelle.\\\passeLigne}%
\end{groupe}%
%% v. 903
\begin{groupe}%
\personnageFr{Mnésiloque}%
\trimIambFr{J'ai honte, après l'outrage imprimé sur mes joues.%\commentaire{903}{}{On n'oublie pas qu'il a perdu sa barbe.}.
\\\passeLigne}%
\end{groupe}%
%% v. 904-905
\begin{groupe}%
\personnageFr{Euripide}%
\trimIambFr{Qu'est-ce là ? Une stupeur me prend.\\
Ô dieux, que vois-je ? Femme, qui donc es-tu?\\\passeLigne}%
\end{groupe}%
%% v. 906
\begin{groupe}%
\personnageFr{Mnésiloque}%
\trimIambFr{Qui donc es-tu toi-même ? « C'est une question qui m'appartient aussi\commentaire{906}{« C'est une question qui m'appartient aussi ».}{J'emprunte ce vers à A. de Musset, \italique{À quoi rêvent...}, acte II, scène 3. (Note de Willems.)} ».\\\passeLigne}%
\end{groupe}%
%% v. 907
\begin{groupe}%
\personnageFr{Euripide}%
\trimIambFr{« Es-tu Hellène, ou bien es-tu de ce pays ? »\\\passeLigne}%
\end{groupe}%
%% v. 908
\begin{groupe}%
\personnageFr{Mnésiloque}%
\trimIambFr{« Hellène. Et à mon tour je veux savoir ta race. »\\\passeLigne}%
\end{groupe}%
%% v. 909
\begin{groupe}%
\personnageFr{Euripide}%
\trimIambFr{« Ô femme, je te vois tout pareille à Hélène. »\\\passeLigne}%
\end{groupe}%
%% v. 910
\begin{groupe}%
\personnageFr{Mnésiloque}%
\trimIambFr{« Toi-même à Ménélas, si j'en juge par ces hardes.%\commentaire{}{}{Le texte est visiblement corrompu. Porson, et bien d'autres après lui, lisent σέ τῷ 'κ ἰφύων, « le Ménélas du marché aux légumes », plaisanterie qui, même dans une parodie, paraît par trop intempestive surtout venant de Mnésiloque. Aussi Droysen la met-il dans la bouche de Critylla.}. 
»\\\passeLigne}%
\end{groupe}%
%% v. 911
\begin{groupe}%
\personnageFr{Euripide}%
\trimIambFr{« Oui, tu l'as reconnu, ce grand infortuné. »\\\passeLigne}%
\end{groupe}%
%% v. 912-916a
\begin{groupe}%
\personnageFr{Mnésiloque}%
\trimIambFr{« Ô toi, le tard venu dans les bas de ta femme »,}
\end{groupe}%
%% v. 912-916a
\begin{groupe}%
\diDochFr{prends, prends-moi, mon époux. Presse-moi dans tes bras.\\
Voyons, que je te baise. Emmène-moi, emmène, emmène, emmène-moi,}
\end{groupe}%
%% v. 912-916a
\begin{groupe}%
\trimIambFr{et au plus vite. (Il descend de l'autel, les femmes l'entourent.)\\\passeLigne}%
\end{groupe}%
%% v. 916b-917
\begin{groupe}%
\personnageFr{Kritylla}%
\trimIambFr{\ligne{916}\espace{2.8cm}Ah! il pleurera, par les deux déesses,\\
celui qui t'emmènera. Je le frapperai de ma torche.\\\passeLigne}%
\end{groupe}%
%% v. 918-919
\begin{groupe}%
\personnageFr{Euripide}%
\trimIambFr{Quoi ? mon épouse à moi, fille de Tyndare,\\
 tu voudrais m'empêcher de la conduire à Sparte ?\\\passeLigne}%
\end{groupe}%
%% v. 920-923
\begin{groupe}%
\personnageFr{Kritylla}%
\trimIambFr{Ah! toi aussi, que tu m'as l'air d'un fripon\\
et d'être son conseiller. Rien d'étonnant que depuis un temps\\
vous jouiez les Égyptiens. Mais celui-là en portera la peine.\\
Car voici venir le prytane et l'archer.\\\passeLigne}%
\end{groupe}%
%% v. 924
\begin{groupe}%
\personnageFr{Euripide}%
\trimIambFr{Voilà qui ne vaut rien. Il est temps d'esquiver.\\\passeLigne}%
\end{groupe}%
%% v. 925a
\begin{groupe}%
\personnageFr{Mnésiloque}%
\trimIambFr{Et moi, malheureux, que vais-je faire ?\\\passeLigne}%
\end{groupe}%
%% v. 925b-927
\begin{groupe}%
\personnageFr{Euripide}%
\trimIambFr{\ligne{925}\espace{4.2cm}Sois en paix.\\
Jamais je ne te trahirai, tant que j'aurai le souffle,\\
à moins que mes mille stratagèmes ne viennent à me faillir.\\\passeLigne}%
\end{groupe}%
%% v. 928
\begin{groupe}%
\personnageFr{Mnésiloque}%
\trimIambFr{\lednopb%
Voilà une amorce qui n'a pas pris.\\\passeLigne}%
\end{groupe}%
%% v. 929-934a
\begin{groupe}%
\personnageFr{Prytane}%
\trimIambFr{Est-ce là le fripon dont nous parlait Clisthène ?\\
Holà, pour quoi courbes-tu le front? Archer, rentre avec lui et attache-le\\
à l'ais. Puis expose-le ici,\\
et garde-le bien. Ne laisse approcher personne,\\
mais frappe de ton fouet\\
quiconque s'avancera.\\\passeLigne}%
\end{groupe}%
%% v. 934b-935
\begin{groupe}%
\personnageFr{Kritylla}%
\trimIambFr{\ligne{934}\espace{2.9cm}Bien, par Zeus, car tout à l'heure\\
un rapetasseur de voiles a failli me l'enlever.\\\passeLigne}%
\end{groupe}%
%% v. 936-938
\begin{groupe}%
\personnageFr{Mnésiloque}%
\trimIambFr{Ô prytane, au nom de cette droite dont tu aimes\\
à tendre le creux (À part :) quand on t'offre de l'argent...\\ 
Accorde-moi une petite faveur, à moi qui suis sur le point de mourir.\\\passeLigne}%
\end{groupe}%
%% v. 939a
\begin{groupe}%
\personnageFr{Prytane}%
\trimIambFr{Quelle faveur ?\\\passeLigne}%
\end{groupe}%
%% v. 939b-942
\begin{groupe}%
\personnageFr{Mnésiloque}%
\trimIambFr{\ligne{939}\espace{2.8cm}Puisque j'ai été mis à nu,\\
ordonne à l'archer de me mettre face à l'ais.\\
Je crains d'être la risée des crocotes et des mitres\\
avec les corbeaux que je vais régaler.\\\passeLigne}%
\end{groupe}%
%% v. 943-944
\begin{groupe}%
\personnageFr{Prytane}%
\trimIambFr{La boulé a décidé qu'on t'attacherait dans ce costume,\\
 pour témoigner aux personnes présentes que tu es un fripon.\\\passeLigne}%
\end{groupe}%
%% v. 945-946
\begin{groupe}%
\personnageFr{Mnésiloque}%
\trimIambFr{Aïe, aïe! aïe, aïe! Ô crocote, que de malheurs tu as causés !\\
 Il n'est plus pour moi le moindre espoir de salut.\\\passeLigne}%
\end{groupe}%
%% v. 947-952
\begin{groupe}%
\personnageFr{Coryphée}%
\tetramAnapFr{Allons, livrons-nous à nos danses, comme c'est ici la coutume des femmes,\\
 quand aux dates saintes nous célébrons les augustes mystères des deux déesses.}%
\end{groupe}%
\begin{groupe}%
\dimAnapFr{Pauson aussi les sanctifie. Il jeûne\\
et souvent prie avec nous que d'année\\
en année il puisse longtemps}
\end{groupe}%
\begin{groupe}%
\dimAnapFr{encore leur témoigner même dévotion.\\\passeLigne}%
\end{groupe}%
%% v. 953-958
\begin{groupe}%
\personnageFr{Chœur}%
\chantChoeurFr{Pars, élance-toi}
\end{groupe}%
%% v. 953-958
\begin{groupe}%
\chantChoeurFr{d'un pied léger, forme le cercle,}
\end{groupe}%
%% v. 953-958
\begin{groupe}%
\chantChoeurFr{les mains enlacées aux mains,}
\end{groupe}%
%% v. 953-958
\begin{groupe}%
\chantChoeurFr{marque le rythme de la danse.}
\end{groupe}%
%% v. 953-958
\begin{groupe}%
\chantChoeurFr{Avance de tes pieds alertes.}
\end{groupe}%
%% v. 953-958
\begin{groupe}%
\chantChoeurFr{Mais il faut que le chœur soit aux aguets}
\end{groupe}%
%% v. 953-958
\begin{groupe}%
\chantChoeurFr{et tourne la vue de toutes parts.\\\passeLigne}%
\end{groupe}%
%% v. 959-961 strophe a
\begin{groupe}%
\chantChoeurFr{Et ensemble célèbre et honore de la voix la gent des dieux olympiens,}
\end{groupe}%
%% v. 959-961 strophe a
\begin{groupe}%
\chantChoeurFr{d'une âme transportée par la danse.\\\passeLigne}%
\end{groupe}%
%% v. 962-965 antistrophe a
\begin{groupe}%
\chantChoeurFr{Que si tel croit\ligne{963} qu'étant femme je dois médire}
\end{groupe}%
%% v. 962-965 antistrophe a
\begin{groupe}%
\chantChoeurFr{des hommes pendant nos rites,\ligne{965} il s'abuse.\\\passeLigne}%
\end{groupe}%
%% v. 966-968 mesodos a
\begin{groupe}%
\chantChoeurFr{Mais d'abord il convient,}
\end{groupe}%
%% v. 966-968 mesodos a
\begin{groupe}%
\chantChoeurFr{c'est la règle,}
\end{groupe}%
%% v. 966-968 mesodos a
\begin{groupe}%
\chantChoeurFr{de procéder aux figures élégantes de la danse cyclique. (Danse.)\\\passeLigne}%
\end{groupe}%
%% v. 969-976 strophe b
\begin{groupe}%
\chantChoeurFr{\ligne{969}Pars du pied en célébrant le dieu à la belle lyre}
\end{groupe}%
%% v. 969-976 strophe b
\begin{groupe}%
\chantChoeurFr{et Artémis armée de l'arc, la chaste souveraine.}
\end{groupe}%
%% v. 969-976 strophe b
\begin{groupe}%
\chantChoeurFr{Salut, dieu tutélaire,}
\end{groupe}%
%% v. 969-976 strophe b
\begin{groupe}%
\chantChoeurFr{et donne-nous la victoire.}
\end{groupe}%
%% v. 969-976 strophe b
\begin{groupe}%
\chantChoeurFr{Puis célébrons, comme il sied,}
\end{groupe}%
%% v. 969-976 strophe b
\begin{groupe}%
\chantChoeurFr{Héra la conjugale,}
\end{groupe}%
%% v. 969-976 strophe b
\begin{groupe}%
\chantChoeurFr{qui se récrée}
\end{groupe}%
%% v. 969-976 strophe b
\begin{groupe}%
\chantChoeurFr{avec tous les chœurs et}
\end{groupe}%
%% v. 969-976 strophe b
\begin{groupe}%
\chantChoeurFr{garde les clefs de l'hymen.\\\passeLigne}%
\end{groupe}%
%% v. 977-984 antistrophe b
\begin{groupe}%
\chantChoeurFr{Et Hermès pastoral,}
\end{groupe}%
%% v. 977-984 antistrophe b
\begin{groupe}%
\chantChoeurFr{et Pan et les Nymphes chéries,}
\end{groupe}%
%% v. 977-984 antistrophe b
\begin{groupe}%
\chantChoeurFr{je les prie de sourire avec bienveillance}
\end{groupe}%
%% v. 977-984 antistrophe b
\begin{groupe}%
\chantChoeurFr{et de se plaire}
\end{groupe}%
%% v. 977-984 antistrophe b
\begin{groupe}%
\chantChoeurFr{à nos danses.}
\end{groupe}%
%% v. 977-984 antistrophe b
\begin{groupe}%
\chantChoeurFr{Allons, commence avec ardeur}
\end{groupe}%
%% v. 977-984 antistrophe b
\begin{groupe}%
\chantChoeurFr{la diple, joie de la danse.}
\end{groupe}%
%% v. 977-984 antistrophe b
\begin{groupe}%
\chantChoeurFr{\ligne{983}Ébattons-nous, femmes, suivant l'usage.}
\end{groupe}%
%% v. 977-984 antistrophe b
\begin{groupe}%
\chantChoeurFr{Aussi bien nous sommes à jeun.\\\passeLigne}%
\end{groupe}%
%% v. 985-989 mesodos b
\begin{groupe}%
\chantChoeurFr{Or sus! marque une autre mesure de ton pied cadencé.}
\end{groupe}%
%% v. 985-989 mesodos b
\begin{groupe}%
\chantChoeurFr{Chante à pleine voix.}
\end{groupe}%
%% v. 985-989 mesodos b
\begin{groupe}%
\chantChoeurFr{Et guide-nous toi-même,}
\end{groupe}%
%% v. 985-989 mesodos b
\begin{groupe}%
\chantChoeurFr{couronné de lierre, seigneur Bacchus.}
\end{groupe}%
%% v. 985-989 mesodos b
\begin{groupe}%
\chantChoeurFr{\ligne{988}Je te célébrerai}
\end{groupe}%
%% v. 985-989 mesodos b
\begin{groupe}%
\chantChoeurFr{en des chants bien rythmés.\\\passeLigne}%
\end{groupe}%
%% v. 990-994 strophe g
\begin{groupe}%
\chantChoeurFr{Évios! fils de Zeus}
\end{groupe}%
%% v. 990-994 strophe g
\begin{groupe}%
\chantChoeurFr{et de Sémélé, ô Bromios,}
\end{groupe}%
%% v. 990-994 strophe g
\begin{groupe}%
\chantChoeurFr{toi que les chœurs délectent,}
\end{groupe}%
%% v. 990-994 strophe g
\begin{groupe}%
\chantChoeurFr{et qui dans les montagnes, parmi les hymnes aimables des Nymphes,}
\end{groupe}%
%% v. 990-994 strophe g
\begin{groupe}%
\chantChoeurFr{Évios ! Évios ! Évoé !}
\end{groupe}%
%% v. 990-994 strophe g
\begin{groupe}%
\chantChoeurFr{\ligne{994}te livres à la danse.\\\passeLigne}%
\end{groupe}%
%% v. 995-1000 antistrophe g
\begin{groupe}%
\chantChoeurFr{Autour de toi}
\end{groupe}%
%% v. 995-1000 antistrophe g
\begin{groupe}%
\chantChoeurFr{l'écho du Cithéron résonne,}
\end{groupe}%
%% v. 995-1000 antistrophe g
\begin{groupe}%
\chantChoeurFr{et retentissent les montagnes ombreuses}
\end{groupe}%
%% v. 995-1000 antistrophe g
\begin{groupe}%
\chantChoeurFr{et les ravins rocheux,}
\end{groupe}%
%% v. 995-1000 antistrophe g
\begin{groupe}%
\chantChoeurFr{tandis qu'à tes côtés le lierre}
\end{groupe}%
%% v. 995-1000 antistrophe g
\begin{groupe}%
\chantChoeurFr{aux belles feuilles pousse sa vrille tout à l'entour.\lednopb%
\\\passeLigne}%
\end{groupe}%
%% v. 1001
\begin{groupe}%
\personnageFr{Archer Scythe}%
\trimIambFr{Ici maintenant, tu vas crier, en plein air.\\\passeLigne}%
\end{groupe}%
%% v. 1002a
\begin{groupe}%
\personnageFr{Mnésiloque}%
\trimIambFr{Archer, je t'en supplie...\\\passeLigne}%
\end{groupe}%
%% v. 1002b
\begin{groupe}%
\personnageFr{Archer}%
\trimIambFr{\ligne{1002}\espace{3.1cm}Ne me supplie pas.\\\passeLigne}%
\end{groupe}%
%% v. 1003a
\begin{groupe}%
\personnageFr{Mnésiloque}%
\trimIambFr{Desserre la cheville.\\\passeLigne}%
\end{groupe}%
%% v. 1003b
\begin{groupe}%
\personnageFr{Archer}%
\trimIambFr{\ligne{1003}\espace{2.9cm}C'est ce que je vais faire\\\passeLigne}%
\end{groupe}%
%% v. 1004
\begin{groupe}%
\personnageFr{Mnésiloque}%
\trimIambFr{Aïe, misère! tu la pousses davantage.\\\passeLigne}%
\end{groupe}%
%% v. 1005a
\begin{groupe}%
\personnageFr{Archer}%
\trimIambFr{Encore plus, veux-tu?\\\passeLigne}%
\end{groupe}%
%% v. 1005b-1006a
\begin{groupe}%
\personnageFr{Mnésiloque}%
\trimIambFr{\ligne{1005}\espace{2.7cm}Aie, aïe, aïe, aie, aïe, aie !\\
 la peste soit de toi.\\\passeLigne}%
\end{groupe}%
%% v. 1006b-1007
\begin{groupe}%
\personnageFr{Archer}%
\trimIambFr{\ligne{1006}\espace{2.5cm}Silence, misérable vieillard.\\
 Allons, je vais, moi, apporter une natte, pour faire ma garde. (Il sort.)\\\passeLigne}%
\end{groupe}%
%% v. 1008-1014
\begin{groupe}%
\personnageFr{Mnésiloque}%
\trimIambFr{Ces douceurs-là j'en ai l'obligation à Euripide.\\
(Il aperçoit dans la coulisse Euripide costumé en Persée.) Hein ! dieux ! Zeus sauveur ! il y a espoir.
L'homme ne semble pas vouloir m'abandonner. Mais il m'a\\
donné à entendre, en apparaissant en Persée,\\
qu'il me faut devenir Andromède. Aussi bien,\\
les chaînes, j'en suis muni. Il est donc clair\\
qu'il viendra me délivrer. Autrement il n'aurait pas dirigé son vol de ce côté.\\\passeLigne}%
\end{groupe}%
%% v. 1015-1055
\begin{groupe}%
\personnageFr{Mnésiloque}%
\paratragiquesFr{\ligne{1015}Vierges, vierges amies, puissé-je}
\end{groupe}%
%% v. 1015-1055
\begin{groupe}%
\paratragiquesFr{partir d'ici et}
\end{groupe}%
%% v. 1015-1055
\begin{groupe}%
\paratragiquesFr{tromper le Scythe !}
\end{groupe}%
%% v. 1015-1055
\begin{groupe}%
\paratragiquesFr{M'entends-tu, ô toi qui en des antres répètes mes chants ?}
\end{groupe}%
%% v. 1015-1055
\begin{groupe}%
\paratragiquesFr{acquiesce, souffre que}
\end{groupe}%
%% v. 1015-1055
\begin{groupe}%
\paratragiquesFr{j'aille retrouver ma femme.}
\end{groupe}%
%% v. 1015-1055
\begin{groupe}%
\paratragiquesFr{Sans merci est celui qui m'enchaîna,}
\end{groupe}%
%% v. 1015-1055
\begin{groupe}%
\paratragiquesFr{des mortels le plus éprouvé.}
\end{groupe}%
%% v. 1015-1055
\begin{groupe}%
\paratragiquesFr{Échappé à grand'peine à une vieille}
\end{groupe}%
%% v. 1015-1055
\begin{groupe}%
\paratragiquesFr{guenipe, je ne laissai pas d'être perdu.}
\end{groupe}%
%% v. 1015-1055
\begin{groupe}%
\paratragiquesFr{Car depuis un temps ce Scythe se tient en faction : il m'a suspendu, brisé, sans amis,}
\end{groupe}%
%% v. 1015-1055
\begin{groupe}%
\paratragiquesFr{en pâture aux corbeaux. Vois-tu, ce n'est pas}
\end{groupe}%
%% v. 1015-1055
\begin{groupe}%
\paratragiquesFr{pour des chœurs de danse}
\end{groupe}%
%% v. 1015-1055
\begin{groupe}%
\paratragiquesFr{ni parmi des jeunes filles de mon âge}
\end{groupe}%
%% v. 1015-1055
\begin{groupe}%
\paratragiquesFr{que je suis ici debout avec cornet et jeton.%\commentaire{1029}{κημὸν ἕστηκ᾽ ἔχουσ᾽ }{A comparer avec le bas-relief de la galerie Gustiniani et le vase d'argent du musée Corsini, représentant le jugement d'Oreste devant l'Aréopage, et où une jeune fille se tient en effet debout derrière l'urne (SAGLIO, I, p.399). Était-ce l'usage reçu à l'Aréopage?}.
}
\end{groupe}%
%% v. 1015-1055
\begin{groupe}%
\paratragiquesFr{Mais enchaînée en d'étroits liens,}
\end{groupe}%
%% v. 1015-1055
\begin{groupe}%
\paratragiquesFr{je me vois exposée à la voracité du monstre Glaucétès.%\commentaire{1031}{Γλαυκέτῃ}{Glaucétès l'omnivore, le glouton, figure ici le monstre marin qui doit dévorer Andromède.}
}
\end{groupe}%
%% v. 1015-1055
\begin{groupe}%
\paratragiquesFr{Lamentez sur moi un péan, non d'hymen, mais de chaînes,}
\end{groupe}%
%% v. 1015-1055
\begin{groupe}%
\paratragiquesFr{ô femmes, car,}
\end{groupe}%
%% v. 1015-1055
\begin{groupe}%
\paratragiquesFr{pitoyable, j'ai pitoyablement souffert.}
\end{groupe}%
%% v. 1015-1055
\begin{groupe}%
\paratragiquesFr{Malheureux que je suis, malheureux !}
\end{groupe}%
%% v. 1015-1055
\begin{groupe}%
\paratragiquesFr{Et entre d'autres tortures impies infligées par mes proches,}
\end{groupe}%
%% v. 1015-1055
\begin{groupe}%
\paratragiquesFr{implorant l'homme}
\end{groupe}%
%% v. 1015-1055
\begin{groupe}%
\paratragiquesFr{d'une lamentation pleine de larmes de l'obscur Hadès.}
\end{groupe}%
%% v. 1015-1055
\begin{groupe}%
\paratragiquesFr{ah ah! ah ah!}
\end{groupe}%
%% v. 1015-1055
\begin{groupe}%
\paratragiquesFr{l'homme qui me rasa d'abord,}
\end{groupe}%
%% v. 1015-1055
\begin{groupe}%
\paratragiquesFr{me revêtit de cette crocote,}
\end{groupe}%
%% v. 1015-1055
\begin{groupe}%
\paratragiquesFr{puis de surcroît me dépêcha vers ce}
\end{groupe}%
%% v. 1015-1055
\begin{groupe}%
\paratragiquesFr{temple, au milieu de femmes.}
\end{groupe}%
%% v. 1015-1055
\begin{groupe}%
\paratragiquesFr{Hélas!}
\end{groupe}%
%% v. 1015-1055
\begin{groupe}%
\paratragiquesFr{démon implacable de mon destin !}
\end{groupe}%
%% v. 1015-1055
\begin{groupe}%
\paratragiquesFr{Ô moi le maudit !}
\end{groupe}%
%% v. 1015-1055
\begin{groupe}%
\paratragiquesFr{Qui, en présence de mes misères, ne tournerait les yeux}
\end{groupe}%
%% v. 1015-1055
\begin{groupe}%
\paratragiquesFr{sur cette inenviable souffrance ?}
\end{groupe}%
%% v. 1015-1055
\begin{groupe}%
\paratragiquesFr{Puisse l'Étoile dévorante de l'éther}
\end{groupe}%
%% v. 1015-1055
\begin{groupe}%
\paratragiquesFr{anéantir un infortuné !%\commentaire{1050-1}{πυρφόρος αἰθέρος ἀστὴρ ...}{Cf. sur le sens de ce passage et le rôle de l'étoile Sirius : Bull. cité, note 26, ci-après.} !
}
\end{groupe}%
%% v. 1015-1055
\begin{groupe}%
\paratragiquesFr{Car il ne m'est plus doux de voir la lumière immortelle,}
\end{groupe}%
%% v. 1015-1055
\begin{groupe}%
\paratragiquesFr{suspendu que je suis,}
\end{groupe}%
%% v. 1015-1055
\begin{groupe}%
\paratragiquesFr{attendant comme terme d'un supplice jugulant%\commentaire{1054}{}{Allusion au carcan qui lui serre le cou. Lisez : δία μένων. Cf. sur ce passage : Bull. cité, note 27, ci-après.}
}
\end{groupe}%
%% v. 1015-1055
\begin{groupe}%
\paratragiquesFr{un trajet rapide vers les morts.%\commentaire{1015-55}{}{Tout ce chant est d'une incohérence voulue.}.
\\\passeLigne}%
\end{groupe}%
%% v. 1056-1057
\begin{groupe}%
\personnageFr{Euripide}%
\trimIambFr{\ligne{1056}(Dans la coulisse et sans être vu). Salut, chère enfant. Mais ton père Cepheus\\
qui t'exposa, que les dieux le confondent.\\\passeLigne}%
\end{groupe}%
%% v. 1058
\begin{groupe}%
\personnageFr{Mnésiloque}%
\trimIambFr{Qui es-tu, toi qui compatis à mon tourment ?\\\passeLigne}%
\end{groupe}%
%% v. 1059-1063a
\begin{groupe}%
\personnageFr{Euripide}%
\trimIambFr{Je suis Écho, la rediseuse goguenarde,\\
celle qui, l'an passé, en ce même lieu,\\
prêtai mon aide à Euripide.\\
Songe, mon enfant, à tenir bien ta partie\\
et à gémir lamentablement.\\\passeLigne}%
\end{groupe}%
%% v. 1063b
\begin{groupe}%
\personnageFr{Mnésiloque}%
\trimIambFr{\ligne{1063}\espace{2.4cm}Et toi, à gémir de plus belle après moi.\\\passeLigne}%
\end{groupe}%
%% v. 1064
\begin{groupe}%
\personnageFr{Euripide}%
\trimIambFr{Je n'y manquerai pas. Commence.\\\passeLigne}%
\end{groupe}%
%% v. 1065-1069a
\begin{groupe}%
\personnageFr{Mnésiloque}%
\anapFr{Ô Nuit sainte,\\
quelle longue carrière tu fournis,\\
en guidant ton char sur le dos constellé\\
de l'éther sacré,\\
à travers le très auguste Olympe.\\\passeLigne}%
\end{groupe}%
%% v. 1069b
\begin{groupe}%
\personnageFr{Euripide}%
\anapFr{\ligne{1069}\espace{4cm}Olympe.\\\passeLigne}%
\end{groupe}%
%% v. 1070-1071a
\begin{groupe}%
\personnageFr{Mnésiloque}%
\anapFr{Pourquoi seule Andromède reçut-elle tant de maux\\
en partage ?\\\passeLigne}%
\end{groupe}%
%% v. 1071b
\begin{groupe}%
\personnageFr{Euripide}%
\anapFr{\ligne{1071}\espace{2.5cm}En partage.\\\passeLigne}%
\end{groupe}%
%% v. 1072a
\begin{groupe}%
\personnageFr{Mnésiloque}%
\anapFr{Malheureuse que la mort...%\commentaire{1072}{θανάτου τλήμων...}{(Attend.)}
\\\passeLigne}%
\end{groupe}%
%% v. 1072b
\begin{groupe}%
\personnageFr{Euripide}%
\anapFr{\ligne{1072}\espace{2.7cm}Malheureuse que la mort.\\\passeLigne}%
\end{groupe}%
%% v. 1073
\begin{groupe}%
\personnageFr{Mnésiloque}%
\anapFr{Tu me tueras, vieille, avec ton babil.\\\passeLigne}%
\end{groupe}%
%% v. 1074
\begin{groupe}%
\personnageFr{Euripide}%
\anapFr{Ton babil.\\\passeLigne}%
\end{groupe}%
%% v. 1075-1076a
\begin{groupe}%
\personnageFr{Mnésiloque}%
\anapFr{Par Zeus, ton ingérence m'excède\\
 au dernier point.\\\passeLigne}%
\end{groupe}%
%% v. 1076b
\begin{groupe}%
\personnageFr{Euripide}%
\anapFr{\ligne{1076}\espace{1cm}Au dernier point.\\\passeLigne}%
\end{groupe}%
%% v. 1077-1078a
\begin{groupe}%
\personnageFr{Mnésiloque}%
\anapFr{Mon brave, laisse-moi chanter ma monodie.\\
Tu m'obligeras : cesse !\\\passeLigne}%
\end{groupe}%
%% v. 1078b
\begin{groupe}%
\personnageFr{Euripide}%
\anapFr{\ligne{1078}\espace{3.5cm}(Criard :) C'est ce... que j'allais dire.\\\passeLigne}%
\end{groupe}%
%% v. 1079a
\begin{groupe}%
\personnageFr{Mnésiloque}%
\anapFr{Va-t'en aux corbeaux.\\\passeLigne}%
\end{groupe}%
%% v. 1079b
\begin{groupe}%
\personnageFr{Euripide}%
\anapFr{\ligne{1079}\espace{2.9cm}Va-t'en aux corbeaux.\\\passeLigne}%
\end{groupe}%
%% v. 1080a
\begin{groupe}%
\personnageFr{Mnésiloque}%
\anapFr{Qu'est-ce qui t'arrive ?\\\passeLigne}%
\end{groupe}%
%% v. 1080b
\begin{groupe}%
\personnageFr{Euripide}%
\anapFr{\ligne{1080}\espace{1.5cm}Qu'est-ce qui t'arrive ?\\\passeLigne}%
\end{groupe}%
%% v. 1080c
\begin{groupe}%
\personnageFr{Mnésiloque}%
\anapFr{\ligne{1080}\espace{3cm}Tu veux rire.\\\passeLigne}%
\end{groupe}%
%% v. 1080d
\begin{groupe}%
\personnageFr{Euripide}%
\anapFr{\ligne{1080}\espace{4.2cm}Tu veux rire.\\\passeLigne}%
\end{groupe}%
%% v. 1081a
\begin{groupe}%
\personnageFr{Mnésiloque}%
\anapFr{Va te promener..\\\passeLigne}%
\end{groupe}%
%% v. 1081b
\begin{groupe}%
\personnageFr{Euripide}%
\anapFr{\ligne{1081}\espace{1.1cm}Va te promener..\\\passeLigne}%
\end{groupe}%
%% v. 1081c
\begin{groupe}%
\personnageFr{Mnésiloque}%
\anapFr{\ligne{1081}\espace{2.3cm}Foin de toi.\\\passeLigne}%
\end{groupe}%
%% v. 1081d
\begin{groupe}%
\personnageFr{Euripide}%
\anapFr{\ligne{1081}\espace{3.5cm}Foin de toi.\\\passeLigne}%
\end{groupe}%
%% v. 1083a
\begin{groupe}%
\personnageFr{Archer}%
\anapFr{\ligne{1083}Holà ! qu'as-tu à bavarder ?\\\passeLigne}%
\end{groupe}%
%% v. 1083b
\begin{groupe}%
\personnageFr{Euripide}%
\anapFr{\ligne{1083}\espace{2.6cm}Holà ! qu'as-tu à bavarder ?\\\passeLigne}%
\end{groupe}%
%% v. 1084a
\begin{groupe}%
\personnageFr{Archer}%
\anapFr{J'appellerai les prytanes.\\\passeLigne}%
\end{groupe}%
%% v. 1084b
\begin{groupe}%
\personnageFr{Euripide}%
\anapFr{\ligne{1084}\espace{2.7cm}J'appellerai les prytanes.\\\passeLigne}%
\end{groupe}%
%% v. 1085a
\begin{groupe}%
\personnageFr{Archer}%
\anapFr{Qu'est-ce qui arrive ?\\\passeLigne}%
\end{groupe}%
%% v. 1085b
\begin{groupe}%
\personnageFr{Euripide}%
\anapFr{\ligne{1085}\espace{1.6cm}Qu'est-ce qui arrive ?\\\passeLigne}%
\end{groupe}%
%% v. 1086a
\begin{groupe}%
\personnageFr{Archer}%
\anapFr{D'où part cette voix ?\\\passeLigne}%
\end{groupe}%
%% v. 1086b
\begin{groupe}%
\personnageFr{Euripide}%
\anapFr{\ligne{1086}\espace{2.4cm}D'où part cette voix ?\\\passeLigne}%
\end{groupe}%
%% v. 1087a
\begin{groupe}%
\personnageFr{Archer}%
\anapFr{Est-ce toi qui bavardes?\\\passeLigne}%
\end{groupe}%
%% v. 1087b
\begin{groupe}%
\personnageFr{Euripide}%
\anapFr{\ligne{1087}\espace{1.8cm}Est-ce toi qui bavardes?\\\passeLigne}%
\end{groupe}%
%% v. 1088a
\begin{groupe}%
\personnageFr{Archer}%
\anapFr{Gare à toi.\\\passeLigne}%
\end{groupe}%
%% v. 1088b
\begin{groupe}%
\personnageFr{Euripide}%
\anapFr{\ligne{1088}\espace{1.4cm}Gare à toi.\\\passeLigne}%
\end{groupe}%
%% v. 1089a
\begin{groupe}%
\personnageFr{Archer}%
\anapFr{Tu te gausses de moi.\\\passeLigne}%
\end{groupe}%
%% v. 1089b
\begin{groupe}%
\personnageFr{Euripide}%
\anapFr{\ligne{1089}\espace{2.2cm}Tu te gausses de moi.\\\passeLigne}%
\end{groupe}%
%% v. 1090
\begin{groupe}%
\personnageFr{Chœur}%
\anapFr{ Non, par Zeus : c'est cette femme, là tout près.\\\passeLigne}%
\end{groupe}%
%% v. 1091
\begin{groupe}%
\personnageFr{Euripide}%
\anapFr{Là tout près.\\\passeLigne}%
\end{groupe}%
%% v. 1092-1094a
\begin{groupe}%
\personnageFr{Archer}%
\anapFr{Où est la canaille ?\\\passeLigne}%
\end{groupe}%
%% v. 1092-1094a
\begin{groupe}%
\personnageFr{Chœur}%
\anapFr{\espace{2.2cm}La voilà qui se sauve.\\\passeLigne}%
\end{groupe}%
%% v. 1092-1094a
\begin{groupe}%
\personnageFr{Archer}%
\anapFr{Où, où te sauves-tu ?\\\passeLigne}%
\end{groupe}%
%% v. 1092-1094a
\begin{groupe}%
\personnageFr{Euripide}%
\anapFr{\ligne{1092}\decalage{Où, où te sauves-tu ?}Où, où voles-tu ?\\\passeLigne}%
\end{groupe}%
%% v. 1092-1094a
\begin{groupe}%
\personnageFr{Archer}%
\anapFr{Il t'en cuira.\\\passeLigne}%
\end{groupe}%
%% v. 1094b
\begin{groupe}%
\personnageFr{Euripide}%
\anapFr{\ligne{1094}\espace{2.1cm}Tu ne m'auras.\\\passeLigne}%
\end{groupe}%
%% v. 1095a
\begin{groupe}%
\personnageFr{Archer}%
\anapFr{Tu maronnes encore ?\\\passeLigne}%
\end{groupe}%
%% v. 1095b
\begin{groupe}%
\personnageFr{Euripide}%
\anapFr{\ligne{1095}\espace{2.5cm}Tu maronnes encore ?\\\passeLigne}%
\end{groupe}%
%% v. 1096a
\begin{groupe}%
\personnageFr{Archer}%
\anapFr{Arrête, la canaille.\\\passeLigne}%
\end{groupe}%
%% v. 1096b
\begin{groupe}%
\personnageFr{Euripide}%
\anapFr{\ligne{1096}\espace{2.3cm}Arrête, la canaille.\\\passeLigne}%
\end{groupe}%
%% v. 1097
\begin{groupe}%
\personnageFr{Archer}%
\anapFr{La bavarde, la maudite femme.\\\passeLigne}%
\end{groupe}%
%% v. 1098-1102a
\begin{groupe}%
\personnageFr{Euripide}%
\trimIambFr{(En Persée :) « En quel pays barbare, ô dieux, ces talonnières\\ 
ont-elles amené Persée ? Au milieu de l'éther\\
me frayant un chemin, prestement je fends l'air d'un pied ailé,\\
portant à la cité argienne la tête de\\
 Gorgone. »\\\passeLigne}%
\end{groupe}%
%% v. 1102b-1103a
\begin{groupe}%
\personnageFr{Archer}%
\trimIambFr{\ligne{1102}\espace{2.2cm}Que dis-tu? Tu portes la tête de Gorgon\\
le scribe ?\\\passeLigne}%
\end{groupe}%
%% v. 1103b-1104a
\begin{groupe}%
\personnageFr{Euripide}%
\trimIambFr{\ligne{1103}\espace{4.4cm}De Gorgone, ai-je dit.\\
ἔγωγε φημί.\\\passeLigne}%
\end{groupe}%
%% v. 1104b
\begin{groupe}%
\personnageFr{Archer}%
\trimIambFr{\ligne{1104}\espace{2cm}Gorgon, parfaitement, je dis de même.\\\passeLigne}%
\end{groupe}%
%% v. 1105-1106
\begin{groupe}%
\personnageFr{Euripide}%
\trimIambFr{« Ciel! Quelle est cette vierge aux déesses pareille amarrée à ce roc,\\
comme un navire au port ? »\\\passeLigne}%
\end{groupe}%
%% v. 1107-1108a
\begin{groupe}%
\personnageFr{Mnésiloque}%
\trimIambFr{« Étranger, prends pitié de moi, la désolée.\\
Détache mes liens. »\\\passeLigne}%
\end{groupe}%
%% v. 1108b-1109
\begin{groupe}%
\personnageFr{Archer}%
\trimIambFr{\ligne{1108}\espace{2.8cm}Veux-tu ne pas bavarder ?\\
Maudit, oses-tu bien ? Tu vas mourir et tu bavardes.\\\passeLigne}%
\end{groupe}%
%% v. 1110
\begin{groupe}%
\personnageFr{Euripide}%
\trimIambFr{« Ô vierge, une pitié me prend quand je te vois suspendue à ce roc. »\\\passeLigne}%
\end{groupe}%
%% v. 1111-1112a
\begin{groupe}%
\personnageFr{Archer}%
\trimIambFr{Ce n'est pas une vierge, c'est un vieux coquin,\\
et un voleur, et un fripon.\\\passeLigne}%
\end{groupe}%
%% v. 1112b-1113
\begin{groupe}%
\personnageFr{Euripide}%
\trimIambFr{\ligne{1112}\espace{3.9cm}Tu plaisantes, ô Scythe.\\
C'est Andromède, fille de Cépheus.\\\passeLigne}%
\end{groupe}%
%% v. 1114
\begin{groupe}%
\personnageFr{Archer}%
\trimIambFr{Regarde son cas : te paraît-il petit ?\\\passeLigne}%
\end{groupe}%
%% v. 1115-1118a
\begin{groupe}%
\personnageFr{Euripide}%
\trimIambFr{Donne-moi la main, que j'atteigne à la jeune fille.\\
Voyons, Scythe « chaque homme a ses faiblesses » :\\
la mienne est d'être féru\\
d'amour pour cette jeune fille.\\\passeLigne}%
\end{groupe}%
%% v. 1118b-1120
\begin{groupe}%
\personnageFr{Archer}%
\trimIambFr{\ligne{1118}\espace{3.3cm}Je ne suis pas jaloux.\\
S'il tournait les fesses par ici,\\
 je ne te refuserais pas d'en passer ta fantaisie.\\\passeLigne}%
\end{groupe}%
%% v. 1121-1122
\begin{groupe}%
\personnageFr{Euripide}%
\trimIambFr{Que ne me permets-tu de la délier, ô Scythe\\
« et de me laisser choir dans le lit nuptial »\\\passeLigne}%
\end{groupe}%
%% v. 1123-1124
\begin{groupe}%
\personnageFr{Archer}%
\trimIambFr{Si tu as tant envie de jouir de ce vieillard,\\
perce-moi l'ais par derrière et enfourne.\\\passeLigne}%
\end{groupe}%
%% v. 1125a
\begin{groupe}%
\personnageFr{Euripide}%
\trimIambFr{Non, par Zeus, je briserai ses liens.\\\passeLigne}%
\end{groupe}%
%% v. 1125b
\begin{groupe}%
\personnageFr{Archer}%
\trimIambFr{\ligne{1125}\espace{3.7cm}Alors je te fouette.\\\passeLigne}%
\end{groupe}%
%% v. 1126a
\begin{groupe}%
\personnageFr{Euripide}%
\trimIambFr{N'importe, je vais le faire.\\\passeLigne}%
\end{groupe}%
%% v. 1126b-1127
\begin{groupe}%
\personnageFr{Archer}%
\trimIambFr{\ligne{1126}\espace{3.2cm}Alors je te coupe la tête\\
avec ce sabre.\\\passeLigne}%
\end{groupe}%
%% v. 1128-1132
\begin{groupe}%
\personnageFr{Euripide}%
\trimIambFr{Aïe, aïe! que faire ? à quelles raisons recourir ?\\
Mais cette nature barbare ne les admettrait pas.\\
« Proposer à des sots des doctrines nouvelles,\\%\commentaire{1130-1}{}{Vers qu'Euripide met dans la bouche de Médée, dans la tragédie de ce nom, v. 298.},\\
c'est perdre sa peine. » Eh bien, dressons quelqu'autre\\
machine, qui soit de sa portée. (Il s'en va.)\\\passeLigne}%
\end{groupe}%
%% v. 1133
\begin{groupe}%
\personnageFr{Archer}%
\trimIambFr{La canaille de renard, comme il pensait m'embabouiner !\\\passeLigne}%
\end{groupe}%
%% v. 1134
\begin{groupe}%
\personnageFr{Mnésiloque}%
\trimIambFr{Souviens-toi, Persée, que tu me laisses en vilaine posture.\\\passeLigne}%
\end{groupe}%
%% v. 1135
\begin{groupe}%
\personnageFr{Archer}%
\trimIambFr{As-tu encore envie d'attraper le fouet? (Il se recouche et s'endort.)\\\passeLigne}%
\end{groupe}%
%% v. 1136-1159
\begin{groupe}%
\personnageFr{Chœur}%
\chantChoeurFr{C'est ma coutume d'appeler parmi nous Pallas,}
\end{groupe}%
%% v. 1136-1159
\begin{groupe}%
\chantChoeurFr{amie du chœur, la jeune fille,}
\end{groupe}%
%% v. 1136-1159
\begin{groupe}%
\chantChoeurFr{\ligne{1138}la vierge libre du joug,\\\passeLigne}
\end{groupe}%
%% v. 1136-1159
\begin{groupe}%
\chantChoeurFr{\ligne{1140}gardienne de notre acropole,}
\end{groupe}%
%% v. 1136-1159
\begin{groupe}%
\chantChoeurFr{seule en possession d'un visible pouvoir,}
\end{groupe}%
%% v. 1136-1159
\begin{groupe}%
\chantChoeurFr{et qui a nom Clédouque.\\\passeLigne}
\end{groupe}%
%% v. 1136-1159
\begin{groupe}%
\chantChoeurFr{Apparais, toi qui détestes les tyrans, comme faire se doit.\\\passeLigne}
\end{groupe}%
%% v. 1136-1159
\begin{groupe}%
\chantChoeurFr{Le peuple des femmes t'appelle.}
\end{groupe}%
%% v. 1136-1159
\begin{groupe}%
\chantChoeurFr{Viens à moi avec}
\end{groupe}%
%% v. 1136-1159
\begin{groupe}%
\chantChoeurFr{la Paix amie des fêtes.\\\passeLigne}
\end{groupe}%
%% v. 1136-1159
\begin{groupe}%
\chantChoeurFr{\ligne{1148}Venez, augustes déesses, bienveillantes et propices,}
\end{groupe}%
%% v. 1136-1159
\begin{groupe}%
\chantChoeurFr{dans votre enceinte sacrée,}
\end{groupe}%
%% v. 1136-1159
\begin{groupe}%
\chantChoeurFr{où vous ne permettez pas aux hommes de contempler}
\end{groupe}%
%% v. 1136-1159
\begin{groupe}%
\chantChoeurFr{vos saintes orgies à la lueur des torches.\\\passeLigne}
\end{groupe}%
%% v. 1136-1159
\begin{groupe}%
\chantChoeurFr{Découvrez vos faces immortelles.\\\passeLigne}
\end{groupe}%
%% v. 1136-1159
\begin{groupe}%
\chantChoeurFr{Venez, approchez, nous vous en conjurons,}
\end{groupe}%
%% v. 1136-1159
\begin{groupe}%
\chantChoeurFr{Thesmophores très vénérées.}
\end{groupe}%
%% v. 1136-1159
\begin{groupe}%
\chantChoeurFr{Si jamais auparavant vous répondîtes à notre appel,}
\end{groupe}%
%% v. 1136-1159
\begin{groupe}%
\chantChoeurFr{\ligne{1159}venez aujourd'hui, sur nos instantes prières, ici au milieu de nous.\\\passeLigne}
\end{groupe}%
%% v. 1160-1163
\begin{groupe}%
\personnageFr{Euripide}%
\trimIambFr{(Déguisé en vieille procureuse et portant une pectis ou harpe lydienne. Il est accompagné d'une danseuse et d'un jeune flûtiste). Femmes, si vous voulez pour le reste du temps
faire la paix avec moi, c'est le moment,\\
à charge que jamais plus à l'avenir je ne dirai de vous\\
le moindre mal. Telles sont mes propositions.\\\passeLigne}%
\end{groupe}%
%% v. 1164
\begin{groupe}%
\personnageFr{Coryphée}%
\trimIambFr{\ligne{1164}Qu'est-ce qui te décide à faire ces ouvertures ?\\\passeLigne}%
\end{groupe}%
%% v. 1165-1169
\begin{groupe}%
\personnageFr{Euripide}%
\trimIambFr{Cet homme attaché à l'ais est mon beau-père.\\
Que je le recouvre, et jamais\\
je ne médirai de vous. Si vous refusez,\\
toutes vos pratiques à huis-clos, je les révélerai\\
à vos maris, à leur retour de l'armée.\\\passeLigne}%
\end{groupe}%
%% v. 1170-1171
\begin{groupe}%
\personnageFr{Coryphée}%
\trimIambFr{Pour ce qui est de nous, sache que nous acquiesçons ; \\
mais ce barbare, c'est à toi de le convaincre.\\\passeLigne}%
\end{groupe}%
%% v. 1172-1175
\begin{groupe}%
\personnageFr{Euripide}%
\trimIambFr{Ça, c'est mon affaire. (À la danseuse ) La tienne, Élaphion,\\
est de songer à t'y prendre comme je t'ai dit chemin faisant.\\
Ainsi donc, passe d'abord à l'autre bout, et reviens en courant sur tes pointes.\\% \commentaire{1074}{ἀνακόλπασων}{Hésych. καλπαζεῖ ὀξυποδεῖ, ou, comme aurait dit Longus (II, 37) : ἐπ' ἄκρων τῶν ὀνύχων τρέχει. G. Hermann, à l'aide du scholiaste, a restitué la vraie leçon ἀνακάλπασων (au lieu de ἀνακόλπασων, qui n'est pas grec).}.\\
(Au flûtiste :) Toi, Térédon, joue-nous sur la flûte une Persique.\\
(Térédon joue de la flûte, pendant que la jeune fille danse.)
\\\passeLigne}%
\end{groupe}%
%% v. 1176
\begin{groupe}%
\personnageFr{Archer}%
\trimIambFr{(Se réveillant.) Qu'est-ce que ce ronflement? c'est un comos qui me réveille ?\\\passeLigne}%
\end{groupe}%
%% v. 1177-1178
\begin{groupe}%
\personnageFr{Euripide}%
\trimIambFr{Cette enfant, archer, pensait s'essayer.\\
Elle s'en va danser chez des particuliers.\\\passeLigne}%
\end{groupe}%
%% v. 1179-1180
\begin{groupe}%
\personnageFr{Archer}%
\trimIambFr{Qu'elle danse, qu'elle s'essaie, je ne m'y oppose pas.\\
Qu'elle est vive! On dirait une puce sur une peau de bique.\\\passeLigne}%
\end{groupe}%
%% v. 1181-1183a
\begin{groupe}%
\personnageFr{Euripide}%
\trimIambFr{Tends-moi ce vêtement, ma fille, en le tirant par en haut.\\
Assieds-toi sur les genoux du Scythe,\\
et avance les pieds, que je les déchausse.\\\passeLigne}%
\end{groupe}%
%% v. 1183b-1185
\begin{groupe}%
\personnageFr{Archer}%
\trimIambFr{\ligne{1183}\espace{4.5cm}Si fait, si,
mets-toi, si fait, ma fillette.\\
\lednopb%
(Pendant qu'Euripide ôte les chaussures à la danseuse.) Ouais! est-il ferme, ce téton ! une vraie rave. (La jeune fille déchaussée et nue danse au son de la flûte.)
\lednopb%
\\\passeLigne}%
\end{groupe}%
%% v. 1186
\begin{groupe}%
\personnageFr{Euripide}%
\trimIambFr{Joue donc plus vite. As-tu encore peur du Scythe ?\\\passeLigne}%
\end{groupe}%
%% v. 1187-1189
\begin{groupe}%
\personnageFr{Archer}%
\trimIambFr{Ah! que voilà de belles fesses %\commentaire{1187}{}{L'archer suit d'un œil émérillonné tous les mouvements de la danseuse et se penche tantôt à droite, tantôt à gauche (ἀνακύπτει καὶ παρακύπτει ἀπεψωλημένος\eff{...autre sens... cf . v. scène de déshabillage + citer Willems}]). Son agitation se trahit d'une manière par trop visible. Le Scythe, qui s'en aperçoit, reprend « gare à toi, etc .. »}
! — Gare à toi si tu ne restes dans ton coin.\\
(Il s'arrange.) — À la bonne heure ! Elle a repris de la tenue, la broquette.\\\passeLigne}%
\end{groupe}%
%% v. 1189-1190a
\begin{groupe}%
\personnageFr{Euripide}%
\trimIambFr{(À la danseuse). C'est assez. Reprends ton vêtement. Il est temps que\\
nous partions.\\\passeLigne}%
\end{groupe}%
%% v. 1190b
\begin{groupe}%
\personnageFr{Archer}%
\trimIambFr{\ligne{1190}\espace{2.1cm}Sans me donner un baiser?\\\passeLigne}%
\end{groupe}%
%% v. 1191a
\begin{groupe}%
\personnageFr{Euripide}%
\trimIambFr{Très volontiers. (À la danseuse.) Baise-le.\\\passeLigne}%
\end{groupe}%
%% v. 1191b-1193a
\begin{groupe}%
\personnageFr{Archer}%
\trimIambFr{\ligne{1191}\espace{3.9cm}Oh! oh! oh! papapapai!\\
Que sa langue est douce. C'est comme du miel attique.\\
Pourquoi ne couche-t-elle pas avec moi ?\\\passeLigne}%
\end{groupe}%
%% v. 1193b-1194a
\begin{groupe}%
\personnageFr{Euripide}%
\trimIambFr{\ligne{1193}\espace{3.7cm}Adieu, archer,\\
\lednopb%
 cela ne se peut.\\\passeLigne}%
\end{groupe}%
%% v. 1194b-1195a
\begin{groupe}%
\personnageFr{Archer}%
\trimIambFr{\ligne{1194}\espace{3.7cm}Si, si, petite vieille,\\
fais-moi ce plaisir.\\\passeLigne}%
\end{groupe}%
%% v. 1195b
\begin{groupe}%
\personnageFr{Euripide}%
\trimIambFr{\espace{3.3cm}Me donneras-tu une drachme ?\\\passeLigne}%
\end{groupe}%
%% v. 1196a
\begin{groupe}%
\personnageFr{Archer}%
\trimIambFr{Si fait, si, je te la donnerai.\\\passeLigne}%
\end{groupe}%
%% v. 1196b
\begin{groupe}%
\personnageFr{Euripide}%
\trimIambFr{\ligne{1196}\espace{2.4cm}Voyons ton argent.\\\passeLigne}%
\end{groupe}%
%% v. 1197
\begin{groupe}%
\personnageFr{Archer}%
\trimIambFr{Je n'en ai pas. Mais prends ma lance. (Il la lui confie.)\\
Je te l'apporte sur-le-champ. Suis-moi, mon enfant.\\
(À Euripide :) Toi, petite vieille, garde ce bonhomme.\\
Mais comment te nommes-tu ?\\\passeLigne}%
\end{groupe}%
%% v. 1200b
\begin{groupe}%
\personnageFr{Euripide}%
\trimIambFr{\espace{3.5cm}Artémisia.\\\passeLigne}%
\end{groupe}%
%% v. 1201
\begin{groupe}%
\personnageFr{Archer}%
\trimIambFr{Je le retiendrai, ce nom : Artamouxia. (Il s'en va avec la danseuse.)\\\passeLigne}%
\end{groupe}%
%% v. 1202-1206
\begin{groupe}%
\personnageFr{Euripide}%
\trimIambFr{Hermès, dieu de la ruse, jusqu'ici tu régis tout à souhait.\\
Toi, mon petit, enfuis-toi, avec ces instruments. (Il lui donne ses pectis.)\\
Pendant ce temps, je le détacherai.\\
(À Mnésiloque :) Quant à toi, dès que tu seras libre, vois à fuir résolument\\
et cours retrouver chez toi, ta femme et tes enfants.\\\passeLigne}%
\end{groupe}%
%% v. 1207
\begin{groupe}%
\personnageFr{Mnésiloque}%
\trimIambFr{J'en fais mon affaire, une fois délivré.\\\passeLigne}%
\end{groupe}%
%% v. 1208-1209a
\begin{groupe}%
\personnageFr{Euripide}%
\trimIambFr{Sois libre. Le reste te regarde : fuis avant que l'archer\\
ne revienne et ne t'empoigne.\\\passeLigne}%
\end{groupe}%
%% v. 1209b
\begin{groupe}%
\personnageFr{Mnésiloque}%
\trimIambFr{\espace{2.9cm}C'est ce que je fais. (Euripide et Mnésiloque s'enfuient.)% chacun de leur côté.)
\\\passeLigne}%
\end{groupe}%
%% v.1210-1216
\begin{groupe}%
\personnageFr{Archer}%
\trimIambFr{(Revenant en scène). Ah! petite vieille, la charmante fillette que tu as là,\\
et pas méchante... gentille! Où est la vieille ?\\
Ah ! je suis perdu! Où s'est éclipsé le vieillard ?\\
Ah! petite vieille, petite vieille ! Ce n'est pas bien, petite vieille.\\
(Appelant.) Artamouxia !\\
Elle m'a trompé, la vieille. (La danseuse s'esquive.) Va-t'en au plus vite, toi.\\
(Reprenant sa lance :) La lance est droite, je descends.\\
Ah ! que vais-je faire ? Où est la petite vieille ?\\
Artamouxia !\\\passeLigne}%
\end{groupe}%
%% v. 1217
\begin{groupe}%
\personnageFr{Coryphée}%
\trimIambFr{Tu demandes après la vieille qui portait les pectis ?\\\passeLigne}%
\end{groupe}%
%% v. 1218a
\begin{groupe}%
\personnageFr{Archer}%
\trimIambFr{Si fait, si. L'as-tu vue ?\\\passeLigne}%
\end{groupe}%
%% v. 1218b-1219
\begin{groupe}%
\personnageFr{Coryphée}%
\trimIambFr{\espace{3.54cm}C'est par là qu'elle a disparu,\\
elle et un vieux qui la suivait.\\\passeLigne}%
\end{groupe}%
%% v. 1220a
\begin{groupe}%
\personnageFr{Archer}%
\trimIambFr{Il avait une crocote, le vieux ?\\\passeLigne}%
\end{groupe}%
%% v. 1220b-1221
\begin{groupe}%
\personnageFr{Coryphée}%
\trimIambFr{\espace{3.9cm}Tout juste.\\
Tu peux encore les rattraper en tirant de ce côté. (Indiquant une autre direction.)\\\passeLigne}%
\end{groupe}%
%% v. 1222a-1222b
\begin{groupe}%
\personnageFr{Archer}%
\trimIambFr{Ô la vieille canaille! Par quel chemin prendre ?\\
\ligne{1222}Artamouxia ! (ll court vers l'entrée de l'orchestre.)\\\passeLigne}%
\end{groupe}%
%% v. 1223-1224
\begin{groupe}%
\personnageFr{Coryphée}%
\trimIambFr{Monte tout droit. Où cours-tu ?\\
Veux-tu bien revenir ici ? Tu cours à contresens.\\\passeLigne}%
\end{groupe}%
%% v. 1225
\begin{groupe}%
\personnageFr{Archer}%
\trimIambFr{Malheureux !... Mais je me sauve. Artamouxia !\\\passeLigne}%
\end{groupe}%
%% v. 1226-1231
\begin{groupe}%
\personnageFr{Coryphée}%
\anapFr{Cours maintenant, cours aux corbeaux, et qu'un bon vent te seconde.\\
Mais nous nous sommes assez ébattues.\\
Il est temps pour chacune de rentrer au logis.\\
Puissent, en retour, les Thesmophores\\
nous accorder\\
leur précieuse faveur !\\\passeLigne}%
\end{groupe}%
\endnumbering
\end{Rightside}%
\end{pages}%
\Pages
\cleardoublepage
\end{document} 